\chapter{Compatible fibrations and collapse assembly}
\label{ch:compatible-fibrations}

\paragraph{Scope.} The compatibility and gluing portion of the selected local-collapse proof.

\paragraph{Chapter structure.}
\begin{enumerate}[leftmargin=*]
\item Adapted coordinates.
\item Submersions.
\item Overlap adjustment.
\item Proper bundle output.
\item Corners.
\item Graph presentation.
\end{enumerate}

\paragraph{Current status.} I05 now has 46 implementation contracts,
complete blueprint coverage and a full chapter audit. Written conditional
proofs and source-qualified producers are distinguished throughout.
The source-supplied and source-omitted proof obligations, inherited
release bindings and eventual Lean implementation remain explicit.
Sections~\ref{sec:fibration-alignment-completion}--\ref{sec:fibration-complete-audit}
give the completed coverage and current disposition. Revision121 adds
Section~\ref{sec:fibration-cloud-local-repair}: five subordinate cloud
arguments, including a source counterexample and a smaller-tube repair.
The historical 46-row audit is retained. Revision122 adds CFS06--CFS10
and refutes the general FC28 claim. Revision123 adds CFS11--CFS15:
a global smoothing replacement with full coverage and projection under
stronger marker-scale hypotheses, plus its parameter-order modulus.
Actual cloud/consumer bindings and global assembly remain open.
Revision124 adds CFS16--CFS20 for original-cloud tubes, perturbed-input
estimates, closed supports and explicit three-stage tolerance budgets.
Uniform cutoff/plateau producers and final-base bindings remain open.
Revision125 adds CFS21--CFS25: buffered uniform cutoff calculus under
a new zero-marker bound, a counterexample to omitting that bound, and
an exact spectral locality criterion. Actual marker/cloud binding remains open.

\section{Compatibility proof worksheet}
For each chosen construction, record its domain and collar data,
regularity, quantitative error bounds, overlap conditions and the
structure retained after modification. Separate the existence of local
maps from compatibility on overlaps, and compatibility from the final
bundle or graph-presentation output. The following sections now specify these tasks from the selected
proof passages; qualified producer contracts remain unimplemented.

\section{Source and handoff boundaries}
I05 and L36 record the initial reading of Kleiner--Lott's Proposition14.1;
the author correction adds the $D^3$ case to part (2). L145--L155
now record the full supplied arguments of Sections12--16 and19--21,
including their explicitly omitted proofs. Generic smooth calculus,
submersion and bundle facts remain inherited reuse candidates. A raw graph
presentation belongs to the local-collapse output; essential torus
refinement is treated in Chapter~\ref{ch:seifert-graph}, with ambient
transport in Chapter~\ref{ch:ambient-reconstruction}.


\section{Receiving the completed Chapter 13 blueprint}
\label{sec:fibration-local-inputs}
Use LC87 as the local input list, with LC80's zero pieces, LC83's
circle charts, LC84's edge disks, LC85's slim fibers and LC88's
boundary augmentation. Their source-qualified producer obligations
remain explicit; a packet definition does not prove existence.
The next development is the actual global map and its adjustments in
KL Sections 12--15, including the auxiliary cutoffs from Section 9,
quantitative overlap compatibility, the common parameter order and
proper bundle restrictions after adjustment. Preserve the corrected
$D^3$ case in Proposition 14.1(2), the tested-domain factor bounds and
LC82's compact isolation. The resulting static graph theorem feeds
LC90's application stage; it is not an input to the local producers.


\section{The global coordinate map and its retained information}
\label{sec:fibration-global-map}
The first new targets of this chapter are FC01--FC07. They expand
KL Sections 12.1--12.3; the beginning of Section 13 identifies their
consumer. All use the original metric and the SAME local witnesses
exported by LC87. In this section the carrier is closed; LC88's
boundary augmentation still requires the separate Section 16 assembly.
No adjusted fibration or graph-manifold conclusion is inferred from
the construction of a map.

\begin{definition}[Finite block map; FC01]
\label{def:fibration-block-map}
Let $M$ be a smooth Riemannian manifold with smooth positive scale
$\rho$. Take finite, disjointly tagged index sets $I_0,I_s,I_e,I_2$
for the supplied zero, slim, edge and two-stratum families. Add the
two separate tags $\rho$ and $E'$. Put
\[
 H=\R_\rho\oplus H_{E'}\oplus H_0\oplus H_s\oplus H_e\oplus H_2,
 \qquad H_a=\bigoplus_{i\in I_a}H_i,
\]
where $H_i=\R^{k_i}\oplus\R$, with $k_i=2$ for $i\in I_2$
and $k_i=1$ otherwise, including $E'$. These are orthogonal direct
sums with the Euclidean $\ell^2$ norm. For $i\in I_0$ let
$R_i=r_i^0$; for $i\in I_s\cup I_e\cup I_2$ let $R_i=\rho(p_i)$.
Only $R_{E'}(p)=\rho(p)$ varies with $p$.

For each tag other than $\rho$, retain an open domain $U_i$, a smooth
coordinate $\eta_i:U_i\to\R^{k_i}$, and a global smooth cutoff
$\zeta_i:M\to[0,1]$ whose CLOSED support is contained in $U_i$.
This last condition is supplied by the buffered local cutoff
construction. Define, by zero extension outside $U_i$,
\[
 \mathcal E^0_\rho(p)=\rho(p),\qquad
 \mathcal E^0_i(p)=\bigl(R_i(p)\eta_i(p)\zeta_i(p),
                             R_i(p)\zeta_i(p)\bigr).
\]
The scalar last component of each block is its marker. Retain
\[
 Q_1=H,\quad Q_2=H_0\oplus H_s\oplus H_e,\quad
 Q_3=H_0\oplus H_s,\quad Q_4=H_0
\]
and the orthogonal projections $\pi_{a,b}:Q_a\to Q_b$ for
$a\le b$. In particular $\pi_{b,c}\pi_{a,b}=\pi_{a,c}$.
\end{definition}

\paragraph{Construction check.}
At a point outside $U_i$, the closed-support hypothesis gives a
neighborhood on which $\zeta_i$ vanishes. Hence the zero extension
is smooth there as well as on $U_i$; no undefined product is evaluated.
Finiteness gives a smooth map $\mathcal E^0:M\to H$. There is no
uniform bound on $\dim H$ as the carrier varies. A cutoff merely
vanishing on the boundary, without smooth zero extension or the
stated support buffer, would not suffice.

Writing a block as $(x'_i,x''_i)$, its image satisfies
$0\le x''_i\le R_i$ and $|x'_i|\le C_i x''_i$ whenever
$|\eta_i|\le C_i$ on its support. For the supplied cutoffs one may use
\[
 C_0=9/10,\quad C_s=9\cdot10^5\Delta,\quad
 C_e=C_{E'}=9\Delta,\quad C_2=9.
\]
For $E'$ the upper bound is $x''_{E'}\le x_\rho$, not a constant
radius. KL (12.4) prints $10^5\Delta$ for the slim block, whereas
(10.2) has support up to $9\cdot10^5\Delta$. We retain the latter,
justified bound; no cutoff is silently changed. This is a project
source comparison, not an entry on the author correction sheet.

\begin{lemma}[Derivative bound from active blocks; FC02]
\label{lem:fibration-active-derivative}
All estimate constants are nonnegative. Suppose $|d\rho|\le\Lambda$.
At every point at most $N$ of the
constant-radius block supports are present. On their supports assume
\[
 |\eta_i|\le C_i,\quad R_i|d\eta_i|\le a_i,\quad
 R_i|d\zeta_i|\le b_i,\quad a_i+(C_i+1)b_i\le B.
\]
For the single $E'$ block suppose instead
$|\eta_{E'}|\le C_E$, $\rho|d\eta_{E'}|\le a_E$ and
$\rho|d\zeta_{E'}|\le b_E$ on its support. Then, putting
$B_E=a_E+(C_E+1)(b_E+\Lambda)$, one has
\[
 \|D\mathcal E^0\|\le
       \bigl(\Lambda^2+NB^2+B_E^2\bigr)^{1/2}.
\]
All norms on the source use the original metric.
\end{lemma}
\begin{proof}
For a constant-radius block differentiate both factors of
$R_i\eta_i\zeta_i$. Its vector derivative has norm at most
$a_i+C_i b_i$ and its marker derivative at most $b_i$; their
Euclidean combined norm is at most their sum. For the variable
block the additional terms are
$(\eta_{E'}\zeta_{E'}\,d\rho,\zeta_{E'}\,d\rho)$, bounded by
$(C_E+1)\Lambda$. Blocks outside their closed supports have zero
derivative. For each unit source vector, square and sum the block
bounds in the orthogonal target, then take the supremum.
\end{proof}

The bound is independent of total chart count. To obtain KL 12.5's
constant depending only on its multiplicity parameter, one must also
verify the local cutoff estimates with $C_i b_i$ bounded uniformly
as $\Delta$ varies, and choose $\Lambda$ with the required
$\Delta$ dependence. FC02 proves the assembly estimate, not that
extra producer. Pointwise multiplicity is sufficient here. Bounding
all supports meeting an entire ball, as needed below, is a stronger
packing assertion and must be supplied separately from LC86.

\section{Markers, localization and the image scale}
\label{sec:fibration-marker-localization}
\begin{lemma}[A marker recovers its own coordinate; FC03]
\label{lem:fibration-marker-recovery}
Use a constant-radius block of FC01. Suppose $p\in U_i$,
$\zeta_i(p)=1$, $|\eta_i(p)|\le A$, and $0<e<1$. If $q\in M$
and $|\mathcal E^0(q)-\mathcal E^0(p)|<eR_i$, then
\[
 \zeta_i(q)>1-e>0,\qquad q\in U_i,\qquad
 |\eta_i(q)-\eta_i(p)|<\frac{(1+A)e}{1-e}.
\]
\end{lemma}
\begin{proof}
Orthogonal projection cannot increase distance. The marker gives
$|\zeta_i(q)-1|<e$, so $q$ belongs to the actual coordinate domain.
The vector block gives
$|\zeta_i(q)\eta_i(q)-\eta_i(p)|<e$. Subtract
$\zeta_i(q)\eta_i(p)$, use the triangle inequality, and divide by
the strictly positive marker. This argument applies to every $q$
on the carrier; it has no prior local-distance assumption.
\end{proof}

For a two-stratum chart, $|\eta_i(p)|<8$ implies
$\zeta_i(p)=1$. If $e<1/10$, the displayed bound is at most
$10e<20e$, giving KL 12.21 with room to spare. In particular, when
$|\eta_i(p)|<7$ and $e=1/100$, every point of the image ball of
radius $R_i/100$ comes from $|\eta_i|<7.2$. The marker prevents
unseen remote sheets from entering this ball. It gives neither a
global embedding nor a bound on intrinsic distance within a fiber.
The lemma also applies after any orthogonal projection which retains
the ENTIRE $i$th block. A projection discarding that marker has no
such implication.

\begin{lemma}[Exact scale on the first image; FC04]
\label{lem:fibration-image-scale}
Let $\widetilde A_1=\bigcup_{i\in I_2}\{|\eta_i|\le8\}$ and
$A_1=\bigcup_{i\in I_2}\{|\eta_i|\le7\}$, with preimages taken
inside their coordinate domains. These closed coordinate thresholds
match KL (12.6); the surrounding domains remain open. Set
$\widetilde S_1=\mathcal E^0(\widetilde A_1)$ and
$S_1=\mathcal E^0(A_1)$. For fixed $\Sigma_1>0$, define on
$\widetilde S_1$
\[
 r_1(x)=\Sigma_1x_\rho.
\]
This is positive, independent of any chosen preimage, and
$\Sigma_1$-Lipschitz for the restricted Euclidean distance. If $M$
is compact, it is bounded above and bounded away from zero.
\end{lemma}
\begin{proof}
Every preimage $p$ of $x$ has $\rho(p)=x_\rho$ by FC01.
Projection to $\R_\rho$ has operator norm one. Compactness and
positivity of $\rho$ give its minimum and maximum on $M$.
\end{proof}

Thus the radius in KL 12.7 requires no arbitrary preimage choice at
the first stage. This observation uses the retained scale coordinate.
It does not descend through $\pi_2$ or $\pi_3$, which discard that
coordinate; the projected edge/slim radii need their own
well-defined choice and comparison estimates. Likewise an adjusted
map must be checked anew if its scale coordinate changes.

\section{Quantitative calculus for the two-stratum model}
\label{sec:fibration-weighted-calculus}
KL 12.12 compares nearby normalized coordinates with affine
isometries followed by orthogonal projections. Its proof uses
absence of three-splittings, compatibility (AC76), changes of
basepoint/scale, adapted-coordinate comparison, the edge annulus
and the zero-shell lower ratio. These inputs still refer to the
SAME selected coordinates. KL 12.17 then composes the comparisons
with cutoffs. The following lemma makes explicit the cancellation
needed when a zero-piece radius is much larger than a chart radius.

\begin{lemma}[Scaled cutoff calculus; FC05]
\label{lem:fibration-scaled-cutoff}
Let $\phi:\R^k\to[0,1]$ be $C^2$, supported in the closed ball
of radius $C>0$, with $\|D\phi\|\le L_1$ and
$\|D^2\phi\|\le L_2$. For $s>0$ put
\[
 F_s(u)=\bigl(u\phi(u/s),\ s\phi(u/s)\bigr),\qquad
 A=1+(C+1)L_1,\quad B=2L_1+(C+1)L_2.
\]
Then $\|DF_s\|\le A$ and $\|D^2F_s\|\le B/s$ everywhere.
On a Riemannian domain $D$ let $u,v:D\to\R^k$ be $C^1$ and
assume
\[
 \sup_D|u-v|\le\varepsilon,\quad
 \sup_D\|Du-Dv\|\le\varepsilon,\quad
 \sup_D\|Dv\|\le L.
\]
With the $C^1$ norm defined as the maximum of the value and
derivative suprema,
\[
 \|F_s\circ u-F_s\circ v\|_{C^1(D)}
        \le (A+BL/s)\varepsilon.
\]
For at most $N$ orthogonal blocks satisfying common bounds and
$s\ge s_*>0$, the combined bound is
$\sqrt N(A+BL/s_*)\varepsilon$.
\end{lemma}
\begin{proof}
Put $z=u/s$. The first component has derivative
$h\phi(z)+zD\phi_z(h)$ and second derivative
\[
 s^{-1}\bigl(hD\phi_z(\ell)+\ell D\phi_z(h)
                       +zD^2\phi_z(h,\ell)\bigr).
\]
The marker has derivative $D\phi_z$ and second derivative
$s^{-1}D^2\phi_z$. Where derivatives of $\phi$ are nonzero,
$|z|\le C$; outside its support all these terms vanish.
The triangle inequality gives $A,B/s$. The mean value estimate on
the Euclidean segment between $u(p)$ and $v(p)$ bounds the values
and $DF_s(u(p))-DF_s(v(p))$. Write the derivative difference as
\[
 DF_s(u)(Du-Dv)+(DF_s(u)-DF_s(v))Dv.
\]
This gives the stated bound and the orthogonal square-sum estimate.
No convexity of the source domain is used.
\end{proof}

For a zero block set $s=R_j/R_i$ and
$u=s\eta_j$, using the source metric $R_i^{-2}g$.
Then $R_i^{-1}\mathcal E^0_j=F_s(u)$. KL 11.5(2) supplies
$s>T_0/20$ in the relevant overlap. A large $s$ improves the Hessian
bound; it must not be estimated as a large error multiplier. There
is no assumption that translated affine coordinates $u$ are small.
For other blocks whose cutoffs depend on several coordinates, use
the same chain rule with their actual jointly bounded first and
second derivatives. In particular $\zeta_{E'}$ depends on
$\zeta_{\mathrm{edge}}$ from (9.28)--(9.29), hence on multiple
edge charts. FC05 alone does not establish that joint derivative
budget or the variable-scale replacement $R_{E'}\rightsquigarrow R_i$.

\begin{lemma}[Rank and singular-value margin; FC06]
\label{lem:fibration-rank-margin}
Let $V,H$ be finite-dimensional inner-product spaces. All maps
$T,L,D$ below are linear. Let $T:\R^2\to H$ have the graph form $Tz=(z,Gz)$ in an orthogonal
decomposition and $\|T\|\le b$. Put $A^0=\operatorname{im}T$.
Suppose $L:V\to\R^2$ satisfies
\[
 \|L\|\le\ell,\qquad |L^*z|\ge a|z|\quad(z\in\R^2),
 \qquad a>0.
\]
If $D:V\to H$ satisfies $\|D-TL\|\le e<a$, then
$P=\pi_{A^0}D$ is SURJECTIVE onto $A^0$,
\[
 \|D-P\|\le e,\qquad
 (a-e)|v|\le |Pv|\le(b\ell+e)|v|
       \quad(v\perp\ker P).
\]
The constants are independent of $\dim H$.
\end{lemma}
\begin{proof}
$T^*T=I+G^*G\ge I$, so $T$ has least singular value at least one.
The adjoint of $TL:V\to A^0$ is $L^*T^*$ and has lower bound $a$.
The adjoint of $P$ differs by at most $e$, hence has lower bound
$a-e>0$. Thus $P$ is onto. In finite dimensions the nonzero
singular values of $P$ and $P^*$ agree, proving the lower bound on
$(\ker P)^\perp$. The upper bound follows from $\|TL\|\le b\ell$.
Finally $(I-\pi_{A^0})TL=0$ and orthogonal projection is a
contraction, giving the normal-error estimate.
\end{proof}

An inequality on $(\ker P)^\perp$ alone does not certify rank two:
the zero map satisfies it vacuously. FC06 requires and proves
surjectivity. In applying it, $L$ is the normalized derivative of
the actual $\eta_i$, whose quantitative rank bound must be supplied.
The graph derivative $T$ obtains a bound from the local block
estimates, not from a uniform bound on the total number of charts.
This replaces the global $\dim H$ assertion at the end of the
printed proof of KL 12.7 by a local or operator-norm bound.

\section{The first image-model producer and its remaining obligations}
\label{sec:fibration-two-cloud-input}
\begin{foundation}[Two-stratum image packet; FC07]
\label{found:fibration-two-cloud}
For the closed-carrier LC87 data, keep $A_1,\widetilde A_1$,
$S_1,\widetilde S_1,r_1$ from FC04. The source-qualified producer
is KL 12.7: after the simultaneous small-parameter restrictions
listed there, this is a $(2,\Gamma_1)$ cloudy two-manifold in $H$.
For every $x\in S_1$ it supplies an affine plane
$A_x=x+A_x^0$ and, at every preimage $p\in A_1$, a projected
derivative onto $A_x^0$ which is onto, with positive lower and
upper singular-value bounds and normal error less than $\Gamma_1$.
For each $i\in I_2$ it supplies a smooth map
$h_i$ from a neighborhood of $\overline B(0,8)\subset\R^2$
to $(H_i')^\perp$ satisfying
\[
 \|Dh_i\|\le\Omega_1R_i,\qquad
 \left\|R_i^{-1}\mathcal E^0-
       (\eta_i,R_i^{-1}h_i\circ\eta_i)\right\|_{C^1}<\Gamma_1
\]
on $\{|\eta_i|\le8\}$, with derivatives measured in $R_i^{-2}g$.
At $x\in S_1$ a choice of $i,p$ with $|\eta_i(p)|\le7$ gives
$A_x^0=\operatorname{im}(I,R_i^{-1}(Dh_i)_{\eta_i(p)})$.
This is an unimplemented source-qualified input, not an accepted
axiom or a proved consequence of the preceding inventory edges.
\end{foundation}

Here a cloudy packet means the following explicit convention from
KL Definition 20.1. It has $S\subset\widetilde S\subset H$ and
$r:\widetilde S\to(0,\infty)$; for all $x,y\in\widetilde S$,
\[
 |r(y)-r(x)|\le C(|x-y|+r(x)),
\]
and at every $x\in S$ the rescaled pointed subset
$(r(x)^{-1}\widetilde S,x)$ is $\delta$-close in pointed
Hausdorff distance to an affine $k$-plane through $x$ with the
same rescaling. Distances are restricted ambient distances, not
intrinsic length distances on the image. A future binding must
state the tested radius and the two coverage directions of that
pointed Hausdorff convention and translate them to Chapter 3's
fixed-error interfaces. We do not replace the enlarged set
$\widetilde S$ by $S$ or silently choose an unrecorded error
convention. FC04 supplies the radius condition directly for the
first image; it does not prove the plane approximation.

\paragraph{Proof decomposition to implement.}
First obtain the enlarged-ball count for all active supports and
the simultaneous coordinate comparisons of KL 12.12. These use
AC76 and the inherited adapted-coordinate calculus; they retain
the edge annular and zero-shell domains. Next reconstruct the local
graph approximation of 12.17, including the multi-coordinate
edge cutoff and the variable-scale terms. FC05 handles the
potentially large zero radius. FC03 localizes the WHOLE image
near a core point. Local graph derivative bounds and a uniform
modulus for their variation, together with actual coordinate-image
coverage, must then give both sides of the pointed Hausdorff
comparison. A bounded first derivative by itself is insufficient
for uniform tangent approximation; a bounded second derivative
or an explicitly proved derivative modulus is needed. Finally use
FC06 and the quantitative rank of the original adapted map to
obtain the projected rank bounds. No step bounds $\dim H$ by
pointwise multiplicity.

\paragraph{Dependency and parameter boundary.}
FC01--FC06 are definitions or written conditional arguments using
inherited finite-dimensional calculus and linear algebra. FC07
remains source-qualified. Its geometric producer uses LC20's
no-three-stratum conclusion, AC76's compatibility and the LC87
local packet; its analytical auxiliaries are FC01--FC06. Given
fixed overlap/support bounds and $\Gamma_1$, the scale
$\Sigma_1$ and comparison errors must be selected in KL 12.7's
order before the local witnesses. The global simultaneous order
with later stages is still the separate KL (15.4) task.

Next the chapter must supply the cloud-to-smooth-submanifold
construction of KL 20.2 and the first adjustment in Section 13.2,
then the projected edge/slim analogues and their successive
compatibility. The author correction to 20.2 requires recentering
all basepoints at the origin and retaining the enlarged sets in
the limit. The corrected $D^3$ alternative in 14.1(2), proper
restrictions and LC82's common compact isolation remain required.
The static output flows to LC90, Chapter 6 refinement and
Chapter 15 transport; neither that output nor LC90 is a premise
of FC01--FC07. Chapter 13's completed blueprint coverage and its
explicit producer gates are unchanged.

\paragraph{Source record.}
L133--L135 and \code{reference_checks_revision97.md} give the exact
KL14 passages, hashes, author corrections and project findings.

\section{Counting all supports near a two-stratum chart}
\label{sec:fibration-ball-packing}
FC08--FC11 discharge elementary parts of the preceding decomposition.
They do not change FC07's source-qualified status. In particular,
a bound for supports containing a single point does not bound those
meeting a ball. The following proof expands KL 8.5, 9.24 and 10.4
in the form needed by 12.12--12.17.

\begin{lemma}[Enlarged-ball packing; FC08]
\label{lem:fibration-ball-packing}
Let $M$ be a complete smooth Riemannian three-manifold and let
$\rho>0$ be $\Lambda$-Lipschitz. Fix $R,a,C>0$ with
$\Lambda\max(R,C)\le1/4$. A finite family has centers $p_j$,
$r_j=\rho(p_j)$, pairwise disjoint OPEN balls $B(p_j,ar_j)$,
and sets $S_j\subset\overline B(p_j,Cr_j)$. For $p\in M$ put
\[
 r=\rho(p),\quad D=B(p,Rr),\quad
 J=\{j:S_j\cap D\ne\varnothing\},\quad
 H=4(R+2C+a).
\]
Assume, for a fixed $\kappa\ge0$ and EVERY $j\in J$, that
\[
 \operatorname{Ric}_g\ge-2\kappa^2r_j^{-2}g
       \quad\hbox{on the WHOLE ball }B(p_j,Hr_j).
\]
With the inherited local Bishop--Gromov comparison theorem supplied,
\[
 \#J\le\frac{V_\kappa(H)}{V_\kappa(a)},\qquad
 V_\kappa(t)=4\pi\int_0^t s_\kappa(u)^2\,du,
\]
\[
 s_0(u)=u,\qquad
 s_\kappa(u)=\frac{\sinh(\kappa u)}{\kappa}\ (\kappa>0).
\]
Moreover $r_j/r\in[1/2,2]$ for every $j\in J$.
\end{lemma}
\begin{proof}
An intersection point gives $d(p,p_j)\le Rr+Cr_j$.
Lipschitz control yields
\[
 \frac{1-\Lambda R}{1+\Lambda C}
 \le\frac{r_j}{r}\le
 \frac{1+\Lambda R}{1-\Lambda C},
\]
which lies in $[1/2,2]$. Write $D_0=R+2C$; all centers lie
within $D_0r$ of $p$. If $J$ is empty there is nothing to prove.
Otherwise choose $j_0\in J$ with smallest core-ball volume.
Every core is contained in
$B(p_{j_0},(2D_0+2a)r)\subset B(p_{j_0},Hr_{j_0})$.
Disjointness and positivity of smooth ball volumes give
\[
 \#J\operatorname{vol}B(p_{j_0},ar_{j_0})
 \le\operatorname{vol}B(p_{j_0},Hr_{j_0}).
\]
Apply volume comparison at $p_{j_0}$ after rescaling the metric
by $r_{j_0}^{-2}$. No absolute lower volume bound is used.
\end{proof}

\paragraph{A bound chosen before $\Delta$.}
For the two-stratum family use $R=10$, $a=1/3$, $C=10$,
$\kappa=1$. For the edge and slim families use respectively
\[
 a=\Delta/3,\qquad C=C_0\Delta,\qquad
 C_0=100\ \hbox{or}\ 10^6,\qquad \kappa=1/\Delta,
 \qquad\Delta\ge1.
\]
The $C_0$ values are safe enclosures by the ACTUAL coordinate
domains of KL 9.18 and 10.1, not the smaller coordinate-value
bound incorrectly printed for the slim block in 12.4. The cores
are those selected in KL 9.5 and 10.2. Since
$V_{1/\Delta}(t)=\Delta^3V_1(t/\Delta)$, FC08 gives
\[
 \#J\le
 \frac{V_1\bigl(4(10+2C_0+1/3)\bigr)}{V_1(1/3)}.
\]
This constant is independent of $\Delta$, the carrier and its
total chart count. The required curvature bound is now
$\operatorname{Ric}\ge-2(\Delta r_j)^{-2}g$ on the WHOLE enlarged
ball. A bound only at the center or only on the core is insufficient.
In the standing sequence, LC09's uniform-in-center curvature
exhaustion supplies this after all these parameters are fixed:
if the available radius is $Q r_j$ with sectional curvature at least
$-(Qr_j)^{-2}$, take $Q\ge\max(H,\Delta)$.
This is the rescaling used in LC05, not a new curvature theorem.

Thus a universal enlarged-ball multiplicity budget can be chosen
before $\Delta$, while $\Lambda$ and the eventual sequence tail
may depend on $\Delta$. This respects the relevant direction of
KL (15.4), where $\mathcal M\prec\Delta\prec\Lambda$.
The full simultaneous order, including later adjustments, is still
open. Applying curvature $-1$ comparison on a radius of order
$\Delta$ without this rescaling would give a bound growing with
$\Delta$ and would not justify that order.

\begin{lemma}[At most one zero support; FC09]
\label{lem:fibration-zero-support}
Let $\rho>0$ be $\Lambda$-Lipschitz on a metric space.
Put $D=B(p,L\rho(p))$, $L>0$. Suppose the selected zero balls
$B(z_j,R_j)$ are pairwise disjoint, $R_j\ge T\rho(z_j)$,
and $S_j\subset\overline B(z_j,\theta R_j)$ with
$0<\theta<1$. Retain the LC62 comparison
\[
 d(z_j,q)\le10R_j\quad\Longrightarrow\quad
                   R_j/\rho(q)\ge T/20.
\]
If $L\Lambda\le1/4$ and
$T>\max\{40L,40L/(1-\theta)\}$, at most one $S_j$ meets $D$.
\end{lemma}
\begin{proof}
For a meeting support, first establish the domain of the LC62
comparison. An intersection point and Lipschitz control give
\[
 d(z_j,p)\le\theta R_j+L\rho(p)
 \le(\theta+L/T)R_j+L\Lambda d(z_j,p).
\]
Consequently $d(z_j,p)/R_j<(1+1/40)/(3/4)<2<10$.
Only NOW use LC62 to obtain $R_j/\rho(p)\ge T/20$.
For every $x\in D$, using the intersection point once more,
\[
 d(z_j,x)<\theta R_j+2L\rho(p)
          \le(\theta+40L/T)R_j<R_j.
\]
Hence $D\subset B(z_j,R_j)$. Since $p\in D$, two meeting
supports would contradict disjointness of the selected open balls.
\end{proof}

LC31 gives the closed-support enclosure with
$\theta=37/40$ when its radial value error is less than $1/40$;
LC80 supplies the selected balls and LC62 their local ratio.
The preliminary distance estimate prevents circularly applying
that ratio before checking its domain. This count does not require
an upper bound on $R_j/\rho(p)$. It also does not supply the lower
annular margin needed to compare the SAME radial coordinate;
that remains part of KL 12.12's zero-shell comparison.
Together FC08--FC09 bound the number of constant-radius supports
meeting $B(p,10\rho(p))$; the scale coordinate and the SINGLE
$E'$ block add two tags. The latter block's cutoff still depends
on all nearby edge coordinates.

\section{The joint edge cutoff and the variable scale}
\label{sec:fibration-joint-cutoff}
For a map on a Riemannian domain use
$\|F\|_{C^1}=\max(\sup|F|,\sup\|DF\|)$.
The next lemma is a finite-dimensional calculus statement, independent
of the geometric existence of the coordinate comparisons.

\begin{lemma}[Joint edge cutoff calculus; FC10]
\label{lem:fibration-joint-cutoff}
Let $n$ be a nonnegative integer, $\Delta\ge1$, and let
$s_j\in[1/2,2]$ be CONSTANT.
Take smooth profiles $f,g,h,\chi:\R\to[0,1]$ with
$f$ supported in $[-9,9]$, $h$ supported in $[1/5,9]$,
and $\chi(0)=0$. Their first and second derivatives are globally
bounded by $P\ge1$. For $U=(u_1,\ldots,u_n,v)$ define
\[
 z_j=f(u_j/\Delta)g(v/\Delta),\qquad
 z_E=h(v/\Delta)\chi\left(\sum_{j=1}^nz_j\right),
\]
\[
 W_j=s_j(u_jz_j,z_j),\qquad W_E=(vz_E,z_E),\qquad
 W=(W_1,\ldots,W_n,W_E).
\]
Use the orthogonal Euclidean norms. Put
\[
 K=10(n+1)^2P^3,\quad
 A_n=\sqrt{n+1}(2+20K),\quad B_n=24\sqrt{n+1}K.
\]
Globally on $\R^{n+1}$,
\[
 |W|\le20\sqrt{n+1}\Delta,\qquad
 \|DW\|\le A_n,\qquad \|D^2W\|\le B_n/\Delta.
\]
If $U,V$ are $C^1$ maps on the SAME domain with
$\|U-V\|_{C^1}\le\epsilon$ and $\|DV\|\le L_0$, then
\[
 \|W\circ U-W\circ V\|_{C^1}
       \le(A_n+B_nL_0/\Delta)\epsilon.
\]
\end{lemma}
\begin{proof}
For each $z_j$, the product rule gives gradient bound
$2P/\Delta$ and Hessian bound $4P^2/\Delta^2$.
For $Z=\sum_jz_j$, sum these bounds before composing with $\chi$.
The chain rule gives bounds $2nP^2/\Delta$ and
$4(n^2+n)P^3/\Delta^2$ for $\chi(Z)$.
Multiplication by $h(v/\Delta)$ gives bounds
$(1+2n)P^2/\Delta$ and $(1+8n+4n^2)P^3/\Delta^2$.
Thus every cutoff has gradient at most $K/\Delta$ and Hessian
at most $K/\Delta^2$.

Where a cutoff or either of its first two derivatives can be
nonzero, the coordinate weighting its block has absolute value at
most $9\Delta$. This uses the support of $f$ for $W_j$ and of
$h$ for $W_E$; no bound on all other coordinates is required.
Differentiating the weighted blocks, including their markers,
gives bounds $2+20K$ and $24K/\Delta$. Their values are at most
$20\Delta$. Square and sum over the $n+1$ blocks.
For the final assertion, integrate $DW$ and $D^2W$ on the straight
segment between $U(x)$ and $V(x)$, and write
\[
 D(W\circ U)-D(W\circ V)
 =DW(U)(DU-DV)+(DW(U)-DW(V))DV.
\]
The global Hessian bound controls this segment even when it leaves
a cutoff support. The value estimate is bounded by $A_n\epsilon$.
\end{proof}

\paragraph{Binding the actual KL profiles.}
In the corrected scalar-cutoff convention retained by LC31, take
\[
 f=\Phi_{-9,-8,8,9},\quad g=\Phi_{8,9},\quad
 h=\Phi_{1/5,3/10,8,9},\quad \chi=1-\Phi_{1/2,1}.
\]
These are precisely the profiles in KL (9.19), (9.28) and (9.29)
after division by $\Delta$. In particular, $z_E$ is a function
of the SUM of the edge cutoffs, not of $v$ alone. The two-stage
chain rule is necessary. The source's corrected profile construction
supplies a finite universal $P$; no numerical value is asserted.
For $n=0$ the sum is zero and $z_E$ vanishes, as required.
On a fixed chart set $s_j=\rho(p_j)/\rho(p_i)$; FC08 gives the
stated ratio interval for all nearby edge supports. Actual
coordinates and every model replacement must be defined on the
SAME domain. A support count does not itself extend a coordinate
across the boundary of its open domain.

\begin{lemma}[Freezing the single variable scale; FC11]
\label{lem:fibration-freeze-scale}
Fix $L>0$ and $\Delta\ge1$. Let $R=\rho(p)$ and
$D=B(p,LR)$, with $\rho$ smooth, positive
and $\Lambda$-Lipschitz, and $L\Lambda\le1/2$.
All derivatives on $D$ now use $R^{-2}g$. Put $\sigma=\rho/R$.
Then
\[
 |\sigma-1|\le L\Lambda,\qquad |d\sigma|\le\Lambda,
 \qquad 1/2\le\sigma\le3/2.
\]
If a $C^1$ block $F$ on $D$ has $|F|\le V\Delta$ and
$\|DF\|\le A$, with nonnegative constants, then
\[
 \|\sigma F-F\|_{C^1}
       \le\Lambda\bigl((L+1)V\Delta+LA\bigr).
\]
If a scalar $u$ on $D$ has $|u|\le C\Delta$ and
$|du|\le D_1$, then
\[
 \|u/\sigma-u\|_{C^1}
       \le\Lambda\bigl((2L+4)C\Delta+2LD_1\bigr).
\]
\end{lemma}
\begin{proof}
The first bounds follow directly from the Lipschitz property and
metric scaling. Differentiate $(\sigma-1)F$; its value is bounded
by $L\Lambda V\Delta$ and its derivative by
$\Lambda V\Delta+L\Lambda A$. Their maximum is bounded by the
stated expression. For the quotient use
$|1/\sigma-1|\le2L\Lambda$ and
$|d(1/\sigma)|\le4\Lambda$ in the same product rule.
\end{proof}

For the actual $E'$ block, $R^{-1}\mathcal E^0_{E'}
=\sigma(\eta_{E'}\zeta_{E'},\zeta_{E'})$.
FC10 supplies its fixed-scale cutoff calculus, and FC11 controls
the additional factor $\sigma$. If a SAME-map comparison is
stated for $u=(\rho/R)\eta_{E'}$, the quotient estimate converts
it to the $\eta_{E'}$ argument used in the joint cutoff, provided
the displayed value and derivative bounds hold on the entire tested
domain. Those bounds require the edge/annular localization; they
are not inferred merely from a coordinate existing somewhere.
The error is generally of order $\Delta\Lambda$, not uniformly
of order $\Lambda$. Choosing $\Lambda$ later than $\Delta$ is
allowed by KL (15.4). Freezing it without these product and quotient
terms would omit derivatives of the actual global map.

\paragraph{Current disposition of FC07.}
FC08--FC11 are written conditional proofs: volume comparison,
selected local witnesses and any actual coordinate comparisons are
supplied inputs, not new Lean results. They provide enlarged-ball
counts, the zero-support count, a joint cutoff Hessian bound and
explicit variable-scale error estimates. FC05 continues to handle
the potentially unbounded zero-radius ratio. The full SAME-coordinate
comparison on buffered domains in KL 12.12 remains to be assembled,
including the zero-shell and edge-annular domains; then these
estimates can be combined into 12.17's local model. Actual two-sided
image coverage, the model derivative modulus and translation to
Chapter 3's pointed-Hausdorff conventions remain open. A Hessian
bound for a cutoff network is not yet a modulus for the assembled
local graph. FC07, KL 20.2 smoothing and Section 13.2 adjustment
are therefore still source-qualified or undeveloped. L136--L138
and \code{reference_checks_revision98.md} record the source checks.

\section{Common domains for the original overlap coordinates}
\label{sec:fibration-common-domains}
The next four targets expand the domain and differential steps of
KL 12.12. Coordinate existence on a whole test ball must be checked
before invoking a $C^1$ estimate there. The functions remain the
SAME ones used in the global map; choosing new adapted coordinates
would not establish compatibility of that map.

\begin{lemma}[Buffered support domain; FC12]
\label{lem:fibration-buffered-domain}
In a metric space $M$ let $r,r_j,L>0$, $c\ge0$, and $b>c$.
Suppose $r_j/r\ge1/2$, $S_j\subset\overline B(p_j,cr_j)$,
$B(p_j,br_j)\subset U_j$, and $b-c>4L$. If
$S_j\cap B(p,Lr)\ne\varnothing$, then
\[
 B(p,Lr)\subset B(p_j,(c+4L)r_j)\subset U_j.
\]
Moreover the distance to $M\setminus U_j$, when nonempty, is at
least $(b-c-4L)r_j$ on this test ball.
\end{lemma}
\begin{proof}
Choose a meeting point $q$. For $x\in B(p,Lr)$,
$d(x,p_j)\le d(x,q)+d(q,p_j)<2Lr+cr_j\le(c+4L)r_j$.
For $y\notin U_j$, $d(y,p_j)\ge br_j$, so the reverse triangle
inequality gives the asserted distance bound. No coordinate is
extended outside its original open domain.
\end{proof}

For a two-stratum chart, KL 8.2 and the support enclosure in 8.5
give $b=200$, $c=10$; FC08 supplies the scale ratio. With $L=10$
the domain margin is $150r_j$. For edge and slim charts, FC08's
safe enclosures $100\Delta r_j$ and $10^6\Delta r_j$ were enough
to COUNT supports, but equal their coordinate-domain radii and
leave no buffer. They cannot be reused as buffered enclosures.
Sufficient stronger bounds are displayed here as implementation
budgets, not as newly proved consequences of the cutoff values:
\[
 \begin{array}{c|c|c}
 \text{family}&\text{coordinate-domain radius}/r_j&
                   \text{sufficient support radius}/r_j\\ \hline
 \text{edge}&100\Delta&20\Delta\\
 \text{slim}&10^6\Delta&950000\Delta
 \end{array}
\]
For $\Delta\ge1$, these give margins greater than $4L=40$.
Their production requires estimates for the ACTUAL local product
map and its residual coordinate, with the map distortion and
adapted-coordinate value error charged. A bound on $|\eta_j|$
alone cannot control distance from $p_j$: it does not control the
residual factor. Retain LC84--LC85's source-qualified geometric
producer for this step. KL 12.12's printed $10^5\Delta r_j$ slim
enclosure also cannot be inferred from the cutoff in (10.2), which
allows coordinate values up to $9\cdot10^5\Delta$.

\begin{lemma}[A buffered zero shell; FC13]
\label{lem:fibration-zero-shell}
On a Riemannian manifold let $\rho>0$ be $\Lambda$-Lipschitz,
$r=\rho(p)$, $L>0$,
and $D=B(p,Lr)$. Retain FC09's selected ball $B(z,R_0)$ and
its local LC62 ratio, with $R_0\ge T\rho(z)$,
$L\Lambda\le1/4$ and $T\ge1600L$.
Suppose a closed support $S$ meets $D$, the SAME radial function
$\eta:M\to\R$ obeys
\[
 |\eta(x)-d(z,x)/R_0|<e,\qquad 0<e<1/40,
 \qquad S\subset\{1/5\le\eta\le9/10\},
\]
and $\eta$ is smooth on $\{1/10<d(z,x)/R_0<10\}$.
Then
\[
 R_0/r\ge T/20,\qquad
 D\subset\{3/20\le d(z,x)/R_0\le19/20\}.
\]
In particular $\eta$ is defined and smooth on a neighborhood of
$D$, in the interior of the CLOSED shell used by LC73.
\end{lemma}
\begin{proof}
The radial error bounds the support distances between
$7R_0/40$ and $37R_0/40$. Thus FC09 applies with $\theta=37/40$:
its required strict bound $T>40L/(1-\theta)$ follows from
$T\ge1600L$. Its domain-first argument gives $R_0/r\ge T/20$.
Choose $q\in S\cap D$. For $x\in D$,
\[
 \frac{|d(z,x)-d(z,q)|}{R_0}
       <\frac{2Lr}{R_0}\le\frac{40L}{T}\le\frac1{40}.
\]
Subtract and add $1/40$ to the two support bounds. The resulting
interval $[6/40,38/40]$ is strictly inside $(1/10,1)$.
\end{proof}

LC31 provides the radial value error and support range; LC73 and
LC80 retain the original radial coordinate, its selected model and
its annular adaptedness. FC13 proves the missing domain buffer,
without an upper bound on $R_0/r$. Adaptedness on a ball normalized
by $10r$, rather than $r$, still requires the explicit rescaling
and test-domain restrictions of KL Remark 4.35. Merely belonging
to the shell does not identify those two adaptedness conventions.

\section{The weak-edge coordinate as an auxiliary argument}
\label{sec:fibration-auxiliary-edge}
KL 12.14 uses every argument needed by each cutoff. This is a
stronger requirement than listing its nonzero OUTPUT blocks.
For example, in FC10's formula a small positive $z_j$ may still
depend on $v$, while $\chi(\sum_jz_j)=0$ makes $z_E=0$.
Thus the definition of the active block set alone does not show
that it contains all cutoff arguments. Either prove an additional
geometric implication or retain an auxiliary coordinate. We use
the latter, with a region where it can be omitted exactly.

\begin{lemma}[Weak-edge localization and cutoff alternative; FC14]
\label{lem:fibration-weak-edge-domain}
Let $E\subset E'$ be nonempty subsets of a metric space, and let
$\rho>0$ be $\Lambda$-Lipschitz. Distances to sets mean distances
to their closures. Let $P\ge0$ be globally $(1+\zeta)$-Lipschitz,
where $0\le\zeta\le1$, and put $v=P/\rho$.
Fix $r=\rho(p)$, $L>0$, $D=B(p,Lr)$ and
\[
 \Delta\ge\max\{1,120L\},\qquad
 \Delta\Lambda\le1/100,\qquad L\Lambda\le1/4.
\]
Suppose $v(q)\le9\Delta$ for some $q\in D$ and
$d(p,E)\le30\Delta r$. Then exactly one of the following
alternatives holds:
\begin{enumerate}[label=(\roman*),leftmargin=*]
\item $v<3\Delta/20$ everywhere on $D$;
\item some $y\in D$ has $v(y)\ge3\Delta/20$, and throughout $D$
\[
 \Delta/10\le v\le181\Delta/20<10\Delta,
 \qquad d(\cdot,E)/\rho<50\Delta.
\]
\end{enumerate}
For the actual profiles of FC10, in case (i) every edge cutoff is
$f(u_j/\Delta)$ and the $E'$ cutoff is identically zero on $D$.
If moreover $M$ is Riemannian and $\rho$ is smooth, then in case
(ii) the smoothing conclusion of KL 9.12, when supplied for $P$,
makes $v$ smooth on a neighborhood of $D$.
\end{lemma}
\begin{proof}
On $D$, $3r/4\le\rho\le5r/4$. If $a,x\in D$ and
$v(a)\le b\Delta$, the quotient identity gives
\[
 |v(x)-v(a)|
 \le\frac{(1+\zeta)+b\Delta\Lambda}{\rho(x)}d(x,a)
 <\frac{8L}{3}(2+b/100).
\]
Use $a=q$, $b=9$ to obtain $|v(x)-v(q)|<6L\le\Delta/20$.
In particular $v(x)\le181\Delta/20$ throughout $D$.
If case (i) fails, choose $y$ as in (ii); repeating the same
estimate with $a=y$, $b=181/20$ still gives an upper bound
$6L$. Thus $v(x)\ge3\Delta/20-\Delta/20=\Delta/10$.
Independently, the strong-edge distance satisfies
\[
 \frac{d(x,E)}{\rho(x)}
 <\frac{30\Delta+L}{1-L\Lambda}
 \le40\Delta+4L/3<50\Delta.
\]
This is the SECOND domain condition in KL 9.12(2); control of
$v$ alone would not establish it. In case (i),
$g(v/\Delta)=1$ and $h(v/\Delta)=0$ for the actual KL profiles.
In case (ii), both displayed domain conditions lie in the
source's smoothing region. Since $\rho$ is positive and smooth
in the Riemannian application, its quotient with the supplied
smooth $P$ is smooth there.
\end{proof}

\paragraph{Application and retained hypotheses.}
In the Riemannian application set $P=\rho_{E'}$ from KL 9.12;
its difference from $d_{E'}$ is $\zeta$-Lipschitz, so it has the
stated Lipschitz constant. Smoothness of $\rho$ is required for
this application; the metric inequalities and the low alternative
do not require it. If an edge cutoff has nonzero values on $D$,
its $g$ factor supplies a point with $v<9\Delta$. A CLOSED edge
support meeting the OPEN ball $D$ also gives such a point, by the
definition of support. If no edge support meets $D$, the sum in
$\chi$ vanishes there and the $E'$ block vanishes as well.

The additional hypothesis $d(p,E)\le30\Delta r$ is an explicit
strong-edge proximity obligation, not a consequence of distance
to $E'$ or of the $E'$ support range. The sufficient $20\Delta r_j$ support enclosure in FC12
would supply it: if $p_j\in E$ and its support meets $D$, then
\[
 \frac{r_j}{r}\le\frac{1+L\Lambda}{1-20\Delta\Lambda}<1.26,
 \qquad
 \frac{d(p,E)}r< L+20\Delta(1.26)<26\Delta<30\Delta.
\]
Here $L\Lambda\le1/12000$ follows from FC14's parameter bounds;
the denominator is positive. Thus strong-edge proximity follows
conditionally from that stronger enclosure, without an attained
nearest point in $E$. The coarse $100\Delta r_j$ enclosure used
only for counting does not give this conclusion. KL 9.7 and LC84's
local half-plane producer are the relevant geometric inputs. We do not discard
the $d_E/\rho$ condition in KL 9.12 or replace $E$ by $E'$.

When case (ii) holds, retain $v$ as a cutoff ARGUMENT even if
the $E'$ OUTPUT block vanishes on $D$. This adds at most one
auxiliary scalar and does not change the global map, its index
sets or the local packing count. In case (i), replace the cutoff
formulas by their exact reduced expressions before making model
substitutions; do not evaluate a possibly nonsmooth $v$ merely
to multiply it by zero. The alternative itself is a local proof
case, not a change of the global cutoff.

\section{From shared directional tests to an affine comparison}
\label{sec:fibration-directional-comparison}
\begin{lemma}[Directional saturation gives $C^1$ control; FC15]
\label{lem:fibration-directional-comparison}
Let $M$ be a complete connected Riemannian manifold,
$D=B(p,L)$ with $L>0$, and $1\le k\le m$ integers.
Take $C^1$ maps $F:D\to\R^k$, $G:D\to\R^m$ and a CONSTANT
linear map $A:\R^m\to\R^k$ with $AA^*=I$.
Assume $0\le\epsilon\le1$, $\|dF_a\|\le1+\epsilon$ for
every scalar component $a$, and $\|DG\|\le1+\epsilon$.
For every $x\in D$ and every $a\in\{1,\ldots,k\}$ suppose
there is a unit $w_{x,a}\in T_xM$ such that
\[
 dF_a(w_{x,a})\ge1-\epsilon,\qquad
 d(AG)_a(w_{x,a})\ge1-\epsilon.
\]
Define $b=F(p)-AG(p)$ and $T(u)=Au+b$. Then
\[
 \|F-T\circ G\|_{C^1(D)}
 \le2\max\{1,L\}\sqrt{k(4\epsilon+\epsilon^2)}.
\]
In particular a prescribed error $\tau>0$ is obtained strictly
when $0\le\epsilon<\min\{1,\tau^2/(20k\max\{1,L\}^2)\}$.
No continuity in the choices of $w_{x,a}$ is required.
\end{lemma}
\begin{proof}
For a linear functional $\ell$ with $\|\ell\|\le1+\epsilon$
and a unit vector $w$ with $\ell(w)\ge1-\epsilon$, the Riesz
vector satisfies
\[
 |\ell^\sharp-w|^2
 =\|\ell\|^2+1-2\ell(w)\le4\epsilon+\epsilon^2.
\]
Every row of $A$ has norm one because $AA^*=I$.
Apply this inequality to $dF_a$ and $d(AG)_a$ with their SAME
$w_{x,a}$, and take the triangle inequality. Square summation
over the $k$ scalar components gives
\[
 \|DF-A\,DG\|\le2\sqrt{k(4\epsilon+\epsilon^2)}.
\]
The centered difference $F-TG$ vanishes at $p$. Integrate its
derivative along a minimizing segment from $p$ to $x\in D$;
that segment stays in $D$ and has length less than $L$.
This proves the value bound and hence the stated max $C^1$ norm.
Since $4\epsilon+\epsilon^2\le5\epsilon$ for
$0\le\epsilon\le1$, the last parameter choice follows.
\end{proof}

\paragraph{What this supplies for KL 12.12.}
Use the fixed metric $r^{-2}g$, $r=\rho(p_i)$, and $L=10$.
For a CONSTANT-radius block take $F=(R_j/r)\eta_j$ and
$G=\eta_i$. Normalization cancels the $R_j/r$ in the derivative
bound: a component of $\eta_j$ with Lipschitz constant
$1+\epsilon$ in $R_j^{-2}g$ gives the same bound for $F$ in
$r^{-2}g$. The affine translation $b$ retains the original value
at $p_i$; it need not be small, especially for a zero shell.
FC05 permits the resulting large zero-radius ratio.

For the high weak-edge alternative the natural normalized
coordinate is instead
\[
 F=P/r=(\rho/r)\eta_{E'}.
\]
Its derivative bound comes directly from $P$'s Lipschitz bound.
The factor $\rho/r$ is VARIABLE. FC11 converts back to the
$\eta_{E'}$ argument of FC10, charging its product and reciprocal
errors. This avoids treating $R_{E'}$ as a constant chart radius.

The remaining geometric input to FC15 is concrete: the SAME
orthogonal projection $A$ and the common tested unit directions
must satisfy the two saturation inequalities on the entire
buffered ball. KL 4.28 and 4.31 construct this from compatible
splittings and adapted-coordinate derivative tests; their supplied
proofs were reread. AC79 handles recentering of the original
splittings, LC20 excludes rank three and AC76 gives compatibility.
But metric compatibility alone is not a $C^1$ theorem. The
quantitative transfer of all derivative-test domains and the
Euclidean factor isometry still needs binding to the actual maps,
including the edge-distance smoothing and scaled radial coordinate.
FC15 proves the calculus implication once those directional data
are supplied; it does not assume the desired $C^1$ conclusion.

\paragraph{Current implementation boundary.}
FC12--FC15 supply a common-domain lemma, an actual zero-shell
buffer, the weak-edge cutoff alternative and a quantitative
$C^1$ comparison from scalar tests. The sufficient edge/slim
support enclosures, strong-edge proximity and geometric production
of the common directions remain explicit obligations. Next bind
those to the original local packet, then assemble KL 12.17's
local graph with FC05 and FC08--FC11. FC07 remains source-qualified;
its two-sided image coverage, assembled-graph derivative modulus
and Chapter 3 convention translation are not supplied here.
The corrected KL 20.2 smoothing, Section 13.2 adjustment and full
parameter order remain downstream. Sources and project inferences
are recorded in L139--L141 and \code{reference_checks_revision99.md}.

\section{Value control for the original adapted coordinates}
\label{sec:fibration-adapted-values}
We now supply the edge/slim buffers requested in FC12, conditional
on the original local approximation and adapted-coordinate data.
No fiber classification or new smooth comparison model is needed
for these metric enclosures. The coordinates used in the cutoffs
must retain their original center values.

\begin{lemma}[Long tests control coordinate values; FC16]
\label{lem:fibration-adapted-values}
Let $M$ be a complete connected Riemannian manifold, $k\ge1$ an
integer, $R>0$, and $\delta,\alpha\ge0$. Take
\[
 s>2R+4\delta,\qquad H>s+R+3\delta,
\]
\[
 \phi=(u,z):B(p,H)\longrightarrow\R^k\times Z,
 \qquad \phi(p)=(0,z_0).
\]
Use the Euclidean product metric. Suppose $\phi$ has distortion at
most $\delta$, and every point of the target ball $B((0,z_0),H-\delta)$
is within $\delta$ of $\phi(B(p,H))$. Let $\eta:B(p,R)\to\R^k$
be $C^1$, with every scalar differential bounded by $1+\alpha$.
For every $x\in B(p,R)$ and $y\in B(p,H)$ with $d(x,y)>R$,
assume that EVERY initial unit tangent $w$ of a minimizing segment
from $x$ to $y$ satisfies the original adapted test
\[
 \left|D\eta_x(w)-\frac{u(y)-u(x)}{d(x,y)}\right|\le\alpha.
\]
Put
\[
 \epsilon=\alpha+\frac{2R+4\delta}{s-R-2\delta}\le1,
 \qquad K=\sqrt{4\epsilon+\epsilon^2}.
\]
Then, on $B(p,R)$,
\[
 |\eta(x)-\eta(p)-u(x)|
 \le\sqrt{k}\left(RK+4\delta+
                  \frac{(R+\delta)^2}{2(s-R-\delta)}\right).
\]
This uses the inherited Riemannian distance-gradient and locally
Lipschitz integration facts; it does not require smoothness or
continuity of the approximation $\phi$.
\end{lemma}
\begin{proof}
For each $a$ choose $q_a\in B(p,H)$ whose image is within
$\delta$ of $(se_a,z_0)$. Distortion at the basepoint gives
$|d(p,q_a)-s|\le2\delta$. For $x\in B(p,R)$,
\[
 d(x,q_a)>s-R-2\delta>R,\qquad
 \frac{u_a(q_a)-u_a(x)}{d(x,q_a)}
 \ge1-\frac{2R+4\delta}{s-R-2\delta}.
\]
Indeed $|u(x)|\le R+\delta$, the numerator is at least
$s-R-2\delta$, and the denominator is at most $s+R+2\delta$.
The adapted test implies $d\eta_a(w)\ge1-\epsilon$ for every
minimizing direction to $q_a$. The Riesz estimate in FC15 gives
$|(d\eta_a)^\sharp-w|\le K$.

The distance function $d_{q_a}$ is locally Lipschitz and, at its
points of differentiability in this ball, has gradient $-w$.
Consequently $|d(\eta_a+d_{q_a})|\le K$ almost everywhere.
The inherited local gradient-to-Lipschitz theorem now bounds its
change along any curve contained in the ball by $K$ times the
curve length. Apply it along a radial minimizing segment from $p$
to $x$, wholly contained in $B(p,R)$, to obtain
\[
 |\eta_a(x)-\eta_a(p)-d(p,q_a)+d(x,q_a)|\le RK.
\]
This step uses the locally Lipschitz theorem, not the unjustified
claim that an ambient almost-everywhere estimate holds on every
chosen geodesic.

Let $D_a=d(\phi(x),(se_a,z_0))$. Distortion and the chosen target
error give $|d(x,q_a)-D_a|\le2\delta$. Pythagoras yields
\[
 0\le D_a-(s-u_a(x))
 \le\frac{(R+\delta)^2}{2(s-R-\delta)}.
\]
Together with $|d(p,q_a)-s|\le2\delta$, this proves the scalar
estimate; square summation proves the assertion.
\end{proof}

For fixed $R,k$ and any requested positive value error, first
choose $s$ sufficiently large, then $\alpha,\delta$ sufficiently
small. The right side tends to zero in that order. The long test
points must lie in the ORIGINAL tested domain: for a radius-$R$
adapted chart, KL 4.21 and Remark 4.35 give tests in
$B(p,R/\alpha)$ with separation greater than $R$. Choose the
parameters so $H<R/\alpha$, as well as inside the original
splitting's distortion and coverage domains. If $\alpha=0$, use
the explicitly supplied test domain in the lemma instead.
Distances, coordinate values and the distortion are all rescaled
together; the differential tolerance is dimensionless.

FC16 controls $\eta-\eta(p)$, not an arbitrary translate of $\eta$.
For a cutoff centered at zero, fix $\eta(p)=0$ when choosing the
local witness (or retain and charge its nonzero center value).
Then use that SAME witness in FC01. Recentring an already assembled
coordinate without also changing its cutoff would change the global
map. Translation invariance in KL 4.21 permits this normalization during
the local KL 4.23 construction; general adaptedness by itself does
not specify that center value.

\section{The padded edge strip and actual support buffers}
\label{sec:fibration-support-enclosures}
The edge estimate needs weak-edge control beyond the original
radius-$100\Delta$ coordinate domain. We supply a padded version
of the relevant part of KL 9.7 by a finite metric argument. It does
not assume that the weak-edge set is closed.

\begin{lemma}[Weak-edge points stay near the bottom side; FC17]
\label{lem:fibration-padded-edge}
There is a universal $\tau\in(0,1/100)$ with the following
property. Let $\Delta>0$, $0<\delta\le\tau\Delta$, and let
$(X,p)$ be a metric space with a pointed map
\[
 Q:B(p,200\Delta)\longrightarrow\R\times[0,C],
 \qquad Q(p)=(0,0),\qquad C>200\Delta.
\]
Suppose $Q$ has distortion at most $\delta$ and covers the target
ball of radius $200\Delta-\delta$ to error $\delta$.
Let $E'\subset X$. For EVERY $z\in E'\cap B(p,120\Delta)$,
suppose there is a map
\[
 W_z:B(z,\Delta)\longrightarrow\R\times[0,\infty),
 \qquad W_z(z)=(0,0),
\]
with distortion at most $\delta$. All these distances and errors
use the SAME normalization. Writing $Q(z)=(u_z,h_z)$, one has
$h_z<\Delta/10$ for every such $z$.
\end{lemma}
\begin{proof}
A Euclidean finite-configuration fact is sufficient: for fixed
$t>0$, the four points $\pm te_1,\pm te_2$, together with their
center, cannot be approximated arbitrarily well by a configuration
in the closed upper half-plane whose center is its boundary origin.
If such approximations existed, the four bounded image vectors
would have a convergent subsequence. In the limit each opposite
pair has norms $t,t$ and distance $2t$, hence consists of opposite
vectors. Nonnegative second coordinates force all four vectors
onto the horizontal axis. The cross distances $\sqrt2t$ are then
impossible. Thus a positive error threshold exists. Only compactness
of a bounded subset of $\R^8$ is used here.

Rescale by $\Delta^{-1}$, and suppose $h_z\ge1/10$.
For $\delta$ small, $|Q(z)|<120+\delta$ and $C>200$ leave all
four target points $Q(z)\pm te_1,Q(z)\pm te_2$, with $t=1/100$,
inside the strip and its covered radius-$200-\delta$ ball.
Choose lifts $x_1,\ldots,x_4$. Distortion gives
$d(z,x_a)\le t+2\delta<1$ and controls pairwise lift distances
by the corresponding cross configuration, with error at most
$3\delta$. Apply $W_z$ to these lifts and to $z$; their distances
then have error at most $4\delta$. Choosing $\tau$ below the
finite-configuration threshold divided by four is a contradiction.
All choices and bounds are independent of $C,X,z$ and $\Delta$.
\end{proof}

In the source application normalize by $\rho(p)^{-1}$. The map
$Q$ is the ORIGINAL composition in KL (9.5), with its pointed
endpoint value retained. If $\Delta\Lambda$ is sufficiently small,
then $\rho(z)/\rho(p)\in[1/2,2]$ on $B(p,120\Delta)$.
The weak-edge map at $z$, rescaled in both source and target by
this ratio, supplies $W_z$ on $B(z,\Delta)$: its original tested
radius need only exceed $2\Delta$, and its distortion is multiplied
by at most two. The strong and weak splitting/interval errors may
be chosen small enough to supply the displayed uniform budgets.
This uses the actual endpoint-marked maps in LC84, not an assertion
that every weak point is a strong point. No global Alexandrov
compactness or no-three-stratum theorem is needed for FC17.

\begin{proposition}[Buffered edge and slim supports; FC18]
\label{prop:fibration-support-enclosures}
Use the actual cutoffs of KL (9.19) and (10.2), extended by zero
from their ORIGINAL domains. Let $\Delta\ge1$, let $\rho>0$ be $\Lambda$-Lipschitz in the
physical metric $d_g$, and put $r_j=\rho(p_j)$. In the following
cases, $d=d_g/r_j$ and all balls use $d$; $P$ and $\rho$ retain
physical units.
\begin{enumerate}[label=(\roman*),leftmargin=*]
\item For a slim chart take $\phi=(u,z):B(p_j,10^6\Delta)\to\R\times Z$,
with $\phi(p_j)=(0,z_0)$. Suppose its distortion is at most $\Delta$,
its ORIGINAL factor
has diameter at most $1000\Delta$, and
$|\eta_j-u|\le\Delta$ on $B(p_j,10^6\Delta)$. Then
\[
 \operatorname{supp}\zeta_j
 \subset\overline B(p_j,901002\Delta)
 \subset B(p_j,950000\Delta).
\]
\item For an edge chart at $p_j\in E\subset E'$, suppose
$\Delta\Lambda\le1/10000$. Take FC17's $Q=(u,h)$ and weak maps,
with distortion $\delta\le\min(\tau\Delta,\Delta/100)$.
Suppose $|\eta_j-u|\le\Delta/100$ on $B(p_j,100\Delta)$, and
retain the SAME nonnegative $P=\rho_{E'}$ with
\[
 |P-d_{g,E'}|\le\zeta\rho,\qquad 0\le\zeta\le1/100.
\]
Then
\[
 \operatorname{supp}\zeta_j\subset\overline B(p_j,15\Delta)
                  \subset B(p_j,20\Delta).
\]
\end{enumerate}
With physical distances restored, multiply all displayed radii by
$r_j$. The value hypotheses are supplied by FC16 after the local
center normalization and its explicit test-domain restrictions.
\end{proposition}
\begin{proof}
For a nonzero slim cutoff, $|\eta_j|<900000\Delta$. The product
basepoint is $(0,z_0)$, so the sum bound for the product distance gives
\[
 d(p_j,x)\le |u(x)|+d_Z(z(x),z_0)+\Delta
             \le901002\Delta.
\]
Thus the closure of the nonzero set has the same bound, strictly
inside the original domain. This uses the diameter of the chosen
factor, not compactness or a diameter inference from a single ball.

For a nonzero edge cutoff at $x$ the ACTUAL two factors in (9.19)
give $|\eta_j(x)|<9\Delta$ and $P(x)/\rho(x)<9\Delta$.
Since $x\in B(p_j,100\Delta)$, slow variation gives
$\rho(x)/r_j\le1.01$. In these normalized units,
\[
 d(x,E')\le(9\Delta+1/100)1.01<9.2\Delta.
\]
Choose $z\in E'$ with $d(x,z)<9.2\Delta$; attainment of the
infimum is unnecessary. Then $d(p_j,z)<109.2\Delta<120\Delta$,
so FC17 gives $h(z)<\Delta/10$. Distortion now yields
\[
 |u(x)|\le9.01\Delta<10\Delta,\qquad
 h(x)<\Delta/10+9.2\Delta+\Delta/100<10\Delta.
\]
Consequently $d(p_j,x)\le|Q(x)|+\delta<15\Delta$.
Take the closure of the nonzero set. Its closed radius-$15\Delta$
enclosure is still strictly inside the ORIGINAL radius-$100\Delta$
domain, so no coordinate is evaluated on an artificial extension.
\end{proof}

\paragraph{What is now supplied.}
Under these source-map budgets, the two sufficient enclosures in
FC12 are proved, not merely requested. FC08 gives $r_j/r\ge1/2$
for every support meeting the radius-$10r$ test ball, once its
original slow-variation restrictions are also imposed. FC12 then
gives coordinate-domain margins at least
$(80\Delta-40)r_j$ and $(50000\Delta-40)r_j$ for edge and slim
charts, respectively. Combining the edge enclosure with FC14's
additional bounds $\Delta\ge120L$ and $\Delta\Lambda\le1/100$
also supplies $d(p,E)<30\Delta r$ and hence its SECOND smoothing
domain restriction. No distance-to-$E'$ estimate substitutes for
this strong-edge bound.

The local selection of errors is explicit: fix $\Delta$ and the
chart radius, prescribe the value error and FC17 tolerance, choose
the long-test length, then choose adapted/splitting errors and
$\Lambda$ small enough. These are finitely many local restrictions.
Checking their compatibility with the full simultaneous order of
KL (15.4) remains a later task. Neither FC18 nor its metric proof
promotes the qualified local fiber, smoothing or model producers.

\section{Producing common directional tests from long endpoints}
\label{sec:fibration-common-endpoints}
We next replace FC15's directional premise by explicit metric and
adapted-test data. This separates the differential step from the
remaining Euclidean alignment and original-domain binding.

\begin{lemma}[Common endpoints give the affine comparison; FC19]
\label{lem:fibration-common-endpoints}
Let $M$ be a complete connected Riemannian manifold and
$1\le k\le m$ integers. Fix $L>0$, $T,\alpha,\delta,E\ge0$,
\[
 s>T+2\delta,\qquad H>L+s+3\delta,
 \qquad \phi=(u,z):B(p,H)\to\R^m\times Z,
 \quad \phi(p)=(0,z_0).
\]
The product map has distortion at most $\delta$ and covers the
target radius-$H-\delta$ ball to error $\delta$. Suppose
$\Psi:B(p,H)\to\R^k$, a CONSTANT linear coisometry
$A:\R^m\to\R^k$, and $b_0\in\R^k$ satisfy
\[
 AA^*=I,\qquad |\Psi-Au-b_0|\le E\quad\hbox{on }B(p,H).
\]
Let $F:B(p,L)\to\R^k$, $G:B(p,L)\to\R^m$ be $C^1$, with
$|dF_a|\le1+\alpha$ for every component and $\|DG\|\le1+\alpha$.
For every $x\in B(p,L)$ and $y\in B(p,H)$ with $d(x,y)>T$,
require the original tests, along every minimizing initial unit
vector $w$ from $x$ to $y$:
\[
 \left|DF_x(w)-\frac{\Psi(y)-\Psi(x)}{d(x,y)}\right|\le\alpha,
 \qquad
 \left|DG_x(w)-\frac{u(y)-u(x)}{d(x,y)}\right|\le\alpha.
\]
If $\epsilon=\alpha+(3\delta+2E)/(s-2\delta)\le1$, then with
$b=F(p)-AG(p)$,
\[
 \|F-(A\,G+b)\|_{C^1(B(p,L))}
       \le2\max\{1,L\}\sqrt{k(4\epsilon+\epsilon^2)}.
\]
\end{lemma}
\begin{proof}
For each $x,a$ put $v_a=A^*e_a$, which is a unit vector.
The target $(u(x)+sv_a,z(x))$ has distance at most
$L+\delta+s<H-\delta$ from the product basepoint. Choose
$y=y_{x,a}$ whose image is within $\delta$ of that target.
Distortion gives
\[
 s-2\delta\le\ell=d(x,y)\le s+2\delta,
 \qquad\ell>T.
\]
The same point and minimizing direction can therefore be used in
BOTH adapted tests. Since $Av_a=e_a$ and $\|A\|=1$,
\[
 (A(u(y)-u(x)))_a\ge s-\delta,\qquad
 \Psi_a(y)-\Psi_a(x)\ge s-\delta-2E.
\]
After division by $\ell$, both ratios are at least
$1-(3\delta+2E)/(s-2\delta)$. The two differential tests give
$dF_a(w),d(AG)_a(w)\ge1-\epsilon$. Apply FC15.
The choices need no continuity and no mutual orthogonality of
the resulting tangent directions.
\end{proof}

\paragraph{Binding the actual overlaps.}
All data in FC19 must use the SAME metric and original coordinates.
With metric $r^{-2}g$, constant-radius coordinates use
$F=(R_j/r)\eta_j$, $G=\eta_i$; the original adapted-test separation
thresholds must both be below $T$, and every test point in the
radius-$H$ ball must lie in BOTH original tested domains. A
coordinate's smooth domain and its much larger derivative-test
domain are different requirements. FC18 supplies the former for
edge/slim overlaps; it does not silently supply the latter.
For the zero shell retain FC13 and LC73's original radial function.
For the high weak-edge branch use $F=P/r$ and convert by FC11
before FC10; source9.12's generalized-gradient tests still need
their original-domain binding.

AC79 and AC76 provide the recentering and metric compatibility
route, with LC20's higher-rank exclusion. The Euclidean factor
approximation in that compatibility must still be aligned with a
SINGLE isometry to obtain $A,b_0,E$ uniformly on the enlarged
radius-$H$ domain, as in KL 4.31. FC19 does not assume a $C^1$
comparison or infer it from value closeness alone: it proves the
implication using the supplied long-endpoint derivative tests.
No new curvature premise is needed by FC16--FC19 once those metric,
calculus and adapted-coordinate data have been supplied.

\paragraph{Current implementation boundary.}
FC16--FC18 give written conditional proofs of the actual edge/slim
support buffers, including the strong-edge proximity required by
FC14. FC19 produces common directional tests from explicit aligned
metric coordinates and common test domains. Remaining tasks are
Euclidean-factor alignment, the full rescaled test-domain binding
(including radial and weak-edge smoothing), and the assembly of
KL 12.17's local graph. FC07 remains source-qualified. Its uniform
graph derivative modulus, two-sided image coverage and Chapter3
pointed-Hausdorff translation remain separate, followed by corrected
20.2 smoothing,13.2 adjustment and the full parameter order.
L142--L144 and \code{reference_checks_revision100.md} record the
source comparisons and these project proof expansions.


\section{Alignment on the required domain}
\label{sec:fibration-alignment-completion}
The remainder of this chapter completes its BLUEPRINT COVERAGE.
It specifies every stage through static graph assembly, with written
conditional arguments and explicitly source-qualified producers.
It does not claim Lean verification or unconditional proof closure.
Earlier progress paragraphs describe their original revision; the
current disposition is the audit at the end of this chapter.

\begin{lemma}[Euclidean anchor alignment; FC20]
\label{lem:fibration-euclidean-alignment}
Let $n\ge1$, $a>0$, and $0\le\delta\le a/(20n)$.
Suppose $f:B(0,2a)\subset\R^n\to\R^n$, $f(0)=0$, satisfies
\[
 \bigl||f(v)-f(w)|-|v-w|\bigr|\le\delta.
\]
There is one orthogonal linear map $Q$ such that
$|f(v)-Qv|\le24n\delta$ for $|v|\le a$.
No continuity or differentiability of $f$ is required.
\end{lemma}
\begin{proof}
Let $V e_i=f(ae_i)/a$ and $G=V^*V$. Polarization of the
squared-distance inequalities gives
$|G_{ij}-\delta_{ij}|\le5\delta/a$. Thus
$\|G-I\|\le\eta=5n\delta/a\le1/4$.
Set $Q=VG^{-1/2}$. The spectral theorem gives
$Q^*Q=I$, $\|V-Q\|\le\eta$ and $\|V^{-*}\|\le2$.
Polarization with $v,0,ae_i$, using $|v|\le a$, also gives
$|(V^*f(v)-v)_i|\le6\delta$. Consequently
\[
 |f(v)-Qv|\le
  2|V^*f(v)-v|+2\|V-Q\|\,|v|
  \le(12\sqrt n+10n)\delta\le24n\delta.
\]
All linear algebra here is an inherited reuse target.
\end{proof}

\begin{lemma}[Raw factor alignment; FC21]
\label{lem:fibration-raw-alignment}
Suppose approximate compatibility supplies, on a specified source
set $D$, raw coordinates $u:D\to\R^m$, $\Psi:D\to\R^k$,
a constant coisometry $P:\R^m\to\R^k$, and
\[
 |\Psi(x)-f(Pu(x)-c)-d|\le e_c,
 \qquad |Pu(x)-c|\le a.
\]
Assume the translated factor map $f$ satisfies FC20 with $n=k$.
Then $A=QP$, $b_0=d-Qc$ satisfy
\[
 AA^*=I,\qquad |\Psi-Au-b_0|\le e_c+24k\delta\quad\hbox{on }D.
\]
\end{lemma}
\begin{proof}
Substitute FC20 and use $PP^*=I$. The translations need not be
small. If the source factor map does not send zero to zero, subtract
its value there and include it in $d$ first.
\end{proof}
This is the explicit Euclidean step used in KL 4.31. AC76 supplies
approximate compatibility; AC79 supplies the original recentering.
Their curvature, exclusion, radius and rescaling hypotheses remain
required. The estimate is uniform on the chosen set $D$, not on an
unrecorded larger ball.

\begin{lemma}[Asymmetric long and short tests; FC22]
\label{lem:fibration-asymmetric-tests}
Let $1\le k\le m$ be integers and $L>0$.
Let $F:D\to\R^k$ and $G:D\to\R^m$ be $C^1$ on
$D=B(p,L)$ in a complete connected Riemannian manifold. Suppose
$|dF_a|\le1+\alpha_F$, $\|DG\|\le1+\alpha_G$.
For every $x\in D$ and component $a$ supply a minimizing unit-speed
segment $\gamma:[0,\ell]\to M$, $\gamma(0)=x$, with
$\ell\ge\ell_0>t>0$. Put $y=\gamma(\ell)$, $z=\gamma(t)$,
$w=\dot\gamma(0)$. Require raw coordinates $\Psi,u$ on their
respective ORIGINAL tested domains, with
\[
 \Psi_a(y)-\Psi_a(x)\ge\ell-\beta,\qquad
 \Psi_a(y)-\Psi_a(z)\le\ell-t+\delta,
\]
\[
 \left|dF_a(w)-\frac{\Psi_a(y)-\Psi_a(x)}\ell\right|\le\alpha_F,
 \qquad
 \left|DG(w)-\frac{u(z)-u(x)}t\right|\le\alpha_G.
\]
Here $\alpha_F,\alpha_G,\beta,\delta,E\ge0$ are uniform.
A constant coisometry $A$ and translation $b_0$ need satisfy
$|\Psi-Au-b_0|\le E$ only at $x,z$, not at $y$.
If
\[
 \epsilon=\max\left\{\alpha_F+\frac\beta{\ell_0},
           \alpha_G+\frac{\beta+\delta+2E}{t}\right\}\le1,
\]
then FC15 gives, with $b=F(p)-AG(p)$,
\[
 \|F-(AG+b)\|_{C^1(D)}
       \le2\max(1,L)\sqrt{k(4\epsilon+\epsilon^2)}.
\]
\end{lemma}
\begin{proof}
Subtract the two raw scalar inequalities to obtain
$\Psi_a(z)-\Psi_a(x)\ge t-\beta-\delta$.
The short alignment gives
$(A(u(z)-u(x)))_a\ge t-\beta-\delta-2E$.
Both original tests therefore saturate on the SAME initial vector
$w$: $dF_a(w),d(AG)_a(w)\ge1-\epsilon$.
Since $\|A\|=1$, the norm hypotheses of FC15 hold.
\end{proof}
A lift of a distant positive axis point of the $F$ product model
supplies $\ell\in[s-2\delta,s+2\delta]$ and $\beta=3\delta$.
Its scalar factor is coarsely $1$-Lipschitz, supplying the second
inequality along the segment. Choose $s>t+2\delta$ above the
$F$ test threshold and $t$ above the $G$ threshold. Only the short
segment must remain in $G$'s original tested domain. FC19 remains
a valid special route when both long domains are actually available.
It is not the default for slim/two-stratum overlaps: KL chooses
$\beta_2$ before $\Delta$, with $\Delta>\beta_2^{-1}$, so one
cannot demand a two-stratum test ball larger than $10^6\Delta$.

\begin{foundation}[Original overlap binding packet; FC23]
\label{found:fibration-original-binding}
For each local model used in FC07 or its projected analogues,
produce a FINITE collection of constant affine comparisons on the
same buffered source domain. Each comparison must provide:
\begin{enumerate}[leftmargin=*]
\item The original metric $r^{-2}g$, exact coordinate and center
normalization, active support set and count, and the smooth domain.
\item Separate original derivative-test domains, separation
thresholds, a permitted short length $t$, permitted long lengths
$s$, and every inclusion needed by FC22. Require $\ell_0>t$;
changing the quality parameter does not enlarge an already chosen
map's domain.
\item The AC76/AC79 compatibility data on the SHORT buffer,
FC21's single coisometry, and explicit error budgets small enough
for FC22. For components arising from radial smoothing, provide
the same directional inequalities directly if a product raw map
is not available.
\item Uniform $C^1$ comparison errors for ALL coordinates entering
the joint cutoff, even if their output block vanishes. Freeze the
scale by FC11 before applying FC10.
\end{enumerate}
The binding is source-qualified by KL 4.21--4.35, 9.12,
11.1 and 12.12. It is an implementation obligation, not a claim
that the common-domain conditions follow just from support overlap.
\end{foundation}
The binding matrix is as follows.
\begin{center}
\begin{tabular}{p{.16\textwidth}p{.34\textwidth}p{.40\textwidth}}
Coordinate & Smooth buffer already specified & Derivative test to retain \\\hline
Two-stratum & Original centered chart; FC12 & Original adapted test; short $t$ before $\Delta$ \\
Edge/slim & FC18 and FC12, with rescaled radii & Own long axis test; shorten before comparing to a smaller chart \\
Zero radial & FC13's original shell and LC73 & Original radial directional estimate, $R_0/r$ retained; no bounded translation assumption \\
Weak edge & FC14 low/high alternative & Low branch omits the argument exactly; high branch retains BOTH smoothing restrictions and the $P/r$ gradient estimate
\end{tabular}\end{center}
This packet closes the specification of the overlap interface.
Its geometric source binding remains qualified; it is not promoted
by the elementary FC20--FC22 proofs.

\section{Local graph models and both coverage directions}
\label{sec:fibration-graph-completion}
\begin{lemma}[A dimension-independent graph modulus; FC24]
\label{lem:fibration-graph-modulus}
On a convex Euclidean domain let a model be a finite orthogonal
sum of blocks $b_j\circ\ell_j$, where $\ell_j$ is affine,
$\|D\ell_j\|\le L_j$, and
$\|Db_j\|\le A_j$, $\|D^2b_j\|\le B_j$ globally.
Suppose at most $N$ blocks have nonzero derivatives at each point
and $A_jL_j\le A$, $B_jL_j^2\le B$.
Then its derivative and Hessian norms are at most
$\sqrt N A$ and $\sqrt N B$, respectively. Its derivative is
$\sqrt N B$-Lipschitz on every contained line segment.
\end{lemma}
\begin{proof}
For unit vectors square and sum the block bounds in the orthogonal
target. Integrate the Hessian on the segment. The identity coordinate
of a graph contributes $1$ to the squared first-derivative bound
and zero to the Hessian. No estimate uses the total number of blocks.
\end{proof}
Build the KL 12.15 model by replacing each normalized coordinate
by its FC23 affine comparison, using the ACTUAL cutoff formulas,
and replacing $\rho/r$ by $1$. FC05 bounds zero blocks even when
$R_0/r$ is arbitrarily large; their Hessians improve as this ratio
grows. FC10 bounds the entire edge network as one block depending
on finitely many affine variables. FC11 controls the replacement
error, including $\Delta\Lambda$. The affine-model active count
must be verified on its domain as well as on the carrier; using
the finite list of supports meeting the original comparison ball
supplies such a count. These statements yield a graph $\Phi(u)=(u,h(u))$
with a uniform Hessian bound, and a $C^1$ approximation to the
normalized image, once FC23's input budgets are delivered.

\begin{lemma}[Graph coverage with identical truncation radii; FC25]
\label{lem:fibration-truncated-coverage}
Let $V$ be a finite-dimensional Euclidean space, $R>0$, $B\ge0$,
and $X\subset\R^k\oplus V$ contain $x$ with first coordinate
$u_0$, and let $\Phi(u)=(u,h(u))$ be $C^2$ on a neighborhood
of $\overline B(u_0,R)$, with $\|D^2\Phi\|\le B$ there.
For $e\ge0$ assume:
\begin{enumerate}[leftmargin=*]
\item Every $x'\in X\cap\overline B(x,R)$ has first coordinate
$u$ and $|x'-\Phi(u)|\le e$; in particular
$|x-\Phi(u_0)|\le e$.
\item For EVERY $u\in\overline B(u_0,R)$, some $x_u\in X$
has first coordinate $u$ and $|x_u-\Phi(u)|\le e$.
\end{enumerate}
Set $T=D\Phi(u_0)$, $A=x+\operatorname{im}T$,
$q=2e+BR^2/2$. Then the Hausdorff distance between
$X\cap\overline B(x,R)$ and $A\cap\overline B(x,R)$ is at most
$3q$, using ambient distances.
\end{lemma}
\begin{proof}
For the first direction $|u-u_0|\le R$. Taylor's estimate gives
$|x'-x-T(u-u_0)|\le q$. If the candidate in $A$ lies outside
the radius-$R$ ball, radial contraction about $x$ moves it by at
most $q$, so the distance is at most $2q$.
For the converse, if $R\le2q$ the common point $x$ suffices.
Otherwise move any $a\in A\cap\overline B(x,R)$ radially to
$a'\in\overline B(x,R-2q)$, by distance at most $2q$.
Write $a'-x=Tv$. Since $T$ has identity first component,
$|v|\le|Tv|\le R$. The supplied point $x_{u_0+v}$ lies within
$q$ of $a'$, hence inside the radius-$R$ ball and within $3q$
of $a$. This proves both directions, including the boundary.
\end{proof}
For the first image FC03 localizes ALL preimages in the test ball;
LC83's actual bundle coordinate coverage supplies the second
hypothesis. It cannot be replaced by approximate metric density
without accounting for the extra error. In $R_i=1$ units let
$r\in[\Sigma/2,2\Sigma]$ and test $R=r/\Gamma$.
The sufficient restrictions
\[
 e\le\Gamma\Sigma/24,\qquad B\Sigma\le\Gamma^3/6
\]
give $3q/r\le\Gamma$. Also require $R$ smaller than the chart
buffer. This is an explicit fixed-error route to Chapter 3's
pointed convention; merely intersecting two close sets with the
same ball would not prove it. Strict source inequalities follow
by taking slack. A different open-ball convention needs the
corresponding margin or limiting argument, not an implicit change.

\begin{lemma}[Radii after discarding the scale coordinate; FC26]
\label{lem:fibration-projected-radii}
Let $X$ be an image retaining the entire constant-radius blocks.
Choose a preimage $p_x$ for each $x\in X$, and a block $i_x$
whose marker equals $R_{i_x}>0$ at $x$. Suppose for ANY preimage
$q$ of ANY image point with positive $i$ marker,
\[
 3R_i/4\le\rho(q)\le5R_i/4.
\]
For $0<\Sigma\le1/2$ define $r(x)=\Sigma\rho(p_x)$.
Then, independently of the choices,
\[
 |r(y)-r(x)|\le2\bigl(|x-y|+r(x)\bigr).
\]
\end{lemma}
\begin{proof}
Put $D=|x-y|$ and $R=\max(R_{i_x},R_{i_y})$.
If $D\ge R/2$, both selected scale values are at most $5R/4$,
so their difference times $\Sigma$ is at most
$(5/2)\Sigma D\le2D$.
Otherwise the full-marker block of radius $R$ has positive marker
at BOTH points, since marker differences are at most $D<R/2$.
The two selected scales are in $[3R/4,5R/4]$, so their ratio is
between $3/5$ and $5/3$. Their difference is at most
$2\rho(p_x)/3$, which suffices.
\end{proof}
FC18 and $\Lambda C\le1/4$ provide the scale hypothesis for edge
and slim support radius $CR_i$. A choice is allowed: KL 20.1 does
not require continuity of $r$. The first image still uses the exact
FC04 coordinate; $Q_2,Q_3$ have discarded it. Positivity and bounds
away from zero follow on the compact carrier. No unproved descent
of $\rho$ through a projected image is asserted.

\begin{foundation}[Projected edge and slim clouds; FC27]
\label{found:fibration-projected-clouds}
For $j=2$ use the edge core/enlargement thresholds $7\Delta,8\Delta$
for BOTH $|\eta_i|$ and $\eta_{E'}$, with $i\in I_e$.
For $j=3$ use the slim thresholds $7\cdot10^5\Delta,8\cdot10^5\Delta$,
with $i\in I_s$. Apply $\pi_j\mathcal E^0$ to obtain
$S_j\subset\widetilde S_j\subset Q_j$, and use FC26's selected
radius $r_j=\Sigma_j\rho(p_x)$.
KL 12.27 and 12.34 are source-qualified producers of
$(2,\Gamma_j)$ cloudy ONE-manifolds, local graph models, normal
errors less than $\Gamma_j$, and surjective projected derivatives
with two positive singular-value bounds.
\end{foundation}
Their proofs are omitted in the source. The implementation must
repeat FC23--FC25 and FC06 in dimension one, localizing the entire
projected image by a retained marker and supplying the appropriate
coordinate coverage from LC84/LC85. The excluded two-stratum blocks
must not be reintroduced. KL (12.32) prints $I_{\rm edge}$; its
section title, thresholds, conclusion and (12.38) identify $I_s$
as the required family. This is a recorded project correction,
not an author-listed erratum. FC07 and FC27 remain source-qualified
until these complete geometric bindings are checked.

\section{Smoothing cloudy images}
\label{sec:fibration-cloud-smoothing}
\begin{foundation}[Historical cloudy smoothing claim; FC28]
\label{found:fibration-cloud-smoothing}
Fix integers $k,K\ge1$, $0<C<\infty$, and $\varepsilon>0$.
KL20.2 claims $\delta=\delta(k,K,C,\varepsilon)>0$,
INDEPENDENT of $\dim H$. For a $(C,\delta)$ cloudy packet
$S\subset\widetilde S\subset H$ with bounded $\widetilde S$ and radius bounded above and away from zero (the compact-carrier application), the claimed output is a
smooth embedded $k$-manifold $W$ and a smooth nearest-point
submersion $P:N_r(S)\to W$, where
$N_r(S)=\bigcup_{x\in S}B(x,r(x))$.
The claimed consumer properties are listed for comparison. CFS06 below
refutes this implication at its stated generality; an application-specific
replacement must be proved before any consumer may use it:
\begin{enumerate}[leftmargin=*]
\item At each $x\in S$, the recentered, $r(x)^{-1}$-rescaled
ENLARGED set and $W$ are $\varepsilon$-close in the fixed pointed
Hausdorff convention. $W\subset N_{\varepsilon r}(\widetilde S)$.
\item $W\cap N_r(S)$ is properly embedded in $N_r(S)$.
On $B(x,r(x))\cap W$ its normal projections are within
$\varepsilon$ of the normal projection to $A_x$.
\item For $1\le q\le K$, on $B(x,r(x))$,
$\|D^q(P-P_{A_x})\|\le\varepsilon r(x)^{1-q}$.
In particular $\|DP-\pi_{A_x^0}\|\le\varepsilon$.
\end{enumerate}
The source also states curvature-derivative estimates. They are
not needed by the $C^1$ adjustment here and are not inferred from
the listed first-derivative conclusions.
\end{foundation}
The author correction to the convergence proof uses
$(\widetilde S_j-x_j)/r(x_j)$ based at ZERO, not the smaller set
or a moving ambient basepoint. KL19.1 must select centers in $S$, where
model planes are supplied; the earlier enlarged-set-cover sentence
is corrected here. Moreover, CFS01 below disproves the unshrunk-cover
multiplicity inference in the source. CFS02 repairs it using smaller
radii and the cloud's $k$-dimensional approximation, not ambient volume. Average the local normal
projections with smooth cutoffs; their spectral gap defines a smooth
normal bundle. The weighted normal-displacement section must be
uniformly transverse to zero. Its zero set gives $W$. Separate
proof targets then give its local coverage, properness, uniqueness
of nearest points, and derivative estimates for $P$.
A spectral gap alone does not prove the last three assertions.

\begin{lemma}[The relevant finite span; FC29]
\label{lem:fibration-active-span}
For $N$ affine $k$-planes $x_i+L_i\subset H$, let
$V=\operatorname{span}\{x_i,L_i:1\le i\le N\}$.
Then $\dim V\le N(k+1)$. If $P_i$ is normal projection to $L_i$
and $O=\sum_i w_iP_i$, where $w_i\ge0$, $\sum_iw_i=1$,
then $O$ preserves $V\oplus V^\perp$, is the identity on
$V^\perp$, and $\operatorname{rank}(I-O)\le Nk$.
If a spectral projection $Q$ of $O$ selects its eigenvalues near
one and the other eigenvalues lie near zero, then $Q$ is also the
identity on $V^\perp$. Consequently
\[
 \operatorname{proj}_{V^\perp}
       \sum_iw_i Q(v-x_i)=\operatorname{proj}_{V^\perp}v.
\]
\end{lemma}
\begin{proof}
$P_i=I-\pi_{L_i}$ and $L_i\subset V$. Sum these identities.
Spectral calculus respects the orthogonal direct sum. The last
identity follows from $x_i\in V$ and $\sum_iw_i=1$.
\end{proof}
The source's bounded-rank assertion for the averaged NORMAL
projection is not literally valid as ambient dimension grows;
its complement has bounded rank. FC29 gives the required local
finite span for the zero set. Extending this reduction to a uniform
neighborhood, all derivative directions and the projection estimates
is part of FC28's source-qualified smoothing implementation.
CFS02--CFS05 below now supply the uniform-neighborhood and ambient-
derivative part for a smaller tube; the full FC28 projection conclusion
remains open. Neither pointwise multiplicity nor this linear lemma alone
proves uniform smoothing. This distinction also repairs the earlier
12.7 temptation to bound the total dimension of $H$.

\subsection{A corrected cover and local smoothing construction}
\label{sec:fibration-cloud-local-repair}
Revision121 develops the first geometric part of FC28. The following
five arguments refine that one parent contract; they do not add parent
inventory rows or accept its full output. The source is KL Definition3.1,
Lemma19.1 and Definition20.1--Lemma20.2 (printed21,94--97;
PDF16,89--92), with the May15,2015 author correction. Two additional
issues below are project findings, not author-listed corrections.

\paragraph{Precise cloud convention.}
Use the radius inequality stated after FC07, and require the actual
truncated-set bounds
\[
 d_H\bigl(\widetilde S\cap B(x,r(x)/\delta),
                A_x\cap B(x,r(x)/\delta)\bigr)\le\delta r(x)
                 \qquad(x\in S).                                      \tag{CS}
\]
Here $A_x=x+L_x$, $\dim L_x=k$, and distance-to-set bounds define $d_H$.
Either open balls throughout or closed balls throughout suffice for
these five arguments: every test below has strict interior slack.
No closedness or attained nearest point of $\widetilde S$ is needed;
witnesses may be chosen with error $2\delta r(x)$. This is the
explicit fixed-error input, not an unproved convention adapter.
Assume $S$ nonempty and bounded and $r$ bounded above and away from
zero on $S$, as follows from the retained compact-carrier assumptions.
Planes are available at centers in $S$, not at arbitrary points of
$\widetilde S$. Covering $\widetilde S$ does not supply those planes.

\begin{example}[The unshrunk overlap claim fails; CFS01]
\label{ex:fibration-cloud-overlap-counterexample}
For fixed $k=1$, $C=1$ and any $0<\delta<1$, there are bounded
cloud packets satisfying (CS) with zero error, for which the
KL19.1 selection using radii $r$ has arbitrarily many cutoffs
$\phi(|v-x|/(10r(x)))$ equal to one at the same point of $N_r(S)$.
The count of selected centers in $B(x,5r(x))$ can also be arbitrarily
large. Thus the asserted bounds on printed96 do not follow from
Definition20.1, even when $r$ is Lipschitz and $\dim H=2$.
\end{example}
\begin{proof}
Identify the horizontal line in $H=\R^2$ with $\R$. For $n\ge1$ put
\[
 S_n=\{0\}\cup\{2^{-i}:0\le i<n\},\qquad
 \widetilde S=[-R,R],\qquad R=2+\delta^{-1},
\]
\[
 e_n=2^{-n}/100,\qquad r(t)=\max\{e_n,|t|/4\}.
\]
This radius is positive, bounded and $1/4$-Lipschitz. It satisfies
the radius inequality with $C=1$. At every $x\in S_n$ the whole
ball tested in (CS) lies inside the interval, so $A_x$ can be the
horizontal line and the error is zero.

The ambient balls $B(x,r(x))$, $x\in S_n$, are pairwise disjoint.
For positive centers $a>b$ one has $b\le a/2$, whence
$a-b\ge a/2>3a/8\ge(a+b)/4$. The ball at zero is disjoint
as well, since $e_n<a/100$ for every positive center. The greedy
selection can choose the centers in decreasing order of radius;
none is removed by an earlier ball, so it selects all $n+1$.
At $v=0$, each positive center has $|v-x|/(10r(x))=2/5$,
and the zero center has argument zero. All cutoffs equal one.
Also every center lies in $B(1,5r(1))=B(1,5/4)$.
These claims hold with fixed $\delta,C,k$ while $n$ grows.
\end{proof}
This disproves the stated packing and multiplicity inference, not
Lemma20.2's smoothing conclusion: this cloud already lies on a line.
The earlier FC29 rank correction alone cannot repair this new defect.
CFS06 below gives a different counterexample to the general theorem
itself; CFS01 alone established only failure of the proof inference.

\begin{lemma}[A small-scale cover with neighborhood multiplicity; CFS02]
\label{lem:fibration-cloud-small-cover}
Put
\[
 \lambda=\frac1{100(C+1)},\qquad A=2(C+1),\qquad D=20A+6,
 \qquad N=\left\lceil(1+2AD)^k\right\rceil.
\]
If $0<\delta\le\lambda/(2A)$, a finite set $T\subset S$ can be
chosen so that the balls $B(x,\lambda r(x))$, $x\in T$, are
pairwise disjoint and
\[
 N_{\lambda r}(S)\subset U:=\bigcup_{x\in T}B(x,5\lambda r(x)).
\]
For any $v\in B(x_0,5\lambda r_0)$, where $x_0\in T$ and
$r_0=r(x_0)$, let $J$ consist of all $x_i\in T$ whose closed
$20\lambda r_i$-ball meets $B(v,\lambda r_0)$. Then
\[
 A^{-1}r_0\le r_i\le Ar_0,\qquad
 |x_i-x_0|\le D\lambda r_0,\qquad |J|\le N.                 \tag{NC}
\]
All constants are independent of $\dim H$, the size of $T$, and
$\inf r/\sup r$.
\end{lemma}
\begin{proof}
Apply the finite greedy selection of KL19.1 to $S$ and radius
$\lambda r$. Its finiteness uses bounded $S$, positive $\inf r$
and finite-dimensional $H$; no uniform count comes from that step.
It gives, for every $x\in S$, a selected $x_i$ with
$|x-x_i|<3\lambda r_i$ and $r(x)\le2r_i$. This proves the
inclusion of tubes.

If $i\in J$, the triangle inequality gives
$|x_i-x_0|\le20\lambda r_i+6\lambda r_0$. Apply the radius
inequality in both orders to obtain
\[
 (1-20C\lambda)r_i\le(C+1+6C\lambda)r_0,
\]
\[
 (1-6C\lambda)r_0\le(C+1+20C\lambda)r_i.
\]
Since $C\lambda<1/100$, these imply the radius bounds in (NC),
and then the center bound. Notice that this argument concerns
\emph{all supports meeting a whole neighborhood}, not merely
cutoffs nonzero at $v$.

One has $D\lambda<1$ and $\delta<1$, so every such center is
strictly inside the test ball at $x_0$. By (CS), its distance to
$A_{x_0}$ is at most $\delta r_0$. Orthogonally project the centers
to $A_{x_0}$. Disjointness of the selected open balls gives original
separation at least $2\lambda r_0/A$; projection decreases this by
at most $2\delta r_0$, leaving separation at least $\lambda r_0/A$.
The projected centers lie in the $k$-ball of radius $D\lambda r_0$.
Disjoint open $k$-balls of radius $\lambda r_0/(2A)$ around them
lie in the ball with that radius added. Comparing \emph{$k$-dimensional}
Euclidean volumes gives $|J|\le(1+2AD)^k$. No ambient-dimensional
volume comparison is used.
\end{proof}

\begin{lemma}[Coherence of every plane in that neighborhood; CFS03]
\label{lem:fibration-cloud-plane-coherence}
In CFS02, after reducing $\delta$ as a function of $C$ only, let
$P_i$ and $P_0$ be the normal projectors to $L_{x_i}$ and $L_{x_0}$.
For every $i\in J$,
\[
 |P_0(x_i-x_0)|\le\delta r_0,\qquad
 \|P_i-P_0\|\le6(A+1)\delta.                              \tag{PC}
\]
These bounds compare the actual chosen cloud planes and use the
ENLARGED set in both directions of (CS).
\end{lemma}
\begin{proof}
The first assertion was proved above. For any unit $u\in L_{x_i}$,
test $y=x_i+r_0u$ in the plane at $x_i$. For
$\delta<1/(4A)$ it is strictly inside that plane's tested ball,
and there is $z\in\widetilde S$ with
$|z-y|<2\delta r_i\le2A\delta r_0$. Since $D\lambda<1/2$,
$|z-x_0|<2r_0$. Shrink $\delta$ below $1/4$; the set-to-plane
test at $x_0$ now gives $\operatorname{dist}(z,A_{x_0})\le\delta r_0$.
Subtracting the normal offset of $x_i$ shows
\[
 |P_0u|\le(2A+2)\delta=:b.
\]
Take $b<1/2$. Projection $L_{x_i}\to L_{x_0}$ has lower norm
$\sqrt{1-b^2}$ and is an isomorphism, since both dimensions are $k$.
For unit $w\in L_{x_0}$ its inverse has normal part of norm at
most $b/\sqrt{1-b^2}\le2b$, so the distance of $w$ to $L_{x_i}$
is at most $2b$. Write the difference of tangent projections as
\[
 (I-\pi_{L_{x_0}})\pi_{L_{x_i}}
       -\pi_{L_{x_0}}(I-\pi_{L_{x_i}}).
\]
The first operator has norm at most $b$ and the second at most
$2b$, by taking its adjoint. Normal projectors have the same
difference norm. This gives $3b=6(A+1)\delta$.
\end{proof}

\begin{lemma}[Uniform spectral and displacement estimates; CFS04]
\label{lem:fibration-cloud-neighborhood-spectral}
Fix a nonnegative smooth cutoff $\phi$ equal to one on $[0,1]$
and zero on $[2,\infty)$. On the buffered open set
$\Omega=\bigcup_{x_i\in T}B(x_i,6\lambda r_i)$ define
\[
 \phi_i(z)=\phi\bigl(|z-x_i|/(10\lambda r_i)\bigr),\quad
 w_i=\phi_i/\sum_j\phi_j,\quad O(z)=\sum_iw_i(z)P_i.
\]
For sufficiently small $\delta=\delta(k,C)>0$ the spectral
projector $Q(z)$ of $O(z)$ for the cluster near one is smooth
and has rank $\dim H-k$. Set
\[
 \mu(z)=\sum_iw_i(z)x_i,\qquad \eta(z)=Q(z)(z-\mu(z)).
\]
For every integer $m\ge1$ there are finite constants $B_q,E_q$,
$0\le q\le m$, depending only on $k,C,m,\phi$, such that on
each neighborhood $B(v,\lambda r_0)$ in CFS02,
\[
 \|D^qw_i\|\le B_qr_0^{-q},\qquad
 \|D^q(Q-P_0)\|\le E_q\delta r_0^{-q},                    \tag{SE}
\]
\[
 \bigl\|D^q\bigl(\eta-P_0(\,\cdot-x_0)\bigr)\bigr\|
                  \le E_q\delta r_0^{1-q}.               \tag{DE}
\]
The norms are multilinear operator norms in \emph{all} ambient
directions. On the whole neighborhood the single space
\[
 V=\operatorname{span}\{x_i-x_0,L_{x_i}:i\in J\},\qquad
 \dim V\le N(k+1),
\]
satisfies $Q|_{V^\perp}=I$ and
$\operatorname{proj}_{V^\perp}\eta(z)
 =\operatorname{proj}_{V^\perp}(z-x_0)$.
\end{lemma}
\begin{proof}
Every neighborhood in CFS02 lies in $\Omega$, and these
neighborhoods cover $\Omega$. The reference cutoff is one throughout
each neighborhood, since $|z-x_0|<6\lambda r_0$. Thus its denominator is at least one.
All other cutoffs outside $J$ vanish identically there. Radial
cutoffs are smooth at their centers because they are constant
near zero. Their derivatives of order $q$ are bounded by constants
times $(\lambda r_i)^{-q}$: differentiate away from the constant
region, where $|z-x_i|\ge10\lambda r_i$. Derivatives of the norm
have dimension-independent multilinear bounds. The product and
reciprocal rules, (NC), and $|J|\le N$ give the first estimate.

Since $\sum_iw_i=1$, every derivative of $O-P_0$ is a sum of
$D^qw_i(P_i-P_0)$. CFS03 gives bounds $c_q\delta r_0^{-q}$,
including $q=0$. Choose $\delta$ so $\|O-P_0\|<1/4$.
Self-adjointness puts the spectrum into disjoint $1/4$-neighborhoods
of zero and one. The near-zero cluster has dimension $k$, for
example by continuous deformation from $P_0$ without closing the gap.

In the complexification, use the circle $|\zeta-1|=1/2$ to write
\[
 Q=\frac1{2\pi\mathrm i}\int_{|\zeta-1|=1/2}
                     (\zeta I-O)^{-1}\,d\zeta.
\]
The resolvent norm is at most four on this contour. The resolvent
identity compares $Q$ with $P_0$ at order zero. Differentiating an
inverse expresses each higher derivative as a finite sum of products
of resolvents and positive-order derivatives of $O$. Each such
product contains a factor $\delta$; for $0<\delta\le1$ the extra
powers can be discarded. Operator norms of products and integrals
therefore give (SE) independently of $\dim H$. This is a use of
inherited finite-dimensional spectral calculus, not a bounded-rank
claim about $O$.

Translate $x_0$ to zero and rescale $r_0$ to one. The centers in $J$
have norm at most $D\lambda$, and $|z|<6\lambda$. Write
\[
 \eta(z)-P_0z=(Q(z)-P_0)(z-\mu(z))-P_0\mu(z).
\]
The derivative bounds for weights bound $\mu$ in every fixed finite
order. Each $P_0x_i$ has norm at most $\delta$, so the same rules
bound every derivative of $P_0\mu$ by a constant times $\delta$.
The product rule proves (DE); rescaling gives its stated powers.
Finally FC29 applies to this \emph{fixed neighborhood list}, after
translation. It proves the $V^\perp$ identities everywhere there.
The weights can depend on the $V^\perp$ coordinate; we have not
silently treated the entire construction as constant in those
directions. Their derivatives were estimated in the full space.
\end{proof}

\begin{proposition}[A local smooth interpolant on the smaller tube; CFS05]
\label{prop:fibration-cloud-local-zero-set}
For sufficiently small $\delta$ depending on $k,C$ and the fixed
cutoff $\phi$, the zero set
\[
 W_0=\{z\in U:\eta(z)=0\}
\]
is a smooth embedded $k$-manifold, closed in $U$ and hence properly
embedded there. For every selected center $x_i$, in the orthogonal
cylinder
\[
 z=x_i+t+n,\quad t\in L_{x_i},\quad n\in L_{x_i}^{\perp},
 \quad |t|<\lambda r_i/4,\quad |n|<\lambda r_i/4,
\]
it is exactly one graph $n=g_i(t)$. For each fixed $m\ge1$,
\[
 \|D^qg_i\|\le F_q\delta r_i^{1-q}\quad(0\le q\le m),
\]
with constants depending only on $k,C,m,\phi$, after the required
smallness choice. This conclusion does not assert that $W_0$ has
a nearest-point projection on the larger tube $N_r(S)$.
\end{proposition}
\begin{proof}
At any $z\in U$, choose a neighborhood and reference $x_0$ as in
CFS02. Since $\|Q-P_0\|<1$, restriction of $P_0$ to
$\operatorname{im}Q$ is injective and hence bijective onto
$\operatorname{im}P_0$, by equal finite dimensions. As $\eta$ is
$Q$-valued, its zero set equals that of the fixed-target map
$f=P_0\eta$. By (DE), the restriction of $Df$ to
$L_{x_0}^{\perp}$ differs from the identity by norm less than
$1/2$. Thus $f$ is a submersion on that neighborhood, proving the
embedded-manifold assertion. The global function $\eta$ is
continuous on $U$, so its zero set is relatively closed; intersecting
with any compact subset of $U$ proves properness.

For the asserted graph take $v=x_i$ and that center as reference.
The closed normal balls in the displayed cylinder lie strictly
inside $B(x_i,\lambda r_i)$. Write $f(t,n)=n+e(t,n)$.
For small $\delta$, (DE) gives $|e|\le\lambda r_i/8$ and
$\|D_ne\|<1/2$. For each $t$ the map $n\mapsto-e(t,n)$
is a contraction from the closed normal ball of radius
$\lambda r_i/4$ into its interior. Completeness gives exactly one
fixed point, and every zero in that cylinder is that fixed point.
The implicit function theorem makes it smooth. Its value is
$O(\delta r_i)$ and
\[
 Dg_i=-(I+D_ne)^{-1}D_te
\]
is $O(\delta)$. Repeated differentiation gives the higher estimates:
the inverse has norm at most two, every derivative of $e$ has
its bound from (DE), and the finite product rule uses only
multilinear operator norms. Induction in $q$ supplies constants
$F_q$ independent of the ambient dimension. These graph descriptions
refer to the same globally defined zero set; they introduce no
additional choices to glue.
\end{proof}

\paragraph{What this iteration supplies and what remains.}
CFS02--CFS05 give a single smooth zero-set construction, uniform
neighborhood multiplicity, affine-plane coherence, and estimates in
all ambient directions. They replace the invalid unshrunk-cover
inference; FC29 is now applicable on a fixed neighborhood rather
than only at one point. Generic Euclidean calculus, finite-dimensional
spectral calculus, the implicit function theorem and the contraction
principle remain inherited foundation bindings.

This construction deliberately uses $\lambda r$ and yields properness
only in $U$. It does not rename that domain $N_r(S)$. FC28 still needs
an interpolant on sufficiently buffered regions, both directions of
its \emph{full} pointed coverage, a separation argument excluding
competing sheets or missing nearest points, and smooth nearest-point
projection estimates on every $B(x,r(x))$, $x\in S$. A small local
graph does not prove those global tube statements. Nor does cloud
information at $S$ supply planes at every point of $\widetilde S$.
In particular the source's restriction of its zero set to $N_r(S)$
cannot alone establish arbitrary $\varepsilon^{-1}r(x)$-radius
pointed coverage when $S$ is small; even $S=\{0\}$ with a long
flat enlarged interval makes that domain issue visible.

The next FC28 iteration must specify the enlarged construction domain
and its transition/coverage buffers, then prove the full-tube nearest
projection conclusion, or prove a precisely weaker interface suffices
for every FC30--FC33 consumer. No weakening has been adopted here.
The original-overlap producer FC23, projected clouds FC27, adjustment
cutoffs FC30, one-sheet submersions FC33, proper bundles and corners
FC35--FC38, topological closure FC41, boundary assembly FC43 and the
complete parameter binding FC44 remain separate obligations. The
new smaller-cover constants depend only on $k,C$ and the fixed cutoff;
they add no dependence on $\dim H$, the later $\Delta$, or a global
radius ratio. This has not yet completed the KL15.4 inequality register
or the Morgan--Tian compatibility gates H1--H4.

\subsection{The full-tube obstruction and an application-specific repair route}
\label{sec:fibration-cloud-full-tube-repair}
The attempted extension of CFS05 reveals a stronger problem than
CFS01. The literal radius hypothesis permits incompatible model planes
on overlapping full-radius balls. Thus the general FC28 implication
must be withdrawn; adding more details to its proof cannot supply it.
The following counterexample concerns the archived statement of KL20.2,
including its corrected enlarged-set convention. It is a project finding,
not an author-confirmed correction. The smaller-tube CFS02--CFS05
construction is unaffected.

\begin{example}[The general smoothing conclusion fails; CFS06]
\label{ex:fibration-cloud-incompatible-full-tubes}
For $k=K=1$, $C=2$ and $\varepsilon=1/10$, no positive cloud
threshold $\delta$ guarantees even properties (1) and (3) of the
printed KL20.2 statement. Equivalently, the approximation and normal
requirements of historical FC28 cannot hold for all its stated inputs.
This remains so with a compact enlarged set and a positive bounded
Lipschitz radius.
\end{example}
\begin{proof}
For \emph{any} $\delta>0$, put
\[
 a=3/4,\quad h=\min\{\delta/10,1/100\},\quad
 s=\frac{\delta h}{100(1+\delta)},\quad R=2+\delta^{-1}.
\]
In $H=\R^2$ take $o=(0,0)$, $p=(a,0)$, $S=\{o,p\}$ and
\[
 \widetilde S=
 \bigl(([-R,a-h]\cup[a+h,R])\times\{0\}\bigr)
       \ \cup\ \bigl(\{a\}\times[-h,h]\bigr).
\]
Set $r(z)=\max\{s,1-2|z|\}$ on $\widetilde S$. This is
$2$-Lipschitz, $s\le r\le1$, $r(o)=1$ and $r(p)=s$.
In particular it satisfies the displayed radius inequality of
Definition20.1 with $C=2$ for every ordered pair in $\widetilde S$.
Choose $A_o$ horizontal and $A_p$ vertical through $p$.

The (CS) test at $p$ is exact: its radius is
$s/\delta=h/[100(1+\delta)]<h$, so it sees only the vertical
segment and no horizontal endpoint. At $o$, the tested radius is
$T=\delta^{-1}<R$. All vertical points in the test have distance
at most $h$ from the horizontal plane, with their horizontal
projections still inside the same test ball. Conversely, a tested
horizontal point outside the removed gap belongs to $\widetilde S$.
For a point in that gap use $(a-h,0)$ whenever the gap meets the
test ball: then $T>a-h$ and this endpoint is in the same ball.
The distance is at most $2h\le\delta/5$. If closed balls are used
and $T=a-h$, that endpoint itself supplies the boundary point.
Thus both truncated-set bounds hold with error at most $\delta$.
The open-ball case has the same strict witnesses. The enlarged set
is compact and all the historical compact-carrier radius bounds hold.

Suppose the asserted $W$ exists. Pointed coverage at $p$ supplies
$w\in W$ with $|w-p|<2\varepsilon s$. Since $s<1/10000$,
this point lies in both $B(p,r(p))$ and $B(o,r(o))$. Let $P_w$
be its normal projector. The two required normal bounds give
\[
 1=\|P_{A_o^{\perp}}-P_{A_p^{\perp}}\|
 \le\|P_{A_o^{\perp}}-P_w\|+
       \|P_w-P_{A_p^{\perp}}\|
 \le2\varepsilon=1/5,
\]
a contradiction. Here $A_x^{\perp}$ means the normal linear
space to its direction, as in the source. No nearest-point theorem
or curvature estimate is used in this contradiction. Since the
construction works for every $\delta>0$, choosing a smaller source
threshold does not fix the statement.

There is also a direct obstruction to the projection conclusion alone.
The point $p$ belongs to both full-radius balls. Evaluating the two
first-derivative bounds there would give
$\|\pi_{A_o^0}-DP_p\|\le\varepsilon$ and
$\|DP_p-\pi_{A_p^0}\|\le\varepsilon$, again implying
$1\le2\varepsilon$. Thus discarding the pointed-coverage requirement
would not rescue the general full-tube derivative statement.
\end{proof}

\paragraph{Contract consequence.}
FC28 is now a \emph{historical source claim refuted at its stated
generality}, with an application-specific replacement still to prove.
Its old inventory row is preserved for provenance; the current overlay
and task status record this disposition. FC30--FC33 may not import it
as a supplied theorem. The old smoothing-modulus lemma remains a valid
\emph{conditional} quantifier argument but its general FC28 premise is
unavailable. None of this refutes the collapse theorem or the particular
coordinate clouds: those clouds retain stronger scale information.

\begin{lemma}[Nearby scale comparison from retained markers; CFS07]
\label{lem:fibration-cloud-marker-scale-comparison}
Let $L\ge1$ and $0<\Sigma\le1/(5L)$.
Under the exact hypotheses of FC26, its arbitrarily selected radius
$r(x)=\Sigma\rho(p_x)$ satisfies, for all image points $x,y$,
\[
 |x-y|\le L\max\{r(x),r(y)\}
 \quad\Longrightarrow\quad
 \frac35\le\frac{r(y)}{r(x)}\le\frac53.                 \tag{MC}
\]
The same conclusion holds for FC04's exact first-image radius,
using its $\Sigma$-Lipschitz property. The constants do not depend
on preimage choices, the global range of $\rho$, or $\dim H$.
\end{lemma}
\begin{proof}
For the FC26 case let $R$ be the larger of the two chosen full-marker
block radii. The scale bounds at those blocks give
$\max\{\rho(p_x),\rho(p_y)\}\le5R/4$. Hence
\[
 |x-y|\le L\Sigma(5R/4)\le R/4<R/2.
\]
The marker equal to $R$ at one point is positive at both, since
orthogonal coordinate projection is $1$-Lipschitz. The hypothesis
for \emph{every} preimage at a positive marker puts both chosen
scales in $[3R/4,5R/4]$, proving (MC). This argument does not require
continuity or descent of the chosen radius.

For FC04, write $r_+=\max(r(x),r(y))$ and $r_-$ for the minimum.
Then $r_+-r_-\le\Sigma|x-y|\le L\Sigma r_+\le r_+/5$.
Thus $r_+/r_-\le5/4$, which is stronger than required.
\end{proof}
This uses the full marker and support-scale data already specified in
FC26, not merely its weaker $(2,\delta)$ cloud consequence. The actual
projected cloud producers and all-marker support bindings remain FC27
obligations. The finite constant $L$ must be fixed \emph{before}
$\Sigma$; an unbounded later choice of $L$ is not allowed by this lemma.

\begin{lemma}[Plane coherence on a prescribed larger scale; CFS08]
\label{lem:fibration-cloud-macroscopic-plane-coherence}
Suppose the actual (CS) tests hold and $x,y\in S$ satisfy
\[
 B^{-1}\le r(y)/r(x)\le B,\qquad
 |x-y|\le L\max\{r(x),r(y)\},\qquad B,L\ge1.
\]
For
\[
 0<\delta<\min\left\{\frac1{4B},\frac1{2(LB+3)},
                            \frac1{4(B+1)}\right\},
\]
the normal projectors obey
\[
 |P_x(y-x)|\le\delta r(x),\qquad
 \|P_y-P_x\|\le6(B+1)\delta.                             \tag{LC}
\]
In particular CFS07 supplies $B=5/3$ for any fixed $L$ allowed
by its prior parameter choice.
\end{lemma}
\begin{proof}
Normalize $x=0$ and $r(x)=1$. Then $|y|\le LB<\delta^{-1}$,
so the set-to-plane test at $x$ gives the offset bound. For unit
$u$ in the direction of $A_y$, the point $y+u$ lies strictly in
its tested plane ball, because $\delta<1/(4B)$. Choose
$z\in\widetilde S$ with $|z-y-u|<2\delta r(y)\le2B\delta$.
Then $|z|<LB+2<\delta^{-1}$, and the other test bounds its
distance to $A_x$ by $\delta$. Subtracting the normal offset of
$y$ gives $|P_xu|\le2(B+1)\delta<1/2$. The equal-dimensional
projection argument in CFS03 now gives the bound $6(B+1)\delta$.
All tests have strict interior slack, regardless of whether (CS)
uses open or closed truncations.
\end{proof}
Thus the two-scale perpendicular-plane obstruction is excluded when
(MC) is supplied. This is a plane-comparison result, not the production
of a large graph or a proof of full pointed coverage.

\begin{lemma}[Nearest points to a buffered graph, with all derivatives; CFS09]
\label{lem:fibration-buffered-graph-nearest-projection}
Let $H=E\oplus F$ be a finite-dimensional orthogonal splitting,
$\dim E=k$, $R>0$, $m\ge1$ and $0<a\le1/100$.
Let $g:B_E(0,4R)\to F$ be $C^{m+1}$ with
\[
 \|D^qg\|\le aR^{1-q}\qquad(0\le q\le m+1).             \tag{GB}
\]
Let $G$ be this whole graph. Suppose a subset $W\subset H$ satisfies
\[
 G\subset W,\qquad W\cap B_H(0,3R)=G\cap B_H(0,3R).     \tag{IS}
\]
Every $z\in B_H(0,R)$ has a unique nearest point in the \emph{whole}
set $W$. The nearest-point map $P$ is $C^m$, has rank $k$, and
\[
 |P(z)-\pi_Ez|\le3aR,\qquad
 \|DP-\pi_E\|\le8a,
\]
\[
 \|D^q(P-\pi_E)\|\le c_q aR^{1-q}\quad(2\le q\le m),   \tag{NP}
\]
where $c_q$ depends only on $q$, independently of $\dim H$.
If $g$ is smooth, so is $P$. No global closedness of $W$ is assumed;
the buffered graph and the exclusion in (IS) are essential inputs.
\end{lemma}
\begin{proof}
Write $z=u+v$ in the splitting. The graph point $u+g(u)$ is at
distance at most $(1+a)R$ from $z$. Any point of $W$ outside
$B(0,3R)$ has distance greater than $2R$ from $z$. A graph point
with $|t|\ge5R/2$ has horizontal distance greater than $3R/2$.
Both bounds are larger than the candidate distance. Consequently
it suffices to minimize on the graph over the compact closed
$E$-ball of radius $5R/2$, and no minimizing point is on its
boundary. This proves existence in $W$ without assuming $W$ closed.

On that convex parameter ball, half the squared distance is
\[
 F_z(t)=\tfrac12\bigl(|t-u|^2+|g(t)-v|^2\bigr).
\]
For any $b\in E$ its Hessian satisfies
\[
 D^2F_z(t)[b,b]=|b|^2+|Dg(t)b|^2+
              \langle g(t)-v,D^2g(t)[b,b]\rangle
 \ge\bigl(1-a(1+a)\bigr)|b|^2.
\]
Thus it is strictly convex and the minimizing parameter $t(z)$ is
unique. Its equation is
\[
 F(t,z):=t-u+Dg(t)^*(g(t)-v)=0.                           \tag{NE}
\]
Put $A=D_tF$. The lower Hessian bound gives $\|A^{-1}\|<2$.
The implicit function theorem, followed by uniqueness on overlaps,
makes $t(z)$ a $C^m$ map on the whole ball, and
\[
 D_zt=A^{-1}[I_E,Dg(t)^*],\qquad
 DP=(I_E,Dg(t))A^{-1}[I_E,Dg(t)^*].                       \tag{PD}
\]
The first bracketed map is onto $E$, so $P$ has rank $k$ onto
its graph. For smooth $g$ the same reasoning is smooth.

Equation (NE) gives $|t-u|\le a(1+a)R$, and $|g(t)|\le aR$;
the value bound follows. Also
\[
 \|A-I_E\|\le a+2a^2,\qquad
 \|A^{-1}-I_E\|\le2(a+2a^2).
\]
Compare (PD) with the product of the inclusion $E\hookrightarrow H$
and its orthogonal projection. The two graph-factor errors have
norm at most $a$ and their factors have norm at most $1+a$.
The product rule for differences bounds the total by
\[
 2a(1+a)+2(a+2a^2)(1+a)+a<8a.
\]

For higher derivatives normalize $R=1$. On the indicated bounded
parameter and target balls, the remainder
$F(t,z)-(t-u)=Dg(t)^*(g(t)-v)$ has derivatives through order $m$
bounded by constants depending only on their order times $a$.
This follows from (GB), the product rule and $|g-v|\le1+a$.
Differentiate $F(t(z),z)=0$ successively. The term containing the
highest derivative of $t$ is $A D^qt$. Every other term for
$q\ge2$ contains a derivative of the remainder and hence a
factor $a$; induction and $\|A^{-1}\|<2$ give
$\|D^qt\|\le c'_q a$. The chain rule for $P=t+g(t)$ gives
(NP). Restoring the scale contributes $R^{1-q}$.
Only multilinear operator norms, inner products and finite product
rules occur; no sum over ambient normal coordinates introduces a
dimension factor. One extra derivative of $g$ is required in (NE),
which is why (GB) goes through $m+1$.
\end{proof}

\begin{corollary}[The full-tube analytic assembly from one global graph certificate; CFS10]
\label{cor:fibration-full-tube-from-buffered-graphs}
Suppose a single smooth embedded $k$-manifold $W\subset H$ is supplied.
For every $x\in S$, suppose that after translating $x$ to zero and
splitting $H=L_x\oplus L_x^\perp$, this \emph{same} $W$ satisfies
(GB) and (IS) with $R=r(x)$, one $a\le1/100$, and $m=K$.
The graph functions are smooth, with the bounds through $K+1$.
Then $W\cap N_r(S)$ is properly embedded in $N_r(S)$ and there is
one smooth nearest-point submersion $P:N_r(S)\to W$. It satisfies
CFS09's scaled value and derivative estimates relative to $P_{A_x}$
on every $B(x,r(x))$. On $W\cap B(x,r(x))$ the normal projections
are within $2a$ of $P_x$.
\end{corollary}
\begin{proof}
CFS09 gives a nearest point in the entire set $W$ for each point
of each ball. Uniqueness in that same set forces the maps to agree
on overlaps, so they define one smooth map, with the stated local
rank and estimates. This is not gluing arbitrarily chosen local
nearest-point maps to different graph patches.

Condition (IS) makes $W$ relatively closed near every point of
$N_r(S)$. Indeed any convergent sequence there can be tested in
one $B(x,3r(x))$ containing a neighborhood of its limit; its graph
parameters have norms below $3r(x)<4r(x)$, so continuity of that
graph includes the limit. Thus $W\cap N_r(S)$ is closed in
$N_r(S)$; intersecting with compact subsets gives properness.
A tangent space is the graph of $Dg_x$, whose norm is at most $a$.
A unit graph vector has normal component at most $a$, while a unit
vector of $L_x$ is within $a$ of that graph. The two-term projection
identity used in CFS03 bounds the difference of tangent, and hence
normal, projectors by $2a$.
\end{proof}
For a prescribed projection/normal accuracy $\varepsilon$, choose $a$
below $1/100$, $\varepsilon/3$, $\varepsilon/8$ and every
$\varepsilon/c_q$. This gives the required small constants without
changing the full tube. The implication is now written, but its
\emph{single global buffered graph certificate} is not yet produced.
CFS05 has much smaller graph domains and does not supply it.

\paragraph{Revised next step and retained consumers.}
The following records the revision122 frontier, advanced by
Section~\ref{sec:fibration-cloud-global-repair} below.
Use the actual marker data to retain (MC) at a fixed sufficiently
large comparison scale before choosing $\Sigma_j$. Combine CFS08
with the CFS02--CFS05 local construction to produce a \emph{single}
interpolant satisfying CFS10's $4r$ graph and $3r$ isolation data
at every required center. One must also prove the requested
pointed enlarged-set coverage and $W\subset N_{\varepsilon r}(\widetilde S)$,
or explicitly prove that a replacement covers every actual consumer.
No such weakening is adopted in this revision. A scale $L$ depending
on a later accuracy cannot be inserted into the prior $\Sigma$ choice
without repairing the parameter order.

The source consumers are concrete: (13.5), the analogous edge projection,
and (13.37) use the full $N_{r_j}(S_j)$ domains;13.6/13.21/13.38 require
cutoff supports inside those domains,13.15/13.34/13.43 use value and
first-derivative errors, and13.46 also uses normal control and actual
one-sheet graph restrictions. These are separate from CFS10's analytic
implication. FC23/FC27's cloud producers, FC30's closed support buffers,
FC33's one-sheet conclusion, the proper bundles/corners and global
assembly remain qualified. The source-stated general FC28 theorem is
no longer a valid premise for these consumers. The actual application
may be repaired using the stronger marker structure, but that repair
has not been completed by merely excluding CFS06's counterexample.

\subsection{A global smoothing theorem with retained marker-scale control}
\label{sec:fibration-cloud-global-repair}
The next construction supplies the global graph certificate left open
in CFS10, and then the full pointed coverage. It uses a stronger,
explicit radius hypothesis. This does not restore the refuted general
KL20.2 statement. The model planes are still supplied only at $S$.

Fix a buffer parameter $b\ge1$, and put $B=5/3$ and $D=165$.
In CFS11--CFS14 assume that $S\ne\varnothing$, that
$S\subset\widetilde S\subset H$ is bounded in a finite-dimensional
Euclidean space, that $r$ on $\widetilde S$ is bounded above and
away from zero, and that the actual (CS) tests hold with quality
$\delta$ and chosen affine $k$-planes $A_x=x+L_x$, $x\in S$.
Require the additional hypothesis
\[
 \begin{gathered}
 x,y\in\widetilde S,\quad |x-y|\le128b\max\{r(x),r(y)\}
 \\
 \Longrightarrow\quad B^{-1}\le r(y)/r(x)\le B.
 \end{gathered}                                    \tag{MCb}
\]
Continuity of $r$ is not assumed. The constants below may depend on
$k,b$, a fixed smooth cutoff and the requested derivative order, but
not on $\dim H$, the number of selected centers or the global radius
ratio. Initially take
\[
 0<\delta\le d_b:=\min\left\{\frac1{8B},
             \frac1{4(128bB+3)},\frac1{8(B+1)}\right\}.
                                                               \tag{DS}
\]
Further smallness choices will be explicit dependencies, not changes
of the supplied planes or of the radius function.

\begin{lemma}[A large buffered cover with intrinsic multiplicity; CFS11]
\label{lem:fibration-cloud-large-cover}
There is a finite $T=\{x_i\}\subset S$ with disjoint balls
$B(x_i,r_i)$, such that every $x\in S$ has a selected $x_i$ with
$|x-x_i|<3r_i$ and $r(x)\le2r_i$. Define
\[
 U_b=\bigcup_i B(x_i,20br_i),\qquad
 \Omega_b=\bigcup_i B(x_i,30br_i),\qquad
 N_b=\left\lceil(1+2BDb)^k\right\rceil.
\]
Then $N_{8br}(S)\subset U_b\subset\Omega_b$.
For either of the reference balls
\[
 V_x=B(x,8br(x))\quad(x\in S),\qquad
 V_i=B(x_i,30br_i)\quad(x_i\in T),
\]
let $J$ be all selected centers whose closed $80br_j$-ball meets
that reference ball. With its center denoted $x$ and radius scale
$r_x=r(x)$, every $j\in J$ satisfies
\[
 B^{-1}r_x\le r_j\le Br_x,\qquad
 |x_j-x|<Db r_x,\qquad |J|\le N_b.                    \tag{LM}
\]
Moreover
\[
 |P_x(x_j-x)|\le\delta r_x,\qquad
 \|P_j-P_x\|\le6(B+1)\delta,                         \tag{LP}
\]
where $P_x$ denotes the normal projector to $L_x$.
\end{lemma}
\begin{proof}
Apply KL19.1 to $S$ with the original radius $r$. Boundedness,
finite-dimensionality and the positive lower radius bound make the
selection finite. Its number need not have a uniform global bound.
The displayed cover and radius dominance are the exact greedy
conclusions. If $z\in B(x,8br(x))$, the corresponding selected center
has
\[
 |z-x_i|<8br(x)+3r_i\le(16b+3)r_i\le19br_i<20br_i.
\]
Thus the large tube is contained in $U_b$.

For either reference ball, a center in $J$ has
\[
 |x_j-x|<80br_j+30br_x\le110b\max(r_j,r_x).
\]
Hypothesis (MCb) gives the radius comparison. Consequently
$|x_j-x|<(80B+30)br_x<Db r_x$.
Conditions (DS) imply both $Db<\delta^{-1}$ and the strict
smallness requirements of CFS08 with $L=128b$. That lemma proves
(LP), using both actual cloud tests at these core centers.

Project the centers of $J$ to $A_x$. Their original separation is
at least $2r_x/B$, by disjointness of the selected balls. Projection
loses at most $2\delta r_x$, so (DS) leaves separation at least
$r_x/B$. All projected centers lie in a $k$-ball of radius $Db r_x$.
The disjoint $k$-balls of radius $r_x/(2B)$ around them lie in its
$r_x/(2B)$ enlargement. Comparing $k$-dimensional volumes gives
(LM). The list includes every support meeting the whole reference
ball, including supports that vanish at a particular point.
\end{proof}

\begin{lemma}[One spectral displacement section on the large domain; CFS12]
\label{lem:fibration-cloud-large-spectral}
Fix a nonnegative smooth $\phi$ equal to one on $[0,1]$ and zero
on $[2,\infty)$. On $\Omega_b$ set
\[
 \phi_i(z)=\phi\bigl(|z-x_i|/(40br_i)\bigr),\quad
 w_i=\frac{\phi_i}{\sum_j\phi_j},\quad
 O=\sum_iw_iP_i,\quad \mu=\sum_iw_ix_i.
\]
For sufficiently small $\delta$ depending only on $k,b,\phi$, the
spectral projector $Q$ of $O$ for the cluster near one is smooth
of rank $\dim H-k$. Define the single smooth section
$\eta(z)=Q(z)(z-\mu(z))$ on $\Omega_b$.
For every fixed $m\ge1$ there are constants $E_q$, $0\le q\le m$,
such that on each reference ball in CFS11,
\[
 \|D^q(Q-P_x)\|\le E_q\delta r_x^{-q},\qquad
 \left\|D^q\bigl(\eta-P_x(\,\cdot-x)\bigr)\right\|
                      \le E_q\delta r_x^{1-q}.       \tag{LS}
\]
These are full ambient multilinear operator norms. The estimates
hold at every core center, not just at selected centers.
\end{lemma}
\begin{proof}
On $V_i$ the $i$th cutoff equals one. On $V_x$ choose the selected
center given by CFS11. Its distance from every point of $V_x$ is
less than $19br_i$, so its cutoff again equals one. In both cases
the denominator is at least one on the entire reference ball.
Only the fixed list $J$ of CFS11 contributes there. Radius comparison,
its bound $N_b$ and the product/reciprocal rules give
$\|D^qw_i\|\le c_q r_x^{-q}$. Radial derivatives have
ambient-dimension-independent bounds, as in CFS04.

Since $\sum_iw_i=1$, (LP) gives $\|O-P_x\|\le6(B+1)\delta$
and corresponding derivative estimates. Choose this norm below
$1/4$. The two spectral clusters are separated, with $k$ eigenvalues
near zero. The same fixed contour and resolvent derivative calculation
as in CFS04 gives the first estimate of (LS), in all directions.
These locally computed projectors coincide: each is the spectral
projector of the same globally defined $O$. No choice of local
spectral frames is glued.

Translate $x$ to zero and rescale $r_x$ to one. On a reference ball,
$|z|<30b$ and every active center has norm less than $Db$.
Write
\[
 \eta(z)-P_xz=(Q(z)-P_x)(z-\mu(z))-P_x\mu(z).
\]
The weights bound $\mu$ and its derivatives. Every $P_xx_i$ has
norm at most $\delta$, by (LP). Product differentiation gives the
second estimate of (LS), and rescaling restores the stated powers.
All constants depend only on the stated fixed data. Alternatively,
FC29 applies to the span of the at most $N_b$ active tangent planes
and translated centers on this whole ball; it bounds the complementary
rank, not the rank of $O$. The derivative proof does not suppress
normal-coordinate dependence of the cutoffs.
\end{proof}

\begin{proposition}[A global interpolant with large graphs and sheet exclusion; CFS13]
\label{prop:fibration-cloud-global-buffered-graphs}
For sufficiently small $\delta$, the single set
\[
 W=\{z\in U_b:\eta(z)=0\}
\]
is a smooth embedded $k$-manifold, properly embedded in $U_b$.
For every $x\in S$ it contains a smooth graph
\[
 G_x=\{x+t+g_x(t):t\in L_x,\ |t|<4br_x\},
 \qquad g_x(t)\in L_x^\perp,
\]
and
\[
 W\cap B(x,3br_x)=G_x\cap B(x,3br_x).                 \tag{GI}
\]
For any fixed $m\ge1$ there are $F_q=F_q(k,b,m,\phi)$ such that
\[
 \|D^qg_x\|\le F_q\delta r_x^{1-q}\quad(0\le q\le m). \tag{GG}
\]
Thus every prescribed common small graph bound through a fixed
finite order can be obtained by decreasing $\delta$, independently
of ambient dimension and the global radius ratio.
\end{proposition}
\begin{proof}
At any $z\in U_b$, choose a selected reference ball $V_i$ containing
it. For $\|Q-P_i\|<1$, restriction of $P_i$ to $\operatorname{im}Q$
is an isomorphism onto $L_{x_i}^{\perp}$. Hence $\eta=0$ is locally
equivalent to $P_i\eta=0$. By (LS), the normal derivative of that
fixed-target map is within $E_1\delta$ of the identity. Taking
$E_1\delta<1/2$ proves it is a submersion. This proves smooth
embeddedness. The global zero set is closed in $U_b$, giving relative
properness.

Fix an arbitrary core center $x$. For $|t|<4br_x$ and
$|n|\le r_x$, the point $x+t+n$ lies in $B(x,8br_x)\subset U_b$.
Write $f_x(t,n)=P_x\eta(x+t+n)=n+e_x(t,n)$. By (LS), choose
$\delta$ so $|e_x|\le r_x/4$ and $\|D_ne_x\|<1/2$ on all
these cylinders. For each $t$, the map $n\mapsto-e_x(t,n)$
is a contraction of the closed normal $r_x$-ball into its interior.
It has exactly one fixed point, of norm at most $r_x/4$.
The isomorphism $P_x:\operatorname{im}Q\to L_x^\perp$ on this
ball makes that fixed point a zero of the original $\eta$, so it
belongs to the same $W$ defined above. The implicit function theorem
makes $g_x$ smooth. Differentiating $f_x(t,g_x(t))=0$ exactly as
in CFS05 proves (GG), with constants independent of dimension.

If $w\in W\cap B(x,3br_x)$, let $t,n$ be its tangential and
normal coordinates at $x$. Since $P_x\eta(w)=0$, (LS) gives
$|n|\le E_0\delta r_x\le r_x/4$. Also $|t|<3br_x<4br_x$.
Thus $w$ is in the contraction cylinder and equals its unique graph
point. The reverse inclusion in (GI) follows from $G_x\subset W$.
This excludes every competing sheet in the indicated ball, not just
sheets visible in one selected-center chart. The graphs are descriptions
of a single global zero set, not independently chosen patches.
\end{proof}

\begin{theorem}[Full smoothing under scale control at the accuracy-dependent buffer; CFS14]
\label{thm:fibration-cloud-controlled-global-smoothing}
Fix $k,K\ge1$ and $0<\varepsilon\le1/10$, and put $b=\varepsilon^{-1}$.
There is $d_*(k,K,\varepsilon)>0$, independent of ambient dimension,
such that the hypotheses above with (MCb) and any
$0<\delta\le d_*$ produce a smooth $W$ satisfying all the listed
consumer conclusions of historical FC28:
\begin{enumerate}[leftmargin=*]
\item At each core center, the $r_x^{-1}$-rescaled ENLARGED set
and $W$ have truncated Hausdorff error at most $\varepsilon$ on
the common radius-$\varepsilon^{-1}$ ball. Both open and closed
truncations work. Moreover $W\subset N_{\varepsilon r}(\widetilde S)$.
\item $W\cap N_r(S)$ is properly embedded in $N_r(S)$, and its
normal projectors on $W\cap B(x,r_x)$ differ from $P_x$ by at most
$\varepsilon$.
\item There is one smooth nearest-point submersion $P:N_r(S)\to W$.
On every $B(x,r_x)$, $|P-P_{A_x}|\le\varepsilon r_x$ and
$\|D^q(P-P_{A_x})\|\le\varepsilon r_x^{1-q}$ for $1\le q\le K$.
\end{enumerate}
This is a strengthened-hypothesis replacement for the required
smoothing interface. It does not claim the source's additional
curvature-derivative conclusion or supply the actual cloud producers.
\end{theorem}
\begin{proof}
Use CFS11--CFS13 with $m=K+1$. Choose $0<a\le1/100$ so
$8a\le\varepsilon$ and $c_q a\le\varepsilon$ for every
$2\le q\le K$, using CFS09's constants. Reduce $d_*$ so the
resulting graphs have (GG) bounded by $a r_x^{1-q}$ through $K+1$,
and so
\[
 \delta\le\varepsilon/16,\qquad
 B(E_0+2)\delta<\varepsilon,                           \tag{AC}
\]
in addition to all earlier strict smallness bounds. Here $E_0$ is
the fixed CFS12 constant at this $b$. All are positive upper-bound
choices depending only on $k,K,\varepsilon,\phi$. Choosing $d_*$
strictly below the relevant thresholds makes them hold for every
$0<\delta\le d_*$.

Restrict each graph to $|t|<4r_x$. Since $b\ge1$, (GI) gives
exactly CFS10's $3r_x$ isolation for that restricted graph.
CFS09--CFS10 give the unique nearest point in the whole $W$, the
single smooth submersion on $N_r(S)$, relative properness, normal
error at most $2a$ and the stated derivative bounds. The value error
is at most $3ar_x\le\varepsilon r_x$.

For global proximity, take $w\in W$ and a selected center with
$|w-x_i|<20br_i$. Estimate (LS) and $\eta(w)=0$ give normal
distance from $w$ to $A_{x_i}$ at most $E_0\delta r_i$.
Its affine projection has distance less than $20br_i$ from $x_i$,
strictly within the cloud test by (DS). The plane-to-set test supplies
$q\in\widetilde S$ at distance less than $2\delta r_i$ from that
projection. Thus $|q-x_i|<21br_i$, and (MCb) implies
$r_i\le Br(q)$. Consequently
\[
 |w-q|<(E_0+2)\delta r_i
       \le B(E_0+2)\delta r(q)<\varepsilon r(q).
\]
This proves the required proximity with the radius AT THE WITNESS
$q$, not an unproved descended radius on $W$.

It remains to check actual same-radius coverage. Translate an arbitrary
$x\in S$ to zero and rescale $r_x$ to one. For
$q\in\widetilde S\cap\overline B(0,b)$, the cloud test gives
$|q-\pi_{L_x}q|\le\delta$. Trim $t=\pi_{L_x}q$ radially to
the closed tangent ball of radius $b-2a$, changing it by at most
$2a$. The graph point over the trimmed vector lies in $W\cap B(0,b)$,
since its height is at most $a$. Its distance from $q$ is at most
$\delta+3a\le7\varepsilon/16<\varepsilon$.

Conversely, a point $w\in W\cap\overline B(0,b)$ belongs to
$G_x$ by (GI), and its normal component has norm at most $a$.
Trim its tangent coordinate to radius $b-3\delta$, moving it by
at most $3\delta$. This plane point is strictly inside the original
cloud test. A plane-to-set witness of error less than $2\delta$
lies in $\widetilde S\cap B(0,b)$ and is at distance less than
$a+5\delta\le7\varepsilon/16$ from $w$. The trimmed radii are
positive. These two arguments prove the closed-ball Hausdorff bounds
and also the open-ball bounds, since both constructed witnesses lie
strictly inside the common ball. No invalid same-radius intersection
rule has been used. Restoring $r_x$ gives conclusion (1).
\end{proof}

\begin{corollary}[A smoothing modulus with the marker radius chosen afterward; CFS15]
\label{cor:fibration-cloud-marker-modulus}
For fixed $k,K$ there are $\theta_1>0$ and a positive function
$\Xi(\Gamma)\to0$ as $\Gamma\downarrow0$ with the following
conditional application. First choose $0<\Gamma<\theta_1$ and set
$\varepsilon=\Xi(\Gamma)$; then choose
\[
 0<\Sigma\le\Xi(\Gamma)/640.                          \tag{MO}
\]
Any actual cloud of quality $\Gamma$ satisfying the bounded-carrier
hypotheses and the exact FC04 or FC26 marker hypotheses with this
$\Sigma$ has all CFS14's conclusions at accuracy $\Xi(\Gamma)$.
The modulus and this extra radius bound are independent of ambient
dimension. This is not a claim that the actual cloud/marker hypotheses
or every other parameter inequality have already been produced.
\end{corollary}
\begin{proof}
For $m\ge1$ put $\varepsilon_m=2^{-(m+3)}$ and choose the
threshold $d_m=d_*(k,K,\varepsilon_m)$ from CFS14. Set
\[
 \theta_m=\min\{2^{-m},d_1/2,\ldots,d_m/2\}.
\]
For $0<\Gamma<\theta_1$, take the largest $m$ with
$\Gamma<\theta_m$ and put $\Xi(\Gamma)=\varepsilon_m$.
The set of eligible $m$ is nonempty and finite. As in the elementary
R104 modulus argument, $\Gamma<\theta_M$ implies $m\ge M$,
so $\Xi(\Gamma)\to0$. Here the thresholds come from the proved
CFS14 implication, not the refuted general FC28 claim.

After this choice set $b=\Xi(\Gamma)^{-1}$ and $L=128b$.
Condition (MO) is exactly $\Sigma\le1/(5L)$, so CFS07 supplies
(MCb) for every pair of enlarged image points. Since
$\Gamma<d_m$, CFS14 applies directly with quality $\delta=\Gamma$.
No monotonicity of an arbitrary fixed-error Hausdorff predicate is
needed. Any finitely many other positive upper bounds on $\Sigma$
can be met by taking their minimum together with (MO), provided their
hypotheses and predecessor dependence are actually supplied.
\end{proof}

\paragraph{Source order, completed analytic part and remaining producers.}
The source's Section13.2 (printed72/PDF67) explicitly uses
$\Xi_1=\Xi_1(\Gamma_1)$; the edge and slim passages have the analogous
form. In (15.4), printed87/PDF82, $\Gamma_j$ precedes the pair
$\{\Sigma_j,\Xi_j\}$. For this repaired smoothing input, refine
that pair to $\Xi_j$ first and $\Sigma_j$ second, with (MO).
This adds a forward choice inside the pair; it does not put a new
accuracy-dependent bound on a radius already chosen. KL12.7's proof
(printed67--68/PDF62--63) also chooses $\Sigma_1$ after $\Gamma_1$.
These exact passages and the author correction sheet were rechecked.
The full simultaneous constraint register FC44 is still an obligation.

CFS11--CFS15 supply a written global smoothing theorem, including full
pointed coverage and the entire projection domain, from explicit
stronger scale control. CFS15 binds that control to the existing
conditional marker contracts. It does not bind or prove FC23/FC27's
actual image coverage, selected planes or all-marker support hypotheses.
Nor does an abstract interpolant establish the cutoff support enclosures
of FC30, the adjustment estimates for the actual maps, or FC33's
original-domain one-sheet and fiber-preservation conclusions.
The next work is to bind this replacement to the three actual clouds
and their successive adjustments, keeping these domains and the new
parameter inequality in every consumer. Proper bundles/corners,
global topology, boundary assembly and Morgan--Tian gates H1--H4
remain separate. Generic spectral calculus, contraction and implicit
functions still require accepted-foundation binding and Lean proofs.

\section{Successive adjustments and preservation of fibers}
\label{sec:fibration-adjustments-completion}
\begin{foundation}[Buffered ambient cutoffs; FC30]
\label{found:fibration-ambient-cutoffs}
For the three clouds choose smoothing accuracy $\Xi_j$ and smooth
functions $\psi_j$ on an OPEN neighborhood of the relevant image,
with values in $[0,1]$. Their CLOSED supports within that
neighborhood must lie in the open domain of $P_j\pi_j$.
On $\mathcal E^{j-1}(M)$ require
$\|d\psi_j\|\le C_j/\rho(p)$ at the point indexed by $p$.
The plateaux equal one on coordinate thresholds $6$, $6\Delta$,
$6\cdot10^5\Delta$, and vanish outside thresholds $7$,
$7\Delta$, $7\cdot10^5\Delta$, respectively; for edge include
the corresponding weak-edge restriction.
\end{foundation}
KL 13.6, 13.21 and 13.38 construct these with buffered marker
ratios, using internal thresholds $6.1,6.5,6.9$. Divisions occur
only where the marker is positive and bounded below. The edge
formula also detects the branch with vanishing $E'$ OUTPUT by
its retained edge markers; it is not just a ratio in that block.
The cutoff is constructed using the previous adjusted image and
its quantitative closeness to $\mathcal E^0$.
A statement only about nonzero cutoff values on the image does
not justify a global extension by zero across the tube boundary.
For implementation, restrict to an open neighborhood of the compact
image with a closed-support buffer, or supply a global cutoff with
that stronger property. Either suffices to define the smooth
composition on $M$. CFS18--CFS19 below now supply the neighborhood extension and fixed-data
support stability. Uniform cutoff stability, derivative bounds and
plateaux in the allowed parameter order remain FC30 obligations.

\begin{lemma}[One adjustment estimate; FC31]
\label{lem:fibration-adjustment-estimate}
Let $f:M\to H=Q\oplus Q^\perp$ be smooth, $0\le\psi\le1$,
and $a,b,L,d,e\ge0$. On a buffered
neighborhood define
\[
 \Psi(x)=x+\psi(x)\bigl(P(\pi_Qx)-\pi_Qx\bigr),\qquad g=\Psi f,
\]
with the displayed vector placed in $Q$. Suppose on the support
along $f$, for $\rho>0$ and an orthogonal projection $\pi_A$ in $Q$,
\[
 |P(\pi_Qf)-\pi_Qf|\le a\rho,\quad
 \|d\psi\|\le b/\rho,\quad \|Df\|\le L,
\]
\[
 \|DP-\pi_A\|\le d,\qquad
 \|(I-\pi_A)D(\pi_Qf)\|\le e.
\]
Then $|g-f|\le a\rho$ and
$\|Dg-Df\|\le abL+dL+e$.
Where $\psi=1$ and $DP\,D(\pi_Qf)$ is onto $TW$,
$\pi_Qg=P\pi_Qf$ is a submersion to $W$.
\end{lemma}
\begin{proof}
Differentiate the displayed formula. The derivative of the cutoff
is composed with $Df$, giving the factor $L$ in its term.
Split $(DP-I)D(\pi_Qf)$ into
$(DP-\pi_A)D(\pi_Qf)-(I-\pi_A)D(\pi_Qf)$.
The support buffer justifies smooth extension by the identity.
For rank, combine the supplied onto margin with a perturbation
smaller than that margin, as in FC06.
\end{proof}
Apply this successively for $j=1,2,3$, keeping the displacement
budgets and derivative budgets summable below $c_{\rm adjust}$.
FC28 supplies displacement bounds from actual image coverage and
projection estimates; closeness of the previous image has to fit
INSIDE the same tube. It is not enough that $\Xi_j$ is small in
isolation: the term $a b L$ and prior errors must also be small.

\begin{lemma}[Postcomposition and retained projections; FC32]
\label{lem:fibration-postcomposition}
For $g=\Psi f$, $f(p)=f(q)$ implies $g(p)=g(q)$ and
$\ker Df\subset\ker Dg$. Equality of fibers or kernels requires
injectivity, respectively injective differential, of $\Psi$ on
the relevant image; postcomposition alone supplies only inclusion.
For the adjustments above,
$\pi_{Q_j^\perp}\Psi_j=\pi_{Q_j^\perp}$.
If $Q_3\subset Q_2$ and $\psi_3$ depends only on $Q_2$
coordinates, then $\Psi_3$ factors over $\pi_2$:
\[
 \pi_2\Psi_3=\Psi_{32}\pi_2.
\]
\end{lemma}
\begin{proof}
The first statements are composition and the chain rule. The
adjustment vector lies in $Q_j$, proving the second. The last
formula follows because both its cutoff and its $Q_3$ projection
are determined by $\pi_2x$, while $Q_2^\perp$ is unchanged.
\end{proof}
This direction matters: later stages may merge earlier fibers.
The local graph and connectedness argument below proves that they
do not merge them on the needed restricted bases. The $H_2$ blocks
are unchanged by stages two and three, and $H_e$ by stage three;
these supply the coordinate maps used to check rank and graphness.

\begin{foundation}[Final bases and submersions; FC33]
\label{found:fibration-final-bases}
Deliver KL 13.1 with $\mathcal E=\Psi_3\Psi_2\Psi_1\mathcal E^0$,
\[
 |\mathcal E(p)-\mathcal E^0(p)|<c_{\rm adjust}\rho(p),\qquad
 \|D\mathcal E-D\mathcal E^0\|<c_{\rm adjust}.
\]
There are embedded bases $W_j\subset Q_j$ of dimensions $2,1,1$
such that $\pi_j\mathcal E:U_j\to W_j$ is a submersion,
where $U_j$ are exactly the original thresholds $5$, $5\Delta$
(including $\eta_{E'}<5\Delta$), and $5\cdot10^5\Delta$.
The bases are the appropriately restricted later images of the
smoothed bases, using marker $>.9R_i$ and coordinate thresholds
$5.5$, $5.5\Delta$, $5.5\cdot10^5\Delta$.
\end{foundation}
KL 13.46 must be implemented, not replaced by the claim that the
image of a manifold under a smooth map is a manifold. Its proof
has four tasks: local coordinate projection is a proper onto local
diffeomorphism; hence a finite covering of a Euclidean ball; local
coordinate fibers on $M$ are connected; path lifting in the ORIGINAL
bundle joins preimages inside the buffered chart and rules out two
sheets. FC03 rules out remote preimages. A retained coordinate
projection then certifies rank after later modifications. Edge and
slim cases require their own connected fibers and properness;
$S^2,T^2,D^2,S^1$ are connected, but the actual witnesses must be
those of LC83--LC85. This source-qualified delivery is strictly
stronger than FC32's fiber inclusion.

\subsection{Original clouds, perturbed inputs and closed supports}
\label{sec:fibration-adjustment-domain-repair}
Revision124 supplies the following consumer arguments for CFS14--CFS15.
The clouds remain images of $\mathcal E^0$. KL's overview at
printed72/PDF67 mentions $\pi_j\mathcal E^{j-1}$, whereas the actual
choices in Sections13.3--13.4 refer back to (12.26) and (12.33),
which use $\pi_j\mathcal E^0$. We use these latter definitions and
prove that the perturbed inputs lie in their projection domains.
These are written project arguments, not new Lean declarations.

\begin{lemma}[Every-preimage scale and a half-tube buffer; CFS16]
\label{lem:fibration-original-cloud-tube}
Let $F=\mathcal E^0$, $x=\pi_jF(p)\in S_j$, and assume FC04 for
$j=1$ or the full-marker and EVERY-preimage hypotheses of FC26 for
$j=2,3$. For the radius actually chosen on the original image,
\[
 \frac35\Sigma_j\rho(p)\le r_j(x)\le\frac53\Sigma_j\rho(p).
\]
For $j=1$ there is equality $r_1(x)=\Sigma_1\rho(p)$.
If $|f(p)-F(p)|\le e\rho(p)$ and
$e\le3\Sigma_j/10$, then $y=\pi_j f(p)$ satisfies
\[
 |y-x|\le r_j(x)/2,\qquad
 B(y,r_j(x)/2)\subset B(x,r_j(x))\subset N_{r_j}(S_j).
\]
The entire segment from $x$ to $y$ is in the latter open ball.
No descent of $\rho$ through $\pi_jF$ is required.
\end{lemma}
\begin{proof}
For $j=2,3$, choose a full-marker block at $x$ with radius $R_i$.
Both the selected preimage defining $r_j(x)$ and the present $p$
have positive $i$ marker. Their scale values lie in
$[3R_i/4,5R_i/4]$, so their ratio lies in $[3/5,5/3]$.
FC04 gives the exact first-image identity. Orthogonal projection
contracts distances, and the stated lower bound gives
$e\rho(p)\le r_j(x)/2$. The open-ball inclusion follows from the
strict inequality for points of $B(y,r_j(x)/2)$.
\end{proof}
This bound applies to arbitrary preimages, including those on cutoff
supports. Comparing only the selected preimages is insufficient.
It uses original core membership; closeness to the ENLARGED image
alone does not provide a core-centered projection ball.

\begin{lemma}[Projection estimates at a perturbed input; CFS17]
\label{lem:fibration-perturbed-projection-estimate}
Keep the preceding original cloud and put $B=5/3$,
$0<\varepsilon\le1/10$, $\Sigma=\Sigma_j$.
Assume CFS14 with accuracy $\varepsilon$ and at least one derivative.
Let $f,F:M\to H$ be smooth, $\|DF\|\le L_0$, and suppose
\[
 |f-F|\le E\rho,\qquad \|Df-DF\|\le H_0,\qquad
 E\le3\Sigma/10.
\]
On the closed cutoff support along $f$, suppose $p$ is in the
original core domain, $\|d\psi\|\le b/\rho(p)$, and the chosen
core plane has orthogonal projection $\Pi_x$ satisfying
\[
 \|(I-\Pi_x)D(\pi_jF)_p\|\le\nu.
\]
Suppose the ambient adjustment is smoothly defined as in FC31,
with $0\le\psi\le1$, and put $g=\Psi f$. Define
\[
 a=B\varepsilon\Sigma+(1+\varepsilon)E.
\]
Then the following bounds hold globally:
\begin{align*}
 |g-F|&\le(E+a)\rho,\\
 \|Dg-DF\|&\le
 ab(L_0+H_0)+\varepsilon(L_0+H_0)+\nu+2H_0.
\end{align*}
At a point of the plateau $\psi(f(p))=1$, suppose
$\Pi_xD(\pi_jF)_p$ has rank $k_j$ and its least nonzero singular
value is at least $\mu>0$. If
\[
 \varepsilon(L_0+H_0)+H_0<\mu,
\]
then $D(P_j\pi_j f)_p$ is onto $T_{P_j\pi_j f(p)}W_j^0$.
On an open plateau this is a submersion statement.
\end{lemma}
\begin{proof}
Write $x=\pi_jF(p)$ and $y=\pi_j f(p)$.
CFS16 places the segment $[x,y]$ in $B(x,r_j(x))$.
CFS14 gives $|P_j(x)-x|\le\varepsilon r_j(x)$ and
$\|DP_j-\Pi_x\|\le\varepsilon$ throughout that ball.
Consequently $\|D(P_j-I)\|\le1+\varepsilon$ there. Integration
along the segment gives
\[
 |P_j(y)-y|\le\varepsilon r_j(x)+(1+\varepsilon)|y-x|
 \le a\rho(p).
\]
The normal derivative of $\pi_j f$ is at most $\nu+H_0$.
Apply FC31 with $L=L_0+H_0$, $d=\varepsilon$ and
$e=\nu+H_0$, then add the preceding error $H_0$.
Outside the closed support the adjustment is locally the identity,
so the same global bounds hold.
For rank, the difference between $DP_j(y)D(\pi_j f)_p$ and
$\Pi_xD(\pi_jF)_p$ has norm at most
$\varepsilon(L_0+H_0)+H_0$. On the $k_j$-dimensional orthogonal
complement of the latter kernel this difference is smaller than
its lower singular-value bound. Hence the former map is injective
on that subspace and has rank at least $k_j$. Its image lies in
the $k_j$-dimensional tangent space to $W_j^0$, proving onto-ness.
\end{proof}
In particular, the displacement at $y$ is not bounded just by
$\varepsilon r_j(x)$: the preceding displacement contributes the
explicit $(1+\varepsilon)E$ term. The plane and radius used here
belong to the ORIGINAL center $x$, not a newly asserted cloud at $y$.

\begin{lemma}[Smooth extension on a neighborhood of the image; CFS18]
\label{lem:fibration-closed-support-neighborhood}
Let $O\subset H$ be open, $K\subset O$, $U\subset O$ open, and
$\psi:O\to\R$ smooth. Suppose
$K\cap\operatorname{supp}_O\psi\subset U$, with CLOSED support
relative to $O$. There is an open neighborhood $V$ of $K$ in $O$
such that $\operatorname{supp}_V(\psi|_V)\subset U$.
For any smooth $h:U\to H$, the formula
\[
 \Psi(z)=
 \begin{cases}z+\psi(z)h(z),&z\in V\cap U,\\z,&z\in V\setminus U\end{cases}
\]
defines a smooth map on $V$. It changes neither $\psi$ nor its
bounds along $K$. If $O=\pi^{-1}(O_Q)$, $U=\pi^{-1}(T)$,
$\psi=\chi\pi$ and $h(z)=\iota h_Q(\pi z)$ for an orthogonal
projection $\pi:H\to Q$ and inclusion $\iota:Q\to H$, $V$ can
be chosen as a preimage of an open subset of $Q$; the resulting
adjustment factors over $\pi$ and fixes $Q^\perp$.
\end{lemma}
\begin{proof}
Take $V=O\setminus(\operatorname{supp}_O\psi\setminus U)$.
The removed set is relatively closed in $O$, and disjoint from
$K$. Closed support in $V$ is contained in
$V\cap\operatorname{supp}_O\psi$, hence in $U$.
Every point of $V\setminus U$ therefore has a neighborhood in
$V$ on which $\psi$ vanishes identically. This proves smoothness
across the boundary as well as equality to the identity there.
In the factored case, use
$V_Q=O_Q\setminus(\operatorname{supp}_{O_Q}\chi\setminus T)$
and $V=\pi^{-1}(V_Q)$. Orthogonal projection is open and onto,
so the support of $\chi\pi$ is the preimage of the support of
$\chi$; thus the hypothesis ensures $K\subset V$.
The formula is then $(q,q^\perp)\mapsto
(q+\chi(q)h_Q(q),q^\perp)$ on its open domain.
\end{proof}
Compactness is not needed in CFS18 itself. It is needed in the next
fixed-data stability argument. A nonzero-value condition alone cannot
replace the closed-support hypothesis: for $O=\R$, $K=\{0\}$,
$U=(0,\infty)$ and $\psi(t)=e^{-1/t^2}$ for $t>0$, zero otherwise,
all nonzero values on $K$ lie in $U$ vacuously, but
$0\in K\cap\operatorname{supp}\psi\setminus U$.
The required domain is a neighborhood of the actual compact image;
no global ambient extension is inferred from the source formulas.

\begin{proposition}[Threshold gates and fixed-data support stability; CFS19]
\label{prop:fibration-cutoff-gate-stability}
Let $M$ be compact, $F:M\to O\subset H$ continuous, $\rho>0$
continuous, and $A\subset M$ open. Let $G:O\to\R$ be
continuous and let a smooth cutoff $\psi$ obey
\[
 \operatorname{supp}_O\psi\subset\{G\ge1/2\},\qquad
 G(F(p))=0\quad(p\notin A).
\]
There is a number $\kappa>0$ such that every map $f$ with
$|f(p)-F(p)|\le\kappa\rho(p)$ has image in $O$ and satisfies
\[
 f(p)\in\operatorname{supp}_O\psi\quad\Longrightarrow\quad p\in A.
\]
This is a FIXED-DATA existence statement: $\kappa$ may depend on
$F,G,O,\rho$ and the entire carrier. If additionally the original
cloud over $A$ satisfies CFS16 and $\kappa\le3\Sigma/10$, then
\[
 f(M)\cap\operatorname{supp}_O\psi
 \subset \pi_j^{-1}N_{r_j}(S_j).
\]
CFS18 therefore supplies the required smooth adjustment near $f(M)$.
\end{proposition}
\begin{proof}
The compact set $F(M)$ has a closed positive-radius Euclidean
neighborhood contained in $O$. On a possibly smaller such compact
neighborhood, $G$ is uniformly continuous. Choose $\delta>0$ such
that $|z-F(p)|\le\delta$ stays in $O$ and changes $G$ by less
than $1/4$. Choose $\kappa>0$ with
$\kappa\max_M\rho<\delta$. If $p\notin A$ then
$G(f(p))<1/4$, so $f(p)$ cannot be in the closed support.
The last inclusion follows from CFS16 at the same original
preimage $p$, and CFS18 applies with the projection tube as $U$
intersected with $O$. Empty carriers are vacuous.
\end{proof}
For KL (13.9), (13.25) and (13.41), $G$ is the finite
aggregate inside the outer function $1-\Phi_{1/2,1}$.
Continuity implies the displayed CLOSED-support inclusion, even
at a point where the outer cutoff itself is zero.
For the first and slim formulas, each summand contains a positive
marker gate and a ratio cutoff. On the original image a positive
summand forces $|\eta_i|<6.5$ or
$|\eta_i|<6.5\cdot10^5\Delta$, respectively. Thus the aggregate
is zero outside the corresponding open threshold-$7$ domain.
For the edge formula the calculation (13.30)--(13.32),
printed76/PDF71, makes EVERY summand zero outside the threshold-
$6.9\Delta$ region in both $|\eta_i|$ and $\eta_{E'}$.
It includes the retained-marker branch where the $E'$ output
vanishes; the ratio alone is not a substitute.
The positive scale coordinate restricts its ambient domain to
$\{x_\rho>0\}$, which contains the original compact image.

At the first stage $f=F$, so these gates give closed-support
localization without a perturbation tolerance. For stages two and
three CFS19 gives a qualitative tolerance for each fixed setup.
It does NOT prove a tolerance depending only on the predecessors
allowed by FC44, nor the required uniform cutoff derivative bound
on the perturbed image. The source plateaux and their stability
also remain separate inputs: support control alone does not prove
that a cutoff remains EXACTLY one. Choosing a new arbitrary cutoff
on $M$ would forfeit the required ambient postcomposition and is
not used here. The slim formula depends only on $Q_3$ coordinates,
so CFS18 preserves its factorization over $Q_2\supset Q_3$.

\begin{proposition}[Explicit three-stage tolerance budgets; CFS20]
\label{prop:fibration-adjustment-budget-order}
For each stage assume the original cloud hypotheses of CFS15,
the normal error $\nu_j\le\Gamma_j$, a rank margin $\mu_j>0$,
and the cutoff/plateau inputs of CFS17--CFS18 with derivative
constant $b_j$. Suppose $\|D\mathcal E^0\|\le L_0$.
Fix a target $0<c_j\le1$ and define
\[
 \alpha_j=\min\left\{
 \frac{c_j}{16(1+b_j)(1+L_0)},
 \frac{\mu_j}{8(1+L_0)}\right\}.
\]
Choose $\Gamma_j$ small enough that
$\Gamma_j\le c_j/16$ and
$\varepsilon_j=\Xi_j(\Gamma_j)\le\min\{1/10,\alpha_j\}$.
Then choose $0<\Sigma_j\le\min\{1/2,\varepsilon_j/640\}$
and require the preceding cumulative value and derivative errors
$E,H_0$ to be at most
\[
 t_j=\min\{3\Sigma_j/10,\alpha_j,1\}.
\]
The adjusted map has cumulative value error at most $3c_j/16$
and derivative error at most $3c_j/8$, measured relative to
$\mathcal E^0$ as in CFS17. It is a submersion to $W_j^0$ on
the supplied open plateau. In particular both errors are strictly
less than $c_j$.
For three stages these numeric conditions can be imposed in the order
\[
 c_3\ ;\ \Gamma_3,\varepsilon_3,\Sigma_3\ ;\ c_2\ ;\
 \Gamma_2,\varepsilon_2,\Sigma_2\ ;\ c_1\ ;\
 \Gamma_1,\varepsilon_1,\Sigma_1,
\]
by requiring $c_2\le t_3$ and $c_1\le t_2$; the first stage
has $E=H_0=0$. Any additional positive cutoff stability tolerance
can also be included by taking a minimum, PROVIDED its dependence
is confined to already chosen data.
\end{proposition}
\begin{proof}
CFS15 allows the stated choice of $\Gamma_j$, since its modulus
tends to zero. Suppress $j$ and write $\alpha=\alpha_j$.
CFS17 gives, using $B=5/3$, $\Sigma\le1/2$ and
$\varepsilon\le1/10$,
\[
 a\le(5/6)\varepsilon+(11/10)E
 \le(29/15)\alpha<2\alpha.
\]
Consequently $E+a\le3\alpha\le3c/16$. Since $H_0\le1$,
\begin{align*}
 ab(L_0+H_0)+\varepsilon(L_0+H_0)+\nu+2H_0
 &\le2\alpha b(L_0+1)+\alpha(L_0+1)+c/16+2\alpha\\
 &\le c/8+c/16+c/16+c/8=3c/8.
\end{align*}
The rank perturbation is at most
$\alpha(L_0+2)\le\mu/4<\mu$, so CFS17 applies.
Now choose stage-three parameters first, then a positive $c_2$
below $t_3$ and any admissible stage-three cutoff bound. Choose
stage-two parameters, then a positive $c_1$ below $t_2$ and its
admissible cutoff bound. Finally choose stage-one parameters.
Execution proceeds in the opposite order. Each completed stage
has cumulative errors smaller than its target, so induction places
the next input in its half-tube and yields its claimed estimates.
All minima used in this numeric argument are of finitely many
positive bounds fixed earlier.
\end{proof}
The new bounds $c_2\le3\Sigma_3/10$ and
$c_1\le3\Sigma_2/10$ respect KL (15.4). In contrast, the
qualitative $\kappa$ from CFS19 depends on a fully constructed
original map, which uses parameters chosen later in (15.4).
It cannot yet be inserted into that order. Thus CFS20 closes the
numeric tube/perturbation subproblem, conditional on uniform
cutoff producers; it does not close FC44's entire register.
Nor does submersion to an intermediate $W_j^0$ prove the final
bases are embedded or preserve fibers after later adjustments.
Those are FC33's one-sheet and retained-coordinate obligations.

\paragraph{Revision124 frontier (advanced below).}
Prove the uniform support and plateau estimates for the explicit
edge/slim gates, with constants fixed at their permitted levels;
finish the actual FC23/FC27 cloud, plane and every-preimage
bindings. Then apply CFS16--CFS20 and complete FC33's original-domain
one-sheet argument, proper restrictions and global assembly.
The historical general FC28 remains refuted, the strengthened
CFS14 remains conditional on its actual producers, and H1--H4,
the accepted-PC foundation binding and Lean implementation remain open.

\subsection{Uniform cutoff formulas and preservation of zero markers}
\label{sec:fibration-uniform-cutoff-repair}
Revision125 supplies uniform cutoff calculus under an explicit new
marker condition and a sufficient spectral locality criterion for that
condition. These project formulas replace the auxiliary expressions
(13.25) and (13.41); they retain FC30's thresholds $6,7$ and the required
ambient factorization. They do not assert that the preceding actual
adjustments already satisfy the new condition.

Fix a smooth nondecreasing $\chi:\R\to[0,1]$, zero on
$(-\infty,0]$, one on $[1,\infty)$, and strictly between zero and
one on $(0,1)$. Put $P\ge\max\{1,\|\chi'\|_\infty\}$ and
\[
 \chi_{a,b}(s)=\chi\bigl((s-a)/(b-a)\bigr),\qquad
 q(s)=1-\chi_{6.2,6.8}(s),\qquad
 T(s)=s\chi_{1/16,1/8}(s).
\]
Thus $T=0$ on $(-\infty,1/16]$, $T(s)=s$ for $s\ge1/8$,
$0\le T(s)\le s$ for $s\ge0$, and $|T'|\le3P$.
For the original edge cutoffs use $1-\chi$ as the source profile.
This uses the endpoint convention of KL Section2.2, printed19/PDF14.
Its literal (2.1), $\phi(a+(b-a)x)$, disagrees with the endpoints
stated immediately below; the explicit rescaling here implements those
endpoints. This is a project correction, not an author-listed erratum.

\begin{example}[What a scale-relative perturbation does not imply; CFS21]
\label{ex:fibration-marker-activation-obstruction}
A bound $|f-F|\le e\rho$, however small the positive uniform $e$,
does not by itself preserve cutoff support after discarding the scale
coordinate. It also need not preserve the exact plateau of KL (13.41),
even for one comparable-radius block.
\end{example}
\begin{proof}
For the first assertion take two points $p,q$ with
$\rho(p)=1$, $\rho(q)=R$, where $0<R<e$. Give a single retained
block the original values $F_i(p)=(0,0)$ and $F_i(q)=(0,R)$.
The original core contains $q$ but not $p$, with coordinate zero at
$q$. Set $f_i(p)=(0,R)$ and $f(q)=F(q)$; leave every other
coordinate, including the scale coordinate, unchanged.
The error bound holds. At the only positive original marker, scale
comparability is exact and the original multiplicity is one.
Yet any cutoff depending only on this retained block and equal to
one at $f(q)$ is also one at $f(p)$. It cannot have the required
support. This is an algebraic counterexample to an inference from
the listed data, not an example of the actual KL geometric adjustments.

For the second assertion use one original block $(0,R)$, and change
only its marker to $(1-h)R$, with $0<h<1/2$. In (13.41) the
axis factor is one, the inner marker factor is
$u=\chi_{1/2,1}(1-h)\in(0,1)$, and the outer factor is
$\chi_{1/2,1}(u)<1$. The perturbation $hR$ can be arbitrarily
small. Thus qualitative closeness and an original value exactly one
are not a proof of an exact perturbed plateau for that literal formula.
\end{proof}
Neither example refutes the collapse theorem or a statement restricted
to the actual adjustments. They show why a new marker hypothesis and
strict plateau buffers must be checked. There is a further ambient
calculus issue: differentiating the untruncated sum
$\sum_i x_i''/R_i$ counts even coordinates whose original markers are
zero. Its differential has coefficients $1/R_i$. The truncated $T$
above makes such coordinates locally inactive in the following proofs.

\begin{lemma}[Uniform one-axis cutoff with a zero-marker bound; CFS22]
\label{lem:fibration-uniform-axis-cutoff}
Let $I$ be finite, $R_i>0$, $\ell\ge1$, and let the original
blocks of $F=\mathcal E^0$ be
$F_i(p)=(R_i\eta_i(p)\zeta_i(p),R_i\zeta_i(p))$, with
$0\le\zeta_i\le1$ and the zero-extension convention of FC01.
Assume at each $p$ at most $N\ge1$ markers are positive, and
\[
 \zeta_i(p)>0\ \Longrightarrow\
 3R_i/4\le\rho(p)\le5R_i/4,\quad |\eta_i(p)|\le9\ell.
\]
Assume $\zeta_i=1$ wherever the original coordinate is defined and
$|\eta_i|<6\ell$. Write $z_i=(u_i,v_i)$ for ambient blocks.
Set
\[
 \kappa=\frac1{1000(N+1)P^2},\qquad
 C=10^4(N+1)^2P^4.
\]
Suppose the preceding map $f$ satisfies
\begin{gather*}
 |f(p)-F(p)|\le(4\kappa/5)\rho(p),\\
 \zeta_i(p)=0\ \Longrightarrow\ |v_i(f(p))|\le R_i/32.
                                                        \tag{ZM}
\end{gather*}
Define, using smooth zero extension when $v_i\le0$,
\[
 a_i(z)=\chi_{1/2,3/4}(v_i/R_i)\,
 q\bigl(|u_i|/(\ell v_i)\bigr),\qquad
 \psi_s(z)=\chi_{1/4,1/2}\left(\sum_i a_i(z)\right).
\]
Then $\psi_s$ is smooth, equals one on $f$ of the original
threshold-$6\ell$ union, and its CLOSED ambient support along
$f(M)$ comes only from the original threshold-$7\ell$ union.
On every segment $z_t=(1-t)F(p)+tf(p)$, $0\le t\le1$,
\[
 \|D\psi_s(z_t)\|\le C/\rho(p).
\]
Constants are independent of $\ell$, total index count and ambient
dimension. The formula depends only on the indicated blocks; for the
slim family $\ell=10^5\Delta$ it factors through $Q_3$.
\end{lemma}
\begin{proof}
The marker gate vanishes on a neighborhood of $v_i\le0$, so
zero extension is smooth. The apparent singularity of $|u_i|$ at
zero is harmless because $q$ is constant near zero.
For an originally positive marker, the error in either block
component is at most $\kappa R_i$, along the whole segment.
For an originally zero marker, (ZM) keeps $v_i/R_i\le1/32$;
its gate and all its derivatives vanish on an ambient neighborhood
of each segment point. Only the at most $N$ original positive
indices can contribute.

For a plateau witness $\zeta_i=1$, $|\eta_i|<6\ell$ gives
$v_i/R_i\ge1-\kappa>3/4$ and
\[
 \frac{|u_i|}{\ell v_i}
 \le\frac{6+\kappa/\ell}{1-\kappa}<6.2.
\]
Thus $a_i=1$ exactly and $\psi_s=1$.
Whenever the marker gate can be nonzero, $v_i/R_i\ge1/2$,
and the reverse ratio estimate is
\[
 \left|\frac{|u_i|}{\ell v_i}-\frac{|\eta_i|}{\ell}\right|
 \le20\kappa.
\]
If $p$ is outside the original threshold-$7\ell$ union, every
originally positive index therefore has either a zero marker gate,
or ratio at least $7-20\kappa>6.8$.
This also holds with a strict margin around the possible marker
transition at $v_i/R_i=1/2$ (use $v_i/R_i\ge0.49$ there).
Together with the strict $1/32$ bound for originally zero markers,
it makes all summands identically zero in an ambient neighborhood of
$f(p)$. Finiteness permits a common such neighborhood. This proves
CLOSED-support localization, not only a zero value at $f(p)$.

For derivatives at a contributing segment point, the marker has
$v_i\ge R_i/2$. Where the axis cutoff or its derivative contributes,
$|u_i|/(\ell v_i)\le6.8$. The ratio differential has norm at most
$16/R_i$, the marker differential at most $4P/R_i$, and
$\|Da_i\|\le36P/R_i\le45P/\rho(p)$.
Summing at most $N$ terms and using the outer derivative $4P$
gives $\|D\psi_s\|\le180NP^2/\rho(p)\le C/\rho(p)$.
Outside these supports the derivative is zero.
\end{proof}
The bound on zero markers in (ZM) is a NEW input. Positive-marker
scale comparability alone does not supply it. It asks only for the
scalar marker, not a bound on the entire block at its possibly tiny
radius; CFS21 shows it cannot be dropped.

\begin{proposition}[Uniform buffered edge cutoff; CFS23]
\label{prop:fibration-uniform-edge-cutoff}
For the edge family retain the positive-marker count and scale bounds
of CFS22, take $\ell=\Delta\ge1$, and assume (ZM) with the same
$N,P,\kappa,C$. On the original image write $t=\eta_{E'}\ge0$,
$m_i^0=\zeta_i$ and
\[
 Z_0=\chi_{1/2,1}\left(\sum_i m_i^0\right),\qquad
 z_0=\zeta_{E'}=h(t/\Delta)Z_0,
\]
where $h$ is the source profile, zero outside $(1/5,9)$ and equal
to $g=1-\chi_{8,9}$ on $[3/10,\infty)$.
The original variable block is $(\rho t z_0,\rho z_0)$.
For $|\eta_i|<8\Delta$ assume the actual edge identity
$m_i^0=g(t/\Delta)$, with $m_i^0=1$ when also $t<6\Delta$.
These identities are those of KL (9.19), (9.28)--(9.29).
The function $t$ need only be evaluated on those original domains;
where $z_0=0$ its ambient block is defined by zero extension.

On $O=\{x_\rho>0\}$ put $m_i=v_i/R_i$, $v=x_{E'}''/x_\rho$,
\[
 Z=\chi_{1/2,1}\left(\sum_i T(m_i)\right),\qquad
 B_E=\chi_{1/8,1/4}(v)\,
 q\bigl(|x_{E'}'|/(\Delta x_{E'}'')\bigr),
\]
where $B_E=0$ near $x_{E'}''\le0$, and define
\[
 d_i=m_iZ-v,\qquad b_i=\chi_{1/8,1/4}(d_i),\qquad
 \psi_e=\chi_{1/4,1/2}\left(\sum_i a_i(B_E+b_i)\right).
\]
Then $\psi_e:O\to[0,1]$ is smooth. On $f(M)$ it equals one on
the original union $\{|\eta_i|<6\Delta,t<6\Delta\}$, and its
CLOSED support comes only from the original union
$\{|\eta_i|<7\Delta,t<7\Delta\}$. On every preceding segment,
\[
 \|D\psi_e\|\le C/\rho(p).
\]
Thus the support, derivative and EXACT plateau conclusions are uniform
before $\Delta$ is chosen, conditional on the displayed original data
and (ZM). The zero-output branch of $E'$ is included.
\end{proposition}
\begin{proof}
The scale coordinate along the segment is between
$(1-\kappa)\rho$ and $(1+\kappa)\rho$, hence positive.
Smoothness follows from the positive-marker gates as in CFS22.
Originally zero edge markers stay below $1/32<1/16$, so $T$ and
$a_i$ vanish on their neighborhoods. At most $N$ indices contribute
to either sum, including its ambient differential.
For originally positive indices, $|m_i-m_i^0|\le\kappa$.
The variable marker satisfies $|v-z_0|\le3\kappa$. If $Z^0$ denotes
the NEW truncated $Z$ evaluated at $F(p)$, then
\[
 Z^0\le Z_0,\qquad |Z-Z^0|\le6NP^2\kappa.
                                                        \tag{ZT}
\]
The first inequality uses $0\le T(m_i^0)\le m_i^0$ and monotonicity;
the second uses $|T'|\le3P$ and $|\chi_{1/2,1}'|\le2P$.
No equality between the truncated and source aggregates is assumed.

On the inner region choose its full-marker witness $i$.
Then $a_i=1$. Also $m_i\ge1-\kappa$ and $T(m_i)=m_i$, so
$Z\ge1-2P\kappa$. If $z_0\ge3/8$, the variable marker has
$v>1/4$, and its ratio differs from $t/\Delta$ by at most
$100\kappa$, hence is less than $6.2$. Thus $B_E=1$.
If $z_0<3/8$, then
\[
 d_i\ge(1-\kappa)(1-2P\kappa)-(3/8+3\kappa)>1/4,
\]
so $b_i=1$. In either case the nonnegative aggregate is at least
one and $\psi_e=1$. This also proves the case $z_0=0$ without
dividing by its zero original marker.

Suppose $p$ is outside the outer union. As in CFS22, indices with
original zero marker or $|\eta_i|\ge7\Delta$ give $a_i=0$
on an ambient neighborhood of $f(p)$. For every remaining index,
$t\ge7\Delta$ and $|\eta_i|<7\Delta<8\Delta$, so
\[
 m_i^0=g(t/\Delta),\qquad z_0=m_i^0Z_0,
 \qquad m_i^0Z^0-z_0\le0.
\]
By (ZT) and the marker estimates,
\[
 d_i\le(4+6NP^2)\kappa<1/8.
\]
Thus $b_i$ vanishes with a strict margin.
If the gate in $B_E$ could contribute, $v\ge1/8$ implies
$x_{E'}''\ge(1-\kappa)\rho/8$. Then $z_0>0$ and
$7\Delta\le t<9\Delta$, and direct division gives
\[
 \left|\frac{|x_{E'}'|}{\Delta x_{E'}''}-\frac t\Delta\right|
 \le100\kappa.
\]
The ratio is therefore greater than $6.8$; near the threshold
$v=1/8$ the same strict conclusion holds with $v\ge1/9$.
Hence $B_E=0$ on a neighborhood as well. Every summand is locally
zero at $f(p)$, proving CLOSED-support localization.

For completeness, all the following differential bounds are in the
FULL ambient operator norm at a segment point, in units $\rho^{-1}$:
\[
 \|Da_i\|\le50P/\rho,\quad \|Dv\|\le4/\rho,\quad
 \|DZ\|\le8NP^2/\rho,\quad \|DB_E\|\le200P/\rho.
\]
For the last bound, on its contributing support
$x_{E'}''\ge(1-\kappa)\rho/8$ and the normalized ratio is at most
$6.8$. Its differential is at most $72/\rho$; multiply by $2P$
and add the marker term $8P\cdot4/\rho$.
Since $|m_i|\le2$ for contributing indices,
\[
 \|Dd_i\|\le22(N+1)P^2/\rho,\qquad
 \|D(a_i(B_E+b_i))\|\le500(N+1)P^3/\rho.
\]
There are at most $N$ such terms. Multiplication by the outer
$4P$ proves $\|D\psi_e\|\le2000N(N+1)P^4/\rho\le C/\rho$.
The powers of $\Delta$ cancel in the normalized ratios; no ambient
count or smallest radius is present.
\end{proof}
This replaces the raw marker sum in (13.24) by a truncated sum and
its signed fallback in (13.25) by a nonnegative saturated gate.
The comparison $Z^0\le Z_0$, not equality, is what excludes false
outer-region activation. The original map, its local cutoffs, cloud
images and final thresholds are unchanged. Only the auxiliary
adjustment cutoffs are replaced.

\begin{lemma}[Exact coordinate locality of the spectral smoothing; CFS24]
\label{lem:fibration-spectral-marker-locality}
Use the single smoothing construction CFS11--CFS14, with
$b=\varepsilon^{-1}$, core center $x$ and radius $r_x$.
Let $J:H\to V$ be orthogonal projection onto a fixed coordinate
subspace. Suppose EVERY selected center $x_i$ whose closed
$80br_i$ cutoff support meets $B(x,8br_x)$ satisfies
\[
 Jx_i=0,\qquad L_i\subset\ker J.
                                                        \tag{ZL}
\]
Then the global section has $J\eta(z)=Jz$ on $B(x,8br_x)$.
Every $w\in W\cap B(x,3br_x)$ has $Jw=0$, and
\[
 JP(z)=0\quad(z\in B(x,r_x)).
\]
Consequently the blended adjustment preserves a zero $J$ coordinate
at any such input. In particular scalar marker projections are allowed.
\end{lemma}
\begin{proof}
For every contributing plane, its normal projector $P_i$ is the
identity on $V$, and $JP_i=J$. The weighted average $O$ is
self-adjoint and is the identity on $V$, with no off-diagonal block
between $V$ and $V^\perp$. Its spectral projector $Q$ for the
cluster near one therefore satisfies $JQ=J$.
The weighted center $\mu$ has $J\mu=0$. Hence
$J\eta=JQ(z-\mu)=Jz$ throughout the stated ball.
The assertion about the zero set follows immediately.
By CFS14, for $z\in B(x,r_x)$,
\[
 |P(z)-x|\le|P(z)-P_{A_x}(z)|+|P_{A_x}(z)-x|
 <(1+\varepsilon)r_x<3br_x.
\]
Thus $JP(z)=0$. Finally, the projected coordinate of
$(1-\psi)z+\psi P(z)$ is $(1-\psi)Jz$, which is zero when
$Jz=0$.
\end{proof}
Condition (ZL) concerns ALL contributors on the whole buffered ball,
including their chosen planes. A zero coordinate at the core center,
a small cloud error, or vanishing at one preimage does not imply it.
For a stage operating in $Q_j$, apply the lemma inside that Euclidean
space; coordinates in $Q_j^\perp$ are retained exactly by FC32.
This gives a sufficient local mechanism for (ZM), not its actual
geometric producer.

\begin{corollary}[Uniform cutoff consumer with an explicit locality gate; CFS25]
\label{cor:fibration-uniform-cutoff-consumer}
Assume the actual original clouds and CFS15 smoothing hypotheses,
the positive-marker multiplicity and scale bounds above, and the
original edge identities. Suppose the actual preceding adjustments
supply (ZM) for the next edge/slim family. Replace their auxiliary
cutoffs by CFS23 and CFS22, respectively. Their support, exact
plateau and derivative requirements in FC30 then hold with a constant
$C=C(N,P)$ fixed before $\Delta$, provided the preceding cumulative
value error is at most $4\kappa/5$ and $3\Sigma_j/10$.
CFS18 supplies a smooth neighborhood of each actual image; the slim
adjustment still factors over $Q_2$. Use $b_j=C$ in CFS20 and add
$4\kappa/5$ to its preceding-error minimum. These numerical choices
respect KL (15.4).
A sufficient way to supply (ZM) is to prove (ZL) for every originally
zero marker at each earlier stage where that marker can change.
\end{corollary}
\begin{proof}
The new cutoffs localize the ORIGINAL point to the appropriate
threshold-$7$ core. CFS16 then puts the perturbed input in its
half-tube. Their closed-support conclusions permit CFS18, without
altering cutoff values or derivatives along the image. CFS22's
formula uses only slim blocks, so the factored case of CFS18 retains
FC32's $Q_2$ factorization. Their uniform differential bound is exactly
the $b_j/\rho$ input to CFS17--CFS20.
Fix $N,P$ first. Then $C$ and $\kappa>0$ are fixed at the source's
multiplicity/cutoff-constant level, before $\Gamma_j,\Sigma_j$ and
$\Delta$. The additional positive bound can be included when choosing
$c_2$ for stage three and $c_1$ for stage two. No smallest chart
radius or compact-carrier modulus enters this choice.
For the sufficient locality route, start with every originally zero
marker equal to zero. Outside an earlier adjustment support the map
is unchanged; inside it, CFS24 with its original core ball preserves
that coordinate under (ZL). A coordinate outside the adjusted summand
is unchanged by FC32. Induction preserves the zero marker, which is
stronger than (ZM).
\end{proof}
The analytic cutoff construction and its parameter dependence are now
written under explicit data, replacing revision124's mere fixed-data
compactness tolerance. The new unresolved producer is concrete:
prove (ZM), or the stronger (ZL), for the actual cloud planes and all
selected centers influencing each preceding projection. This must be
bound to FC23/FC27's original-image localization, not inferred from
$C^1$ closeness. Those cloud producers, FC33's final one-sheet/fiber
preservation, proper restrictions, the full FC44 register and global
assembly remain open. No source theorem or parent contract is marked
complete, and no accepted-PC or Lean binding is claimed.

\subsection{Actual support scales and coordinate-preserving cloud planes}
\label{sec:fibration-actual-marker-control}
Revision144 removes the separate zero-marker hypothesis in CFS25 by
fixing the choice of the auxiliary graph planes. The construction below
uses the original local packets of revision143 and their actual domains.
It does not alter $F=\mathcal E^0$, its cutoffs, image sets or radii.
The rough graph approximation and coverage required by FC23/FC27 are
still explicit producers; once supplied, the plane choice below preserves
their estimates and needs no additional locality assumption. Thus this
is progress on the actual geometric producer, not a promotion of the
whole projected-cloud or final-fibration contract.

\begin{lemma}[Support-scale bounds for every original preimage; CFS26]
\label{lem:fibration-actual-support-scale}
Use the actual closed-carrier local packets, and put
$I_+=I_2\cup I_e\cup I_s$, $R_i=\rho(p_i)$ and
$D_*=10^6\Delta$, where $\Delta\ge1$.
Choose the Lipschitz constant of $\rho$ so that
$\Lambda D_*\le1/4$. On the entire original coordinate domain
$U_i$, and on the closed support of its zero-extended cutoff, one has
\[
                 3R_i/4\le\rho\le5R_i/4.                 \tag{AS}
\]
This holds for every preimage of an image point with positive $i$
marker. Every point of each original core or enlarged cloud has a
retained full constant marker. Thus FC26 and CFS07 apply to ALL
preimage choices in each projected cloud; for any two such choices
over the same image their scales have ratio in $[3/5,5/3]$.
\end{lemma}
\begin{proof}
The actual domains have radii $200R_i$, $100\Delta R_i$ and
$10^6\Delta R_i$ for the two-stratum, edge and slim families,
respectively. Each is at most $D_*R_i$. The Lipschitz inequality
$|\rho(q)-R_i|\le\Lambda d(q,p_i)$ gives (AS) on the domain.
The local packets place the closed cutoff supports inside those
domains; alternatively continuity gives the same inequalities on
their closures. Positive markers mean positive original cutoffs,
so this argument concerns every preimage, not a selected one.
The threshold-$8$ enlargement in the two-stratum family lies where
its cutoff is one. In the edge family BOTH coordinates are below
$8\Delta$, and both factors are one. In the slim family the
threshold is $8\cdot10^5\Delta$, again in the plateau. These
blocks are retained in the respective $Q_j$.
Apply (AS) at their full marker and then FC26 and CFS07.
For $Q_1$ the scale coordinate gives the stronger exact identity.
For $Q_2,Q_3$ it has been discarded: no descent of $\rho$ is used.
\end{proof}

\begin{lemma}[Pruning identically zero blocks before choosing planes; CFS27]
\label{lem:fibration-pruned-model-planes}
For a reference chart $a\in I_+$ define the orthogonal coordinate
projection $K_a$ to delete every retained $I_+$ block $i$ with
$R_i\le R_a/2$. Keep all other coordinates, including the reference
block and every zero-stratum or variable-radius block that is present.
Then on $U_a$,
\[
                         K_a\pi_jF=\pi_jF.              \tag{PR}
\]
Use $R_a$ units and a local comparison patch where the reference
marker is full. If FC23/FC27 supplies a graph
$\Phi_a(u)=(u,h_a(u))$, replace it, BEFORE smoothing, by
$\widehat\Phi_a=K_a\Phi_a$. This is still a graph over the same
reference coordinate. Value and $C^1$ comparison errors on that patch,
and upper derivative and Hessian bounds, do not increase.
At $x=\pi_jF(q)$ choose the actual affine plane to be
\[
 A_x=x+L_x,\qquad
 L_x=\operatorname{im}D\widehat\Phi_a(\eta_a(q)).       \tag{PP}
\]
Every deleted block is zero at $x$ and on $L_x$.
Graph-to-cloud and derivative-rank estimates may be established using
this graph with the same error budgets as before pruning.
\end{lemma}
\begin{proof}
For $q\in U_a$, (AS) gives $\rho(q)\ge3R_a/4$.
A positive deleted marker would instead give
$\rho(q)\le5R_i/4\le5R_a/8$, a contradiction. Its entire block,
which is multiplied by that cutoff, is therefore zero. This proves
(PR) on the open domain, including equality of derivatives there.
The strict scale gap also gives local vanishing at any domain point;
this is stronger than vanishing only at the selected center.

The reference block is not deleted, so the first component remains
the identity. In particular the new tangent map is injective and
has lower norm bound one. Orthogonal projection contracts the
original comparison error and all of its derivatives, by (PR).
It contracts the upper derivative and Hessian norms of $\Phi_a$.
For clarity, this does not assert that projecting an arbitrary plane
preserves its rank: rank follows from the retained identity component.

Every source point used in the local comparison lies in $U_a$,
so its image is fixed by $K_a$. Consequently the two pointwise
comparison hypotheses of FC25 remain valid with $\widehat\Phi_a$
and the SAME image points; its proof gives the new tangent-plane
cloud test. All-preimage localization and exact coordinate coverage
remain prerequisites of that application, not consequences of (PR).
For the differential estimate write
$D(\pi_jF)=T D\eta_a+E$, where
$T=D\widehat\Phi_a$, $\|E\|\le e$ in these normalized units.
Then the component normal to $\operatorname{im}T$ has norm at
most $e$. If $D\eta_a$ has least nonzero singular value $\lambda$,
$T$'s identity component implies the same lower bound for
$TD\eta_a$. Perturbing by at most $e<\lambda$ gives a rank-$k$
projected derivative with lower bound $\lambda-e$; the upper bound
uses the unchanged upper bounds. This proves the claimed preservation
of the estimates used by FC06, rather than assuming invariance of
orthogonal projections. Finally (PR) fixes $x$, and differentiation
of $K_a\Phi_a$ puts its tangent in every deleted coordinate kernel.
\end{proof}
The rule is a project choice of auxiliary models and planes. In the
first source model, KL (12.15) already omits blocks whose supports
miss the comparison ball. The displayed rule makes the corresponding
vanishing explicit for every projected model as well. It must be
applied before CFS11--CFS14 construct their section and projection;
one cannot silently change the planes of an already constructed $W$.

\begin{lemma}[Every contributing center and plane has small-marker locality; CFS28]
\label{lem:fibration-all-contributors-locality}
Use the plane choice (PP) in stage $j$, with each reference chart
chosen at a preimage where its retained marker is full.
Let $b\ge1$, $0<\Sigma_j\le1/(640b)$,
$x=\pi_jF(p)\in S_j$, and $r_j=\Sigma_j\rho(p_x)$ with
arbitrary selected preimages, as above.
For ANY retained $I_+$ block $i$ with $R_i<\rho(p)/16$,
every selected smoothing center $u$ whose closed $80br_j(u)$-ball
meets $B(x,8br_j(x))$ satisfies
\[
                  J_i u=0,\qquad L_u\subset\ker J_i,    \tag{AL}
\]
where $J_i$ projects onto the whole $i$ block. This includes all
contributors on that buffered ball, not just those with positive
cutoff at one point, and allows different preimages for a radius
and for a graph model at the same center.
\end{lemma}
\begin{proof}
The meeting condition implies
\[
 |u-x|<80br_j(u)+8br_j(x)
       \le88b\max\{r_j(u),r_j(x)\}.
\]
CFS07 with $L=128b$ and CFS26 give
$r_j(u)/r_j(x)\ge3/5$.
Write $\sigma_x,\sigma_u$ for the two scales chosen in the radii.
The present preimage $p$ and the chosen one at $x$ satisfy
$\sigma_x\ge(3/5)\rho(p)$. If $q_u$ is the preimage chosen
for the graph model at $u$, then
$\rho(q_u)\ge(3/5)\sigma_u$.
The full reference marker $a=a_u$ at $q_u$ gives
$R_a\ge(4/5)\rho(q_u)$. Hence
\[
 R_a\ge\frac45\left(\frac35\right)^3\rho(p)
       =\frac{108}{625}\rho(p),\qquad
 \frac{R_i}{R_a}<\frac{625}{1728}<\frac12.             \tag{AM}
\]
The $i$ block is therefore deleted by $K_a$.
CFS27 proves both its vanishing at $u=\pi_jF(q_u)$ and its
vanishing on $L_u$. Every center in the entire indicated support
list obeys the same argument. No uniqueness of preimages, ambient
dimension bound, smallest chart radius or continuity of the chosen
radius was used. The exact first-image radius only improves the
inequalities; the uniform constants cover all three stages.
\end{proof}

\begin{lemma}[Actual preservation of small originally zero markers; CFS29]
\label{lem:fibration-small-marker-preservation}
Construct each stage with the planes (PP), the CFS15 smoothing,
and $\Sigma_j\le1/(640b_j^{\rm sm})$, where
$b_j^{\rm sm}=\varepsilon_j^{-1}$ is its smoothing buffer.
Suppose its cutoff localizes the ORIGINAL point to $A_j$ on its
closed support, and the preceding value error is at most
$3\Sigma_j\rho(p)/10$.
If an input map has zero $i$ marker whenever
$\zeta_i(p)=0$ and $R_i<\rho(p)/16$, then its blended adjustment
has the same property. The threshold uses the original $\rho(p)$.
\end{lemma}
\begin{proof}
Outside the cutoff support the adjustment is the identity.
On its support use $x=\pi_jF(p)$ from the original core.
CFS16 puts $\pi_jf(p)$ in $B(x,r_j(x))$.
For a retained block CFS28 verifies EVERY center and plane
hypothesis of CFS24, with $J$ the scalar marker projection.
Thus the projected output marker is zero. The blend of two zero
markers is zero. A block outside $Q_j$ is unchanged by FC32.
This proves the assertion at every point. No condition is placed
on the possibly changed scale coordinate of the input map.
\end{proof}

\begin{corollary}[The zero-marker bound from the actual plane rule; CFS30]
\label{cor:fibration-actual-zero-marker-bound}
Start with $F$ and perform any initial segment of the three stages
with CFS29's support and error bounds. If the cumulative value error
of the resulting map $f$ satisfies $|f-F|\le E\rho$,
$E\le1/512$, then for EVERY $i\in I_+$,
\[
 \zeta_i(p)=0\quad\Longrightarrow\quad
                         |v_i(f(p))|\le R_i/32.        \tag{AZM}
\]
The same bound holds on the straight segment from $F(p)$ to $f(p)$.
It has no dependence on the number of charts, on $\Delta$ or on
the global range of their radii.
\end{corollary}
\begin{proof}
If $R_i<\rho(p)/16$, induction with CFS29, starting from the
original zero marker, gives exact vanishing at every stage.
If $R_i\ge\rho(p)/16$, orthogonal projection onto that scalar
coordinate and $v_i(F(p))=0$ give
$|v_i(f(p))|\le E\rho(p)\le16E R_i\le R_i/32$.
The bound is preserved on the segment by convexity of the scalar
interval $[-R_i/32,R_i/32]$. These two cases exhaust the indices.
No analogous claim is made for the zero-stratum or variable-radius
blocks, which are not the inputs to CFS22--CFS23.
\end{proof}

\begin{proposition}[Sequential cutoff construction without a marker assumption; CFS31]
\label{prop:fibration-marker-cutoff-sequence}
Assume the rough original graph, coverage, normal and rank estimates
required by FC23/FC27 and CFS20, and use the pruned models of CFS27.
Retain the actual local cutoffs and original edge identities, with
the support multiplicity bound $N$ and normalized profile bound $P$.
The first source cutoff, followed by the auxiliary edge and slim
cutoffs of CFS23 and CFS22, gives all of CFS25's cutoff conclusions
without assuming (ZM) or (ZL). Their common ambient differential
bound is $C/\rho$, $C=10000(N+1)^2P^4$.
In CFS20 take its derivative constants $b_j=C$ and impose
\[
 \begin{gathered}
 c_3\le1/512,\qquad
 c_2\le\min\{t_3,4\kappa/5,1/512\},\\
 c_1\le\min\{t_2,4\kappa/5,1/512\},\qquad
 \kappa=\frac1{1000(N+1)P^2}.
 \end{gathered}                                        \tag{MB}
\]
These choices and $\Lambda10^6\Delta\le1/4$ are compatible
with the previously stated order. This closes the marker and cutoff
subproblem relative to the rough-cloud producers; it does not prove
those producers, FC33, proper boundary restrictions or the full FC44.
\end{proposition}
\begin{proof}
For $i\in I_2$ write its block $(u_i,v_i)$ and use the increasing
profiles of CFS22 to define the source first cutoff
\[
 a_i=\chi_{1/2,1}(v_i/R_i)
             [1-\chi_{6,13/2}(|u_i|/v_i)],\qquad
 \psi_1=\chi_{1/2,1}\left(\sum_{i\in I_2}a_i\right).
\]
Each summand is zero-extended near $v_i\le0$ and is smooth,
including at $u_i=0$, where the radial factor is constant.
At $F(p)$ an original $|\eta_i|<6$ supplies a full marker and
$a_i=1$, proving the exact plateau. If $F(p)$ is in the CLOSED
ambient support, some summand can contribute arbitrarily nearby;
finiteness gives an index with $v_i\ge R_i/2$ and
$|u_i|/v_i\le13/2$. At the original image this means
$|\eta_i|\le13/2<7$. Thus $p\in A_1$ with a strict buffer.
An originally zero marker has a neighborhood where its summand
vanishes, and only the at most $N$ originally positive markers can
contribute to the differential. When a derivative of $a_i$ is
nonzero, $v_i\ge R_i/2$ and the ratio is at most $13/2$;
the ratio differential is at most $15/R_i$.
The two profile derivatives give $\|Da_i\|\le32P/R_i$.
Using (AS) and the outer derivative bound $2P$ yields
$\|D\psi_1\|\le80NP^2/\rho\le C/\rho$ along $F(M)$.
CFS18 supplies its smooth adjustment domain.

Execute stage one with CFS20's budgets. CFS29 and then CFS30
give (AZM) for its output; its cumulative value error is less
than $c_1$. Therefore CFS23, now with its zero-marker input
PROVED, supplies the second cutoff with its closed-support,
exact-plateau and differential conclusions. The $t_2$ bound
places the second input in its projection tube. Execute stage two
and apply CFS29--CFS30 again. Its error is less than $c_2$, so
CFS22 supplies the third cutoff and $t_3$ supplies its tube.
Execute stage three; its cumulative error is less than $c_3$,
and CFS29--CFS30 remain applicable. This is a sequential induction:
no later cutoff is used to prove an earlier marker bound.
CFS18 handles each closed support; the slim formula still factors
through $Q_3\subset Q_2$ as required by FC32.

Fix $N,P$ first, so $C,\kappa$ and $1/512$ are available before
any $\Gamma_j,\Sigma_j$ or $\Delta$. CFS20 already requires
$\Sigma_j\le\varepsilon_j/640$, exactly the smoothing-buffer
bound in CFS29. Add (MB) to its positive finite minima, then choose
$\Delta$ and finally $\Lambda$ small enough for (AS) as well as
the previously required inequalities. These additional choices have
no later-dependent tolerance. The rough graph and rank moduli and
the remaining geometric register still require their own bindings.
\end{proof}

\paragraph{Current Chapter14 frontier after revision144.}
The every-center/every-plane locality issue is resolved by the actual
support-scale argument and a specified admissible choice of the model
planes. Ordinary $\rho$-scaled closeness is used only for comparable
markers; arbitrarily small markers are preserved exactly. The next
producer to bind is FC23/FC27's rough graph approximation, including
all original derivative-test domains and both coverage directions.
Then FC33 must exclude additional sheets and preserve whole fibers,
followed by proper restrictions, compatible faces and the full parameter
register. No parent contract, accepted-PC binding or Lean theorem is
promoted by this revision. This is an intermediate development, not
the requested final audit of the static assembly goal.

\subsection{Binding the original block map and the one-sheet bases}
\label{sec:fibration-original-map-and-one-sheet}
Revision145 first binds the actual profiles and the uniform derivative
constant to the completed local packets. It then proves the one-sheet
and later-image assertions from the rough graph estimates still required
by FC23/FC27. The latter remain prerequisites. The proof uses CFS13's
exclusion of all competing sheets in a buffered ball, so it does not
assume that a marked smoothed patch already equals an adjusted image.

\begin{proposition}[The source profiles on the actual local maps; CGP01]
\label{prop:fibration-source-profile-binding}
In LPA06's closed-carrier construction choose the global-map cutoffs
using the increasing profile $\chi$ of CFS22. Set
\[
 f(s)=1-\chi_{8,9}(|s|),\quad g(s)=1-\chi_{8,9}(s),\quad
 h(s)=\chi_{1/5,3/10}(s)g(s),\quad t=P/\rho,
\]
where $P$ is the ONE nonnegative physical distance smoothing supplied
by LFR33/LFR38. For the two-stratum, slim and edge families use
\[
 \begin{aligned}
 \zeta_i&=1-\chi_{8,9}(|\eta_i|)&& (i\in I_2),\\
 \zeta_i&=f(\eta_i/(10^5\Delta))&& (i\in I_s),\\
 \zeta_i&=f(\eta_i/\Delta)g(t/\Delta)&& (i\in I_e).
 \end{aligned}
\]
Use the original annular profile for $I_0$, and put
\[
 Z_0=\chi_{1/2,1}\left(\sum_{i\in I_e}\zeta_i\right),
 \qquad \zeta_{E'}=h(t/\Delta)Z_0.
\]
All blocks of FC01 extend smoothly by zero, with their closed supports
strictly inside the appropriate smooth domains. The original identities
required by CFS23 hold exactly. These choices preserve the local covering
plateaux and all local maps, models and proper bundle domains.
They are made BEFORE defining $F=\mathcal E^0$ and its clouds.
\end{proposition}
\begin{proof}
Take the local coordinate value tolerances at most $\Delta/100$ for
edge/slim charts and $1/10$ for circle charts. In the one common
physical smoothing choose its value error additionally below
$\min_M\rho/1000$. This is an allowed positive error on the fixed
compact carrier; it changes no earlier geometric quality threshold.
The Lipschitz-difference error is at most $1/100$.
LPA03/LFR07 put the circle cutoff support in its radius-$10$ ball.
FC18 puts the closed slim and edge supports in the radius
$950000\Delta R_i$ and $20\Delta R_i$ balls, respectively,
strictly inside the ORIGINAL coordinate domains.
Its value and coarse-border hypotheses are supplied by LFR19,
LFR20 and LFR32; decrease their already permitted errors as needed.
The annular support is contained in $\eta_i\in[1/5,9/10]$ and
its prescribed smooth shell, by LC67 and LPA05.

For an edge block, possible nonsmoothness of $t$ at the edge core
does not enter the formula: $g(t/\Delta)=1$ on a neighborhood
there. Where its derivative or transition matters, the support
enclosure, slow scale variation and the value bound for $P$ put
the point in LFR34's enlarged smooth set $C^+$.
For the $E'$ block, every point in its closed support has
$t/\Delta\in[1/5,9]$ and
$\sum\zeta_i\ge1/2$, hence belongs to a positive edge cutoff.
For such an index it lies in $B(p_i,20\Delta R_i)$; imposing
$100\Delta\Lambda<10^{-4}$ gives
\[
 .19\Delta<d_{\overline{E'}}/R_i<9.1\Delta.
\]
These bounds lie strictly inside LFR34's smooth height interval
$[\Delta/25,10.5\Delta]$. The sharper radius-$15\Delta R_i$
enclosure in FC18 leaves a neighborhood inside its spatial buffer.
Thus $t$ is smooth near this support. Outside that closed support
the block is locally zero. No nonsmooth distance is differentiated
merely because it is multiplied by a zero output.

All coordinate plateaux remain at their original threshold $8$.
For $|\eta_i|<8\Delta$ one has $\zeta_i=g(t/\Delta)$, and
$\zeta_{E'}=h(t/\Delta)Z_0$ is its exact shared expression.
Nonnegativity of $P$ gives $t\ge0$. These are CFS23's identities,
including the branch where the $E'$ block is zero.
The local-packet proofs sometimes chose a transition ending at $8.9$
to obtain an additional coordinate margin. Here the CLOSED level-$9$
support is still inside the larger metric-domain buffer just proved.
We choose these source profiles at the initial construction, rather
than replacing a cutoff of an already constructed global map.
The same argument proves their smooth zero extension.
\end{proof}

\begin{proposition}[An early derivative constant for the actual global map; CGP02]
\label{prop:fibration-actual-global-derivative}
Let $N\ge1$ be LPA06's early nonzero-family multiplicity bound, and
$P_0\ge1$ bound the first derivatives of the fixed unit profiles.
Retain the local-packet convention $\Delta\ge1$.
The original map in CGP01 satisfies
\[
                 \|DF\|\le L_0:=1000(N+2)P_0^2.        \tag{GD}
\]
This constant is fixed before $\Delta$, noncollapse, the total number
of charts and the global scale range. It is a produced input to
CFS20, not an additional derivative assumption.
\end{proposition}
\begin{proof}
Choose the adapted and radial Lipschitz errors below one, so that
$R_i|d\eta_i|\le2$ for each constant-radius coordinate; use
operator norm for the two-component coordinates.
On the edge supports,
\[
 \rho|dt|\le |dP|+t|d\rho|\le2,\qquad R_i|dt|\le3,
\]
after decreasing $\Delta\Lambda$, using CFS26. These estimates
apply where the corresponding factors vary; elsewhere they are
locally constant. For a circle or slim block, put $\ell=1$ or
$10^5\Delta$. Then $R_i|d\zeta_i|\le2P_0/\ell$ and
$|\eta_i|\le9\ell$ on its support. For an edge block the analogous
bound is $R_i|d\zeta_i|\le5P_0/\Delta$, with $|\eta_i|\le9\Delta$.
For a zero block its annular profile has derivative at most $20P_0$,
so $R_i|d\zeta_i|\le40P_0$ and $|\eta_i|\le9/10$.
Thus FC02's constant-radius bound can be taken to be
$B=2+80P_0$. There are at most $N+1$ active such blocks, since
the zero supports are disjoint.

On the $E'$ support, $\rho|dt|\le2$ and $0\le t\le9\Delta$.
Since $|h'|\le11P_0$, $|\chi_{1/2,1}'|\le2P_0$ and
$\rho/R_i\le5/4$ at every positive edge marker,
\[
 \rho|d\zeta_{E'}|
 \le\frac{22P_0+(25/2)NP_0^2}{\Delta}
 \le\frac{40(N+1)P_0^2}{\Delta}.
\]
The same bound includes closed-support points by smoothness.
Take $\Delta\Lambda\le1/10$. FC02's variable-block constant is
at most $500(N+1)P_0^2$. Its square-sum bound, $\Lambda\le1$,
and $B\le82P_0$ give (GD). Only the pointwise multiplicity enters
this estimate; the total target dimension is unrestricted.
\end{proof}

\begin{lemma}[Actual sections through the inner coordinate balls; CGP03]
\label{lem:fibration-actual-inner-sections}
Put $\ell_i=1,\Delta,10^5\Delta$ for $i\in I_2,I_e,I_s$,
respectively. For each such $i$ there is a continuous section
\[
 s_i:B(0,\tfrac{23}{4}\ell_i)\longrightarrow U_i,\qquad
 \eta_i s_i(a)=a,
\]
whose image lies in the original threshold-$6$ plateau of its
adjustment. For edge charts it also has $t<\Delta/100$.
All points of this image have a full $i$ marker and scale in
$[3R_i/4,5R_i/4]$.
\end{lemma}
\begin{proof}
For circles, take a fixed point of the fiber in the proper trivial
bundle of LPA03/LFR07 over $B(0,100)$. For slim charts use LFR20's
proper trivial bundle over $(-9\cdot10^5\Delta,9\cdot10^5\Delta)$.
Their restrictions give the stated sections and stay inside the
threshold-$6\ell_i$ plateau.

For an edge chart do not cite the disk bundle over $(-4\Delta,4\Delta)$
as if its base were larger. Use the SAME buffered product embedding
in LFR28's proof, on a neighborhood of
$[-6\Delta,6\Delta]\times\overline B_Z(z_0,9\Delta)$.
Its estimates (LFR28.2)--(LFR28.5), with the original sufficiently
small parameters, give along the core line $z=z_0$
\[
 |\eta_i j(a,z_0)-a|<2\mu\Delta,\quad
 \partial_a(\eta_i j)>3/4,\quad 0\le t j(a,z_0)<\Delta/100.
\]
Here $\eta_i$ is the tangential coordinate, whereas $t=P/\rho$
is the shared height. The last estimate follows also directly
from (LFR28.2), the coarse-border distance bound, the physical
smoothing error and $r(z_0)=0$. Decrease those permitted errors
once, BEFORE constructing the packets. Endpoint values at
$a=\pm(59/10)\Delta$ then enclose
$[-(23/4)\Delta,(23/4)\Delta]$. Strict monotonicity and the
intermediate value theorem give a continuous inverse in that
line segment. Composing with $j$ produces $s_i$.
The finite-model embedding has ample regularity; only continuity
is used here. These are extra fixed-buffer estimates in the SAME
compactness proof, with the same analytic data and sufficiently
small strong splitting error, not a new disk-bundle theorem on a
silently enlarged base. More explicitly, add the existence of this
section to the conclusions whose failure defines LFR28's counterexample
sequence. Its same common tail satisfies the three displayed estimates
and the endpoint inequalities, and hence constructs the section above.
It contradicts that failure just as it contradicts failure of the
bundle conclusion. This proves one uniform threshold for the augmented
packet. The resulting points have full cutoffs
by CGP01. CFS26 gives the scale bounds.
\end{proof}

\begin{lemma}[No tiny positive marker anywhere on the smoothed locus; CGP04]
\label{lem:fibration-smoothed-locus-marker-exclusion}
Use CFS11--CFS15 with CFS27's planes. Write $r(x)=\Sigma\sigma_x$,
where $\sigma_x$ is the original scale at its selected preimage.
If $y$ is a selected center, then for every retained $I_+$ block
with $R_i<\sigma_y/16$,
\[
 J_iw=0\quad\hbox{for all }w\in W\cap B(y,20br(y)).   \tag{WE}
\]
Consequently a positive $i$ marker at such a point forces
$\sigma_y\le16R_i$. This applies to the whole $W$, since these
selected-center balls cover its defining domain $U_b$.
\end{lemma}
\begin{proof}
Use the LARGER reference ball $B(y,30br(y))$ of CFS11.
Any contributing center $u$ has
$|u-y|<80br(u)+30br(y)\le110b\max\{r(u),r(y)\}$.
CFS07 still applies with $L=128b$. The proof of CFS28, with the
chosen preimage at $y$, gives the same lower bound
$R_{a_u}\ge108\sigma_y/625$ for its model-reference radius.
Thus every contributing center and plane has zero $i$ block.
The spectral calculation in CFS24 is pointwise on this entire
reference ball: $J_iQ=J_i$, $J_i\mu=0$ and $J_i\eta(z)=J_i z$.
At a zero of the section it gives $J_iw=0$.
No nearest-point assertion on this larger ball is needed.
\end{proof}

\begin{lemma}[Every marked smoothed point has an original full-marker witness; CGP05]
\label{lem:fibration-marked-smoothed-witness}
Fix a stage and its family $i$, with $\ell_i$ as in CGP03.
Assume $0<\varepsilon\le1/10$, $\Sigma\le\varepsilon/640$.
For every
\[
 w\in V_i^0:=\{w\in W:v_i(w)>.9R_i,\
                              |u_i(w)|<5.5\ell_iR_i\}
\]
there is an original point $q\in\widetilde A_j\cap U_i$ with
\[
 \zeta_i(q)=1,\quad |\eta_i(q)|<6.2\ell_i,\quad
 |w-\pi_jF(q)|<\frac{25}{12}\varepsilon\Sigma R_i.    \tag{MW}
\]
For the edge stage $t(q)\le8\Delta$; strict enlargements give a
strict inequality. This is a statement about EVERY point of $V_i^0$.
\end{lemma}
\begin{proof}
Choose a selected $y$ with $|w-y|<20br(y)$ and a CFS14 witness
$z\in\widetilde S_j$ with $|w-z|<\varepsilon r(z)$.
The triangle inequality gives
$|z-y|<(20b+\varepsilon)\max\{r(z),r(y)\}<128b\max\{r(z),r(y)\}$.
Thus $\sigma_z\le(5/3)\sigma_y$. CGP04 and $v_i(w)>0$ imply
$\sigma_y\le16R_i$. Initially, therefore,
\[
 |w-z|<\frac{80}{3}\varepsilon\Sigma R_i
                 \le R_i/2400<R_i/1000.
\]
Choose ANY preimage $q\in\widetilde A_j$ of $z$.
The marker at $z$ is greater than $.899R_i$, so its original
cutoff is positive and $q\in U_i$. Division of its vector block
by its marker gives
\[
 |\eta_i(q)|<\frac{5.5\ell_i+.001}{.899}<6.2\ell_i.
\]
For circles and slim charts this is in the full plateau.
For an edge chart, membership in $\widetilde A_2$ supplies
$t(q)\le8\Delta$ for the SAME global height function, so its
two cutoff factors are both one as well. The scale at this $q$
is at most $5R_i/4$; comparing it with the preimage chosen for
$r(z)$ gives $\sigma_z\le(5/3)(5R_i/4)=25R_i/12$.
This improves the proximity estimate to (MW). All steps tolerate
different radius and image preimages, and require no scale coordinate
in $Q_2$ or $Q_3$.
\end{proof}

\begin{lemma}[A retained coordinate is injective on the whole local graph; CGP06]
\label{lem:fibration-retained-coordinate-graph}
Let $H$ be Euclidean, $\dim L=k$, and $\pi:H\to\mathbb R^k$
an orthogonal coordinate projection. Suppose
$|\pi v|\ge m|v|$ on $L$, $m>0$. If $g$ is smooth and
$G=\{x+t+g(t):t\in B_L(0,R)\}$ with
$\|Dg\|\le a<m$, then $\pi|_G$ is injective, is a local
diffeomorphism, and its inverse on its open image is smooth.

Here is a sufficient way to produce $m$ for the actual cloud at
$x=\pi_jF(p)$, in reference $R_i$ units. Suppose
\[
 T=D\Phi_i,\quad \pi T=I,\quad \|T\|\le\Omega,\qquad
 \|D(\pi_jF)-T D\eta_i\|\le e,
\]
where $D\eta_i$ is onto with a right inverse of norm at most two,
and $\|(I-\Pi_{L_x})D(\pi_jF)\|\le\nu$.
In this criterion $D(\pi_jF)$ abbreviates the derivative of
$R_i^{-1}\pi_jF$, and source norms use $R_i^{-2}g$.
If $\Omega\ge1$ and $\nu+e\le1/(48\Omega)$, then
$|\pi v|\ge |v|/(2\Omega)$ on $L_x$.
\end{lemma}
\begin{proof}
The segment between any two graph parameters stays in the convex
parameter ball. Integration gives
\[
 |\pi(t_1-t_2+g(t_1)-g(t_2))|
                 \ge(m-a)|t_1-t_2|.
\]
The same inequality for the differential proves invertibility,
and the inverse function theorem proves the first assertion.
For the sufficient criterion choose a right inverse $B$ of
$D\eta_i$. Then
$\|(I-\Pi_{L_x})T\|\le2(\nu+e)$.
Since $\pi T=I$, its inverse on its image has norm at most one.
The equal-dimensional projector estimate used in CFS03 bounds
the projector difference between $L_x$ and $\operatorname{im}T$
by $\theta=6(\nu+e)\le1/(8\Omega)$.
For a unit $v\in L_x$, its projection $v_0$ onto
$\operatorname{im}T$ has norm at least $1-\theta$, and
$|\pi v_0|\ge|v_0|/\Omega$. Therefore
$|\pi v|\ge(1-\theta)/\Omega-\theta\ge1/(2\Omega)$.
Only the reference identity component and the normal error were used;
the chosen plane at $x$ need not have been chosen from chart $i$.
\end{proof}

\begin{theorem}[The entire marked base patch is one sheet; CGP07]
\label{thm:fibration-marked-base-one-sheet}
For each reference $i$ suppose the actual rough-graph producer gives
a smooth normalized graph $\Phi_i$ on $B(0,8\ell_i)$, with
its own vector and marker components $(a,1)$, $\|D\Phi_i\|\le\Omega$,
and value and derivative error at most $e$ for the actual
$R_i^{-1}\pi_jF$ wherever $|\eta_i|<8\ell_i$ and, for edges,
$t\le8\Delta$. Derivative norms use $R_i^{-2}g$.
Retain the original coordinate rank bound in CGP06.
Use the actual clouds, CFS27's planes and the sequential adjustments
of CFS31, with normal errors at most $\nu$ and cumulative errors
at most $1/1000$. Impose
\[
 \Omega\ge1,\quad e\le\Sigma/1000,\quad
 \nu+e\le1/(48\Omega),\quad
 \varepsilon\le\frac1{1000(\Omega+1)}.                \tag{OS}
\]
Retain CFS13's graph bounds $|g_x|\le\varepsilon r_x$ and
$\|Dg_x\|\le\varepsilon$ on their full convex domains, as in
the construction of CFS14. Then
\[
 u_i/R_i:V_i^0\longrightarrow B(0,5.5\ell_i)
\]
is a diffeomorphism. In particular it is proper and onto, and
$V_i^0$ is the image of points in the original threshold-$6$
plateau under $f_j=\pi_j\mathcal E^j$.
The conclusion supplies the previously unproved one-sheet statement,
relative to the displayed rough-graph inputs.
\end{theorem}
\begin{proof}
Fix $a\in B(0,5.5\ell_i)$ and put $p_a=s_i(a)$, using CGP03.
The point $x=\pi_jF(p_a)$ is an ORIGINAL core center, and
$r_x\ge(3/5)\Sigma\rho(p_a)\ge(9/20)\Sigma R_i$.
For any $w\in V_i^0$ with $u_i(w)/R_i=a$, CGP05 gives a
full-marker $q$ in the rough-graph domain. Projection of (MW) gives
$|\eta_i(q)-a|\le(25/12)\varepsilon\Sigma$.
Both coordinates and the segment between them belong to
$B(0,8\ell_i)$. The graph approximation and its Lipschitz bound yield
\[
 |w-x|\le
 \left(2e+\frac{25}{12}(1+\Omega)\varepsilon\Sigma\right)R_i
 <\frac{\Sigma R_i}{100}<r_x/4.                       \tag{SN}
\]
CFS13 puts EVERY such $w$ in the SAME graph $G_x$ over $L_x$.
CGP06 gives a lower retained-coordinate bound $1/(2\Omega)$
on that plane. Since $\varepsilon<1/(2\Omega)$, the first part
of CGP06 proves injectivity on the whole graph, including any two
candidate points. It also proves that the coordinate map is a local
diffeomorphism at every point of $V_i^0$. Thus no additional sheet
can enter from another selected-center patch.

To prove existence, use the section $s_i$ on the closed ball of
radius $1/100$ about $a$, which lies in its domain and in the
original threshold-$6$ plateau. Let
$H(u)=u_i(f_j(s_i(u)))/R_i$. Cumulative value control and CFS26 give
$|H(u)-u|\le(5/4)/1000<1/100$.
The continuous map $u\mapsto a+u-H(u)$ carries that closed ball
into itself. For a two-dimensional ball the Brouwer fixed-point theorem
gives $H(u)=a$ for some $u$; in dimension one the endpoint inequalities
and the intermediate value theorem do so directly.
At this point the stage's exact plateau
puts $f_j(s_i(u))$ in $W$, and its marker is greater than
$(1-5/4000)R_i>.9R_i$. Hence this is a point of $V_i^0$.
The only foundational use here is finite-dimensional Brouwer on a
ball, not a new assertion about the possibly nonproper global map.
This standard foundation is an inherited reuse candidate: its accepted
Lean declaration remains to be bound with the completed foundation.

The bijective local diffeomorphism is a diffeomorphism. Its inverse
is the graph parametrization, so compact subsets of the target have
compact full preimages in $V_i^0$. The constructed preimages are
in the actual plateau. By uniqueness they account for ALL of $V_i^0$,
which proves the stated image assertion without assuming it first.
\end{proof}

The fixed-point input was checked in Lee's archived second-edition
TeX, \code{c17.tex}, lines1925--1961, with the preceding extension-degree
argument at lines1896--1923. Its two-dimensional application has
connected sphere boundary. The author errata dated January3,2026
contains no Brouwer entry; the corrected noncompact top-degree argument
for Theorem17.32 is not used here. No published pagination is inferred
from those archived TeX lines.

\begin{theorem}[Later adjustments preserve the marked bases and restricted fibers; CGP08]
\label{thm:fibration-final-bases-no-merging}
Under CGP07's actual inputs put $V_j^0=\bigcup_{i\in\mathcal I_j}V_i^0$,
where $\mathcal I_1=I_2$, $\mathcal I_2=I_e$, $\mathcal I_3=I_s$.
Define
\[
 W_1=(\Psi_3\Psi_2)(V_1^0),\qquad
 W_2=\Psi_{32}(V_2^0),\qquad W_3=V_3^0.
\]
These are embedded manifolds of dimensions $2,1,1$, and the
displayed maps from $V_j^0$ are diffeomorphisms onto them.
The final maps $\pi_j\mathcal E:U_j\to W_j$ are submersions
on FC33's EXACT original threshold-$5$ domains, with its cumulative
value and derivative bounds after the CFS20 target is decreased.

Let $B_j$ be the original threshold-$6$ plateau, including the
height restriction for edges, and put
$D_j=B_j\cap f_j^{-1}(V_j^0)$. On $D_j$ the final and
stage-$j$ maps have exactly the SAME fibers and kernels.
This does not assert fiber equality on arbitrary points outside
that carrier, or equality to the original local-coordinate fibers.
\end{theorem}
\begin{proof}
CGP07 puts $V_1^0$ inside $\mathcal E^1(M)$, where $\Psi_2$
is defined on a neighborhood. Its image lies in $\mathcal E^2(M)$,
where $\Psi_3$ is defined. Similarly $V_2^0\subset\pi_2\mathcal E^2(M)$,
and CFS18/FC32 supply the factored map $\Psi_{32}$ on a neighborhood.
Thus the displayed images do not evaluate an adjustment outside its
proved domain. No global ambient extension is needed.

In a stage-one patch, both later adjustments retain the entire
$i\in I_2$ block. In a stage-two patch, the third retains the
$i\in I_e$ block. A later image of a graph over $u_i/R_i$ is
therefore still a smooth graph over the SAME ball, with the same
coordinate inverse. If two points from different earlier patches
have the same later image, choose a marker patch at the first.
The unchanged vector and marker put the second point in that very
same earlier patch, so CGP07 makes the two points equal.
Moreover the marked ambient conditions are open and are retained;
their intersection with the entire later image is exactly the image
of that earlier patch. These graphs give the subspace topology and
smooth embedded charts on $W_j$. The global inverse is smooth in
those charts. The third-stage assertion is immediate.

For $p\in U_j$, an original index has its full marker and
$|\eta_i(p)|<5\ell_i$, including $t(p)<5\Delta$ for the edge
case. The cumulative error puts $f_j(p)$ inside $V_i^0$ with strict
slack, and the stage's threshold-$6$ plateau makes it a submersion
to the intermediate $W_j^0$, by CFS20/CGP01--CGP02.
Composition with the just proved diffeomorphism gives the final
submersion. The cumulative bounds remain those of CFS20; choose its
last target below $c_{\rm adjust}$ and $1/1000$.
Finally on $D_j$, the final map factors as this diffeomorphism
composed with $f_j$. Injectivity and the chain rule give equality
of its fibers and kernels there. The factorization uses $\mathcal E^1$
at stage one, not the unadjusted $\mathcal E^0$ appearing in the
source's following rank paragraph.
\end{proof}

\paragraph{Parameter and consumer boundary after revision145.}
CGP01--CGP03 supply actual original profiles, the early global derivative
bound and the inner sections from the existing local packets. The tiny-marker
exclusion and witness argument in CGP04--CGP05 apply to the actual strengthened
smoothing with its specified planes. CGP07--CGP08 prove one-sheetness and
the later-image conclusion once FC23/FC27's rough graph and normal/rank
estimates are supplied. Their extra constants depend only on the early
graph derivative bound $\Omega$; $\varepsilon,\Gamma$ can be chosen first,
then $\Sigma$, and then $e\le\Sigma/1000$ through the later local errors.
This is no proof yet that ALL rough-model moduli satisfy that order.
The latter production, proper whole-preimage restrictions, boundary faces,
the exhaustive parameter register and the nearly cuspidal boundary version
remain. In particular $D_j$ is a specified restricted carrier, not an
exclusion of extra preimages in the later decomposition. No parent contract
or accepted-PC/Lean binding is promoted, and this is not the final audit.

\subsection{Actual slim graph models and the third cloudy image}
\label{sec:fibration-actual-slim-cloud}
Revision146 supplies the slim part of FC27 from the actual local
packets. There is a domain distinction: an overlapping slim coordinate
need not be defined on the entire large reference ball. Its original
RAW splitting is defined there. We prove that the actual or model
cutoff can be nonzero only inside the smooth coordinate domain.
This replaces the unnecessarily strong common-smooth-domain demand
of FC23 for this projected family; it does not assert that demand
was verified. The first and edge clouds remain separate producers.

Throughout this subsection $M$ is the CLOSED carrier of LPA04--LPA06,
$F$ is CGP01's original block map, and the actual slim coordinates
are chosen by LFR19--LFR20 at the initial construction. We are allowed
to decrease their quality and their separate value tolerance before
choosing the common late tail. No coordinate is changed after $F$
is defined. Let $P_*\ge1$ bound the first two derivatives of all
fixed unit profiles, including the annular profile.

\begin{lemma}[The actual slim comparison list; SGP01]
\label{lem:fibration-slim-comparison-list}
There is one numerical bound $N_*$, chosen before $\Delta$ and
noncollapse, with the following properties. Put
\[
 L=10^6\Delta,\qquad \ell=10^5\Delta,\qquad
 D_i=B_g(p_i,.95LR_i),\qquad R_i=\rho(p_i),\quad i\in I_s.
\]
Assume $L\Lambda<10^{-5}$ and take the sufficiently late common
tail of LPA01. The list
\[
 J_i=\{j\in I_s:\operatorname{supp}\zeta_j\cap D_i\ne\varnothing\}
\]
contains $i$, has cardinality at most $N_*$, and satisfies
\[
 .99<s_j:=R_j/R_i<1.01,\qquad d_g(p_i,p_j)<2LR_i.
                                                               \tag{SL}
\]
With the zero-radius parameter $T_0\ge1600L$, at most one zero
support meets $D_i$. If it does, its radius $R_0$ satisfies
$R_0/R_i\ge T_0/20$ and the whole $D_i$ is inside FC13's buffered
closed radial shell. Every point of the original set
$\{|\eta_i|\le8\ell\}$ lies strictly inside $D_i$.
\end{lemma}
\begin{proof}
CGP01/FC18 place each slim closed support in its radius-$.95LR_j$
ball. A meeting point gives
$d(p_i,p_j)<.95L(R_i+R_j)$.
Slow variation then bounds $s_j$ between
$(1-.95L\Lambda)/(1+.95L\Lambda)$ and its reciprocal, proving
(SL). The selected cores $B(p_j,\Delta R_j/3)$ are pairwise
disjoint. Apply FC08 with
\[
 R=.95\cdot10^6\Delta,\quad C=10^6\Delta,\quad
 a=\Delta/3,\quad\kappa=1/\Delta.
\]
Its whole-ball Ricci hypothesis is supplied by LPA01 on one later
tail. Since its $H/\Delta$ and $a/\Delta$ are numerical constants,
$V_{1/\Delta}(H)/V_{1/\Delta}(a)=V_1(H/\Delta)/V_1(a/\Delta)$
has a numerical upper bound. Choose $N_*$ larger than this number.
This is a bound on the WHOLE list, not just pointwise multiplicity.
It may be included in the initial multiplicity parameter.

FC09 and FC13, with comparison radius $.95L$, give the zero
assertions. The zero radial value error is chosen below $1/40$
as in LPA05. For the last assertion use the SAME slim product map,
whose entire residual factor has diameter at most $1000\Delta$.
Choose its distortion below $\Delta/100$ and the smooth coordinate
value error below $\Delta/100$. Pointed product distance gives
\[
 R_i^{-1}d(p_i,q)
 \le\sqrt{(8\ell+\Delta/100)^2+(1000\Delta)^2}
                      +\Delta/100<.81L<.95L.
\]
The cutoff at $p_i$ is one, so $i\in J_i$.
\end{proof}

\begin{lemma}[One affine alignment of each original raw coordinate; SGP02]
\label{lem:fibration-slim-raw-alignment}
Fix $\Delta,\beta_2$ and $0<\delta,E<1$. By choosing the later
$\beta_1$ sufficiently small, and then the zero-packet and common
tail parameters, the following holds simultaneously for every $i$.
All distances in this lemma use $R_i^{-2}g$. Set $H=30L$ and
write $u_i$ for the original real coordinate of the slim splitting.
For every $j\in J_i$ there is a sign $a_j\in\{-1,1\}$ such that
\[
 |s_j u_j(x)-a_j u_i(x)-c_j|<E
 \quad(x\in B(p_i,H)),\qquad c_j=s_j u_j(p_i).
                                                               \tag{RA}
\]
Take $a_i=1,c_i=0$. Every original raw tested ball contains the
radius-$100L$ ball about its own center. The reference product
map has distortion at most $\delta$ on that ball and actual
lifts in $B(p_i,21L)$ to error less than $\delta$ for target
points of radius less than $20L$.
If the zero support in SGP01 is present, there is also a sign
$a_0$ such that
\[
 |d(p_0,x)-d(p_0,p_i)-a_0u_i(x)|<E
                     \quad(x\in B(p_i,H)).                     \tag{R0}
\]
The distances in (R0) are still in $R_i$ units. No smooth
coordinate of chart $j$ is evaluated at $p_i$ in (RA).
\end{lemma}
\begin{proof}
The reference center is in the one-stratum, so it admits no
normalized $(2,\beta_2)$-splitting. Rescale each original $j$
splitting by $s_j$ in BOTH source distances and product distances.
It is an arbitrarily good rank-one splitting in reference units,
uniformly for $s_j\in[.99,1.01]$: distortion is multiplied by
$s_j$, and its tested and covered radii tend uniformly to infinity
as $\beta_1\to0$. MC10--MC11 restrict with a buffer, and AC79
recenters it at $p_i$, whose displacement is less than $2L$.
The actual real coordinate becomes $s_j(u_j-u_j(p_i))$.
Neither its values away from the basepoint nor its residual
coordinate are replaced.

Apply AC76 with ranks $j=k=1$, dimension three, and fixed
exclusion quality $\beta_2$. Its growing-ball lower curvature
hypothesis is supplied by LPA01 after the requested tolerances
are fixed. In AC71's compatibility witness the real factor map
is a pointed approximation $\mathbb R\to\mathbb R$.
Choose its accuracy $\tau$ so that
\[
 \tau^{-1}>4(H+1),\qquad 25\tau<E,\qquad
                         \tau<(H+1)/20.
\]
FC20 with radius $H+1$ replaces that factor map on the needed
interval by a sign, with error at most $24\tau$. The source
commutator costs at most $\tau$. Since
$|u_i(x)|\le H+\delta<H+1$, this proves (RA), with exactly the
translation displayed. All choices depend on $H,E,\beta_2$
and not on the number of charts. Decreasing the original
$\beta_1$ further, in particular below
$\min\{1/(1000L),\delta/100\}$, supplies the reference distortion
and actual coverage witnesses to error less than $\delta$.
The original coverage witnesses lie in its original tested ball;
pointed distortion puts those for the radius-$20L$ target ball
inside $B(p_i,21L)$. No source infimum is assumed attained.

For the zero coordinate, SGP01 puts $p_i$ in the required shell.
LC73/LCP04 give an actual $(1,\beta_1)$-splitting there whose real
coordinate is EXACTLY $d(p_0,\cdot)-d(p_0,p_i)$ in reference units.
Apply the same AC76/FC20 argument to it and $u_i$. The real
function is defined on all of $M$; its good product tests have the
chosen large radius. LC73's cone, radial and radius thresholds
are imposed after $\beta_1$, as in LPA04--LPA05. No upper bound
on $R_0/R_i$ is used.
\end{proof}

\begin{lemma}[Binding the original derivative tests and model supports; SGP03]
\label{lem:fibration-slim-actual-affine-comparison}
Given $0<\theta<1$, choose at the initial construction the slim
qualities $\sigma_s$ and value tolerances $v_s$ sufficiently small
depending on $\theta,L$, and the radial adapted quality $\zeta_0$
and difference-Lipschitz error $\epsilon_0$ sufficiently small
depending on them. Then choose $\delta,E,\beta_1,T_0$ and the common
tail as in SGP01--SGP02. For $j\in J_i$ put
\[
 \lambda_j(a)=a_j a+c_j,\qquad U_j=s_j\eta_j.
\]
At every $x\in D_i\cap B_g(p_j,LR_j)$, these SAME
coordinates have
\[
 |U_j-\lambda_j(\eta_i)|<\theta,\qquad
       \|DU_j-a_jD\eta_i\|<\theta.                           \tag{SC}
\]
Derivative norms use $R_i^{-2}g$. Moreover if
\[
 f\bigl(\lambda_j(\eta_i(x))/(s_j\ell)\bigr)>0,
\]
then $x$ lies strictly inside $B(p_j,.91LR_j)$, where $\eta_j$
is defined on a neighborhood. Thus (SC) applies whenever the actual
OR model $j$ cutoff is positive.

For a meeting zero support, with $s_0=R_0/R_i$, use
\[
 U_0=s_0\eta_0,\qquad
 \lambda_0(a)=a_0a+U_0(p_i).
\]
The analogous two bounds hold on ALL of $D_i$.
\end{lemma}
\begin{proof}
Here are sufficient ADDITIONAL strict budgets, alongside all
retained local-packet thresholds; they also specify the order.
Take
\[
 \begin{gathered}
 0<\sigma_s,\zeta_0<\min\{1/100,\,\theta^2/10^6,\,1/(100L)\},\\
 0<v_s,E<\theta/100,\qquad
 0<\epsilon_0<\theta/(100L),\qquad
 0<\delta<\min\{1/100,\,\theta^2/10^6\}.
 \end{gathered}                                                \tag{SB}
\]
In addition decrease $E$ below $\theta^2/10^6$.
Use LFR19 with this separate $v_s$, not just LFR20's default
$\Delta/100$ bound. LFR20 applies to EVERY coordinate meeting
its weaker bounds, so its topology and enclosures are retained.
The original slim test threshold is $L$ in its own units and its
outer tested ball has radius $L/\sigma_s\ge100L$.

Fix $x\in D_i\cap B_g(p_j,LR_j)$. In the ORIGINAL
reference product choose a lift $y$ of
$(u_i(x)+10L a_j,v_i(x))$ to error less than $\delta$.
Its target radius is less than $12L$: the whole residual diameter
is at most $1000\Delta=.001L$.
Pointed distortion puts $y$ in $B(p_i,13L)$ and gives
\[
 10L-2\delta<d(x,y)<10L+2\delta,\qquad
 a_j(u_i(y)-u_i(x))>10L-\delta.
\]
These points lie inside the raw alignment ball. A minimizing
segment from $x$ to $y$ therefore has initial unit vector $w$
on which BOTH original adapted tests apply. For chart $j$ its
length exceeds $s_jL$, and
$d(p_j,y)<15L< s_jL/\sigma_s$; for chart $i$ its length exceeds
$L$ and $y$ is in its original outer tested ball. The derivative
comparison error stays $\sigma_s$ on changing units, since
both coordinate differences and distances acquire the same factor.
Consequently
\[
 DU_j(w),\ a_jD\eta_i(w)\ge1-\varepsilon,\qquad
 \varepsilon=\sigma_s+\frac{3\delta+2E}{10L-2\delta}.
\]
Both covectors have norm at most $1+\sigma_s$. The Riesz-vector
calculation in FC15 gives
$\|DU_j-a_jD\eta_i\|\le2\sqrt{4\varepsilon+\varepsilon^2}<\theta$.
This is a POINTWISE use of the common-endpoint argument: it does
not assume that $\eta_j$ exists elsewhere in $D_i$.
The value estimate follows separately from (RA) and the two
original value bounds:
$|U_j-\lambda_j(\eta_i)|<s_jv_s+E+v_s<\theta$.

If the model cutoff is positive, its argument has absolute value
less than $9$. The same raw comparison, which is valid even
before $\eta_j$ is known to exist, gives
$|u_j(x)|<9\ell+1$. The point $x$ lies within $3L$ in reference
units of $p_j$, so it is inside the ORIGINAL raw distortion ball.
The entire $j$ residual factor has diameter at most $1000\Delta$.
With its distortion made less than $\Delta/100$, pointed distance
therefore gives
\[
 R_j^{-1}d_g(p_j,x)
 <\sqrt{(9\ell+1)^2+(1000\Delta)^2}+\Delta/100<.91L.
\]
This proves the required smooth-domain buffer. An actual positive
cutoff already lies in its original domain. No undefined
coordinate was used to prove the model-support enclosure.

For the zero coordinate, LC73 applies with basepoint $x$ at EVERY
point of $D_i$, by SGP01. Put $s_x=\rho(x)/R_i\in[.99,1.01]$.
Its prescribed radial function has an original all-direction test
for distances greater than $s_x$, up to radius $s_x/\zeta_0$
about $x$, in reference units. The same lift of the reference
axis with sign $a_0$ has length about $10L$ and fits this test
because of (SB). The raw radial coordinate at $x$ differs from
the one in (R0) only by translation and change of units.
Thus the same saturation calculation gives the derivative bound.
For values, LC67's GLOBAL difference-Lipschitz estimate gives
\[
 |U_0(x)-U_0(p_i)-d(p_0,x)+d(p_0,p_i)|
                 \le\epsilon_0 d(p_i,x)<\epsilon_0 L.
\]
Combine this with (R0) and $|\eta_i-u_i|<v_s$. The result is
less than $\theta$. In particular the error is NOT the uncentered
radial value error multiplied by the arbitrarily large $s_0$.
All radial choices are made before the joint zero witnesses.
\end{proof}

\begin{theorem}[Actual smooth slim graph with an early modulus; SGP04]
\label{thm:fibration-actual-slim-graph}
Set
\[
 C_*=1000(N_*+2)(P_*+1).
\]
For every $0<e<1/100$, the permitted choices in SGP01--SGP03
produce, for each actual slim reference $i$, a smooth map
$\Phi_i:\mathbb R\to Q_3$ whose own block is $(a,1)$, such that
\[
 \|D\Phi_i\|,\ \|D^2\Phi_i\|\le C_*,\qquad
 \|R_i^{-1}\pi_3F-\Phi_i\eta_i\|_{C^1}<e
                    \quad\hbox{on }\{|\eta_i|\le8\ell\}.       \tag{SG}
\]
The first and second derivative constants are chosen BEFORE
$\Delta,\beta_2$ and noncollapse. The approximation uses the actual
global zero extensions and all source points in its stated set.
The planes may be chosen by CFS27's prescribed pruning rule.
\end{theorem}
\begin{proof}
Take $\theta<e/(10C_*)$. For $j\in J_i\setminus\{i\}$ use the
block
\[
 \left(\lambda_j(a)
 f\left(\frac{\lambda_j(a)}{s_j\ell}\right),\
 s_j f\left(\frac{\lambda_j(a)}{s_j\ell}\right)\right).
\]
Use $(a,1)$ for $i$. If the zero support is present use
$F_{s_0}(\lambda_0(a))$ with its fixed annular profile, as in FC05.
All other $Q_3$ blocks are zero. This is a finite sum of globally
smooth functions; no actual source coordinate is extended.

For a slim block apply FC05 with the profile
$z\mapsto f(z/\ell)$, support radius $9\ell$, first derivative
bound $P_*/\ell$ and second derivative bound $P_*/\ell^2$.
Since $\ell\ge1$ and $s_j>.99$, its first and second derivatives
are bounded by $50(P_*+1)$. The zero block has the same generous
bound because $s_0\ge T_0/20>1$; its Hessian decreases as $s_0$
increases. Every affine input has derivative a sign. There are
at most $N_*+1$ blocks, including the reference identity block.
Orthogonal square summation proves both bounds in (SG) with
the displayed $C_*$, regardless of total target dimension.

On the original set in (SG), the own block is exactly
$(\eta_i,1)$ and SGP01 puts the point inside $D_i$. For any other
listed slim block, whenever either cutoff is positive SGP03
allows evaluation of the SAME smooth coordinate and bounds its
value and derivative difference from the affine input.
The global first and second derivative bounds on the Euclidean
block map, with $\|D(\lambda_j\eta_i)\|\le2$, give composition
error at most $200(P_*+1)\theta$.
If both cutoffs vanish, both weighted blocks and their
derivatives vanish. This uses smoothness and nonnegativity of
the cutoffs; outside an original domain its closed support
buffer makes the zero extension locally zero.
The zero block is treated directly by FC05 and SGP03 on all of
$D_i$. Every omitted block vanishes on this open ball by the
definition of the list. Square summation gives error less than $e$.

No listed nonzero-family radius is at most $R_i/2$, by (SL).
All other such blocks are already zero in this model. Hence
CFS27's pruning leaves this model unchanged, preserves the
identity block and imposes exactly its required kernel rule.
\end{proof}

\begin{lemma}[Actual projected rank at every preimage; SGP05]
\label{lem:fibration-slim-all-preimage-rank}
Choose for each $x\in S_3$ an actual core witness $i,p$, with
$|\eta_i(p)|\le7\ell$, and put
\[
 T_x=D\Phi_i(\eta_i(p)),\qquad
 A_x=x+\operatorname{im}T_x.
\]
For EVERY $q\in M$ with $\pi_3F(q)=x$, the projection
$P_xD(\pi_3F)_q$ onto $A_x^0$ is onto, with normal error less
than $e$. On the orthogonal complement of its kernel its
singular value lies between $1/2$ and $3C_*$.
These are original-metric operator estimates. Thus this choice
of planes supplies actual normal and rank bounds, not a
vacuous estimate on a possible zero-rank map.
\end{lemma}
\begin{proof}
The full marker of $i$ at $x$ forces $\zeta_i(q)=1$ and
$q\in U_i$. Its vector block then gives
$\eta_i(q)=\eta_i(p)$. Thus SGP04 applies at EVERY such $q$
with the SAME $T_x$.
The original reference-axis test used in SGP03, now with its
positive sign, gives
$\|R_i d\eta_i(q)\|>9/10$ in the original metric, while its
upper bound is at most two. The lift fits all original test
domains by exactly the bounds already checked there.
The identity component of $T_x$ gives least singular value
at least one and $\|T_x\|\le C_*$.

Consequently $D(\pi_3F)_q$ differs from
$T_x(R_i d\eta_i(q))$ by less than $e$ in the original metric:
the factor $R_i^{-1}$ in the target and the factor $R_i$
in a unit source vector for $R_i^{-2}g$ cancel.
Its normal part is less than $e$.
On the one-dimensional target $A_x^0$ the adjoint of the
comparison map has norm at least $9/10$. After perturbation
it still has lower bound $9/10-e>1/2$, hence is injective.
The projected map is therefore onto. Equality of nonzero
singular values for a map and its adjoint gives the claimed
lower bound on the kernel complement. The upper bound is
$2C_*+e<3C_*$. This is FC06's elementary rank argument in
dimension one, with its onto conclusion proved explicitly.
\end{proof}

\begin{theorem}[The actual third cloudy packet, with both coverage directions; SGP06]
\label{thm:fibration-actual-third-cloud}
For $0<\Gamma<1$, choose before $\Delta$
\[
 0<\Sigma<\min\{\Gamma/200,\ \Gamma^3/(100C_*)\},
 \qquad 0<e<\min\{1/100,\ \Gamma\Sigma/100,\ \Sigma/1000\}.
                                                               \tag{CP}
\]
Use SGP04 with that $e$, FC27's ORIGINAL slim core and enlargement,
and $r(x)=\Sigma\rho(p_x)$ with any chosen original preimage.
If $I_s\ne\varnothing$, then $(S_3,\widetilde S_3,r)$ is an actual $(2,\Gamma)$ cloudy
one-manifold in the required recentered truncated-Hausdorff
convention. Its planes are those of SGP05. They have normal
error less than $\Gamma$, rank one, the uniform singular-value
bounds there, and the whole-chart graph approximation (SG).
For CFS14 with $b=\varepsilon^{-1}$ and its smoothing threshold
on $\Gamma$, also choose $\Sigma\le1/(640b)$;
CFS07's required all-center scale
comparability and all remaining geometric cloud inputs hold.
If $I_s=\varnothing$, the third adjustment is instead the identity.
\end{theorem}
\begin{proof}
The family is finite, every original threshold-$8\ell$ set has
a compact interior enclosure in its smooth domain, and $M$ is
compact. Thus the enlarged image is bounded. Its chosen radii
are positive, bounded above and bounded away from zero.
Every enlarged image point has a full slim marker. CFS26 supplies
the support-scale comparison for EVERY preimage, so FC26 gives
the stated radius inequality. With the additional bound on
$\Sigma$, exactly the same actual data meet CFS07 with $B=5/3$.

Fix a core center $x=\pi_3F(p)$ and its SGP05 reference $i$.
Use $R_i$ units for the target. Write $\widehat r=r(x)/R_i$,
$a=\eta_i(p)$ and $R=\widehat r/\Gamma$. CFS26 gives
\[
 3\Sigma/4\le\widehat r\le5\Sigma/4,\qquad R<1/100.
\]
If $z\in\widetilde S_3$ and $|z-x|\le r(x)/\Gamma$, every
original preimage $q$ of $z$ has positive $i$ marker, and FC03 gives
\[
 |\eta_i(q)-a|
 \le\frac{(1+7\ell)R}{1-R}<\ell/10.
\]
Hence $|\eta_i(q)|<8\ell$ and that marker is in fact full.
Its vector coordinate is therefore exactly $\eta_i(q)$.
This localizes ALL of the truncated enlarged image, including
points initially supplied by some other chart.

For every $u\in\overline B(a,R)$ the original proper trivial
bundle in LFR20 has a point $q_u$ with $\eta_i(q_u)=u$,
since $|u|<8\ell<9\ell$. Such a point belongs to the original
enlargement and has full marker. Its actual image approximates
$\Phi_i(u)$ by less than $e$, by SGP04. This is exact coordinate
coverage, not an assumption of metric density.
FC25 now applies with Hessian bound $C_*$ to the normalized
ENLARGED image. Its error is at most
\[
 3q,\qquad q=2e+C_*R^2/2.
\]
The coarser bounds $\Sigma/2\le\widehat r\le2\Sigma$ and (CP)
give $3q/\widehat r<\Gamma/2<\Gamma$. Thus both same-radius
truncated coverage directions hold, with strict slack. The proof
also gives the open-ball distance-to-set convention: in its
plane-to-set step the radial trimming puts the witness strictly
inside the ball; in the other direction an arbitrarily small
further radial contraction uses the available strict slack.
This is the required cloud test at radius $r(x)/\Gamma$.

SGP05 supplies normal error, surjectivity and both positive
singular-value bounds for the very same planes at every
preimage. All affine models use CFS27's stipulated recipe.
If $I_s$ is empty, both images are empty and the third adjustment
is the identity, since its marker sum is zero; no nonempty cloud
or base is invented.
\end{proof}

\paragraph{Consumer and parameter boundary after revision146.}
SGP01--SGP06 produce the previously source-omitted SLIM part
of FC27, including the actual graph, early derivative/Hessian
modulus, all-preimage rank and both coverage directions.
Here $N_*,P_*,C_*$ are early constants; $\Gamma,\Sigma,e$ are
chosen before $\Delta$; the scalar qualities and value errors
are chosen after $\Delta$; $\Lambda$ may be chosen there as well;
the raw compatibility thresholds reduce the later $\beta_1$;
the zero radius/cone choices and the common curvature tail follow.
The uncentered zero values and the total number of charts do not
enter an early constant. No choice makes $\beta_2^{-1}$ dominate
a later slim test radius. The normal/rank and rough-model inputs
of CGP07 are now produced for the third cloud, with its further
accuracy minima imposed among the displayed early choices.

This does not yet construct the first or edge rough clouds,
the completed sequence of all three adjustments, proper bundle
restrictions, compatible faces/corners, the full simultaneous
parameter register, or the nearly cuspidal boundary assembly.
Those remain within the active goal. No parent contract as a
whole, accepted foundation binding or Lean declaration is
promoted, and this is not the final mathematical audit.

\subsection{Actual edge graph models and the second cloudy image}
\label{sec:fibration-actual-edge-cloud}
Revision147 supplies the edge part of FC27. The carrier, maps and
single distance smoothing are the SAME as in CGP01 and SGP01--SGP06.
All extra quality requirements below are imposed at their initial
construction, jointly with the previous requirements. No map is
altered after $F$ is defined. Write $t=P/\rho$ for the shared height,
$L=10^6\Delta$, and retain the unit-profile bound $P_*$.
The first cloudy image and the global assembly remain separate.

\begin{lemma}[Actual edge centers exclude a two-splitting; EGP01]
\label{lem:fibration-edge-exclusion}
Choose the early $\beta_2<10^{-6}$ and the later strong endpoint
and splitting qualities $s,b<10^{-6}$, in addition to LFR29.1.
At every actual strong edge center $p$, in its own scale, there is
no $(2,\beta_2)$-splitting.
\end{lemma}
\begin{proof}
If such a splitting existed, lift its four targets
$(\pm e_1,y_0),(\pm e_2,y_0)$ with error less than $2\beta_2$.
The actual witnesses have pointed radii less than $1+3\beta_2<2$.
Their distances to $p$ differ from $1$ by less than $3\beta_2$,
and their pairwise distances differ from the unit-cross distances
by less than $5\beta_2$. These conclusions use the original tested
ball, not assumed continuity of the approximation.
The actual composite edge chart $Q=(u,Gv)$ of LFR30 is pointed
at the BORDER $(0,0)$, has image in the upper half-plane and
distortion at most $b+s$ on this radius-two ball.
Its four images would satisfy LFR29 with error
$5\beta_2+b+s<10^{-3}$, a contradiction.
This proves the exclusion directly. It does not assume that a
strong edge center satisfies the later $\beta_1$ one-stratum
predicate. It is the quantitative adapter for KL9.6 needed here.
\end{proof}

\begin{lemma}[Whole edge comparison lists and original domains; EGP02]
\label{lem:fibration-edge-comparison-list}
For $i\in I_e$ put $R_i=\rho(p_i)$ and
$D_i=B_g(p_i,20\Delta R_i)$. Let $J_e,J_s$ consist of the
edge and slim indices whose ACTUAL closed supports meet $D_i$.
There is an early numerical $N_\dagger$, independent of $\Delta$,
noncollapse and the total number of charts, bounding
$|J_e|+|J_s|$. If $L\Lambda<10^{-5}$, then $i\in J_e$ and,
for every listed $j$, $.99<s_j=R_j/R_i<1.01$.
Furthermore
\[
 \begin{array}{ll}
 j\in J_e:&d_g(p_i,p_j)<36\Delta R_i,\quad
                         D_i\subset B_g(p_j,57\Delta R_j),\\
 j\in J_s:&d_g(p_i,p_j)<.92LR_i,\quad
                         D_i\subset B_g(p_j,.91LR_j).
 \end{array}                                                   \tag{EL}
\]
Thus every listed smooth coordinate is defined on a neighborhood
of $D_i$. With $T_0\ge1600L$ there is at most one meeting zero
support, and all of $D_i$ lies in its buffered shell.
The original edge enlargement
$\{|\eta_i|\le8\Delta,\ t\le8\Delta\}$ lies in $D_i$.
\end{lemma}
\begin{proof}
FC18 gives support radii $15\Delta R_j$ and $901002\Delta R_j$.
A meeting point therefore gives
$d_g(p_i,p_j)<20\Delta R_i+C\Delta R_j$ with
$C=15$ or $901002$. Slow variation bounds $s_j$ between
$(1-20\Delta\Lambda)/(1+C\Delta\Lambda)$ and
$(1+20\Delta\Lambda)/(1-C\Delta\Lambda)$, proving the ratio bounds.
For edge indices the center distance is less than
$(20+15\cdot1.01)\Delta R_i<36\Delta R_i$; adding the
reference radius and dividing by $.99$ gives the first line of
(EL). For slim indices the center bound is
$(20+901002\cdot1.01)\Delta R_i<.92LR_i$.
For any $x\in D_i$ the sharper own-scale bound is
$d_g(p_j,x)/R_j<(40/.99+901002)\Delta<.91L$.
Both are strictly inside the ORIGINAL smooth domains.

Apply FC08 separately to the disjoint selected $\Delta R_j/3$
cores, with $R=20\Delta$, $a=\Delta/3$, $\kappa=1/\Delta$
and $C=15\Delta$ or $10^6\Delta$. All dimensionless volume
ratios are numerical. LPA01 supplies their whole-ball curvature
hypotheses on a common later tail. Their sum gives $N_\dagger$,
which can be included in the initial multiplicity constant.
FC09/FC13 give the zero assertions with comparison radius
$20\Delta$. At an enlargement point the original $i$ cutoff is
one, so FC18 gives distance at most $15\Delta R_i<20\Delta R_i$.
At $p_i$ its two factors are one, proving $i\in J_e$.
\end{proof}

\begin{lemma}[Fixed raw edge alignments on a short common ball; EGP03]
\label{lem:fibration-edge-raw-alignment}
Fix $\Delta,\beta_2$ and $0<\delta,E<1$.
Choose $\beta_E$ sufficiently small, then the later $\beta_1$
and zero parameters, and the common tail. In reference $R_i$
units, on $B(p_i,600\Delta)$ each $j\in J_e\cup J_s$ has
one sign $a_j$ and translation $c_j=s_j u_j(p_i)$ satisfying
\[
 |s_j u_j-a_j u_i-c_j|<E.                                  \tag{ER}
\]
Here $u_j$ is the ORIGINAL raw real coordinate, and
$a_i=1,c_i=0$. Every original edge/slim raw tested ball contains
its own radius-$100L$ ball and has distortion at most $\delta$
there; actual lifts of targets of radius at most $20L$ have
error less than $\delta$ and pointed radius less than $21L$.
For a meeting zero support the SAME normalized distance function
has, on the short common ball, a sign $a_0$ with
\[
 |d(p_0,x)-d(p_0,p_i)-a_0u_i(x)|<E.                         \tag{ER0}
\]
\end{lemma}
\begin{proof}
EGP01 supplies the exclusion at the reference center.
Rescale each original splitting by $s_j$ and recenter at $p_i$
using MC10--MC11 and AC79. The largest displacement is less than
$L$, by EGP02. This changes the real function precisely to
$s_j(u_j-u_j(p_i))$. AC76 with ranks one and one, followed by
FC20 as in SGP02, gives (ER): take its factor tolerance $\tau$
with $\tau^{-1}>4(600\Delta+1)$ and $25\tau<E$.
The original qualities can be made uniformly small enough for
all these fixed buffered domains, including the radius-$100L$
distortion and coverage assertions. A coverage witness is first
chosen inside the original tested ball and then localized by
pointed distortion; this proves the stated radius bound.

There is no dependence of the $\beta_E$ bound on the later
$\beta_1$: choose both below the required compatibility threshold,
in their prescribed order. The later common curvature tail
supplies AC76 on its growing comparison ball. For the zero
coordinate use LC73/LCP04 at $p_i$ in the buffered shell.
Its real function is exactly the displayed distance difference.
The same compatibility and factor alignment prove (ER0).
The original edge RAW balls may be large here; its smooth
derivative test radius remains exactly $1000\Delta$.
\end{proof}

\begin{lemma}[The actual asymmetric edge comparisons; EGP04]
\label{lem:fibration-edge-actual-comparison}
For every $0<\theta<1$, choose the original tangential qualities,
separate value errors and radial smoothing errors sufficiently
small, then use EGP03 with sufficiently small $\delta,E$.
On ALL of $D_i$, with derivative norm in $R_i^{-2}g$, put
\[
 U_j=s_j\eta_j,\qquad \lambda_j(a)=a_ja+c_j
                      \quad(j\in J_e\cup J_s).
\]
Then the SAME smooth coordinates satisfy
\[
 |U_j-\lambda_j(\eta_i)|<\theta,\qquad
                  \|DU_j-a_jD\eta_i\|<\theta.                \tag{EC}
\]
For a meeting zero block use $U_0=s_0\eta_0$ and
$\lambda_0(a)=a_0a+U_0(p_i)$; (EC) holds for it too.
\end{lemma}
\begin{proof}
Here are sufficient additional budgets, alongside the retained
local-packet bounds. Choose the edge/slim qualities and radial
adapted quality below
$\min\{\theta^2/10^8,1/(1000L)\}$.
Use LFR19 with separate value errors less than $\theta/100$.
Take the radial difference-Lipschitz error less than
$\theta/(100L)$, and $\delta,E<\theta^2/10^8$.
EGP03 is then imposed through the later splitting qualities and
zero choices. These budgets do not replace the existing local
topology or smoothing thresholds.

First let $j\in J_e$. For $x\in D_i$, choose an ACTUAL lift in
the reference raw product of
$(u_i(x)+400\Delta a_j,v_i(x))$ to error less than $\delta$.
Its target radius is less than $421\Delta$.
Its witness $y$ lies in $B(p_i,430\Delta)$, has
$400\Delta-2\delta<d(x,y)<400\Delta+2\delta$, and its
signed raw gain exceeds $400\Delta-\delta$.
Both $x,y$ are in (ER)'s ball. They lie in both ORIGINAL
edge outer test balls: for chart $j$ the pointed radius of $y$
is less than $466\Delta<1000s_j\Delta$.
Their separation exceeds both $100\Delta$ and $100s_j\Delta$.
The two original all-direction tests on a common minimizing
initial unit vector $w$ give
\[
 DU_j(w),\ a_jD\eta_i(w)\ge1-\epsilon,\qquad
 \epsilon\le\max\{\varsigma_{\rm edge},\varsigma_{\rm slim}\}
                       +(3\delta+2E)/(400\Delta-2\delta).
\]
Both covectors have norm at most one plus their prescribed quality.
FC15's pointwise Riesz calculation gives their difference less
than $\theta$.

For $j\in J_s$, a common long endpoint would FAIL the original
edge test. Instead, in $j$ units lift
$(u_j(x)+10L,v_j(x))$ to error less than $\delta$.
EGP02 puts $x$ within $.91L$ of $p_j$; the target has radius
less than $11L$ and its witness $y$ has radius less than $12L$.
The segment length $\ell=d(x,y)$ in reference units satisfies
$s_j(10L-2\delta)<\ell<s_j(10L+2\delta)$.
This exceeds the original slim separation $s_jL$, and $y$ is
inside its original outer test radius $s_jL/\varsigma_{\rm slim}$.
On a minimizing segment from $x$ to $y$ let $z$ be the point at
reference distance $200\Delta$. Then $z\in B(p_i,220\Delta)$;
the edge test to $z$ is within $1000\Delta$ and exceeds
$100\Delta$. The original slim raw distortion applies to
$x,z,y$, which all remain within its radius-$100L$ ball.
It gives
\[
 \begin{split}
 s_j(u_j(y)-u_j(x))&\ge\ell-3s_j\delta,\\
 s_j(u_j(y)-u_j(z))&\le\ell-200\Delta+s_j\delta .
 \end{split}
\]
The raw alignment (ER) is needed only at $x,z$.
Thus the slim long test and the edge short test, with the SAME
initial vector, saturate their covectors with error at most
\[
 \max\{\varsigma_{\rm slim},\varsigma_{\rm edge}\}
                   +(4s_j\delta+2E)/(200\Delta).
\]
This is FC22's pointwise calculation, with every original domain
checked. FC15 again proves the differential assertion.
The value assertion in both cases follows independently from
(ER) and the two original value errors, at cost
$E+(1+s_j)\theta/100<\theta$.

For the zero block use the reference axis lift of length
$400\Delta$ with sign $a_0$. LC73 at EVERY $x\in D_i$ has
the same raw distance function up to translation and units.
Its own radius is $\rho(x)/(R_i\zeta_0)>990L$ in reference
units, while its separation threshold is $\rho(x)/R_i<1.01$.
The lift fits these original tests and (ER0), proving the
derivative estimate. LC67 bounds the centered value error by
$\epsilon_0 d(p_i,x)<20\Delta\epsilon_0$.
Together with (ER0) and the reference value error this is less
than $\theta$. The arbitrary zero radius never multiplies an
uncentered smoothing error.
\end{proof}

\begin{lemma}[An exact edge section for cloud coverage; EGP05]
\label{lem:fibration-edge-cloud-section}
At the initial construction the common local packets may also
be required to have a continuous section
\[
 s_i:(-8.5\Delta,8.5\Delta)\longrightarrow U_i,\qquad
             \eta_i s_i(a)=a,\quad 0\le t s_i(a)<\Delta/100.
\]
Its image lies in $B_g(p_i,10\Delta R_i)$. The threshold for
this augmented conclusion has the same permitted dependencies
as LFR28. No disk bundle over a larger base is asserted.
\end{lemma}
\begin{proof}
Use LFR28's SAME product embedding, on a fixed neighborhood of
$[-9\Delta,9\Delta]\times\{z_0\}$, in normalized units.
LFR14 provides this fixed buffer and its distance and metric
comparisons. In its counterexample-sequence proof, tests against
$(a+200\Delta,z_0)$ have separation tending to $200\Delta$,
greater than $100\Delta$, and pointed radii less than
$210\Delta<1000\Delta$. LFR18 and the original LFR19 test
therefore give uniformly along the WHOLE line segment
\[
 |\eta_i j_i(a,z_0)-a|<2\mu\Delta,\qquad
                     \partial_a(\eta_i j_i)>3/4 .
\]
The source radius is less than $10\Delta$ on the common tail.
LFR28.2 is valid on this entire segment, since it lies in the
fixed radius-$100\Delta$ model ball. Here $r(z_0)=0$.
Combine that bound with the coarse-border distance estimate,
the same physical smoothing value error and slow variation.
With the retained sufficiently small $\tau,\mu,\Lambda$ these
give $0\le(P/\rho)j_i<\Delta/100$. This is a value estimate;
no differentiability of the height at the core is used.

Take $2\mu<1/100$. Values at $a=\pm8.9\Delta$ enclose the
closed interval $[-8.5\Delta,8.5\Delta]$. Strict monotonicity
and the intermediate value theorem supply a continuous inverse
there, which composed with $j_i$ gives the asserted section.
For uniformity, add failure of this section to LFR28's failure
conjunction, keeping all earlier tolerances and analytic data
fixed. Its same compactness limit and common tail produce the
three displayed bounds and the inverse, contradicting that
failure. This proves a single allowed threshold for the augmented
local packet, not a separate choice depending on each center.
\end{proof}

\begin{theorem}[Actual edge graph with an early modulus; EGP06]
\label{thm:fibration-actual-edge-graph}
Set $C_\dagger=1000(N_\dagger+2)(P_*+1)$ before $\Delta$.
For every $0<e<1/100$, permitted choices in EGP01--EGP05
produce a smooth $\Phi_i:\mathbb R\to Q_2$ with own block
$(a,1)$ and
\[
 \|D\Phi_i\|,\|D^2\Phi_i\|\le C_\dagger,\qquad
 \|R_i^{-1}\pi_2F-\Phi_i\eta_i\|_{C^1}<e
       \quad\hbox{on }\{|\eta_i|\le8\Delta,\ t\le8\Delta\}.
                                                               \tag{EG}
\]
This is the original global zero-extended map on the entire
stated set. The affine models obey CFS27's pruning rule.
\end{theorem}
\begin{proof}
Choose $\theta<e/(10C_\dagger)$ in EGP04.
Use the own block $(a,1)$. For $j\ne i$ in the actual lists use
\[
 \left(\lambda_j(a)
         f\left(\frac{\lambda_j(a)}{s_j\ell_j}\right),
       s_j f\left(\frac{\lambda_j(a)}{s_j\ell_j}\right)\right),
 \qquad
 \ell_j=\begin{cases}\Delta&j\in J_e,\\10^5\Delta&j\in J_s.\end{cases}
\]
For the possible zero block use FC05's $F_{s_0}(\lambda_0(a))$.
All other blocks are zero; $Q_2$ contains neither the weak-edge
height block nor the scale or two-stratum blocks.

On the stated original set, EVERY edge height factor
$g(t/\Delta)$ equals one and has zero differential, including
at $t=8\Delta$. The composed factor is smooth there by CGP01;
its derivative is zero also on the possibly nonsmooth core,
where it is locally constant. Thus the actual edge blocks reduce
exactly to the tangential factors used above. EGP02 places this
whole set in $D_i$ and every listed coordinate in its original
smooth domain. For unlisted blocks the actual closed support
misses the open $D_i$, so both the block and its derivative vanish.

FC05's cutoff derivative bounds apply with $\ell_j\ge1$ and
$s_j\in[.99,1.01]$. Each model block has its first two derivatives
bounded by $50(P_*+1)$; the zero radius satisfies
$s_0\ge T_0/20>1$, so its Hessian bound only improves.
Orthogonal square summation over at most $N_\dagger+1$ blocks
plus the identity gives $C_\dagger$.
Using (EC), the corresponding actual/model $C^1$ difference
per block is at most $200(P_*+1)\theta$, as in SGP04.
This proves (EG). All listed nonzero-family radii exceed
$.99R_i$; all other such model blocks already vanish.
Consequently CFS27's deletion of radii at most $R_i/2$
leaves this model unchanged.
\end{proof}

\begin{theorem}[The actual second cloudy packet; EGP07]
\label{thm:fibration-actual-second-cloud}
For $0<\Gamma<1$ choose, BEFORE $\Delta$,
\[
 0<\Sigma<\min\{\Gamma/200,\Gamma^3/(100C_\dagger)\},
 \qquad 0<e<\min\{1/100,\Gamma\Sigma/100,\Sigma/1000\}.
                                                               \tag{EP}
\]
Use EGP06 with this $e$ and FC27's ORIGINAL edge core and
enlargement. If $I_e\ne\varnothing$, their projected images
$(S_2,\widetilde S_2,r)$, with $r(x)=\Sigma\rho(p_x)$ for any
chosen original core preimage, are an actual $(2,\Gamma)$
cloudy one-manifold in the identical-radius truncated convention.
For each $x\in S_2$ choose an original core witness $i,p$,
write $a=\eta_i(p)$ and set
\[
 T_x=D\Phi_i(a),\qquad A_x=x+\operatorname{im}T_x .
\]
At EVERY $q\in M$ with $\pi_2F(q)=x$, projection of
$D(\pi_2F)_q$ to $A_x^0$ is onto, its nonzero singular value
lies in $[1/2,3C_\dagger]$, and its normal error is less than
$e<\Gamma$. These planes and radii satisfy CFS14's geometric
inputs when its smoothing threshold on $\Gamma$ and
$\Sigma\le1/(640b)$ are included. If $I_e$ is empty, the
second adjustment is the identity.
\end{theorem}
\begin{proof}
Every enlarged image point has an original full edge marker.
At ANY of its preimages that marker gives
$f(\eta_j/\Delta)g(t/\Delta)=1$. Each factor is in $[0,1]$,
and the fixed profile is strictly between zero and one on
$(8,9)$. Hence both factors equal one, and in particular
$t\le8\Delta$ at EVERY preimage. CFS26 and FC26 give the
radius inequality independently of the selected preimages.
Finite support enclosures and compactness give bounded enlarged
image and positive upper and lower radius bounds. CFS07 with
$B=5/3$ supplies the additional scale input.

At a core image $x$, its full $i$ marker gives
$|\eta_i(q)|\le8\Delta$ and its vector block gives
$\eta_i(q)=a$ at every preimage. Thus (EG) holds there with
the SAME $T_x$. The reference positive-axis test in EGP04 gives
$\|R_i\,d\eta_i\|>9/10$ in the original metric, while its upper
bound is two. The identity component of $T_x$ gives lower norm
one and its upper norm is $C_\dagger$. Restoring physical units
in (EG) cancels $R_i$ between source and target exactly as in
SGP05. The projected adjoint therefore has lower norm
$9/10-e>1/2$, proving actual surjectivity; the upper norm is
less than $3C_\dagger$. The normal part has norm less than $e$.

For cloud coverage fix the same core witness and use $R_i$
units. Write $\widehat r=r(x)/R_i$, $R=\widehat r/\Gamma$.
Then $3\Sigma/4\le\widehat r\le5\Sigma/4$ and $R<1/100$.
If $z\in\widetilde S_2$ with $|z-x|\le r(x)/\Gamma$,
its positive $i$ marker and FC03 give, at every preimage $q$,
\[
 |\eta_i(q)-a|
       \le\frac{(1+7\Delta)R}{1-R}<\Delta/10.
\]
As just proved, every such $q$ also has $t\le8\Delta$.
Thus its full $i$ marker, its exact vector coordinate and (EG)
apply; this localizes the ENTIRE truncated enlarged image.

Conversely, for each $u\in\overline B(a,R)$, one has
$|u|<8\Delta$. EGP05 supplies an actual point with
$\eta_i=u$ and $t<\Delta/100$. It belongs to the ORIGINAL
enlargement and its image is within $e$ of $\Phi_i(u)$.
FC25 now proves both coverage directions with error
$3q$, $q=2e+C_\dagger R^2/2$. The coarser bounds
$\Sigma/2\le\widehat r\le2\Sigma$ and (EP) give
$3q/\widehat r<\Gamma/2$. This has the same strict slack as
SGP06, and its radial trimming also proves the open-ball
distance-to-set convention. The already proved normal and rank
estimates are for exactly these planes and all preimages.
If the family is empty, both images are empty and no artificial
base or nonempty cloud is required.
\end{proof}

\paragraph{Consumer and parameter boundary after revision147.}
EGP01--EGP07 and SGP01--SGP06 now produce BOTH projected
clouds of FC27 under the retained foundations, including the
original rough graphs, early derivative/Hessian bounds,
all-preimage rank and both coverage directions. The bounds
$N_\dagger,C_\dagger$ are early numerical constants.
The no-two-splitting requirement on $\beta_2$ is early; the
endpoint, coordinate and raw qualities are later minima in
their prescribed order. Large slim tests are never placed in
the original edge derivative domain. EGP05 augments a local
section, not the proper disk-bundle base.

This leaves the first cloud, the binding of the whole three-stage
adjustment, proper WHOLE preimage restrictions and original-fiber
isotopies, compatible faces/corners, the exhaustive simultaneous
parameter register and nearly cuspidal boundary assembly.
FC33 and the parent static theorem are not promoted.
No accepted-PC or Lean binding is claimed, and this is not the
requested final full-goal mathematical audit.

\subsection{Actual two-stratum graph models and the first cloudy image}
\label{sec:fibration-actual-first-cloud}
Revision148 completes the original-image producers for FC07 and
FC27 in the CLOSED standing setting. Retain LC20's absence of
three-splittings at its fixed early quality $\beta_3$, the same
local families of LPA04--LPA06, and CGP01's single map $F$.
The following are simultaneous initial requirements on those data,
not modifications of a previously constructed map.
Let $P\ge1$ bound the first two derivatives of all fixed unit
profiles, including the radial two-dimensional circle cutoff.

\begin{lemma}[Whole two-stratum comparison lists and early buffers; TCP01]
\label{lem:fibration-first-comparison-list}
For $i\in I_2$ set $R_i=\rho(p_i)$ and
$D_i=B_g(p_i,10R_i)$. Let $J_2,J_e,J_s$ be the indices whose
ACTUAL closed supports meet $D_i$. Their total cardinality is
bounded by an early numerical $N$, independent of $\Delta$ and
noncollapse. For $L=10^6\Delta$ and $L\Lambda<10^{-5}$ all
listed ratios $s_j=R_j/R_i$ belong to $(.99,1.01)$.
Every listed original smooth coordinate is defined on $D_i$.
There is at most one meeting zero support, with $s_0\ge T_0/20$,
and $T_0\ge1600L$ buffers its entire shell around $D_i$.
Every original point $|\eta_i|\le8$ lies in $D_i$.

At the EARLY circle-coordinate construction, for any prescribed
small $\gamma>0$, one may require the original LFR06/LFR07
coordinate to have quality at most $\gamma$, value error
less than $201\gamma$, and
\[
 \|D\eta_i(D\eta_i)^*-I\|<\gamma/4
\]
in its own normalized metric on $B(p_i,200)$.
Its original derivative separation is $201$; its original far-test
ball may be required to contain $B(p_i,1000)$.
These choices, and the needed smaller $\beta_2$, precede $\Delta$.
\end{lemma}
\begin{proof}
Use the original circle, edge and slim support radii
$10R_j$, $15\Delta R_j$ and $901002\Delta R_j$.
The meeting-point and slow-scale calculation of EGP02 gives the
ratio bounds. Circle centers are less than $21R_i$ from $p_i$,
so $D_i$ is within radius $32$ in their own units, less than
their original smooth radius $200$. For edge and slim indices,
the sharper respective own-scale bounds are
$20/.99+15\Delta<36\Delta$ and
$20/.99+901002\Delta<.91L$. Their original domains therefore
also contain $D_i$. FC08 bounds the entire circle list with
$(R,C,a,\kappa)=(10,10,1/3,1)$. For edge/slim use
$(10\Delta,15\Delta,\Delta/3,1/\Delta)$ and
$(10\Delta,10^6\Delta,\Delta/3,1/\Delta)$.
The source comparison ball of radius $10$ is contained in the
one of radius $10\Delta$; hence these bounds apply to its lists.
LPA01 supplies all whole-ball curvature hypotheses on one later
tail. All the resulting volume ratios are numerical.
FC09/FC13 give the zero assertions.
LFR07's original residual bound and value error $<1/10$ give
\[
 d_g(p_i,q)/R_i<\sqrt{(8+1/10)^2+1}+1/10<10
                     \quad\hbox{if }|\eta_i(q)|\le8 .
\]

For the final assertion use exactly LFR06's construction after
the fixed rescaling by $201$. Its proof bounds the Gram matrix
on a larger smooth ball, so this estimate remains valid on the
stated original radius-$200$ domain. Choose its auxiliary tested
radius greater than $\max\{100,\gamma^{-1}\}$ before choosing
the splitting threshold. Returning to original units multiplies
values and tested lengths by $201$; differential and Gram bounds
are unchanged. LPA03 permits this smaller $\beta_2$ without
noncollapse data. No later $\Delta$ enters this choice.
\end{proof}

\begin{lemma}[Raw comparisons with an early reference quality; TCP02]
\label{lem:fibration-first-raw-alignment}
Fix small positive $E,\delta$ BEFORE $\Delta$. The reference
$\beta_2$ can be chosen before $\Delta$ so that, after the later
edge, slim, zero and common-tail choices, each listed constant-radius
coordinate has, on $B_{R_i^{-2}g}(p_i,1000)$, a fixed coisometry
$A_j:\mathbb R^2\to\mathbb R^{k_j}$ and a translation $c_j$ with
\[
 |s_j u_j-A_j u_i-c_j|<E,\qquad c_j=s_j u_j(p_i).          \tag{TR}
\]
Here $k_j=2$ for circles and $1$ otherwise, and every $u_j$ is
the ORIGINAL raw Euclidean coordinate. One takes $A_i=I,c_i=0$.
The original circle raw balls contain radius $5000$ and have
distortion less than $\delta$ there, with actual coverage to
error less than $\delta$ on buffered target balls.
The later edge/slim raw maps have these properties on every
radius-$100L$ ball used below. A meeting zero block has a unit
row $A_0$ with
\[
 |d(p_0,x)-d(p_0,p_i)-A_0u_i(x)|<E
                                  \quad(d(p_i,x)<1000),     \tag{TR0}
\]
where distances are in $R_i$ units.
\end{lemma}
\begin{proof}
For a listed circle chart, the center displacement is bounded
by $21$ independently of $\Delta$. MC10--MC11 and AC79
rescale and recenter its original two-splitting at $p_i$.
AC76, with reference rank two and supplied no-three-splitting
quality $\beta_3$, makes these splittings compatible. FC20--FC21
replace the Euclidean factor map by an orthogonal map.
All target accuracies and the radius $1000$ are early constants.
Thus the required reference and other circle quality bounds can
all be imposed on $\beta_2$ BEFORE $\Delta$.

For a listed edge/slim chart its displacement can grow with
$\Delta$, but its OWN quality is chosen later. Rescale/recenter
it with that displacement and apply AC76 with ranks one and two.
Its unit Euclidean factor, including the orthogonal projection,
is the stated coisometry. The reference threshold in AC76 uses
only the early desired compatibility accuracy and $\beta_3$;
only the recentered chart's threshold depends on the displacement.
For either rank choose the factor accuracy $\tau$ with
$\tau^{-1}>8000$ and $100\tau<E$. FC20 on a radius containing
all $u_i(B(p_i,1000))$ bounds the factor error by $48\tau$,
and the compatibility error by $\tau$. This proves (TR).
The required buffered distortion and coverage follow by decreasing
the same original qualities, in their own stages.
Pointed distortion localizes actual coverage witnesses, as in
SGP02; no attained infimum is assumed.

For (TR0) apply LC73/LCP04 at the reference point in the
buffered zero shell. Its real coordinate is exactly the
normalized distance difference. Apply the rank-one/rank-two
comparison just proved. Growing-ball curvature requirements
are supplied on the later uniform tail, not inferred from
a splitting quality alone.
\end{proof}

\begin{lemma}[Actual constant-radius block comparisons; TCP03]
\label{lem:fibration-first-constant-comparison}
For any early $0<\theta<1$, choose the early circle quality and
raw accuracies sufficiently small, then the later scalar qualities,
value tolerances and radial parameters. On all of $D_i$,
with norms in $R_i^{-2}g$, the SAME smooth coordinates satisfy
\[
 \|s_j\eta_j-\lambda_j\eta_i\|_{C^1}<\theta,\qquad
             \lambda_j(a)=A_ja+c_j .                       \tag{TC}
\]
For the meeting zero block use
$\lambda_0(a)=A_0a+s_0\eta_0(p_i)$; (TC) also holds for it.
All circle choices required here precede $\Delta$.
\end{lemma}
\begin{proof}
Take, for example, the circle quality
$\gamma<\theta^2/10^{12}$ and the early raw errors
$E,\delta<\theta^2/10^8$. These make $201\gamma<\theta/100$.
Later take the edge/slim qualities below $\theta^2/10^8$
and their separate LFR19 value errors below $\theta/100$.
The radial adapted quality may be made less than
$\min\{\theta^2/10^8,1/2000\}$ and its difference-Lipschitz
error less than $\theta/1000$. Retain all other local bounds.

For another circle block and a component $a$, lift the reference
product target $(u_i(x)+400A_j^*e_a,v_i(x))$.
For $x\in D_i$ its target radius is less than $412$, and its
actual witness lies in $B(p_i,420)$.
The segment length lies between $400-2\delta$ and
$400+2\delta$. It exceeds both original separation thresholds
$201$ and $201s_j$, and fits both original far-test balls.
The raw comparison (TR) applies at both endpoints.
The two component covectors saturate on the SAME minimizing
initial unit vector with error bounded by
$\max\{\gamma,\varsigma_{\rm edge},\varsigma_{\rm slim}\}
 +(8\delta+2E)/398$.
FC15's rowwise Riesz calculation, with $k\le2$, gives the
desired differential error.

For an edge chart use its OWN positive raw-axis lift of length
$200\Delta$ in its units; for a slim chart use length $10L$.
TCP01 puts $x$ within $36\Delta$ or $.91L$ of its center.
The edge endpoint lies within $250\Delta<1000\Delta$,
and the slim endpoint within $12L<L/\varsigma_{\rm slim}$
after the allowed later quality choice. Both satisfy their
original separation thresholds. Take $\Delta\ge10^6$.
On this long minimizing segment let $z$ be the point at
reference distance $400$ from $x$.
Then $z\in B(p_i,410)$, so the original CIRCLE test to $z$
fits its unchanged domain and exceeds $201$.
Coarse Lipschitzness of the long chart's raw real coordinate
transfers its almost maximal long gain to $z$ with loss at
most $8\delta$ in reference units. The alignment (TR) is
needed only at $x,z$, in the fixed radius-$1000$ ball.
FC22 and FC15 therefore give the same differential bound.
Neither original circle quality nor its tested radius depends
on the later long length.

The values follow independently from (TR) and the two separate
value errors, at cost less than $E+(1+s_j)\theta/100<\theta$.
For the zero block use a reference lift of length $400$ in
direction $A_0^*1$. LC73 at EVERY $x\in D_i$ has the same
raw distance difference, and the stated radial quality puts
this endpoint in its original all-direction test ball.
Its separation exceeds the original unit threshold.
The same saturation argument proves the differential bound.
LC67's centered difference-Lipschitz error is at most
$10\epsilon_0$ on $D_i$; together with (TR0) and the reference
value error this proves the value bound. No zero radius
multiplies an uncentered error.
\end{proof}

\begin{lemma}[Binding the SAME high-height coordinate; TCP04]
\label{lem:fibration-first-height-comparison}
If $J_e$ is empty, the weak-edge block vanishes on $D_i$.
Otherwise the actual data have either the low branch
$t<3\Delta/20$ throughout $D_i$, or a high branch on which,
for every prescribed early $0<\theta<1$, the permitted initial
choices produce a fixed unit row $B:\mathbb R^2\to\mathbb R$ with
\[
 \|t-\lambda_t\eta_i\|_{C^1(D_i)}<\theta,\qquad
                \lambda_t(a)=t(p_i)+Ba .                   \tag{TH}
\]
In the high branch the SAME $t=P/\rho$ is smooth near $D_i$,
$\Delta/10\le t<10\Delta$, and $\|Dt\|\le2$.
In the low branch all height factors are exactly $g=1,h=0$;
no smoothness or affine comparison of $t$ is required.
\end{lemma}
\begin{proof}
If no edge support meets the open $D_i$, the sum in the
weak-edge cutoff is zero there, so the whole block vanishes.
Otherwise a point of an actual positive edge cutoff lies in
$D_i$ and has $t<9\Delta$. TCP01 and the radius-$15\Delta$
support enclosure give $d(p_i,E)<30\Delta R_i$.
With $\Delta\ge1200$ and the later scale minima, FC14 applies.
It gives exactly the stated low branch or
$\Delta/10\le t\le181\Delta/20$ throughout $D_i$,
including the required strong-edge proximity.

In the high branch choose any $j\in J_e$ and an actual positive
support point $q\in D_i$. Then $|\eta_j(q)|<9\Delta$.
The original LFR19 Lipschitz bound gives
$|\eta_j(p_i)|<9\Delta+20/.99<10\Delta$ for large $\Delta$.
TCP01 already puts $p_i$ in its original radius-$100\Delta$
domain. Thus $p_i$ satisfies the ENTIRE band LFR36.1 of this
actual strong edge chart, including the height bounds.
Apply LFR38 AT THIS POINT, with its own metric $R_i^{-2}g$.
It provides an actual pointed $(2,\beta_2)$-splitting $\Psi$
and the SAME pair $(\eta_j,t)$ adapted on $B(p_i,100)$
of an arbitrarily small EARLY prescribed quality $\gamma_h$.
No height function or smoothing is replaced.

Choose $\gamma_h<\theta^2/10^{12}$ before $\beta_2$ and
require $\beta_2<\gamma_h/1000$, as well as TCP02's
compatibility threshold. LFR35--LFR38 then choose the later
$\Delta$ and endpoint, smoothing and tangential errors.
This order also retains their full quotient derivative for
$t=P/\rho$. In particular its gradient bound is a conclusion
about $t$, not one about $P$ mistakenly applied to a quotient.

Compare $\Psi$ with the original circle splitting using AC76,
the same no-three-splitting hypothesis and FC20.
Both are based at $p_i$ in exactly $R_i$ units.
There is a fixed orthogonal $A$ with
$|\Psi-Au_i|<E$ on the fixed radius-$1000$ ball; impose
$E,\delta<\theta^2/10^8$ at the earlier reference choice.
Put $B=e_2^*A$. Lift the reference raw-axis target in direction
$B^*1$ at length $400$. LFR35.1's actual scale-$100$ test
allows every endpoint within $100/\gamma_h$ at separation
greater than $100$. The original circle test allows this same
endpoint at separation greater than $201$. Both apply at every
$x\in D_i$. Their scalar saturation on the SAME initial vector
and the rowwise FC15 estimate give
$\|Dt-BD\eta_i\|<\theta/10$ after the displayed strict budgets
are decreased if necessary. The reference coordinate vanishes
at $p_i$. Integrating along a minimizing radial segment in
$D_i$ gives the centered value error less than $\theta$.
This proves (TH), and the adapted Lipschitz bound gives
$\|Dt\|<1+\gamma_h<2$ there.

Only the HEIGHT row of the pair has been used here.
Its tangential member $\eta_j$ is normalized at $p_j$, whereas
both displayed adapted tests use the point's OWN scale $R_i$.
The constant-radius block is still handled separately by
$s_j\eta_j$ in TCP03. No unnoticed constant scale factor is
inserted in the shared height component.
\end{proof}

\begin{theorem}[The actual first graph and its early modulus; TCP05]
\label{thm:fibration-actual-first-graph}
Let $A=\max_{0\le n\le N}A_n$ and
$B_0=\max_{0\le n\le N}B_n$, using FC10's constants with
the fixed profile bound $P$. Put
\[
 C=10^4(N+2)^2(P+1+A+B_0).
\]
This is fixed before $\Delta$ and noncollapse. For every early
$0<e<1/100$, the permitted choices above produce a smooth
$\Phi_i:\mathbb R^2\to H$ with own block $(a,1)$ and
\[
 \|D\Phi_i\|,\|D^2\Phi_i\|\le C,\qquad
 \|R_i^{-1}F-\Phi_i\eta_i\|_{C^1}<e
                      \quad\hbox{on }\{|\eta_i|\le8\}.       \tag{TG}
\]
It is an approximation to the original global zero-extended
map, with CFS27's prescribed pruned plane recipe.
\end{theorem}
\begin{proof}
Choose $\theta<e/(100C^2)$ in TCP03--TCP04 and impose
$\Delta\Lambda<e/(1000C)$ at the later scale choice.
All required circle/collar qualities and raw reference errors
therefore remain EARLY constants. Use the own block $(a,1)$,
the constant scale coordinate $1$, and for the other circle
and slim blocks substitute $\lambda_j(a)$ into their original
scaled cutoff formulas. Use $F_{s_0}(\lambda_0(a))$ for the
possible zero block. Unlisted model blocks are zero.

If there are no edge supports, or the low branch of TCP04
holds, set the weak-edge block to zero and use the independent
tangential edge blocks with their height factors exactly one.
In the high branch use the WHOLE joint network of FC10 with
affine inputs
\[
 u_j(a)=\lambda_j(a)/s_j\quad(j\in J_e),\qquad
 v(a)=\lambda_t(a).
\]
Its actual fixed-scale inputs are $(\eta_j,t)$.
The actual sum cutoff includes no other index on $D_i$:
every omitted edge block and cutoff vanish there. This network
is therefore the same coupled cutoff as the original $F$.

Each affine row has norm at most two, so its combined linear
part has norm at most $2\sqrt{N+1}$.
FC10 bounds its first derivative by $2\sqrt{N+1}A$
and Hessian by $4(N+1)B_0/\Delta$.
The remaining scaled blocks have first and second derivative
bounds at most $50(P+1)$ by FC05 and $s_j>.99$.
The zero radius is at least one and only improves its Hessian
bound. Square summation, including the identity, is less than
$C$. These estimates are global on the affine input plane,
independent of translations or target dimension.

On $D_i$ the actual/model input difference of the joint network
is at most $2\sqrt{N+1}\theta$ in $C^1$, while the derivative
of its model input composed with $\eta_i$ is at most
$4\sqrt{N+1}$. FC10, FC05 and TCP03--TCP04 therefore bound
all fixed-scale block errors by less than $C^2\theta$.
The remaining variable output factor is
$\sigma=\rho/R_i$, multiplying ONLY the weak-edge block
$(t\zeta_{E'},\zeta_{E'})$; the scale component is itself
$\sigma$. On $D_i$ one has
$|\sigma-1|\le10\Lambda$, $|D\sigma|\le\Lambda$.
FC10 and the actual input derivative bounds give the fixed-scale
weak-edge block value at most $C\Delta$ and derivative at most
$C$. FC11 thus bounds this product error by
$\Lambda(11C\Delta+10C)$; the scale-coordinate error is at
most $10\Lambda$. The total is less than $30C\Delta\Lambda$.
This explicitly retains the derivative of the variable scale.
TCP04 already compared the actual QUOTIENT $t$; no further
unproved comparison of $P/\rho$ is being inferred from values.

The sum $C^2\theta+30C\Delta\Lambda$ is less than $e$.
On the entire original set $|\eta_i|\le8$, TCP01 puts every
point in $D_i$ and its own block is exactly $(\eta_i,1)$,
including its derivative. Thus these estimates prove (TG).
Every listed nonzero-family radius is greater than $.99R_i$;
all other model blocks already vanish. CFS27's deletion of
radii at most $R_i/2$ leaves this actual graph model unchanged.
\end{proof}

\begin{theorem}[The actual first cloudy packet, with all-preimage rank; TCP06]
\label{thm:fibration-actual-first-cloud}
For $0<\Gamma<1$ choose before $\Delta$
\[
 0<\Sigma<\min\{\Gamma/200,\Gamma^3/(100C)\},\qquad
 0<e<\min\{1/100,\Gamma\Sigma/100,\Sigma/1000\}.              \tag{TP}
\]
Use TCP05 and FC04/FC07's ORIGINAL core and enlargement.
If $I_2\ne\varnothing$, $(S_1,\widetilde S_1,r)$ is an
actual $(2,\Gamma)$ cloudy TWO-manifold in the identical-radius
truncated convention, with the exact radius
$r(x)=\Sigma x_\rho$.
At $x=F(p)\in S_1$ choose an original core witness $i,p$ and set
\[
 T_x=D\Phi_i(\eta_i(p)),\qquad A_x=x+\operatorname{im}T_x.
\]
At EVERY $q$ with $F(q)=x$, the derivative projected to $A_x^0$
is onto, its two nonzero singular values lie in $[1/2,3C]$,
and its normal error is less than $e<\Gamma$.
These planes, radii and graphs satisfy CFS14's geometric inputs
when its smoothing threshold on $\Gamma$ and
$\Sigma\le1/(640b)$ are imposed. If $I_2$ is empty, the
first adjustment is the identity.
\end{theorem}
\begin{proof}
FC04 makes the radius well-defined on the ENTIRE enlarged
image and gives its Lipschitz radius inequality directly from
the actual scale coordinate. The finite original support
enclosures and compact carrier give bounded image and upper
and positive lower radius bounds. Every enlarged image point
has a full circle marker, so CFS26 and CFS07 with $B=5/3$
also provide the all-center scale comparison used by CFS14.

At every preimage of a core point the full $i$ marker forces
membership in its original domain and the vector block gives
$\eta_i(q)=\eta_i(p)=a$, with $|a|\le7$.
Thus (TG) holds at every such $q$ with exactly the SAME $T_x$.
TCP01's Gram bound gives least singular value greater than
$9/10$ and operator norm less than two for $R_iD\eta_i$
in the physical metric. The identity component of $T_x$
gives its least singular value at least one and upper norm $C$.
Apply FC06 to (TG), canceling $R_i$ between source and target.
Its projected adjoint has lower norm $9/10-e>1/2$,
so the projection is actually onto the two-plane.
Its upper norm is less than $3C$ and normal error less than $e$.
This proves rank, not just an inequality away from a possibly
full kernel.

For coverage use $R_i$ units, $\widehat r=r(x)/R_i$ and
$R=\widehat r/\Gamma$. Full markers give
$3\Sigma/4\le\widehat r\le5\Sigma/4$, hence $R<1/100$.
For any $z\in\widetilde S_1$ with $|z-x|\le r(x)/\Gamma$,
FC03 at EVERY preimage $q$ gives
\[
 |\eta_i(q)-a|\le 8R/(1-R)<1/10.
\]
Consequently its $i$ marker is full and (TG) applies with
exact vector coordinate $\eta_i(q)$. Conversely every
$u\in\overline B(a,R)$ has $|u|<8$; LFR07's original
proper trivial circle bundle over $B(0,100)$ supplies an actual
point with $\eta_i=u$. It belongs to the original enlargement,
and its image is within $e$ of $\Phi_i(u)$.
FC25 now gives both coverage directions with error
$3q$, $q=2e+CR^2/2$. The bounds (TP) imply
$3q/\widehat r<\Gamma/2$ exactly as in SGP06.
The same radial trimming and strict slack give the open-ball
distance-to-set convention as well. This proves the cloud
for precisely the planes already used for normal and rank
estimates. Empty families require only the identity adjustment.
\end{proof}

\paragraph{Consumer and parameter boundary after revision148.}
TCP01--TCP06 supply the actual first-cloud producer, and
SGP01--SGP06/EGP01--EGP07 supply both projected producers.
All three original cloudy packets, rough graphs, early moduli,
all-preimage rank and exact coverage directions are now written
under the retained foundations. The high-height argument reuses
LFR38's SAME pair and actual splitting; it does not repeat the
completed local smoothing theory or introduce a new function.
All fixed short radii and circle/collar accuracies precede
$\beta_2$ and $\Delta$. Later long-coordinate tests remain on
their original domains, and $\Delta\Lambda$ is charged explicitly.

The complete three-stage adjusted-map binding, proper WHOLE
preimage restrictions and isotopy to original fibers, simultaneous
faces/corners, the exhaustive FC44 register and nearly cuspidal
boundary assembly are still separate unfinished parts of the
active goal. The parent static theorem and FC33 are not promoted.
No accepted-PC or Lean binding is claimed. This is an intermediate
development, not the final full-goal mathematical audit.

\subsection{Binding the adjustments, exact full markers and whole closed fibers}
\label{sec:fibration-actual-adjustment-and-closed-fibers}
Revision149 uses the three actual original-cloud producers.
Throughout this subsection $M$ is the CLOSED smooth carrier of the
standing construction, $F=\mathcal E^0$ is CGP01's same map, and
all retained geometric and smooth-topological foundations remain
explicit. Write $\mathcal I_1=I_2$, $\mathcal I_2=I_e$,
$\mathcal I_3=I_s$, with $k_j=2,1,1$, respectively, and put
$\ell_i=1,\Delta,10^5\Delta$ in these three families.
Let $u_i,v_i$ denote the vector and marker projections of a
retained block. The source domains and all original maps are
unchanged.

\begin{proposition}[One actual choice for the three adjustments; GAF01]
\label{prop:fibration-actual-adjustment-choices}
For every early target $c_{\rm adjust}>0$ there is one choice,
in the retained order, supplying the actual hypotheses of CFS31,
CFS20 and CGP07 for ALL three stages. One may additionally require
\[
 \Sigma_j\le\varepsilon_j/10000,\qquad
 c_3<\min\{c_{\rm adjust},1/1000,1/512\}.                 \tag{JA}
\]
These choices use the actual SGP, EGP and TCP producers, not
extra rough-graph or all-preimage rank assumptions.
This binds the adjustment and marked-base consumers; constraints
from later boundary extraction are not yet included.
\end{proposition}
\begin{proof}
Choose one early $N$ dominating the original pointwise multiplicity
and the three whole-support counts, and one early $P\ge1$ dominating
all fixed profile constants. Let $C_j$ be the actual graph moduli
in TCP05, EGP06 and SGP04, and take
$\Omega=\max\{1,C_1,C_2,C_3\}$.
CGP02 gives an early global derivative bound $L_0\ge1$.
Use the common projected rank margin $\mu=1/2$ from
TCP06, EGP07 and SGP05. Their original coordinate maps also
have right inverses of norm less than two on the needed domains.
The actual cutoff constants are
\[
 b_{\rm cut}=10000(N+1)^2P^4,\qquad
                  \kappa=\frac1{1000(N+1)P^2}.
\]
These numbers all precede every cloud accuracy and $\Delta$.

Start with $c_3$ as in (JA). At stage $j$, whose target $c_j$
is now fixed, put
\[
 \alpha_j=\min\left\{
 \frac{c_j}{16(1+b_{\rm cut})(1+L_0)},
 \frac{1}{16(1+L_0)}\right\}.
\]
Use CFS15 to choose $\Gamma_j$ so small that
$\Gamma_j\le c_j/16$ and
\[
 \varepsilon_j=\Xi_j(\Gamma_j)
 <\min\{1/10,\alpha_j,1/[1000(\Omega+1)]\}.
\]
Choose
\[
 \begin{split}
 0<\Sigma_j&<
 \min\{1/2,\varepsilon_j/10000,\Gamma_j/200,
                              \Gamma_j^3/(100C_j)\},\\
 0<e_j&<
 \min\{1/100,\Gamma_j\Sigma_j/100,\Sigma_j/1000,
                                                1/(200\Omega)\},\\
 t_j&=\min\{3\Sigma_j/10,\alpha_j,1\}.
 \end{split}                                                \tag{JB}
\]
After the stage-three choices take a positive
$c_2\le\min\{c_3,t_3,4\kappa/5,1/1000,1/512\}$;
after the stage-two choices take
$c_1\le\min\{c_2,t_2,4\kappa/5,1/1000,1/512\}$.
All these are minima of positive earlier quantities.

Now impose the actual local-producer requirements. TCP01--TCP06
choose the early circle/collar accuracies and $\beta_2$ below
their required thresholds, before $\Delta$.
Choose $\Delta$ large enough for those fixed short buffers and
the completed local collar and section constructions.
At the later coordinate, endpoint, smoothing and scale stages,
take the joint minima for SGP01--SGP06, EGP01--EGP07,
TCP01--TCP06 and CGP01--CGP03, using the requested $e_j$.
In particular retain their separate value errors,
$\Lambda10^6\Delta<10^{-5}$ and their charged
$\Delta\Lambda$ bounds. The later $\beta_E$, then $\beta_1$,
zero witnesses and uniform curvature tail meet the remaining
requirements. The individual producer proofs verify that these
bounds use only allowed predecessors; none makes the circle
quality depend on a later long length.

Thus the ACTUAL clouds have rough error $e_j$, normal error
at most $e_j<\Gamma_j$, and the stated positive rank margin.
CFS15 applies with the stronger radius bound (JA).
CFS31's marker and cutoff hypotheses are now supplied by the
actual map, profiles and planes. CFS20 has exactly the displayed
$\alpha_j,t_j,b_{\rm cut},L_0,\mu$.
Moreover $2e_j<1/(48\Omega)$, and (JB) supplies all of
CGP07's (OS) bounds. The CFS14 construction retains its full
convex graph bounds as required there.
No compact-carrier-specific cutoff tolerance has been inserted
before the carrier is constructed. Empty families are assigned
the identity stage; their unused positive numerical choices
cause no difficulty.
\end{proof}

\begin{theorem}[The actual adjusted map and embedded bases; GAF02]
\label{thm:fibration-actual-final-submersions}
With GAF01's choices there are smooth adjustments on open
neighborhoods of the successive actual images and a smooth map
\[
 E=\mathcal E^3=\Psi_3\Psi_2\Psi_1F:M\longrightarrow H
\]
such that
\[
 |E-F|<c_3\rho,\qquad \|DE-DF\|<c_3<c_{\rm adjust}.
                                                               \tag{AE}
\]
All intermediate cumulative errors are less than their $c_j$.
For every $i\in I_+$ and every intermediate map $f$,
\[
 \zeta_i(p)=0\ \Longrightarrow\ |v_i(f(p))|\le R_i/32 .
                                                               \tag{AM0}
\]
The same holds along the segment from $F(p)$ to $E(p)$.

FC33's EXACT original threshold-$5$ domains $U_j$ have smooth
submersions $\pi_jE:U_j\to W_j$ onto their images in embedded
bases of dimensions $2,1,1$. Each marked patch has the coordinate
diffeomorphism of CGP07 over $B(0,5.5\ell_i)$.
The later maps are diffeomorphisms on the entire earlier marked
bases and preserve fibers and kernels on CGP08's specified
threshold-$6$ carriers. No whole-preimage claim is implicit in
this last restricted statement.
\end{theorem}
\begin{proof}
Execute the three stages in increasing order, using the original
clouds supplied in GAF01. CFS31 first constructs the source
cutoff, then the edge cutoff and then the slim cutoff. At each
step its CLOSED support localizes the original point to the
corresponding original core. CFS16 puts the perturbed input
inside that core's half-tube. CFS18 therefore supplies an actual
smooth neighborhood on which the required postcomposition is
defined; a global extension to all of $H$ is unnecessary.

CFS20 gives the cumulative value and derivative bounds and
submersions on the open threshold-$6$ plateaux.
Its bounds are less than $c_j$, and the preceding target is at
most $t_j$, so induction verifies the next domain and estimate.
CFS29--CFS30 give (AM0), including the straight-segment assertion.
For an empty family the stage is the identity; the previous
error is already below its permitted target.

All actual rough graphs, rank estimates, original inner sections
and (OS) inequalities required by CGP07 are now present.
That theorem gives each WHOLE marked stage base as a single
graph over the full $5.5\ell_i$ ball, exhausted by actual
threshold-$6$ source images. Apply CGP08 to their actual later
postcompositions. It gives the embedded $W_j$, including
injectivity across different marker patches, and the stated
threshold-$5$ submersions. Its factorization proves restricted
fiber and kernel equality. The final estimate is CFS20's
stage-three estimate, which is stronger than (AE).
This proves the written FC33 conclusion under the retained
foundations, rather than leaving its rough-cloud inputs as
additional assumptions. Original fiber type and full-preimage
restrictions still require the arguments below.
\end{proof}

\paragraph{Source comparison for the edge base.}
KL (13.45), printed78/PDF73, puts conditions on $y_r,y_{E'}$
in a set declared to lie in $Q_2$. But the source's own
Section12.4, printed69/PDF64, and FC01 omit those coordinates
from $Q_2$. We use CGP08's well-defined marker-selected $W_2$;
the original height condition stays in the SOURCE domain $U_2$.
GAF02 proves the requested submersion with that explicit
translation. It does not evaluate discarded coordinates or
claim identity with the ill-typed displayed source set.
This is a project comparison, not an author-listed erratum.

\begin{lemma}[Affine marker locality of the spectral construction; GAF03]
\label{lem:fibration-affine-marker-locality}
In CFS11--CFS14 let $J$ be orthogonal projection onto a
coordinate subspace, and let $c$ be a fixed vector in that
subspace. Suppose EVERY selected center $y$ whose closed
$80br_y$ support meets $B(x,8br_x)$ satisfies
$Jy=c$ and $L_y\subset\ker J$. Then
\[
 J\eta(z)=Jz-c\quad(z\in B(x,8br_x)),\qquad
 Jw=c\quad(w\in W\cap B(x,3br_x)),
\]
and $JP(z)=c$ for $z\in B(x,r_x)$.
\end{lemma}
\begin{proof}
As in CFS24, each contributing normal projector is the identity
on the coordinate subspace and has no mixed block. Therefore
the spectral projector satisfies $JQ=J$. The normalized weights
sum to one, so their mean center has $J\mu=c$.
Thus $JQ(z-\mu)=Jz-c$. Its zero set gives the assertion about
$W$. CFS14 puts $P(z)$ within $(1+\varepsilon)r_x<3br_x$
of $x$, proving the last assertion. This is the affine,
constant-coordinate version of CFS24; it still requires EVERY
contributing center and plane.
\end{proof}

\begin{lemma}[Actual full-marker control at every contributor; GAF04]
\label{lem:fibration-actual-full-marker-contributors}
Fix stage $j$ and $i\in\mathcal I_j$. Let
$x=\pi_jF(p)$ with $|\eta_i(p)|\le7\ell_i$ and, for an
edge index, $t(p)\le7\Delta$. Thus $x$ is an original core
center with full $i$ marker. With $\Sigma_j\le\varepsilon_j/10000$,
EVERY selected smoothing center $y$ whose closed $80br_y$
support meets $B(x,8br_x)$ satisfies
\[
 v_i(y)=R_i,\qquad v_i|_{L_y}=0.                         \tag{FM}
\]
Here $b=\varepsilon_j^{-1}$ and $L_y$ is the ACTUAL plane
chosen by the SGP, EGP or TCP producer and CFS27.
Consequently GAF03 applies with $J=v_i$ and $c=R_i$.
\end{lemma}
\begin{proof}
The support intersection gives
$|y-x|<80br_y+8br_x\le88b\max\{r_y,r_x\}$.
CFS07 with its actual marker hypotheses gives
$r_y/r_x\le5/3$. The preimage selected for $r_x$ has the
same full $i$ marker, so CFS26 bounds its scale by $5R_i/4$.
Hence
\[
 |y-x|<
 \left(80\frac53+8\right)b\Sigma_j\frac54R_i
 =\frac{530}{3}b\Sigma_jR_i<R_i/50 .                     \tag{FD}
\]
This controls the WHOLE contributor list, independently of the
preimages selected for radii and for models.

Let $q_y$ be the actual original core preimage used to choose
the model-reference chart $a$ at $y$.
Its $i$ marker exceeds $.98R_i$, so its original cutoff is
positive and its original $i$ coordinate exists. Marker division
and (FD) give
\[
 |\eta_i(q_y)|
 <\frac{7\ell_i+1/50}{1-1/50}
 \le\frac{351}{49}\ell_i<7.2\ell_i .                     \tag{FV}
\]
For circles and slim charts this is in the full plateau.
For edge charts $q_y$ belongs to the original edge CORE, so
the SAME height has $t(q_y)\le7\Delta$ regardless of its
model-reference index. Both $i$ cutoff factors are therefore
one. This proves $v_i(y)=R_i$ exactly.

The actual core point $q_y$ lies inside the open reference
comparison ball $D_a$ of the corresponding producer.
The positive $i$ cutoff therefore places $i$ in its ACTUAL
whole support list. Its ratio $s_i=R_i/R_a$ is in
$(.99,1.01)$. If $i=a$, the model marker is identically one.
Otherwise the actual affine comparison used in
SGP04, EGP06 or TCP05 gives at its model parameter
$a_y=\eta_a(q_y)$
\[
 \big|\lambda_i(a_y)/s_i-\eta_i(q_y)\big|<1/50 .
\]
Indeed their original input comparison tolerances are less
than $1/100$, and $s_i>.99$.
By (FV) this affine cutoff argument has norm less than
$7.22\ell_i<8\ell_i$. The model $i$ marker is thus the
CONSTANT $s_i$ on a neighborhood of $a_y$.
In the edge model its height factor has already been removed
exactly on the original enlarged carrier; the model marker
here is its tangential factor as proved in EGP06.
Therefore its derivative at $a_y$ vanishes.

CFS27 does not delete this block, since $s_i>.99$.
The physical tangent plane, obtained from this actual model,
lies in the scalar marker kernel. This proves the second part
of (FM). The argument treats every contributing center and its
actual chosen affine plane, not merely the reference chart's
own identity component.
\end{proof}

\begin{proposition}[Exact full markers on the bases and source plateaux; GAF05]
\label{prop:fibration-exact-interior-markers}
Every marked stage patch $V_i^0$ of CGP07 satisfies
$v_i=R_i$ identically. Its later image $V_i\subset W_j$
has the same property and
\[
 a=u_i/R_i=u_i/v_i:
                V_i\longrightarrow B(0,5.5\ell_i)
\]
is a diffeomorphism. Moreover, on the ORIGINAL threshold-$6$
source plateau for index $i$ (including $t<6\Delta$ for edges),
both its stage-$j$ image and its final image have marker
exactly $R_i$.
\end{proposition}
\begin{proof}
For $w\in V_i^0$ put $a=u_i(w)/R_i$ and use CGP03's
actual inner section to set $x=\pi_jF(s_i(a))$.
This is an original core point with full $i$ marker and
$|\eta_i(s_i(a))|<5.5\ell_i$; the edge section has height
less than $\Delta/100$. CGP07's estimate (SN) puts
$w$ in $W\cap B(x,r_x/4)$. GAF04 and GAF03 imply
$v_i(w)=R_i$. CGP07 already gives its coordinate
diffeomorphism, so replacing $R_i$ in the denominator by
this exact marker changes neither the map nor its derivative.
Later stages retain the entire earlier family's block.
CGP08 therefore gives the identical conclusion on $V_i$.

For an original threshold-$6$ source point, $x=\pi_jF(p)$
again meets GAF04. The preceding perturbed input is in
$B(x,r_x)$ by GAF02/CFS16. Its actual cutoff equals one,
so the stage output is $P_j$ of that input. GAF03 gives
marker $R_i$, and later stages retain it.
No approximate differentiation of a quotient is used:
the marker is CONSTANT on the indicated smooth base patch.
In particular no error of order $\ell_i c_3$ is incurred.
\end{proof}

\begin{lemma}[Whole-preimage localization along the global segment; GAF06]
\label{lem:fibration-whole-ratio-preimage-localization}
Put $H_\tau=(1-\tau)F+\tau E$, $0\le\tau\le1$.
For ANY original $i\in I_+$ and ANY $p\in M$, if
\[
 v_i(H_\tau(p))>.9R_i,\qquad
 \frac{|u_i(H_\tau(p))|}{v_i(H_\tau(p))}\le4\ell_i ,
                                                               \tag{RP}
\]
then its original coordinate is defined and
\[
 |\eta_i(p)|<4.01\ell_i .
\]
For circles and slim charts its original cutoff is consequently
one. For edge charts this assertion controls the tangential
coordinate only; an additional height restriction is still needed.
\end{lemma}
\begin{proof}
If the original cutoff were zero, GAF02's segment estimate
would give $|v_i(H_\tau(p))|\le R_i/32$, contrary to (RP).
Thus it is positive, and CFS26 gives
$3R_i/4\le\rho(p)\le5R_i/4$. In $R_i$ units the vector and
marker errors from their original values are at most
$\delta=(5/4)c_3<1/800$.
Writing $\zeta=\zeta_i(p)$, we obtain
$\zeta>.9-\delta>719/800$. The ratio inequality implies
\[
 |\eta_i(p)|\zeta
 \le4\ell_i(\zeta+\delta)+\delta .
\]
Since $\ell_i\ge1$,
\[
 |\eta_i(p)|
 <4\ell_i+\frac{4\ell_i+1}{719}
 \le(4+5/719)\ell_i<4.01\ell_i .
\]
This excludes every remote component of the indicated full
preimage, uniformly in time; it did not assume membership in
an earlier restricted carrier.
\end{proof}

\begin{theorem}[Proper whole circle and slim bundles and their original fibers; GAF07]
\label{thm:fibration-whole-closed-fiber-bundles}
For $j=1,3$ define the actual open base
\[
 B_j=\bigcup_{i\in\mathcal I_j}
       \{w\in W_j:v_i(w)>.9R_i,\ |u_i(w)|/v_i(w)<4\ell_i\},
 \qquad X_j=(\pi_jE)^{-1}(B_j).
\]
Then
\[
 \bigcup_{i\in\mathcal I_j}\{|\eta_i|\le3.5\ell_i\}
             \subset X_j\subset U_j .
\]
The map $\pi_jE:X_j\to B_j$ is a proper onto smooth bundle
with connected fibers. For $j=1$ these are smooth circles;
for $j=3$ they are smooth $S^2$ or $T^2$ fibers.
Each WHOLE fiber is smoothly isotopic, inside its original
buffered chart, to a fiber of the SAME original coordinate
$\eta_i$. Later adjustments preserve the corresponding
stage-$j$ fiber as a WHOLE subset of $M$, not just its
intersection with a previously specified carrier.

These are the proper open-base restrictions before removal of
zero domains and before corner extraction. No edge-disk or
boundary-face conclusion is asserted here.
\end{theorem}
\begin{proof}
First check that the ratio-defined pieces introduce no extra
base points. Every $w\in W_j$ has an actual final preimage,
by CGP07--CGP08. If $v_i(w)>.9R_i$, (AM0) forces that
preimage's original $i$ cutoff to be positive. CFS26 and (AE)
therefore give $v_i(w)<(1+1/800)R_i$. The ratio threshold
then implies $|u_i(w)|<4(1+1/800)\ell_iR_i<5.5\ell_iR_i$,
so $w$ belongs to the original marked patch $V_i$.
By GAF05 each base piece is consequently exactly
$\{w\in V_i:|u_i(w)/R_i|<4\ell_i\}$, a coordinate ball
in the embedded base. Thus $B_j$ is open.
At an original $|\eta_i|\le3.5\ell_i$ point, GAF02 puts
the final image in $W_j$, GAF05 makes its marker $R_i$,
and (AE) changes its normalized vector by less than $1/800$.
Its coordinate is therefore less than $4\ell_i$, proving the
first inclusion. For EVERY point in the full preimage of $B_j$,
choose an index witnessing base membership. GAF06 puts it in
the original $|\eta_i|<4.01\ell_i<5\ell_i$ domain, proving
$X_j\subset U_j$ without discarding any source component.

GAF02 now gives a submersion on this entire open set.
Each marked base point is the final image of an actual source
point by CGP07--CGP08, proving onto-ness.
For a compact $K\subset B_j$, its inclusion in the finite
Euclidean target is compact and closed. The FULL preimage
$(\pi_jE)^{-1}(K)$ is therefore closed in compact $M$,
and it lies in $X_j$. This proves properness of the restriction.
The inherited smooth proper-submersion theorem supplies its
local trivializations.

We identify the WHOLE fiber, including connectedness.
Fix $w\in B_j$ and an index $i$ with
$a=u_i(w)/R_i$, $|a|<4\ell_i$.
On the original open domain
$Y_i=\{|\eta_i|<5\ell_i\}$ define
\[
 g_i=u_i(\pi_jE)/R_i,\qquad
                  h_\tau=(1-\tau)\eta_i+\tau g_i .
\]
Every point of $Y_i$ has full original marker, so (AE) gives,
with source norm in $R_i^{-2}g$,
\[
 |g_i-\eta_i|<1/800,\qquad
                        \|Dg_i-D\eta_i\|<c_3 .
\]
The original coordinate has least singular value greater than
$9/10$, by TCP01 or the actual slim axis test of SGP03--SGP05.
Every $h_\tau$ is therefore a submersion on $Y_i$.
If $h_\tau(p)=a$, then
$|\eta_i(p)|<4\ell_i+1/800<4.01\ell_i$.
The SAME compact set
\[
 Q_i=\{p\hbox{ in the original chart}:|\eta_i(p)|
                                                   \le4.01\ell_i\}
                   \Subset Y_i
\]
contains ALL these inverse images, with a strict interior
buffer. Its compactness follows from the original proper
circle bundle of LFR07 or the whole proper slim bundle of
LFR20. In particular the entire time-dependent trace is
compact, not merely its endpoint fibers.
Apply the inherited smooth compact isotopy theorem in FC34
to $h_\tau$ and the point $\{a\}$ (equivalently, the smooth
case of LFR03's explicit transport construction).
There is no target boundary in this application.
It identifies $h_1^{-1}(a)$ with the original smooth fiber,
which is $S^1$, or $S^2$ or $T^2$, respectively.

It remains to check that this is the desired WHOLE fiber of
$\pi_jE$. GAF06 already puts every preimage of $w$ in
$Y_i$, where its vector coordinate is $a$.
Conversely, if $p\in Y_i$ and $g_i(p)=a$, GAF02 puts
its final image in $W_j$ and GAF05 gives marker $R_i$.
It is in $V_i$ and has coordinate $a$, so the coordinate
diffeomorphism in GAF05 forces its image to equal $w$.
Thus $h_1^{-1}(a)=(\pi_jE)^{-1}(w)$ as whole subsets of
$M$. This proves the smooth fiber type and connectedness.

Finally let $f_j=\pi_j\mathcal E^j$ and let
$\Theta_j:V_j^0\to W_j$ be CGP08's later diffeomorphism.
On the whole stage image, later postcomposition gives
$\pi_jE=\Theta_j f_j$ wherever $f_j$ is in this marked base;
the ambient later map itself is defined on every stage image.
If $f_j(p)=\Theta_j^{-1}(w)$, that postcomposition gives
$\pi_jE(p)=w$, without a prior restriction on $p$.
In the reverse direction GAF06 puts every such final preimage
inside the original threshold-$5$ chart, where CGP08 puts
$f_j(p)$ in $V_j^0$. Injectivity of $\Theta_j$ then forces
$f_j(p)=\Theta_j^{-1}(w)$.
These two inclusions prove equality of the WHOLE stage and
final fibers. The same factorization gives kernel equality there.
\end{proof}

\paragraph{Boundary after revision149.}
GAF01--GAF02 bind the actual three-stage adjustments and the
written FC33 conclusion. GAF03--GAF05 prove exact full markers
from ALL actual contributors, removing the large-$\Delta$
ratio-derivative problem on the marked bases.
GAF06--GAF07 give proper WHOLE circle/slim restrictions and
smooth isotopy to the original local fibers, before the successive
piece removals. Their uniform thresholds follow the actual
early/late order; no per-manifold closeness tolerance replaces it.

Zero-domain extraction, the height-bounded proper edge disks,
compatibility with the removed regions and all faces/corners,
global graph presentation, the remaining exhaustive parameter
register and nearly cuspidal boundary assembly are unfinished.
These remain part of the active goal. No accepted-PC or Lean
binding is claimed, and this is not its final mathematical audit.

\section{Proper restrictions, topology and corners}
\label{sec:fibration-proper-restrictions}
\begin{lemma}[Compact transverse transport interface; FC34]
\label{lem:fibration-compact-transport}
Let $Y,N$ be smooth manifolds without boundary and
$X\subset N$ a smooth submanifold with boundary. Use a smooth
family $f_t:Y\to N$, transverse to both $X$ and
$\partial X$ for every $t\in[0,1]$. Require the ENTIRE set
$\{(y,t):f_t(y)\in X\}$ to be compact, in one common interior
buffer. The inherited compact isotopy theorem transports
$f_0^{-1}X$ to $f_1^{-1}X$. For a perturbation, uniform $C^1$
smallness supplies transversality on that buffer only after a
positive quantitative margin is fixed; a separate exclusion of
new preimages outside it supplies the common compact set.
\end{lemma}
\begin{proof}
The transverse total preimage is a manifold with corners,
including its time-end faces and the face over $\partial X$. Its projection to $[0,1]$ is a proper submersion:
surjectivity of the normal derivative solves for a tangent lift
of $\partial_t$. Integrate a boundary-tangent lift on the compact
set, using the inherited isotopy-extension result when an ambient
isotopy is required. LC82 records the indispensable isolation
condition. No global perturbation claim follows from closeness
only on an arbitrary compact neighborhood.
\end{proof}
This is the safe application interface for KL 21.1--21.3.
The generic calculus, proper submersion and extension theorems
are inherited audit targets, not newly constructed foundational
geometry. For a domain with corners, require tangency and
transversality on EVERY face and their intersections.

\begin{foundation}[Zero-domain extraction; FC35]
\label{found:fibration-zero-extraction}
For each zero center retain KL (14.3)'s core ball
$B(p_i,.35R_i)$ together with the adjusted marker sublevel
$x''_i\ge.9R_i$, $x'_i/x''_i\le.4$.
The resulting compact domains are disjoint, have smooth boundary,
and are diffeomorphic to the corresponding LC80 models.
They contain $B(p_i,.38R_i)$ and lie in $B(p_i,.42R_i)$,
with the small-error restrictions of KL 14.4.
\end{foundation}
On the original shell $\eta_i\in[.3,.8]$, KL 12.39 says the only
nonzero $Q_4$ block is $(R_i\eta_i,R_i)$. Interpolate the
adjusted ratio with the SAME original radial function, extend
inside/outside the shell, and apply FC34 at level $.4$.
The core ball in the definition is necessary because the cutoff
vanishes near the center. Keep the source and author-corrected
alternatives: $D^3$, solid torus, twisted interval bundles over
$\mathbb{RP}^2$ and the Klein bottle, and the closed
$S^1\times S^2$, $\mathbb{RP}^3\#\mathbb{RP}^3$, orientable
flat and spherical space forms. Each alternative retains its
boundary and category, rather than an undifferentiated finite list.

\begin{foundation}[Proper slim bundles; FC36]
\label{found:fibration-slim-bundles}
Truncate $W_3$ by the union of marker conditions $x''_i>.9R_i$
and $|x'_i/x''_i|<4\cdot10^5\Delta$ and take its FULL preimage
under $\pi_3\mathcal E$. KL 14.7 supplies a proper onto
submersion, containing the $3.5\cdot10^5\Delta$ local regions
and contained in $U_3$, with fibers $S^2$ or $T^2$.
Remove the interiors of zero domains and choose a compact
one-dimensional base submanifold with boundary covering the
required smaller regions. Its saturated preimage defines
$M^{\rm slim}$.
\end{foundation}
Properness is a statement about compact subsets of the target,
not that an open coordinate inequality is compact. In the closed
case, a full preimage of a compact target set under the GLOBAL
continuous map is closed in compact $M$; localization puts it
inside the smooth submersion region. FC34 and LC85 identify the
fiber, while graphness of the base rules out another sheet.
The zero boundary condition uses only $Q_4\subset Q_3$ and is
therefore constant on these fibers. A connected zero boundary
is a single fiber, once that fiber identification is supplied.

\begin{foundation}[Proper edge disk bundles with faces; FC37]
\label{found:fibration-edge-bundles}
KL 14.8--14.11 truncate the edge base by marker $\ge.9R_i$
and ratio $<4\Delta$. The vertical restriction is the UNION
\[
 \{\eta_{E'}\le.35\Delta\}\ \cup\
 \mathcal E^{-1}\{x_\rho>0,\ x'_{E'}/x_\rho\le4\Delta\}.
\]
Intersect the full preimage of that base with this union to obtain
$U'_2$. After removing zero and slim interiors, the remaining
$M_2\cap U'_2$ is compact and is a union of disk fibers of the
proper restriction $\pi_2\mathcal E$. It is $M^{\rm edge}$.
\end{foundation}
Both vertical inequalities are CLOSED, so the disk boundary is
retained. The base ratio inequality is OPEN. Retain the low-height
branch even when the auxiliary output vanishes. KL 14.10 omits its proof; implementing it requires the
actual LC84 disk cores, FC34's common compact isolation, and
transversality on vertical and horizontal boundary faces.
For 14.11 the only potential escape through a base threshold is
excluded by the already removed slim region and KL 9.26's
replacement edge chart with a strict smaller threshold.
Thus each connected edge piece is $S^1\times D^2$ or
$I\times D^2$, in the orientable category, with its disk-bundle
structure and labeled horizontal and vertical faces. Proper
submersion for manifolds with boundary requires its face condition;
an unrestricted application of the boundaryless theorem is invalid.

\begin{foundation}[The remaining circle bundle; FC38]
\label{found:fibration-circle-bundles}
Truncate $W_1$ by marker $>.9R_i$ and ratio $<4$, using FULL
preimages under $\mathcal E$. KL 14.13 gives the proper circle
bundle on $U'_1$, containing the local $3.5$ regions and contained
in $U_1$. The complement of the interiors of zero, slim and edge
pieces is contained in $U'_1$, is saturated for this bundle, and
is a compact manifold with corners of depth at most two.
Call it $M^{2\text{-stratum}}$.
\end{foundation}
KL (14.12) prints marker $>.9$ without $R_i$; we use the
scale-consistent $>.9R_i$, as in (13.45) and the preceding stages.
This is another recorded project comparison, not an author erratum.
The marker comparison, actual LC83 circle fibers, compact
preimages, FC34 transport and saturation of the previously cut
faces are separate proof obligations. The normal derivatives of
two meeting face equations must be independent. KL 14.13's
omitted proof does not authorize omitting these from an implementation.
The induced circle fibration on the vertical boundary of an edge
disk bundle agrees with the one on $M^{2\text{-stratum}}$;
commuting projections together with the actual rank and connected
fiber identifications establish this agreement.

\subsection{Actual zero domains and saturated compact slim pieces}
\label{sec:fibration-actual-zero-slim-extraction}
Revision150 continues with GAF01's SAME closed-carrier construction.
Write $T_0$ for its common lower zero-radius parameter, so LPA05
and the three actual support-list lemmas give the ratio $T_0/20$.
Its existing bound $T_0\ge1600\cdot10^6\Delta$ in particular
implies $T_0\ge200$. At the ORIGINAL joint zero-packet choice
require radial value error $e_0<1/1000$ and difference-Lipschitz
error $\varepsilon_0<1/100$. These are permitted later minima in
LPA04--LPA05 and the actual producers; no selected radial function,
model, radius or already constructed map is subsequently replaced.
Keep GAF01's $c_3<1/1000$.

\begin{lemma}[Actual zero-block isolation and a global fixed-radius error; ZSP01]
\label{lem:fibration-actual-zero-block-isolation}
For every $i\in I_0$, every intermediate map $f$ and every point $p$,
\[
 \rho(p)>200R_i/T_0\quad\Longrightarrow\quad
                             J_i f(p)=J_iF(p)=0,              \tag{ZI}
\]
where $J_i$ projects onto the WHOLE zero block.
Consequently, with $\delta_0=200c_3/T_0<1/1000$,
\[
 |J_i(f(p)-F(p))|<\delta_0R_i                              \tag{ZE}
\]
for all $p$. The same bound holds on the straight segment from
$F$ to $E$. These estimates do not require a global bound on
$\rho/R_i$ or a new pruning rule for zero blocks.
\end{lemma}
\begin{proof}
A positive original $i$ cutoff places $p$ in the selected zero
ball's closed ten-radius neighborhood, by its annular support
and the radial value estimate. LPA05 then gives
$\rho(p)\le20R_i/T_0$. Thus the original block vanishes in (ZI).

Induct through the actual stages. Outside the stage cutoff support
there is no change. On that support put $x=\pi_jF(p)$; it belongs
to the original core by GAF02. For EVERY selected smoothing center
$y$ whose closed $80br_y$ support meets $B(x,8br_x)$, the scale
calculation (AM) in CFS28 gives, for its actual model reference $a$,
\[
 R_a\ge(108/625)\rho(p)>\frac{864}{25}\frac{R_i}{T_0}
                                  >20R_i/T_0.                \tag{ZL}
\]
If the zero support met that reference's actual comparison ball
$D_a$, SGP01, EGP02 or TCP01 would instead give
$R_i/R_a\ge T_0/20$. This contradicts (ZL).
The zero block is therefore ABSENT from the actual support list.
The model used to choose the plane has that entire block identically
zero; its actual core preimage lies in $D_a$, so $J_i y=0$ as well.
Thus EVERY contributing center and plane satisfies CFS24's
whole-block zero hypothesis. CFS16 puts the perturbed input in
$B(x,r_x)$, and the projection has zero $i$ block. The blend of
this with the inductively zero input is zero. All zero blocks are
retained in every $Q_j$, so this proves (ZI) at all stages.

If (ZI)'s inequality fails, GAF02 gives
$|J_i(f-F)|<c_3\rho(p)\le200c_3R_i/T_0$; in the other case
this difference is zero. Convexity gives the segment assertion.
Unlike the nonzero-family argument in CFS27, this proof uses
absence from the ACTUAL support list, not deletion of a zero block
from an already chosen plane.
\end{proof}

\begin{theorem}[The actual compact zero domains; ZSP02]
\label{thm:fibration-actual-zero-domains}
For $i\in I_0$ define, exactly with FC35's core branch,
\[
 Z_i=B_g(p_i,.35R_i)\ \cup\
 E^{-1}\{v_i\ge.9R_i,\ u_i/v_i\le.4\}.                     \tag{ZD}
\]
These are pairwise disjoint compact smooth domains and
\[
 \overline B_g(p_i,.38R_i)\subset\operatorname{int}_M Z_i,
 \qquad Z_i\subset B_g(p_i,.42R_i).                         \tag{ZB}
\]
Each is smoothly ambient isotopic to the SAME original radial
sublevel $\{\eta_i\le.4\}$ and has its LPA05/LFR54 smooth type,
including the boundary. Its boundary has the GLOBAL description
\[
 \partial Z_i=E^{-1}\{v_i\ge.9R_i,\ u_i=.4v_i\}.             \tag{ZF}
\]
On a neighborhood of this entire boundary $v_i>.99R_i$ and
$u_i/v_i-.4$ is a smooth defining function with nonzero differential.
Every nonempty boundary is a connected $S^2$ or $T^2$.
Empty-boundary cases retain their original compact types.
\end{theorem}
\begin{proof}
Fix $i$, put $R=R_i$, $d=d_g(p_i,\cdot)/R$, $\eta=\eta_i$,
and $\zeta=\zeta_i$. In $R$ units ZSP01 gives errors less than
$\delta_0\le c_3<.001$ in both scalar components, globally.
If $v_i(E)\ge.9R$, then $\zeta>.9-\delta_0>.899$, so the
original radial coordinate is in its annular support.
If also $u_i(E)/v_i(E)\le.4$, then
\[
 \eta\le .4+\frac{1.4\delta_0}{.9-\delta_0}<.402,
 \qquad d<.403.                                             \tag{ZR}
\]
If the ratio equals $.4$, the same calculation in both directions
gives $|\eta-.4|<.002$ and hence $.397<d<.403$.
In particular these conclusions exclude EVERY remote component
of the globally defined sets, including regions where $\rho/R$
is large.

On $.32\le d\le.38$ one has $.319<\eta<.381$, so the original
annular cutoff is exactly one. Here the adjusted marker exceeds
$.999R$ and the ratio is less than $.382/.999<.4$ with strict
slack. This whole closed annulus is in the interior of the second
set in (ZD). Together with the core ball this proves the first
inclusion in (ZB); (ZR) proves the second. Moreover (ZD) equals
the union of the closed ball of radius $.34R$ and its globally
closed second set. It is compact. The selected open unit-radius
zero balls are disjoint by LPA05, so these compact domains are
pairwise disjoint.

On $\eta\in(.3,.5)$ the original block is $(R\eta,R)$.
Set $r=u_i(E)/v_i(E)$ there. The marker is positive and the
normalized original gradient obeys
$1-\varepsilon_0\le\|D\eta\|\le1+\varepsilon_0$ by LC29.
All differential norms in this paragraph use $R^{-2}g$.
The physical estimate $\|DE-DF\|<c_3$ becomes
\[
 \|D(u_i(E)/R)-D\eta\|<c_3,\qquad
                             \|D(v_i(E)/R)\|<c_3.
\]
The quotient rule, $|r|<.502$ and $v_i(E)/R>1-\delta_0$, give
\[
 |r-\eta|<2\delta_0,\qquad
 \|Dr-D\eta\|\le
 \frac{(1+.502)c_3+(1+\varepsilon_0)\delta_0}{1-\delta_0}
                                    <3c_3.                  \tag{ZC}
\]
Hence $r-.4$ has nonzero differential at its zero set.
The enclosing annuli show that (ZD) is locally its sublevel
there, with no marker-boundary face and no core-sphere face.
They also prove both inclusions of (ZF). The domain is therefore
smooth, and $v_i>.999R$ near its boundary after a possible
shrinking of the neighborhood.

For completeness use GLOBAL smooth functions for the isotopy.
Choose a smooth nondecreasing $\psi$ equal to the identity on
$[.31,.49]$, constant $.3$ below $.30$ and constant $.5$ above
$.50$, with $\psi(s)\le.4$ exactly when $s\le.4$.
Choose a smooth $\chi$ equal to one on $[.34,.46]$ with support
in $(.32,.48)$. Put
\[
 h_\tau=\psi(\eta)+\tau\chi(\eta)(r-\eta),
                   \qquad0\le\tau\le1,                     \tag{ZH}
\]
extending the second term by zero. Both terms are smooth on ALL
of $M$: their nonconstant supports are in LC67's original smooth
shell, with strict buffers, and the first is constant near the
possibly nonsmooth radial core. By (ZC), any level $h_\tau=.4$
lies where $|\eta-.4|<.002$. There $\psi(\eta)=\eta$ and
$\chi=1$, and $Dh_\tau$ is within $3c_3$ of $D\eta$.
Thus the entire family is transverse to $.4$.
Its full sublevel trace is a closed subset of compact
$M\times[0,1]$, not just a collection of compact endpoint fibers.
The enclosing annuli and (ZR) give
\[
 h_0^{-1}(-\infty,.4]=\{\eta\le.4\},\qquad
                              h_1^{-1}(-\infty,.4]=Z_i.
\]
Apply FC34's smooth ambient isotopy with target half-line.
LPA05 and LFR54 identify the actual original sublevel and boundary;
the four noncompact-core alternatives have one $S^2$ or $T^2$
boundary, and the compact alternatives have none. This proves the
asserted types without a new metric or a replacement radial function.
The argument also covers the empty-level compact case.
\end{proof}

\begin{proposition}[Zero faces are saturated and descend to the bases; ZSP03]
\label{prop:fibration-zero-faces-saturated}
Put $Z=\bigcup_iZ_i$ and $M_1=M\setminus\operatorname{int}_M Z$.
For every $j=1,2,3$, every WHOLE fiber of $\pi_jE$ meeting
$\partial Z_i$ lies in $\partial Z_i$.
For the proper circle/slim restrictions $f_j=\pi_jE:X_j\to B_j$
of GAF07, $M_1\cap X_j$ is saturated and equals $f_j^{-1}(C_j)$,
where $C_j$ is a closed relative subset of $B_j$ and a smooth
domain with boundary. The boundary equation descends from the
retained $Q_4$ block, with nonzero differential on $B_j$.

For $j=3$, if $\partial Z_i$ meets $X_3$, the ENTIRE connected
surface $\partial Z_i$ is one whole fiber in $X_3$.
Thus $\partial C_3$ is a finite set of base points, one for each
such zero boundary. These assertions include actual face equalities,
not only an inclusion of tangent kernels.
\end{proposition}
\begin{proof}
Every $Q_j$ retains $Q_4$, so equality of $\pi_jE$ preserves the
$u_i,v_i$ equations in (ZF). Its global equality proves the first
assertion, without restricting the second point to an original chart.

The fibers in GAF07 are connected. A fiber not meeting $\partial Z$
cannot have points on both sides, since the two open sides would
disconnect it. A fiber meeting the boundary is entirely there.
Thus each whole fiber is either in the interior of $Z$, in $M_1$,
or in their common boundary, proving saturation. The proper map
$f_j$ is closed, so the image $C_j=f_j(M_1\cap X_j)$ is closed
in $B_j$ and its full preimage is $M_1\cap X_j$.
Near a boundary point use the smooth base function
\[
 b_i(w)=u_i(w)/v_i(w)-.4.
\]
The marker has a positive buffer. Its pullback is the defining
function in ZSP02 and has nonzero differential; since $f_j$ is a
submersion, $db_i$ is nonzero. Submersion charts give exactly
$C_j=\{b_i\ge0\}$ locally. Zero domains are disjoint, so two
such zero faces never meet. This proves the asserted smooth domain
and its boundary equation.

For $j=3$, a fiber meeting $\partial Z_i$ is a compact embedded
connected two-manifold contained in that connected two-manifold.
The inclusion is a local diffeomorphism, hence has open image;
compactness makes it closed. It therefore equals $\partial Z_i$.
In particular no part of this boundary can leave $X_3$.
Finiteness of the selected zero family proves finiteness of
$\partial C_3$. Properness and the base coordinate also give the
full local product over the corresponding boundary half-interval.
\end{proof}

\begin{theorem}[A compact saturated slim piece with all horizontal faces; ZSP04]
\label{thm:fibration-actual-compact-slim-piece}
There is a compact smooth one-dimensional submanifold with boundary
$K_3\subset B_3$ such that
\[
 f_3\left(\bigcup_{i\in I_s}
             \{|\eta_i|\le3.5\cdot10^5\Delta\}\right)
                       \subset\operatorname{int}_{B_3}K_3,
 \qquad \partial K_3\cap\partial C_3=\varnothing.            \tag{SK}
\]
Set $D_3=K_3\cap C_3$ and
\[
 M^{\rm slim}=f_3^{-1}(D_3)=M_1\cap f_3^{-1}(K_3).
\]
Then $D_3$ is a compact one-manifold with boundary and
$M^{\rm slim}\to D_3$ is a proper onto smooth bundle with the
SAME whole $S^2$ or $T^2$ fibers. It is a compact smooth domain
in $M_1$ and a smooth manifold with boundary. In particular
\[
 \begin{split}
 \partial D_3&=(\partial K_3\cap\operatorname{int}C_3)
                   \ \sqcup\ (\operatorname{int}K_3\cap\partial C_3),\\
 \partial M^{\rm slim}&=f_3^{-1}(\partial D_3),\qquad
 M^{\rm slim}\cap\partial M_1
                  =f_3^{-1}(K_3\cap\partial C_3).
 \end{split}                                                \tag{SF}
\]
The last set is a union of ENTIRE zero-boundary components and
has a collar inside $M^{\rm slim}$ relative to $M_1$.
Every required smaller slim region surviving in $M_1$ is contained
in this piece. Empty families give the empty piece.
\end{theorem}
\begin{proof}
Each original closed $3.5\cdot10^5\Delta$ slab is compact by
LFR20's whole proper bundle. Their finite union is compact and
lies in $X_3$ by GAF07. Its image is thus compact in $B_3$.
Cover that image by finitely many open coordinate intervals with
compact closures in $B_3$. Choose slightly larger closed coordinate
intervals still inside $B_3$, whose endpoints avoid the finite set
$\partial C_3$, and whose interiors cover the compact image.
Their finite union is a compact one-dimensional submanifold with
boundary; any touching endpoints simply become interior points.
A whole circle component, if covered, is retained as such. This
constructs $K_3$ and (SK), using the actual embedded base.

ZSP03 makes $C_3$ a domain with finitely many boundary points.
The avoidance in (SK) implies that its intersection with $K_3$
is a compact one-manifold with precisely the boundary in (SF);
no isolated endpoint or double boundary constraint occurs.
Restrict GAF07's whole proper bundle to this base. Its product
charts restrict to intervals or half-intervals. Properness makes
the total space compact and gives the two face equalities (SF).
Near a point over $\partial C_3$, the $K_3$ restriction is absent
on a full neighborhood, so the whole inward collar of $M_1$ is
included in $M^{\rm slim}$. ZSP03 shows that these shared faces
are entire connected zero boundaries. This proves the collar and
coverage statements, including all their source points.

On each interval component of $D_3$, the inherited smooth bundle
transport trivializes the piece as $I\times S^2$ or $I\times T^2$,
with its actual endpoint fibers. On a circle component it gives
the actual surface mapping torus. The total space inherits the
orientation of $M$. No assertion identifying its monodromy with
an isometry is used, and the later recognition input stays FC41.
\end{proof}

\begin{proposition}[The remaining smooth carrier after zero and slim removal; ZSP05]
\label{prop:fibration-zero-slim-complement-faces}
Use the RELATIVE interior for the second removal:
\[
 M_2=M_1\setminus\operatorname{int}_{M_1}(M^{\rm slim}).       \tag{RC}
\]
Then $M_2$ is a compact smooth manifold with boundary, and
\[
 \begin{split}
 M&=Z\cup M^{\rm slim}\cup M_2,\\
 \partial M_2&=(\partial M_1\setminus M^{\rm slim})
          \ \sqcup\ (\partial M^{\rm slim}\setminus\partial M_1),\\
 M^{\rm slim}\cap M_2
       &=\partial M^{\rm slim}\setminus\partial M_1.
 \end{split}                                                \tag{RF}
\]
All displayed faces are actual embedded closed surfaces with collars.
The three domains have disjoint ambient interiors. Each zero/slim
shared face is retained by those two pieces; it is removed from $M_2$.
A closed connected zero or slim component, if present, is the entire
connected carrier $M$.
\end{proposition}
\begin{proof}
By (SF), $M^{\rm slim}$ meets $\partial M_1$ along whole connected
components and contains their relative inward collars. At those
components its relative interior includes the face itself, so (RC)
removes a full relative neighborhood, with no isolated surface left
in the complement. At a different boundary component of
$M^{\rm slim}$, the boundary lies inside $\operatorname{int}_M M_1$
and is a full regular fiber. Its product collar cuts off one side
and leaves the other as the smooth boundary of $M_2$.
At the remaining components of $\partial M_1$, the slim piece is
absent on a neighborhood. These disjoint local cases prove that
$M_2$ is a smooth domain and give every equality in (RF).
Its compactness follows because (RC) removes a relatively open set
from compact $M_1$. The union and disjoint-interior statements follow
in the same local charts. An embedded compact three-dimensional
component without boundary is open and closed in connected $M$,
which proves the last assertion.

This specifies the interior convention left implicit in the source's
successive-removal notation. Removing only the AMBIENT interior of
$M^{\rm slim}$ from $M_1$ would retain a shared zero/slim surface as
an isolated lower-dimensional remnant. That operation is not used.
The later edge removal must likewise state its relative face convention.
\end{proof}

\paragraph{Boundary after revision150.}
ZSP01--ZSP02 close the actual zero-domain extraction, including
fixed-radius global isolation and a smooth whole-trace isotopy.
ZSP03--ZSP05 supply the saturated zero/slim faces, compact proper
slim restrictions and the precise smooth remaining carrier.
Circle saturation at zero faces is also available from ZSP03.
The height-bounded proper edge disks, their vertical/horizontal
transversality, circle remainder, final compatible decomposition,
full parameter register and nearly cuspidal boundary assembly remain
required work. This is an intermediate written development under the
retained foundations, not a final mathematical audit or Lean completion.


\subsection{The actual edge disks and their two kinds of faces}
\label{sec:fibration-actual-edge-disk-restriction}
Revision151 uses the SAME $F,E,P,t=P/\rho$ and original edge
coordinates of GAF and ZSP. Put $s=E_\rho$ and $A=u_{E'}(E)$.
The scale and weak-edge blocks are unchanged after stage one.
This subsection constructs the proper edge restriction before
proving compactness at the remaining open base endpoints.
All geometric and smooth foundations retain their stated bindings.

\begin{lemma}[The actual adjusted scale varies slowly; EDP01]
\label{lem:fibration-actual-adjusted-scale-derivative}
There is a finite constant $C_\rho\ge100$, fixed after the first
smoothing parameters but BEFORE $\Delta$ and $\Lambda$, such that
\[
 |s-\rho|\le C_\rho\Lambda\rho,\qquad
                           \|Ds\|_g\le C_\rho\Lambda.        \tag{SD}
\]
In the permitted order we may require
\[
 c_3<\min\{10^{-5},[10^5(P_0+1)]^{-1}\},\qquad
                  \vartheta:=C_\rho\Delta\Lambda<10^{-6},   \tag{SE}
\]
where $P_0$ denotes GAF01's early profile bound. In particular
$s>0$ on all of $M$. The additional smallness of $c_3$ is an
EARLY choice; the additional bound on $\Lambda$ is a LATER choice.
No condition $\Delta c_3\ll1$ is imposed.
\end{lemma}
\begin{proof}
In EVERY actual TCP05 model the scale component is the constant
one. CFS27 retains it. Thus every first-cloud plane has zero scale
component, and its normal projector is the identity on that scalar
coordinate. CFS12's spectral projector has the same property, so
on its zero set $W_1^0$ one has
\[
 w_\rho=\mu_\rho(w)=\sum_y w_y(w)\rho(q_y),                 \tag{SM}
\]
where $q_y$ is any actual original preimage of the selected center.
At stage one the scale coordinate descends exactly, so its value
is independent of that choice.

Write $\chi=\psi_1\circ F$. At any point of its closed support,
CFS31 supplies an original circle index $i$ with
$|\eta_i(p)|\le6.5<7$. Set $x=F(p)$ and $r_x=\Sigma_1\rho(p)$.
GAF04 proves that EVERY center contributing on $B(x,8br_x)$ has
full $i$ marker. Its original preimage has $|\eta_i|<7.2$ and
therefore lies in $B(p_i,10R_i)$ by TCP01 and the original circle
support enclosure. The point $p$ lies in the same ball. Consequently
\[
 |\rho(q_y)-R_i|\le10\Lambda R_i,\qquad
                          |\rho(p)-R_i|\le10\Lambda R_i.    \tag{SC}
\]
This applies to the ENTIRE contributing list, including centers whose
weight vanishes at one point.

Let $N_b$ be its bound in CFS11 and let $c_w$ be CFS12's bound
$\|Dw_y\|\le c_w/r_x$ on the whole reference ball.
Subtract the SAME constant $R_i$ in (SM); since $\sum_yDw_y=0$,
\[
 \|D\mu_\rho\|\le10N_bc_w\Lambda R_i/r_x
                     \le\frac{40N_bc_w}{3\Sigma_1}\Lambda. \tag{SW}
\]
The projection $P_1(x)$ lies in this reference ball, has derivative
norm less than two, and obeys (SM). With $z_\rho=(P_1F)_\rho$,
$\|DF\|\le L_0$ and (SC) therefore give
\[
 |z_\rho-\rho|\le20\Lambda R_i\le(80/3)\Lambda\rho,
 \qquad \|Dz_\rho\|\le
                   \frac{80N_bc_wL_0}{3\Sigma_1}\Lambda.
\]
The actual blend is $s=(1-\chi)\rho+\chi z_\rho$.
CFS31 gives $\|D\chi\|\le b_{\rm cut}L_0/\rho$.
The product rule bounds its derivative by
\[
 \left(1+\frac{80N_bc_wL_0}{3\Sigma_1}
                    +\frac{80}{3}b_{\rm cut}L_0\right)\Lambda.
\]
Outside the closed support $s=\rho$. The preceding estimates hold
on neighborhoods of support points, so also at their boundary.
One may take
\[
 C_\rho=100(L_0+1)
                (1+b_{\rm cut}+N_bc_w/\Sigma_1).
\]
All its ingredients are early parameters. Empty first families have
$s=\rho$ and may use $C_\rho=100$. Later adjustments retain $s$.
Choose $c_3$ as in (SE) before GAF01's reverse choices, and then
$\Lambda$ below the additional positive bound after $\Delta$.
This preserves every earlier producer requirement. The value estimate
and $C_\rho\Lambda<10^{-6}$ prove positivity.
\end{proof}

\begin{lemma}[The actual edge base and all height-qualified preimages; EDP02]
\label{lem:fibration-actual-edge-height-localization}
Define
\[
 \begin{split}
 B_2&=\bigcup_{i\in I_e}
       \{w\in W_2:v_i(w)>.9R_i,\ |u_i(w)|/v_i(w)<4\Delta\},\\
 T&=A/s,\qquad
 V=\{t\le.35\Delta\}\ \cup\ \{s>0,\ T\le4\Delta\},\\
 X_2&=(\pi_2E)^{-1}(B_2)\cap V.
 \end{split}                                                \tag{ED}
\]
The base $B_2$ is open in $W_2$; each indicated patch is exactly
$V_i\cap\{|u_i|/R_i<4\Delta\}$ and has $v_i=R_i$.
Replacing the marker $>.9R_i$ by $\ge.9R_i$ produces the SAME base.
For EVERY $p\in X_2$ and every witnessing base index $i$,
\[
 |\eta_i(p)|<4.01\Delta,\qquad t(p)<4.01\Delta,
 \qquad p\in U_2.                                         \tag{ELoc}
\]
Conversely all original closed smaller sets
$\{|\eta_i|\le3.5\Delta,\ t\le3.5\Delta\}$ lie in the
ambient interior of $X_2$.
On the full preimage of $B_2$, the low branch in (ED) implies
$T<4\Delta$, so $X_2=(\pi_2E)^{-1}(B_2)\cap\{T\le4\Delta\}$.
This equality is proved for the actual base; the original union
is retained in the definition.
\end{lemma}
\begin{proof}
For the strict marker condition repeat GAF07's base argument:
every $w\in W_2$ has an actual final preimage, (AM0) excludes a
zero original marker, and CFS26 bounds the scale there. Hence
$v_i(w)<(1+5c_3/4)R_i$ and the ratio condition implies
$|u_i(w)|<5.5\Delta R_i$. Thus $w\in V_i$, and GAF05 gives
the exact marker and coordinate patch. The latter is an open interval.

Here is the endpoint check for $v_i(w)=.9R_i$. The same upper bound
and vector bound hold with this nonstrict hypothesis. Pull $w$ back
through CGP08's later diffeomorphism to $w^0\in W_2^0$; the $i$
block is unchanged. CGP05's proof still supplies (MW): its first
proximity is STRICTLY less than $R_i/1000$, so the original marker
is still greater than $.899R_i$, and every subsequent inequality is
unchanged. Put $a=u_i(w^0)/R_i$. Compare with the actual inner
section at $a$ exactly as in CGP07's proof of (SN). This places
$w^0$ in the same $r_x/4$ ball about that original core point.
GAF04 and GAF03 force $v_i(w^0)=R_i$, a contradiction to $.9R_i$.
This proves the equality of the two marker conventions; it did not
assume the point already belonged to the strict patch.

Now take ANY whole preimage of such a base point. GAF06 gives
$|\eta_i|<4.01\Delta$ and a positive original cutoff, with
$3R_i/4\le\rho\le5R_i/4$. Since the final marker is exactly $R_i$,
\[
 \zeta_i>1-d,\qquad d=5c_3/4.
\]
The tangential cutoff equals one, so $\zeta_i=g(t/\Delta)$ and
$t<9\Delta$. Put $Z_0=\chi_{1/2,1}(\sum_{j\in I_e}\zeta_j)$.
This is not yet asserted to equal one. Its actual profile derivative
bound gives $Z_0\ge1-2P_0d$, since the sum is at least $1-d$.
Thus
\[
 g(t/\Delta)Z_0>1-\beta,\qquad
                   \beta=(1+2P_0)d<1/1000.                 \tag{EZ}
\]
If $t\ge.3\Delta$, the original weak-edge vector is
$\rho t g(t/\Delta)Z_0$. From $T\le4\Delta$, (AE) and
$s/\rho<1+c_3$ one obtains
\[
 t(1-\beta)<4\Delta(1+c_3)+c_3
                              <4.01\Delta(1-\beta).
\]
If instead $t<.3\Delta$ or the first branch of $V$ applies,
height localization is immediate. This proves (ELoc) for ALL
preimages. On the low branch the height cutoff is one; together
with the tangential cutoff this gives $Z_0=1$ exactly. The original
weak-edge vector is between zero and $\rho t$, so
$T\le(.35\Delta+c_3)/(1-c_3)<4\Delta$. This proves the stated
set equality without losing the low branch in the construction.

For an original smaller-set point, the index's marker is one and
both coordinates are below the threshold-$5$ domain. GAF02 and
GAF05 put its image in $W_2$ with marker $R_i$. Its normalized vector
has size at most $3.5\Delta+5c_3/4<4\Delta$.
Here the original weak-edge vector is at most $3.5\Delta\rho$, so
$T\le(3.5\Delta+c_3)/(1-c_3)<4\Delta$ with strict slack.
All defining base and vertical inequalities hold on a neighborhood,
which proves the interior inclusion.
\end{proof}

\begin{lemma}[A smooth original buffer and the full height quotient estimate; EDP03]
\label{lem:fibration-edge-original-buffer-and-height}
For $i\in I_e$ let $U_i=B(p_i,100\Delta R_i)$ be its original
coordinate domain, and put
\[
 Y_i=\{p\in U_i:|\eta_i(p)|<5\Delta,\ t(p)<5\Delta\},
 \quad H_0=\Delta\psi(t/\Delta),\quad g_i=u_i(E)/R_i,
\]
using LFR27's SAME constant-core profile $\psi$.
The functions $\eta_i,H_0,g_i,T$ are smooth on $Y_i$.
There is a compact subset $Q_i\Subset Y_i$ containing
\[
 \{p\in U_i:|\eta_i(p)|\le4.1\Delta,\ t(p)\le4.1\Delta\}
                                                               \tag{EBuf}
\]
in its interior. All of $Y_i$ lies in $B(p_i,8\Delta R_i)$.
With differential norms in $R_i^{-2}g$, throughout $Y_i$,
\[
 |g_i-\eta_i|<5c_3/4,\qquad
 \|Dg_i-D\eta_i\|<c_3,\qquad \|D\eta_i\|>.99.             \tag{ETan}
\]
On its band $.3\Delta\le t<5\Delta$,
\[
 |T-t|,\ \|DT-Dt\|<h_*:=50(c_3+\vartheta)<1/1000.          \tag{EH}
\]
For $t<.3\Delta$, $T<.31\Delta$. On the collar
$3.9\Delta<t<4.1\Delta$, the original pair $(\eta_i,t)$ has
least singular value greater than $.9$. All these margins can be
imposed in the existing parameter order.
\end{lemma}
\begin{proof}
Work in $R_i$ units and put $q=\rho/R_i$ and
$\lambda=100\Delta\Lambda$. The original scale estimate gives
$|q-1|\le\lambda$ on $U_i$. Keep the original LFR28 bounds
$\lambda,\mu,\tau\le10^{-8}$.
For a point with $|\eta_i|\le a\Delta$, $t\le a\Delta$, $a\le5$,
the original raw tangential coordinate $u$ and physical distance to
$\overline{E'}$ satisfy
\[
 |u|\le(a+\mu)\Delta,\qquad
 d_{\overline{E'}}/R_i\le\{a(1+\lambda)+\mu\}\Delta.
\]
A nearest border point initially lies inside radius $106\Delta$,
so the ORIGINAL coarse-border test applies before any smaller-ball
conclusion. As in LFR28.6 it yields
\[
 d(p_i,p)/R_i\le
 \Delta\sqrt{(a+\mu)^2+
           \{a(1+\lambda)+\mu+2\tau\}^2}+\tau\Delta.
                                                               \tag{EDist}
\]
For $a=5$ this is less than $8\Delta$; for $a=4.2$ it is less
than $6\Delta$. Consequently one may take
\[
 Q_i=\{p\in\overline B(p_i,7\Delta R_i):
                  |\eta_i(p)|\le4.2\Delta,\ t(p)\le4.2\Delta\}.
\]
It is compact, contained in $Y_i$, and (EBuf) lies in its interior
by the strict radius and coordinate margins. This uses only the
original smooth source domain and the same coarse-border chart.
It does not assert a disk bundle on a larger interval or change a
finite model embedding.

The enclosure also puts $Y_i$ in LFR27's smooth constant-core
domain, proving smoothness of $H_0$ even where $t$ itself is not
known smooth. The final functions are smooth because $E$ is smooth
and $s>0$. Original cutoffs on $Y_i$ are full, so (AE) proves the
first two inequalities of (ETan). For the last, at any such point
lift $(u(p)+200\Delta,v(p))$ through its ORIGINAL splitting.
The lift lies in the original radius-$1000\Delta$ tested domain;
its distance is $200\Delta+O(\beta_E)>100\Delta$. The actual
LFR19 all-direction estimate gives
$D\eta_i(v_+)\ge1-\sigma-O(\beta_E/\Delta)>.99$.
These ordered smallness bounds are already permitted. No derivative
of a coarse map is used.

On $.3\Delta\le t<5\Delta$ the original tangential and height
cutoffs are one, $Z_0=1$, and the weak-edge vector is EXACTLY $P$.
Put $S=s/R_i$, $B=A/R_i$ and $p_*=P/R_i=qt$. Then
\[
 \begin{gathered}
 .99<q<1.01,\quad S=q(1+\epsilon_s),\quad
 |\epsilon_s|\le C_\rho\Lambda,\quad
 \|DS\|\le C_\rho\Lambda,\quad\|Dq\|\le\Lambda,\\
 |B-p_*|<c_3q,\quad \|DB-Dp_*\|<c_3,\quad\|Dp_*\|<2.
 \end{gathered}
\]
These are normalized versions of (SD), (AE) and the SAME physical
smoothing's Lipschitz bound. The exact identities
\begin{align*}
 T-t&=\frac{B-p_*+t(q-S)}{S},\\
 D(T-t)&=\frac{DB-Dp_*}{S}
       +\left(\frac1S-\frac1q\right)Dp_*
       -\frac{T}{S}DS+\frac{t}{q}Dq.
\end{align*}
give respectively bounds $2c_3+10\vartheta$ and
$2c_3+20\vartheta$, using $t<5\Delta$, $S>.98$ and $\Delta\ge1$.
Both are less than $h_*$, proving (EH), with the FULL quotient rule.
For $t<.3\Delta$, the original auxiliary vector is between zero
and $\rho t$; the same value bound gives $T<.31\Delta$.

Finally use LFR35's actual positive-anchor Gram estimate within
$1/100$ and LFR36's comparison of the SAME pair with those directions
within $1/1000$. These are choices before $\Delta$, followed by
the already allowed later local accuracies and curvature tail.
They give least singular value greater than $.95$ in the point's
own scale on the stated collar. Changing to $R_i$ units divides by
$q<1.01$, leaving the asserted $.9$. The full band is within
LFR36.1 by the proved enclosure. Thus the needed rank is quantitative
and refers to the original pair, not newly chosen reference functions.
\end{proof}

\begin{theorem}[The whole proper edge disk bundle; EDP04]
\label{thm:fibration-actual-whole-edge-disk-bundle}
The ACTUAL restriction
\[
 f_2=\pi_2E:X_2\longrightarrow B_2
\]
is an onto proper smooth bundle with connected fiber the closed
smooth disk. The boundary of its total space is exactly
\[
 \partial X_2=X_2\cap\{T=4\Delta\},                       \tag{EV}
\]
and its boundary restriction is an onto submersion with circle
fibers. Here the boundary symbol denotes the manifold boundary;
excluded endpoints of the open base are treated separately. Each WHOLE disk fiber is smoothly ambient isotopic inside
some $Y_i$ to a fiber of the SAME original LFR28 disk packet.
The third adjustment preserves the whole stage-two disk fibers on
this SAME vertical restriction, without merging sheets.
\end{theorem}
\begin{proof}
EDP02 puts every point of $X_2$ in an original threshold-$5$ domain.
GAF02 therefore supplies its submersion to $W_2$. On a witnessing
patch $g_i$ is its actual coordinate by GAF05.
For the fiber identification fix $w\in B_2$ and such an $i$, and
write $a=u_i(w)/R_i$, $|a|<4\Delta$.
On the ORIGINAL open smooth source $Y_i$ consider
\[
 \mathcal H_\tau=
 ((1-\tau)\eta_i+\tau g_i, (1-\tau)H_0+\tau T),
 \qquad X_a=\{a\}\times(-\infty,4\Delta].                 \tag{EI}
\]
Every inverse image of $X_a$ has
$|\eta_i|<4\Delta+5c_3/4<4.01\Delta$.
If $t\le2\Delta$, its height is already below $4.01\Delta$.
If $t>2\Delta$, then $H_0=t$ and (EH) gives
$t\le4\Delta+h_*<4.01\Delta$.
Thus the ENTIRE time-dependent inverse image lies in the interior
of the SAME compact $Q_i$ from EDP03. It is closed in
$Q_i\times[0,1]$, and hence compact.

The first differential is nonzero everywhere by (ETan).
On the inverse image of $\partial X_a$, one cannot have
$t\le2\Delta$: both $H_0$ and $T$ are then less than
$2\Delta+h_*$. Thus $|t-4\Delta|<h_*$ there. On that collar
$H_0=t$; the differential pair in (EI) is within $c_3+h_*<.002$
of $D(\eta_i,t)$, whose least singular value exceeds $.9$.
This proves transversality to BOTH $X_a$ and $\partial X_a$
throughout the family, with a uniform positive margin.
Apply FC34 to this smooth source family and its whole compact trace.
The initial fiber is exactly the original
$\{\eta_i=a,t\le4\Delta\}$ in its radius-$100\Delta$ domain,
by LFR27. It is LFR28's closed smooth disk.

At the final time the fiber is precisely $f_2^{-1}(w)\cap X_2$.
One inclusion follows from (ELoc) and the equality of vertical sets
in EDP02. Conversely a point of $Y_i$ with $g_i=a,T\le4\Delta$
has final image in $W_2$ and marker $R_i$, by GAF02 and GAF05.
Its base coordinate $a$ puts it in $V_i$, where coordinate uniqueness
forces its image to be $w$. It then belongs to $X_2$.
This proves equality of the WHOLE fibers, including all components,
and gives their smooth type and onto-ness.

The same rank argument at the final time shows that $(g_i,T)$ has
rank two along $T=4\Delta$. EDP02 makes the low branch strictly
interior there. Hence $X_2$ is a smooth manifold with boundary (EV),
and $f_2$ restricts to a submersion on that boundary as well.
For a compact $K\subset B_2$, its whole inverse image under the
GLOBAL continuous $\pi_2E$ is closed in compact $M$; intersecting
with the globally closed set $V$ remains compact and equals the
inverse image under this restriction. This proves properness.

To obtain smooth bundle charts with the boundary retained, on a base
coordinate interval solve $Dg_i(X)=1$ and, near (EV), also $DT(X)=0$.
The rank-two margin permits both equations. In the interior solve only
the first, and combine the solutions by a smooth partition of unity.
The resulting field is tangent to the vertical boundary. Properness
traps its trajectories over each compact base subinterval, so forward
and reverse flows give smooth product trivializations, including the
closed disk boundary. Their boundary fibers are the connected $S^1$
boundaries of the just identified disks. This checks the boundary
condition rather than applying a boundaryless theorem to $X_2$.

Finally $s,A$, and therefore $V$, are unchanged by the third stage.
Let $f_2^0=\pi_2\mathcal E^2$ and let $\Theta_2:V_2^0\to W_2$
be CGP08's later diffeomorphism. If $f_2^0(p)=\Theta_2^{-1}(w)$
and $p\in V$, the actual ambient postcomposition gives $f_2(p)=w$.
Conversely $f_2(p)=w$, $p\in V$ implies by (ELoc) that $p\in U_2$,
where $f_2^0(p)\in V_2^0$ and injectivity of $\Theta_2$ gives the
reverse equality. These are both inclusions of WHOLE restricted
fibers, not only fiber inclusion on an assumed chart.
\end{proof}

\begin{proposition}[Saturated edge restriction and transverse horizontal faces; EDP05]
\label{prop:fibration-actual-edge-horizontal-faces}
For ZSP05's actual smooth $M_2$, the set
$M_2\cap X_2=f_2^{-1}(C_2)$ is saturated, where $C_2$ is a
closed relative subset of $B_2$ and a smooth one-dimensional domain
with boundary. This restriction is a proper disk bundle over $C_2$
and a smooth manifold with corners of depth at most two. Its faces
are exactly
\[
 H=\partial M_2\cap X_2=f_2^{-1}(\partial C_2),\qquad
                    V_e=M_2\cap\partial X_2.                \tag{EF}
\]
At $H\cap V_e$ their defining differentials are independent.
Every horizontal disk is contained in an actual zero or slim boundary.
Compactness at the open endpoints of $B_2$ is not asserted here.
\end{proposition}
\begin{proof}
Each boundary component of $M_2$ is either an entire zero boundary
or a new entire slim fiber, by (RF). A disk fiber meeting a zero
boundary lies there by ZSP03's GLOBAL $Q_4$ equation.
If it meets a slim boundary, its constant $\pi_2E$ makes $\pi_3E$
constant; GAF07 and ZSP04 identify the WHOLE latter fiber with that
slim boundary. Thus the entire disk lies in the same component of
$\partial M_2$. Connectedness of the disk then implies saturation
of $M_2\cap X_2$, just as in ZSP03. Properness of $f_2$ implies
that its image $C_2$ is closed in $B_2$.

Near a horizontal face, an ambient defining function for $M_2$
descends to a smooth function $b$ on the edge base. For a zero face
it is the retained ratio in (ZF). For a slim face use a local
coordinate defining its chosen endpoint in $D_3$, composed with
$\pi_{2,3}$ on the base. This latter composition has its required
domain: the whole compact edge disk lies in the open set $X_3$;
properness of $f_2$ gives a base neighborhood whose whole disks
stay in $X_3$. Their onto-ness puts the projected base in $B_3$.
The original submersion neighborhoods make these compositions
smooth and give the defining-function identity there.

The ambient defining differential for $M_2$ is nonzero. Since it
is $D(b\circ f_2)$, one has $db\ne0$ on the one-dimensional
base, and $C_2=\{b\ge0\}$ locally. On a coordinate patch,
$db$ is a nonzero multiple of $dg_i$. EDP04 gives independence
of $dg_i$ and $dT$ at the vertical boundary. Therefore $d(b\circ f_2)$
and $dT$ are independent at their intersection. Their actual
sublevel charts identify the intersection with a quadrant times
an open line. At points on just one face the corresponding
half-space chart applies. This proves (EF), depth at most two,
and the claimed face transversality. Restricting the proper bundle
to the closed base domain supplies its disk charts over intervals
or half-intervals. These statements do not assume that $C_2$ is compact.
\end{proof}

\begin{proposition}[The actual vertical circle fibers agree; EDP06]
\label{prop:fibration-actual-vertical-circle-agreement}
The whole vertical boundary $\partial X_2$ lies in GAF07's circle
region $X_1$. Its disk-boundary circles are EXACTLY whole fibers of
$E:X_1\to B_1$. The same holds after restricting to $V_e$ in (EF).
Near a meeting of vertical and horizontal edge faces, the circle
base has local coordinates $(g_i,T)$; the two faces are a coordinate
level and a transverse boundary curve in this base. In particular
both face equations are constant on the actual circle fibers.
\end{proposition}
\begin{proof}
At a point of (EV), (ELoc) and (EH) give
$|\eta_i|<4.01\Delta$ and $|t-4\Delta|<h_*$ for a witnessing
edge index. LFR38 applies to the SAME original pair on this full
band and makes the point two-stratum, using the retained no-three
result. LPA06's original circle cover supplies a circle chart with
$|\eta_a|<2$. GAF07 therefore puts the point in $X_1$.

A whole circle fiber through this point has the same value of the
GLOBAL $E$. It therefore has the same edge base point and the same
$T=4\Delta$; all its points belong to $X_2$ and, by EDP04, to the
boundary of that ONE whole disk fiber. This inclusion of a compact
connected smooth circle into the disk's connected smooth boundary
circle is open by its equal-dimensional tangent inclusion and closed
by compactness. It is the entire boundary circle.

The same global equations show that every whole circle fiber meeting
a component of $\partial M_2$ lies in it: use (ZF) for zero components
and the whole $\pi_3E$ fiber for new slim components. Connectedness
therefore makes $M_2\cap X_1$ saturated. Thus the preceding equality
of circles survives exactly on $V_e$.

Finally $(g_i,T)$ is a smooth function of $E$ on this neighborhood.
Its source differential has rank two by EDP04, and $E$ is a
submersion to the two-dimensional $B_1$. Its descended differential
on $B_1$ is therefore invertible; the inverse function theorem gives
these actual local base coordinates. The vertical face is $T=4\Delta$.
EDP05's horizontal equation is $b(g_i)=0$, with $b'\ne0$.
This proves their independence and equality of the induced circle
fibrations, not merely compatibility of coordinate projections.
\end{proof}

\paragraph{Boundary after revision151.}
EDP01--EDP04 supply the actual slow scale coordinate, whole height
localization and proper closed-disk fibers, with compact transverse
isotopies on the original smooth source. EDP05--EDP06 prove the
horizontal/vertical face independence and equality with the actual
circle fibers. These fill KL14.10's omitted local proof in the
present repaired construction. Compactness of $M_2\cap X_2$ at
its open base threshold, the remaining circle-domain coverage,
relative final removal and global compatible decomposition still
require proof. The exhaustive parameter register, final graph
presentation and SEPARATE nearly cuspidal boundary route remain
part of the active goal. No final full-goal audit or Lean completion
is claimed.


\subsection{Compact edge pieces and the actual remaining circle region}
\label{sec:fibration-actual-final-decomposition}
This subsection retains the closed carrier and the SAME assignment,
map, source coordinates and relative interiors of GAF, ZSP and EDP.
It supplies the geometric decomposition for that assignment. The
exhaustive parameter register and the separate boundary carrier still
have to be treated. All inherited recognition inputs remain explicit.

\begin{lemma}[Replacement edge charts after the earlier removals; FDC01]
\label{lem:fibration-actual-replacement-edge-chart}
Suppose $q\in M_2$ and an original selected edge index $i$ satisfies
\[
 q\in U_i,\qquad |\eta_i(q)|\le4.01\Delta,
                         \qquad t(q)\le4.01\Delta.
\]
Then $p_i$ is a nonslim one-stratum point for the ACTUAL later
$\beta_1$ predicate. There is a selected edge index $j$ with
\[
 |\eta_j(q)|<2\Delta,\qquad \zeta_j(q)=1,\qquad
 v_j(\pi_2E(q))=R_j,\qquad
                |u_j(\pi_2E(q))|/R_j<3\Delta.             \tag{Repl}
\]
In particular $\pi_2E(q)\in B_2$. The hypotheses need not include
this last assertion or strict membership in the original $i$ base.
\end{lemma}
\begin{proof}
EDP03's original enclosure, with $a=4.01$, gives
$d(q,p_i)<6\Delta R_i<10\Delta R_i$.
First exclude the zero stratum at $p_i$. If $p_i$ were zero-stratum,
LPA05's actual tenth-radius cover would put it in $B(z,R_0/10)$
for a selected zero center. The same shell scale comparison gives
$\rho(p_i)\le20R_0/T_0$. Therefore
\[
 d(q,z)<(1/10+200\Delta/T_0)R_0<.38R_0.
\]
ZSP02 puts this ball in $\operatorname{int}Z$, contradicting
$q\in M_2\subset M\setminus\operatorname{int}Z$.
EGP01 excludes a two-splitting at $p_i$, and LC18 excludes the
three-stratum. The exhaustive LC16 stratification therefore makes
$p_i$ a one-stratum point. This step is necessary: its strong-edge
quality by itself need not be the later $\beta_1$ quality.

If $p_i$ were slim, LPA06's actual slim cover would put it in
$B(p_k,\Delta R_k)$ for a selected slim center. Slow variation
then gives $d(q,p_k)<12\Delta R_k$. Its original value estimate
puts $|\eta_k(q)|<13\Delta<3.5\cdot10^5\Delta$.
GAF07 and (SK) put $q$ over $\operatorname{int}K_3$.
Since $q\in M_1$, the restricted bundle charts show that it lies
in $\operatorname{int}_{M_1}M^{\rm slim}$, even on a retained
old boundary face. This contradicts the definition of $M_2$.
Thus LFR44's nonslim hypothesis holds at $p_i$.

Let $u_i^0$ be the original raw tangential function in $R_i$ units.
LFR32's ACTUAL border coverage supplies
$q'\in E'$ with
\[
 |Q_i(q')-(u_i^0(q),0)|<\tau\Delta,
 \qquad d(q',p_i)<4.1\Delta R_i,
 \qquad |\eta_i(q')-\eta_i(q)|<.01\Delta.                 \tag{WB}
\]
Here the radius bound follows from pointed distortion, not from
an assumed continuity of $Q_i$. The coordinate bound uses the
SAME smoothing's value error $\mu\Delta$ at both points, so its
error is at most $(\tau+2\mu)\Delta$. All these estimates use
LFR32's original radius-$200\Delta$ domain. The corresponding
height bound and distortion also give
$d(q,q')<4.2\Delta R_i$.

Apply LFR44 item2 to $p_i,q'$. It gives $p\in E$ with
$d(p,q')<\rho(p)$. The SAME finite maximal family, by LFR44.1,
has a center $p_j$ with
\[
 d(p,p_j)<.7\Delta R_j,\qquad
 d(q',p_j)<(.7\Delta+1.01)R_j<.72\Delta R_j.
\]
The last inequality uses the already available $\Delta>100$.
Its original coordinate estimate gives $|\eta_j(q')|<.8\Delta$.
Moreover $q'\in E'$, so the SAME distance smoothing gives tiny
$t(q')$; both original $j$ cutoffs are full. Consequently
$j\in J_e(i)$, since $q'\in D_i$. EGP02--EGP04 apply to this
actual selected index and to both $q,q'\in D_i$. With their
permitted $\theta<1/100$, they give
\[
 |\eta_j(q)-\eta_j(q')|
 \le\frac{|\eta_i(q)-\eta_i(q')|+2\theta}{.99}
 <.04\Delta.
\]
Thus $|\eta_j(q)|<2\Delta$. The distance estimate just obtained
also puts $q$ in its original smooth $j$ domain. The common height
is below $4.01\Delta$, so $\zeta_j(q)=1$. The original $i$
coordinates put $q\in U_2$. GAF02 and GAF05 therefore put its
final image in $W_2$ and give its exact full $j$ marker. Finally
(AE) bounds the normalized vector by
$|\eta_j(q)|+5c_3/4<3\Delta$. This proves (Repl) without any
open-base assumption at the original index.
\end{proof}

\begin{theorem}[The actual compact edge piece; FDC02]
\label{thm:fibration-actual-compact-edge-piece}
The set
\[
 M^{\rm edge}=M_2\cap X_2=f_2^{-1}(C_2)
\]
is compact. Its base $C_2$ is a compact smooth one-manifold with
boundary, hence a finite disjoint union of intervals and circles.
The restriction is the actual proper whole disk bundle of EDP04--05,
with the horizontal and vertical faces in (EF). In particular no
extra face occurs at an excluded endpoint of the open base $B_2$.
\end{theorem}
\begin{proof}
Let $q_n\in M_2\cap X_2$ converge in compact $M$ to $q$.
Then $q\in M_2$ and $q\in V$, because $M_2,V$ are closed;
$V$ is globally closed since $s>0$ on compact $M$.
After a subsequence, one selected index $i$ witnesses the base
condition for every $q_n$. EDP02 gives
\[
 v_i(\pi_2E(q_n))=R_i,\qquad
                |u_i(\pi_2E(q_n))|<4\Delta R_i.
\]
At the limit these become equality of the marker and a weak vector
inequality. The proof of GAF06 applies to these global coordinate
inequalities with $\le4\Delta$ in place of $<4\Delta$:
(AM0) still forces a positive original cutoff, and the strict error
margin still gives $|\eta_i(q)|<4.01\Delta$ in the original domain.
The height calculation (EZ), using $q\in V$, likewise gives
$t(q)<4.01\Delta$. Neither calculation assumes that the limit
already belongs to $B_2$ or $W_2$.
FDC01 now gives a replacement index with STRICT base inequalities.
Thus $\pi_2E(q)\in B_2$ and $q\in X_2$.
This proves closedness in compact $M$, hence compactness.
Its continuous image $C_2$ is compact; EDP05 already identifies its
smooth domain structure and every boundary point. A compact
one-manifold has finitely many interval/circle components, proving
the remaining statements. Empty edge families give the empty piece.
\end{proof}

\begin{theorem}[The actual relative remainder is a compact circle bundle; FDC03]
\label{thm:fibration-actual-circle-remainder}
Use the precise relative removal
\[
 M_3=M_2\setminus\operatorname{int}_{M_2}M^{\rm edge}.
                                                               \tag{Last}
\]
Then $M_3\subset X_1$ and
\[
 M^{2\text{-stratum}}:=M_3=E^{-1}(C_1)\cap X_1,
                      \qquad C_1=E(M_3)\subset B_1
\]
is an onto proper smooth circle bundle over a compact smooth surface
with corners of depth at most two. All fibers are the SAME whole
circles of GAF07. Its total space has depth at most two and
\[
 M_2=M^{\rm edge}\cup M_3,\qquad
 M^{\rm edge}\cap M_3=V_e,\qquad
 M_3\cap\partial M_2=
 \partial M_2\setminus\operatorname{int}_{\partial M_2}H.   \tag{LastFaces}
\]
Every vertical edge boundary is a circle-bundle face of this
remainder, with exactly the same circles and attaching inclusion.
\end{theorem}
\begin{proof}
Compactness of $M_3$ follows from (Last). Near a horizontal edge
face, EDP05 gives $M_2=\{h\ge0\}$ and the edge piece is locally
$\{h\ge0,T\le4\Delta\}$, with independent differentials at their
intersection. Along the horizontal face away from the vertical
face the edge piece contains a full relative neighborhood in $M_2$.
Thus (Last) leaves no isolated horizontal surface. At the vertical
face it leaves $T\ge4\Delta$; at a meeting it leaves the actual
quadrant $h\ge0,T\ge4\Delta$. FDC02 excludes additional limiting
faces at open base endpoints. Elsewhere it leaves the original
smooth domain. These charts prove (LastFaces), disjoint relative
interiors and depth at most two.

Here is the required coverage, before using a circle fibration.
A zero-stratum point is in an original selected tenth-radius ball
and hence in $\operatorname{int}Z$; it cannot be in $M_3$.
A slim point in $M_1$ is in a selected slim covering ball and,
by (SK), in $\operatorname{int}_{M_1}M^{\rm slim}$; it cannot
be in $M_2$. Any nonslim one-stratum point lies, by LFR44.2 and
LPA06, in a selected edge ball with
$|\eta_i|<3.1\Delta$ and $t<3.1\Delta$.
EDP02 puts a full ambient neighborhood of that point in $X_2$.
If the point lies in $M_2$, it is therefore in
$\operatorname{int}_{M_2}M^{\rm edge}$ and is removed in (Last).
There is no three-stratum. Consequently every point of $M_3$ is
two-stratum. LPA06's ORIGINAL circle cover has $|\eta_a|<2$ at
such a point, and GAF07 puts it in $X_1$.

EDP06 proves that $M_2\cap X_1$ is saturated by whole connected
circle fibers. Membership in $X_2$ is constant on each such fiber:
its base point $\pi_2E$ and its height $T$ are both functions of $E$,
and EDP02 identifies $X_2$ exactly by those conditions.
For a point of $M^{\rm edge}$ the relative interior in (Last) is
characterized by $T<4\Delta$. This follows from the preceding
charts, including at the horizontal face, and from the absence of
other boundary faces in FDC02. Its truth is constant on the whole
circle fiber. Thus $M_3$ is saturated, proving the stated full
inverse-image equality.

It remains to check the base and boundary structure, rather than
infer it from saturation alone. Each zero/slim face of $M_2$
has a smooth ambient defining function descending to $B_1$:
use the global retained zero ratio or the local slim endpoint
coordinate as in EDP05. On a compact whole circle fiber in the
latter face, properness of GAF07 supplies a neighborhood over which
that coordinate is defined. Its nonzero ambient differential and
the submersion $E$ give a nonzero base differential.
On the vertical face use $T-4\Delta$. EDP06 gives actual base
coordinates $(g_i,T)$ there. At a meeting, the horizontal equation
is $b(g_i)$ with $b'\ne0$. Hence the two base differentials are
independent. These local equations describe $C_1$ as a smooth
domain, half-space or quadrant in $B_1$.
Its compactness follows from that of $M_3$. Restrict the actual
proper circle-bundle charts to these base domains; this gives its
smooth bundle, its face charts, and its whole connected fibers.
The vertical equality is precisely EDP06, including at the corners.
\end{proof}

\begin{theorem}[Actual closed-carrier decomposition for the common assignment; FDC04]
\label{thm:fibration-actual-closed-decomposition}
The domains $Z,M^{\rm slim},M^{\rm edge},M^{2\text{-stratum}}$
just constructed supply every geometric field of FC39 on the
ORIGINAL closed carrier. They are finite unions of compact embedded
pieces, cover the carrier, and have disjoint ambient interiors.
Their common faces, their bundle restrictions and their attaching
maps are the actual inclusions proved above. Given the retained
FC41 recognition and closure library, FC42 therefore produces a
finite KL graph presentation on this same carrier.
This conclusion is for the one actual assignment satisfying the
listed producer inequalities; the exhaustive simultaneous register
and the separate boundary construction are not claimed here.
\end{theorem}
\begin{proof}
Combine (RF) and (LastFaces) for the union and every complement
identity. ZSP02 and ZSP04 give compact smooth zero/slim pieces;
FDC02 and FDC03 give compact edge/circle pieces and their full
base descriptions. Compact manifold bases have finitely many
components, so all the lists are finite. The original zero topology
list is unchanged, including the corrected $D^3$ alternative.

Every horizontal edge disk is a whole $f_2$ fiber in one component
of $\partial M_2$, hence in an actual zero or slim boundary by
(RF) and EDP05. Distinct disks are disjoint as fibers. Their boundary
circles are the actual vertical-circle fibers by EDP06. A remaining
part of any such zero/slim boundary is a circle-bundle surface:
FDC03's descended horizontal defining function is a regular base
curve, with its endpoints exactly at the transverse vertical face.
There is no additional face or component between these sets.
The zero/slim shared surfaces are the ENTIRE common boundary
components in (SF), and do not reappear in the remainder. These
are precisely the hypotheses needed for FC40's Euler count and
FC42's finite assembly, not just abstract incidence assertions.

For the category of the final presentation, retain the usual
smooth corner-rounding operation in FC41's collared assembly.
It can be chosen compatible with the produced circle bundle:
the corners of the compact base $C_1$ are finitely many points;
in pairwise disjoint base rectangles round each quadrant boundary
by a smooth arc agreeing with its two axes outside a smaller
rectangle. Lift each replacement by the actual bundle product
chart over that rectangle. Thus the local change is the product
of the same circle with a base rounding. The complementary side
uses the same arc, so the attaching subsets agree. Away from
these rectangles the inclusions and fibration are unchanged.
After FC42 groups the adjacent ball/handle pieces, their rounded
boundary components and the rounded circle region meet along
whole smooth tori, with collars. Their identifications are induced
by these actual inclusions. The retained recognition/gluing
library supplies their finite graph witnesses and the sphere
reconstruction records. No assertion of essentiality of the cut
tori, or of a new PC theorem binding, follows from this construction.
\end{proof}

\paragraph{Source comparison and remaining work.}
FDC01 supplies the replacement needed by KL9.26 and14.11 while
explicitly proving the later one-stratum hypothesis at the selected
edge center. FDC02 fills the open-end compactness step. FDC03
expands KL14.13's omitted remainder proof with relative interiors,
whole fibers and actual independent base equations. FDC04 binds
these constructed pieces to KL14.1/15.1--15.3 and FC39--FC42.
The graph class is EXACTLY KL1.2; its torus and sphere closure
inputs retain FC41 and Chapter6's separate comparison gates.
The full parameter register, nearly cuspidal boundary augmentation
and final full-goal mathematical audit remain required work.


\begin{definition}[Compatible decomposition certificate; FC39]
\label{def:fibration-decomposition-certificate}
The static certificate on a compact orientable carrier consists
of the four domains just constructed, with disjoint interiors
and union the whole carrier; actual inclusions and corner charts;
all bundle projections, bases and fiber identifications; and
labeled common boundary identifications. Zero and slim domains
have smooth boundary. Edge and two-stratum pieces have corner
depth at most two. Each horizontal edge disk lies in a zero or
slim boundary, and each vertical edge circle bundle agrees with
the adjacent two-stratum circle bundle. The zero topology list
includes the author-corrected $D^3$ alternative.
\end{definition}
FC35--FC38, with their supplied source bindings, produce this
certificate by successive removal of interiors. Record complement
and common-face equalities as part of the output; the slogan
``compatible projections'' is not a substitute for those equalities.
This is the precise consumer of KL Proposition14.1 in the next
static topological assembly.

\section{Static graph assembly and the boundary route}
\label{sec:fibration-static-assembly}
\begin{lemma}[Number of disk faces; FC40]
\label{lem:fibration-euler-faces}
Let a connected closed surface $Y$ be either $S^2$ or $T^2$.
Suppose $Y=A\cup B$, where $A$ is a finite disjoint union of
$d$ closed disks, $B$ is a compact circle-bundle surface, and
$A\cap B=\partial A=\partial B$ is a disjoint union of fiber
circles. Then $d=2$ for $S^2$ and $d=0$ for $T^2$.
\end{lemma}
\begin{proof}
The inherited Euler-characteristic gluing formula and bundle
formula give $\chi(Y)=\chi(A)+\chi(B)-\chi(A\cap B)=d$.
Here a circle-bundle surface has Euler characteristic zero,
including the case of a base interval. Use $\chi(S^2)=2$ and
$\chi(T^2)=0$.
\end{proof}
These are KL 15.1--15.2's boundary counts. Disk interiors and all
common circles must be actual embedded faces of FC39. An abstract
incidence list without those embeddings is insufficient.

\begin{foundation}[Topological closure library for the static output; FC41]
\label{found:fibration-topological-closure}
Use the KL Definition1.2 graph class: a compact orientable
three-manifold, possibly with boundary, cut along FINITELY many
disjoint interior tori into circle-bundle pieces, with actual
collared gluing maps. Supply the following category-correct
closure and recognition results used in KL Section15:
\begin{enumerate}[leftmargin=*]
\item Gluing along entire torus boundary components and connected
sum preserve this graph class, with finite presentation witnesses.
\item Solid tori, twisted interval bundles over the Klein bottle,
$S^1\times S^2$, $\mathbb{RP}^3$, orientable flat and spherical
space forms, and $T^2$ bundles over $S^1$ have such presentations.
The last case is cut at a fiber. An orientable $S^2$ bundle over
$S^1$ is recognized as $S^1\times S^2$.
\item A punctured $\mathbb{RP}^3$ may be replaced by a ball with
an explicit connected-sum reconstruction; a sphere cylinder may
be cut and capped, with reconstruction either a connected sum
of components or addition of $S^1\times S^2$.
\item A finite cyclic assembly of three-balls with exactly two
disk faces and product disk handles is an orientable disk bundle
over $S^1$, hence a solid torus, with its actual attaching maps.
\end{enumerate}
\end{foundation}
These are supplied topological inputs, not consequences of the
word ``graph.'' Their exact inherited declarations and smooth
category must be audited against the completed PC release.
Chapter 6's existing Matveev comparison and relative graph-refinement
contracts retain the distinction between this raw graph class,
Seifert presentations and essential-torus outputs. No new edition
equivalence or essentiality is asserted here. Schoenflies remains
an inherited candidate; it is not labeled a new missing theorem.

\begin{proposition}[Assembly from the decomposition; FC42]
\label{prop:fibration-static-assembly}
Given FC39 and the topological library FC41, a closed connected
orientable carrier has a finite KL graph presentation.
\end{proposition}
\begin{proof}
A closed zero component already is the whole connected carrier;
use its recognition input. A slim component over a circle is
likewise closed and exhausts the carrier; use the sphere or torus
bundle case. Otherwise slim components are products over intervals.
Group torus-boundary zero pieces and torus slim cylinders. Their
components are individual solid tori or twisted Klein-bottle
interval bundles, pairs glued along a torus, or torus cylinders;
FC41 makes this a graph submanifold. Remove its interior, retaining
torus attaching maps.

For the remaining sphere-boundary pieces replace each punctured
$\mathbb{RP}^3$ by a ball and retain its connected-sum record.
A chain through sphere slim cylinders gives a ball, a sphere
cylinder or a closed sphere component. Cap sphere-cylinder cuts,
recording whether reconstruction joins components or contributes
$S^1\times S^2$. There are finitely many pieces, so these explicit
operations terminate. Remove any edge component whose base is a circle; it is a solid
torus. Retain its torus gluing.

The remaining zero/slim vertices are balls, and the edge pieces
are product disk handles. FC40 gives degree two at each vertex,
counting loop incidences twice. A finite graph all of whose vertices
have degree two is a disjoint union of cycles, allowing loops and
parallel edges. FC41 identifies each corresponding embedded union
as a solid torus. The rest is the supplied circle bundle
$M^{2\text{-stratum}}$. Glue these pieces along their entire torus
boundaries, then reverse the recorded torus removals and sphere
operations using FC41. This yields the required finite presentation.
\end{proof}
The argument proves neither incompressibility of its tori nor a
geometric metric on every piece. Those are downstream obligations.
Its elementary incidence argument must not replace the supplied
three-dimensional recognition or gluing witnesses.

\begin{foundation}[Nearly cuspidal boundary augmentation; FC43]
\label{found:fibration-boundary-assembly}
Use LC88's separate KL16.1 input on a compact connected orientable
manifold with boundary: small torus boundary diameters; actual
$C^{K+1}$ nearly isometric cusp collars of depth $100$; the volume
condition only where $d(p,\partial M)>10$; and the stated
whole-ball curvature-derivative bounds, for the SAME metric.
Adjoin boundary coordinate blocks $H_\partial$ to EVERY $Q_j$,
including $Q_4$. The smoothed inward collar coordinate has support
between levels $20,90$ and plateau between $30,80$.
Its block is $(\eta_i\zeta_i,\zeta_i)$, with physical collar
units, not a local $\rho(p_i)$ factor.
\end{foundation}
KL16.4--16.6 supply the collar adapted coordinate and separation
from zero balls. Handle the exceptional whole carrier
$I\times T^2$ first. Otherwise repeat the adjustments while
retaining these blocks like zero blocks, then extract a torus
collar core by $N_{35}(\partial_iM)$ and the adjusted marker
$x''_i\ge.9$ with ratio $\le40$.
It is $I\times T^2$ with the original boundary label preserved.
Treat it in FC42 as a torus-boundary piece without a zero core.
The extra $r_\partial$ is chosen after $\Upsilon'_0$ in the
source parameter order. This is a source-qualified boundary
implementation, not a reduction to the closed theorem by doubling.
The phrase ``closed'' in Standing Assumption16.3 is incompatible
with its boundary hypotheses; retain the manifold-with-boundary
reading of Theorem16.1 and record the discrepancy.

\subsection{Intrinsic boundary scales before the collar construction}
\label{sec:fibration-boundary-scale-adapter}
The closed construction is already bound in
Section~\ref{sec:fibration-closed-parameter-binding}. Its boundary
counterpart first needs the following adapters. All distances and
balls below use the intrinsic length metric of the SAME smooth
metric on the compact manifold with nonempty boundary. We use
$\mathrm{II}(v,v)=\langle\nabla_v\nu,v\rangle$ for the outward
unit normal. A near-isometric cusp chart means the $C^{K+1}$
embedding and pulled-back $C^K$ metric control of KL3.4/FC43,
with its relative norms measured against
$h=dz^2+e^{-z}g_T$; $g_T$ is flat on the reference torus.
The constants below are independent of its lattice and aspect ratio.

\begin{lemma}[Cusp first-exit and curvature-scale buffers; BSA01]
\label{lem:fibration-cusp-first-exit-buffers}
There are numerical $\delta_*>0$ and $C_0=1000$ such that
the FC43 collars with metric error at most $\delta\leq\delta_*$
and boundary-component diameter at most $\delta$ have these
properties. One may interpret this diameter even in the ambient
intrinsic distance, the weaker convention.
The external boundary is strictly convex, with
$\mathrm{II}\geq\frac14g|_{T\partial M}$. On the part
$0\leq z\leq98$ of each collar,
\[
 -\tfrac12\leq\sec_g\leq-\tfrac18,\qquad
 1/1.01\leq |dz|_g\leq1.01.
 \tag{BSA01.a}
\]
If $p$ is in this collar with $z(p)\leq95$ and $0<r\leq1$,
then its ENTIRE intrinsic ball stays in this collar, and
\[
 \Vol_g B_M(p,r)\leq C_0\delta^2 r,\qquad 1\leq R_p<3.
 \tag{BSA01.b}
\]
Every point with $d(p,\partial M)\leq10$ belongs to such a collar
with $z(p)<11$. At every $p\in M$ the curvature scale is finite;
writing $d=d(p,\partial M)$ gives
\[
 R_p\leq d+3.                                      \tag{BSA01.c}
\]
The same first-exit assertions hold for curves of length below
their indicated buffer, not just for chosen minimizing curves.
\end{lemma}
\begin{proof}
For the reference metric, $\sec=-1/4$ and the boundary has
outward normal $-\partial_z$ and second fundamental form
$g_T/2$. The $C^2$ portion of the supplied relative control
gives the displayed bounds and $1.01$ bi-Lipschitz control
after decreasing the numerical $\delta_*$. They do not depend
on a choice of flat lattice coordinates.

The image of $T^2\times[0,99)$ is relatively open in $M$ at
its external boundary; this uses that the map is an embedding
of pairs onto the ENTIRE indicated boundary component.
Its only exit inside $M$ is through the compact outer level.
Along a curve which has not exited, variation of $z$ is at
most $1.01$ times its length. A curve from height at most95
of length at most1 therefore cannot reach level98, hence
cannot leave the collar. A curve of length at most10 from
the boundary similarly stays below11. Compactness and the
intrinsic length metric supply minimizing curves, including
from a closest boundary point. The argument applies before
any use of their geodesic regularity.

To bound the reference torus diameter, join any two points
of the given boundary component by curves in $M$ of length
arbitrarily close to their ambient intrinsic distance, at
most $\delta$. These curves stay in the same collar below
$1.01(\delta+o(1))$. Projection onto the reference torus has
length at most
$1.01e^{1.01(\delta+o(1))/2}(\delta+o(1))<2\delta$.
Thus $\operatorname{diam}(T^2,g_T)\leq2\delta$ even with the
weaker diameter convention. A Dirichlet cell for its flat
lattice is contained in the Euclidean disk of this radius:
each cell vector realizes the distance from its torus class
to the origin. Consequently $\operatorname{Area}(T^2,g_T)
\leq4\pi\delta^2$. The reference volume form is
$e^{-z}dA_T\,dz$. A ball as above lies in a height interval
of length at most $2.02r$, and the actual volume form is at
most $1.01^3$ times the reference form. Integration over
the WHOLE torus and this interval proves (BSA01.b)'s
volume bound, with substantial slack.

The sectional lower bound on $B(p,1)$ gives $R_p\geq1$;
the sectional upper bound at $p$ gives
$R_p\leq\sqrt8<3$. For an arbitrary $p$, choose a closest
boundary point $b$ and the point $q$ of height1 above it
in that component's chart. Then $d(p,q)\leq d+1.01$ and
$\sec(q)\leq-1/8$. If a permissible curvature radius were
larger than $d+3$, its ball would contain $q$ and its
required lower bound would exceed $-1/9$, a contradiction.
The same negative point rules out an infinite curvature
scale. Taking the supremum of permissible radii proves
(BSA01.c), without treating infinity as a volume radius.
\end{proof}

\begin{lemma}[Local relative volume comparison with convex boundary; BSA02]
\label{lem:fibration-convex-boundary-volume-comparison}
Let $M$ be a compact smooth Riemannian three-manifold with
strictly convex smooth boundary. For $p\in M$, assume
$\operatorname{Ric}\geq-2\kappa^2g$ on the intrinsic ball
$B_M(p,L)$. With $V_\kappa$ as in the LC02 comparison
interface, for $0<a\leq b<L$ one has
\[
 \frac{\Vol B_M(p,b)}{V_\kappa(b)}
 \leq\frac{\Vol B_M(p,a)}{V_\kappa(a)}.             \tag{BSA02.a}
\]
The conclusion also holds at an endpoint controlled by
volume continuity. It includes centers ON the boundary.
This is the boundary adapter to the inherited radial
Jacobi/measure foundations, not an application of the
boundaryless completeness hypothesis in GSM77.
\end{lemma}
\begin{proof}
First, an intrinsic minimizing curve cannot touch the
boundary at an interior time unless it is constant.
Here is the local domain check. Extend the metric smoothly
across a boundary chart only for this local argument.
The inward signed distance $u$ satisfies $u\geq0$ in $M$
and $\operatorname{Hess}u|_{T\partial M}=-\mathrm{II}<0$.
Replacing $u$ by $u-Cu^2$, for suitable $C$, makes its
Hessian negative definite in a small collar chart, without
changing its sign. Choose a smaller ambient normal-convex
neighborhood in that chart. Its sufficiently short
geodesics with endpoints in $u\geq0$ stay there, since
the defining function is concave along them. Thus an
intrinsic minimizer is locally an ambient geodesic at a
putative boundary contact. At such a contact its velocity
is tangent and the defining function would have a
nonnegative local minimum, contradicting its strictly
negative second derivative. This also shows that a
nonconstant minimizer starting on the boundary has
strictly inward initial velocity. No metric on $M$ has
been altered.

Use the sphere of initial directions if $p$ is interior,
and its inward open hemisphere if $p\in\partial M$.
For each direction stop the ray at its first boundary
exit or when it ceases to minimize, whichever occurs
first. Tangential initial directions at a boundary
center cannot enter $M$ as minimizing rays and have
angular measure zero in any case. Except on a set of
volume zero, every interior point is reached uniquely
before this stopping time. To see the measure
qualification in this setting, distance from $p$ is
locally Lipschitz on the interior and hence differentiable
almost everywhere. Two different minimizing terminal
directions violate differentiability by the first
variation formula; coinciding terminal directions give
the same ray by uniqueness of the geodesic equation.
Critical values of the exponential map have measure
zero by the inherited smooth measure/Sard argument.
Away from these two exceptional sets the inverse
function theorem supplies the polar chart. The smooth
boundary itself has three-dimensional volume zero.
These are the usual inherited polar-measure facts on
interior segments, with the exit domain now explicit.

Let $J_\theta(t)$ be the radial Jacobian and put
$s_\kappa(t)=\sinh(\kappa t)/\kappa$, or $t$ if
$\kappa=0$. The inherited Riccati/Jacobi comparison gives
that $J_\theta(t)/s_\kappa(t)^2$ is nonincreasing while
the ray minimizes within $B_M(p,L)$. Extend this
nonnegative ratio by zero after its stopping time.
It remains nonincreasing; a boundary exit therefore
causes no adverse term. Integrating over the FIXED
initial sphere or hemisphere shows that
\[
 f(t)=\frac{\int_{\theta:\,t<\mathrm{stop}(\theta)}
                       J_\theta(t)\,d\theta}
                    {4\pi s_\kappa(t)^2}
\]
is nonincreasing for almost every $t<L$. The volume
ratio in (BSA02.a) is, by polar integration, the
weighted average of $f$ on $(0,t)$ with
positive weight $4\pi s_\kappa^2$. Such averages are
nonincreasing: the average on an earlier interval
is at least every subsequent value, and adjoining
a later interval cannot increase it. This proves
(BSA02.a). The same integral expression gives
volume continuity in the radius and the endpoint
assertion. The factor $4\pi$ is retained in the
boundary case; its small-radius ratio is $1/2$,
which does not affect this relative inequality.
\end{proof}

\begin{lemma}[Capped first scales and smooth scalar approximation; BSA03]
\label{lem:fibration-boundary-first-scales-and-smoothing}
Under BSA02 put $w_{\mathrm{cap}}=\omega_3/4$.
For EVERY $p\in M$ and $0<w<w_{\mathrm{cap}}$, the first
volume scale $r_p(w)$ is positive, finite, attained,
and strictly decreasing in $w$, with all the LC01
first-root inequalities. No first scale for all
$w<\omega_3$ is asserted at boundary centers.
For any $L_0$-Lipschitz scalar function $f$ on this
compact manifold and $\varepsilon>0$, there is a
function smooth up to the boundary with
\[
 \|\widetilde f-f\|_\infty<\varepsilon,\qquad
 \operatorname{Lip}_M(\widetilde f)<L_0+\varepsilon.
 \tag{BSA03.a}
\]
\end{lemma}
\begin{proof}
Smooth interior charts give
$\Vol B_M(p,r)/r^3\to\omega_3$ at interior centers.
A smooth boundary half-chart, after linear
normalization of the metric, gives limit
$\omega_3/2$ at a boundary center; the flattening
and metric errors tend to zero under dilation.
Both limits exceed $w_{\mathrm{cap}}$. Volume
continuity follows from the polar formula in BSA02,
and finite total volume makes this quotient tend
to zero as $r\to\infty$. The first crossing of
the value $w$ is consequently positive, finite,
and attained, with larger quotient at every
smaller radius. Continuity forces the first
crossing of a larger value to precede it strictly.
This is exactly the LC01 argument in the capped
range; it uses no Euclidean-normalized volume
monotonicity.

For the scalar assertion cover $M$ by finitely
many sufficiently small interior charts and
boundary half-charts. After linear normalization,
each chart metric is as close to the Euclidean
metric as prescribed. In a boundary half-chart,
reflect the scalar function across its flat
coordinate boundary. Reflection preserves the
Euclidean Lipschitz bound: folding a straight
segment into the half-space does not increase
its length. On smaller charts convolve this
extension with a smooth positive mollifier.
Interior charts use the same convolution without
reflection. The local approximants have
arbitrarily small uniform error and differential
norm at most $L_0+\varepsilon/3$, after first
choosing the chart distortion small enough.
This reasoning also covers $L_0=0$ (a constant
on each connected component).

Choose a smooth partition of unity $\chi_j$
subordinate to these smaller charts. In the
differential of $\sum_j\chi_jf_j$, the extra
term is $\sum_j(f_j-f)d\chi_j$, because
$\sum_jd\chi_j=0$. There are finitely many
charts and the sum of their differential norms
is bounded on compact $M$. Choose the local
uniform errors small enough that this term has
norm below $\varepsilon/3$ and the final
uniform error is below $\varepsilon$.
This gives a smooth function up to the boundary
with gradient norm strictly below
$L_0+\varepsilon$. Integrating along curves
in the intrinsic metric proves (BSA03.a).
The scalar reflection is an approximation
device; it neither doubles the carrier for
the collapse theorem nor smooths its metric.
\end{proof}

\begin{proposition}[Boundary counterexamples give uniform small-scale data; BSA04]
\label{prop:fibration-boundary-standing-sequence-adapter}
Fix $K\geq10$ and the positive finite function
$A:(0,\omega_3)\to(0,\infty)$ in the same-metric
static hypotheses of FC43/LC88. If that static
conclusion has no positive threshold, choose
counterexamples with
\[
 \delta_n=\min\{\delta_*,\omega_3/8,1/(16n^4)\},
 \qquad n\longrightarrow\infty.
\]
They satisfy, on one uniform tail and at EVERY
point including the boundary,
\[
 R_p>2n\,r_p(1/n).                              \tag{BSA04.a}
\]
All auxiliary first scales below are restricted
to $w<w_{\mathrm{cap}}$. For $C>0$ set
$D=\max\{1,C\}$ and define ONCE
\[
 A'(C,w)=\max_{0\leq k\leq K}
                A(D^{-3}w)D^{-k-2}.
 \tag{BSA04.b}
\]
For $0<C<n$ and $1/n\leq w<w_{\mathrm{cap}}$,
on the whole ball $B_M(p,Cr_p(w))$,
\[
 |\nabla^k\Rm|\leq A'(C,w)\,r_p(w)^{-k-2}
                 \quad(0\leq k\leq K).         \tag{BSA04.c}
\]
The function $A'$ is independent of $n,p$ and
the boundary component. This corrects the
auxiliary range in the literal boundary
reading of KL Standing Assumption16.3.
\end{proposition}
\begin{proof}
For $d(p,\partial M)>10$ apply the static
volume assumption at the finite radius $R_p$.
Since $\delta_n<1/n$, continuity gives the
first $1/n$ scale $r<R_p$. Monotonicity of
ball inclusion, not of volume ratio, gives
\[
 r^3/n=\Vol B(p,r)\leq\Vol B(p,R_p)
                  \leq\delta_nR_p^3.
\]
Thus $R_p/r\geq(1/(n\delta_n))^{1/3}>2n$.
For $d(p,\partial M)\leq10$, BSA01 applies.
Put $a_n=2\sqrt{C_0n}\,\delta_n$. For all
large $n$, $a_n\leq1$ and $a_n<1/(2n)$;
then
\[
 \frac{\Vol B(p,a_n)}{a_n^3}
 \leq\frac{C_0\delta_n^2}{a_n^2}
 =\frac1{4n}<\frac1n.
\]
The first $1/n$ scale is less than $a_n$.
Since $R_p\geq1$ here, (BSA04.a) follows.
All estimates were uniform in the centers.

Fix $C,w$ in the asserted range and write
$r=r_p(w)\leq r_p(1/n)$. For $n>1$,
$D<n$ and $Dr<R_p$. Also
$\Vol B(p,Dr)\geq\Vol B(p,r)=wr^3
=(D^{-3}w)(Dr)^3$. The triggering volume
parameter $\widehat w=D^{-3}w$ satisfies
$\delta_n<1/n^4<\widehat w<\omega_3$.
The original whole-ball derivative
hypothesis at radius $Dr$ applies with
this $\widehat w$. Multiplying its bound
by $r^{k+2}$ gives
$A(D^{-3}w)D^{-k-2}$, hence (BSA04.c).
The use of $D=\max\{1,C\}$ includes
$C<1$ without reversing ball inclusion.
No supremum of $A$ over a compact interval
and no regularity of $A$ has been assumed.
\end{proof}

\begin{proposition}[One smooth scale on the whole boundary carrier; BSA05]
\label{prop:fibration-boundary-modified-scale}
For the sequence in BSA04 fix $\Lambda>0$,
$0<w<w_{\mathrm{cap}}$, and
$w'=w/[2(1+2\Lambda^{-1})^3]$.
On one sufficiently late tail there is a
smooth positive $\Lambda$-Lipschitz function
$\rho$ on ALL of $M$, up to its boundary,
such that
\[
 \tfrac12 r_p(w)\leq\rho(p)\leq2r_p(w').
 \tag{BSA05.a}
\]
For every fixed $\varepsilon_\partial>0$,
on a further tail it also satisfies
$\rho<\varepsilon_\partial$ throughout
ALL collar regions $0\leq z\leq95$.
The tail is uniform over components and
over every admissible choice of $\rho$.
\end{proposition}
\begin{proof}
Set $l(p)=r_p(w)$, $u(p)=r_p(w')$.
The entire LC02 envelope comparison
works for intrinsic balls: its two ball
inclusions use only the triangle inequality,
BSA04.a supplies the curvature bound
uniformly on their fixed outer buffer,
and BSA02 supplies the relative comparison
even when its center is on the boundary.
Explicitly, a violation of
$l(p)-(\Lambda/2)d(p,q)\leq u(q)$ would
give, with $S=1+2/\Lambda$ and
$x=u(q)/l(p)<1$,
\[
 \frac{x^3}{J_c(x)}\geq\frac{2S^3}{J_c(S)}.
\]
Choose $c$ using
$T=1+2\max\{\Lambda,\Lambda^{-1}\}$ so
$1\leq3J_c(t)/t^3<3/2$ on $(0,T]$.
The left side is at most3 and the right
side is greater than4. The outer
curvature buffer is available once
$2n>1+4/\Lambda$ and $2n>1/c$,
as well as $1/n\leq w$. This explicitly
retains LC02's corrected interval.

MC19, which is a metric envelope argument,
now gives $l\leq f\leq u$ and
$\operatorname{Lip}(f)\leq\Lambda/2$.
It is positive everywhere and continuous,
so compactness gives $\min f>0$.
Apply BSA03 with error less than
$\min\{\min f/2,\Lambda/2\}$ to obtain
(BSA05.a). For the last assertion,
choose any fixed $0<a<\min\{1,
\varepsilon_\partial/2\}$. BSA01 gives
$\Vol B(p,a)/a^3\leq C_0\delta_n^2/a^2
<w'$ on one uniform tail. Thus
$r_p(w')<a$ and $\rho(p)<2a$ for
every indicated collar point. This
argument does not require a collapse
volume hypothesis near the boundary.
\end{proof}

\begin{proposition}[Boundary analytic data and interior exhaustion; BSA06]
\label{prop:fibration-boundary-analytic-interior-buffers}
Fix the parameters of BSA05 and define
$v_*=w'/(24I(1))$, $\mathcal A(R)
=2^{K+2}A'(2R+2,w')$, and $H_n=n/4$.
The volume lower bound, sectional bound,
and whole-ball derivative conclusions
of LPA01 hold on the intrinsic scaled
balls at EVERY point of the boundary
carrier, with $\alpha=n$ and these same
constants. For every center with
$d(p,\partial M)>10$ one additionally has
\[
 \frac{d(p,\partial M)}{\rho(p)}>n/2,
 \tag{BSA06.a}
\]
on the fixed-parameter tail. Hence every
fixed scaled buffer at such centers,
and every geodesic tested inside a
smaller fixed buffer, lies in the
smooth INTERIOR on one uniform tail.
The interior local-model producers may
use these exhausting domains; their
boundary-point analogues are not being
asserted.
\end{proposition}
\begin{proof}
First-scale attainment and BSA02 give
exactly the LC03 upper and LC04 lower
volume computations, with their same
integrals and constants. In particular
$\Vol_{\rho^{-2}g}B(p,1)\geq v_*$.
For $u=r_p(w')$, BSA04 gives
$R_p>2nu$ once $1/n\leq w'$,
whereas $\rho\leq2u$. Thus
$H_n\rho<R_p/2$, and the defining
curvature bound yields the LPA01
sectional conclusion after rescaling.
BSA04.c at $C=2R+2<n$ supplies the
derivative bound on the WHOLE ball;
multiplication by $\rho^{k+2}$ costs
at most $2^{K+2}$. The proof does not
require monotonicity of $A'$. In
particular $n>2B+2$ handles every
real $0<R<B$ on a single tail.

Write $d=d(p,\partial M)>10$.
BSA01.c and BSA04.a give
$2nu<R_p\leq d+3<1.3d$.
Combining with $\rho\leq2u$ yields
$d/\rho>n/1.3>n/2$. A path from $p$
to the boundary must have length
at least $d$. Thus a fixed ball of
scaled radius $B$, together with
all curves tested in it, has an
interior buffer tending to infinity.
For example, minimizing segments
between points of its radius-$B$
subball remain in the radius-$3B$
ball and are interior once $3B<n/2$.

For clarity, the locality needed by
the existing metric and finite-category
model arguments is as follows. Local
volume and triangle comparisons use
only those bounded interior segments.
Their finite nets and packing tests
lie in a larger fixed buffer. Smooth
compactness and its coordinate
extraction use the whole-ball
derivative and noncollapse estimates
on each fixed buffer; a diagonal
exhaustion then gives a complete
boundaryless limit because the
available radii tend to infinity.
The cut-locus, harmonic-coordinate
and finite-category foundations
retained there are unchanged.
All their later compact test
domains, model embeddings and
isotopy enclosures are contained
in one such fixed buffer once the
parameters are fixed. This supplies
their interior-source domain check,
rather than applying a theorem
whose source is required to be a
closed manifold directly to a
boundary point. It makes no claim
that such a boundary point has a
boundaryless smooth blow-up.
\end{proof}

\paragraph{Scope and source comparison.}
KL16.1--16.4, printed87--89/PDF82--84,
give the boundary input, its stated
standing sequence and the elementary
collar claims. Their literal auxiliary
range $w<\omega_3$ at every boundary
point is not used here. BSA03 restricts
it to $w<w_{\mathrm{cap}}$; LC02--LC04,
LC09, LC18 and LPA01 need only the two
fixed small parameters $w,w'$.
In their threshold choices intersect
the allowed interval with
$(0,w_{\mathrm{cap}})$ before fixing
these parameters. Derivative triggers
in BSA04 retain the original larger
range in the STATIC hypotheses.
GSM77, lines6387--6409 and6796--6900,
states boundaryless comparison and
gives its radial mechanism, leaving
relative monotonicity as an exercise.
BSA02 supplies that monotonicity and
the boundary exit translation under
the inherited analytic foundations.
KL3.14--3.15, printed26/PDF21, omits
the mollification details; BSA03
supplies the scalar half-chart
adapter needed here.
Actual adapted cusp coordinates,
collar separation, augmented clouds,
retained torus cores and their
parameter rows remain to be written.
These six arguments alone do not
complete FC43 or the full user goal.

\subsection{Actual cusp coordinates and separated original collars}
\label{sec:fibration-actual-cusp-collars}
We now use the boundary sequence and smooth scale of BSA04--BSA06.
The quantities below are uniform in the component, the eligible
center and the member of a sufficiently late tail. Finite
approximations on an individual compact carrier may be chosen
after that carrier is given; their errors are bounded by the
specified uniform tolerances. Write $z_i$ for the original
finite-category height in the $i$th embedded cusp.

\begin{lemma}[Smooth original collar height with controlled geometry; BCP01]
\label{lem:fibration-smooth-original-cusp-height}
On every original collar there is a smooth function $\eta_i$
on the band $1<z_i<99$ such that, on $2\leq z_i\leq98$,
\[
 |\eta_i-z_i|+\|d(\eta_i-z_i)\|_g
       +\|\nabla^2(\eta_i-z_i)\|_g<\delta_n.       \tag{BCP01.a}
\]
The functions may be chosen simultaneously for all components.
On this band, after enlarging the tail,
\[
 .99<|d\eta_i|_g<1.01,\qquad
 |\nabla^2\eta_i|_g<2,\qquad
 .99<\partial_{z_i}\eta_i<1.01.                 \tag{BCP01.b}
\]
Every level $\eta_i=t$, $3\leq t\leq97$, is an
actual smooth torus, with second fundamental form
of norm less than3 and intrinsic diameter at most
$10\delta_n$. The level bands have smooth product
trivializations preserving $\eta_i$.
Moreover $\rho\to0$ uniformly on $2\leq z_i\leq96$.
The component
\[
 B_i\subset M\setminus\{\eta_i=90\}
 \quad\hbox{containing }\partial_iM
 \tag{BCP01.c}
\]
is precisely the original inner collar cut by this
level. Its closure is smoothly diffeomorphic as a
pair to $(T^2\times[0,1],T^2\times\{0\})$, with the
original external boundary label retained.
\end{lemma}
\begin{proof}
The finite-category height is at least $C^3$ on
the original smooth interior carrier. On the
compact band $2\leq z_i\leq98$, finite smooth
charts, convolution and a partition of unity
give arbitrarily close $C^2$ scalar approximation,
defined on a slightly larger open band.
All partition derivatives are bounded on the
finite cover and the local errors can be made
arbitrarily small in every needed derivative.
Use covariant norms of the SAME smooth metric
when selecting these errors. Compactness gives
only finitely many boundary components, so this
can be done simultaneously. This smooths the
height, not the metric or its comparison chart.
The first-exit and volume proof of BSA01 also applies
for $z_i\leq96$ and radii at most1: such curves stay
below $96+1.01<98$. Repeating BSA05's
fixed small test-radius argument proves the extra
uniform smallness of $\rho$.

For the reference metric, $|dz_i|=1$ and
$\nabla^2z_i=-\tfrac12(g-dz_i^2)$.
The actual metric is $C^2$ close and (BCP01.a)
holds, proving (BCP01.b). Along each original
vertical segment $\eta_i$ is strictly increasing.
Every value $t\in[3,97]$ has exactly one preimage
on that segment, with $|z_i-t|<\delta_n$. Thus
the level is a graph over the original reference
torus. Its graph slope, measured in the
reference level metric, is $O(\delta_n)$:
the tangential derivative of $\eta_i$ differs
from that of $z_i$ by (BCP01.a), and the
vertical derivative is bounded below.
The graph metric is consequently bounded by a
fixed multiple of $e^{-t}g_T$. BSA01's reference
diameter bound gives the asserted $10\delta_n$
bound. Its second fundamental form is
$\nabla^2\eta_i/|d\eta_i|$ on tangent vectors,
so its norm is less than3.

The original graph parametrization is a
$C^2$ diffeomorphism to this smooth level.
LFR04 supplies its smooth torus type.
The smooth field
$\nabla\eta_i/|d\eta_i|^2$ flows between
any two levels in $[3,97]$: each finite
level slab is compactly contained in the
original collar and the field raises
$\eta_i$ at speed one. Forward and backward
flows give the claimed smooth product.
At level90 the original reference subgraph,
including the external boundary, is a
compact $C^2$ product with boundary.
Its image is a smooth compact submanifold
because its two boundary components are
the smooth external boundary and a regular
smooth level. Apply LFR04 as a pair; its
approximation is small enough to retain
the two separate boundary components.
The inner region is open and closed in
the complement of the level, and connected,
so it is exactly (BCP01.c).
\end{proof}

\begin{proposition}[The actual cusp splitting with its torus factor; BCP02]
\label{prop:fibration-actual-cusp-product-splitting}
Fix any $0<\beta<\gamma<1/10$ and any fixed
additional normalized buffer $L<\infty$.
At EVERY point with $5\leq\eta_i(p)\leq95$,
on one uniform tail, $(M,\rho(p)^{-2}g,p)$
has an ACTUAL pointed $(1,\beta)$-splitting
whose Euclidean coordinate on all its tested
balls is
\[
 \Phi_p(q)=\frac{\eta_i(q)-\eta_i(p)}{\rho(p)}.
 \tag{BCP02.a}
\]
The function $\Phi_p$ is an adapted coordinate
of quality $\gamma$ on the unit ball in
KL4.21's convention, including ALL minimizing
initial directions in its long test.
The same height is available on the additional
radius-$L$ buffer. The complete residual factor
is the flat torus with metric
$\rho(p)^{-2}e^{-z_i(p)}g_T$; it is NOT
replaced by a point. In particular taking
$\beta=\beta_1$ (or a smaller parameter)
excludes the zero stratum throughout this band.
\end{proposition}
\begin{proof}
Write $r=\rho(p)$, $z_0=z_i(p)$ and
$H=\max\{10/\beta,10/\gamma,10L,100\}$.
BSA05 makes $rH$ arbitrarily small uniformly
in this band: for late $n$ it lies in
$4<z_i<96$, where that assertion applies.
Also $\delta_n\to0$. Freeze the reference
warping factor at $z_0$ and use
$\eta_i-\eta_i(p)$ as the vertical coordinate.
By (BCP01.a), its differential differs
arbitrarily little from $dz_i$.
On the entire physical radius-$Hr$ buffer,
comparison with the frozen product has
bi-Lipschitz distortion
\[
 1+\epsilon,\qquad
 \epsilon\leq C(\delta_n+Hr),                 \tag{BCP02.b}
\]
for a numerical $C$. This follows directly
from $e^{-(z-z_0)}=1+O(Hr)$ and the
near-isometric pullback; it does not
depend on the flat lattice or on $r^{-1}$
times its diameter.

The actual map on that buffer is (BCP02.a)
together with the original torus coordinate.
Its basepoint is the origin and the original
torus basepoint. Intrinsic distances on the
smaller tested balls have the same estimates:
any competing curve short enough to improve
a tested distance stays in this larger
collar buffer by first exit, and product
minimizers also stay there. Hence differences
of tested distances are at most
$C\epsilon H$. Conversely each product point
in the tested radius $\beta^{-1}-\beta$
ball has a unique height preimage, by
$\partial_z\eta_i> .99$, and the distance
estimate places it in the source
$\beta^{-1}$ ball once $C\epsilon H<\beta/10$.
This proves both KL3.2 requirements, with
strict slack, by further increasing $n$.
Extend the map arbitrarily outside its
tested buffer to obtain the globally
defined pointed approximation. The target
$\mathbb R$ times the frozen flat torus
is complete; no small-diameter assertion
for that torus at scale $r$ was used.

The adapted differential test can be
checked directly on the SAME map. In the
scaled metric, $\|\nabla^2\Phi_p\|\leq2r$.
For a unit-speed minimizing segment of
scaled length $\ell$ from $m\in B(p,1)$
to $m'\in B(p,\gamma^{-1})$, the entire
segment is inside the buffered collar.
Integrating its second derivative gives
\[
 \left|D\Phi_p(v)
 -\frac{\Phi_p(m')-\Phi_p(m)}{d(m,m')}\right|
 \leq r\ell\leq r(1+\gamma^{-1})<\gamma.
 \tag{BCP02.c}
\]
This holds for EVERY such segment and thus
every permitted initial direction; in
particular it holds when $d(m,m')>1$ as
in KL4.21. Smoothness is actual smoothness,
and (BCP01.b), with the sharper errors
tending to zero, gives Lipschitz constant
less than $1+\gamma$ on the unit ball.
Minimizing joins between its points remain
in the larger smooth buffer. Vertical
segments give coordinate values on both
sides of the origin within $\gamma$ of
$\pm1$, at points of the unit ball:
take the target values $\pm(1-\gamma/2)$
and use (BCP02.b). The Lipschitz bound,
chosen less than $1+\gamma/2$, gives the
other Hausdorff inclusion. This verifies
the image clause of adapted coordinates.
All required bounds are strict upper
bounds on $\delta_n$ and $r$, after the
fixed parameters. They hold on ONE tail,
uniformly in $p$ and $i$.
\end{proof}

\begin{theorem}[Actual collar overlap is the whole torus product; BCP03]
\label{thm:fibration-cusp-collar-product-or-disjoint}
For all sufficiently late members of the boundary
sequence, either $M$ is smoothly diffeomorphic
to $T^2\times[0,1]$, preserving its two external
boundary labels, or the enlarged relative open
collars $B_i^+$ cut at $\eta_i=92$ are pairwise
disjoint. In the latter case the original
$B_i$ are pairwise disjoint and
\[
 d_M(\overline B_i,\overline B_j)\geq1
                          \quad(i\ne j).       \tag{BCP03.b}
\]
The enlargement is only a separation buffer;
the source's level90 collar and block support
are kept.
\end{theorem}
\begin{proof}
Suppose $B_i^+\cap B_j^+\ne\varnothing$ for
distinct components. Put $T_i=\{\eta_i=92\}$.
The original inner product for $B_i^+$ contains
no external boundary other than $\partial_iM$.
A path in $B_j^+$ from $\partial_jM$ to the
intersection must therefore cross $T_i$ at
some $q\in B_j^+$. Thus $\eta_j(q)<92$.

Every curve from $T_i$ to a different
external boundary must first leave the
$i$th original collar through height99.
It costs more than6 in length, by BSA01
and $|z_i-92|<\delta_n$. On the other hand
the original vertical $j$-curve from its
boundary to $q$ has length at most
$1.01z_j(q)$. Hence $z_j(q)>5.9$,
while $z_j(q)<93$.
BCP01 gives $\operatorname{diam}_{T_i}T_i
\leq10\delta_n$. Short intrinsic paths on
$T_i$ starting at $q$ cannot exit the
$j$th band $5<\eta_j<95$, by first exit
and the uniform height derivative bound.
Consequently the ENTIRE $T_i$ lies in
that band and
$\operatorname{osc}_{T_i}\eta_j\leq
20\delta_n$.

Let $F=\eta_j|_{T_i}$. The ambient Hessian
bound in BCP01, its gradient bound and the
second fundamental form bound for $T_i$
give $|\nabla^2_{T_i}F|\leq10$.
For any unit tangent direction with
$b=|dF(v)|$, choose its sign positive and
follow the complete intrinsic geodesic for
time $b/10$. Taylor integration gives an
increase at least $b^2/20$. Therefore
\[
 \|dF\|^2\leq20\,\operatorname{osc}_{T_i}F
                  \leq400\delta_n.             \tag{BCP03.a}
\]
No injectivity-radius lower bound is
needed to follow this geodesic; the
Taylor estimate remains valid if it
passes through a loop.

Use the smooth $\eta_j$ flow product
on its compact levels from5 to95.
The projection of $T_i$ to any one
of these reference levels is a local
diffeomorphism: (BCP03.a) and
$|d\eta_j|>.99$ say its flow direction
is transverse to $T_i$. It is proper
and onto, hence a finite covering
of the connected torus. Embedded
covering sheets in a torus times
an interval are ordered by their
$\eta_j$ heights. Local continuation
preserves their strict order, so
monodromy cannot permute them.
Ordering the finitely many heights
thus defines separate global
smooth graphs. Since $T_i$ is
connected, there is just ONE graph.

Let $N_i=\overline{B_i^+}$ and let
$N_j'$ be the $j$th region from its
external boundary up to that graph.
Each is a smooth torus product, by
BCP01 and the $\eta_j$ flow, and
they meet exactly along $T_i$.
Indeed the interior of $N_j'$ is
connected, avoids $T_i$, and meets
$\partial_jM$, which is outside
$N_i$; it cannot enter the component
$B_i^+$. The two regions occupy the
opposite sides of $T_i$, so their
union is relatively open in $M$,
as well as compact and hence
closed. It is nonempty; connectedness
gives $M=N_i\cup N_j'$.
Gluing two torus collars along their
inner boundary diffeomorphism is
a torus product: absorb this
diffeomorphism into the base
coordinate of one collar. Its two
remaining boundary components are
exactly the original labeled ones.
This proves the alternative using
actual collars, not only their
homotopy classes.

In the nonproduct case a curve from
$\overline B_i$ to $\overline B_j$
must leave $B_i^+$ before it
reaches $B_j^+$. Its passage
through the $i$th height band
from90 to92 costs at least
$2/1.01>1$. This proves
(BCP03.b), including endpoints
on either level90 torus.
\end{proof}

\begin{proposition}[Restricted interior packets and the actual remaining cover; BCP04]
\label{prop:fibration-boundary-restricted-local-packets}
The local producer proofs LPA02--LPA06 and
the overlap comparisons for the original
interior blocks admit the
following boundary version, with the
same allowed parameter dependencies.
Define the ranks on the interior region
$d(p,\partial M)>5$. Select circle, slim,
strong-edge and zero-ball centers only
where $d(p,\partial M)>10$.
One common assignment and uniform tail
supply their SAME original packets,
compact supports, enclosures, pointwise
zero witnesses and early multiplicity
bounds. Their actual covering domains
cover every point with
$d(p,\partial M)\geq20$.
The zero-ball family has its original
tenth-radius cover of the eligible zero
stratum; each selected ball MEETS that
stratum, and its center need not be in it.
No assertion that these interior families
cover the external boundary is made.
\end{proposition}
\begin{proof}
First check the producer domains, before
making any finite selection. BSA04--BSA05
in fact give at every point with $d>0$
\[
 \frac d{\rho(p)}>\frac{nd}{d+3}.                \tag{BCP04.a}
\]
For $d>5$ this exceeds $5n/8$.
Thus every fixed scaled ball needed
by a producer at such a point, and
every larger fixed comparison buffer
around it, is interior on a uniform
tail. BSA06 supplies its noncollapse,
sectional and derivative data.
In the metric-model contradiction,
the finite nets, local comparison
and local positive-volume argument
of LC05--LC15 take place on these
exhausting balls. They therefore
give the same low-dimensional
limit and a small-volume threshold
with the SAME allowed predecessors
as LC09. Intersect that threshold
with the closed one and
$(0,w_{\mathrm{cap}})$ before
choosing the common $w$.
This is a boundary-local version
of the proof, not an application
of LC08's literal boundaryless
source hypothesis.

For the finite-category models,
apply the harmonic-coordinate
compactness argument to each
fixed interior buffer and take
its usual diagonal exhaustion.
BSA06 provides every derivative
test on one tail for a bounded
set of radii. The limits are
complete boundaryless manifolds,
and retain exactly the finite
metric category of LFR14/LFR49.
The same cone, classification,
coordinate and compact-fiber
arguments then prove the local
packets and uniform joint-witness
conclusion LPA02. Their prescribed
test domains are compact and
bounded after the parameters are
fixed. In particular their
resulting finite upper zero
factor $V$ is still an OUTPUT,
not a bound chosen before
$\Lambda$. No assumption
$V\Lambda\ll1$ occurs.

For centers with $d>10$, enlarge
$n$ so every originally required
domain of radius $C\rho(p)$ has
$C\rho(p)<d/2$; this follows
from BSA06 for every one of
the finitely many fixed buffer
constants, including $400V$.
These entire domains lie in
$d>5$. All shell comparisons,
same-map transfer arguments and
whole-fiber enclosures therefore
use the valid rank region, even
if an auxiliary point is not
eligible as a selected center.
The original interior-block
comparison proofs use these
same domains and rank tests.
Adding boundary blocks still
requires its separate model
and locality argument.

Fix once the zero witnesses at
eligible centers and run the
LC62--LC64 selection with this
restricted index set. Positivity
of $\min\rho$ and finiteness of
$V\max\rho$ still give termination.
Every eligible zero point is an
eligible center candidate, so
the cover proof loses none of
its intended points. The shell
exclusion applies throughout
each tested ball as just checked;
it gives the same tenth-radius
cover, radial packets and
at-most-one-end conclusion.
The ball-meets-zero premise,
not a zero-center premise, is
what that proof uses.

Choose the circle and slim
families maximal among eligible
centers, exactly as LPA06.
The disjoint-small-ball argument
and scale comparability cover
every eligible point of their
respective strata. For strong
edges choose the maximal family
among eligible strong points.
A nonslim one-stratum point
with $d\geq20$ has the LFR44
strong witness within a fixed
$C\Delta\rho(p)$ of it.
By (BCP04.a), choose the
uniform tail so this distance
is less than $d/2$; that
witness is then eligible.
Maximality supplies the SAME
edge covering and the original
$3\Delta$ plateau as in
LFR44/LPA06. The source-domain
buffers also contain every
auxiliary weak-border witness
used later by FDC01.
Absence of three-strata and
these four cases give the
asserted cover of $d\geq20$.

Finally, packing uses only the
disjoint selected balls, their
positive volumes and the
larger controlled ball around
one selected center. It is
unchanged when the candidate
set is restricted. Its uniform
bound still precedes $\Delta$
and $w'$, as in LPA06; the
necessary larger curvature
and interior buffers are tail
conditions. After selection,
use ONE physical distance
smoothing for the edge family
with the previously prescribed
finitely many tolerances.
All cutoffs remain supported
strictly inside their original
smooth domains. This completes
the restriction argument
without extending the metric
or identifying $M$ with a
closed manifold.
\end{proof}

\begin{theorem}[Actual zero balls are disjoint from the cusp collars; BCP05]
\label{thm:fibration-cusp-zero-ball-separation}
In the nonproduct case of BCP03, take
the actual zero family from BCP04.
On one sufficiently late tail the
collection consisting of all $B_i^+$
and all selected open zero balls
$B(z,R_z^0)$ is pairwise disjoint.
This includes the actual selected
centers that are not zero-stratum
points themselves.
\end{theorem}
\begin{proof}
The $B_i^+$ are pairwise disjoint
by BCP03, and selected zero balls
are pairwise disjoint by their
unchanged LC64 selection.
Suppose one zero ball meets
$B_i^+$ at $x$. Its radius obeys
$R_z^0\leq V\rho(z)$.
Since $R_z>n\rho(z)$, take
$n>2V$ to put this whole ball
inside its curvature-scale
ball. At the meeting point,
BSA01 gives $\sec(x)\leq-1/8$,
whereas the curvature-scale
lower bound is $-R_z^{-2}$.
Thus $R_z\leq\sqrt8<3$,
and
\[
 \rho(z)<3/n,\qquad R_z^0<3V/n.                 \tag{BCP05.a}
\]
Take $n$ so large that the last
quantity is less than $1/100$.
First exit along a short join
from $x$ to $z$ puts $z$ in
the same original collar below
height93. Since $d(z,\partial M)>10$,
its height there is greater
than $10/1.01>9$. Another
first-exit test puts the
ENTIRE zero ball inside
$5<\eta_i<95$.
BCP02, with splitting error
less than $\beta_1$, excludes
the zero stratum at EVERY
point of this ball. That
contradicts the actual
selection premise that it
meets the eligible zero
stratum. It is unnecessary,
and would be incorrect, to
replace that premise by the
claim that $z$ is zero.
All new bounds are lower
bounds on the sequence index
after the finite $V$ is fixed;
none requires $V\Lambda$
to be small.
\end{proof}

\paragraph{What has and has not been assembled.}
BCP01--BCP05 supply the actual geometric
claims behind KL16.4--16.5, including its
omitted collar-overlap step, in the
source's metric and finite-category
conventions. BCP02 uses KL4.21's full
adapted-coordinate test, printed34/PDF29,
not merely a gradient-length estimate.
The product exception is already a
graph piece under FC41. In the other
case there are actual disjoint original
collars, valid restricted interior
families and separated zero balls.
The physical boundary blocks in EVERY
cloud stage, their marker locality,
the adjusted torus cores, saturated
boundary faces and the full boundary
parameter register are the next
obligations. No final goal audit or
Lean verification is claimed here.

\subsection{Actual boundary cloud blocks and adjusted torus cores}
\label{sec:fibration-actual-boundary-clouds}
Work in BCP03's nonproduct case. At the initial construction
use BCP04's restricted interior families, with
one additional eligibility buffer: circle, slim and zero centers
still have distance greater than10, while the strong-edge family
is maximal among centers of distance greater than20. BCG01
verifies the needed cover and collar consequences. Add one copy
$\mathcal H_b=\mathbb R^2$ for each external boundary component and
retain their direct sum in EVERY $Q_j$, including $Q_4$.
Define the fixed physical profile and block
\[
 \chi_\partial(t)=\chi_{20,30}(t)(1-\chi_{80,90}(t)),\qquad
 \mathcal{B}(t)=(t\chi_\partial(t),\chi_\partial(t)),\qquad
 F_b=\mathcal{B}(\eta_b),
 \tag{BCG.0}
\]
extended by zero outside the original smooth collar band.
The original interior blocks of $F$ keep their source formulas.
Choose the early profile constant $P$ to bound also
$\|\mathcal{B}'\|_\infty,\|\mathcal{B}''\|_\infty$ and all original
profiles; these are fixed numerical bounds. Add one to the
early whole-list and multiplicity constants. This is the
initial BOUNDARY assignment, not a modification of an already
constructed closed-carrier map. The late physical parameter
$r_\partial$ is chosen after all fixed normalized test radii
and the output zero factor $V$.

\begin{lemma}[Whole boundary support lists and their original domains; BCG01]
\label{lem:fibration-boundary-whole-support-lists}
Retain $L=10^6\Delta$ and $L\Lambda<10^{-5}$.
Use the ORIGINAL three comparison domains
\[
 D_a=B(p_a,C_aR_a),\quad R_a=\rho(p_a),\qquad
 C_a=10,\ 20\Delta,\ .95L
 \tag{BCG01.a}
\]
for circle, edge and slim references, respectively.
For $r_\partial<1/(1000L)$ and a sufficiently late
common tail, at most ONE closed boundary support
meets $D_a$. If it is the $b$th support, then
\[
 R_a<2r_\partial,\qquad
 D_a\subset\{19<\eta_b<91\}\subset B_b^+.
 \tag{BCG01.b}
\]
No zero support then meets $D_a$. Every boundary
block is smooth on the whole carrier, and
$\|DF\|\leq1000(N+2)P^2$ with the augmented
early multiplicity $N$. The estimate is independent
of the total number of boundary components.
The revised edge selection still gives the entire original
four-family cover on $d(p,\partial M)\geq35$.
Moreover every selected edge packet's original
radius-$100\Delta R_a$ domain lies in $d>10$, so its ENTIRE
LFR38 circle collar is covered by the selected circle family.
\end{lemma}
\begin{proof}
A meeting point $x$ has $20\leq\eta_b(x)\leq90$,
so BCP01 and BSA05 give $\rho(x)<r_\partial$.
Slow variation gives
$|\rho(x)-R_a|\leq C_a\Lambda R_a<10^{-5}R_a$;
hence $R_a<2r_\partial$. The diameter of $D_a$
is less than $4Lr_\partial<1/250$.
First exit and the height derivative bound
put its whole domain in the SAME smooth
collar band $19<\eta_b<91$.
BCP03 says that the closures of distinct
level90 collars have distance at least1.
Since every closed boundary support is in
such a closure, two supports cannot meet
this $D_a$. Its inclusion in $B_b^+$
and BCP05 exclude every zero support.
These arguments include closed-support
points where the cutoff itself is zero.

The boundary profiles and all their
nonconstant parts have compact supports
strictly within their smooth collar bands,
so extension by zero is smooth.
At any point at most one boundary block
or its derivative can contribute, by
the same separation. On its support
$|d\eta_b|<1.01$, whence
$\|d\mathcal{B}(\eta_b)\|\leq2P$.
Use CGP02's square-sum calculation,
counting this as one additional block
with derivative bounded by $82P$.
The old weak-edge calculation is unchanged.
Its generous bound $1000(N+2)P^2$ therefore
bounds the augmented map as well.

For the eligibility assertion,
BCP04.a makes every fixed
$C\Delta\rho$ displacement less
than a prescribed fraction of
the distance to the external
boundary on one uniform tail.
A nonslim one-stratum point
with $d\geq35$ therefore has
its LFR44 strong witness at
distance greater than20.
Maximality among those strong
centers proves the same original
$3\Delta$ edge-cover conclusion.
The other three selections and
their cover arguments are unchanged.
Conversely, for a selected edge
center $d_a>20$, require
$100\Delta R_a<d_a/2$.
Its WHOLE original edge domain
then lies in $d>10$. Every
two-stratum point in its LFR38
circle collar is an eligible
circle-center candidate, so
circle maximality covers it.
All additional requirements are
uniform lower bounds on the
sequence index after the fixed
parameters. No already chosen
map or family is replaced.
\end{proof}

\begin{lemma}[The actual physical height has an affine comparison; BCG02]
\label{lem:fibration-boundary-actual-affine-height}
For any requested $0<\theta<1/100$ in one of the
three rough-graph constructions, the original
reference qualities can be chosen in their
permitted order so that the following holds.
If boundary support $b$ meets $D_a$, there is
a unit row $A_b:\mathbb R^{k_a}\to\mathbb R$
(a sign when $k_a=1$) such that, on ALL $D_a$,
\[
 \begin{split}
 U_b&=\frac{\eta_b-\eta_b(p_a)}{R_a},\\
 |U_b-A_b(\eta_a-\eta_a(p_a))|&<\theta,\\
 \|DU_b-A_bD\eta_a\|&<\theta .
 \end{split}                                             \tag{BCG02.a}
\]
Differential norms use $R_a^{-2}g$.
For circles all requested reference accuracies
are fixed BEFORE $\Delta$; they do not depend
on a later long test length. The remaining
cusp accuracy and physical buffer bounds are
satisfied by $r_\partial$ and a uniform tail,
using the SAME smooth heights $\eta_b$.
\end{lemma}
\begin{proof}
BCG01 puts the reference center in the
original cusp band. BCP02 supplies an
arbitrarily accurate actual rank-one
splitting there, whose Euclidean
coordinate is exactly $U_b$. It retains
its frozen torus factor. It may be
tested on any prescribed fixed ball
after the late physical and sequence
choices. On the reference side use
the ORIGINAL raw coordinate $u_a$.

For a circle reference apply AC76 with
ranks two and one, using the original
no-three-splitting quality $\beta_3$.
FC20--FC21 replace the Euclidean factor
map by a unit row, exactly as in TCP02.
This gives
$|U_b-A_bu_a|<E$ on the radius1000
ball, with early prescribed $E$.
The reference distortion and coverage
are on its original radius5000 ball.
All these are EARLY fixed radii;
the quality of the cusp splitting
can be supplied later independently.
For edge references EGP01 supplies
absence of a $(2,\beta_2)$-splitting;
for slim references this is their
one-stratum property. Use AC76 and
FC20 with ranks one and one as in
EGP03 and SGP02, obtaining a sign
and the same raw comparison on
the required radius-$30L$ ball.
These reference qualities may depend
on the already chosen $\Delta$.
The common curvature tail supplies
each compatibility buffer.
Translations are fixed by the
centered $U_b(p_a)=u_a(p_a)=0$.
The original smooth reference value
errors, including their value at
$p_a$, give the value estimate in
(BCG02.a) once $E$ and those errors
are less than $\theta/4$.

Here is a direct check of its
differential conclusion, preserving
the original derivative tests.
For $x\in D_a$, lift a reference
product point in the $A_b^T$
direction from $u_a(x)$, keeping
the residual coordinate fixed.
Use length $400$ for a circle,
$400\Delta$ for an edge, and
$10L$ for a slim reference.
The original coverage and distortion
give a witness $y$ and a minimizing
initial unit vector $v$ from $x$
to $y$. In the circle case $y$
is within radius500 and its
separation exceeds201; in the
edge case its radius is less
than $430\Delta$ and separation
exceeds $100\Delta$; in the
slim case its radius is less
than $12L$ and separation
exceeds $L$. Thus every endpoint
lies in the ORIGINAL appropriate
far-test ball (1000, $1000\Delta$,
or $L/\varsigma_{\rm slim}$).
The raw comparison above applies
at both endpoints. The original
adapted test consequently makes
$A_bD\eta_a(v)$ arbitrarily close
to1 from below.

Choose $r_\partial$ smaller, after
these fixed radii, so the larger
physical balls and all intervening
segments stay in $5<\eta_b<95$.
In reference units
$\|\nabla^2U_b\|\leq2R_a$.
Writing $\ell=d(x,y)$ in those
units, Taylor integration yields
\[
 \left|DU_b(v)-\frac{U_b(y)-U_b(x)}{\ell}\right|
                      \leq R_a\ell.              \tag{BCG02.b}
\]
The raw alignment and axis lift
make the quotient close to1.
For example choose the reference
adapted qualities, raw errors and
cusp norm errors below
$\theta^2/10^8$, and impose
$2r_\partial(12L+1000)<\theta^2/10^8$.
The endpoint distortions divided
by their positive separations
then give saturation errors
less than $\theta^2/10^5$
for both covectors. Their norms
are at most one plus their
prescribed small errors.
FC15's pointwise Riesz estimate
gives $\|DU_b-A_bD\eta_a\|<\theta$.
In the circle case only its
fixed short derivative test was
used; the late bound involving
$L$ is imposed solely on
$r_\partial$, not on its
earlier quality. All comparisons
use actual minimizing directions
and the originally selected maps.
\end{proof}

\begin{theorem}[Actual augmented clouds and all three adjustments; BCG03]
\label{thm:fibration-actual-augmented-boundary-clouds}
With the enlarged EARLY numerical constants
specified above, the augmented map $F$ has
the actual TCP, EGP and SGP rough-graph,
cloud, all-preimage rank and section
properties. Its chosen planes follow
CFS27 on the interior blocks and retain
all boundary blocks. GAF01's ordered
construction therefore gives actual
adjustments and $E:M\to H$ satisfying
GAF02--GAF07, including the original
zero-marker bound, one-sheet bases,
no merging and proper whole circle
and slim fibers, in the boundary
setting. All required source domains
and compact fiber traces lie in
$\operatorname{int}M$.
The first adjusted scale also retains
EDP01's early derivative constant.
\end{theorem}
\begin{proof}
For a reference $a$, leave its original
interior rough-graph formulas unchanged.
If no boundary support meets $D_a$,
set every boundary model block to
zero. If support $b$ meets it, put
\[
 h_b(z)=\eta_b(p_a)+R_aA_b(z-\eta_a(p_a)),\qquad
 \Phi_{a,b}(z)=R_a^{-1}\mathcal{B}(h_b(z)).
 \tag{BCG03.a}
\]
All other boundary model blocks
are identically zero. This formula
is smooth on the whole parameter
space, though the actual coordinate
is evaluated only in its original
collar domain. BCG02 gives the
actual/model value and differential
comparison on ALL $D_a$.
Since $R_a<2r_\partial<1$,
\[
 \|D\Phi_{a,b}\|\leq P,\qquad
 \|D^2\Phi_{a,b}\|\leq R_aP<P.             \tag{BCG03.b}
\]
Indeed differentiating $h_b$ cancels
the prefactor $R_a^{-1}$ once and
leaves a factor $R_a$ after the
second differentiation. The possibly
large constant value $1/R_a$ never
enters these derivative bounds.

More explicitly,
$|\eta_b-h_b(\eta_a)|<R_a\theta$
and the same normalized derivative
error is less than $R_a\theta$.
Using $\|D\eta_a\|\leq2$, the chain
rule and the global bounds on
$\mathcal{B}',\mathcal{B}''$ give normalized
value error at most $P\theta$ and
differential error at most
$(P+2R_aP)\theta\leq3P\theta$.
For an unlisted block the ACTUAL
block and its derivative vanish
throughout $D_a$, as does its model.
Thus choose $\theta$ to make this
extra error less than half the
prescribed rough error and ask the
original interior producer for
the other half. Square summation
with one extra block proves the
same early modulus formulas with
augmented $N,P$. The own reference
block is still $(z,1)$ and the
first model scale coordinate
is still constant1.

The whole-chart enclosures and
exact coordinate sections are the
original ones. Original-marker
division therefore localizes
EVERY near image preimage to those
same enclosures, even in the larger
ambient target. The proofs of
TCP06, EGP07 and SGP05--SGP06
now apply to these actual enlarged
graphs: their retained identity
component gives the positive rank
margin, and their derivative error
gives the normal estimate for the
SAME planes. Their section and
Hessian arguments give BOTH
truncated cloud coverage directions.
Their radius comparison uses the
original retained markers and
scale, so is unchanged by adding
coordinates. CFS27 deletes only
the specified small interior
blocks; (BCG03.a) is never pruned
after planes have been chosen.

BCG01 supplies the early global
derivative bound. The source cutoff
formulas of CFS31 use only the
original interior and weak-edge
blocks; their smooth supports,
exact plateaux and derivative
estimates remain valid. Their
constants can be bounded by the
same enlarged early $N,P,L_0$.
Choose the three targets, graph
moduli and cloud scales in GAF01's
reverse stage order, and then the
reference qualities in BCG02's
permitted order. After all fixed
normalized tests and $V$, choose
the positive $r_\partial$ below
the finitely many physical upper
bounds and increase the common
sequence tail. No cloud scale
depends on $r_\partial$.

CFS28--CFS31 still verify every
contributing-center and plane
condition for the ORIGINAL
zero-marker assertions: their
pruning and marker calculations
are unchanged, and all boundary
blocks are merely retained
coordinates. CFS20 and CFS18
therefore perform the actual
three-stage postcompositions on
neighborhoods of the successive
images. CGP07--CGP08 give their
one-sheet marked bases and
injectivity across patches.
GAF04--GAF06 give exact original
full markers and global full-
preimage localization. Original
circle and slim fiber homotopies
have their entire compact traces
in the original buffered coordinate
domains. BCP04 puts these in
$\operatorname{int}M$, so FC34
applies there; extend its compactly
supported flow by identity to
the boundary. This verifies the
domain change in GAF07's proper
whole-fiber conclusions, rather
than assuming external-boundary
fibers have the same type.

Finally every first model still
has scale coordinate1. Hence
every selected normal projector
is the identity on that scalar
coordinate, with no mixed block.
EDP01's normalized-mean formula
and its subtraction of the
COMMON reference radius in
the partition derivative apply
unchanged, with the augmented
early constants. Its resulting
$C_\rho$ is still chosen before
$\Lambda$ and $\Delta$; impose
$C_\rho\Delta\Lambda<10^{-6}$
at the later scale choice.
This does not impose
$\Delta c_3\ll1$ or
$V\Lambda\ll1$.
\end{proof}

\begin{lemma}[Global physical boundary-block isolation; BCG04]
\label{lem:fibration-global-boundary-block-isolation}
For every boundary index $b$, every intermediate
map $f$ in BCG03 and every $p\in M$,
\[
 \begin{gathered}
 \rho(p)>20r_\partial\ \Longrightarrow\
                 J_bf(p)=J_bF(p)=0,\\
 |J_b(f-F)(p)|<\epsilon_\partial,\qquad
 \epsilon_\partial=20c_3r_\partial .
 \end{gathered}
 \tag{BCG04.a}
\]
where $J_b$ projects onto the WHOLE
physical boundary block. The error
bound also holds along the segment
from $F$ to $E$. In particular it
requires no global upper bound on
$\rho$.
\end{lemma}
\begin{proof}
A nonzero original block implies
$20<\eta_b<90$ and $\rho<r_\partial$,
so the original block is zero in
the asserted large-scale region.
Induct over the actual adjustments.
On a stage cutoff's closed support,
BCG03 puts $x=\pi_jF(p)$ in the
ORIGINAL stage core. For EVERY
smoothing center $y$ whose closed
$80br_y$ support meets
$B(x,8br_x)$, CFS28's scale
calculation (AM) gives, for its
actual chosen model reference $a$,
\[
 R_a\geq\frac{108}{625}\rho(p)
             >\frac{2160}{625}r_\partial
             >2r_\partial .
\]
BCG01 shows that boundary support
$b$ cannot meet $D_a$. Its model
block in (BCG03.a) is thus absent,
the actual core preimage used at
$y$ lies in $D_a$, and
$J_by=0$, $L_y\subset\ker J_b$.
This checks the entire contributor
list, including distinct radius
and model preimage choices.
CFS24 and CFS16 give zero block
for the projected output at the
perturbed input. Its blend with
the inductively zero input is
zero. Outside the stage cutoff
there is no change. Every $Q_j$
retains the block, proving the
first assertion through all stages.

If the displayed large-scale
inequality fails, BCG03 gives
$|J_b(f-F)|<c_3\rho(p)\leq
20c_3r_\partial$; otherwise
the difference is zero.
Convexity gives the segment
assertion. This is physical
smallness proved from actual
support lists, not an estimate
which tacitly assumes $\rho$
small at an arbitrary preimage.
\end{proof}

\begin{lemma}[Exact boundary marker at every contributing plane; BCG05]
\label{lem:fibration-actual-boundary-full-markers}
Require $r_\partial<10^{-4}$ and
$\Sigma_j\leq\varepsilon_j/10000$.
At every source point with
$32\leq\eta_b(p)\leq78$, the
boundary scalar marker of EVERY
intermediate map is exactly1.
At each active stage this follows
from the exact marker1 and
zero marker derivative of EVERY
actual contributing center and
plane. In particular the final
marker is locally constant1
on $32<\eta_b<78$.
\end{lemma}
\begin{proof}
The original marker is1 on the
given band. If the current stage
cutoff vanishes, the map does not
change. Otherwise use the original
core point $x=\pi_jF(p)$.
Write $b_{\rm sm}=\varepsilon_j^{-1}$
for the smoothing buffer, to
distinguish it from the boundary
index. CFS07 and CFS26 give
$r_y\leq5r_x/3$ and
$r_x\leq(5/3)\Sigma_j\rho(p)$
for every contributing center.
Consequently
\[
 |y-x|<
 250b_{\rm sm}\Sigma_j\rho(p)
 \leq.025\rho(p)<r_\partial.                    \tag{BCG05.a}
\]
This includes every center whose
closed support meets the buffered
ball, not only centers active at
the input itself.

Let $q_y$ be the actual original
core preimage used to choose its
model reference. At $x$ the
physical boundary block is
$(\eta_b(p),1)$.
The marker bound in (BCG05.a)
makes $\zeta_b(q_y)>1-r_\partial$.
Thus its ORIGINAL coordinate is
defined, and marker division gives
\[
 |\eta_b(q_y)-\eta_b(p)|
 \leq\frac{79r_\partial}{1-r_\partial}<.01 .
\]
Therefore $31<\eta_b(q_y)<79$
and $\zeta_b(q_y)=1$ exactly.
This proves the actual marker
of $y$ is1.

The point $q_y$ lies in its
chosen reference's $D_a$.
Hence $b$ is in that ACTUAL
whole support list and
$R_a<2r_\partial$.
At the model parameter
$z_y=\eta_a(q_y)$,
BCG02 gives
$|h_b(z_y)-\eta_b(q_y)|<R_a\theta$.
The model height is strictly
between30 and80. Its boundary
marker is consequently locally
constant $1/R_a$, and the
physical model plane lies in
the scalar marker kernel.
CFS27 retains this block.
GAF03 with constant $c=1$
now gives marker1 for the
spectral projection at the
perturbed input, which lies
in $B(x,r_x)$ by CFS16.
Induction starts from marker1,
and each blend of marker1
with marker1 preserves it.
This proves both exact claims
without assuming that a
nearby image point has a
preselected local preimage.
\end{proof}

\begin{theorem}[The actual adjusted torus collar and its internal frontier; BCG06]
\label{thm:fibration-actual-adjusted-boundary-core}
Take $\epsilon_\partial<10^{-6}$
and $c_3<10^{-5}$ in BCG03--BCG05.
Write $(u_b,v_b)=J_bE$ in PHYSICAL
units and define
\[
 C_b=N_{35}(\partial_bM)\ \cup\
        E^{-1}\{v_b\geq.9,\ u_b\leq40v_b\}.
 \tag{BCG06.a}
\]
Then $C_b$ is the WHOLE compact smooth
inner collar up to a single smooth
torus, and is smoothly diffeomorphic
to $T^2\times[0,1]$, retaining the
original external boundary label.
Its internal frontier
$H_b=\partial_M C_b$ has the GLOBAL
description
\[
 H_b=E^{-1}\{v_b\geq.9,\ u_b=40v_b\}.
 \tag{BCG06.b}
\]
It lies in $39.99<\eta_b<40.01$,
where $v_b=1$ exactly and
$u_b-40$ has nonzero differential.
The full manifold boundary of $C_b$
is $\partial_bM\sqcup H_b$; equation
(BCG06.b) describes only its relative
frontier in $M$.
The $C_b$ are pairwise disjoint and
disjoint from all the actual zero
domains of ZSP02. Replacing $\geq.9$
by $>.9$ in (BCG06.a) or (BCG06.b)
does not change these sets.
\end{theorem}
\begin{proof}
BCG04 gives global errors less
than $\epsilon=\epsilon_\partial$
in both scalar components.
If $v_b(E(p))\geq.9$, then its
original marker is greater than
$.9-\epsilon$, so $p$ is in
the original support and its
height is defined. Marker
division gives
\[
 \begin{split}
 u_b/v_b\leq40
    &\ \Longrightarrow\
 \eta_b\leq40+\frac{41\epsilon}{.9-\epsilon},\\
 u_b/v_b=40
    &\ \Longrightarrow\
 |\eta_b-40|\leq
           \frac{41\epsilon}{.9-\epsilon}<.001 .
 \end{split}                                             \tag{BCG06.c}
\]
These bounds apply to EVERY
point of the full preimage,
including any possible remote
sheet; they assume no prior
distance localization.

A path of length below35
from $\partial_bM$ stays
in the original collar
below $z_b<35.35$ by first
exit. Thus $N_{35}(\partial_bM)$
is below $\eta_b<35.36$
on its part in the height
band. Conversely the entire
original inner region up
to $\eta_b=34$ is in this
neighborhood: its vertical
path has length at most
$1.01(34+\delta_n)<35$.
On $34\leq\eta_b\leq38$
the original and adjusted
markers are1 by BCG05,
and $u_b<38+\epsilon<40$.
So there is no gap between
the neighborhood branch
and the marked branch.

On the compact band
$38\leq\eta_b\leq42$
we have exactly $v_b=1$,
and BCG03--BCG04 give
\[
 |u_b-\eta_b|<\epsilon,\qquad
 |d(u_b-\eta_b)|_g<c_3 .
\]
Along BCP01's smooth
height field
$X_b=\nabla\eta_b/|d\eta_b|^2$,
\[
 du_b(X_b)>1-1.02c_3>.99 .
\]
Every flow line across this
band has exactly one point
where $u_b=40$, between
heights $39.99$ and $40.01$.
The implicit function theorem
makes these points a smooth
graph over the original
level40 torus. Both ends
of every line are on the
specified opposite sides.

By (BCG06.c), no other
internal boundary or remote
marked component is possible.
Below height34 the neighborhood
branch fills the whole inner
collar; from34 to38 the
marked branch fills it; from38
onward its upper boundary is
exactly this graph. Every
point in the marked branch
below34 is already in the
neighborhood branch. Thus
(BCG06.a)'s set is precisely
the WHOLE original inner
subgraph, including its
external boundary. This also
proves its compactness and
smooth-boundary structure.

For an explicit type
identification, move the
original level40 torus to
this graph by its smooth
height-flow product, using
an interpolation supported
in $38<\eta_b<42$.
The displacement in the
height variable is less
than $.01$, so a smooth
cutoff with bounded derivative
gives a strictly increasing
height map; its inverse is
smooth by the implicit
function theorem. Extend
by identity outside this
band. This is an ambient
diffeomorphism fixed near
the external boundary.
BCP01's graph-and-LFR04
argument at level40 gives the
original inner subgraph smooth type
$T^2\times[0,1]$, so the
same holds for $C_b$.

The frontier equation follows
in both directions from the
global equality case in
(BCG06.c) and the graph
description. Its differential
is nonzero by the displayed
height derivative, and its
marker is1 by BCG05.
Each collar core lies inside
$B_b$, and BCP03--BCP05
separate these from one
another and from the selected
zero balls containing ZSP02's
domains. Finally a point of
the marked branch with
$v_b=.9$ cannot have
$\eta_b\geq32$, since the
preceding bound puts it
below41 and BCG05 would
give $v_b=1$. It must
lie below32, already in
$N_{35}(\partial_bM)$.
The equality frontier always
has marker1. These facts
prove the strict-marker
equivalence as well.
\end{proof}

\begin{proposition}[Retained internal torus faces and relative removal; BCG07]
\label{prop:fibration-boundary-frontiers-saturated}
The actual circle, slim and edge whole-fiber
restrictions and compatible vertical faces
of GAF07/EDP01--EDP06 hold for BCG03's
boundary construction on their original
interior domains. Put
\[
 Z=\bigcup_i Z_i,\qquad C=\bigcup_b C_b,\qquad
 M_1=M\setminus\operatorname{int}_M(Z\cup C).
 \tag{BCG07.a}
\]
Then $M_1$ is a compact smooth manifold
with boundary, contained in
$\{d(p,\partial M)\geq35\}$; its boundary
consists only of the original internal
zero faces and the internal cusp tori
$H_b$. No external boundary component
remains as an isolated face.
Every WHOLE fiber of $\pi_jE$ meeting
$H_b$ lies in $H_b$. For the proper
circle and slim restrictions,
$M_1\cap X_j=f_j^{-1}(D_j)$ is saturated
over a closed relative base domain with
smooth boundary, whose cusp equation
is $u_b/v_b-40=0$.
If $H_b$ meets $X_3$, the ENTIRE connected
torus $H_b$ is one whole slim fiber.
\end{proposition}
\begin{proof}
The proper restrictions use the
original charts and compact
source enclosures in BCP04.
Their homotopy traces lie in
$\operatorname{int}M$ by the
global marker localizations.
Consequently the GAF07 and
EDP whole-fiber proofs, including
the two target strata for disk
fibers, apply on that open
smooth carrier with the SAME
enclosures. BCG03 verified
the adjusted-scale derivative
needed in the full height
quotient. The original full
collar circle family and
its weak-border witnesses
remain valid on their
interior buffers. BCG01's
extra edge-center margin
makes EVERY such circle-collar
point eligible for the actual
selected circle cover. EDP06's
open-and-closed argument
therefore identifies each
whole disk-boundary circle
with its whole circle fiber.
Nothing here requires the
external boundary to be
a regular fiber.

Likewise ZSP01--ZSP02's
zero-block isolation still
uses the actual old support
lists and retained blocks.
Its entire radial level
trace is in the original
interior zero ball. Apply
its ambient isotopy there
and extend by identity
near $\partial M$.
Thus these are the actual
$Z_i$, not separately chosen
replacement cores.

BCG06 gives disjoint compact
smooth $C_b$ with disjoint
internal regular frontiers,
and separates them from
all zero domains. Removing
RELATIVE interiors in $M$
removes each external
boundary together with its
neighborhood. Standard
sublevel charts then give
exactly the smooth internal
boundary components asserted
for $M_1$. Each $C_b$
contains $N_{35}(\partial_bM)$
in its relative interior,
so all remaining points have
distance at least35 from
the original boundary.
Compactness follows since
$M_1$ is closed in compact $M$.

Every $\pi_j$ retains the
WHOLE physical boundary
block. Equality of $\pi_jE$
therefore preserves the
GLOBAL equation (BCG06.b);
a whole fiber meeting it
is wholly in that frontier.
For the proper circle and
slim restrictions the
connected-fiber argument
of ZSP03 now applies to
both sorts of disjoint
internal face: a fiber
avoiding all frontiers
cannot lie on both sides.
Their complement restriction
is thus saturated. Properness
makes its base image closed.
Near a cusp frontier the
base function $u_b/v_b-40$
pulls back to the actual
regular defining function
in BCG06. The submersion
forces its base differential
to be nonzero, and gives
the local base domain on
the required side. Distinct
cusp and zero faces do
not intersect.

If a slim fiber meets
$H_b$, it is a compact
connected surface inside
the same two-dimensional
smooth connected torus.
The regularity just proved
makes its inclusion locally
open; compactness makes
it closed. Hence it is
ALL of $H_b$. As in
ZSP03, this also puts the
whole torus in $X_3$ and
identifies its one base
boundary point. These are
actual face equalities,
not only tangent-kernel
inclusions.
\end{proof}

\paragraph{Remaining assembly and parameter binding.}
The augmented graphs, every-contributor
boundary marker control and actual torus
cores are now constructed from the
original data. The final boundary
decomposition must still use these
new internal torus faces in the
compact slim-base choice and in
FDC's face partition. Its Euler
count and torus gluings, the full
boundary parameter register and
static threshold conclusion are
the next obligations. The requested
ONE final full-goal mathematical
audit has not yet occurred.

\subsection{Actual boundary decomposition and graph assembly}
\label{sec:fibration-actual-boundary-assembly}
Retain the SAME boundary assignment, map $E$, local packets and
cores of BCG01--BCG07. In particular the edge family is selected
at distance greater than20; the other centers at distance greater
than10. This subsection constructs the remaining compact pieces
and their actual face partition. Its conclusion still requires
one simultaneous boundary parameter binding, supplied separately.
The retained topological recognition library is FC41.

\begin{theorem}[The compact slim restriction with retained cusp faces; BCF01]
\label{thm:fibration-boundary-compact-slim-restriction}
For $M_1$ in (BCG07.a), let
\[
 D_3=f_3(M_1\cap X_3)\subset B_3.
\]
One can choose a compact one-dimensional submanifold with boundary
$K_3\subset B_3$ containing in its interior BOTH the image of
all original closed $3.5\cdot10^5\Delta$ slim slabs and the finite
set $\partial D_3$, with $\partial K_3\cap\partial D_3=\varnothing$.
Then
\[
 S=f_3^{-1}(K_3\cap D_3),\qquad
 M_2=M_1\setminus\operatorname{int}_{M_1}S                 \tag{BCF01.a}
\]
are compact smooth domains. The first is the actual proper whole
$S^2$ or $T^2$ bundle over $K_3\cap D_3$.
Their common faces and the remaining boundary satisfy
\[
 \begin{split}
 M&=Z\cup C\cup S\cup M_2,\\
 S\cap M_2&=\partial S\setminus\partial M_1,\\
 \partial M_2&=(\partial M_1\setminus S)
                \sqcup(\partial S\setminus\partial M_1).
 \end{split}                                               \tag{BCF01.b}
\]
All these internal faces have collars. Every cusp torus meeting
$X_3$ is an ENTIRE torus boundary fiber of $S$; its full inward
collar in $M_1$ belongs to $S$, and it does not remain in $M_2$.
All points of $M_2$ have distance at least35 from $\partial M$.
\end{theorem}
\begin{proof}
BCG07 and ZSP03 identify $D_3$ as a closed relative smooth
domain in the actual one-dimensional base. Every boundary point
comes from an entire connected zero face or cusp torus; there
are finitely many such faces. Conversely their regular defining
functions give the corresponding base boundary points. Thus
$\partial D_3$ is finite. The original closed slim slabs are
compact by their original proper bundle and lie in $X_3$;
their image and this finite set form one compact subset of $B_3$.
Use finitely many relatively compact coordinate intervals,
enlarge them slightly with endpoints avoiding $\partial D_3$,
and take their union, retaining any whole circle component.
This is exactly ZSP04's construction with the finite face set
also included in the set to be covered. In particular no new
quantitative tolerance is chosen from these endpoints.

Intersecting with $D_3$ gives a compact one-manifold whose
boundary is
\[
 (\partial K_3\cap\operatorname{int}D_3)
                  \sqcup(\operatorname{int}K_3\cap\partial D_3).
\]
Restrict the actual proper bundle. Its product charts prove
compactness, the smooth boundary, and the whole-fiber identities.
At any shared old face $K_3$ imposes no restriction on a full
base neighborhood, so its entire inward collar is retained.
BCG07 identifies a cusp torus meeting $X_3$ with its whole
connected torus fiber. Our extra finite-set inclusion puts this
fiber in $S$, even if it was outside the required original slim
slabs. This establishes the assertion for every such cusp torus.

The three disjoint local cases in ZSP05's proof now apply:
a retained old zero or cusp face is removed with its relative
collar; a new slim endpoint is an interior regular level; and
an old face not meeting $S$ is unchanged. They give every
identity in (BCF01.b) and smoothness of $M_2$ without an isolated
old surface in the complement. Closedness in compact $M_1$
gives compactness. Disjoint ambient interiors follow in the
same charts. The distance assertion follows from $M_2\subset M_1$.
On each base interval, bundle transport gives the actual
product with its endpoint fibers. A closed zero or slim component
would be open and closed in connected $M$ and lie in its interior;
in the present nonempty-boundary case it cannot occur.
\end{proof}

\begin{theorem}[The compact edge and circle pieces in the boundary carrier; BCF02]
\label{thm:fibration-boundary-edge-circle-decomposition}
With the original edge region $X_2$ and map $f_2$ of BCG07, set
\[
 P=M_2\cap X_2,\qquad 
 R=M_2\setminus\operatorname{int}_{M_2}P .                 \tag{BCF02.a}
\]
Then $P$ is a compact proper whole disk bundle over a compact
one-manifold, and $R\subset X_1$ is a compact proper whole
circle bundle over a compact surface with corners of depth at
most two. Their actual common face is the full vertical face
$V_e=M_2\cap\partial X_2$. Their other faces are
\[
 H_e=\partial M_2\cap X_2,\qquad
 R\cap\partial M_2=
       \partial M_2\setminus\operatorname{int}_{\partial M_2}H_e.
                                                               \tag{BCF02.b}
\]
Every horizontal disk lies in a single zero, slim or retained
cusp boundary component. Its entire boundary circle is the SAME
whole circle fiber of $R$. All horizontal/vertical meetings have
independent defining differentials. No face is lost at an open
edge-base endpoint or at the original external boundary.
\end{theorem}
\begin{proof}
First port the horizontal-face argument, including its new
hypothesis. Each component of $\partial M_2$ is an entire zero
face, a retained cusp torus, or a new whole slim endpoint fiber,
by (BCF01.b). A whole disk meeting the first is contained in it
by the retained zero equation, and one meeting the second is
contained in it by the GLOBAL equation (BCG06.b). For the third,
constancy of $\pi_2E$ makes $\pi_3E$ constant, so the whole disk
lies in that same whole slim fiber. Connectedness then gives
saturation of $P=f_2^{-1}(A_2)$; properness makes $A_2$ relatively
closed in $B_2$.

Each defining equation descends smoothly to the edge base.
For a cusp face use $u_b/v_b-40$. Its pullback is the actual
nonzero defining differential of BCG06. For a slim face use
its endpoint coordinate, on a neighborhood of the compact whole
disk furnished by properness; for a zero face use its retained
ratio. Thus the descended function has nonzero differential.
In a one-dimensional edge coordinate $g_i$ it is $b(g_i)$ with
$b'\ne0$. EDP04's independence of $dg_i,dT$ proves actual
horizontal/vertical transversality. This is the full EDP05
argument for all three types of faces, not just saturation.
BCG01 and BCG07 already guarantee the entire vertical circle
cover and exact equality of the whole boundary circles.

Here is the eligibility check for compactness at an open base
endpoint. In FDC02's limiting argument the limiting point $q$
lies in $M_2$, so $d(q,\partial M)\geq35$. Its witnessing
selected edge center $p_i$ has distance greater than20.
Therefore the zero tenth-cover and slim-cover exclusions in
FDC01 are both eligible at $p_i$ and prove that it is an actual
nonslim one-stratum point, not merely a strong-edge center.
The larger choice of $K_3$ still contains all the original
closed slim slabs in its interior, as required in that exclusion.
All the actual weak-border witnesses and comparisons in FDC01
are on the unchanged original domains.

It remains to check the strong witness used to find a replacement
selected edge, since this witness must now have distance greater
than20. The same estimates give
\[
 d(q,p_i)<6\Delta R_i,\qquad 
 d(q,p)<(4.2\Delta+1.01)R_i<5\Delta R_i,\qquad 
 R_i<2\rho(q).
\]
The last inequality follows from slow variation, without assuming
that $q$ already lies in the open edge base. Require on the
uniform tail that $10\Delta\rho(q)<d(q,\partial M)/3$ for all
$q$ with distance at least35; BCP04.a supplies this fixed-buffer
bound. Then
\[
 d(p,\partial M)>\tfrac23d(q,\partial M)
                    \geq70/3>20.
\]
Maximality of the SAME restricted edge family therefore applies
to this witness. FDC01's actual comparisons produce the strict
replacement marker and $|u_j|/R_j<3\Delta$. FDC02's subsequence
argument now proves $P$ closed in compact $M$ and its base
compact, with no additional open-end face. Its proper disk
bundle and complete face charts were proved in the preceding
paragraphs. A compact one-manifold has only finitely many
interval or circle components.

For the relative complement in (BCF02.a), the local charts
are exactly a domain, half-space or quadrant
$h\geq0,T\geq4\Delta$. Its horizontal surface away from
the vertical face is removed with a full relative neighborhood.
These charts give (BCF02.b), compactness of $R$, and the common
face $V_e$. Its coverage is also actual: every point under
consideration is at distance at least35. The BCG01 restricted
four-family cover therefore excludes zero points by their
interior cores, slim points by the interior of $K_3$, and
nonslim one-stratum points by a strict interior edge chart.
There are no three-stratum points. All points of $R$ are thus
two-stratum and eligible for the ORIGINAL selected circle cover,
which puts them in $X_1$.

Every face of $M_2$ is saturated by whole circle fibers, using
the same global zero/cusp equations or the whole slim fiber.
Membership in $X_2$ and the strict height inequality specifying
its relative interior are functions of $E$. Hence $R$ is
saturated by the SAME whole connected circle fibers. Its base
is compact. The cusp equation, like the zero or slim equation,
descends with nonzero differential to this two-dimensional
base. At a vertical meeting its horizontal equation is $b(g_i)$,
and $(g_i,T)$ are actual base coordinates by EDP06. They give
independent base equations and the entire half-space/quadrant
charts. Restrict the proper circle-bundle charts to these
compact domains. This proves every asserted base, fiber and
face property, including the depth bound and actual inclusions.
\end{proof}

\begin{proposition}[Actual face partition and exclusion of cusp disk faces; BCF03]
\label{prop:fibration-boundary-torus-face-partition}
The five domains $Z,C,S,P,R$ cover $M$, with disjoint ambient
interiors and the actual face identities above. Every connected
component $Y$ of $\partial M_2$ has the exact partition
\[
 Y=A_Y\cup B_Y,\qquad
 A_Y=Y\cap P,\quad B_Y=Y\cap R,\qquad 
 A_Y\cap B_Y=\partial A_Y=\partial B_Y .                 \tag{BCF03.a}
\]
Here $A_Y$ is a finite disjoint union of whole horizontal disks,
$B_Y$ is a compact circle-bundle surface, and their common
circles are the actual fiber circles. Consequently a torus
component has no horizontal disks; a sphere component has
exactly two. In particular every cusp core $C_b$ attaches
along its ENTIRE internal torus $H_b$ either to a torus slim
piece or to one boundary component of the circle piece $R$.
No cusp core enters the ball-and-disk-handle incidence graph.
\end{proposition}
\begin{proof}
The union and relative-interior statements follow by combining
(BCG07.a), (BCF01.b) and (BCF02.a)--(BCF02.b).
A horizontal disk is one whole fiber over an endpoint of the
compact edge base. The finite base list gives finitely many
such disks, disjoint even at their boundaries because distinct
fibers are disjoint. BCF02 places each disk in one connected
component $Y$. The remaining part of $Y$ is exactly $Y\cap R$.
Its regular horizontal equation descends to the circle base;
at the vertical meetings the independent base equations give
precisely its boundary endpoints. Thus this is a compact circle
bundle over a compact one-manifold with boundary. Its boundary
circles are exactly the disk boundaries, by the WHOLE-circle
identity in BCG07. There is no missing limiting face, by the
compactness part of BCF02. These are all the embedded partition
hypotheses of FC40, which now gives the asserted counts.

If $H_b$ meets $X_3$, BCF01 already identifies it with a whole
torus boundary fiber of $S$, and its relative collar is removed
from $M_2$. Otherwise it is disjoint from $S$ and remains an
entire component of $\partial M_2$. Apply the partition just
proved to this connected torus. The disk count is zero.
An intersection with $P$ would contain a WHOLE disk fiber in
$H_b$ and hence a horizontal disk, even if the original point
was on the vertical edge boundary. Therefore $H_b\cap P$ is
empty, and $H_b\subset R$. The base regularity in BCF02 and
absence of vertical corners make it a whole smooth boundary
component of $R$. Connectedness and the local boundary charts
rule out a proper subset of another component. This proves the
stated dichotomy without assuming that cusp points are slim or
that the original cusp band contains no edge charts.
\end{proof}

\begin{theorem}[The actual boundary graph presentation; BCF04]
\label{thm:fibration-actual-boundary-graph-presentation}
Under the retained foundations, in particular FC41, the SAME
boundary assignment produces a finite KL Definition1.2 graph
presentation on the original compact connected orientable carrier
$M$. Its external boundary components retain their original
labels. Its internal torus gluings, bundle pieces and corner
roundings come from the actual constructed domains and inclusions.
Together with BCP03's whole-product alternative, this covers
both geometric cases of the nearly cuspidal boundary construction.
This is a static conclusion for an admissible assignment; the
uniform boundary parameter register and static threshold remain
to be bound. No essentiality, flow application or accepted Lean
binding is asserted by this theorem.
\end{theorem}
\begin{proof}
In the whole-product case, $M=T^2\times I$ is itself a circle
bundle over an annulus; use its actual labeled boundary
identification. In the other case use BCF01--BCF03.
There are finitely many compact pieces and base components.
All zero types remain those of ZSP02, including the corrected
three-ball alternative. Slim components over intervals have
the actual $S^2$ or $T^2$ product type; there is no closed
zero or slim component in this nonempty-boundary carrier.

Group the torus-boundary zero pieces, torus slim cylinders and
all cusp cores. Each is already a graph piece: the first uses
FC41's solid-torus or twisted Klein-bottle recognition, and
each cylinder is a circle bundle over an annulus. Common faces
within this group are ENTIRE tori. BCF03's zero disk count
shows that every other internal boundary of this group is a
whole torus boundary of $R$. Finite torus gluing in FC41 gives
a graph presentation of every grouped component. Its external
boundary components, which occur only on the cusp cores, remain
unglued and keep their original labels. Remove the RELATIVE
interior of this group when considering its complement; its
external boundary and collar then leave no isolated remnants.
If it is the entire connected carrier the proof is complete.

For the other components follow FC42's finite sphere-and-handle
argument with the actual face partition now established.
Replace punctured projective three-spaces by balls, recording
their connected-sum reconstruction. Cut and cap each remaining
sphere cylinder, recording whether reconstruction joins two
components or adds $S^1\times S^2$. These operations are supported
in the interior, away from the torus group and the original
external boundary. Closed components produced during this step
are handled by the same FC41 recognition inputs. Remove a disk
bundle over a circle as a solid torus, retaining its entire
torus attachment. All remaining vertices are balls with exactly
two actual disk faces by BCF03 and FC40, and the remaining
handles are the actual product disk bundles over intervals.
The finite incidence graph has degree two at every vertex,
including loops counted twice. It is a disjoint union of cycles;
FC41 identifies each corresponding embedded collared union with
a solid torus. There are no additional vertices or ends coming
from a cusp core, by BCF03.

The complement is exactly the actual circle bundle $R$.
For its finitely many base corners, use FDC04's compatible
rounding: in disjoint base rectangles round the quadrant axes
by smooth arcs, and lift by the actual circle-bundle product
charts. Use the SAME arc for the complementary adjacent piece.
This gives identical attaching subsets and smooth collared
torus boundaries. A cusp internal torus has no such corner,
and all original external boundaries lie on the untouched cusp
cores. Thus these local roundings preserve their labels and
neighborhoods. No fiber matching across an arbitrary torus
attaching diffeomorphism is presumed: FC41 supplies the graph
class closure with the actual collared gluing maps.

Glue the resulting solid tori, $R$, and the retained torus group
along their whole internal torus boundaries, and reverse the
finite recorded sphere operations. FC41 gives the finite graph
presentation of this same carrier with its original external
boundary. The source's instruction after16.6 to treat a cusp
core like a torus zero piece ``without a core'' is realized here
by an actual torus-product piece and its one internal attaching
face, with its other face retained externally.
\end{proof}

\paragraph{Scope after boundary assembly.}
BCF01--BCF04 finish the actual geometric decomposition and
static graph assembly for one admissible boundary assignment.
The remaining quantitative task is to bind ALL BSA, BCP, BCG
and BCF choices to one uniform order and tail, and then apply
the boundary counterexample reduction to obtain a single static
threshold. The requested final full-goal mathematical audit
follows that task. Lean and accepted-PC verification remain
separate from these written arguments.

\section{Parameter order and cross-chapter interfaces}
\label{sec:fibration-parameter-audit}
\begin{lemma}[Finite acyclic parameter selection; FC46]
\label{lem:fibration-parameter-selection}
Suppose finitely many parameters are ordered in an acyclic directed
graph. At each node, after predecessor values are fixed, require
finitely many strict positive upper bounds and no lower bound
other than positivity, OR finitely many finite lower bounds and
no upper bound. Then all constraints have a simultaneous solution.
For a node with both types require a specified nonempty admissible
interval for every allowed predecessor tuple.
\end{lemma}
\begin{proof}
Choose parameters in a topological order. A positive number below
the minimum of the finitely many positive upper bounds exists;
a number above the maximum of finitely many finite lower bounds
exists. In the mixed case use the supplied nonempty interval.
Continue through the finite list. A cycle or a mixed interval
without its nonemptiness proof is not covered.
\end{proof}

\begin{foundation}[The complete source order and binding gate; FC44]
\label{found:fibration-parameter-order}
KL (15.4), printed87/PDF82, orders choices as follows; braces
indicate a layer, not mutually dependent choices. The source's
multiplicity $\mathcal M$ and three-stratum exclusion $\beta_3$
are fixed first:
\begin{align*}
 \{\mathcal M,\beta_3\}&\prec\{c_{\rm slim},\Omega_i,\Omega'_i\}
 \prec\Gamma_3\prec\{\Sigma_3,\Xi_3\}\prec c_{\rm edge}
 \prec\Gamma_2\prec\{\Sigma_2,\Xi_2\},\\
 &\prec c_{2\text{-stratum}}\prec\Gamma_1
 \prec\{\Sigma_1,\Xi_1\}\prec\varsigma_{2\text{-stratum}}
 \prec\beta_2\prec\Delta,\\
 &\prec\{\varsigma_{\rm edge},\varsigma_{E'},\varsigma_{\rm slim}\}
 \prec\varsigma_{0\text{-stratum}}
 \prec\{\beta_{E'},\sigma_{E'}\}\prec\sigma_E,\\
 &\prec\{\sigma,\Lambda\}\prec\bar w\prec w'
 \prec\beta_E\prec\beta_1\prec\{\Upsilon_0,\delta_0\}
 \prec\Upsilon'_0.
\end{align*}
The boundary route appends $r_\partial$.
Deliver a register of every consumer inequality and its allowed
predecessors, including positivity/finiteness of each bound.
Only then apply FC46. A printed ordering is not by itself a proof
that newly expanded quantitative inequalities respect it.
\end{foundation}
CFS20 now supplies explicit tube and perturbation bounds in this order,
including $c_2\le3\Sigma_3/10$ and $c_1\le3\Sigma_2/10$.
Its cutoff hypotheses still need uniform producers.
The implementation register must keep these specific tests:
\begin{itemize}[leftmargin=*]
\item FC08's multiplicity is independent of $\Delta$; choose the
curvature comparison normalization $\kappa=1/\Delta$ and retain
all whole-ball requirements. Constants fixed before $\Delta$
cannot secretly grow with it.
\item FC24--FC25 choose graph modulus and cloud scales before the
later witness errors. A Hessian bound depending on a later
$\Delta$ is not enough unless that dependence is canceled by the
scaled cutoff estimates. The source's $\Omega_i(\mathcal M)$
uniformity remains an explicit geometric binding requirement.
\item FC22 short tests avoid asking $\beta_2^{-1}$ to dominate a
later $10^6\Delta$. The long domain belongs to the chart whose
accuracy is chosen later. Record all actual rescaled thresholds.
\item FC18 and FC26 require $\Lambda\Delta$ small, and slim
supports require its $10^6$ multiple small. This is compatible
with choosing $\Lambda$ after $\Delta$.
\item FC28 smoothing tolerances and FC30 cutoff derivative bounds
must fit FC31's $abL+dL+e$ budget at each of three stages.
The earliest errors are chosen small enough for later stages;
$c_{\rm adjust}=c_{\rm slim}$ is the final tolerance.
\item The compact isolation and face transversality tolerances are
chosen before the perturbations and before the local witnesses.
Any finite-$C^K$ to smooth-category adapter remains LC81's gate.
\end{itemize}
Thus the order and every type of constraint are now specified;
the complete numeric/source binding register is still a qualified
producer obligation. No global simultaneous-existence theorem is
claimed from a collection of independently satisfiable local bounds.

\subsection{Binding the entire closed construction to one assignment}
\label{sec:fibration-closed-parameter-binding}
This register concerns the CURRENT CLOSED route, using the stronger
actual-cloud hypotheses that replace the false unrestricted FC28.
The unused FC19 common-domain variant and CFS19's arbitrary
carrier-dependent cutoff tolerance are not imposed. Boundary FC43
and the flow-application rows CP17--CP18 retain their separate scope.
Write $\omega_3$ for Euclidean unit-ball volume, to distinguish it
from the third adjustment tolerance $c_3$.

\begin{proposition}[Closed register and simultaneous admissibility; PBR01]
\label{prop:fibration-closed-register-admissibility}
Fix $K\ge10$, the standing derivative function $A'$ and the retained
foundations. All parameter requirements of the actual closed route
LPA--SGP--EGP--TCP--GAF--ZSP--EDP--FDC admit one assignment.
The following register is ordered by dependency. A named producer's
qualitative threshold means its proved POSITIVE threshold with the
earlier arguments specified below; its geometric conclusions are
not inserted as additional assumptions.
\end{proposition}
\begin{proof}
\emph{PR01--PR03: early constants.}
Fix $c_{\rm adjust}=1/1000$, the profiles and numerical domain/fiber margins. Let $N$ dominate
LPA06, SGP01, EGP02 and TCP01's WHOLE support-list counts.
Their packing uses $\kappa=1/\Delta$ with numerical dimensionless
ratios; its whole-ball curvature is supplied on the final tail.
Thus $N$ precedes $\Delta$. Let $P\ge1$ dominate the required
profile derivatives. Fix $C_j,\Omega,L_0,b_{\rm cut},\kappa$ as in
GAF01. FC05/FC10 and the three graph producers make these finite
functions of $N,P$, independent of $\Delta$, noncollapse and target
dimension. Retain the actual rank margin $1/2$ and source
right-inverse bound two.

\emph{PR04--PR09: reverse stage choices.}
Choose $c_3$ to satisfy BOTH (JA) and (SE), using $P_0\le P$.
Then perform GAF01's reverse choices, keeping every bound in (JB),
$\Xi_j$ from CFS15, $c_2\le t_3$, $c_1\le t_2$ and
$c_1,c_2\le4\kappa/5$. In particular,
\[
 \begin{gathered}
 e_j<\min\{1/100,\Gamma_j\Sigma_j/100,\Sigma_j/1000,1/(200\Omega)\},\\
 \Sigma_j<\min\{1/2,\varepsilon_j/10000,\Gamma_j/200,
                                      \Gamma_j^3/(100C_j)\}.
 \end{gathered}                                               \tag{RegStage}
\]
The modulus tends to zero, so $\Gamma_j$ exists before $\Sigma_j$.
Every subsequent minimum contains finitely many positive earlier
numbers. These are CFS20's cumulative tube/rank budgets, CFS31's
actual cutoff/plateau bounds and CGP07's one-sheet inequalities.
The forward execution is justified by their cumulative estimates.

\emph{PR10: the derived scale constant.}
After these stages CFS11--CFS12 supply finite $N_b,c_w$. Fix
\[
 C_\rho=100(L_0+1)(1+b_{\rm cut}+N_bc_w/\Sigma_1).
\]
It is derived now, before $\Delta,\Lambda$. It imposes a later
bound on $\Lambda$, not a backwards bound on $c_3$.
Empty families retain harmless positive choices and execute identity
stages.

\emph{PR11--PR13: early circle data, then long lengths.}
Fix comparison requests
$\theta_s<e_3/(10C_3)$,
$\theta_e<\min\{e_2/(10C_2),1/100\}$ and
$\theta_2<e_1/(100C_1^2)$.
TCP02--TCP04 choose their circle/collar qualities and raw reference
errors BEFORE $\Delta$; retain, for example,
$\gamma<\theta_2^2/10^{12}$ and
$E,\delta<\theta_2^2/10^8$ in TCP03, together with the TCP01
Gram and LFR35--LFR38 anchor requests. Choose $\beta_3$ below
LC18's obstruction and $\beta_2$ below LPA03's
noncollapse-independent circle threshold, TCP02's compatibility
threshold and $10^{-5}$ for EGP01.
Choose $\Delta$ afterwards, larger than $10^6$, $100/\beta_2$,
and all fixed short-buffer lower bounds in TCP01--TCP04, FC22
and LFR35--LFR38. This is a finite lower-only choice.
No earlier circle test radius has to exceed $10^6\Delta$:
TCP03 uses the short point on the scalar chart's long segment.
Put $L=10^6\Delta$ and $\ell=10^5\Delta$.

\emph{PR14--PR18: coordinate, border and endpoint errors.}
Choose the ORIGINAL edge/slim qualities and separate value errors
for LFR19--LFR20, LFR27--LFR28, LFR34--LFR38, CGP01--CGP03,
EGP05, SGP03, EGP04 and TCP03--TCP04 at the fixed requests.
Include SGP03's (SB), EGP04's qualities below
$\min\{\theta_e^2/10^8,1/(1000L)\}$ and value errors below
$\theta_e/100$, and TCP03's corresponding $\theta_2$ bounds.
The radial adapted quality and difference-Lipschitz error meet
all three CENTERED comparisons, LC67/LC73's smoothing buffers,
and ZSP's $e_0<1/1000$, $\varepsilon_0<1/100$.
These are requests before the joint cone/radial witnesses; no
constructed function is subsequently replaced.

Choose $\tau,\mu<10^{-8}$ below the original edge enclosure,
full-collar and EGP05 SECTION thresholds, with
$4\tau\Delta<\beta_2$. EGP05 does not enlarge the original disk
bundle. Choose $b',s'<\min\{W,10^{-8}\}$ with $W$ in LFR29.1.
Then choose the strong ENDPOINT quality
$s<10^{-5}\min\{b',s'\}$ and the original endpoint thresholds.
LFR44 supplies $a_*$ before the strong SPLITTING quality $b$.
Fix $\sigma_{\rm col}$ as in LPA04.
The long scalar raw errors $E,\delta$ are fixed requests here,
to be produced by the later splitting tolerances. These choices
do not use the later $w'$ or its analytic data.

\emph{PR19--PR20: scale variation and analytic data.}
Choose $\Lambda>0$ below all the support, comparison and smoothing
bounds in LPA04 and the three actual producers. Retain LFR29.1's
$s'/(10^8\Delta^2)$ bound and, explicitly,
\[
 \begin{gathered}
 L\Lambda<10^{-5},\quad100\Delta\Lambda<10^{-8},\quad
 C_\rho\Delta\Lambda<10^{-6},\\
 \Delta\Lambda<e_1/(1000C_1),\qquad
                       100(2\cdot10^6)\Delta\Lambda<1 .
 \end{gathered}                                              \tag{RegScale}
\]
Any smaller bound requested in TCP04 or the directional comparisons
is included in this SAME choice, with its earlier arguments.
Every bound is positive. Apply LC09 with
$\sigma_{\rm col},\Lambda$; choose ONE $w>0$ below its threshold
and $\omega_3$, and define
\[
 w'=\frac{w}{2(1+2\Lambda^{-1})^3},\qquad
 v_*=\frac{w'}{24I(1)},\qquad
                 \mathcal A(R)=2^{K+2}A'(2R+2,w').
\]
Neither $w$ nor $w'$ is reduced after fixing a strong splitting.

\emph{PR21--PR23: strong splittings and joint zero witnesses.}
Choose $b=\beta_E$ below LFR29.1 and all LFR28/LFR38/EGP05
thresholds at these $v_*,\mathcal A$, including EGP03/TCP02's
actual raw alignment and coverage requests.
Then choose $\beta_1$ below LFR20's threshold,
LFR44's $b_{1,*}(b)$, the remaining SGP02/EGP03/TCP02 bounds,
$\beta_2$ and the slim quality. Test ranges such as
$\beta_1^{-1}$ are now finite; their curvature bounds come later.

Apply LPA04's last step to LC73's reference quality, cone error
and lower zero scale, including $T_0\ge1600L$ and every
shell/normal-flow lower bound. Keep the PR14 radial requests.
LPA02 supplies a finite $V\ge T_0$ and a uniform tail for ONE
joint witness at each center, with scale in $[T_0,V]$.
This is a proved existential OUTPUT, not a free mixed interval
invoked without a nonemptiness proof. No condition
$V\Lambda\ll1$ is imposed.

\emph{PR24--PR25: common tail and actual finite choices.}
Take a finite maximum of the uniform tails from LPA04 and the
actual comparison producers. For each fixed upper test range $B$,
LPA01 controls ALL $0<R<B$ simultaneously. Include
$H_\alpha>400V$ and every entire packing/test ball.
There is no supremum of a possibly nonmonotone $A'$ over a real
interval and no single tail for all unbounded $R$.
LPA05--LPA06 select the finite families and their SAME pointwise
witnesses. LFR33/LFR38 choose one physical distance smoothing
for the selected edge family; its physical value error may use
the positive minimum of the selected scales, as their proofs allow.
That post-selection choice changes no previously fixed quality,
predicate, model or early modulus. Use the prescribed affine
planes and execute GAF02's actual maps.

\emph{PR26--PR28: extraction without backwards choices.}
GAF03--GAF07 use the existing markers, rank, inner sections and
$c_3$ margins for whole fibers. ZSP01--ZSP03 use $T_0$,
the initial radial errors and $c_3<10^{-5}$ for GLOBAL isolation
and the whole compact radial trace.
ZSP04 chooses its actual $K_3$ AFTER the map: endpoints avoid a
finite set and interiors cover a compact smaller-slab image.
This choice on a fixed smooth base places no earlier perturbation
tolerance. ZSP05 supplies the relative removal.

EDP01 has (RegScale); EDP02--EDP04 have
$\beta<1/1000$, $h_*<1/1000$, the original $4.01\Delta$ enclosure,
tangential margin $.99$ and joint singular margin $.9$.
Their entire trace lies in an ORIGINAL source buffer, so no
carrier-dependent isotopy smallness is inserted into an early layer.
EDP05--EDP06 use actual descending face differentials and whole-fiber
identities, not another perturbation.
FDC01 uses the existing $\theta_e<1/100$, $\Delta>100$ and
$T_0$ bound. FDC02--FDC04 are exact compactness, coverage,
relative-face and finite assembly arguments. The final finite
base rectangles for corner rounding are chosen afterwards.

This accounts for every free accuracy, scale, buffer and tail
choice in the CURRENT closed construction. FC01--FC18 and
FC20--FC26's algebraic/domain consumers receive their actual
comparisons and enclosures in PR01--PR25.
Those rows also bind the strengthened CFS11--CFS15 cloud/plane/scale
hypotheses, CFS16--CFS31's actual cutoffs, and CGP01--CGP08's
one-sheet/factorization inputs through their produced maps.
PR26--PR28 bind FC34--FC42.
The unused FC19 and unrestricted FC28 add no requirements.
Every upper-only bound is positive and every lower-only bound
finite, with only earlier arguments. FC46 applies to those
choices; the specified LPA02 and base-selection arguments supply
the remaining existential steps. This proves simultaneous admissibility.
\end{proof}

\begin{theorem}[Actual static certificates on the closed standing tail; PBR02]
\label{thm:fibration-closed-standing-static-certificates}
For fixed $K,A'$ and the retained foundations, every sufficiently
late member of the closed standing construction has an actual
FC39 decomposition and finite KL graph presentation on its SAME
smooth carrier. The parameter assignment is independent of the
center and of the member on that tail.
\end{theorem}
\begin{proof}
PBR01 produces one assignment and tail. LPA supplies local packets;
SGP, EGP and TCP supply their actual cloudy images; GAF supplies
the sequential adjustments and whole fibers. ZSP, EDP and FDC
supply the compact restrictions and actual relative faces for those
SAME data. FDC04 and FC42, with FC41's retained library, give
the graph presentation. Member-dependent choices, such as a finite
selected family or compact base intervals, occur AFTER the uniform
parameters and have the properties already proved.
No finite subcover over the sequence of manifolds is used.
\end{proof}

\begin{theorem}[Closed static collapse with one threshold; PBR03]
\label{thm:fibration-closed-static-collapse-threshold}
Retain the stated geometric, analytic and topological foundations.
For every integer $K\ge10$ and positive finite function
$A:(0,\omega_3)\to(0,\infty)$ there is $w_0\in(0,\omega_3)$
such that the following holds. Let $(M,g)$ be a closed connected orientable
smooth Riemannian three-manifold. For its finite curvature scales suppose
\[
 \Vol B(p,R_p)\le w_0R_p^3.
\]
For EVERY $p$, $w\in[w_0,\omega_3)$, $0<r<R_p$ and
$0\le k\le K$ with $\Vol B(p,r)\ge wr^3$, suppose
\[
 |\nabla^k\Rm|\le A(w)r^{-(k+2)}
                           \quad\hbox{on ALL of }B(p,r).
\]
Then $M$ has a finite smooth KL graph presentation. The globally
nonnegative branch is supplied directly by the retained compact
classification and FC41; no volume is evaluated at radius infinity.
The threshold depends on $K,A$ and the retained foundations,
not on $M$, a point or a later local parameter.
\end{theorem}
\begin{proof}
Treat the nonnegative branch as stated. If no threshold exists in
the finite-curvature branch, choose for each sufficiently large
integer $n$ a nongraph counterexample at
$w_n=\min\{\omega_3/4,1/(16n^4)\}$.
If LC01's first volume scale $r_p(1/n)$ were at least $R_p$,
its first-root property would imply
$\Vol B(p,R_p)\ge R_p^3/n>w_nR_p^3$, a contradiction.
Ball-volume monotonicity now gives
\[
 \frac{r_p(1/n)^3}{n}\le\Vol B(p,R_p)\le w_nR_p^3,
 \qquad R_p/r_p(1/n)\ge(1/(nw_n))^{1/3}>2n .
\]
Define ONCE, for $C>0$ and $0<w<\omega_3$,
\[
 D=\max\{1,C\},\qquad
 A'(C,w)=\max_{0\le k\le K}
                   \{A(D^{-3}w)D^{-(k+2)}\}.                \tag{Reduce}
\]
This is positive, finite and independent of $n$. For $n>1$,
$w\in[1/n,\omega_3)$ and $0<C<n$, one has $D<n$ and
\[
 D r_p(w)\le n r_p(1/n)<R_p,\qquad
 \Vol B(p,D r_p(w))\ge w r_p(w)^3.
\]
The triggering ratio is $\widehat w=D^{-3}w>1/n^4>w_n$.
Apply the assumed WHOLE-ball derivative bound at radius $D r_p(w)$
and ratio $\widehat w$, and restrict to $B(p,C r_p(w))$.
Equation (Reduce) gives every order in the closed standing hypothesis
with the SAME function $A'$. No maximum over infinitely many $w$
is taken, and $A$ need not be monotone.
PBR02 makes every sufficiently late counterexample a graph manifold,
a contradiction. This proves the threshold in the required order,
with the original metric and carrier.
\end{proof}

\paragraph{Scope after the closed register.}
This supplies FC44's CLOSED binding and the closed static endpoint
under retained foundations. KL5.1--5.2, printed39--40/PDF34--35,
and15.2/(15.4), printed87/PDF82, provide the source comparison.
(Reduce) also makes the small-$C$ case explicit. Boundary augmentation,
its additional late constraints and its own threshold reduction
remain required. The accepted-PC/Lean boundary is unchanged.
This is not the final full-goal mathematical audit.


\subsection{Boundary parameter binding and the static collapse theorem}
\label{sec:fibration-boundary-parameter-binding}
The closed register PR01--PR28 remains unchanged. The boundary
register below reuses its actual producer inequalities, with the
boundary domain adapters proved in BSA--BCP and the actual augmented
construction proved in BCG--BCF. It is not an application of a
closed-source theorem at a boundary point. All first-volume parameters
used in this branch lie in $0<w<w_{\rm cap}=\omega_3/4$.
The historical source order and the flow-application rows keep their
separate roles. Only the static construction is being bound here.

\begin{proposition}[Boundary register and simultaneous admissibility; BBR01]
\label{prop:fibration-boundary-register-admissibility}
Fix $K\geq10$, the standing derivative function $A'$ supplied by
BSA04, and the retained geometric, analytic and topological foundations.
All requirements of the actual nonproduct boundary construction
BSA--BCP--BCG--BCF, including its interior local producers and
adjustments, admit one assignment independent of the sequence member,
point and boundary component. The register BR00--BR31 below accounts
for the permitted dependencies and the common tail. Its qualitative
thresholds are the proved positive thresholds of the named producers,
not assumptions of their desired geometric conclusions.
\end{proposition}
\begin{proof}
\emph{BR00--BR03: fixed data and augmented early constants.}
Fix BSA01's numerical $\delta_*,C_0$, the auxiliary cap
$w_{\rm cap}=\omega_3/4$, and the already supplied $A'$.
Use the same early whole-list packing constants as PR01, adding
one for the boundary slot. BCG01 proves this WHOLE-list bound,
independent of the number of boundary components. Fix the original
profiles together with $\mathcal B$ in (BCG.0); let $P\geq1$
bound their required derivatives and both
$\|\mathcal B'\|_\infty,\|\mathcal B''\|_\infty$.
All are finite numerical constants on fixed support intervals.

Take the PR03 graph, cutoff, rank and derivative constants for
these augmented bounds, increasing them by a fixed numerical
factor if necessary. BCG03 proves the required graph bounds
without $r_\partial^{-1}$ in any derivative constant. Its use
of $R_a<2r_\partial<1$ is a later geometric requirement; it
places no early dependence on $r_\partial$. Retain the actual
rank margin $1/2$, right-inverse bound two and the same fixed
source-domain margins. Empty families retain positive harmless
bounds and execute identity stages.

\emph{BR04--BR10: adjustments and the slow adjusted-scale constant.}
Choose $c_3$ with every (JA)/(SE) bound, including $c_3<10^{-5}$.
Follow PR05--PR09 in reverse stage order with these augmented
constants, keeping all (RegStage) bounds. In particular
$\Sigma_j<\varepsilon_j/10000$ is already included and supplies
BCG05's every-contributor buffer. Each graph modulus tends to
zero, so the required choices exist before the later source
qualities. Derive $N_b,c_w,C_\rho$ as in PR10 from this first
smoothing. This still precedes $\Delta$ and $\Lambda$.

\emph{BR11--BR13: the extra early circle requests.}
Choose each rough target so the original interior producer uses
less than half its available error. For the boundary block
choose fixed $0<\vartheta_j<1/100$ with
$3P\vartheta_j<e_j/4$, further decreased if required by the
square-sum budget. These are early requests for BCG02--BCG03.
Add to PR11 the fixed short circle-axis raw/value/norm errors
below $\vartheta_1^2/10^8$ and the original short test radii
5000/1000. Choose $\beta_3$ as before, and then $\beta_2$
below BOTH the old thresholds and AC76/FC20's circle-versus-
cusp alignment threshold at the fixed exclusion $\beta_3$.
The cusp splitting quality needed in that comparison is only
recorded as a request and will be supplied later. No early
circle parameter depends on $L$ or a later long test radius.
Now choose $\Delta$ with every PR13 lower bound and set
$L=10^6\Delta$, $\ell=10^5\Delta$.

\emph{BR14--BR19: the original and boundary comparison budgets.}
Retain ALL PR14--PR18 coordinate, radial, border, section and
strong-endpoint requirements. Add BCG02's edge/slim original
adapted qualities, raw alignment and value requests at the
already fixed $\vartheta_2,\vartheta_3$ and long radius $30L$.
For their one-vector test take errors below
$\vartheta_j^2/10^8$; the actual strong splitting thresholds
which produce those raw comparisons are imposed in BR21--BR22.
These remain requests before any source function is constructed.
The no-two exclusion is the already fixed $\beta_2$.
LC73's radial and the weak-edge distance smoothing remain the
SAME eventual functions for all consumers.

Choose $\Lambda$ with every PR19 inequality, explicitly retaining
(RegScale), $L\Lambda<10^{-5}$ and
$C_\rho\Delta\Lambda<10^{-6}$. All their arguments are now
fixed and every bound is strictly positive. No backwards
$\Delta c_3$ or $V\Lambda$ condition occurs.

\emph{BR20--BR23: analytic data and finite joint witnesses.}
Before choosing $w$, intersect LC09's closed threshold with the
boundary-local threshold proved in BCP04 and $w_{\rm cap}$.
Choose ONE $w$ in that positive interval, define
$w'=w/[2(1+2\Lambda^{-1})^3]$, and use BSA06's
$v_*=w'/(24I(1))$ and
$\mathcal A(R)=2^{K+2}A'(2R+2,w')$.
Both volume parameters stay in the restricted auxiliary range.
No supremum of $A'$ on a real interval is required.

Choose the strong splitting quality $b$ and then $\beta_1$
with all PR21--PR22 thresholds at these fixed analytic data.
Include the edge and slim AC76/FC20 requests just prescribed.
BCP04 supplies the SAME local producer proof on its exhausting
interior domains, with exactly these allowed predecessors.
The circle request was already imposed before $\Delta$.
Then choose the last cone/radial/lower-zero-scale requests of
PR23, with $T_0\geq1600L$. The joint producer outputs a finite
$V\geq T_0$ and a uniform local tail. We do not replace this
proved existential output by a freely chosen mixed interval.
Neither $w$ nor $w'$ is decreased after $b,\beta_1$ are fixed.

\emph{BR24: late cusp accuracy and physical scale.}
There are now only finitely many cusp splitting, norm and
buffer requests: the three raw affine comparisons in BCG02,
the no-zero splitting in BCP05, the all-direction tests of
BCP02, and the normalized domains in BCG01--BCG03. Choose a
positive cusp splitting quality below their thresholds and
below $\beta_1$, and choose its adapted/norm error below the
requested fixed margins. These choices affect no reference
quality. Let $H_\partial$ be a finite common upper bound for
all the resulting cusp test radii and their first-exit buffers.
The product estimate in BCP02 is uniform in the torus lattice
and gives positive physical bounds after those qualities are fixed.

Put $\vartheta=\min_j\vartheta_j$. Choose $r_\partial>0$
below those product/first-exit bounds and, explicitly, below
\[
 \min\left\{\frac1{1000L},\frac1{10^4},
       \frac1{100H_\partial},
       \frac{\vartheta^2}{4\cdot10^8(12L+1000)},
       \frac{10^{-6}}{20(c_3+1)}\right\}.                 \tag{BRegPhysical}
\]
All entries are strictly positive and have earlier arguments.
This gives BCG02's Hessian integral bound, BCG01's complete
support-domain buffer, BCG05's exact-marker localization and
$20c_3r_\partial<10^{-6}$ for the actual torus core.
The parameter is chosen AFTER the finite $V$, as in the source
order, and no earlier modulus is changed.

\emph{BR25: one uniform tail, not separate pointwise tails.}
Take the maximum of the finitely many local/overlap tails and
those in BSA01--BSA06, BCP01--BCP05 and BCG01--BCG07.
Include $1/n\leq w'$, the envelope curvature buffers in BSA05,
$H_n=n/4>400V$, all finite splitting/inverse-quality and packing
radii, and BCP02's errors with the qualities in BR24.
The estimates $\delta_n\to0$ and the uniform collar smallness
of $\rho$ through raw height96 supply all the fixed physical
bounds, independently of the selected centers or boundary count.
They apply to every admissible scale and the SAME smoothed heights.
In particular the expanded level92 collar separation, no-zero
contradiction and all original comparison domains hold together.

For the selection buffers impose, for the finitely many fixed
constants $C$ used by the producers and witness comparisons,
\[
 C\rho(p)<\tfrac12d(p,\partial M)\quad(d>10),
 \qquad
 10\Delta\rho(p)<\tfrac13d(p,\partial M)\quad(d\geq35).
                                                               \tag{BRegInterior}
\]
BCP04.a makes these uniform: $d/\rho>5n/8$ on $d>5$,
so it suffices respectively that $n>16C/5$ and $n>48\Delta$.
Include $C=100\Delta$ for the WHOLE edge circle collar and
every fixed zero/cone buffer. Thus selected edge centers at
$d>20$ have their full circle-collar domains in $d>10$, while
all strong replacement witnesses from the remaining carrier
are also at $d>20$, as BCF02 requires. The four-family cover
of $d\geq35$ follows from BCG01 on this SAME tail.
BSA06 bounds every derivative test for ALL radii below each
fixed finite upper range at once. There is no tail asserted
for all unbounded radii, and no bounded-on-compacts assumption
on the arbitrary function $A$ or $A'$.

\emph{BR26--BR31: actual choices after the uniform parameters.}
On a fixed member, choose the smooth scale by BSA03--BSA05
and the smooth original cusp heights by BCP01 with their
prescribed errors. The estimates above hold for these actual
choices. Choose once the original pointwise local witnesses;
select circle/slim/zero centers in $d>10$ and strong edges in
$d>20$. Their finite maximal selections and uniform packing
bounds are BCP04/BCG01, not an additional covering assumption.
The shared physical weak-edge smoothing may use the positive
minimum of this finite family's scales, as in PR25. This
post-selection operation meets the preset errors without
changing the original metric, earlier qualities or local predicates.
Use BCG03's actual affine blocks and CFS27's actual chosen
planes, retaining every boundary block, and execute the three
actual adjustments in their forward order.

BCG04--BCG06 give global isolation, exact markers and the whole
adjusted torus cores. ZSP01--ZSP02 apply on the original compact
interior traces. These are BR27's $Z,C,M_1$.
For BR28, BCF01 chooses $K_3$ after the map to include the
compact smaller-slab image AND all finite zero/cusp base face
points in its interior. This gives the actual compact slim
piece and relative $M_2$ without an earlier perturbation request.
BR29 uses BCG07/BCF02's proper disks, regular descended face
equations, exact vertical circles and eligible strict replacement
to make the edge piece compact. BR30 uses the actual relative
circle remainder and its full base corner charts. BR31 uses
BCF03's embedded face partition BEFORE the Euler count, and
BCF04's finite collared gluing/rounding under FC41, preserving
every original external boundary label. The latter finite base
rectangles and topology witnesses are chosen after the map;
they impose no earlier smallness bound.

These rows retain all PR01--PR28 inequalities for the original
interior blocks and account for every new boundary free choice.
BSA02 and the scalar smoothing/envelope arguments supply the
boundary analytic conventions. BCP supplies their actual local
domain adapter; BCG supplies all actual added cloud hypotheses;
BCF supplies the new whole faces and finite assembly. Upper-only
choices have positive bounds with earlier arguments and lower-only
choices have finite bounds. FC46 therefore applies to the scalar
choices; the proved joint-witness and finite base constructions
supply the explicitly separated existential steps. This proves
simultaneous admissibility, with no circular feedback from the
physical boundary scale to an early graph or cutoff constant.
\end{proof}

\begin{theorem}[Actual boundary certificates on the standing tail; BBR02]
\label{thm:fibration-boundary-standing-static-certificates}
For the fixed $K,A'$ and retained foundations, every sufficiently
late member of BSA04's boundary counterexample construction has
a finite smooth KL graph presentation on its SAME carrier, with
its original labeled external boundary. One parameter assignment
and one uniform tail work for all its points and components.
\end{theorem}
\begin{proof}
Use BBR01. BSA supplies the intrinsic scale and analytic data;
BCP supplies actual cusp coordinates, the product-or-separated
alternative and the restricted interior packets. The whole-product
case is already a graph piece. Otherwise BCG gives the actual
augmented clouds, every-contributor marker control, adjustments,
whole proper fibers and adjusted torus cores. BCF gives the actual
compact pieces, independent face equations, whole circle equality,
embedded surface partition and finite graph assembly. All use
the SAME selected maps and source metric. Member-dependent finite
families, scalar approximations, compact base intervals and rounding
rectangles occur only after the uniform choices, with the proved
properties. No compactness or finite subcover of the sequence of
carriers is assumed. The graph conclusion contradicts the chosen
nongraph property of any member on this tail.
\end{proof}

\begin{theorem}[Static collapse with nearly cuspidal boundary; BBR03]
\label{thm:fibration-boundary-static-collapse-threshold}
Retain the stated geometric, analytic and topological foundations.
For every integer $K\geq10$ and positive finite function
$A:(0,\omega_3)\to(0,\infty)$ there is
$w_0\in(0,\omega_3)$ with the following property.
Let $(M,g)$ be a compact connected orientable smooth Riemannian
three-manifold with NONEMPTY smooth boundary. Suppose each
boundary component has diameter at most $w_0$ and admits the
actual depth100 $C^{K+1}$ cusp embedding of pairs with relative
pulled-back $C^K$ metric error at most $w_0$, in the precise
normalization and intrinsic metric convention preceding BSA01.
At every point with $d(p,\partial M)>10$ suppose
\[
 \Vol B_M(p,R_p)\leq w_0R_p^3.
\]
For EVERY $p\in M$, $w\in[w_0,\omega_3)$, $0<r<R_p$
and $0\leq k\leq K$ with $\Vol B_M(p,r)\geq wr^3$, suppose
\[
 |\nabla^k\Rm|\leq A(w)r^{-k-2}
                            \quad\hbox{on ALL of }B_M(p,r).
\]
Then $M$ has a finite smooth KL graph presentation retaining
its original external boundary labels. The threshold depends
only on $K,A$ and the retained foundations.
Together with PBR03, one common smaller threshold covers the
closed and nonempty-boundary cases, with their stated conventions.
\end{theorem}
\begin{proof}
Restrict prospective thresholds below BSA01's $\delta_*$;
then every curvature scale in this nonempty-boundary case is
finite by that lemma. If no positive threshold exists, choose
nongraph counterexamples with
$\delta_n=\min\{\delta_*,\omega_3/8,1/(16n^4)\}$.
BSA04 proves at EVERY point, including the external boundary,
$R_p>2n r_p(1/n)$ and the required whole-ball derivative
estimates for ONE function
\[
 A'(C,w)=\max_{0\leq k\leq K}
    A(D^{-3}w)D^{-k-2},\qquad D=\max\{1,C\}.
\]
It handles $d\leq10$ by the actual cusp volume estimate,
not by extending the volume-collapse premise to unassumed points.
It handles $C<1$ with $D=1$, and every triggering ratio is
strictly above $\delta_n$. Thus the reduction uses exactly
the stated same-metric hypotheses, with no monotonicity or
regularity imposed on $A$. Its auxiliary first scales are only
in $0<w<w_{\rm cap}$, so it does not import the false full
Euclidean small-ball range at an external boundary point.

BBR02 supplies a graph presentation on each sufficiently late
counterexample, a contradiction. Hence the asserted positive
threshold exists. Labels are retained by the actual BCF04
construction, not added to an unlabeled recognition afterward.
For the empty-boundary case use PBR03, including its separately
recognized globally nonnegative branch; no distance to an empty
boundary or volume at radius infinity is evaluated. Taking the
minimum of the two positive thresholds preserves all hypotheses:
the geometric errors and collapse ratios become smaller, while
the derivative assumption covers the larger required parameter
interval. This gives the common static threshold without a
metric doubling or a flow-to-static inference.
\end{proof}

\paragraph{Scope of the completed static construction.}
PBR03 and BBR03 are the closed and nearly cuspidal boundary
static endpoints under the retained foundations. The actual
geometric constructions and simultaneous parameters now supply
the ACTIVE written FC23--FC45 consumers in these two scopes;
the false unrestricted FC28 assertion remains excluded as above.
The requested ONE final full-goal mathematical audit is the next
step; no completion claim is made before that audit and any
necessary repairs. Flow application, incompressibility, global
reconstruction and accepted-PC/Lean bindings remain distinct.

\subsection{Final mathematical audit of the static construction}
\label{sec:fibration-final-static-goal-audit}
Revision159 records the requested ONE final mathematical audit of
the Chapter14 goal, including repairs within that same review.
The review used the actual statement/proof bodies and their original
domains, rather than treating source citations, parameter hashes or
successful document builds as proofs. Its seventeen consumer checks
and source records are in \code{AUDIT_CHAPTER14_REVISION159.md} and
\code{chapter14_goal_audit_revision159.csv}.
The following small numerical binding correction is part of the
ACTIVE route; historical statements and registers remain preserved.

\begin{lemma}[Explicit edge-exclusion bounds in both registers; CAA01]
\label{lem:fibration-final-edge-exclusion-register-repair}
In PBR01 and BBR01 require explicitly
\[
 \beta_2<10^{-6},\qquad s<10^{-6},\qquad b<10^{-6},
 \tag{AuditEdge}
\]
at the original choices of the early two-splitting quality, strong
endpoint quality and strong splitting quality, respectively.
Both simultaneous assignments and all their conclusions still hold.
In particular the hypotheses of EGP01 are now explicitly bound.
\end{lemma}
\begin{proof}
The displayed PR12 bound $\beta_2<10^{-5}$ alone is weaker than
EGP01's stated $10^{-6}$ hypothesis. Replace that displayed upper
bound, in the active reading of PR12 and its BR12 counterpart,
by $10^{-6}$, retaining every other positive upper bound.
This occurs before $\Delta$, the endpoint choices and the analytic
data. The finite positive minimum is still positive. Choose
$\Delta$ only AFTER this refined value, including its existing
$100/\beta_2$ lower bound.

At the original later choices of $s$ and $b$, also intersect
their allowed intervals with $(0,10^{-6})$. The existing
$s<10^{-5}\min\{b',s'\}$ and $b',s'<10^{-8}$ already imply
the requested endpoint bound; its repetition makes the consumer
literal. The strong splitting bound is another fixed positive
upper bound at PR21/BR21, before $\beta_1$ and the joint zero
witness. It creates no backwards dependence. All subsequent
thresholds, witnesses and tails are chosen from these refined
predecessors, exactly as in the original registers; no previously
constructed coordinate or selected family is modified.

The actual four-target error in EGP01 is now
$5\beta_2+b+s<7\cdot10^{-6}<10^{-3}$, within LFR29's
obstruction. Thus its stated hypotheses as well as its proof
margin hold. All other requirements remain finite positive upper
bounds or finite lower bounds in the same order. FC46, together
with the already proved existential witness steps, proves both
refined assignments. PBR02--PBR03 and BBR02--BBR03 consequently
use these assignments. This is a project register repair, not
an asserted correction to the source's qualitative Lemma9.6.
\end{proof}

\paragraph{What the audit checks.}
The detailed record has the following grouped conclusions.
\begin{enumerate}[leftmargin=*]
\item SGP, EGP and TCP produce the actual finite comparison lists,
original-coordinate affine comparisons, early graph moduli,
all-preimage rank and both cloud coverage directions. SGP03's
model-support test supplies its needed smooth domains; the stronger
unused common-domain demand is not silently assumed.
\item CFS26--CFS31 check EVERY contributing center and plane for
actual zero markers. CFS11--CFS15 give one global smoothed manifold,
the full tube and whole-set nearest projection. CFS16--CFS20 and
the actual uniform cutoffs supply each sequential adjustment.
The refuted unrestricted FC28 and the unused carrier-specific
CFS19 tolerance are not premises of this argument.
\item CGP04--CGP08 first localize every marked base point, prove
injectivity on the whole relevant graph and exhaust it by actual
source images. Later retained blocks exclude merging even across
patches. GAF and EDP prove BOTH inclusions for whole fibers and
trap the ENTIRE isotopy trace in original compact source buffers.
Their properness arguments concern full preimages of compact
target sets, including the disk-boundary restriction.
\item EDP01 bounds the adjusted scale derivative before choosing
$\Lambda$. EDP03 retains the full quotient derivative. The actual
horizontal and vertical defining functions are independent, and
EDP06 identifies each entire disk-boundary circle with an entire
circle fiber. ZSP and FDC use relative removals, strict replacement
charts and open-end compactness to produce the actual closed
decomposition. The embedded face partition precedes the Euler
count, and the same base arcs round both sides of a corner.
\item BSA proves the intrinsic boundary scale adapters, including
convex-boundary volume comparison and the capped first-volume
range. BCP retains the whole frozen torus, verifies all adapted
tests, proves the product-or-disjoint collar alternative and
uses the actual zero-ball-meets-zero premise. Interior packets
use exhausting interior domains, not boundary-point smooth models.
\item BCG constructs the physical boundary blocks and checks all
their contributing planes, global isolation and exact markers.
Its torus core is the whole inner subgraph; its internal frontier
is distinguished from its external boundary. BCF verifies the
restricted edge eligibility, all whole faces and compact pieces.
The cusp torus has no disk faces only AFTER its actual partition
is proved. All external boundary labels survive the torus gluings.
\item PBR and BBR, with (AuditEdge), bind every active choice in
one permitted order and one uniform tail. The late physical scale
never enters an early graph constant. There is no $\Delta c_3$
or $V\Lambda$ demand, nor a boundedness assumption on the arbitrary
derivative function. The exact small-$C$ failure reduction supplies
the two static thresholds on the SAME original metrics.
\end{enumerate}

\begin{theorem}[Completion of the written static assembly goal; CAA02]
\label{thm:fibration-written-static-goal-completion}
Under the explicitly retained geometric, analytic and topological
foundations, the active Chapter14 construction now proves the
closed static conclusion PBR03 and the nearly cuspidal boundary
conclusion BBR03. In both scopes it produces the actual finite
smooth KL graph presentation on the original carrier, with its
collared attaching maps and, when present, original external
boundary labels. Its actual cloud hypotheses, marker control,
one-sheet and whole-fiber assertions, proper restrictions,
compatible faces and simultaneous parameters are supplied by
the written producers above, including CAA01.
\end{theorem}
\begin{proof}
Use CAA01's refined PBR01 assignment. The local packets feed the
actual SGP/EGP/TCP clouds; CFS/CGP/GAF supply the adjustments and
whole fibers; ZSP/EDP/FDC supply the compact compatible pieces.
FC42, with FC41, gives the closed standing conclusion. PBR03's
exact failure reduction then gives its static threshold, including
the separately recognized nonnegative-curvature branch.

For nonempty boundary use CAA01's refined BBR01 assignment.
BSA supplies the intrinsic scales and analytic data; BCP supplies
the actual collars and restricted interior packets. Its product
alternative is already a graph piece. Otherwise BCG supplies the
augmented actual construction and BCF its compact labeled graph
assembly. BBR03 gives the boundary static threshold. Its common
minimum with the closed threshold retains the stated hypotheses.
Neither implication uses a flow application, essential-torus
refinement or later ambient reconstruction. These are exactly
the two static outputs, with their distinct input conventions,
required for FC45's next consumer.
\end{proof}

\paragraph{Disposition and retained foundations.}
The requested written Chapter14 goal and its ONE final mathematical
audit are complete, with the numerical finding repaired above.
This disposition supersedes earlier revision-specific frontier
paragraphs for the ACTIVE static route; those paragraphs remain
as historical records. It does not restore an unused stronger
contract or the false unrestricted cloudy-smoothing assertion.

The inherited Euclidean and spectral calculus, smooth approximation,
ODE, isotopy and bundle foundations remain supplied. The geometric
inputs and category adapters retained in Chapters3--4 and13,
including their Alexandrov and model-classification foundations,
keep their stated qualifications. FC41's category-correct
recognition and closure library remains an explicit topological
input. Binding these foundations to the accepted completed PC
release, elaborating Lean declarations and checking their axioms
are separate tasks. No such binding or Lean proof is certified
by this written mathematical audit. The output is the raw KL
graph class; incompressible tori, flow application and final
Geometrization reconstruction retain their separate gates.

\begin{definition}[Static-to-global handoff; FC45]
\label{def:fibration-static-handoff}
The Chapter14 output, conditional on the qualified producers and
FC44's binding, is a finite graph presentation on the EXACT
closed carrier or nearly cuspidal boundary carrier, in the required
category, with its collared embeddings and external boundary labels.
Its input is LC87 or LC88 with LC81's category data. Its output
is consumed by LC90's flow application, Chapter6's relative graph
refinement, and Chapter15's ambient transport. It neither supplies
Chapter12's late carriers nor accepts H1--H4 or the endpoint.
\end{definition}
The proof stages are
\begin{gather*}
 \text{local packets}\longrightarrow\text{clouds and adjustments},\\
 \text{clouds and adjustments}\longrightarrow\text{static graph assembly},\\
 \text{static graph assembly}\longrightarrow
       \text{flow application and relative refinement}.
\end{gather*}
Chapter3 supplies metric extraction and the fixed-error conventions;
Chapter4 supplies the geometric compatibility used in FC23, with
its own existing producer gates. Generic linear algebra, smooth
calculus, compact isotopy, bundle classification and elementary
topology stay in the inherited-foundation audit. No current PC
snapshot is modified, no new declaration probes are run, and no
mathlib environment is pinned. Chapter16 still consumes either
late thick/thin data OR explicit terminal exceptional data;
finite-time extinction is not added to the default route.

\section{Historical Chapter 14 specification audit (revision101)}
\label{sec:fibration-complete-audit}
The following preserves the revision101 specification audit.
For the current written completion and its retained-foundation limits,
see Section~\ref{sec:fibration-final-static-goal-audit} above.
The chapter now has 46 implementation contracts and complete
blueprint coverage from adapted-coordinate compatibility through
closed and boundary static graph assembly. This is a specification
and source audit, not a claim that all 46 outputs have been proved
in Lean or that the qualified producers are discharged.
The row-by-row record is \code{chapter14_audit_revision101.csv};
the full findings are in \code{AUDIT_CHAPTER14_REVISION101.md}.
The audit checks statements, quantifier order, uniformity, original
domains, category, source provenance, dependencies and consumers.

\begin{enumerate}[leftmargin=*]
\item FC01--FC19 retain their earlier mathematical statements.
FC19's common large domain is conditional; FC22 supplies the
asymmetric alternative needed by the actual parameter order.
FC16's centered values and FC18's support bounds still require
the ORIGINAL coordinate normalizations and source error budgets.
\item FC20--FC26 supply explicit elementary alignment, comparison,
graph modulus, identical-radius coverage and selected-radius
arguments. FC23 remains the precise geometric overlap binding.
FC07 and FC27 remain qualified cloud producers; citations and
inventory edges do not close them.
\item FC28--FC33 separate cloud smoothing, its local span repair,
ambient cutoff buffers, the full chain-rule error, fiber inclusion
and the stronger graph/connectedness proof for final bases.
Normal-projection rank and total ambient dimension are not confused.
\item FC34--FC39 require common compact isolation, full preimages,
proper face-respecting bundle restrictions, the corrected $D^3$
case, and actual boundary equalities. Source-omitted proofs are
recorded as implementation tasks rather than supplied proofs.
\item FC40--FC43 distinguish the elementary Euler count and finite
assembly argument from their inherited topological recognition
inputs. The boundary route uses its actual retained collar blocks;
no doubling, essential tori or ambient carrier identification is
inferred from the static graph certificate.
\item FC44--FC46 expose the full source order, its application
register and the conditional parameter-selection lemma. The stage
graph is acyclic; no local/cloud/static producer depends on LC90,
Chapter12's late carriers, Chapter15 transport or Chapter16 assembly.
\end{enumerate}
Author errata and project findings are separately recorded. The
former include the $D^3$ case and corrected20.2 recentering/enlarged
sets. The latter include the slim index and support bound, the
normal-projection rank issue, the truncated-ball coverage margin,
unequal test domains, cutoff support buffers, and the common
compactness requirement already identified in LC82.

There are no remaining unexpanded Chapter14 TOPIC placeholders.
The remaining work is the explicitly listed source-qualified
producer proofs, inherited release bindings and eventual Lean
implementation. These are not presented as completed mathematics.
The natural next blueprint topic is Chapter15's actual ambient
transport, using this exact static certificate and retaining the
existing cap, relative-boundary and endpoint gates.
\section{A quantitative register for the collapse parameters}
\label{sec:fibration-constraint-register}

\subsection{Weakening before selecting a smoothing modulus}
The order in KL (15.4) chooses the cloud quality before the smoothing
error; KL Section 13.2 accordingly writes the latter as a function of
the former. The following two elementary arguments make that use of
Lemma 20.2 precise. They do not prove its geometric smoothing input.

\begin{lemma}[Slack-aware restriction toward an affine plane]
\label{lem:fibration-cloud-quality-weakening}
Let $H$ be a real normed vector space, $x\in S\subset H$,
and let $A\subset H$ be an affine subspace containing $x$.
Suppose $R>0$, $e>0$, and the ambient Hausdorff distance between
$S\cap\overline B(x,R)$ and $A\cap\overline B(x,R)$ is at most $e$.
For every $0<r\le R$, the ambient Hausdorff distance between
$S\cap\overline B(x,r)$ and $A\cap\overline B(x,r)$ is at most $6e$.
No closedness of $S$ or attained nearest point in $S$ is assumed.
Here a Hausdorff bound means the two corresponding distance-to-set
bounds, so it is meaningful for these nonempty sets without
assuming that they are closed.
\end{lemma}
\begin{proof}
Take $z\in S\cap\overline B(x,r)$. The hypothesis supplies
$a\in A\cap\overline B(x,R)$ with $|z-a|<2e$.
If necessary move $a$ radially toward $x$ to $a'$ at radius $r$.
This moves it by at most $2e$, so $|z-a'|<4e$.

Conversely take $a\in A\cap\overline B(x,r)$.
If $r\le4e$, use $x\in S\cap\overline B(x,r)$.
If $r>4e$, move $a$ toward $x$ to radius at most $r-4e$,
obtaining $a'$ with $|a-a'|\le4e$.
The original coverage supplies $z\in S\cap\overline B(x,R)$
with $|z-a'|<2e$. Consequently $|z-x|<r-2e<r$ and
$|z-a|<6e$. These two estimates give the asserted Hausdorff bound.
The strict $2e$ witness bounds use only the defining infimum.
\end{proof}

The same $6e$ bound holds with open balls throughout. For $r\le4e$
the common point $x$ suffices in both directions. For $r>4e$, in
the first direction trim $a$ to radius at most $r-2e$, moving it
by less than $4e$ and obtaining total error less than $6e$.
The second direction above already uses radius $r-4e$, so its
selected point lies strictly inside the smaller ball. For different
open/closed input and output conventions take $r<R$; the same
strict buffers apply. Thus the following use does not silently
identify the two source conventions.

KL Definition 3.1 introduces its convention for closed subspaces,
whereas the later cloud notation names subsets. An implementation
must retain its actual truncated-set bounds or a proved
closure/convention adapter. This lemma uses the actual bounds;
it supplies neither closedness nor such an adapter by naming
the sets cloudy. It needs no finite-dimensional compactness.

For a fixed cloudy packet, apply the lemma at every cloud center
after translating that center to zero and dividing distances by its
positive radius. If its quality is $\Gamma$ and
$0<\Gamma<\delta/8$, then it also satisfies the fixed pointed
Hausdorff tests of quality $\delta$: use $R=\Gamma^{-1}$,
$r=\delta^{-1}$, and $e=\Gamma$. Radius comparability and all
chosen affine planes remain the same. This is a weakening statement
with explicit slack, not equality of fixed-error predicates. In the
notation of FC28 the tested set is the ENLARGED set $\widetilde S$.

\begin{lemma}[A smoothing modulus in the source's choice order]
\label{lem:fibration-smoothing-modulus-selection}
Fix $k,K,C$ and retain FC28 as a supplied theorem. There are
$\theta_1>0$ and a function
$\Xi:(0,\theta_1)\to(0,\infty)$ with
$\Xi(\Gamma)\to0$ as $\Gamma\downarrow0$, such that every
admissible $(C,\Gamma)$ cloudy packet has the FC28 conclusions
with accuracy $\Xi(\Gamma)$. Here admissible retains FC28's bounded
enlarged set and radius function bounded above and away from zero, as
well as its tested-set and source-convention hypotheses. The modulus is independent of the
ambient dimension whenever FC28's input threshold is.
\end{lemma}
\begin{proof}
For each integer $m\ge1$, apply the supplied FC28 theorem with
accuracy $2^{-m}$, obtaining a threshold $\delta_m>0$.
Set
\[
 \theta_m=\min\{2^{-m},1/8,\delta_1/8,\ldots,\delta_m/8\}.
\]
These numbers are positive, nonincreasing and tend to zero.
For $0<\Gamma<\theta_1$, the set
$\{m\ge1:\Gamma<\theta_m\}$ is nonempty and finite, since
$\theta_m\le2^{-m}$. Let $m(\Gamma)$ be its largest element,
and define $\Xi(\Gamma)=2^{-m(\Gamma)}$.
The preceding weakening lemma gives the
$\delta_{m(\Gamma)}$ cloudy hypothesis from
$\Gamma<\theta_{m(\Gamma)}\le\delta_{m(\Gamma)}/8$.
FC28 therefore supplies the claimed output.

For every $\eta>0$, choose $M$ with $2^{-M}<\eta$.
If $0<\Gamma<\theta_M$, then $m(\Gamma)\ge M$, and hence
$\Xi(\Gamma)\le2^{-M}<\eta$. This proves the limit assertion.
Only $k,K,C$ and the supplied threshold choices enter the modulus.
\end{proof}

CFS14--CFS15 now supply a strengthened replacement and its own modulus.
The displayed older argument remains conditional on its stated premise.
For any valid replacement supplying the required thresholds, a positive
smoothing-error budget determined by predecessor parameters can be met
by choosing $\Gamma$ sufficiently small,
and then choosing $\Xi=\Xi(\Gamma)$. If an application requests
another smaller tolerance after the cloud has been fixed, it must
provide the stronger cloud hypothesis again; the modulus grants no
unconditional arbitrary improvement of an existing cloud. The full
FC44 inequality register, FC23's actual domain binding, and FC28's
actual cloud/consumer bindings remain open. CFS06 refutes the general
source-stated premise, so the displayed conditional modulus cannot be
used to import that premise. CFS15 reuses the quantifier method with
CFS14's stronger scale hypothesis and the additional radius bound.

\subsection{The bounded constraint register and what is still missing}

The companion \code{formalization/collapse_constraints.csv} contains
19 checked constraint rows. Each records fixed inputs, the next choices,
predecessor rows, units, tested domains, quantifier order, feasibility
status, exact source locators and a failure case. It is a starting subset
of KL Sections 5--16, not a claim that every inequality has been extracted.
In particular FC46's backward-selection lemma is still conditional on a
complete, simultaneously admissible register.

The register covers the original-map identities and overlap tests;
uniform local derivative-block bounds; the graph-to-cloud scale;
the new weakening and smoothing modulus; cutoff support inside actual
projection tubes; the three accumulated adjustment budgets; common
compact isolation of entire inverse images; finite-regularity, boundary
and whole-ball derivative requirements; and the H1--H4 application gate.
For example the cutoff chain rule requires
\[
 \|Dg-Df\|\leq abL+dL+e,
\]
including the factor $L$ from the original map derivative. Every stage
must fit both the rank margin and the next tube. Likewise isotopy needs
one compact isolating set for the entire inverse image throughout the
homotopy, not only small perturbations near one selected component.

The unresolved mixed scale intervals are explicit: a valid acyclic
dependency order alone supplies no point in an empty interval.
Uniformity in the number of ambient coordinate blocks is not supplied
by a bound on one local chart. A center-only curvature estimate is not
the whole-ball estimate of KL16.1. The limit atlas/metric's finite
regularity is not an automatic smooth soul or normal-flow interface.
These remain producer or adapter obligations. The two preceding
lemmas remove a quantifier-order bottleneck while preserving all these
geometric qualifications.

Source checks used KL3.1, 3.4--3.5, 13.2, 15.2/(15.4), 16.1,
17.3 and 20.1--20.2, together with the official May 15, 2015
correction sheet. In particular the cloud is tested on the enlarged
set, and the Lemma20.2 correction concerns its affine projection
conclusion. Exact printed/PDF locators and retained source identities
are in \code{reference_checks_revision104.md}; unchanged hypotheses
reuse the earlier recorded source checks.


\chapter{Ambient transport and reconstruction}
\label{ch:ambient-reconstruction}

\paragraph{Scope.} Actual maps from local/factor data to late and original carriers.

\paragraph{Chapter coverage.}
\begin{enumerate}[leftmargin=*]
\item Finite collared gluing.
\item Sphere reconstruction.
\item Cap avoidance.
\item Actual torus maps.
\item Cusp injection.
\item Finite history.
\item Terminal data.
\end{enumerate}

\paragraph{Coverage and proof status.} Revision 102 completes the detailed
blueprint and full audit of this chapter: AT01--AT31 cover generic marked
transport, actual late prime carriers, finite event reconstruction and
the terminal alternative. Seven source-qualified producer/comparison
contracts remain explicit, one of them optional on the existential route.
No new Lean proof or foundation acceptance is claimed. The final audit
states the current disposition of the historical arguments below.

% BEGIN migrated B051
\section{Finite gluing injection and sphere-reconstruction transport}
\label{sec:gluing-injection-transport}

Revision 53 supplies a written injection argument for the good-block
normal form, then applies the same argument to sphere reconstruction.
These are mathematical proofs from stated smooth-topology interfaces,
not new Lean declarations. All geometric G, J and R implementation statuses remain
unproved. The proof uses the actual gluing maps; an abstract group
isomorphism is insufficient for R06/G05.

\paragraph{Source and route comparison.}
The archived Hatcher Corollary 3.3 and its proof, printed 48/PDF 49,
remove circles from a disk map by replacing a subdisk whose boundary is
nullhomotopic in the two-sided surface. We adapt that operation to a
finite seam system. The Loop Theorem is needed earlier for R05's
conversion of disk incompressibility to group injection, not for the
replacement operation once injection is supplied. The relative smooth
approximation, transversality and collar interfaces used in that proof
remain inherited-foundation binding obligations.

For comparison, \cite{HatcherATChapter1} Section 1.B, printed 91--95/PDF
71--75 of the separately downloaded chapter, constructs graphs of groups.
Theorem 1B.11 proves vertex-group injection when every edge map is
injective. Examples 1B.12--1B.13 distinguish amalgamation from an HNN
extension. The explicit HNN presentation is left to Exercise 4 on printed
96/PDF 76; it is not a supplied presentation proof. The theorem concerns
constructed classifying spaces, so applying it to manifold gluings still
needs a comparison. We use the disk-map proof below, which needs no
asphericity assumption on either blocks or seams. The author's cumulative
errata correct the covering argument on printed 94 and add an injection
hypothesis to Exercise 9 on printed 96; both corrections are present in
the inspected chapter. Source identities, actual reading and correction
scope are recorded in \code{reference_checks_revision53.md}.

\paragraph{Finite collared gluing lemma.}
Let $X$ be a connected smooth compact three-manifold obtained from a
finite family of connected blocks $B_v$ by pairing distinct boundary
ports along specified diffeomorphisms. The paired ports are connected
closed surfaces; every port is used at most once. Store disjoint collars
and assume that the seam images in $X$ are embedded and two-sided, with
the expected cut-block reconstruction. Edges may join two ports of one
block; cycles and multiple edges are allowed. Suppose each of the two
port maps of every paired surface induces an injection into the
corresponding block group. Then every reconstruction map
$q_v:B_v\to X$ induces an injection on $\pi_1$.

If ports of one block are paired, $q_v$ need not be injective as a map of
the closed cut block: two boundary copies are identified. One may instead
use a core obtained by shrinking boundary collars, which is embedded in
$X$ and homotopy equivalent to $B_v$. Fix a basepoint in this core.
All statements about groups use these actual maps and the chosen
basepoint paths. No ambient irreducibility is assumed in this lemma.

\emph{Proof.} Represent a class in the kernel of $(q_v)_*$ by a loop in
the interior core. Choose a disk map into $X$ extending that loop. Apply
smooth approximation and transversality relative to its boundary; the
boundary stays away from the finite seam system. The inverse image of
the seams is a compact one-dimensional submanifold of the disk interior,
hence a finite union of disjoint circles. If there are none, the disk
misses the seams and lifts to the same cut block as its boundary loop,
giving the required nullhomotopy in $B_v$.

Otherwise choose an innermost inverse-image circle $C$, bounding a
subdisk $D$. The interior of $D$ misses all seams and lifts to a single
cut block. Transversality specifies which boundary port contains its
boundary loop $\gamma$. The disk $D$ proves that this loop is trivial in
that block group. Injectivity of that particular port map now proves
that $\gamma$ is nullhomotopic in the seam surface itself. This step uses
a disk \emph{map}, not an embedded compression disk; the embedded-disk
criterion alone would not justify it.

Choose a narrow annulus just outside $C$ in the domain. Its image lies
in the seam collar on the outside side of $D$. Projection along the
collar connects its outer boundary loop to $\gamma$. Replace the map on
$D$ together with this annulus by a disk extension lying entirely on
that outside collar side, using the seam nullhomotopy pushed a small
positive distance off the seam. Keep the map fixed on the annulus's
outer boundary. The collar misses every other seam, so the replacement
removes $C$ and introduces no new inverse-image circle. Relative
smoothing can be performed away from the seams. The new disk need not
be homotopic to the old disk relative to its boundary: only existence
of a disk with that same boundary loop is required.

The finite circle count decreases. After repeating, the disk misses all
seams and lifts to the original block, as in the zero-circle case.
For a self-pairing, its two cut-boundary ports remain distinct throughout
the argument, even if the inside and outside block names coincide.
Thus the proof covers self-pairings as well as edges between different
blocks, and establishes the asserted injection.

\paragraph{Seam and external-boundary consequences.}
Push a seam torus into either adjacent block collar. Its map to $X$ is
homotopic to the port map followed by $q_v$, so its induced map is
injective by the lemma. For a retained external torus $T\subset\partial
B_v$, the same conclusion holds \emph{provided} its inclusion into
$B_v$ is injective. These conclusions concern injection into the whole
reassembled component, not just into the adjacent blocks.

In Section~\ref{sec:graph-prime-normal-form}'s good-block terminal case,
R05 supplies every required port and external-boundary injection.
Together with the written irreducibility proof there, the surviving
seams therefore give a finite essential-torus cutting of that component
into its recorded Seifert blocks. Keeping extra parallel tori is permitted
by E1; canonical minimality is not needed. This is a conditional
mathematical consequence of the supplied relative normal-form witnesses.
It does not declare the original raw seams incompressible. A terminal
solid torus, for example, has compressible external boundary and does
not meet the good-block hypothesis. Closed exceptional factors are
handled by their separate ledger, with no seam asserted merely from
a name or a group classification. G09's complete model metrics and the
relative Lean implementation remain separate.

\paragraph{R04: capping sphere sides preserves the actual group map.}
Let $c_a:R_a\to P_a$ attach finitely many three-balls to the new sphere
sides of one connected punctured component. For a single attachment,
use a collared open cover: one set deformation retracts to the old
component, the other to the cap ball, and the intersection to $S^2$.
The sets and intersection are path connected. The ball and the
intersection are simply connected, so van Kampen identifies the
\emph{inclusion-induced} map with an isomorphism. Iterating over the
finite caps gives
\[
                (c_a)_*:\pi_1(R_a)\ \cong\ \pi_1(P_a).
\]
Keep a basepoint in the punctured core and record the paths used when
applying van Kampen at a cap collar. This is the application of Theorem
1.20, printed 43 with proof on 44--46/PDF 23--26 in the inspected
Hatcher chapter. Attaching a solid torus along a torus has a different
intersection group and is not covered by this argument.

\paragraph{R06: injection through sphere reconstruction.}
Assume the R02 data give the original region $W$ by reassembling the
actual punctured components along a finite paired sphere system, with
collars and its reconstruction diffeomorphism. Each sphere port has
trivial fundamental group, so its map to either adjacent component is
injective without an irreducibility hypothesis. Apply the finite gluing
lemma to these sphere ports. It gives injection of
$(j_a)_*:\pi_1(R_a)\to\pi_1(W)$, including sphere self-pairings.
The reconstruction map may be represented using a shrunken embedded
core as above. Nonseparating sphere edges are retained; they are not
deleted by capping or mistaken for empty factors.

Consequently the group map already specified for R06 is injective:
\[
 \kappa_a=(j_a)_*\circ((c_a)_*)^{-1}:
             \pi_1(P_a)\longrightarrow\pi_1(W).
\]
This is a group map, not an asserted embedding of $P_a$ into $W$.
Given a retained torus map $i:T\to R_a$, functoriality yields, with
compatible basepoints,
\[
                 (j_a i)_*=\kappa_a\circ(c_a i)_*.
\]
Thus injection of $T$ into the capped factor implies injection of its
reconstructed image into $W$. Choices of connecting paths change the
comparison by explicit conjugations, which preserve injectivity; they
do not authorize replacement of the geometric maps by arbitrary group
isomorphisms.

This conclusion requires the finite torus system to have been moved off
all cap balls simultaneously, so that the maps $i$ actually exist and
remain disjoint with the required collar markings. The group argument
neither constructs that isotopy nor supplies a relative prime-factor
comparison. The next subsection supplies the cap-avoidance producer at the written
level; relative prime-factor comparison remains an R06/J02 prerequisite.

\paragraph{The outer carrier and the next implementation boundary.}
For the actual ambient map $k:W\to M$, the remaining map is
\[
                 (k j_a i)_*=k_*\circ\kappa_a\circ(c_a i)_*.
\]
Injection into $W$ alone does not prove injection into $M$. An injection
witness for $k_*$ suffices; alternatively prove this torus-specific
composite injective. One possible sufficient condition is an actual
finite collared gluing description of $M$ containing $W$ as a block,
with \emph{all} paired port maps injective on both sides. The same lemma
then supplies $k_*$. Do not assume that such a description or its
hypotheses are available for a thin region. Surgery-history transport
remains another G05 obligation, governed by its actual reconstruction
maps and the release gate.

The next subsection constructs simultaneous cap avoidance with retained
boundary data and composes it with these group maps, conditional on the
supplied relative prime normal form. J02/R07's
comparison with separately specified prime factors, A28's original
raw-cut application and the full formal J05 contract remain open.
F04's general bundle exercise and H1--H4 are unchanged. No unfinished-PC
adapter or probe is expanded by this argument.

Required regression specifications include the zero-seam case; an edge
between two blocks; two paired ports of one block; a nullhomotopy meeting
several seams; a seam-nullhomotopy that is not embedded; a sphere
self-pairing; and compatible basepoint changes. Reject replacing a disk
map by an embedded-disk premise, deducing asphericity from this lemma,
using torus filling as sphere capping, omitting the cap-avoidance witness,
or inferring injection into $M$ solely from injection into $W$. These
geometric cases are not executed Lean tests. The inventory rejection
controls only check that their required hypotheses and evidence limits
remain recorded.
% END migrated B051

% BEGIN migrated B052
\section{Simultaneous cap avoidance and the relative torus output}
\label{sec:cap-avoidance-output}

Revision 54 constructs the missing cap-avoidance witness at the written
mathematical level. It then composes the actual torus maps with
Section~\ref{sec:gluing-injection-transport}'s group comparison. The
geometric Lean implementations remain open, and the required smooth
isotopy interfaces must be bound to the audited foundation at release.
This is an application of that foundation, not a second construction of
smooth topology or a claim that its PC declarations are missing.

\paragraph{Source scope and a dimensional qualification.}
The local Hatcher Lemma 1.3 proof and connected-sum discussion, printed
3--4/PDF 4--5, use isotopy extension and moving submanifolds off cap
balls. Martelli's smooth isotopy and collar conventions and Theorems
1.1.10 and 1.1.14, printed 12--14/PDF 20--22, give the compact isotopy
extension and single-ball statements. To check the support needed for a
relative construction, we inspected \cite{Hirsch1994Scan}, Chapter 8:
Section 1, printed 177--180/PDF 95--97, and Section 3, printed 185--186/PDF
99--100 of the two-page-per-image scan. Its copyright page identifies
the corrected fifth printing, 1994. Theorems 1.2--1.4 explain compact
support; Theorem 3.1 gives the equioriented disk isotopy, with a supplied
proof; Theorem 3.2 supplies the two-disk induction pattern.

We use that pattern only in dimension three, proving the connectedness
of the punctured carrier at each step. The printed Theorem 3.2 does not
state the dimension restriction needed for its unrestricted labelled-pair
claim: in the oriented line, exchanging two disjoint intervals while
preserving both parametrizations' orientations cannot be achieved by a
diffeomorphism isotopic to the identity. This is our scope observation,
not an author-issued correction. No all-dimensional version is imported.
The publisher page and the available third-party correction list were
checked; that list covers Chapters 1--7, not Chapter 8. Exact provenance,
proof passages and correction limits are in
\code{reference_checks_revision54.md}.

\paragraph{Data and the boundary that is fixed.}
Let $P$ be a connected compact oriented smooth three-manifold. Suppose
$B_1,\ldots,B_m$ are pairwise disjoint smoothly embedded closed
three-balls in its interior, with their cap parametrizations, and put
\[
           R=P\setminus\bigcup_{i=1}^{m}\operatorname{int}(B_i).
\]
The balls are supplied caps, not balls inferred from an arbitrary
embedded sphere. Let $t_\lambda:T_\lambda\to\operatorname{int}(P)$ be a
finite disjoint family of embedded two-sided tori with specified disjoint
collars. No incompressibility is needed for the avoidance lemma.

Restrict the given external-boundary collars to a sufficiently small
closed width and call their union $K$. Choose this width so that $K$
is disjoint from the caps and torus collars and $U=P\setminus K$ is
connected. Such a choice exists by compactness, finiteness and collar
shrinking; the interior of the remaining collar core is connected.
For closed $P$ take $K$ empty. Shrink the torus collar widths if needed
so that their closed union has nonempty open complement in $U$.
Record these positive width restrictions and the unchanged boundary
parametrizations. The conclusion fixes this recorded smaller collar
pointwise; it is not a promise to fix an arbitrary larger region that
could obstruct the motion.

Choose $m$ disjoint small parametrized closed balls $D_i$ in a coordinate
ball in that open complement. Their parametrizations have the same
orientation sign as the corresponding $B_i$ parametrizations. They avoid
the whole retained torus collars, not just their middle surfaces. This
choice uses only that finitely many two-dimensional surfaces and a
sufficiently small collar system leave room in a three-manifold; no
metric estimate or curvature normalization is involved.

\paragraph{The simultaneous ball motion.}
A connected three-manifold without boundary remains connected after
removing finitely many disjoint embedded closed balls in its interior.
For this elementary fact, take a path between two remaining points,
make it transverse to the finitely many ball boundaries, and replace
its excursions through each ball by paths on a slightly larger parallel
sphere in the exterior collar. Such spheres are connected. Choose the
exterior collars disjoint and away from the endpoints. The resulting
path avoids the closed balls. This is the specific dimension-three fact
needed in the induction below; no irreducibility premise is used.

Induct over the labelled caps. Suppose an ambient isotopy has already
sent the first $r$ parametrized balls to $D_1,\ldots,D_r$. Its current
images of the remaining balls stay disjoint from those targets, since
the same ambient diffeomorphism acts on the entire family. Work in
\[
                   U_r=U\setminus\bigcup_{i=1}^{r}D_i.
\]
This is a connected oriented manifold without boundary by the preceding
path argument. Both the current $(r+1)$st ball and its target lie in
$U_r$. Apply the compact-support part of Hirsch's single-disk theorem
there. Extend its ambient isotopy by the identity over the complement
of $U_r$. Compact support inside $U_r$ makes this extension smooth and
fixes a neighborhood of $K$ and every earlier target throughout the
isotopy. Later balls may move, but remain a disjoint family.

Concatenate the finitely many isotopies after smooth time
reparametrizations constant near their endpoints. Their supports have
compact union inside $U$. We obtain an orientation-preserving ambient
isotopy $H_t$ of $P$, starting at the identity, fixed near $K$, such that
$H_1(B_i)=D_i$ for every label, with the chosen parametrizations. If
$m=0$, use the identity isotopy. The proof does not require that the
original cap positions be disjoint from the target positions.

Set $h=H_1^{-1}$ and use the inverse isotopy $H_t^{-1}$. For each torus
and its restricted collar $e_\lambda$, set
\[
                 t'_\lambda=h\circ t_\lambda,
                 \qquad e'_\lambda=h\circ e_\lambda.
\]
These are simultaneously disjoint from every original cap ball: a point
of $B_i\cap h(\operatorname{im}e_\lambda)$ would map under $H_1$ to a
point of $D_i\cap\operatorname{im}e_\lambda$, contrary to the choice of
$D_i$. Thus the entire moved torus collars lie in the interior of $R$.
One ambient diffeomorphism preserves all disjointness, two-sidedness,
orientations and labels. External boundary collars are fixed pointwise.

The isotopy takes place in $P$. It need not preserve $R$ or any original
cap ball at intermediate times, and it is not asserted to extend to an
ambient isotopy of the original region $W$. Only the final torus collars
are being restricted to $R$. Fixing the caps pointwise while removing
an existing intersection with them would be an invalid stronger claim.

\paragraph{The actual maps and injection into the reconstructed region.}
Apply the construction separately in each supplied capped factor $P_a$,
with cap inclusion $c_a:R_a\to P_a$ and sphere reconstruction map
$j_a:R_a\to W$. Restrict $h_a t_{a\lambda}$ to its image in the punctured
component to obtain $i_{a\lambda}:T_{a\lambda}\to R_a$. The defining
identity and the reconstructed torus map are
\[
          c_a i_{a\lambda}=h_a t_{a\lambda},\qquad
          u_{a\lambda}=j_a i_{a\lambda}.
\]
The collars avoid the sphere ports, so reconstruction embeds them into
$W$. Distinct punctured-component interiors have disjoint images; within
each component the common $h_a$ preserves disjointness. The finite
label set is the disjoint union of the original factor label sets, with
no new torus created and none silently discarded. Self-paired sphere
ports cause no identification on these interior collars.

For compatible basepoints, the group-map identity is now
\[
 (u_{a\lambda})_*
       =\kappa_a\circ(h_a)_*\circ(t_{a\lambda})_*,
 \qquad
 \kappa_a=(j_a)_*\circ((c_a)_*)^{-1}.
\]
Section~\ref{sec:gluing-injection-transport} supplies injection of
$\kappa_a$, and $h_a$ induces an isomorphism. Hence every torus known
injective in $P_a$ has an injective reconstructed image in $W$. Record
the basepoint tracks under the isotopy and the connecting paths used in
reconstruction. They yield the required change-of-basepoint maps; do not
replace $(h_a)_*$ by a literal identity at a fixed basepoint merely
because $h_a$ is isotopic to the identity. For $k:W\to M$, injection of
$k_*\circ(u_{a\lambda})_*$ still needs its separate witness.

\paragraph{Cut pieces, fibrations and cap assignments.}
Suppose the original tori cut $P_a$ into the specified Seifert carriers
$Q_v$. Transport the cuts and their collars by $h_a$. The induced
maps $\widetilde h_v:Q_v\to Q'_v$ are actual diffeomorphisms of the cut
carriers. Transport the Seifert fibration by composing its projection
with $\widetilde h_v^{-1}$. In particular, transport fiber slopes and
external boundary markings through these maps. If an old port pairing
is $\phi:b\to b'$, the new pairing is
\[
       \phi'=\widetilde h_{b'}\circ\phi\circ\widetilde h_b^{-1}.
\]
This applies to paired ports of the same piece as well. It preserves the
pairing identities and opposite induced boundary orientations.

Every cap ball is connected and avoids the new cut collars, so it lies
in exactly one $Q'_v$. Store this cap-to-piece assignment, writing $I(v)$ for the labels
assigned to $Q'_v$. Cutting $R_a$
along the moved tori gives the carriers
\[
       Q'_v\setminus
       \bigcup_{i\in I(v)}
                    \operatorname{int}(B_i).
\]
Reattaching the same assigned caps recovers $Q'_v$ with its transported
Seifert structure. The punctured carrier can have sphere boundary and
is not thereby a Seifert piece in the required torus-boundary sense.
Regluing sphere ports in $W$ can join several such punctured carriers;
no conclusion that $W$ cut only along these tori is Seifert follows.
Retain the prime/sphere ledger alongside the torus ledger. The endpoint
cuts live on the capped prime factors; the maps into $W$ serve the
relative reconstruction and incompressibility obligations. Sphere
products and other previously isolated exceptional factors keep their
separate records.

\paragraph{Implementation targets and remaining comparison.}
R06 now has a written producer for simultaneous cap avoidance and a
written torus-family output into $W$, conditional on the supplied relative
prime reconstruction and the torus injection witnesses in its factors.
Its concrete obligations are: admissible smaller collars; the finite
compact-support ball isotopy; inverse transport of the entire collar
family; actual cut-map and fibration transport; the cap assignment; and
composition with the sphere-reconstruction maps. Reuse completed
foundation results for each smooth interface. This does not implement
R02/J02's relative prime construction or identify it with a separately
specified decomposition.

Revision 56 compares the outer map $W\to M$ in
Section~\ref{sec:ambient-history}. The actual late decomposition has a
written conditional proof using the separate cusp-injection input;
arbitrary thin-region inclusions are not thereby injective. History
transport, G09's model metrics, A28's original raw-cut application,
F04 and H1--H4 remain open.
No PC-specific adapter or probe is expanded.

Regression specifications include no caps, no tori, several labelled
caps, self-paired sphere ports, a torus meeting an original cap, and a
shrunk fixed external collar. Reject inconsistent orientation signs,
noncompact support across a protected boundary, independent isotopies
that lose joint disjointness, an isotopy of $W$ inferred from one of
$P_a$, a fixed-basepoint identity without its path, or a Seifert claim
for a sphere-punctured piece. The line-interval counterexample guards
against an unrestricted all-dimensional multiple-disk rule. These are
geometric specifications, not executed Lean tests.
% END migrated B052

% BEGIN migrated B053
\section{Actual late-slice ambient injection and finite history}
\label{sec:ambient-history}

Revision 56 resolves the next written comparison for the actual late
thick/thin decomposition. The result is conditional on a separate global
cusp-incompressibility theorem and the existing gluing and cap-avoidance
arguments. It is neither an implementation nor a blanket injection rule
for thin subsets. The five targets AT01--AT05 are recorded in
\code{ambient_transport_contracts.csv}; source records L33--L35 and L39
identify the passages actually read.

\paragraph{The global input is not part of local collapse.}
In the archived February 20, 2013 version of Kleiner--Lott's 2008 notes,
Proposition 91.2 and Lemma 91.4, printed pages 2817--2819 (PDF pages
231--233), address the homomorphism from an actual cusp torus into the
late post-surgery slice. The remainder of Section 91, through printed
page 2823 (PDF page 237), uses least-area disks, neck avoidance, boundary
estimates, an area differential inequality and a contradiction. Its
boundary estimate refers to Hamilton's detailed argument. These are
substantial global flow inputs, not consequences of the definition of a
hyperbolic cusp. The proof is source-read; its analytic interfaces are
not thereby verified against PC.

Theorem 17.1(3) of the local-collapse monograph uses this ambient
incompressibility in the application to flow (printed pages 90--91,
PDF pages 85--86). The nearby reference to Proposition 90.1 of the notes
must be supplemented by Proposition 91.2 for this particular conclusion.
The local Riemannian Theorem 16.1 does not supply it. In the archived
2012 Morgan--Tian body, Proposition 2.25 (printed page 29, PDF page 33)
is likewise global; Proposition 5.4 (printed pages 55--56, PDF pages
59--60) \emph{assumes} ambient incompressibility. Its graph-manifold
hypothesis cannot be used to discharge that assumption. The 2014
published-body crosswalk remains unverified.

MSM206's TeX proof labelled \code{lmm bdry incompressible}, lines
16236--16294, proves peripheral injection in the hyperbolic model, using
a component of the lifted cusp boundary. This is a useful static
infrastructure target. It does not prove injection of the cusp's image
in an ambient flow slice. All three carriers must remain named.

\paragraph{AT01--AT02: ports first, then whole carriers.}
Let $N=M_t^+$ be a selected late slice and let $H_i$ be the supplied
finite disjoint compact hyperbolic cores. Put
\[
 W=\overline{N\setminus\bigcup_i\operatorname{int}(H_i)}.
\]
Use the actual embedded seam tori $T_e$, collars and pairings which
reconstruct $N$ from the cores and components of $W$. Work in each
connected component of $N$. AT01 supplies injection of every seam map
$s_e:T_e\to N$ on fundamental groups.

For either port $b:T_e\to V$ of a seam and its actual carrier inclusion
$k_V:V\to N$, the composite $k_Vb$ is $s_e$. Include the paths relating
the chosen basepoints. Since $(k_V)_*b_*$ is injective, $b_*$ is
injective. This cancellation does \emph{not} prove that $(k_V)_*$ is
injective. Apply instead the finite collared gluing argument of
Section~\ref{sec:gluing-injection-transport}, now with every paired
port injective, to obtain injection of each actual carrier map
$(k_V)_*$. Its simultaneous finite-seam formulation includes multiple
seams and cycles. MSM206's separating-surface van Kampen argument,
labelled \code{lmm m12 incomp}, lines 16167--16228, is a comparison for
the separating case; it is not an induction that silently assumes every
later seam separates.

If there are no hyperbolic cores in a connected slice, $W=N$ and the
map is the identity. A closed hyperbolic component has no cusp ports
and is treated as its own component. Neither case calls for artificial
torus boundaries or a nonempty seam hypothesis.

% BEGIN REVISION167 ACTUAL AT01 BINDING
\begin{corollary}[Actual hyperbolic--thin seam injections; IMS11]
\label{cor:ambient-actual-at01-binding}
For the actual TCF01/TCF07 late decomposition, every seam inclusion
$s_e:T_e\to N=M_t^+$ is injective on fundamental groups. The injections
are attached to these actual maps, collars and basepoint paths. Hence
AT01 is supplied as written mathematics under IMS10's retained
foundation, and AT02 yields injection of the actual core and thin
carrier inclusions into their ambient slice components.
\end{corollary}
\begin{proof}
IMS10 first proves injection for fixed reference cusp tori and then
transports it along the actual cusp product to each TCF01 cut torus.
At the common time chosen in TCF07, these are exactly the boundary
maps used to define $W$ and its labeled gluings. For either port,
its composition with its carrier inclusion is the corresponding
$s_e$, up to the recorded basepoint change, so that port is injective.
Apply the finite collared gluing theorem to all ports simultaneously
to get injection of the carrier maps, including cycles in the gluing
graph. This last inference uses the normal-form theorem; it is not
cancellation of the carrier map from the composite. For a slice
component with no cusp seams, use the identity inclusion of its
whole carrier. Positive constant metric rescaling does not change
any of these maps or fundamental-group homomorphisms.
\end{proof}
This binding does not place the carriers in prime factors or reconstruct
the initial manifold. Those subsequent AT interfaces retain their own
obligations. It also does not claim an accepted Lean implementation of
Meeks--Yau, the Loop Theorem or the minimal-surface spectral foundation.
% END REVISION167 ACTUAL AT01 BINDING

\paragraph{AT03: compose the actual retained-torus maps.}
For a retained factor torus from
Section~\ref{sec:cap-avoidance-output}, the map into $W$ satisfies
\[
 (u_{a\lambda})_*
   =\kappa_a\circ(h_a)_*\circ(t_{a\lambda})_*.
\]
Consequently the actual map into the late slice has
\[
 (k_Wu_{a\lambda})_*
   =(k_W)_*\circ\kappa_a\circ(h_a)_*\circ(t_{a\lambda})_*.
\]
Each factor is injective under the stated hypotheses. This is the
written outer-composite proof for the supplied late decomposition.
Retain the paths for changes of basepoint, the common isotopy of the
capped factor, actual collar transport, and all cap-to-piece labels.
There is still no embedding of a capped prime factor inferred from its
punctured reconstruction, and a sphere-punctured Seifert carrier is
still not an endpoint Seifert piece. The prime/sphere ledger remains
part of the output. Construction and comparison of that prime ledger
(R02/J02/R07), and the geometric Lean bindings, remain open.

\paragraph{AT04: reconstruct across a finite history.}
Lemma 73.4 of the notes (printed pages 2761--2762, PDF pages 175--176),
with Lemmas 67.5 and 67.13 and their supplied proofs, reconstructs the
pre-surgery diffeomorphism type using finitely many operations: connected
sums of remaining components, connected sums with $S^1\times S^2$ or
$\mathbb{RP}^3$, and extra components of the sphere-product or spherical
quotient types. In particular $\mathbb{RP}^3\#\mathbb{RP}^3$ contributes
two factors. Remark 73.3 explicitly retains some nonnegative-scalar
components that Perelman's convention discards. An event ledger must
translate the chosen convention rather than silently equating them.

Select a finite time prefix with its audited smooth-stage identifications
and surgery events. Iterate the event reconstruction, retaining
orientation, parametrized caps, sphere ports and exceptional factors.
The cited diffeomorphism-type statement motivates this producer; its
upgrade to the required oriented, labelled reconstruction is an adapter
obligation. Do not posit a global reverse-time map $M_t\to M_0$, or an
embedding of each later capped factor in the original manifold.

Assemble the endpoint from certified late factors and the recorded
exceptional factors by the actual connected-sum reconstruction. Use
cap avoidance and torus transport where these representatives need to
be compared. If the public endpoint has a prescribed prime decomposition,
R02/J02/R07 must compare the factors and carry their markings. Neither
ambient injection in $N$ nor an abstract connected-sum identity supplies
that comparison automatically. No new endpoint geometry or production
proof is accepted here.

\paragraph{AT05: explicit terminal data.}
The terminal branch can use an explicitly empty terminal slice, a finite
reconstruction ledger and certified exceptional models to recover the
original manifold without a late hyperbolic/thin contribution. The mere
failure of arbitrarily late nonempty slices supplies none of these data.
The producer of \code{LongTimeOutcome} must provide the witnesses.
Finite-time extinction remains outside the default route; this branch
consumes explicit terminal data if it is produced.

\paragraph{Scope and next work.}
The outer-ambient comparison is now written for the actual sourced late
decomposition. It remains conditional on cusp injection, the previously
written finite gluing and relative prime construction. H1--H4, the
full graph refinement, model metrics and the history implementation
remain open. No PC adapter or declaration probe has been expanded.
The next independent implementation candidate and the topic-by-topic
blueprint program are given in Section~\ref{sec:independent-infrastructure}.
% END migrated B053

\section{Marked carriers and naturality of the actual maps}
\label{sec:transport-marked-carriers}

Revision 102 completes this chapter's blueprint coverage. The preceding
three sections are retained written arguments; their historical ``next''
paragraphs are superseded by the current audit at the end of this chapter.
The new targets refine AT01--AT05, rather than replace the established
manifold wrappers or commission another differential-topology foundation.
The source record is \code{reference_checks_revision102.md}; the full
contract and review inventories are \code{ambient_transport_contracts.csv}
and \code{chapter15_audit_revision102.csv}.

\begin{definition}[Marked carrier and reconstruction data; AT06]
\label{def:transport-marked-carrier}
A carrier is a compact oriented smooth three-manifold, allowed to be a
finite disjoint union. Its markings consist of finite boundary-port labels,
the actual parametrized boundary embeddings and disjoint collars, a
fixed-point-free partial pairing involution, and orientation-reversing
pairing diffeomorphisms. Each port is used at most once; sphere ports,
torus ports and retained external ports have distinct types. A paired
port is different from its mate even in a self-pairing. A reconstruction
includes a diffeomorphism from the collared quotient to the specified
ambient carrier. Interior decorations are finite disjoint embedded spheres
or tori with disjoint collars, their labels and any attached structures.
A morphism records component/port bijections and actual orientation-preserving
diffeomorphisms with commuting marking squares. Being relative to a
protected external collar means pointwise equality there, not merely a
permutation of boundary components.
\end{definition}
All collars may be restricted to recorded smaller positive widths. A
homeomorphism alone does not satisfy this smooth contract. Connected
components and empty carriers are handled before choosing basepoints.
No definition imports a Ricci-flow estimate or a final geometric certificate.

\begin{lemma}[Gluing along commuting collar maps; AT07]
\label{lem:transport-gluing-naturality}
Suppose two finite AT06 diagrams have bijections of vertices and ports
and orientation-preserving diffeomorphisms $f_v$ of their blocks.
Assume these maps have product
form on the chosen collars, with boundary maps $f_b$, and every pairing
satisfies
\[
                 \phi'_{b}=f_{\bar b}\phi_b f_b^{-1}.
\]
Then the $f_v$ induce an orientation-preserving diffeomorphism of the
collared quotients, commuting with every reconstruction map. It is
relative to any protected collars on which the given maps are identities.
\end{lemma}
\begin{proof}
The displayed equation sends each generating identified pair to an
identified pair. Thus the block maps descend to the quotient; their
inverses descend by the inverse equation. The resulting inverse maps are
smooth in interior charts and in the product collar charts across each
seam. Orientation agrees across a seam because both pairings reverse
induced boundary orientation. Finite disjoint collars permit these checks
simultaneously, including loops and parallel edges. Compose with the
supplied ambient reconstruction diffeomorphisms.
\end{proof}
Without product collar form, an inherited collar-straightening theorem
must first supply a compatible representative; equality only on boundary
points is insufficient to glue arbitrary derivatives smoothly. Hirsch
Chapter 8, Section 2, printed 184--185/PDF 99, supplies the source comparison
for smooth gluing and isotopic attaching maps. It does not make arbitrary
boundary markings extendible.

\begin{lemma}[Cut, cap and structure transport; AT08]
\label{lem:transport-cut-naturality}
Let $f:X\to Y$ be a supplied oriented diffeomorphism and $e$ a finite
collared surface system in $X$. Use $f e$ as the target collar system.
There is an induced diffeomorphism of the cut carriers, with both side
labels retained. Sphere caps can be transported using the same
parametrized balls and transported attaching maps. Seifert projections
transport as $\pi'=\pi\widetilde f^{-1}$ on the actual cut pieces.
\end{lemma}
\begin{proof}
Outside the collars use $f$; in each split collar use the same product
coordinates and retain the sign of the cut side. These formulas agree on
the overlaps and have the corresponding inverse. Attach the cap balls
by conjugating their attaching maps, using identity on their abstract
parametrizing balls. AT07 gives the capped equivalence. The displayed
projection has exactly the transported fibers and local Seifert models.
Pairings, fiber slopes and orientations are transported by these maps,
not by an abstract equality of piece names.
\end{proof}
A target cap parametrization fixed in advance requires a supplied smooth
ball-extension/marking comparison. It is not an extra conclusion of AT08.
The construction preserves every external label; a punctured Seifert
piece is still only a punctured piece.

\begin{lemma}[Based composition and kernel cancellation; AT09]
\label{lem:transport-based-maps}
For composable maps of connected pointed carriers, use the actual induced
fundamental-group homomorphisms and record each basepoint path. Composition
with an isomorphism preserves injectivity; if $b a$ is injective then $a$
is injective, but $b$ need not be. If $b$ is injective, then
$\ker(b a)=\ker(a)$. A homotopy gives equality of the induced maps after
the change-of-basepoint isomorphism along its basepoint track.
\end{lemma}
\begin{proof}
The first assertions follow by applying injectivity to equal images.
For the last assertion, the homotopy of a loop is a square; its other two
edges are the basepoint track and its reverse. Its boundary gives the
claimed conjugation. Such conjugations are isomorphisms. These are
applications of the inherited fundamental-group functor and homotopy
interface, not a new construction of that functor.
\end{proof}
This isolates the logical distinction used in AT02: seam-to-slice injection
first gives port injection; the finite gluing argument, not cancellation,
then gives carrier injection. An isotopy does not induce the literal
identity at a fixed basepoint without its track.

\section{Smooth stages, common cores and one selected late slice}
\label{sec:transport-stages}

\begin{lemma}[Transport on a supplied smooth stage; AT10]
\label{lem:transport-smooth-stage}
Suppose a compact carrier $X$ parametrizes a surgery-free stage through
supplied orientation-preserving diffeomorphisms $F_t:X\to M_t$. Then
$F_bF_a^{-1}:M_a\to M_b$ transports all AT06 markings by AT08 and
preserves group injections by AT09. For transport only on an open
unscathed region, restrict the conclusion to that region. A flow
construction from a time-dependent field must provide existence of all
required trajectories on the whole closed time interval.
\end{lemma}
\begin{proof}
The inverse is $F_aF_b^{-1}$, and composition gives the cocycle identity
at intermediate times. Apply AT08 to each marking. In a field-based
construction, inherited existence/uniqueness and a common compact
trajectory buffer, or another actual completeness hypothesis, supply the
maps. Compactness separately at each time is not that hypothesis.
\end{proof}
Hirsch Chapter 8, Section 1, printed 178--180/PDF 96--97, explicitly uses
the support of the entire time-dependent field. LC82/FC34's common
compact isolation and all boundary-face conditions remain necessary for
their level-set applications. None of this produces a diffeomorphism
across a surgery event.

\begin{lemma}[Kernel preservation through a common core; AT11]
\label{lem:transport-common-core}
Let $C$ be a connected common core, $j_-:C\to X_-$ and $j_+:C\to X_+$
actual maps inducing injections, and $a:T\to C$ the same torus map.
After recording basepoint paths,
\[
       \ker((j_-a)_*)=\ker(a_*)=\ker((j_+a)_*).
\]
The same holds when the two core maps of $T$ are homotopic with their
tracks recorded. It does not follow that either $j_\pm$ is surjective,
or that there is an ambient reverse map $X_+\to X_-$.
\end{lemma}
\begin{proof}
Apply AT09 to each injective outer map and then to the homotopy.
\end{proof}
This is the algebraic operation in KL Lemma 91.4, printed 2818--2819/PDF
232--233. Its geometric core inclusions and the unscathed torus still
need actual surgery data. It does not prove the global cusp-injection
producer AT01, whose Section 91 argument also uses minimal disks, neck
avoidance and Hamilton's boundary estimates in Chapter 8's analytic route.

\begin{definition}[Exact static-to-late input; AT12]
\label{def:transport-static-input}
Fix a single admissible late slice $N$ and its metric normalization.
The input consists of actual disjoint compact hyperbolic cores $H_i$
with embeddings and cusp collars, the actual compact complement $W$,
and their collared reconstruction of $N$. Each connected component of
$W$ has the exact FC45 static graph certificate, with its external
boundary labels identified with these same seams. If a chosen static
representative differs, include its oriented diffeomorphism to $W$ and
transport by AT07--AT08. AT01 supplies seam-to-$N$ injection separately.
\end{definition}
LC90 supplies the flow-to-static application only after Chapter 12's
hypotheses, including H1--H4, hold. Neither graphness nor model peripheral
injection supplies AT01. Positive constant metric rescaling changes none
of these underlying maps or group injections; it does not preserve a
numerical curvature or volume normalization without its separate scaling
formula. Empty hyperbolic families give $W=N$; closed hyperbolic components
have no artificial cusp ports. Empty $W$ and disconnected $N$ are allowed.

\begin{lemma}[A common selected time for finite eventual requirements; AT29]
\label{lem:transport-common-late-time}
Fix the truncation/collapse parameters in their required order. Suppose
there are finitely many predicates $P_i(t)$ with supplied thresholds
$T_i$ such that every admissible regular $t\ge T_i$ satisfies $P_i(t)$.
If admissible nonempty regular slices occur arbitrarily late, one slice
satisfies all $P_i$ simultaneously. At that chosen time supply the actual
finite component, seam and surgery-prefix data.
\end{lemma}
\begin{proof}
Choose an admissible $t$ beyond the maximum of the finitely many thresholds
and any requested lower bound. For an empty predicate family use just
that lower bound. The input availability supplies $t$, and every predicate
follows from its threshold. This elementary selection neither proves the
thresholds nor changes their parameter dependence.
\end{proof}
In particular, ``for every sufficiently small $w$, for all sufficiently
large $t$ depending on $w$'' does not give one threshold for every $w$.
An infinite list of derivative orders is not a finite maximum; retain the
already chosen finite order and uniform source function. A merely discrete
set of times need not be finite in a compact interval. A supplied finite
prefix must cover every actual event up to $t$, or a proved local-finiteness
statement must produce such a list. No global eventual absence of surgery
is assumed. Perelman's Section 7 opening explicitly retracts the proposed
argument for that stronger conclusion.

\section{Relative prime pieces and the late topological certificate}
\label{sec:transport-relative-primes}

\begin{foundation}[Relative prime and category input; AT13]
\label{found:transport-relative-primes}
From the exact raw graph carrier $W$, Chapters 5--6 must produce a finite
interior sphere system disjoint from its retained torus collars, the
actual punctured components $R_a$, cap inclusions $c_a:R_a\to P_a$, and
an oriented collared reconstruction of $W$. Each capped component is
prime or a recorded $S^3$ identity; the sphere ledger retains its
nonseparating edges. Graph presentations on these same $P_a$, and all
external boundary correspondences, are part of the relative R02/R07/J02
output, not consequences of an unmarked factor-type list.
\end{foundation}
Hatcher Theorem 1.5, printed 5--8/PDF 6--9, supplies oriented prime
existence/uniqueness; Matveev second edition Proposition 2.4.3 and
Corollary 2.4.4, printed 85--87/PDF 97--99, supply the graph comparison.
Their implementation in this relative smooth category is still qualified.
The caps are actual balls, not a new smooth Schoenflies gap. R01--R07 and
J01--J05 retain their existing formal status.

\begin{lemma}[Cap avoidance for the whole decoration; AT14]
\label{lem:transport-decoration-avoidance}
In the cap-avoidance construction of Section~\ref{sec:cap-avoidance-output},
replace the torus family by any finite disjoint family of smooth embedded
closed spheres and tori, with collars in the interior. After restricting
collar widths, one orientation-preserving ambient isotopy of the capped
carrier, fixed near the recorded smaller external collar, moves the whole
decoration off every supplied cap ball. It preserves all labels and the
joint disjointness of all collars.
\end{lemma}
\begin{proof}
A finite union of embedded surfaces has empty interior. Pick a point in
its complement and restrict the finitely many collar widths so a coordinate
ball about that point misses them. Choose equioriented target balls there.
The existing three-dimensional induction moves every labelled cap to its
target, fixing earlier targets by working in their connected complement.
Invert its final diffeomorphism. The same set-intersection argument now
moves every sphere and torus collar off the original caps simultaneously.
Transport pairings and structures by AT08. The proof requires no
connectedness of the complement of the surfaces: only the complement of
previously placed balls is used for the isotopy.
\end{proof}
The isotopy is in the capped carrier, may move the old caps during its
course, and need not be an isotopy of $W$. Compact support is common for
the finite concatenation. Hirsch's unrestricted multiple-disk statement
is not imported in all dimensions; the earlier line-interval qualification
remains. This extension is needed to flatten nested sphere reconstructions,
not just to move endpoint tori.

\begin{lemma}[Disjoint sphere operations commute with torus gluing; AT15]
\label{lem:transport-sphere-torus-commutation}
Suppose supplied finite sphere-cut/cap data are disjoint from every
retained torus collar used in a gluing. No prime or graph-manifold
hypothesis is needed for this operation lemma. Cutting these spheres and
capping before or after the torus gluings gives naturally diffeomorphic
capped carriers, with the same torus ports, orientations and external
markings. Later torus cuts disjoint from those caps have the analogous
commutation, with each cap assigned to its actual cut piece.
\end{lemma}
\begin{proof}
Away from the disjoint operation neighborhoods both constructions are the
same. On a sphere neighborhood use the identical cut and cap model, and
on a torus neighborhood use the identical gluing chart. These maps agree
on every overlap, so AT07 glues them with their inverses. A cap ball is
connected and misses the torus collars, hence belongs to exactly one
component of the cut complement. This yields the assignment. If an
already jointly disjoint surface decoration meets attachment cap balls,
AT14 moves that whole decoration off the balls first. It does not remove
arbitrary intersections between spheres and tori; the stated joint
disjointness is an input, not an extra conclusion.
\end{proof}
Thus the spheres originally in $W$ also form an actual sphere system in
$N$. The new capped carriers are obtained by gluing $H$ to the $P_a$
through their unchanged external torus ports. They are not all individual
$P_a$, and no embedding of a capped $P_a$ in $N$ is asserted.

\begin{foundation}[Relative good-block refinement input; AT16]
\label{found:transport-good-blocks}
Given AT13 and injections of the external ports of $W$, the Chapter 6
refinement must supply, on the same capped factors, finite essential torus
cuts with actual collared reconstruction. In a component with boundary,
all resulting blocks used in the gluing have irreducible carriers and
injective boundary ports; closed irreducible Seifert factors and closed
sphere-product or spherical exceptions have separate records. The inherited
hyperbolic-core input supplies irreducibility and peripheral injection on
its actual compact cores. Model realization is not included here.
\end{foundation}
The Chapter 6 good-block terminal argument and R05 supply the conditional
route, not an unconditional application of Matveev 2.4.7 to raw seams.
External port injection passes to $R_a$ by cancelling its actual map to
$W$, and to $P_a$ by the sphere-cap group isomorphism R04. A solid-torus
factor cannot carry such an injective external torus port. Noncanonical
extra parallel tori are allowed under E1. The source's optional replacement
of a torus by a Klein bottle in Morgan--Tian Proposition 5.4 is not used.
All relative existence/category bindings in AT13/AT16 remain qualified.

\begin{proposition}[Essential tori in the actual capped late carriers; AT17]
\label{prop:transport-essential-late-tori}
Assume AT01, AT12, AT13 and AT16. In each carrier obtained by AT15, take
the cusp seams and the supplied interior Seifert-factor tori. They form a
finite disjoint collared torus system, injective on fundamental groups in
that carrier. Its cut pieces are exactly the supplied hyperbolic and Seifert
blocks. Separately, AT03 gives the retained factor tori's injection into
$N$ after cap avoidance and reconstruction.
\end{proposition}
\begin{proof}
AT02 first proves injection of each actual $W$ component in $N$. The
port cancellation and cap isomorphisms just described give external port
injection in $P_a$. AT16 gives all other port injections. Apply the finite
collared gluing lemma to the final hyperbolic/Seifert blocks: their
reconstruction maps inject, so their seams do too. The disjoint collar
construction gives the claimed cut pieces, including self-pairings.
For the distinct assertion about actual maps into $N$, use AT03 with
its original cap-avoidance square. Neither assertion assumes the other
carrier is embedded as a capped submanifold of $N$.
\end{proof}

\begin{proposition}[Prime late factors from good blocks; AT18]
\label{prop:transport-prime-late-factors}
Under AT17, every closed connected carrier assembled entirely from the
irreducible good blocks is irreducible, hence prime, using the inherited
sphere/disk and collar interfaces. Include each separately supplied prime
exceptional component and each sphere-product contribution of the sphere
ledger. This gives a chosen prime decomposition of $N$ with reconstruction;
$S^3$ identities may be retained explicitly.
\end{proposition}
\begin{proof}
Apply the written good-block sphere-removal argument of
Section~\ref{sec:graph-prime-normal-form}. Put a sphere transverse to the
finite seams; an innermost circle bounds a disk in one cut block. Port
injection and the surface disk criterion give a disk on that seam; the
block's irreducibility supplies the ball parallelism. The whole-sphere
isotopy removes the circle and any other circles in that disk, and can
avoid external collars. Repeat to place the sphere in one block, where
it bounds a ball. No ambient irreducibility is assumed during this
argument. Zero seams and loop edges are covered. Irreducibility gives
primeness by the inherited connected-sum criterion. AT15 supplies the
sphere reconstruction, and AT20 below counts its sphere-product factors.
\end{proof}
This expands the topological step of Morgan--Tian Proposition 5.4,
printed 55--56/PDF 59--60. It does not use the stronger eventual
irreducibility result on printed 57/PDF 61, and it does not assume every
late component is already prime. The good-block proof is a conditional
application, not a fresh proof of the Chapter 6 producer or hyperbolic
irreducibility.

\section{Sphere reconstruction with labels and cycles}
\label{sec:transport-sphere-ledger}

\begin{definition}[Finite oriented sphere ledger; AT19]
\label{def:transport-sphere-ledger}
A sphere ledger has a finite multigraph with vertices labelled by connected
compact oriented capped carriers $P_v$. Each incident half-edge names a
distinct parametrized ball in the interior of its vertex; all such balls
are disjoint. Remove their interiors and pair the resulting sphere ports
by the recorded orientation-reversing diffeomorphisms and collars. Store
an oriented diffeomorphism from this reconstruction to the intended
carrier. Decorations avoid all ball closures, and external collars are
retained. Loops and parallel edges are allowed. Connected components of
the graph are identified with components of the reconstructed carrier.
\end{definition}
The graph is not required to be a tree. Each cycle has one representation:
an uncontracted graph cycle, or an explicit sphere-product vertex after
normalization, never both. An empty graph reconstructs the
empty carrier. The empty connected-sum list representing $S^3$ is a
separate convention: it does not turn an empty disjoint union into $S^3$.
Each deleted $S^3$ vertex needs its actual identity reconstruction, and
its ball/edge markings must first be accounted for.

\begin{lemma}[Sphere graph cycle count; AT20]
\label{lem:transport-sphere-cycle-count}
For a nonempty connected AT19 graph $G$, assume the inherited oriented
sphere-pair gluing and sphere-boundary extension interfaces. Its
reconstructed carrier is diffeomorphic to
\[
       \mathop{\#}_{v\in V(G)} P_v\ \#\
           \mathop{\#}^{\,|E(G)|-|V(G)|+1}(S^1\times S^2).
\]
The output includes the finite sequence of actual gluing identifications,
not only the count. Apply this formula componentwise for a disconnected
graph. The total number of cycle factors is $|E|-|V|+c(G)$.
\end{lemma}
\begin{proof}
Choose a spanning tree in each nonempty connected component. Glue its
edges one at a time, contracting each edge: its ends currently belong to
distinct carriers, so the operation is their oriented connected sum.
Transport every remaining cap marking through the actual equivalence.
There are $|V|-1$ tree edges. Each remaining edge then pairs two distinct
sphere ports in the same connected carrier, giving one $S^1\times S^2$
summand by the inherited nonseparating sphere-pair construction. There are
$|E|-|V|+1$ such edges, including any original loop. Orientation excludes
the twisted sphere bundle. Hirsch's gluing comparison and smooth extension
over a ball justify replacing an attaching parametrization only with its
actual comparison map. The count is independent of the tree, while the
chosen maps remain part of the witness. To normalize the ledger, replace
each converted cycle edge by its labelled sphere-product vertex and a
connected-sum tree edge, transporting all unused ports. The old cycle
edge is removed. Thus a subsequent use of the formula counts only
remaining cycles, not a sphere product already recorded as a vertex.
\end{proof}
Hatcher's oriented connected-sum discussion, printed 4 and 8/PDF 5 and 9,
and Matveev Corollary 2.4.4, printed 86--87/PDF 98--99, are the source
basis. Different origins of sphere products remain labelled: a cycle in
the late sphere system, an event tube, and a discarded sphere-product
component are distinct contributions. Counting one does not erase another,
and converting a cycle does not create a second copy of that same factor.

\section{Event adapters and finite history substitution}
\label{sec:transport-event-history}

\begin{foundation}[Oriented parametrized event topology; AT21]
\label{found:transport-event-adapter}
For each event in the supplied finite prefix, Chapter 11's audited surgery
interface must give a common retained compact region, its actual sphere
ports/collars, the post-event cap balls, and the pre-event tubes and cap
pieces. Include all discarded components with their actual reconstruction
witnesses. From these data produce an AT19 description of every pre-event
component in terms of post-event components and certified exceptions.
The maps are oriented, label-preserving in the stated relative sense and
compatible with the adjacent smooth-stage identifications.
\end{foundation}
KL Lemmas 67.5, 67.13 and 73.4, printed 2740--2741, 2745, 2761--2762/PDF
154--155, 159, 175--176, supply the source classification and reconstruction
route. The exact adapter has the following proof tasks:
\begin{enumerate}[leftmargin=*]
\item restrict the original neck parametrizations to disjoint sphere collars
and identify both retained copies through the actual common region;
\item identify each removed tube with its two recorded sphere ports, each
ordinary cap with a ball, and each projective cap with a punctured
$\mathbb{RP}^3$, with orientation and attaching maps;
\item add the discarded connected components to the correct pre-event
component list; $\mathbb{RP}^3\#\mathbb{RP}^3$ has two labelled factors;
\item use AT07 and AT20 to reconstruct, keeping the ball-extension
comparisons and all smooth-stage interfaces.
\end{enumerate}
A diffeomorphism-type statement does not already supply these data. This
is a source-qualified adapter against inherited surgery/topology, not a
claim that PC lacks the source theorem. Analytic surgery existence and
nonaccumulation remain Chapter 11 inputs.

\begin{foundation}[Discard-policy comparison; AT22]
\label{found:transport-discard-policy}
Fix the actual flow convention. Every additional removal relative to the
selected KL convention needs a certified topological reconstruction of
that removed component and, before final assembly, the appropriate
prime/model witnesses. Store its time, component provenance and disjoint
union or connected-sum role. A comparison of conventions must neither
lose that component nor count it twice.
\end{foundation}
KL Remark 73.3 omits Perelman's extra deletions of almost-round and
nonnegative-scalar-curvature components. The latter property alone is
not a spherical-factor certificate. For the default KL-type route use
exactly the exceptions of Lemma 73.4. If the inherited flow follows a
stronger deletion convention, its extra components stay a separate
adapter gate until their certified reconstruction is supplied. Do not
silently invoke finite-time extinction to validate them. The primary
source is Perelman's surgery paper, Section 4.4 and the opening of
Section 7; its Section 7.4 summarizes the later topological application.

\begin{lemma}[Substituting one marked reconstruction; AT23]
\label{lem:transport-ledger-substitution}
Suppose a vertex carrier in an AT19 ledger has another finite oriented
sphere reconstruction, with a supplied diffeomorphism to that vertex.
The inner decorations are supplied jointly disjoint from its sphere
system, and any retained outer decorations lie in unreplaced vertices.
One can substitute the second reconstruction, retaining the outer port
labels, this whole decoration, and the exact ambient reconstruction.
All cap-to-piece assignments and new sphere pairings must be stored.
No claim is made for an independently prescribed intersecting decoration.
\end{lemma}
\begin{proof}
Represent the inner reconstruction by its actual disjoint interior sphere
system in the outer vertex, using its supplied diffeomorphism. Combine
these spheres with the inner torus decoration. Apply AT14 to move this
whole system off the outer attachment balls, fixed on any protected
external collars. Transport all inner cut maps and pairings by AT08.
Each outer ball now lies in exactly one inner punctured component; record
that assignment. Its removal commutes with the inner cuts by AT15.
The old outer pairing joins the corresponding assigned new ports, while
inner pairings remain unchanged. AT07 gives the quotient equivalence.
Compose with the original reconstruction. If the inner sphere system is
empty this is just carrier transport. Every label is a tagged outer or
inner label, so no unrelated ports are identified.
\end{proof}
This proof may move the interior decorations. Their final images and
conjugated pairings, rather than an unsupported identity on those
surfaces, are the output. It fixes only the explicitly protected external
collars. Parametrized incoming cap balls are kept fixed in the outer
carrier while the inverse ball-motion moves the decorations.

\begin{proposition}[Flattening a finite surgery history; AT24]
\label{prop:transport-finite-history}
Suppose a finite ordered list covers every surgery event from the initial
slice to one selected slice, with all AT10 smooth-stage identifications,
AT21 event reconstructions and AT22 convention comparisons. A supplied
finite sphere reconstruction of the selected slice can be flattened to
an AT19 reconstruction of the initial carrier whose leaves are the
selected-slice factors and the recorded event exceptions.
\end{proposition}
\begin{proof}
Induct backwards over the finite list. Compose with the last smooth-stage
identification and substitute each known post-event reconstruction into
the corresponding vertex of the event reconstruction by AT23. Add the
certified discarded components in their recorded positions. Repeat for
the preceding event. At each step the vertex, edge and decoration labels
are finite dependent disjoint unions, tagged by event and component;
there is no identification merely because two factors have the same type.
Composition of AT07 equivalences proves the reconstruction identity.
For a zero-event list use the smooth-stage map. Disconnected slices are
handled with their component assignments, even if several post-event
components join one pre-event component.
\end{proof}
\paragraph{The resulting initial-carrier torus maps.}
If a retained torus $t:T\to P$ in a leaf is injective on fundamental
groups, move its entire collar off the final attachment caps by AT14.
Let $c:R\to P$ be that leaf's cap inclusion, $j:R\to M_0$ its actual
sphere reconstruction map, $h$ the final avoidance diffeomorphism and
$c i=h t$. The sphere-cap argument gives $(c)_*$ an isomorphism and the
finite sphere-gluing lemma gives $(j)_*$ injective. Therefore the actual
embedded collar map $u=j i$ satisfies
\[
                 u_* = j_* c_*^{-1} h_* t_*
\]
with the recorded basepoint paths, and is injective by AT09. Apply this
simultaneously to the finite leaf decoration. Interior collars are not
identified at sphere ports, so their images remain disjoint. This is a
forward map of the final transported tori into the initial carrier,
not an embedding of $P$, a map on the whole late slice, or a canonical
backward trajectory for the original torus. Endpoint cut pieces remain
on the capped prime representatives.

In particular, this constructs finite topological history data without an
infinite-time limit or a global map $M_t\to M_0$. An infinite event list
with no finite-prefix certificate does not satisfy the premise. Event
spheres and endpoint tori are made jointly disjoint at each substitution;
independent torus isotopies would not justify this step.

\section[Exceptions and metric transport]{Exceptional factors, representative changes and metric transport}
\label{sec:transport-models-and-representatives}

\begin{foundation}[Actual exceptional factors; AT25]
\label{found:transport-exceptional-factors}
Each exceptional entry used by AT21/AT22 must name a connected oriented
carrier and an actual diffeomorphism to the classified model, with its
prime status or a finite prime reconstruction. For the KL convention,
the list consists of sphere products and spherical space forms, with
$\mathbb{RP}^3\#\mathbb{RP}^3$ split into two entries. A spherical entry
includes a finite free isometric action on the round $S^3$, not merely
a finite abstract fundamental group. Chapter 7 and the inherited PC
adapter supply the actual complete model metrics and prime witnesses.
\end{foundation}
The oriented model can be given the orientation transported by the supplied
diffeomorphism; reversing a convention is not an assumed orientation-reversing
self-diffeomorphism of a chiral spherical space form. The $S^1\times S^2$
model uses the $S^2\times\R$ geometry. A whole $S^3$ identity component,
a removable identity summand and the empty carrier remain distinct.
No fresh PC probe or reproof of spherical recognition is introduced.

\begin{lemma}[Transporting a complete interior metric; AT26]
\label{lem:transport-interior-metric}
Let $f:Q\to Q'$ be a smooth diffeomorphism of the actual compact cut
carriers, and let $g$ be a complete Riemannian metric on all of
$\operatorname{int}(Q)$ with a specified Thurston model. Then
\[
           g'=(f^{-1})^*g\quad\hbox{on }\operatorname{int}(Q')
\]
has the same model, is complete, and has exactly the same total volume
(including the possibility of infinite volume). Thus finite volume is
preserved when the chosen tag is hyperbolic.
\end{lemma}
\begin{proof}
The defining pullback equation makes the interior restriction of $f$ an
isometry. Curve lengths agree in both directions, hence intrinsic distances
agree. A Cauchy sequence pulls back to a Cauchy sequence, converges by
completeness, and pushes forward to a limit. Local model charts compose
with $f^{-1}$. The inherited Riemannian change-of-variables theorem equates
the volume measures, and hence their values on the whole interiors.
\end{proof}
The metric is not claimed to be the restricted late Ricci-flow metric.
That restricted metric on an open compact core can be incomplete. This
lemma transports a supplied full-interior metric; it neither constructs
one from a Seifert tag nor closes Chapter 7's model realization. E1 still
requires finite volume for hyperbolic tags only. Smoothness is essential;
a bare topological equivalence is not a tensor pullback interface.

\begin{foundation}[Optional prescribed-prime comparison; AT27]
\label{found:transport-prescribed-primes}
If a consumer prescribes a different prime decomposition, supply a
permutation of nonidentity factors, actual oriented factor diffeomorphisms,
an $S^3$ identity comparison and, wherever required, the commuting
external-boundary/collar and reconstruction squares. Use AT08 and AT26
to transport cuts, pairings, injections and metrics to those factors.
\end{foundation}
Hatcher's oriented uniqueness theorem supplies the unmarked source route;
it is not already this boundary-marked comparison. An arbitrary boundary
diffeomorphism need not extend: on a solid torus, a map sending a meridian
to a nonmeridional longitude cannot extend since it would not preserve
the kernel of the boundary inclusion on fundamental groups. Thus labels
and arbitrary preassigned parametrizations are different requirements.
The public certificate in Chapter 1 is existential. It may choose the
prime decomposition constructed here, so AT27 is \emph{optional} on that
route; AT13/AT16 and actual prime witnesses are still mandatory. There
is no need to demand canonical JSJ cuts or uniqueness before proving the
existential theorem.

\section{Terminal data and the exact Chapter 16 handoff}
\label{sec:transport-terminal-handoff}

\begin{proposition}[Evaluation of explicit terminal history; AT28]
\label{prop:transport-terminal-evaluation}
Given an explicitly empty selected terminal slice, a covering finite
history with AT21/AT22 data, and AT25 certificates for every exception,
AT24 reconstructs the initial manifold entirely from those entries and
the sphere-product cycle factors. If the initial carrier is nonempty,
a purported record with no remaining components or factors cannot be
accepted as its reconstruction.
\end{proposition}
\begin{proof}
Use the empty disjoint-union ledger for the selected slice in AT24.
Backward event substitution adds exactly the certified entries. Apply
AT20 componentwise. The reconstruction diffeomorphism identifies components,
so an empty result cannot reconstruct a nonempty initial carrier. Retained
$S^3$ identity components are actual vertices, even when a normalized
connected-sum notation later suppresses their summand names.
\end{proof}
This implements the use of explicit AT05 data. It does not derive the
terminal branch from the failure of late-time availability and does not
prove finite-time extinction. Disconnected initial carriers can be treated
componentwise, although the public theorem assumes connectedness.

\begin{definition}[Topological reconstruction export; AT30]
\label{def:transport-reconstruction-export}
The Chapter 15 export consists of:
\begin{enumerate}[leftmargin=*]
\item a finite chosen prime family, any retained $S^3$ identities and the
actual oriented connected-sum reconstruction of the initial carrier;
\item on every prime, the actual finite collared incompressible torus
system, its cut carriers and all pairing/reconstruction maps;
\item exact identifications of these cut carriers with supplied hyperbolic,
Seifert or exceptional model carriers, including orientation, external
boundary and cap provenance;
\item the finite history, sphere-cycle and discarded-factor ledger, and the
actual basepoint comparisons underlying any stated group injection.
\end{enumerate}
Model metrics are separate inputs from Chapters 7--8, transported by AT26
once supplied. The export contains neither a proof of long-time existence
nor a hidden completed {\upshape\code{Geometrizes}} certificate as an input.
\end{definition}
For a late branch, AT12/AT29 consume the same selected slice and AT18 its
chosen prime data with a normalized sphere ledger; for a terminal branch, AT28 supplies the reconstruction.
Chapter 16 attaches actual full-interior model witnesses and packages the
public endpoint. A local static graph certificate alone supplies none of
this analytic availability, history, prime refinement or model geometry.

\begin{theorem}[Conditional Chapter 15 assembly; AT31]
\label{thm:transport-conditional-assembly}
Supply the inherited smooth/topological interfaces, the covering finite
history and all its AT21/AT22/AT25 data. In the late branch also supply
AT12's exact decomposition, AT01, AT13/AT16, and one admissible selected
time as in AT29. In the terminal branch supply AT28's explicit empty
slice. Then the constructions above give AT30 on the initial carrier,
with its model-identification fields conditional on the corresponding
actual model carriers. If a prescribed prime representative is requested,
include AT27; it is unnecessary for the chosen existential output.
\end{theorem}
\begin{proof}
In the late branch use AT15--AT18 to obtain chosen prime factors and their
actual cuts. In the late branch use the normalized AT18 ledger, whose converted
cycle factors are already explicit vertices. In the terminal branch
start with the empty selected-slice ledger as in AT28. Flatten the
finite history by AT24; the
sphere graph formula AT20 records only the remaining additional cycles. Resolve each exceptional entry using AT25.
The finite substitution and cap-avoidance constructions transport the
joint decoration and pairings by AT07--AT08. Prime factors themselves are
kept as capped carriers with their original essential cuts: puncturing
them for reconstruction does not redefine the endpoint pieces. The
finite tagged unions enumerate the resulting primes, torus labels and
cut pieces. Apply AT09 for map comparisons and, if requested, AT27 for a
new prime representative. AT26 handles any supplied full-interior metrics.
These constructions supply exactly AT30's fields without invoking the
Chapter 16 conclusion.
\end{proof}

\section{Complete Chapter 15 source and dependency audit}
\label{sec:transport-complete-audit}

All 31 AT targets, including the earlier AT01--AT05 and the full gluing,
cap-avoidance and late-injection arguments, have been reviewed. There are
no remaining unexpanded Chapter 15 topic placeholders. The following
source-qualified producers or comparisons remain: AT01, AT13, AT16,
AT21, AT22, AT25 and optional AT27. The other entries are definitions or
written conditional arguments, not checked Lean declarations. Inherited
smooth/group/category bindings and upstream Chapter 14 producers remain
explicit premises. Source citation is not proof acceptance.

\paragraph{Mathematical findings.}
\begin{enumerate}[leftmargin=*]
\item Smooth gluing needs compatible collar charts; an equality on the
boundary alone does not match normal derivatives. Every cut, cap and
pairing is transported by an actual map. Basepoint tracks remain explicit.
\item The former torus-only cap-avoidance argument extends to the whole
finite sphere/torus decoration. Its motion is in the capped carrier,
fixes only the recorded smaller external collar, and does not require
the surface complement to be connected.
\item Sphere cuts commute with cusp gluings only after disjointness is
established. The capped late carriers are gluings of hyperbolic cores and
capped thin factors, not alleged embedded copies of every capped factor.
\item Sphere graphs may have loops and parallel edges. The count is
$|E|-|V|+c(G)$, with actual gluing witnesses. Empty disjoint unions,
$S^3$ identities and separate sphere-product origins are not conflated.
\item A finite prefix covers every event; discreteness alone is not local
finiteness, and no eventual cessation of surgery is required. KL and
Perelman discard conventions need an explicit component ledger.
\item Model peripheral injection, global cusp injection, late carrier
injection and finite-history reconstruction are distinct statements.
Neither an abstract group isomorphism nor a backward topological
classification is a global reverse-time map.
\item The public endpoint permits a chosen prime decomposition. Arbitrary
prescribed boundary markings are not supplied by unmarked uniqueness.
Complete metrics live on the entire actual cut-piece interiors and are
transported by diffeomorphisms, not by restricting the flow metric.
\end{enumerate}

\paragraph{Stage order and chapter relations.}
The verified dependency record separates the following roles:
\begin{enumerate}[leftmargin=*]
\item inherited interfaces and generic marked transport supply relative
prime and sphere operations;
\item late analytic data and static graph output supply the late
topological certificate;
\item finite event adapters and a chosen-slice certificate supply initial
reconstruction, followed by Chapter 16 packaging.
\end{enumerate}
Chapter 2 owns the inherited interfaces. Chapters 3--4 remain upstream
of local geometry and never require this transport chapter. Chapters
5--6 provide relative prime/refinement producers; their generic use of
this chapter's finite gluing or marked transport is distinguished from
its later flow application, so chapter numbers do not create a cycle.
Chapter 7 supplies model realization; Chapter 8 supplies hyperbolic-core
facts and the separate global cusp route. Chapter 11 supplies actual
events and finite prefixes; Chapter 12 supplies a selected late slice
or explicit terminal data. Chapters 13--14 provide static geometry on
that same carrier, with LC90 downstream of static production. Generic
transport uses no LC90 or final assembly premise. Chapter 16 consumes
the export and actual model witnesses; it is not a producer for this
chapter. H1--H4, E1, the late-or-terminal alternative and the exclusion
of finite-time extinction from the default route are unchanged.

\paragraph{Verification scope.}
The source checks include Perelman's actual surgery/discard passages,
KL's event and full cusp-injection passages, the archived Morgan--Tian
2012 assembly (not a verified 2014 body), Hatcher's oriented prime/gluing
arguments, Matveev's second-edition graph comparison and Hirsch's
support-controlled isotopy/gluing proofs. Author corrections and project
scope findings remain distinct. Inventories, rejection controls and
independent dependency sorting check consistency, not mathematical truth.
No new PC inspection, release acceptance, Lean invocation or environment
pinning accompanies this chapter. The next blueprint topic is Chapter 16;
the listed qualified producers remain implementation obligations.


% BEGIN REVISION171 ACTUAL GRAPH TRANSPORT BINDING
\section{The completed graph-refinement input on actual late carriers}
\label{sec:graph-audited-transport-binding}
\markright{\thesection. Actual graph transport and metrics}

\begin{corollary}[Actual AT13/AT16 and graph-region GA06 binding; GAU03]
\label{cor:graph-audited-late-binding}
On the common late-slice branch supplied by TCF07 and IAU03, the actual
thin graph carriers admit AT13's chosen relative prime data and AT16's
good-block/essential-cut data, with complete E1 model witnesses on their
resulting cut interiors. With the already supplied good hyperbolic cores,
AT17--AT18 give the corresponding essential tori and chosen prime
carriers of the actual late slice. GA05 and GM06 give all thick and
refined-graph metrics on these cut interiors, with their marked
reconstruction. This supplies the graph-region application of GA06.
\end{corollary}
\begin{proof}
TCF07 retains the actual smooth graph presentation, compact domains,
external labels and maps; BCF04's internal cuts, smooth bundle charts
and collar roundings specify the category used by RG03. Their retained
FC41 foundations remain part of the earlier static theorem. IAU03 gives
injection of each actual hyperbolic--thin interface into the late slice
and injection of each thin carrier. In particular each interface
injects into its thin carrier by factoring the actual map to that slice.
Thus RG05 and GAU01 apply to each component with the required external
injections. They supply the same capped factors and the exact AT13/AT16
maps, and exclude a terminal solid torus carrying an external port.
For a closed thin component there is no such port condition; the closed
factor and sphere-product branches are already included.

AT15 commutes the sphere operations with the actual cusp gluings.
The hyperbolic-core peripheral maps are injective in their cores;
RG05 supplies all the other block-port injections. AT17's finite-gluing
argument therefore proves essentiality in the actual capped late
carriers, and AT18 proves their irreducibility outside the explicitly
recorded prime exceptions. Normalize the sphere ledger as in GAU01,
retaining sphere products and $S^3$ identities. GM06 and GAU02 give
complete metrics on every refined graph piece; GA05 transports the
complete finite-volume hyperbolic metric to every actual thick piece.
Their boundary markings and reconstruction maps are those already
stored. These metrics need not glue isometrically across a torus.
The claim concerns the selected late slice; no step substitutes an
unproved event-history adapter to reach the initial manifold.
\end{proof}
% END REVISION171 ACTUAL GRAPH TRANSPORT BINDING

\chapter{Global assembly and Geometrization}
\label{ch:global-assembly-topic}

\paragraph{Scope.} The structural certificate assembly and final analytic integration.

\paragraph{Planned sections.}
\begin{enumerate}[leftmargin=*]
\item Package fields.
\item Normalization.
\item Global torus system.
\item Cut reconstruction.
\item Model witnesses.
\item Certificate packaging.
\item Public theorem.
\end{enumerate}

\paragraph{Existing material and present task.} M1 field/assembly/module contracts and contract-phase vertical slice. Refine exact handoffs from topic chapters while preserving non-circular structural/analytic separation.

\paragraph{Current status.} Sections16.8 onward complete the assembly
blueprint and qualify all upstream producers. Earlier implementation
orders are historical schedules, not mathematical prerequisites for
independent chapter writing or permission to pin the environment.

% BEGIN migrated B012
\section{M1 field-level contract}
\label{sec:m1-field-contract}

The immediate deliverable is a field table precise enough to guide a Lean
skeleton. The names below remain provisional, but their mathematical content
is fixed for the first implementation pass.

\setlength{\tablebodywidth}{\dimexpr\textwidth-4\tabcolsep\relax}
\begin{longtable}{@{}P{0.300000\tablebodywidth}P{0.433333\tablebodywidth}P{0.266667\tablebodywidth}@{}}
\textbf{Field} & \textbf{Required content} & \textbf{Main obligation}\\\hline
\endhead
Prime index and factors & A finite index type and oriented connected prime
carriers, together with a connected-sum reconstruction equivalence to $M$.
& Prime decomposition and orientation bookkeeping.\\
Torus system & For each prime factor, a finite family of pairwise disjoint,
two-sided embedded tori with explicit incompressibility proofs.
& Embedding, disjointness, two-sidedness, and $\pi_1$-injectivity.\\
Cut pieces & A finite family of compact carriers obtained by cutting along the
tori, with each boundary component labelled by a torus side.
& Cut-space construction and boundary pairing.\\
Geometric witnesses & For each cut piece, a model-geometry tag, an interior
complete locally homogeneous metric, a finite-volume proof when the tag is
hyperbolic, and cusp or boundary collar data.
& Hyperbolic/locally homogeneous geometry API.\\
Global reconstruction & Explicit equivalences or diffeomorphisms identifying
the assembled prime-and-piece object with the original $M$, preserving
orientation and the boundary gluings.
& Topological assembly theorem.\\
\end{longtable}

The \code{ThickThinCollapsePackage} has a deliberately different field profile:
it records finite hyperbolic thick-region certificates, finite graph-manifold
thin-region certificates, common boundary tori, embeddings or correspondence
maps into a sufficiently late flow slice, and the isotopy/topology data that
transports those regions back through the surgery history. Its fields are
analytic output plus topology transport. The package supplies the reference
prime decomposition; \code{assemble_geometrization} retains it and constructs
the final torus, piece and geometric certificate data.

The representation choices for M1 are also fixed provisionally:
finite index types rather than unstructured lists, explicit embeddings rather
than informal subsets, a separate compact cut piece and open interior, and
explicit homeomorphisms or diffeomorphisms for reconstruction. Quotient
implementations may be introduced behind these interfaces only after the
existing manifold APIs have been audited. The public theorem remains an
existence statement and does not require canonical choices.

M1 is accepted with four small endpoint examples: the standard 3-sphere, a
flat 3-torus, a closed hyperbolic 3-manifold, and $\Sph^2\times\Sph^1$.
These examples exercise the spherical, Euclidean, hyperbolic, and product
geometries without invoking finite-time extinction. The examples are API
tests for the certificate and assembly layers, not substitutes for the GC
proof.

The first integration theorem is conditional:
\begin{equation}
\begin{aligned}
 &\texttt{geometrization\_of\_ricci\_flow\_package} : \\
 &\quad \texttt{PoincareAdapter} \\
 &\quad\longrightarrow \texttt{RicciFlowSurgeryPackage M} \\
 &\quad\longrightarrow \texttt{Geometrizes M}.
\end{aligned}
\end{equation}
\code{RicciFlowSurgeryPackage} is an interface whose fields are eventually proved;
it is not a permanent axiom.
% END migrated B012

% BEGIN migrated B013
\section{M1 assembly specification}
\label{sec:m1-assembly-specification}

Revision 14 named the fields. Revision 15 specifies the first structural
proof that consumes them. The assembly theorem is intentionally a theorem
about finite certified regions and their topology; it does not inspect the
Ricci-flow equation, choose surgery scales, or prove a compactness result.
Those analytic tasks belong to the theorem that constructs the
\code{ThickThinCollapsePackage}.

\subsection{The assembly input}

The package must contain enough information to prevent the assembly proof
from silently assuming the conclusion. Its fields are divided into five
groups:

\setlength{\tablebodywidth}{\dimexpr\textwidth-4\tabcolsep\relax}
\begin{longtable}{@{}P{0.263736\tablebodywidth}P{0.461538\tablebodywidth}P{0.274725\tablebodywidth}@{}}
\textbf{Group} & \textbf{Certified data} & \textbf{Consumed by}\\\hline
\endhead
Underlying topology & A prime decomposition of $M$, a finite set of region
indices, and orientation-preserving embeddings of the certified regions into
the corresponding prime factors. & Prime and connected-sum reconstruction.\\
Thick regions & For each thick index, a compact truncated core, its boundary
torus labels, a complete finite-volume hyperbolic witness on its interior,
and the inclusion/collar data needed to compare it with the late flow slice.
& Hyperbolic-piece constructor.\\
Thin regions & For each thin index, a compact graph-manifold certificate,
including a finite Seifert decomposition or a Sol/product model and the
boundary tori exposed to the thick cores. & Graph-manifold-to-geometry
conversion.\\
Compatibility & Pairwise disjointness, common boundary parametrizations,
orientation signs, and a matching relation saying which thick and thin
boundary sides are glued. & Torus system and gluing reconstruction.\\
Transport & Isotopy or homeomorphism data transporting the regions through
the surgery history to the original manifold, together with coverage and
complement statements. & Global reconstruction and incompressibility.\\
\end{longtable}

The coverage field is essential. It states that, after deleting the certified
boundary tori and the prescribed collars, the listed thick and thin regions
and the exceptional prime pieces cover each prime factor. A list of good
regions without this field could prove only that some geometric pieces occur;
it would not prove geometrization of $M$.

The boundary field is also explicit about sides. Every boundary torus has a
region index, a torus index, and a side label. The matching relation is an
involution on boundary sides with no fixed points, and its gluing map is
orientation reversing as a map between the induced boundary orientations.
This is the data used by the reconstruction theorem; it is not inferred from
a picture of the decomposition.

\subsection{Assembly lemma stack}

The proof of \code{assemble_geometrization} is split into small lemmas so that
topology and geometry can be developed independently:

\begin{enumerate}[label=\textbf{A\arabic*:},leftmargin=*]
\item \textbf{Region normalization.} Truncate each finite-volume hyperbolic
  core along its certified cusp collars and turn its boundary labels into
  compact cut-piece data. Pull back the complete metric from the entire
  hyperbolic manifold along a specified diffeomorphism from the cut-piece
  interior; restriction to a finite truncated core is not complete.
  The core embedding and the complete-interior identification are separate maps.
\item \textbf{Thin conversion.} Refine each graph-manifold certificate into
  finitely many Seifert, product, or Sol pieces and attach the corresponding
  model-geometry witnesses. This lemma includes the convention for twisted
  $I$-bundles and torus bundles, rather than treating them as informal
  exceptions.
\item \textbf{Exceptional prime conversion.} Use the standard spherical,
  Euclidean, or $Sph^2\times\R$ models for the prime factors not represented
  by a thick or thin region. The simply connected spherical input is consumed
  through \code{PoincareAdapter}; no finite-extinction theorem is used.
\item \textbf{Torus-system formation.} Collect the distinct boundary tori,
  prove pairwise disjointness and two-sidedness, and prove the required
  $\pi_1$-injectivity. The last obligation is recorded as a theorem field or
  an explicit lemma; it may not be hidden in the definition of a region.
\item \textbf{Gluing reconstruction.} Glue the compact carriers along the
  matching boundary maps and construct an orientation-preserving
  homeomorphism (or diffeomorphism, after the smooth API is fixed) from the
  glued object to $M$. Coverage and transport are used here.
\item \textbf{Packaging.} Assemble the prime decomposition, torus system,
  cut pieces, boundary maps, and model metrics into
  \code{GeometrizationCertificate M}.
\end{enumerate}

The intended dependency chain is therefore
\[
\begin{gathered}
\text{thick hyperbolic conversion},\quad\text{thin graph conversion},\\
\text{exceptional conversion}\\
\Longrightarrow\;\text{torus system and reconstruction}\\
\Longrightarrow\;\texttt{GeometrizationCertificate}.
\end{gathered}
\]
No lemma in this chain has a Ricci-flow hypothesis. The analytic producer
must prove the hypotheses of the package before this chain is invoked.

\subsection{Schematic Lean boundary}

The following is a non-executable interface sketch. It records the shape of
the handoff while leaving concrete manifold and finite-type choices to the
Lean API audit:

\begin{quote}
\code{structure ThickThinCollapsePackage (M) where}\\
\quad\code{primeData : PrimeDecomposition M}\\
\quad\code{thickData : FiniteThickHyperbolicRegions M}\\
\quad\code{thinData : FiniteGraphManifoldRegions M}\\
\quad\code{boundaryData : CompatibleBoundaryGluing thickData thinData}\\
\quad\code{transportData : RegionsTransportToM ...}\\
\quad\code{coverage : CertifiedRegionsCoverM ...}\\[2pt]
\code{def assemble_geometrization}\ :\\
\quad\code{PoincareAdapter -> ThickThinCollapsePackage M ->}\\
\quad\code{GeometrizationCertificate M}
\end{quote}

The names in this display are contracts, not declarations to paste into the
blueprint. The authoritative types will live in Lean modules and will be
changed only through a reviewed API update. In particular, \code{coverage} and
\code{transportData} must not be replaced by an unrestricted existential over a
new certificate, since that would make the assembly theorem circular.

\subsection{M1 acceptance tests and work split}

The assembly layer is accepted when the following tests are proved against
the same definitions used by the theorem:

\begin{itemize}[leftmargin=*]
\item a single spherical prime factor produces a spherical certificate;
\item a flat 3-torus produces a Euclidean certificate with an empty torus
  system;
\item a closed hyperbolic manifold produces one hyperbolic piece with empty
  boundary;
\item a graph-manifold example, including a torus bundle, produces the
  selected Seifert/Sol pieces and verifies the boundary matching;
\item the structural theorem has no dependency on \code{RicciFlowSurgeryPackage}
  or finite-time extinction.
\end{itemize}

This yields four concrete handoffs: (i) topology defines the cut and gluing
objects, (ii) the hyperbolic workstream defines truncated-core witnesses,
(iii) the graph-manifold workstream defines the thin conversion, and (iv) an
assembly workstream proves the six lemmas above. The flow and surgery teams
can implement their producer against these handoffs without deciding the
endpoint representation again.
% END migrated B013

% BEGIN migrated B014
\section{M1 Lean module map and acceptance matrix}
\label{sec:m1-module-map}

The next implementation decision is module ownership. The endpoint should be
assembled from small modules whose imports follow the mathematical dependency
graph. This prevents a convenient analytic definition from leaking into the
topological certificate and makes failed interfaces visible early.

\subsection{Provisional module tree}

The following paths are names for the first GC project layout, not a claim
about the final repository organization:

\setlength{\tablebodywidth}{\dimexpr\textwidth-4\tabcolsep\relax}
\begin{longtable}{@{}P{0.329670\tablebodywidth}P{0.428571\tablebodywidth}P{0.241758\tablebodywidth}@{}}
\textbf{Module} & \textbf{Contents} & \textbf{May import}\\\hline
\endhead
\code{GC/Topology/Basic} & The chosen category of compact oriented
3-manifolds, boundary components, embeddings, homeomorphisms, and orientation
conventions. & Existing manifold API only.\\
\code{GC/Topology/Prime} & Prime factors, connected-sum reconstruction, and
the finite prime-index contract. & Basic.\\
\code{GC/Topology/Torus} & Embedded two-sided tori, disjointness,
incompressibility, boundary sides, and matching involutions. & Basic.\\
\code{GC/Topology/Cut} & Compact cut carriers, interiors, boundary labels,
collars, and reconstruction from matched sides. & Prime; Torus.\\
\code{GC/Geometry/Models} & The eight model tags and the common locally
homogeneous metric-witness interface. & Existing Riemannian API.\\
\code{GC/Geometry/Hyperbolic} & Hyperbolic finite-volume cores, cusp
truncation, and the conversion to a geometric-piece witness. & Models; existing
Riemannian API.\\
\code{GC/Geometry/GraphManifold} & Seifert, product, and Sol certificates
and their conversion to model-geometry witnesses. & Models; Topology/Cut.\\
\code{GC/Assembly/Package} & \code{ThickThinCollapsePackage}, coverage,
transport, boundary compatibility, and the package invariants. & Cut;
Hyperbolic; GraphManifold.\\
\code{GC/Assembly/Normalization} & Region normalization and thin/exceptional
conversion lemmas. & \code{Assembly/Package}; Geometry modules.\\
\code{GC/Assembly/Reconstruction} & Torus-system formation, global gluing,
and certificate packaging. & Normalization; Prime; Cut.\\
\code{GC/Assembly/Examples} & The four endpoint examples and deliberately
invalid package tests. & Reconstruction.\\
\end{longtable}

No foundational topology or geometry module in this tree imports Ricci flow,
surgery, or Perelman estimates. The exceptional-prime lemma in
\code{GC/Assembly/Normalization} may import the narrow \code{PoincareAdapter} interface
needed for the spherical branch, but it may not import PC flow or surgery
implementation modules. The sole purpose of this restriction is architectural:
the structural theorem should remain usable if the analytic implementation or
its time-slice representation is changed.

\subsection{Representation decisions that must be explicit}

The first Lean skeleton should settle the following choices in one place:

\begin{enumerate}[label=\textbf{D\arabic*:},leftmargin=*]
\item \textbf{Finite indexing.} Use finite types with decidable equality for
  prime factors, regions, boundary components, and model tags. Lists may be
  used for computation, but the public certificate must expose the finiteness
  proof independently of list ordering.
\item \textbf{Boundary sides.} Represent a side as a region/boundary index
  together with its induced orientation. The matching relation must carry its
  involution and no-fixed-point proofs, rather than being an untyped pairing
  of names.
\item \textbf{Cut versus interior.} Keep the compact cut carrier and its
  interior as separate objects linked by an inclusion. This lets a complete
  metric live on the open piece while the gluing theorem works with compact
  boundary data.
\item \textbf{Incompressibility.} Choose one definition of the torus
  inclusion (the intended $\pi_1$-injectivity formulation) and prove
  equivalence to any alternate definition used by a source. The certificate
  must store the selected proof, not just a proposition name whose meaning
  varies by module.
\item \textbf{Geometry witness.} Separate the model tag from the metric and
  local-homogeneity proof. This allows a single assembly theorem to treat all
  eight geometries while model-specific analysis remains modular.
\item \textbf{Reconstruction map.} Store an explicit orientation-preserving
  equivalence from the glued object to $M$. Do not encode reconstruction only
  as equality of underlying sets or as an existence proposition hidden inside
  \code{GeometrizationCertificate}.
\end{enumerate}

These decisions are the last representation choices before writing theorem
stubs. Changing one later is possible, but it should be treated as an API
change affecting every assembly consumer.

\subsection{M1 milestone order}

The first implementation pass is divided into milestones with no analytic
placeholders:

\setlength{\tablebodywidth}{\dimexpr\textwidth-4\tabcolsep\relax}
\begin{longtable}{@{}P{0.168539\tablebodywidth}P{0.561798\tablebodywidth}P{0.269663\tablebodywidth}@{}}
\textbf{Gate} & \textbf{Deliverable} & \textbf{Exit evidence}\\\hline
\endhead
M1a & Basic manifold, finite-index, orientation, and boundary-side types. &
Four small type-checking examples.\\
M1b & Prime and torus systems, including disjointness, two-sidedness, and
incompressibility fields. & Construction and projection lemmas; no quotient
implementation required.\\
M1c & Cut carriers and compact/interior relationship. & Boundary-label and
cut/reconstruction lemmas.\\
M1d & Model tags, hyperbolic-core witness, and graph-manifold witness. &
Spherical, Euclidean, hyperbolic, product, and Sol witness examples.\\
M1e & Package invariants: coverage, matching, transport, and compatibility. &
An invalid package cannot be constructed without supplying the missing proof.\\
M1f & Six assembly lemmas and \code{assemble_geometrization}. & Kernel-checked
conditional theorem and four endpoint examples.\\
\end{longtable}

M1e is an important negative test. It should be impossible to obtain the
endpoint certificate from two unrelated lists of regions: the package must
force a proof that the regions cover the relevant prime factors and that all
boundary sides are accounted for. This is the point at which the blueprint
guards against a formally vacuous geometrization theorem.

\subsection{Acceptance matrix and independent ownership}

Each M1 row has one owner and one downstream consumer. The proposed split is:

\begin{itemize}[leftmargin=*]
\item \textbf{Topology owner:} Basic, Prime, Torus, and Cut; handoff is the
  fully typed cut-and-gluing API.
\item \textbf{Geometry owner:} Models, Hyperbolic, and GraphManifold;
  handoff is one certified geometric witness per compact piece.
\item \textbf{Package owner:} Assembly/Package; handoff is a non-circular
  \code{ThickThinCollapsePackage} with coverage and transport proofs.
\item \textbf{Assembly owner:} Normalization, Reconstruction, and Examples;
  handoff is \code{assemble_geometrization} plus the negative package tests.
\end{itemize}

The owners may work independently after the shared choices D1--D6 are
approved. The K0 finite package needs no PC release; concrete geometric packages
consume the audited inherited types after K1. The
analytic producer team begins only after M1f, because its output type is the
package rather than an informal description of a flow slice.

The machine-readable handoff for this map is
\code{m1_module_map.csv}, maintained beside the theorem and reference inventories.
It records the owner, imports, primary outputs, and acceptance evidence for
each proposed module. A module is not considered started merely because its
row exists; its status changes only after the corresponding Lean target
builds.

The exact M1 input/output obligations are also mirrored in
\code{m1_contracts.csv}. This file is deliberately separate from the theorem
inventory: the inventory tracks the project-wide theorem status, while the
contract file records the reviewable boundary of the first vertical slice.

The M1a wrapper audit is recorded separately in \code{m1_boundary_api.csv}, with
one row per field or B-lemma. Its status changes only after an exact Lean
declaration and axiom audit have been recorded.
% END migrated B014

% BEGIN migrated B015
\section{Explicit M1 theorem contracts}
\label{sec:m1-explicit-contracts}

Revision 16 exposed one dependency that must not remain implicit. The
exceptional-prime branch uses the completed Poincar\'e result to recognize the
spherical cases. Therefore the default assembly theorem takes the narrow
\code{PoincareAdapter} as an explicit argument. This is a contract correction, not
a change to the mathematical route.

\subsection{The adapter boundary}

\code{PoincareAdapter} is an import interface to the pinned completed PC release.
It exposes only the consequences used by the endpoint layer:

\begin{itemize}[leftmargin=*]
\item the simply connected closed 3-manifold conclusion in the chosen
  category;
\item recognition of the spherical-space-form or spherical-prime geometric
  witness required by the certificate;
\item the category, orientation, and smooth/topological conversion lemmas
  needed to transport that witness into the GC endpoint types.
\end{itemize}

It does not expose the PC proof's finite-extinction history, internal surgery
data, or implementation-specific namespaces. In particular, no theorem in
the assembly layer may obtain a certificate by importing a PC theorem whose
conclusion is already \code{Geometrizes M}.

The corrected signatures are:
\begin{equation}
\begin{aligned}
 &\texttt{assemble\_geometrization} : \\
 &\quad \texttt{PoincareAdapter} \\
 &\quad\longrightarrow \texttt{ThickThinCollapsePackage M} \\
 &\quad\longrightarrow \texttt{GeometrizationCertificate M},
\end{aligned}
\end{equation}
and
\begin{equation}
\begin{aligned}
 &\texttt{geometrization\_of\_ricci\_flow\_surgery} : \\
 &\quad \texttt{PoincareAdapter} \\
 &\quad\longrightarrow \texttt{RicciFlowSurgeryPackage M} \\
 &\quad\longrightarrow \texttt{Geometrizes M}.
\end{aligned}
\end{equation}

If a later package design stores independently proved geometric witnesses for
all exceptional factors, a pure theorem with no adapter argument may be added
as a derived convenience. It is not the M1 contract, because the default
route should make the PC dependency visible at the point where it is used.

\subsection{Contract catalog}

The following catalog is the precise content of the first theorem stubs. The
names remain provisional until the Lean namespace audit, but each row states
the input, output, and obligations that the eventual declaration must expose.

\setlength{\tablebodywidth}{\dimexpr\textwidth-4\tabcolsep\relax}
\begin{longtable}{@{}P{0.294118\tablebodywidth}P{0.364706\tablebodywidth}P{0.341176\tablebodywidth}@{}}
\textbf{Contract} & \textbf{Input and output} & \textbf{Required obligation}\\\hline
\endhead
\code{RegionNormalization} & A finite-volume hyperbolic core with certified cusp
collars; outputs a compact cut carrier and an interior metric witness. & The
truncation boundary has the stated torus labels; an interior diffeomorphism
to the full hyperbolic manifold pulls back its complete finite-volume metric.\\
\code{ThinGraphToGeometricPieces} & A compact graph-manifold certificate; outputs
finite Seifert, product, or Sol pieces with model tags and metrics. & Twisted
$I$-bundles, torus bundles, and boundary parametrizations are covered by the
same finite output contract.\\
\code{ExceptionalPrimeToGeometricPiece} & An exceptional prime factor and the
\code{PoincareAdapter}; outputs a spherical, Euclidean, or
$\Sph^2\times\R$ witness as appropriate. & The adapter conclusion is
transported into the selected orientation and smooth/topological category.\\
\code{TorusSystemFormation} & Normalized thick pieces, converted thin pieces,
exceptional pieces, and package compatibility data; outputs one finite torus
system with side labels. & Pairwise disjointness, two-sidedness,
$\pi_1$-injectivity, and no unaccounted boundary side are proved.\\
\code{GluingReconstruction} & The torus system, compact carriers, matching maps,
coverage, and transport data; outputs an orientation-preserving equivalence
between the glued object and $M$. & The reconstruction uses the supplied
coverage and transport proofs and does not infer them from the conclusion.\\
\code{CertificatePackaging} & Prime decomposition, torus system, cut pieces,
geometry witnesses, and reconstruction equivalence; outputs
\code{GeometrizationCertificate M}. & Every certificate field is filled by a named
previous contract; no existential certificate is introduced as a shortcut.\\
\end{longtable}

The assembly theorem is the composition of these contracts together with the
adapter call for exceptional factors. Its proof should be short enough that a
reviewer can inspect every imported theorem without opening the Ricci-flow
implementation.

\subsection{Non-circularity and dependency tests}

The M1 package and theorem stubs must satisfy these syntactic and semantic
checks:

\begin{enumerate}[label=\textbf{N\arabic*:},leftmargin=*]
\item \code{ThickThinCollapsePackage} contains no field of type
  \code{GeometrizationCertificate}, \code{Geometrizes M}, or an unrestricted existential
  over either one.
\item The coverage and transport fields refer to the listed region carriers
  and to $M$ itself; they cannot be proofs whose target is the final
  certificate.
\item \code{TorusSystemFormation} proves incompressibility from an explicit
  hypothesis or an explicit preceding lemma. It cannot obtain it by
  projecting a field named \code{incompressible} from a certificate being built.
\item The adapter is imported only through its public interface. A dependency
  on PC finite-extinction or internal surgery namespaces fails the M1 audit.
\item Removing the thick-region data from a test package must make the
  hyperbolic-piece constructor fail, and removing coverage must make the
  reconstruction theorem fail. These are negative tests of the contract.
\end{enumerate}

The checks are deliberately stronger than a successful positive example. A
formally elegant certificate theorem that can be proved from a circular
package is not progress toward GC.

\subsection{M1 implementation sequence}

The first Lean implementation should now proceed in this order:

\begin{enumerate}[label=\textbf{S\arabic*:},leftmargin=*]
\item Freeze the adapter signature and the six representation decisions in
  Section~\ref{sec:m1-module-map}.
\item Implement the structures in \code{GC/Topology} and prove their projection
  and reconstruction lemmas.
\item Implement the model and region-witness interfaces without proving the
  Ricci-flow convergence theorems.
\item Implement \code{ThickThinCollapsePackage} and its non-circularity tests.
\item Prove the six catalog contracts, then compose them into
  \code{assemble_geometrization} with the adapter argument.
\item Run the four endpoint examples, the invalid-package tests, and the
  public-theorem \code{#print axioms} audit.
\end{enumerate}

Only after S6 should the long-time analytic producer be attached. This keeps
the next analytic question precise: construct a package satisfying the fields
already consumed by the assembly theorem.
% END migrated B015

% BEGIN migrated B021
\section{M1e thick\slash thin\slash collapse package contract}
\label{sec:m1e-collapse-package}

M1e fixes the output type that the long-time Ricci-flow workstream must
produce. The package is deliberately weaker than a \code{GeometrizationCertificate}:
it contains certified thick and thin regions, their geometry/topology data,
and their transport and coverage relations, but it does not contain the final
prime/torus/piece certificate as a field. This keeps the analytic producer
from assuming the theorem it is meant to imply.

\subsection{Package field groups}

\code{ThickThinCollapsePackage M} has six field groups:

\setlength{\tablebodywidth}{\dimexpr\textwidth-4\tabcolsep\relax}
\begin{longtable}{@{}P{0.266667\tablebodywidth}P{0.477778\tablebodywidth}P{0.255556\tablebodywidth}@{}}
\textbf{Group} & \textbf{Required data} & \textbf{Why it is needed}\\\hline
\endhead
Reference decomposition & A finite prime decomposition of $M$ and the
orientation-preserving connected-sum reconstruction. & Keeps sphere and torus
cutting in separate layers.\\
Late-slice regions & A finite thick index, a finite thin index, and abstract
late-slice carriers with embeddings into the relevant surgery slice. & Gives
the analytic output a finite target and prevents informal region names.\\
Thick certificates & For each thick index, a normalized hyperbolic limit/core
certificate, cusp/end data, boundary-side labels, and a conversion map to an
M1c cut carrier. & Supplies the hyperbolic branch without yet assembling the
global certificate.\\
Thin certificates & For each thin index, a compact graph-manifold
certificate, its Seifert/product/Sol refinement data, and boundary-side
labels. & Supplies the collapsed branch and its eventual geometric pieces.\\
Compatibility and coverage & Pairwise disjointness, boundary matching,
orientation signs, and a proof that thick, thin, and exceptional regions cover
each prime factor after the selected collars are removed. & Prevents a list of
local conclusions from being mistaken for a global decomposition.\\
Transport and invariance & Topological equivalences or isotopies carrying the
late-slice regions through the finite surgery history to the original prime
carriers, together with preservation of labels and incompressibility. &
Connects the late-time analytic output to the initial manifold $M$.\\
\end{longtable}

The package may refer to a late surgery slice through an abstract
\code{LateSliceCarrier} interface. It must not expose a particular implementation
of time, surgery-event indexing, or Ricci-flow parametrization to the
topological assembly theorem.

\subsection{Thick-region certificate}

The thick branch stores a \code{HyperbolicRegionCertificate}, not merely a model
tag. It includes:

\begin{itemize}[leftmargin=*]
\item a compact core carrier and its open interior;
\item a complete finite-volume hyperbolic metric on that interior;
\item the normalized pointed-limit or convergence witness from the late
  slice, to be supplied by the analytic theorem;
\item cusp coordinates or an equivalent finite-volume end certificate;
\item boundary-side maps into the M1c cut decomposition; and
\item the incompressibility proof for every retained boundary torus.
\end{itemize}

The conversion theorem \code{ThickRegionToGeometricPiece} turns this data into a
\code{GeometricPieceCertificate}. It is a genuine theorem because it must compare
the core, the M1c carrier, and the complete interior metric; the package does
not obtain the result by storing a copy of the target certificate.

\subsection{Thin-region certificate}

The thin branch stores a \code{GraphManifoldRegionCertificate} with:

\begin{itemize}[leftmargin=*]
\item a compact graph-manifold carrier and its embedding into the late slice;
\item a finite auxiliary circle-bundle/Seifert presentation with its
  actual embedded tori, collars and pairings; twisted-bundle and
  torus-bundle cases remain refinement obligations;
\item boundary-side maps and orientation data compatible with the thick
  branch;
\item the incompressibility or boundary-control hypotheses required for the
  selected collapse theorem; and
\item relative carrier and boundary maps for later M1c conversion.
\end{itemize}

The auxiliary cuts carry no assumed incompressibility. In particular,
this raw record does not take \code{TorusCutSystem},
\code{EmbeddedTorusData} (which includes injectivity), or a final
\code{CutDecomposition} as input. KL14 Definition 1.2, printed 8/PDF 3,
only supplies embedded auxiliary cuts into circle bundles. G01--G09
must construct the retained essential system and the final M1c pieces.
The generic \code{SmoothBoundaryGluing} interface requires boundary
carriers and pairings, not a pre-existing incompressible torus system.

The conversion theorem \code{ThinRegionToGeometricPieces} produces finitely many
\code{GeometricPieceCertificate} values and the additional internal tori needed
to refine the graph-manifold region. Its conclusion must state that these
internal tori are disjoint from the thick-region boundary system and that all
new boundary gluings preserve the global orientation convention. Revision 46
decomposes this proposed conversion into G01--G10 in
Section~\ref{sec:graph-refinement-contracts}; production acceptance is
still awaiting the geometric proofs G01--G09 under the revision-47 G10
policy. Closed torus-bundle
classification must retain Nil as well as Euclidean and Sol cases.

\subsection{Coverage and transport are separate fields}

Coverage and transport must not be merged into one informal ``correspondence''
field:

\begin{enumerate}[label=\textbf{P\arabic*:},leftmargin=*]
\item \code{PackageCoverage} states that the listed late-slice regions and
  exceptional prime carriers cover each prime factor after deleting exactly
  the selected collars. It includes both a pointwise statement and the finite
  index bookkeeping needed to use it in Lean.
\item \code{PackageBoundaryCompatibility} states that every exposed boundary side
  occurs once, is paired by the global involution, and has the correct induced
  orientation.
\item \code{PackageTransport} carries each region and boundary label through the
  surgery history to the original prime carrier. It preserves disjointness,
  incompressibility, and the side matching needed by M1c.
\item \code{PackageExceptionalData} identifies the prime factors not represented by
  a thick or thin region. Their geometric witnesses are supplied later by the
  \code{PoincareAdapter} and standard-model lemmas, not hidden in coverage.
\end{enumerate}

This separation lets the analytic workstream improve its transport
representation without changing the coverage theorem and lets the topology
workstream test reconstruction using finite mock packages.

\subsection{Package theorem contracts}

The M1e theorem stubs are:

\setlength{\tablebodywidth}{\dimexpr\textwidth-4\tabcolsep\relax}
\begin{longtable}{@{}P{0.344828\tablebodywidth}P{0.425287\tablebodywidth}P{0.229885\tablebodywidth}@{}}
\textbf{Contract} & \textbf{Statement shape} & \textbf{Consumer}\\\hline
\endhead
\code{ThickRegionToGeometricPiece} & A certified thick region and its M1c transport
produce one \code{GeometricPieceCertificate}. & Assembly normalization.\\
\code{ThinRegionToGeometricPieces} & A graph-manifold region and its refinement
produce finitely many model-geometry witnesses and internal torus data. &
Assembly normalization.\\
\code{PackageBoundaryCompatibility} & All thick, thin, and exceptional boundary
sides have a global orientation-compatible matching. & Torus-system formation.\\
\code{PackageCoverage} & The finite region family covers each prime carrier after
the selected collars are removed. & Gluing reconstruction.\\
\code{PackageTransport} & Late-slice regions and labels transport to $M$ while
preserving the topological invariants required by M1c. & Global reconstruction.\\
\code{PackageNonCircular} & The package contains no \code{GeometrizationCertificate} or
unrestricted existential over one. & Trust-boundary audit.\\
\end{longtable}

The analytic producer theorem is then a separate declaration:
\[
\texttt{RicciFlowSurgeryPackage M}
\longrightarrow
\texttt{LongTimeOutcome M}.
\]
Only the late branch converts to a thick/thin package; the terminal branch
supplies explicit exceptional data and reconstruction obligations.
Its proof includes the long-time estimates, hyperbolic limits, local-collapse
theorem, graph-manifold output, and surgery-history transport. The package
contract does not decide which of Morgan--Tian or Kleiner--Lott supplies each
analytic lemma.

\subsection{M1e acceptance tests}

The package interface is accepted with finite mock packages before the actual
flow theorem exists:

\begin{itemize}[leftmargin=*]
\item a closed hyperbolic example has one thick region, no thin region, and
  empty exposed boundary;
\item a flat 3-torus has one Euclidean exceptional/model region and no
  hyperbolic thick region;
\item a graph-manifold example has a thin region whose Seifert/Sol refinement
  produces the expected finite piece list;
\item a package with an omitted region fails \code{PackageCoverage};
\item a package with one reversed boundary orientation fails
  \code{PackageBoundaryCompatibility}; and
\item a package that attempts to store a final certificate in a field fails
  the non-circularity audit.
\end{itemize}

No test invokes finite-time extinction. The package is an endpoint contract
for the long-time route and can be developed while the completed PC release is
being prepared.

\subsection{Handoff to M1f}

M1f composes the package conversions, Poincar'e adapter, prime data, and M1c
reconstruction into \code{assemble_geometrization}. Once M1f is stable, the
analytic workstream has exactly one target: construct a package satisfying the
fields already consumed by the assembly theorem.

The field and theorem audit for this layer is recorded in
\code{m1_package_api.csv}.
% END migrated B021

% BEGIN migrated B022
\section{M1f assembly-composition specification}
\label{sec:m1f-assembly-composition}

M1f is the first complete structural vertical slice. It composes the M1a
boundary wrapper, M1b torus system, M1c cut decomposition, M1d geometry
witnesses, M1e package invariants, and the explicit Poincar'e adapter. Its
proof does not import Ricci flow, surgery, Perelman estimates, or finite-time
extinction. Those results will later be used only to construct the M1e
package.

\subsection{Assembly input and output}

The public M1f theorem has the single contract
\begin{equation}
\begin{aligned}
 &\texttt{assemble\_geometrization} : \\
 &\quad \texttt{PoincareAdapter} \\
 &\quad\longrightarrow \texttt{ThickThinCollapsePackage M} \\
 &\quad\longrightarrow \texttt{GeometrizationCertificate M}.
\end{aligned}
\end{equation}

The corresponding predicate theorem is a short wrapper:
\begin{equation}
\begin{aligned}
 &\texttt{geometrizes\_of\_package} : \\
 &\quad \texttt{PoincareAdapter} \\
 &\quad\longrightarrow \texttt{ThickThinCollapsePackage M} \\
 &\quad\longrightarrow \texttt{Geometrizes M}.
\end{aligned}
\end{equation}

It applies \code{assemble_geometrization} and introduces the resulting certificate
as the witness for the existential predicate. The theorem does not quantify
over a new certificate independently of the package.

\subsection{Assembly proof stages}

The proof is organized into seven named stages. Each stage has a distinct
output and can be tested before the final structure is constructed.

\setlength{\tablebodywidth}{\dimexpr\textwidth-4\tabcolsep\relax}
\begin{longtable}{@{}P{0.190476\tablebodywidth}P{0.511905\tablebodywidth}P{0.297619\tablebodywidth}@{}}
\textbf{Stage} & \textbf{Construction} & \textbf{Required evidence}\\\hline
\endhead
F1: normalize thick & Apply \code{ThickRegionToGeometricPiece} to every thick index
and transport its core to the M1c cut carrier. & Finite indexing, boundary labels,
complete hyperbolic witness, and incompressibility.\\
F2: normalize thin & Apply \code{ThinRegionToGeometricPieces} to every thin index
and collect its internal tori and model witnesses. & Finite refinement, disjoint
internal tori, orientation-compatible boundary data.\\
F3: exceptional factors & Apply the \code{PoincareAdapter} and standard-model
lemmas to every exceptional prime factor. & No unrepresented factor and no
finite-extinction input.\\
F4: combine torus systems & Form the union of the transported thick-boundary
and thin-refinement tori. & Pairwise disjointness, two-sidedness, and global
$\pi_1$-injectivity.\\
F5: form cut decomposition & Construct the global finite cut-piece family and
its side labels from the M1c interfaces. & Exact boundary incidence and
coverage of each prime carrier.\\
F6: reconstruct & Glue all cut pieces and use package transport to construct
an orientation-preserving diffeomorphism to $M$. & \code{PackageCoverage},
\code{PackageBoundaryCompatibility}, and \code{PackageTransport}.\\
F7: package certificate & Store prime data, torus data, cut pieces, geometry
witnesses, and reconstruction in \code{GeometrizationCertificate M}. & Every field
comes from F1--F6; no new existential is introduced.\\
\end{longtable}

The order matters. In particular, F4 cannot be performed by taking a union of
unlabelled subsets, and F6 cannot be proved from local embeddings without the
coverage and transport fields.

\subsection{Intermediate theorem contracts}

The implementation should expose the following intermediate declarations:

\begin{enumerate}[label=\textbf{F\arabic*:},leftmargin=*]
\item \code{normalize_thick_regions} returns a finite family of M1c cut pieces and
  M1d hyperbolic witnesses indexed by the thick data;
\item \code{normalize_thin_regions} returns a finite family of M1c cut pieces,
  internal torus data, and model witnesses indexed by the thin data;
\item \code{convert_exceptional_primes} returns the remaining model witnesses using
  only the public \code{PoincareAdapter} and standard-model lemmas;
\item \code{combine_global_torus_system} proves the disjointness, side, and
  incompressibility fields for the union of all retained tori;
\item \code{global_cut_decomposition} constructs the finite cut-piece family and
  proves its side incidence and complement coverage;
\item \code{global_reconstruction} constructs the orientation-preserving
  diffeomorphism from the assembled carrier to $M$; and
\item \code{certificate_of_assembled_data} fills the endpoint structure and proves
  that its fields agree with the previous six outputs.
\end{enumerate}

The names may change during the API audit, but the separation must remain.
Combining all seven stages into one opaque proof would make it impossible to
identify whether a later failure comes from graph-manifold conversion,
incompressibility, or global coverage.

\subsection{Dependency and trust boundary}

The M1f module may import only the M1 topology, geometry, package, and PC
adapter interfaces. It must not import the future module that proves
\[
  \texttt{RicciFlowSurgeryPackage M}
  \longrightarrow
  \texttt{LongTimeOutcome M},
\]
otherwise the structural theorem would be coupled to the analytic producer
and the intended dependency direction would be lost.

The final M1f proof receives a \code{#print axioms} audit. Temporary wrapper
interfaces may be present in a sandbox, but the release target must replace
them with proofs or explicitly approved foundational declarations. In
particular, \code{PackageNonCircular} is a design check and not a permanent axiom.

\subsection{M1f acceptance matrix}

M1f is accepted when the following are kernel-checked:

\begin{itemize}[leftmargin=*]
\item the four standard model examples and the graph-manifold mock package;
\item the four negative package tests for missing coverage, incorrect
  orientation, unmatched boundary side, and circular certificate fields;
\item the explicit \code{PoincareAdapter} argument appears in both structural
  theorem signatures;
\item the assembly module has no Ricci-flow, surgery, Perelman, or
  finite-extinction import; and
\item the theorem inventory records the exact declaration, imports, and axiom
  closure for every F-stage.
\end{itemize}

The result is a conditional geometrization theorem: every future analytic
producer that constructs an M1e package automatically obtains \code{Geometrizes M}.
This is the correct handoff to the long-time Ricci-flow workstream.

The field and theorem audit for this layer is recorded in
\code{m1_assembly_api.csv}.
% END migrated B022

% BEGIN migrated B064
\section{Contract-phase GC vertical slice}
\label{ch:m1-slice}

The topology, geometry, and package contracts in this slice can be developed
while the remaining PC proofs are being finished. The final exceptional-prime
conversion waits for the pinned public adapter, but no PC implementation
module or finite-time extinction theorem is needed to design the interface.
The purpose is to make the endpoint and analytic output precise enough that
the adapter has a fixed place to enter.

\subsection{Four interfaces}

\begin{description}[leftmargin=3.4cm,style=nextline]
\item[Endpoint] \code{GeometrizationCertificate M} contains a finite prime-piece
  record, a finite disjoint system of embedded incompressible tori, the cut
  pieces and boundary identifications, and a geometry witness for each piece.
  The witness identifies one of the eight model geometries and supplies the
  complete locally homogeneous metric on the interior, with finite volume
  required when its model tag is hyperbolic.
\item[Assembly input] \code{ThickThinCollapsePackage M} contains finite thick
  pieces with hyperbolic metric data, finite thin pieces with graph-manifold
  data, embeddings or correspondence maps into the relevant flow slice, and
  the boundary and gluing data needed to reconstruct the original manifold.
  It does not itself assert the final certificate; the assembly theorem must
  construct that certificate.
\item[Analytic producer] \code{RicciFlowSurgeryPackage M} contains the initial
  metric, smooth flow intervals, typed surgery events, continuation and
  topology maps and audited analytic prerequisites. The separate theorem
  \code{long_time_outcome} constructs \code{LongTimeOutcome M}; it is not
  a stored output interface. The late branch yields a thick/thin package;
  the terminal branch supplies explicit exceptional data.
\item[PC adapter] \code{PoincareAdapter} is an independent import boundary. It
  exposes only the completed Poincar\'e theorem and the spherical consequence
  actually consumed by the endpoint assembly.
\end{description}

\subsection{The two structural theorems}

The first production assembly theorem to implement after K1 is:
\begin{equation}
\begin{aligned}
 &\texttt{assemble\_geometrization} : \\
 &\quad \texttt{PoincareAdapter} \\
 &\quad\longrightarrow \texttt{ThickThinCollapsePackage M} \\
 &\quad\longrightarrow \texttt{GeometrizationCertificate M}.
\end{aligned}
\end{equation}
It depends on the definitions of prime pieces, incompressible torus cuts,
graph manifolds, Seifert pieces, and the eight geometries, but not on Ricci
flow estimates. The adapter argument is used only for exceptional prime
conversion.

The second theorem is the analytic integration theorem:
\begin{equation}
\begin{aligned}
 &\texttt{geometrization\_of\_ricci\_flow\_surgery} : \\
 &\quad \texttt{PoincareAdapter} \\
 &\quad\longrightarrow \texttt{RicciFlowSurgeryPackage M} \\
 &\quad\longrightarrow \texttt{Geometrizes M}.
\end{aligned}
\end{equation}
Its proof pattern-matches \code{LongTimeOutcome}: use assembly on the
late package and exceptional/model plus reconstruction lemmas on terminal
data, supplying the endpoint adapter where needed. This prevents the final theorem from depending on the internal shape of
the Ricci-flow proof.

\subsection{Independent work chunks}

The slice creates four parallel tasks with explicit handoff points:
\begin{enumerate}[label=\textbf{V\arabic*:},leftmargin=*]
\item \textbf{Topology contract.} Define the endpoint certificate, prime
  pieces, torus systems, graph-manifold and Seifert predicates, and the
  eight-geometry witnesses.
\item \textbf{Collapse-output contract.} Read the selected Morgan--Tian and
  Kleiner--Lott statements and express the finite thick/thin package,
  boundary collars, gluing maps, and graph-manifold output without proving
  the analytic theorem yet.
\item \textbf{Flow/surgery contract.} Define the typed global history and
  the exact fields required from the already available PC Ricci-flow and
  standard-cap APIs. Mark every field as imported, adapter, or new.
\item \textbf{Assembly proof.} Prove \code{assemble_geometrization} from the
  package fields, the explicit adapter argument, and standard topological
  lemmas after K1. The standalone K0 checks precede this production target.
\end{enumerate}

The slice is accepted when the interfaces and the two structural theorems
compile in a clean Lean environment, the theorem inventory lists every field
dependency, and no interface hides a finite-extinction assumption. The later
PC release then fills the adapter and audited analytic fields without forcing
us to redesign the endpoint.
% END migrated B064

\section{Current assembly input and exact scope}
\label{sec:assembly-current-input}

Sections16.1--16.7 retain the M1 specification history. The GA contracts
below give its current detailed realization and supersede ambiguous
informal phrases such as transporting an entire capped prime into the
initial manifold. They preserve the existing endpoint and module boundary.
Their inventory is {\upshape\code{global_assembly_contracts.csv}}. Written arguments
are conditional mathematics; no production Lean implementation is claimed.

\begin{definition}[Pure marked assembly data; GA01]
\label{def:assembly-marked-input}
For the fixed closed connected oriented smooth initial manifold $M$,
take the \emph{supplied topological fields} of AT30: a finite chosen
family of oriented closed prime carriers $P_i$; their actual oriented
connected-sum reconstruction of $M$; on each $P_i$ a finite collared
incompressible torus system; the actual compact cut carriers $Q_{i,j}$;
their boundary-side pairings and reconstruction maps; and specified
smooth identifications with the model carriers used for geometry.
All maps refer to these same indexed carriers. The record contains no
{\upshape\code{GeometrizationCertificate}}, no {\upshape\code{Geometrizes}} field, and no
function returning either conclusion.
\end{definition}
This record forgets the flow, its times and event provenance after those
have produced the topological fields. The fields can also be supplied
without a flow. A \emph{use} of supplied AT30 data is not a dependency on
the theorem producing AT30. Its production in a late or terminal branch
is recorded separately below. Complete model metrics are still separate
inputs. Essential torus maps live in the actual capped primes $P_i$;
punctured prime carriers and their embeddings used in connected sums do
not replace those endpoint objects.

\begin{lemma}[Finite dependent indexing without accidental identification; GA02]
\label{lem:assembly-finite-indexing}
If $I$ is finite and every cut index type $J_i$ is finite, the tagged
family $\coprod_{i\in I}J_i$ is finite. Relabelling each index by a
bijection transports all dependent carrier, map and witness fields.
Equal topological types do not identify two indices.
\end{lemma}
\begin{proof}
Enumerate the finite outer type, and append the tagged finite enumerations
of its fibers. For reindexing, pull each field back along the inverse
bijection, with the corresponding carrier equality or specified
diffeomorphism. In particular two copies of $\mathbb{RP}^3$ remain two
entries, and sphere-product factors coming from different remaining
cycles keep their provenance. No sorting of model names changes the
connected-sum multiplicities.
\end{proof}
Before optional removal of identity summands, retain an actual $S^3$
carrier in the spherical singleton case. The empty disjoint union is
not an $S^3$ component. An empty thick family or empty torus family is
allowed and does not imply an empty prime or cut-piece family.

\section{Torus systems, sides and exact cut carriers}
\label{sec:assembly-tori-and-cuts}

\begin{lemma}[Paired sides and one seam label; GA03]
\label{lem:assembly-side-orbits}
Let the finite boundary-side set $S$ have a fixed-point-free involution
$\sigma$, with inverse pairing diffeomorphisms
$\phi_{\sigma s}=\phi_s^{-1}$. Its orbits form a finite seam set, every
orbit having two sides. An actual seam embedding and its two collar
restrictions are supplied by GA01; the side relation alone does not
construct an embedding or prove its incompressibility.
\end{lemma}
\begin{proof}
The equivalence class of $s$ is exactly $\{s,\sigma s\}$ by involutivity
and has two elements by the no-fixed-point condition. Pass to this finite
quotient, transporting the actual seam and collar maps from the given
cut record. Replacing $s$ by $\sigma s$ exchanges the sides and inverts
the pairing. A seam may join two different boundary components of the
same cut carrier: this is a loop in the incidence graph, not a fixed
side. Parallel edges are likewise retained.
\end{proof}

\begin{lemma}[The finite essential system on the chosen primes; GA04]
\label{lem:assembly-essential-system}
GA01's indexed seam maps, together with its relative internal refinement
maps, form a finite disjoint two-sided incompressible torus system on
each chosen $P_i$, provided the supplied comparison identifies duplicate
descriptions of the same seam and all distinct seam collars are disjoint.
\end{lemma}
\begin{proof}
Use GA03's single seam label for the two descriptions of an external
boundary. Retain each internal refinement label once, with its given
map. The joint collar-disjointness hypothesis separates all distinct
images and supplies two-sidedness. The actual based inclusion maps
already have injectivity proofs from AT17 or its supplied structural
analog. Compose them with the recorded diffeomorphisms using AT09.
There is no inference from injection in a model or from a graph tag.
\end{proof}
This instantiates M1 F4. Independently disjoint thick and thin lists do
not imply disjointness between the lists. Two indices with the same
image cannot be retained as distinct members of an embedded torus system.

\begin{lemma}[The component meaning of the cut family; GA11]
\label{lem:assembly-cut-components}
Suppose the supplied cut/reconstruction comparison identifies
$\coprod_j Q_{i,j}^{\circ}$ with the complement of the selected tori in
$P_i$, and each displayed interior is nonempty and connected. Then the
images are exactly the connected components of that complement.
\end{lemma}
\begin{proof}
Each summand is open and closed in the disjoint union; under the
comparison its image is an open-and-closed connected subset. A connected
component meeting it cannot leave it, and connectedness prevents the
image from meeting two components. Surjectivity accounts for every
point, so no component is omitted.
\end{proof}
The cut interface supplies this actual complement comparison. If an
earlier construction deletes a positive collar, use its specified collar
comparison diffeomorphism first; that excised subset is not literally
$P_i$ minus its tori. A local list of embeddings supplies neither the
comparison nor coverage. This is F5's precise topological obligation.

\section{Assigning geometry to the actual interiors}
\label{sec:assembly-metric-assignment}

\begin{lemma}[Thick normalization; GA05]
\label{lem:assembly-thick-normalization}
For a supplied hyperbolic core $C$ with HG02 data and an actual
diffeomorphism $d:Q_{i,j}\to C$ compatible with the recorded cut labels,
the metric $(\psi\circ d|_{Q_{i,j}^{\circ}})^*h$, where $\psi$ is HG16,
is a complete finite-volume hyperbolic metric on all of $Q_{i,j}^{\circ}$.
\end{lemma}
\begin{proof}
The composite is a diffeomorphism to the entire complete model $H$.
Pullback preserves completeness, local model charts and total volume.
Its domain is the whole interior, not the image of a finite truncation
inside a late slice. Boundary matching uses $d$ and the core collars;
the map $\psi$ is not substituted for the core's late embedding.
\end{proof}
The proof consumes supplied static cusp data. It does not use HG08's
flow producer, HG15's analytic proof, or Mostow rigidity. Those are
dependencies of producing the thick input, not of transporting its metric.

% BEGIN REVISION171 GA06 CURRENT BINDING
\paragraph{Current graph application binding.}
The historical GA06 producer below is supplied for the selected E1
application by Theorem~\ref{thm:graph-marked-geometric-endpoint} and
Corollary~\ref{cor:graph-audited-late-binding}. The actual chosen prime
and good-block cuts have complete metrics with all markings retained.
Optional uncut Sol/semibundle classification and canonical coarsening
are unnecessary on this route. GA07 for other exceptional producers
and the event-history inputs are not supplied by this graph result.
% END REVISION171 GA06 CURRENT BINDING

\begin{foundation}[Thin refinement and model realization; GA06]
\label{found:assembly-thin-models}
On each supplied graph region, the Chapter5--6 relative producers must
give the actual good-block/essential-cut refinement with its retained
external labels and reconstruction. Chapter7 G09 must give actual
complete model metrics on the resulting entire interiors. Include
Seifert cases and the Euclidean, Nil, Sol, product and twisted-bundle
alternatives with their precise hypotheses.
\end{foundation}
Raw auxiliary graph seams are not accepted as essential cuts. Closed
Seifert classification is not a theorem about every open Seifert
interior. E1 is unchanged: no extra Klein-bottle cut is made, and only
hyperbolic tags mandate finite volume. GA06 is a qualified producer,
not a proof that every named graph block already has the desired metric.

\begin{foundation}[Exceptional realization on actual factors; GA07]
\label{found:assembly-exceptional-models}
For every exceptional prime in the actual reconstruction, use AT25 and
the narrow {\upshape\code{PoincareAdapter}}/Chapter7 interfaces to obtain its
specified orientation, prime witness and actual complete model geometry.
Spherical space forms require a genuine free isometric action or the
proved equivalent geometric recognition, not just a finite group.
\end{foundation}
The simply connected Poincar\'e conclusion is not silently strengthened
to all spherical-space-form recognition. The latter is a separately
audited interface. The sphere product uses $S^2\times\R$ geometry.
No finite-extinction theorem or internal PC flow module is an input to
this conversion. Prescribed orientation changes do not assume that a
chiral model admits an orientation-reversing self-diffeomorphism.

\begin{definition}[A typed family of geometric witnesses; GA08]
\label{def:assembly-geometry-family}
For every actual cut index $(i,j)$ choose a model tag $G_{i,j}$ and a
metric $g_{i,j}$ on all of $Q_{i,j}^{\circ}$, together with completeness,
local charts to that model, the required end/collar comparison and
\[
 G_{i,j}=\text{\normalfont\ttfamily Hyperbolic}\ \Longrightarrow\
       \operatorname{Vol}_{g_{i,j}}(Q_{i,j}^{\circ})<\infty.
\]
All witnesses refer to this same metric and this same carrier. The role
assignment to thick, refined thin and exceptional inputs is exhaustive;
if a region description is duplicated, its comparison maps are supplied.
\end{definition}
This does not impose finite volume on the other seven tags, and does
not permit attaching the finite volume of one metric to another metric's
completeness proof. Boundary gluing maps are topological/smooth maps;
they need not be isometries of these independently assigned metrics.

\begin{lemma}[Transport along the actual comparison; GA09]
\label{lem:assembly-model-transport}
Given an actual diffeomorphism $e:Q_{i,j}\to B$ and a complete model
metric $h_B$ on all of $B^\circ$, set $g_{i,j}=(e|_{Q_{i,j}^{\circ}})^*h_B$.
This transports the geometry, completeness and total volume. If boundary
labels or orientations are prescribed, their compatibility is a separate
part of the marked comparison.
\end{lemma}
\begin{proof}
Apply AT26 with the inverse direction of the diffeomorphism. Its length,
Cauchy and measure arguments give exactly these claims. Transport the
specified model charts by composition. A bare homeomorphism is not a
pullback map for a smooth tensor field; the inherited category adapter
must first produce the needed smooth comparison.
\end{proof}

\begin{lemma}[Finite witness selection and normalization; GA10]
\label{lem:assembly-finite-witness-choice}
If GA05--GA07 or other proved model constructors supply a nonempty set
of correctly typed witnesses for every cut index, a single GA08 family
can be chosen without any uniqueness assumption.
\end{lemma}
\begin{proof}
Induct over a finite enumeration from GA02, choosing one witness for the
next index and keeping the previous choices. Transport through GA09
whenever its actual carrier comparison is needed. Each fiber is tied to
its index, so the result is one dependent family, not an unlabelled list
of independently chosen metrics. Classical finite choice suffices.
\end{proof}
This completes F1--F3 conditionally on their input constructors. It does
not prove GA06 or GA07 by choosing from an unproved nonempty set.

\section{Smooth reconstruction and certificate construction}
\label{sec:assembly-reconstruction-proof}

\begin{lemma}[Cut-side gluing and prime reconstruction; GA12]
\label{lem:assembly-prime-reconstruction}
The actual pairing maps of GA01 reconstruct each chosen $P_i$ from its
compact cut carriers. If the carriers have been replaced by marked
diffeomorphic representatives, conjugate each pairing by those maps and
transport its collar. The resulting reconstruction is an oriented
smooth diffeomorphism.
\end{lemma}
\begin{proof}
On each carrier use its recorded comparison. The conjugation relation
$\phi'_s=e_{\sigma s}\phi_s e_s^{-1}$ makes these maps agree on each
identified boundary pair, so they descend to the quotient. The inverse
maps descend as well. In the transported product collars this is a
smooth map with smooth inverse across the seam; orientation follows
from the oriented carrier maps and opposite induced boundary signs.
Compose with the supplied cut reconstruction to $P_i$. This is AT07.
\end{proof}
Equality of boundary values alone does not identify normal derivatives.
For preassigned collars, supply the inherited straightening/comparison
theorem; it cannot be omitted from a smooth gluing claim. No global
Riemannian metric across the seams is asserted.

\begin{lemma}[Connected-sum reconstruction of the initial carrier; GA13]
\label{lem:assembly-global-reconstruction}
Compose GA12's prime reconstructions with GA01's actual oriented
connected-sum reconstruction to $M$, preserving the sphere-port ledger
and the joint disjoint decoration. This produces the endpoint's global
reconstruction witness on the fixed initial carrier.
\end{lemma}
\begin{proof}
Pull the selected attachment balls and their parametrizations back along
the prime diffeomorphisms. Naturality of the collared sphere gluing
gives a diffeomorphism of the connected-sum constructions; compose with
the supplied one to $M$. If the attachment data were preassigned instead,
use their explicit marked comparison. AT14's simultaneous avoidance
ensures that retained tori and collars miss the selected balls.
Identity summands and remaining cycle factors are handled by the
already normalized AT20/AT24 ledger, not counted again. Every composition
preserves the specified orientation.
\end{proof}
The metric-bearing endpoint pieces remain cuts of the capped primes.
Puncturing those primes for connected sums does not restrict or replace
their complete interior metrics. No embedding of an entire capped prime
into $M$, and no global reverse-time flow map, is required.

\begin{theorem}[Certificate from assembled data; GA14]
\label{thm:assembly-certificate}
Given GA01 and the realized GA08 metric family on its actual cut carriers,
the constructions GA02--GA13 produce {\upshape\code{GeometrizationCertificate M}}.
All source-qualified producers used to obtain those inputs remain outside
this theorem's conclusion.
\end{theorem}
\begin{proof}
Use the supplied finite prime family and its prime proofs. GA04 supplies
the collared essential torus system; GA11 identifies the cut interiors
with the actual complementary components. Retain GA03's complete side
pairing, GA12's prime reconstruction and GA13's global reconstruction.
Insert the GA08 metrics, completeness, model charts and conditional
hyperbolic volume proofs, transported by GA09 where necessary. Every
endpoint field is now a named supplied or constructed value. There is
no use of a pre-existing endpoint certificate or an existential over it.
\end{proof}

\begin{corollary}[The public predicate wrapper; GA15]
\label{cor:assembly-predicate}
The same data imply {\upshape\code{Geometrizes M}}.
\end{corollary}
\begin{proof}
Use GA14's constructed certificate as the witness in the defining
existential predicate. No additional choice of a decomposition is made.
\end{proof}

\begin{lemma}[Reindexing and chosen representatives; GA16]
\label{lem:assembly-reindexing}
Finite reindexing, together with actual oriented carrier diffeomorphisms
and their commuting collar/pairing/reconstruction squares, transports
GA14's input and output. In particular the public existential statement
does not require a canonical prime ordering, minimal torus cuts or
uniqueness of the geometric decomposition.
\end{lemma}
\begin{proof}
Apply GA02 to the finite indices, GA09 to the metrics and GA12--GA13 to
the actual reconstruction maps. Incompressibility transports by AT09.
The fields then satisfy the same endpoint definition. If only a chosen
prime family is required, use it directly. Changing to arbitrarily
prescribed boundary-marked primes instead requires AT27; ordinary prime
uniqueness alone does not supply those markings.
\end{proof}

\section{The M1 package theorem and the two outcome branches}
\label{sec:assembly-package-and-branches}

\begin{foundation}[Structural normalization on the reference primes; GA17]
\label{found:assembly-package-normalization}
For a supplied {\upshape\code{ThickThinCollapsePackage M}} and the narrow
{\upshape\code{PoincareAdapter}}, the proved relative refinement and geometry
constructors must produce GA01 and GA08 on the package's reference prime
family. Use its actual coverage, boundary compatibility, carrier
comparisons and exceptional data. GA06/GA07 and the Chapter5--7
conversions remain qualified until proved or bound to accepted imports.
\end{foundation}
This is the precise remaining producer behind structural normalization.
Its output is the explicit intermediate data, not a hidden endpoint
certificate. It does not use a flow theorem. A package constructed from
AT30 chooses that same prime family as its reference; an unrelated
prescribed family would require the separate AT27 comparison.

\begin{theorem}[Conditional M1 assembly; GA18]
\label{thm:assembly-m1-composition}
Once GA17's structural conversion lemmas are proved, the default M1
theorem has the retained signature
\[
\begin{aligned}
&\text{\normalfont\ttfamily PoincareAdapter}\longrightarrow
\text{\normalfont\ttfamily ThickThinCollapsePackage M}\\
&\hspace{2em}\longrightarrow
\text{\normalfont\ttfamily GeometrizationCertificate M}.
\end{aligned}
\]
\end{theorem}
\begin{proof}
Apply GA17 to construct the intermediate input and its realized metrics,
then apply GA14. GA15 gives {\upshape\code{geometrizes_of_package}}. The explicit
adapter is used at exceptional conversion, not to read PC flow internals.
The theorem is conditional on the availability of the stated library
proofs; none becomes an unreviewed axiom in a release.
\end{proof}
Pure GA14 also applies to independently supplied model witnesses without
an adapter argument. That convenience does not erase the default GA18
dependency or assert that exceptional realization has been completed.

\begin{proposition}[Late-outcome conversion; GA19]
\label{prop:assembly-late-conversion}
Given a produced late outcome with its raw region ledger, constructed
AT30 fields and actual comparison/coverage data, form the existing
{\upshape\code{ThickThinCollapsePackage M}}. Use AT30's chosen primes as the
reference decomposition and the supplied thick and thin region records
as its region fields.
\end{proposition}
\begin{proof}
Copy the constructed prime, carrier and metric data with their exact
indices; retain the raw late-region correspondence and the already
proved boundary/coverage witnesses. This is a conversion of supplied
data. It does not rerun LP08, manufacture late regions from AT30 alone,
or project {\upshape\code{PackageTransport}} before the package exists. GA18
then applies when its qualified structural constructors are available.
\end{proof}

\begin{proposition}[Terminal certificate without a fake late package; GA20]
\label{prop:assembly-terminal-certificate}
Given LP11's explicit terminal data, AT28/AT30 reconstruction and the
realized exceptional metrics, GA14 gives the endpoint certificate
directly. No late slice, thin-collapse theorem or thick hyperbolic
producer is required on this branch.
\end{proposition}
\begin{proof}
Use the actual terminal prime factors and their model identifications;
the standard spherical and sphere-product factors admit the empty torus
systems in this reconstruction. Take the realized whole-factor model
metrics from GA07 and transport them by GA09. AT28 supplies the exact
global reconstruction; GA14 fills the fields. The empty selected terminal
slice is not declared to be the nonempty initial carrier. Do not create
a fictitious {\upshape\code{LateThickThinOutcome}} with missing late-region fields.
\end{proof}

\begin{theorem}[Outcome elimination; GA21]
\label{thm:assembly-outcome-elimination}
After the conditional branch constructors are available, a supplied
{\upshape\code{LongTimeOutcome M}} and {\upshape\code{PoincareAdapter}} yield the endpoint
certificate, and hence {\upshape\code{Geometrizes M}}.
\end{theorem}
\begin{proof}
Eliminate the outcome disjunction. For a late value use GA19 and GA18;
for a terminal value use GA20. A particular application supplies one
branch value. Proof dependencies include both branch functions, not a
premise that both outcomes occur.
\end{proof}

\section{Analytic integration and the public theorem}
\label{sec:assembly-public-integration}

\begin{theorem}[Integration for a supplied admissible flow; GA22]
\label{thm:assembly-flow-integration}
After LP12 has been proved against the accepted global-flow interface,
and the GA21 branch constructors have been proved, the existing
{\upshape\code{geometrization_of_ricci_flow_package}} follows for each supplied
{\upshape\code{RicciFlowSurgeryPackage M}} and the narrow adapter.
\end{theorem}
\begin{proof}
Invoke the proved LP12 producer on the actual supplied flow. Apply GA21
to the resulting outcome and then GA15's predicate wrapper. The producer
is a theorem dependency, not an output oracle in the flow record. The
proof uses no assumption that late nonempty slices always exist.
\end{proof}
The historical name {\upshape\code{geometrization_of_ricci_flow_surgery}} denotes
this same intended boundary after the exact signature is aligned; it
does not name a second independent Geometrization theorem.

\begin{theorem}[Public Geometrization composition; GA23]
\label{thm:assembly-public-geometrization}
For every closed connected orientable smooth three-manifold $M$, once
the accepted inherited foundation and the qualified geometric/analytic
producers listed here are proved, the public theorem
{\upshape\code{geometrization}} constructs {\upshape\code{Geometrizes M}} without an extra
flow-existence or late-availability hypothesis in its public statement.
\end{theorem}
\begin{proof}
Choose the orientation and the initial smooth Riemannian metric through
the inherited interfaces, with their compact normalization argument.
Invoke {\upshape\code{GlobalRicciFlowSurgery}} to obtain an actual admissible flow
package and the finite-prefix/global-domain evidence required by LP01.
Apply GA22 with the accepted narrow adapter. All these are proof
dependencies of the public theorem. A conditional implication from a
supplied flow alone does not establish the required existence.
\end{proof}
This is a written composition proof with qualified upstream producers,
not a completed proof of Geometrization or an accepted Lean theorem.
It neither pins mathlib nor restarts the inherited foundation. M2A stays
release-gated and only surviving M3 PDE gaps are new analytic work.

% BEGIN REVISION172 IMPLEMENTATION EVIDENCE
\section{Current implementation evidence}
\label{sec:implementation-revision172}

The active task view is versioned independently of historical prose.
Revision172 retains the refutation of the unrestricted FC28 statement
and selects CAA02's written static replacement. Its audited written
completion does not imply a Lean implementation. The old FC28 status
is preserved explicitly as history, rather than offered as the current
work frontier. Every exported packet now hashes its actual three source
files and all attached evidence records; recorded Lean evidence also
binds the external implementation files and build outputs.

The immutable public PC source is tag \code{v0.1.3} of
\url{https://github.com/qinz1yang/differential-geometry}, commit
\code{7a48598d35109aa99d1cc678e2724c213cdf4ff3}.
It pins Lean 4.33.1 and Mathlib commit
\code{0df444a360eaa60ab8c11dca51a86af692955474}.
A fresh build of its public PC target was started in an isolated checkout.
The build is still pending at the revision172 evidence freeze; cached
upstream Mathlib artifacts were used. This is not yet a reproduced PC
endpoint or a general global-flow binding. Simply connected surgery
exports must not be silently applied to the general GC manifold.

The separate modules \code{PointedBallApprox.lean} and
\code{PointedBallConsumers.lean}, in
\code{GC_PROOF_PILOTS/revision172}, implement the existing closed-ball
contracts on the pinned Mathlib metric-space representation. MC06 stores
the actual map, positive strict error, exact basepoint, pairwise distortion
and actual coverage. MC07 proves that the selected coverage witnesses
survive restriction. MC08 checks every intermediate image's domain.
MC09 selects one lift into the original source ball, forces the exact
basepoint, and proves both inverse bounds and the sharper
$3\varepsilon$ distortion before packaging the reverse approximation.
No continuity or compactness assumption was added to these contracts.

The downstream constructor takes two actual approximations from a common
source at radius $4j+4$, with $j\geq1$ and $10\varepsilon<j$.
It obtains MC24's comparison at radius $j$ and error $10\varepsilon$ by
MC09 followed by MC08, checking all domain inequalities in Lean.
A real-line instance, a closed target-boundary coverage instance and an
impossibility proof at equal error and radius accompany it. These are
bounded proofs and consumers, not a formalization of the subsequent
proper-limit argument, collapse, or Geometrization.

Both modules compile successfully. The ten queried declarations have
exactly the three standard classical axioms; there is no added geometric
axiom or admitted pilot proof. Three failed consumer attempts are retained
as elaboration and real-number normalization examples; none required
weakening a mathematical hypothesis. The versioned implementation record
keeps these accepted atoms separate from unimplemented parent tasks.
The exact source files, commands, hashes, axiom outputs and scope review
are recorded in \code{formalization/implementation_bindings.json} and
\code{REVISION172_CHECKS.md}. Static consistency and provenance checks
are evidence about the records, not substitutes for these Lean runs.
% END REVISION172 IMPLEMENTATION EVIDENCE

% BEGIN REVISION173 ACTIVE GRAPH
\section{Selected proof graph and interface audit}
\label{sec:active-graph-revision173}

The consolidated graph in \code{formalization/active_graph.py} retains
all earlier catalogs and promotes the later MPR, MGL, LFR, CFS, actual
cloud, long-time, ambient and graph-realization arguments to individual
nodes. Its selected-route record gives each reviewed integration edge
its consumer premise, producer output, transport requirements and source.
Other inherited edges remain explicitly unreviewed; a source-qualified
leaf with no recorded predecessors is an open obligation, not a theorem
accepted by default. The current census has 913 nodes and 2443 edges,
with 649 nodes on the selected route and 124 selected open leaves.
These are planning counts, not counts of independent mathematical gaps.
Five atomic implementation entries have separate checked evidence.

The first backward pass makes the following distinctions explicit.
GA23 requires an actual global-flow existence producer, initial metric
choice and the narrow PC adapter, in addition to GA22's conditional
composition. LP12 and GA21 depend on both conditional branch functions;
a particular outcome contains only one branch value. The LP11/GA20
terminal route needs the actual finite history and its exceptional
factors, with no collapse or ambient-cusp premise. GAU03 supplies only
late-slice data; LP08 must still use AT31 and the actual AT21/AT22/AT25
inputs to reconstruct the initial manifold.

QRG09 has two successive conclusions. Its topological part constructs
the global diffeomorphism before the final invocation of QRG03.
The graph therefore gives QRG00 the proof fragment
\code{QRG09.topology}, whose dependencies exclude QRG03. The packaged
QRG09 then uses QRG03. A test that substitutes the entire packaged
QRG09 as QRG00's prerequisite detects the resulting artificial cycle.
No new topological proof or stronger quantitative hypothesis is asserted.

The old FC33 dependency list still led through the refuted unrestricted
FC28. On the selected actual construction, GAF02 expressly supplies
FC33's original threshold-$5$ bases and submersions, and is now its
producer. Separately, BSA01 uses FC43's \emph{given collar hypotheses};
it does not use a completed boundary assembly theorem to prove its
first-exit estimate. That edge is recorded as definition use. FC28,
CFS19's unused tolerance and optional AT27 are excluded from the
selected proof ancestry. These repairs change the current dependency
interpretation, preserving the historical proof text and old records.

A second checked metric experiment imports the public library's
\code{SmoothRiemannianMetric}, \code{scaleMetric} and
\code{riemannianClosedBallOf}. It reuses the proved tensor/distance
scaling theorem and obtains, for $r>0$,
\[
 \overline B_{r^{-2}g}(p,R)=\overline B_g(p,rR),
\]
with the corresponding padded-domain inclusion when $R\leq S$.
For the ordinary radius interpretation take $R,S\geq0$; the upstream
extended-distance definition uses \code{ENNReal.ofReal} for radii.
The same actual metric and manifold occur on both sides. Five queried
upstream/consumer declarations have only the three standard classical
axioms. This binds a narrow MC10/HG01 interface, not the full recentering
contract or an intrinsic-versus-ambient ball comparison.

The exact imported source statements and bodies were checked in the
pinned release: \code{Distance/Basic.lean} for extended distance,
\code{Comparison/DistanceScaling.lean} for the square-root factor,
\code{Distance/Ball.lean} for the closed-ball scale equality, and
\code{Scaling.lean} for the tensor multiplier. No simply-connected,
compactness or completeness hypothesis was needed by this consumer.
General-flow existence, the one-profile time uniformity and the marked
finite-history and exceptional-model bindings require their own source
comparisons and Lean consumers. Acyclicity and these bounded successes
do not finish the whole-book logical or implementation audit.
% END REVISION173 ACTIVE GRAPH

% BEGIN REVISION174 MARKED CONSUMERS
\section{Checked marked-map consumers and release comparison}
\label{sec:marked-consumers-revision174}

The five modules in \code{GC_PROOF_PILOTS/revision174} implement narrow
parts of AT09, AT11, GA02 and GA03 using the pinned PC release's actual
fundamental groups and Mathlib finite quotients. The same build also
instantiates revision173's metric normalization on the PC Euclidean
smooth metric on $\R^3$. Its twenty-nine queried declarations use only
subsets of the three standard classical axioms. The records distinguish
these accepted atomic outputs from their larger unimplemented parents.

For AT09 the target marking is an actual path from the chosen target
basepoint to the image of the source basepoint. The induced map is followed
by the inherited change-of-basepoint isomorphism. Functorial composition,
inner-map cancellation, kernel equality under an injective outer map,
and homotopy invariance with the homotopy's actual basepoint track are
checked. Thus the two torus inclusions into $T^2\times\R$ give the
same based map after the recorded vertical transport. Their injectivity
is proved using actual continuous retractions. The consumer also exhibits
a nontrivial torus loop; it does not rely on a simply connected test
carrier or assume the inclusions' injectivity.

For GA02--GA03 the side involution yields a finite quotient whose fiber
at a side is exactly its two-element pair. Tagged dependent indices
retain their actual map payloads under bijective reindexing. Two
same-type torus seams retain different labels, even when both sides of
a seam belong to the same carrier. The supplied pairing can be the
nonidentity interchange of the two circle factors. Commuting pairing
maps make the images descend to seam labels; they do not make the raw
parametrizations literally equal. The torus-cylinder consumer checks
image disjointness and the induced injections for both descriptions.
Smooth collar gluing, prime recognition and actual surgery-event
production remain separate obligations.

\paragraph{The inspected noncollapsing restriction.}
The PC source's general spatial-canonical volume lemma covers its neck,
cap and positive-volume alternatives under the recorded
\code{requiresVolume} condition. The inspected extension to the round
alternative uses simple connectivity of the component; the public
small-scale surgery producer invokes that extension. This identifies
a precise comparison task, not a declaration-absence result.
Perelman's Section4.4 removes almost-round spherical quotient components,
whereas KL Remark73.3 retains them. Removing those components is therefore
a possible adapter route only with actual removal events, AT22
provenance, AT25 witnesses and the required time/profile control.
No deletion convention is changed in this revision, and no conditional
volume lemma is counted as general flow existence.

\paragraph{Existing exceptional-model data.}
The release already defines a spherical factor by a finite free
orientation-preserving orthogonal action on the round three-sphere
and an actual oriented model diffeomorphism. Its standard connected-sum
presentation records the factor list and a smooth reconstruction map.
These are stronger reuse candidates than an abstract finite fundamental
group. The general cut/cap reconstruction theorem inspected here supplies
an actual homeomorphism with capped-component labels and sphere-product
multiplicity. The relative smooth collar/port adapter and the complete
model metrics on the endpoint's exact carriers still need their own
bindings. Exact source lines, hashes, comparison limits and remaining
consumers are recorded in \code{formalization/public_release_comparisons.json}
and \code{reference_checks_revision174.md}.
% END REVISION174 MARKED CONSUMERS

% BEGIN REVISION175 TOWER PROFILE
\section{Actual observation towers and the common-profile obligation}
\label{sec:tower-profile-revision175}

The three Lean modules in \code{GC_PROOF_PILOTS/revision175} consume the
public PC release's actual \code{ObservationTower}, including its
smooth carriers, metrics and initial identifications. Twelve transitive
axiom queries use only the three standard classical axioms. These are
bounded accepted outputs under an explicitly supplied tower; the
\code{GlobalRicciFlowSurgery} producer is still unaccepted.

\paragraph{Finite histories on the same flow.}
For every $b\geq0$, the inherited observation at $b$ has exactly the
event set $E\cap(0,b]$, and that set is finite. For every $B\in\R$,
the open interval $(\max\{B,0\},\max\{B,0\}+1)$ contains a non-event
time $t$. If the last recorded event time of the observation were $t$,
then $t$ would belong to its event set; the zero-event case is excluded
by $t>0$. Consequently its actual final smooth slab has positive time
length. The constructed consumer retains the tower's actual
\code{InitialIdentification}, rather than choosing a diffeomorphic
replacement manifold. Extending an observation horizon preserves the
slice carrier and its metric, with the inherited heterogeneous equality
recording the carrier transport.

Classical case analysis now gives either an empty observed carrier that
remains empty at all later horizons, with no strictly later events, or
nonempty regular slices beyond every prescribed bound. An event at the
empty observation's horizon is allowed. This supplies a narrow LP01
calendar/observation adapter and the corresponding availability
alternative. It supplies neither the input tower nor the exceptional
geometries needed for LP11, and does not claim extinction for general
initial data.

\paragraph{Two different parameter indices.}
In the inspected public source, \code{CutoffParameters} already contains
time functions such as $\delta(t)$ and the neck radius. However,
\code{HasUniformCutoffRecords} permits a family of such packages indexed
by the observation horizon $b$. Its record at an event of time $t$
therefore evaluates a potentially horizon-dependent profile $\delta_b(t)$.
The definition alone does not equate those profiles as $b$ varies.
This is a comparison of exact definitions, not a claim that no stronger
producer exists elsewhere in the library. Nor must every auxiliary
model parameter be fixed in order to bind the particular common-profile
fields that GC actually uses.

The following elementary obstruction is kernel checked. For each
$B\geq1$, the constant profile $\delta_B(t)=1/B$ satisfies
$0<\delta_B(t)\leq1/B$ on $[0,B]$. There is no single profile with
these properties for every $B$: at $t=1$, put
$B=\max\{1,2/\delta(1)\}$ to obtain
$2\leq B\delta(1)\leq1$. This refutes the quantifier shortcut,
not any geometric existence theorem or its actual bounds.

\paragraph{The source supplies the correct construction problem.}
KL Proposition77.2 and its discussion, printed2770--2773/PDF184--187,
choose nonincreasing parameter sequences and allow a profile below the
resulting bounds before constructing the flow. Printed2772 explains
why repeatedly estimating from time zero would force changes to an old
surgery parameter. Section80, printed2784--2787/PDF198--201, uses
successive time intervals and the preceding interval's noncollapsing
control. The implementation comparison must therefore bind that
successive-interval analytic construction, its actual event finiteness
and coherent restrictions; independent finite-horizon choices are
insufficient. KL Remark86.9's later smallness condition and IAU02's
fixed-profile convention are retained. MT Theorem1.2, archived
printed12--13/PDF16--17, remains the global comparison and reconstruction
spine, with its explicit topological hypotheses and discard list.

The inspected PC \code{FiniteTerminatingTower} constructs an absorbing
tower from supplied finite events and an empty last carrier. Its
\code{TimedEventTower} instead takes a chain with one event in each
specified unit interval. Neither conditional constructor alone is the
general GC existence theorem. Exact source locators, inspected bodies,
versions, hashes and remaining obligations are recorded in
\code{reference_checks_revision175.md} and the updated public-release
comparison. No discard convention or inherited theorem is replaced.
% END REVISION175 TOWER PROFILE

% BEGIN REVISION176 SPHERICAL CONSUMERS
\section[Spherical metrics and orientation correction]{Actual spherical metrics and factor-slot orientation correction}
\label{sec:spherical-consumers-revision176}

The pinned public PC library supplies more than the simply connected
endpoint. The three modules in \code{GC_PROOF_PILOTS/revision176}
bind its actual smooth quotient metrics, curvature naturality,
Riemannian completeness, positive-space-form recognition and orientation
correction. Four bounded implementation atoms have fourteen transitive
axiom queries, all within the three standard classical axioms. Their
parent AT25, GA07 and event-reconstruction contracts remain unaccepted.

\paragraph{The metric on the actual carrier.}
Given a supplied finite free round quotient model $Q$ and an actual
smooth diffeomorphism $f:M\to Q$, define $g=f^*g_Q$. The checked inner
product formula uses the differential of this very map. Curvature
naturality and the descended round metric give, for all tangent vectors,
\[
 \mathrm{Rm}_g(v,w,w,v)
   =g(v,v)g(w,w)-g(v,w)^2.
\]
Thus the normalization is sectional curvature one, not merely an
unspecified positive constant. Compactness gives completeness for the
closed quotient carriers. More generally, the inherited equality of
Riemannian distances under a whole-carrier diffeomorphism transports
completeness between Hausdorff sigma-compact smooth carriers, without a
compactness assumption. A map defined only on a compact core does not
satisfy that statement. The adapter allows source and target carriers
to inhabit different Lean universes.

The public theorem \code{sphericalSpaceFormCovering_holds} takes a
connected closed oriented three-manifold with an actual smooth metric
of constant positive sectional curvature. It produces a finite free
orientation-preserving orthogonal action on the round sphere and an
actual oriented quotient diffeomorphism. Applying the preceding
consumer gives a complete curvature-one metric on the original
carrier. No simple-connectivity premise occurs. A separate checked
consumer constructs the actual antipodal quotient, with this metric
and a nontrivial fundamental group on the same carrier. This tests the
quotient case rather than only the three-sphere. Arbitrary finite
fundamental group, primeness and discarded-component recognition are
not substituted for the geometric input.

\paragraph{Orientation without losing factor slots.}
Let $P$ be the inherited standard connected-sum presentation of a
connected closed oriented carrier $M$, with its actual reconstruction
map and finite factor list $L$. Orientation dichotomy for that map gives
either an oriented presentation with factor list $L$, or one with the
pointwise opposite list $L^{\mathrm{op}}$. The inherited finite-sum
opposite theorem supplies the latter reconstruction. Spherical factors
remain standard by reflection-conjugating their finite actions; the
argument does not assume that a chiral quotient has an orientation
reversing self-diffeomorphism. The inherited sphere-product reversal
handles the other standard factor type.

The new consumer retains the length of $L$ and, after the explicit
\code{Fin.cast} index transport, equality of the underlying carrier of
every factor slot. Repeated diffeomorphism types remain distinct slots.
This is the correct input for later dependent markings, but it does not
construct sphere/cap port maps or prove their commutative diagrams.
Actual event data must supply those maps. Nor does a complete model
metric by itself prove that the carrier is prime.

\paragraph{Evidence and remaining production.}
Exact source definitions, proof passages, line locators and hashes are
recorded in \code{reference_checks_revision176.md}. Three small upstream
modules were freshly compiled into separate audit artifacts; the
concurrent public endpoint build's files were not overwritten. The
fresh consumer build, fourteen axiom outputs and auxiliary provenance
are recorded in \code{formalization/spherical_bindings.json}. Nine
negative controls include rejection of an omitted auxiliary artifact.
The next producer comparisons remain the actual discarded carriers,
prime and product-model witnesses, marked event reconstruction and the
general coherent flow with its common profile. This iteration does not
assert whole-book dependency closure or the Geometrization endpoint.
% END REVISION176 SPHERICAL CONSUMERS

% BEGIN REVISION177 EVENT RECONSTRUCTION
\section[Actual surgery reconstruction]{Actual smooth event reconstruction and component labels}
\label{sec:event-reconstruction-revision177}

The public source comparison now binds the stronger concrete theorem
\code{exists_diffeomorph_finiteConnectedSum_capComponents}. Given the
inherited actual spherical cut/cap transition $E:M\leadsto Q$, it
constructs a finite family of smooth connected sums and a smooth
diffeomorphism from the actual source $M$ to their disjoint union.
This supplies a bounded AT21 reconstruction producer; it does not take
\code{cutComponentGluing} or \code{graphSumRealization} as additional
premises. Older conditional wrappers retain their own meanings, but
their hypotheses must not be reported as an absence of this stronger
concrete result elsewhere in the library.

\paragraph{Exact retained output.}
Let $C$ denote the actual capped manifold. The output includes a finite
index type $W$, an assignment from $\pi_0(C)$ to $W$, lists $L_w$ of
actual capped-component labels and nonnegative integers $k_w$. It
satisfies
\[
 K\in L_w\quad\Longleftrightarrow\quad a(K)=w,
 \qquad L_w\text{ has no repetitions},\qquad L_w\ne\varnothing,
\]
and an actual smooth diffeomorphism
\[
 M\longrightarrow\coprod_{w\in W}
 \mathop{\#}\bigl((C_K)_{K\in L_w},
            \underbrace{S^2\times S^1,\ldots,S^2\times S^1}_{k_w}\bigr).
\]
The notation denotes the library's concrete finite connected-sum
carrier. The result handles disconnected sources. The inherited proof
processes a finite paired-ball gluing: a loop adds one sphere product;
a merge concatenates the two label lists and adds their existing
counts. Its quotient atlases cover both core and seam charts, and
smooth uncapping identifies the resulting quotient with $M$.

The new consumer proves that the literal slots
$\coprod_w\operatorname{Fin}(|L_w|)$ biject with $\pi_0(C)$.
The no-repetition proof is used explicitly. No quotient by manifold
diffeomorphism type is taken. This is stronger bookkeeping than mere
membership coverage and prevents repeated-type factors from being lost.
The exported $k_w$ are actual finite multiplicities. Their identification
with the particular cycle ledger required by AT20 remains a separate
comparison; no canonical count is inferred from an existential theorem.

\paragraph{Actual retained and discarded components.}
The transition already stores a smooth oriented presentation
$C\to Q\sqcup D$, where $D$ is its discarded carrier. Its induced map
on connected components is a bijection, with inverse induced by the
actual inverse presentation. Every capped component is therefore
assigned to a retained or discarded component, and the inherited
restriction maps give an oriented diffeomorphism to that component.
This argument uses an arbitrary point of the capped component; it does
not assume that every discarded component meets a retained core.
The two target tags are disjoint even when the carriers are diffeomorphic.
Composing with the literal-slot bijection gives exact coverage of all
retained and discarded component labels.

If $Q$ is actually empty, every capped component is assigned to $D$.
The source of a surgery transition is nonempty. Its smooth reconstruction
therefore has a nonempty connected-sum group, whose list contains a
capped component; hence $D$ is nonempty. This is a checked single-event
terminal safeguard. It neither proves extinction nor supplies a
geometric classification of $D$.

\paragraph{Application to the same observed events.}
The conversion from a smooth transition with its inherited completion
retains the original presentation and core-inclusion maps. Applying the
reconstruction to every event in an observation of a supplied retained-core
tower uses the library's boundary-frame condition to obtain those
completions. The consumer quantifies over every event of the very same
observation. No event list, new flow, or replacement source carrier is
chosen independently. The general tower and its analytic controls
remain upstream inputs.

\paragraph{The remaining relative interface.}
The smooth finite-sum export does not state that its final reconstruction
preserves orientation or prescribed sphere/cap markings. Oriented
component comparisons do not alone establish those properties of the
whole map. Source-component assignments through the tube graph, exact
port/collar equations and decoration transport through successive
substitutions still need their own binding. Thus this iteration does
not accept full AT21, AT22 or AT24. It narrows their implementation
frontier using actual inherited geometry, without weakening the endpoint.

Three modules in \code{GC_PROOF_PILOTS/revision177} freshly compile;
eighteen transitive axiom queries use only the standard classical
axioms. The three accepted atoms, exact source hashes and eight negative
acceptance controls are recorded in \code{formalization/event_bindings.json}
and \code{reference_checks_revision177.md}. Failed dependent-index,
noncomputability and explicit-argument attempts are retained as API
repair evidence. Full proof closure and Geometrization remain unaccepted.
% END REVISION177 EVENT RECONSTRUCTION

% BEGIN REVISION178 PUBLIC PC ACCEPTANCE
\section{Accepted public PC endpoints and the GC reuse boundary}
\label{sec:public-pc-acceptance-revision178}

The separate public-source checkout is pinned to PC tag \code{v0.1.3},
commit \code{7a48598d35109aa99d1cc678e2724c213cdf4ff3}, with Lean4.33.1
and Mathlib commit \code{0df444a360eaa60ab8c11dca51a86af692955474}.
The public root \code{DifferentialGeometry.Topology.ThreeManifold.Poincare}
built successfully with 22,902 reported jobs. Project artifacts were absent
at the beginning; the pinned official Mathlib dependency cache was used.
This was a fresh public project-target build, not a fresh source rebuild
of Mathlib. All 17,874 tracked source files still match their original
manifest, and tracked Git status is clean.

\paragraph{Exact accepted statements.}
The topological consumer takes an actual compact Hausdorff simply
connected carrier with charts modeled on real Euclidean three-space,
and obtains a homeomorphism to the unit sphere in real Euclidean
four-space. The smooth consumer additionally uses the actual smooth
three-manifold structure and obtains a smooth diffeomorphism to that
sphere. These are exact downstream applications of the public exports,
not replacement predicates named after PC. Seven queried declarations
include those two exports, the two consumers and three surgery exports.
Every transitive axiom set is contained in
\code{propext}, \code{Classical.choice} and \code{Quot.sound}.

The source proof of the topological endpoint installs the inherited
smooth structure and applies the smooth endpoint. The smooth source
proof uses controlled extinction under its explicit simple-connectivity
hypothesis and composes with the actual standard-sphere comparison.
This accepted PC proof does not add extinction to the default general
Geometrization route.

\paragraph{What the surgery queries establish.}
The inspected strong uniform surgery-step export has no
simple-connectivity typeclass. Its expanded statement fixes a finite
horizon, then consumes a compatible history and singular incoming
slab with the stated pinching, derivative, canonical-neighborhood and
noncollapsing controls. Its output includes an actual event, cutoff
records, the boundary-frame condition, standard discarded components
and a volume debit. It is substantial inherited input, but it does not
produce its analytic premises or one coherent global flow merely by
being available as an export.

The inspected small-scale noncollapsing and strong canonical-neighborhood
exports still explicitly assume simple connectivity. Their successful
axiom checks do not remove that hypothesis. Revision175's common-profile
quantifier obstruction is unchanged. General initial data require the
remaining exact producer comparison or extension; no absence claim for
other library declarations is inferred from these particular exports.

\paragraph{Evidence-based current acceptance.}
The current view may now mark the public PC endpoints accepted only
through \code{formalization/public_pc_acceptance.json} and its validator.
The validator checks the original pinned source manifest, successful
root build, exact probe artifacts and seven actual log-derived axiom
sets. The graph exports the corresponding hashes and two scoped release
bindings. General GC flow and full Geometrization remain explicitly
unaccepted. Historical planning rows and revision172's earlier bounded
acceptance record keep their original status; they are not retroactively
rewritten as completed production packets.

The exact command, source meanings, inspected definitions, two corrected
probe-invocation errors and retained evidence are described in
\code{reference_checks_revision178.md}. This completes the public-PC
endpoint reproducibility gate. It does not complete the whole approved
program, the whole-library GC reuse audit or the final proof.
% END REVISION178 PUBLIC PC ACCEPTANCE

% BEGIN REVISION179 GENERAL FLOW INTERFACES
\section[Component volume and a common profile]{Component volume and a common positive surgery profile}
\label{sec:general-flow-interfaces-revision179}

Two precise interfaces in the general-flow comparison now have checked
Lean proofs. They separate inherited metric estimates and elementary parameter
selection from the remaining uniform analytic producers. Neither supplies a
coherent general surgery flow by assuming that flow as a record field.

\paragraph{Normalized component volume is the missing quantitative input.}
Let $W$ be the actual library spatial canonical witness centered at $x$,
with $Q=R_g(x)>0$, constants $C_1,C_2$, and whole-component domain
$U=\operatorname{Comp}(x)$. Suppose one has the quantitative bound
\[
 \operatorname{Vol}_g(U)\geq \frac{\nu}{Q\sqrt Q},\qquad \nu>0.
\]
Put $A=\max\{C_1,1\}$ and $B=\max\{C_2,1\}$. Then for every $r>0$
with $r^4|\operatorname{Rm}|_g^2(x)\leq1$ the checked estimate is
\[
 \operatorname{Vol}_g B_g(x,r)\geq
 \frac{e^{-18A\sqrt B}\nu}{216A^3}\,r^3.
\]
Here the center curvature bound supplements, rather than replaces, the
witness's curvature control throughout $U$. The ball and volume are the
actual ambient Riemannian ball and measure on the compact oriented stage.
No simple-connectivity hypothesis is required by this estimate.

For clarity, the proof uses $L=3A/\sqrt Q$. The inherited scalar contraction
bound gives $r\sqrt Q\leq3$, hence $r\leq L$. Every finite-radius
Riemannian ball centered at $x$ stays in its component; the witness's
radius and inclusion bounds give $U\subset B_g(x,L)$. Consequently the
whole $L$-ball has the witness's curvature bound. With
$q=3\sqrt{BQ}$ this gives the required Ricci lower bound, and completeness
follows from compactness of the actual stage. The inherited local
Bishop--Gromov estimate yields
\[
 \operatorname{Vol}_g B_g(x,r)\geq
 e^{-2qL}\left(\frac r{2L}\right)^3
 \operatorname{Vol}_g B_g(x,L).
\]
Substitution of the component-volume hypothesis proves the displayed
constant. All normalizations, including the squared curvature norm, agree
with the actual library statement.

The existing simply-connected round-model theorem supplies
$\nu=4\sqrt3\,\pi$ and is checked as a downstream application.
A second checked consumer combines the generalized estimate with the
inherited neck, cap and positive-component cases. It obtains a uniform
ball constant from a \emph{uniform} $\nu$ for the round alternatives,
without assuming simple connectivity. This does not produce that uniform
$\nu$ for every actual general-flow component. Finding its appropriate
analytic supplier, or proving a compatible alternative treatment of round
components, remains an explicit obligation. The discard convention is
unchanged.

\paragraph{One profile from locally constrained positive budgets.}
Let $b_n>0$ be an already chosen sequence of interval bounds. Define
\[
 d_0=\min\{1/2,b_0/2\},\qquad
 d_{n+1}=\min\{d_n,b_{n+1}/2\},\qquad
 \delta(t)=d_{\lfloor t\rfloor_++1},
\]
where $\lfloor t\rfloor_+$ is the natural-number floor, equal to zero
for negative $t$. Finite induction gives $d_n>0$, $d_n\leq1/2$,
$d_n\leq b_n/2$, and antitonicity. Thus $\delta$ is positive,
nonincreasing and less than one. For $n\leq\lfloor t\rfloor_++1$,
\[
 \delta(t)<b_n.
\]
In particular, whenever $t\geq n$, both $b_n$ and $b_{n+1}$ are
respected, including interval endpoints. On every bounded time interval
$t\leq T$, the positive number $\delta(T)$ is a lower bound.
Two registers can be combined by replacing $b_n$ with
$\min\{b_n,c_n\}$; the corresponding four interval inequalities are
also checked. No continuity or positive lower bound as $t\to\infty$
is asserted.

This construction takes finite minima at each time. It neither exchanges
$\forall B\exists\delta_B$ with $\exists\delta\forall B$ nor takes a
potentially zero infinite infimum. The interval bounds must be produced
with the required local dependence before the flow is selected.
Kleiner--Lott's Proposition77.2 and Section80 allow adjacent-interval
induction and adjustment of the preceding bound; the present integer
indexing is an elementary selector, not a claim that those analytic
constants have already been constructed on the PC representation.

\paragraph{Actual cutoff-record application.}
The selector is inserted into the actual library \code{CutoffParameters},
retaining supplied auxiliary model and radius data. A record for these
chosen parameters gives, at its actual event time $s$, the recorded neck
accuracy $\delta_\alpha\leq\delta(s)$. Hence $s\geq n$ implies
$\delta_\alpha<b_n,b_{n+1}$. Its nominal surgery radius also retains
the actual inequality $h<\delta(s)^2\rho(s)$. Both applications compile. A deliberate downstream attempt to reuse an old
record after changing its profile is rejected by Lean with a type mismatch.
Records for a different profile are not transported by this construction:
the compatible geometric records and their analytic hypotheses must still
be produced for the chosen profile.

Three modules in \code{GC_PROOF_PILOTS/revision179} and fourteen transitive
axiom queries establish these bounded interfaces and consumers. The general
canonical-neighborhood estimate uses only the permitted standard classical
axioms. Exact source bodies and remaining quantifiers are recorded in
\code{reference_checks_revision179.md}. The active task graph retains the
uniform component-volume input, analytic interval-budget production and
coherent global flow as open requirements; full GC is not accepted.
% END REVISION179 GENERAL FLOW INTERFACES

% BEGIN REVISION180 MODEL VOLUME SUPPLIER
\section[Actual model volume for round components]{Actual ancient-model volume and round-component consumers}
\label{sec:model-volume-revision180}

Revision179 isolated normalized component volume as an input. The new
checked supplier produces it from an actual \code{WindowedModelWitness}
with a fixed positive ancient-model noncollapsing constant $\kappa$.
Its round-model specialization constructs the required model canonical
witness and proves a bound on the actual target component. The remaining
history-wide model producer is distinguished below from this local theorem.

\paragraph{Use the model's quantitative noncollapsing.}
Let $F$ be the witness's actual ancient $\kappa$-solution, normalized by
$R_F(o,0)=1$, and let $K$ be an actual model canonical witness with
constants $C_1,C_2$ and domain $U$. The model definition includes spatial
noncollapsing at all scales. Set $\rho=C_2^{-1}$. The witness's center
scalar bounds imply $C_2\geq1$, and its radius lower bound gives
$B_F(o,\rho)\subset U$. Throughout this ball the squared curvature norm
is at most $C_2^2$, so
\[
 \rho^4|\operatorname{Rm}|^2\leq\rho^4 C_2^2
 =C_2^{-2}\leq1.
\]
Applying the model's spatial noncollapsing and monotonicity of measure gives
\[
 \operatorname{Vol}_{F(0)}(U)\geq\kappa C_2^{-3}.
\]
This applies even when the canonical alternative is round and its ordinary
\code{requiresVolume} field is false. The proof does not invoke
noncollapsing of the target flow, and does not obtain a whole-ball bound
from a center estimate alone.

\paragraph{Transport on the actual embedding.}
Let $W$ be the actual windowed model witness at $(x,t)$ in $S$, with
$Q=R_S(x,t)>0$, accuracy $0<\delta<1$, and the buffer
$2C_1\leq\operatorname{modelRadius}(\delta)$. Its metric comparison,
compact canonical domain and actual embedding give, for every supplied
model-domain bound $\operatorname{Vol}_{F(0)}(U)\geq\mu$,
\[
 \operatorname{Vol}_{S(t)}(W(U))\geq
 \frac{\sqrt{(1-\delta)^3}\,\mu}{Q\sqrt Q}.
\]
The Lean proof retains the inherited volume transport and dimension-three
metric scaling. If $U=\operatorname{Comp}_F(o)$, the inherited compact
open-image argument identifies $W(U)$ with the \emph{actual}
$\operatorname{Comp}_S(x)$. Hence $\delta\leq1/2$ gives the convenient
weaker bound
\[
 \operatorname{Vol}_{S(t)}(\operatorname{Comp}_S(x))\geq
 \frac{\kappa C_2^{-3}/4}{Q\sqrt Q}.
\]
No simple-connectivity hypothesis is used. The component, embedding and
metric in this conclusion are those of the supplied windowed witness.

\paragraph{Produce the round model witness.}
For an actual shrinking spherical-space-form ancient model, the inherited
\code{RoundModelWitness} construction supplies a canonical witness with
whole-model domain and constants
\[
 C_1^{\mathrm{mod}}=2\pi/\sqrt{1/6}+1,\qquad
 C_2^{\mathrm{mod}}=2.
\]
The model is connected, so its whole domain is its basepoint component.
With $2C_1^{\mathrm{mod}}\leq\operatorname{modelRadius}(\delta)$ and
$\delta\leq1/2$, the preceding theorem now gives
\[
 \operatorname{Vol}_{S(t)}(\operatorname{Comp}_S(x))\geq
 \frac{\kappa/32}{Q\sqrt Q}.
\]
The actual model record includes a diffeomorphism to a round quotient and
its shrinking metrics for all ancient times; a geometry name is not used
as a substitute. The model canonical witness and volume bound are
outputs here, rather than additional premises.

For an actual target spatial canonical witness on this whole component,
put $A=\max\{C_1^{\mathrm{src}},1\}$ and
$B=\max\{C_2^{\mathrm{src}},1\}$. Revision179 then supplies
\[
 \operatorname{Vol}_{S(t)} B(x,r)\geq
 \frac{\kappa e^{-18A\sqrt B}}{6912A^3}\,r^3
 \quad\text{if }r>0,\quad r^4|\operatorname{Rm}|^2(x)\leq1.
\]
The checked consumer quantifies one positive ball constant before the
actual stage, time, model and target witnesses. Its dependence on the
fixed $\kappa$ is essential; arbitrary spherical quotients have not been
given a uniform volume bound independent of their geometry.

\paragraph{What the history induction must still retain.}
The inspected \code{SpatiallyCanonicalBefore} interface produces only a
spatial canonical witness and its cap-neck condition. The inspected
\code{CanonicalWitness.toSpatial} projection retains the conditional
volume field but no underlying ancient model or its $\kappa$. Thus the
local theorem cannot be applied to this interface by inventing a forgotten
windowed witness. A general-flow induction must produce and retain suitable
model data, or the resulting normalized-volume reserve, uniformly on the
relevant interval. Its quantifier order must derive the model's $\kappa$
from already available induction data, without assuming the current-flow
noncollapsing estimate it seeks to prove. This is an explicit remaining
producer obligation, not a claim of absence from the entire PC library.

Four modules in \code{GC_PROOF_PILOTS/revision180} compile, and seven
transitive axiom queries use only the standard classical axioms. Exact
source-body locators and the inspected interface boundary are recorded in
\code{reference_checks_revision180.md}. One bounded model-volume supplier
is added to the accepted graph. General coherent flow, analytic interval
budgets, marked history reconstruction and full GC remain unaccepted.
% END REVISION180 MODEL VOLUME SUPPLIER

% BEGIN REVISION181 RESERVED CANONICAL INDUCTION
\section[Retaining quantitative canonical data]{Retaining quantitative canonical data through continuation}
\label{sec:reserved-canonical-revision181}

The actual round-model input now produces a complete target packet, not
just a volume estimate conditional on a separately supplied target witness.
A second checked argument retains this quantitative data through the
strict-before-time continuation proof. These are bounded producers and
adapters; the analytic continuation premises and history-wide uniformity
remain separate obligations.

\paragraph{A record that does not forget the reserve.}
The new \code{ReservedCanonicalWitness} contains the actual library
spatial canonical witness $V$, its cap-neck charts, and the implication
\[
 \neg V.\operatorname{requiresVolume}\quad\Longrightarrow\quad
 \operatorname{Vol}_g(\operatorname{Comp}(x))\geq
 \frac{\nu}{R_g(x)\sqrt{R_g(x)}}.
\]
The negation selects the round alternative of the stored witness. It
neither demands extra volume for necks and caps nor substitutes another
canonical witness on the same point. The checked nonround constructor
uses the existing volume alternative. Enlarging the canonical constants
preserves the actual domain, cap charts and reserve; reducing $\nu$
preserves the same witness. Case analysis then supplies a uniform positive
ambient-ball constant from $\nu>0$, with the same center curvature
condition as revision179. In the round case the stored domain equals the
actual connected component, so the reserve applies to that domain.

\paragraph{An actual round-packet producer.}
For each $0<\epsilon<1/11$, the checked producer chooses $C\geq1$ and
$0<\delta_0<1$ before the stage, flow, time, model and $\kappa$.
An actual round \code{WindowedModelWitness} with
$\delta\leq\delta_0$, the Ricci-flow equation and its prescribed regular
backward window produces
\[
 \operatorname{ReservedCanonicalWitness}
       (g(t),\epsilon,C,C,\kappa/32,x).
\]
The inherited round buffered-canonical construction supplies the actual
canonical neighborhood and cap-neck charts. Its tolerance is weakened to
$\epsilon$ and its spatial projection is retained. Revision180 supplies
the component reserve. The accuracy threshold is the finite minimum of
the inherited threshold, $1/2$, and $(2A)^{-2}$, where
$A=2\pi/\sqrt{1/6}+1$; this last term supplies the model-radius buffer.
An additional checked consumer produces a ball-volume estimate directly
from these round-window hypotheses, with no target canonical witness
among the inputs. No simple-connectivity assumption occurs.

\paragraph{Strict-before-time continuation retains quantitative data.}
For an actual incoming slab $G$ on $[a,s)$, let
$\operatorname{ReservedBefore}(t_0)$ mean that every point with
$a<t<t_0$ and $R_G(y,t)>q_s$ has the preceding record with one fixed
$\nu$. Define the corresponding short forward condition using
$t_0\leq t<t_0+\eta$ and $t<s$. Forgetting the reserve gives the
inherited \code{SpatiallyCanonicalBefore} condition.

The checked extension of the existing supremum argument takes the actual
canonical-neighborhood, scalar-derivative, scalar-gradient and reserved
conditions together. Its premises supply the noncollapsing proposition
$N(t_0)$ from those before-time conditions, $N(a)$, and a positive local
continuation interval from the same data and $N(t_0)$. It proves all four
conditions before $s$. At the supremum $T$ of good times, each tested
$t<T$ already belongs to an earlier good interval, so its actual witness
and quantitative reserve are retained. If $T<s$, local continuation
contradicts maximality. No limiting witness at $t=T$ is asserted.
A downstream Lean consumer applies this conclusion to obtain the uniform
ambient-ball estimate throughout the high-curvature part of the slab.
This proof does not construct either analytic premise from a field that
already assumes its conclusion.

\paragraph{A rejected constant-selection shortcut.}
For the particular explicit bound exhibited in revision180, put
\[
 k(\kappa,A,B)=
 \frac{\kappa e^{-18A\sqrt B}}{6912A^3}.
\]
Lean verifies that $\kappa>0$ and $A\geq1$ imply
$k(\kappa,A,B)<\kappa$. Thus this bound alone cannot close an induction
requiring $\kappa\leq k(\kappa,A,B)$ with the same positive constant on
both sides. This rejects that specific substitution, not the possibility
of another proof of noncollapsing. A distinct earlier model constant or
another independently produced reserve must be justified where this
supplier is used.

Kleiner--Lott's Lemma79.12 and Remark79.13 require the next interval's
noncollapsing constant to be independent of the new canonical radius
$r_+$; the surgery-accuracy bound may depend on $r_+$. Section80 uses
this order in its contradiction sequence. The implementation must retain
that order. Merely adding a reserve field does not prove the required
history-wide analytic supplier. A further source-comparison question is
kept explicit: Definition69.1 omits the ordinary volume bound for its
round alternative, while Sublemma79.23 invokes canonical volume control.
The selected application needs a justified treatment of that alternative;
this targeted observation is not a counterexample to the source theorem.

Four modules in \code{GC_PROOF_PILOTS/revision181} compile; eleven
transitive axiom queries are standard. Two bounded proof bindings and one
checked numerical obstruction are added. Sources, exact reading limits,
constant order and the remaining comparison are recorded in
\code{reference_checks_revision181.md}. General flow, global analytic
reserve production and final GC remain unaccepted.
% END REVISION181 RESERVED CANONICAL INDUCTION

% BEGIN REVISION182 ACTUAL DECLARATION GRAPH
\section{Actual declaration dependencies and composition audit}
\label{sec:actual-declaration-graph-revision182}

The planning graph of Section~\ref{sec:active-graph-revision173} and the
actual declaration graph serve different purposes. A planning edge records a
mathematical requirement; a declaration edge records a constant occurring in a
checked type or proof value. Neither an empty prerequisite list nor a successful
axiom query establishes that all premises of a theorem have been produced.
Revision182 adds the second graph while preserving the first and its open leaves.

\paragraph{Exact extraction scope.}
The exporter loads the immutable public PC release and the already checked
positive proof modules from revisions172--181. Its 127 roots are exactly the
union of their registered declaration/consumer and axiom-query names. It follows
references recursively inside the selected pilot modules, including private
helpers, and retains external references as an explicit boundary. The resulting
288 declarations have 1,440 distinct external boundary references. The export
is not the full transitive body graph of PC or Mathlib.

For every captured declaration, the stored record contains its origin module,
kind, universe parameters, closed type expression, readable type, direct type
constants, direct proof-value constants, projection-structure references,
inductive-constructor links and Lean's transitive axiom list. Types retain all
implicit and instance binders and use bound-variable indices; shared expressions
are stored once in an indexed expression graph. Only nonsemantic expression
metadata is omitted. Constructor types expose the fields of the captured
records instead of presenting only the uninformative outer declaration
\code{Type}. This is inspected tooling output from a checked environment, not
an additional mathematical theorem about Lean's kernel.

The exporter explicitly requests theorem and opaque values. A plain lookup that
omits them would hide actual proof dependencies. Module identity, rather than a
public-name prefix, determines the local traversal; the present closure includes
one private declaration. Separate projection references avoid treating a
structure projection as though it had no structure dependency. The exporter
and its driver freshly elaborate; all 127 root axiom lists agree with the earlier
proof-build evidence, and all 288 captured lists use only the three accepted
standard axioms or a subset. Earlier mathematical proofs are not rebuilt by
this extraction.

\paragraph{Cycles and the meaning of an edge.}
The actual value-reference subgraph of these captured declarations is acyclic.
The combined type/constructor graph has four nontrivial strongly connected
components. Each pairs one of the following records with its own constructor:
\[
\begin{gathered}
\texttt{PointedBallApprox},\quad \texttt{PairedSides},\\
\texttt{CutCapSumData},\quad \texttt{ReservedCanonicalWitness}.
\end{gathered}
\]
These are legitimate specification references. They are retained;
deleting them to force a DAG would discard field information. The condensation
graph can be used for display. No claim of mathematical circularity follows
from these four components.

Even a proof-value reference is a syntactic observation: a constant can occur
in an annotation or a definition unfolded by a later proof. A path is evidence
of use, not a complete review of which output satisfies which premise on which
carrier. Semantic premise matching remains a separate acceptance obligation.
Absence of a path likewise does not refute either individual theorem.

\paragraph{Six corrected composition interpretations.}
Of the twenty-eight registered producer/consumer pairs, twenty-two have an
actual syntactic value path. Source inspection explains the six others:
\begin{enumerate}[label=\arabic*.]
\item \code{MC07.restriction}: the registered composition constructor performs
its own domain control and does not call the restriction constructor. The
separate \code{realRestriction} definition is a genuine existing application,
but it was not the registered consumer or an original axiom-query root.
\item \code{MC10.euclidean_consumer}: the unit-ball equality and buffered-ball
inclusion are parallel applications of scaling. They provide two checked
specializations, not the recorded serial dependency.
\item \code{GA03.paired_images}: disjointness of the example seam images uses
their actual cylinder locations. It does not consume the two-element seam-fiber
theorem; image descent and the fiber description are distinct facts.
\item \code{LP01.nonempty_or_empty}: the no-later-event theorem takes an empty
slice and uses the tower's event-calendar field. It does not eliminate the
previously produced late-or-empty disjunction.
\item \code{AT21.actual_component_slots}: the target-slot bijection composes
actual component equivalences. It does not consume the separate oriented
retained-or-discarded classification theorem.
\item \code{LP01.common_profile_selector}: the actual neck-budget theorem uses
the explicit single-register profile and its interval inequalities. It does not
consume the existential two-register profile selector. Both positive registers
and actual cutoff records for the selected profile remain required inputs.
\end{enumerate}

The existing declarations keep their individual kernel and axiom evidence.
The active graph now carries these composition findings explicitly, rather than
upgrading a shared topic or related example into a checked theorem application.
The next repair must bind a real existing application where appropriate or
construct the intended consumer with its actual inputs. Adding a theorem that
merely restates the desired conclusion as a premise would not repair the chain.

The actual round-model/volume, enriched canonical-packet/ball and
reserve-continuation/ball chains do have the registered value paths. Their
analytic interval/history suppliers remain open as stated in
Section~\ref{sec:reserved-canonical-revision181}. No mathematical parent,
selected-route leaf, general-flow producer or final GC endpoint is accepted by
this audit. Exact source locations, extraction limits, reproducible commands and
negative evidence controls are retained in
\code{reference_checks_revision182.md} and
\code{formalization/lean_declaration_evidence.json}.
% END REVISION182 ACTUAL DECLARATION GRAPH

% BEGIN REVISION183 ACTUAL COMPOSITION REPAIRS
\section[Actual marked consumers]{Actual marked consumers and parameter-register applications}
\label{sec:actual-composition-revision183}

Revision182's six findings concern the accuracy of registered composition
claims, not counterexamples to the individual checked theorems. Five now have
actual applied consumers, and the remaining pair is explicitly classified as
parallel specializations. The historical proofs and audit remain unchanged;
current binding overlays identify the corrected consumers and their evidence.
No original independent proof is assigned a dependency that it does not have.

\paragraph{Restriction and paired sides.}
The existing \code{realRestriction} applies the general restriction constructor
to the actual real-line identity approximation. The new named consumer uses
its coverage at the closed target endpoint: source radius $3$, error $2$, and
target point $1=3-2$. It produces a source point in that same restricted closed
ball with image distance strictly below $2$. This replaces the inaccurate
registration of the separate composition constructor as a use of restriction.

The paired-side consumer applies the general seam-fiber theorem to the actual
four labels $\operatorname{Fin}(2)\times\{\mathrm{false},\mathrm{true}\}$.
For each $i$, the fiber over the seam of $(i,\mathrm{false})$ consists exactly
of $(i,\mathrm{false})$ and $(i,\mathrm{true})$. The output retains these
multiplicities together with disjointness of the actual two torus images in
the cylinder. This is an actual topology example with actual maps; it does
not construct the boundary faces or geometric data of a general graph region.

The two Euclidean scaling statements require a different correction. The
unit-ball equality applies the general normalized closed-ball theorem, and
the buffer inclusion applies the general normalized buffer-inclusion theorem.
Both actual applications are recorded. There is no mathematical requirement
for either specialization to use the other, and no artificial serial theorem
is inserted merely to manufacture such an edge.

\paragraph{One reconstruction with classified literal slots.}
For a supplied actual spherical cut/cap transition $E$, the new consumer first
produces one \code{CutCapSumData} record $R$ using the inherited actual
reconstruction theorem. It then proves simultaneously that the literal slots
\[
  \coprod_{w\in R.\mathrm{Group}}\operatorname{Fin}(|R.\mathrm{factors}(w)|)
\]
map bijectively to the actual retained-or-discarded target components, and
that every such slot has the corresponding oriented smooth component
comparison. Both assertions use that same $R$, its actual capped-component
label, and the original event presentation. The input does not assume the
classification or an arbitrary preassembled reconstruction package.

The output preserves the full scope of $R$'s actual smooth reconstruction
map. The separate oriented component comparisons do not automatically orient
that global map, fix prescribed ports, or provide the full finite-history
cycle and reconstruction data. Those remain distinct obligations. The new
application closes the registered local composition, not the whole AT21 parent.

\paragraph{The same initial identification in both outcomes.}
The late-or-empty consumer now actually eliminates the supplied observation
tower's established dichotomy. On the empty side it retains a horizon $a$, all
later empty carriers, the event-time bound, and a data value
\[
 A:\mathrm{InitialIdentification}(P,g,T.\mathrm{observe}(a)),\qquad
 A=T.\mathrm{observeInitial}(a).
\]
The displayed abbreviation suppresses the nonnegative-time proof arguments,
which are retained in Lean. On the late side, for every lower time bound it
produces a later regular nonempty prefix, the strict final-slab time inequality,
and the corresponding identification equal to the same tower's stored value.

The equality to the actual stored identification is intentional. Mere
nonemptiness of the type would retain an isometric oriented comparison but
would not by itself identify a later chosen witness with the tower's marked
map. Here the exact comparison is part of the result in both branches.
This consumes one already supplied coherent tower; it constructs neither
general surgery flow nor exceptional geometry for an empty continuation.

\paragraph{One selected profile, followed by actual cutoff records.}
Given positive registers $b_n,c_n$ and auxiliary cutoff parameters $p$, the
new consumer applies the existential two-register selector and produces one
actual parameter record $q$ before any history or geometric cutoff record is
quantified. Its neck-radius and protected-radius functions, fixed cap scaffold,
model radius, order, accuracy and recentering constant equal those of $p$.
Its accuracy profile is positive, below $1$ on nonnegative times, and antitone.
For every $n$ and $t\geq n$, the output retains the direct inequalities
\[
 q.\delta(t)<b_n,\quad q.\delta(t)<c_n,\quad
 q.\delta(t)<b_{n+1},\quad q.\delta(t)<c_{n+1}.
\]

Only after choosing $q$ does the theorem quantify an actual history $H$, event
$i$, and \code{GeometricCutoffRecord} for those exact parameters. The record's
own inequality between neck accuracy and $q.\delta$ yields the same four
bounds for each actual neck at that event time. No record built for an old
accuracy profile is silently reused. Keeping the direct profile bounds matters:
record-only conclusions would not supply a profile estimate on an interval
with no surgery events. No decay at infinity is inferred from arbitrary
positive registers; a suitable decaying register is still to be specified when
that conclusion is needed. Analytic register production and a coherent flow
carrying compatible records remain input obligations.

\paragraph{Checked scope and current graph.}
The three proof modules in \code{GC_PROOF_PILOTS/revision183} freshly compile
with seven standard-axiom queries, including the existing restriction
application. The declaration-export driver also freshly compiles. The updated
bounded export has 134 roots and 302 declarations; all root axiom lists agree
with the prior checked roots and these new queries. Of the twenty-eight
registered bindings, twenty-seven now have their actual value paths and one
is explicitly a pair of parallel specializations with its two upstream
applications. The six revision182 registration findings are resolved in this
bounded sense. Global semantic premise closure is not claimed.

The planning graph keeps every previously open selected-route leaf. Sources,
exact carriers, retained maps, quantifier order and reproducible evidence are
recorded in \code{reference_checks_revision183.md} and
\code{formalization/composition_resolution_evidence.json}. The next work returns
to the exact endpoint and the remaining general analytic and marked-history
producers; the full approved program is still active.
% END REVISION183 ACTUAL COMPOSITION REPAIRS

% BEGIN REVISION184 FIXED MODEL ATLASES
\section[Fixed models and geometric atlases]{Fixed models and actual complete geometric atlases}
\label{sec:fixed-model-atlases-revision184}

The model table in Section~\ref{sec:m1d-model-geometry} determines fixed standard
metrics. Its implementation must preserve that requirement: an arbitrary
complete homogeneous metric cannot be assigned a standard-model tag merely by
putting it in a generic record. Revision184 binds three elementary cases to
actual metrics in the accepted PC release. The bounded
\code{ElementaryModel} registry has spherical, Euclidean and spherical-product
constructors; \code{ElementaryGeometry} matches each constructor to its fixed
metric by definition. This registry is not a replacement for the eight-tag
endpoint. The remaining five model bindings, model homogeneity proofs and
full cut-piece certificate are still required.

\paragraph{Actual charts and whole-carrier transport.}
For actual smooth metrics $g$ on $M$ and $h$ on the fixed model $X$,
\code{ModelAtlas g h} says that for every $x\in M$ there is a smooth partial
diffeomorphism $e:U\subset X\to V\subset M$ with $x\in V$ and
\[
 g_{e(y)}(D e_y v,D e_y w)=h_y(v,w)\qquad(y\in U).
\]
Both domains are open and the inverse is smooth. The charts are stored in the
inverse direction from the informal interior-to-model convention; each
partial diffeomorphism already contains its inverse. No orientation sign is
imposed on these metric charts. The model $h$ is fixed outside every point and
chart choice. \code{CompleteModelAtlas} adds completeness of this same $g$
using the inherited metric-dependent Riemannian distance.

If $f:M'\to M$ is an actual whole-carrier diffeomorphism, the new atlas uses
$f^{-1}\circ e$, and the chain rule proves the metric identity for $f^*g$.
The inherited distance-pullback theorem gives completeness for that same
metric. Hausdorff and sigma-compact hypotheses are explicit in the Lean
completeness interface. This does not replace an entire interior by a
truncated core, nor attach boundary labels or orient a reconstruction map.

\paragraph{Three fixed metrics and spherical quotients.}
The spherical metric is the induced round metric on the unit sphere in
$\R^4$; the Euclidean metric is the library's standard inner product on
$\R^3$. The spherical-product metric is the actual
\code{unitCylinderMetric}. A checked equality identifies it with the product
of the unit round $\Sph^2$ metric and the Euclidean line metric. Compactness
and the inherited Euclidean/product completeness proofs supply completeness
of these three fixed metrics. Identity charts give their own model atlases;
this is not asserted to prove model homogeneity.

For an actual \code{SphericalSpaceFormQuotientModel}, revision176 already
constructed a complete round metric on the supplied carrier. The new proof
uses the inherited quotient projection's actual isometric partial charts,
then the stored whole-carrier diffeomorphism. It thereby supplies an atlas
for exactly revision176's metric with target the fixed round $\Sph^3$.
The supplied finite action/model comparison is still an input; this theorem
neither produces it from a discarded component nor proves primeness.

\paragraph{The actual sphere-times-circle factor.}
The construction begins on
$\Sph^2\times(\R/\mathbb Z)$, with the unit round sphere metric and the
inherited flat circle metric whose parameter tangent has length one.
The quotient map from $\R$ is a local diffeomorphism and has the exact
pullback metric $dt^2$. Its actual local charts, multiplied by the sphere
identity, give charts from the fixed $\Sph^2\times\R$ model. Compactness gives
completeness on this carrier.

The inherited additive-circle-to-unit-sphere diffeomorphism then transports
this metric to the precise \code{SphereTwoTimesCircle} type used by PC
reconstruction, namely the product of the unit spheres in $\R^3$ and $\R^2$.
The transported circle has period-length one; it is not claimed to equal the
radius-one induced circle metric, whose circumference is $2\pi$. The local
model remains the Euclidean line, so this normalization is legitimate.
Any supplied factor diffeomorphism now yields the metric and complete
fixed-model atlas on its actual source carrier.

\paragraph{Standard-factor and presentation consumers.}
The theorem \code{standard_factor_complete_geometry} eliminates the actual
library predicate \code{isStandardFactor}. Its spherical branch consumes
the oriented finite-quotient comparison; its other branch consumes the actual
sphere-times-circle factor map. It returns a metric on that factor and its
corresponding fixed elementary tag.

The downstream presentation theorem first applies revision176's orientation
correction, retaining an actual oriented presentation $Q$ whose list is the
original list or its pointwise opposite. For each literal slot
$i\in\operatorname{Fin}(|Q.\mathrm{factors}|)$, it then supplies the metric
and tag on $(Q.\mathrm{factors}[i]).\mathrm{Carrier}$. Repeated factor types
remain separate slots and $Q$ retains its actual oriented connected-sum map.
If the list is empty, this per-slot assertion is vacuous; the sphere unit
convention still needs its explicit treatment in final prime reconstruction.
The theorem does not assert that a connected sum itself has a homogeneous
metric, that its listed factors are certified prime, or that its global map
agrees with prescribed surgery ports.

\paragraph{Sources, checks and remaining scope.}
Scott's definition on printed403/PDF3 and product discussion on459/PDF59
were reopened. The proof uses the pinned PC metric, circle, product,
round-quotient and presentation APIs, with exact inspected locations in
\code{reference_checks_revision184.md}. Author errata URLs remained
unavailable; no new errata clearance or eight-geometry classification is
claimed. Four proof modules, thirteen standard-axiom queries, one newly
compiled unchanged upstream module, and an actual declaration export are
checked. Three registered producer/consumer paths occur in the compiled
values. Static evidence controls and document checks remain separate from
these Lean proofs. No general-flow, marked-history, full M1d/AT25 or final
GC parent is accepted by this iteration.
% END REVISION184 FIXED MODEL ATLASES

% BEGIN REVISION185 DISCARDED GEOMETRY
\section[Geometry on discarded components]{Actual discarded components and nonempty geometric presentations}
\label{sec:discarded-geometry-revision185}

Revision184 supplies complete fixed-model metrics on each literal factor of an
actual oriented standard presentation. The next application must use the
surgery event's actual discarded component and must account for the convention
that the empty connected-sum list represents $\Sph^3$. These are checked here
without upgrading a connected sum to a geometric manifold or changing the
endpoint's prime and boundary requirements.

\paragraph{The sphere unit and the original map.}
The inherited \code{finiteConnectedSum} sends the empty list to
\code{standardThreeSphereLift}; a singleton list is its actual factor.
The inherited theorem \code{isStandardFactor_standardThreeSphereLift}
recognizes this sphere through the trivial finite positive free action and an
actual oriented quotient comparison. Define
\[
 N(L)=\begin{cases}[\Sph^3],&L=[],\\L,&L\ne[].\end{cases}
\]
The new normalization sends an actual presentation $P$ to a presentation with
list $N(P.\mathrm{factors})$. Lean checks that the list is nonempty and that its
smooth diffeomorphism is heterogeneously equal to the original one. In the empty
case both finite sums reduce to the same sphere carrier. Nonempty lists retain
their order and repeated entries. This does not erase a sphere vertex with
prescribed ports and does not identify an empty graph carrier with a sphere.

A \code{StandardGeometricPresentation} now consists of an actual oriented
standard presentation, a nonempty factor list, and a fixed elementary model tag
and complete geometric metric for every literal factor index. Its producer
applies revision184 to the normalized presentation. The final list is $N(L)$
or its pointwise opposite, retaining revision176's explicit orientation
correction. The normalization preserves the map before this correction; the
correction may compose the inherited opposite-sum comparison. A concrete
empty-list sphere example and nonemptiness of the resulting factor index are
also kernel checked.

\paragraph{Actual event and observation consumers.}
The library predicate \code{MetricCutCapEvent.poincareStandardDiscarded}
asserts a standard presentation for each connected component of that event's
actual discarded carrier. It does not ask for a new unlabelled carrier or for
all possible cut/cap completions. The new consumer eliminates this predicate
at the supplied component and constructs its nonempty geometric presentation.
The analogous spherical cut/cap consumer uses \code{poincareControlled}.
Neither predicate introduces a simply-connected premise.

For an actual observed history $H$, retain the dependent indices
\[
\begin{gathered}
i\in\operatorname{Fin}(H.\mathrm{eventCount}),\qquad
C\in\pi_0((H.\mathrm{event}(i)).\mathrm{discarded}),\\
j\in\operatorname{Fin}(|G_{i,C}.\mathrm{factors}|).
\end{gathered}
\]
The history consumer supplies the presentation at each such $i,C$; its factor
index is nonempty. A zero-event history has no event slots and is legitimate.
The tower consumer applies the release's actual restriction theorem at each
observation time, preserving its dependent event identification. It assumes
the tower and its discarded-component classification. It does not produce the
general tower, its analytic estimates or an extinction theorem.

\paragraph{One reconstruction, actual discarded labels.}
The new marked consumer first obtains revision183's one actual
\code{CutCapSumData} $R$ and its bijection from literal reconstruction slots to
retained-or-discarded target components. Each retained branch keeps its actual
oriented component comparison. Each discarded branch stores a point $d$ in
the actual discarded carrier, the exact target label, an oriented map
\[
 f:C_{\mathrm{cap}}\longrightarrow C_{\mathrm{discarded}},
\]
and the produced geometric presentation $G$ of that same discarded component.
The constructed presentation at the capped component uses
\[
 C_{\mathrm{cap}}\xrightarrow{f}C_{\mathrm{discarded}}
 \xrightarrow{G.\mathrm{diffeomorph}}
 \mathop{\#}_{j}G.\mathrm{factors}[j].
\]
The checked \code{atCap_data} identity proves this actual composition and retains
$G$'s factor list, tag function and metric function literally. Thus factor
geometry is attached to the same component slot classified by $R$; repeated
geometric types do not merge their labels.

The global map $R.\mathrm{reconstruct}$ retains its previously accepted smooth
scope. This construction does not make it oriented or compatible with arbitrary
prescribed ports. It does not flatten the nested connected sums, certify their
factors prime, or realize additional sphere-product cycle factors in $R$.
Those are subsequent reconstruction obligations.

\paragraph{The empty-retained branch is nonvacuous.}
Suppose the actual retained carrier is empty. Source nonemptiness and
$R.\mathrm{reconstruct}$ give an actual reconstruction group; its stored
nonempty capped-component list gives an actual slot. The retained alternative
is impossible. The discarded alternative therefore gives a geometric
presentation with a nonempty factor index, together with the actual composite
map above. This argument concerns a supplied event with empty retained carrier;
it makes no assertion that every flow reaches such an event.

\paragraph{Evidence and remaining producers.}
The source comparison checks the pinned PC presentation, finite-sum convention,
trivial spherical action, actual event predicate and restriction proofs, and
reuses the unchanged fixed-metric and component-reconstruction checks from
revisions176--184. Exact locators are in
\code{reference_checks_revision185.md}. Three proof modules and twelve axiom
queries are checked, followed by actual declaration-path extraction. No new
printed classification or errata clearance is claimed.

The actual analytic classification supplier, full general-flow construction,
prime/port certification and finite marked-history substitution remain open.
This bounded implementation supplies geometry once the library's actual
standard-component classification is present; it is not acceptance of AT25,
AT28, the complete eight-model endpoint or Geometrization as a whole.
% END REVISION185 DISCARDED GEOMETRY

% BEGIN REVISION186 FINITE GEOMETRIC HORIZONS
\section[Finite geometric horizons]{Actual metric events and finite geometric horizons}
\label{sec:finite-geometric-horizons-revision186}

The discarded geometry of revision185 is now attached to the actual analytic
surgery event. A further checked construction supplies a finite history at
each prescribed positive time, conditional on the general strong canonical-
neighborhood package. This isolates that analytic premise and the subsequent
coherent-tower problem while reusing the accepted PC continuation proofs.
No simply-connected assumption or extinction conclusion is inserted.

\paragraph{From the actual metric event.}
Given a \code{MetricCutCapEvent} with boundary-frame reversal and its actual
discarded-component classification, use the inherited smooth capping completion.
The converted spherical cut/cap event retains the original presentation and
core inclusion. \code{ActualMetricEventGeometry} records these identities,
revision185's actual labelled reconstruction, and the complete fixed-model
factor geometry of every actual discarded component. Its global reconstruction
map retains the earlier smooth scope; prime recognition and prescribed-port
compatibility are still separate.

\paragraph{The source already resolves the parameter order.}
The uniform-debit step chooses its surgery accuracy $\varepsilon$ after receiving
the time-derivative constant $C_{\mathrm{time}}$, whereas a canonical-neighborhood
package can return that constant after receiving an accuracy. The inherited
integration proof uses two calls, rather than interchanging these quantifiers.
First apply the canonical theorem at
$\min(1/22,\overline\varepsilon)$ to obtain a time-derivative bound. Feed this
constant into the surgery selector to choose $\varepsilon$. Apply the canonical
theorem again at that $\varepsilon$ to obtain gradient, canonical, spatial,
pinching and noncollapsing bounds.

Take the maximum of the two scalar thresholds, the minimum of the permitted
small parameters, and the maximum of the model radii and orders. The first
call supplies the time-derivative estimates; the second supplies the other
analytic estimates. The inherited threshold-monotonicity lemmas and these
common restrictions satisfy the actual event producer's premises. The proof
retains the original order of every choice.

The new \code{GeometricHorizonProduction} is the source's finite-horizon step
contract enriched by the actual geometric output above. For each $B>0$ it
chooses base parameter data, a positive debit $v$, and a history invariant
before receiving the history or incoming slab. It preserves actual incoming-
slab equality, the appended cutoff record, all five fixed parameter equalities,
boundary reversal, component classification and the original compact terminal-
volume inequality. The public PC \code{uniformDebitSurgeryStepStrong} supplies
its step input without a simply-connected hypothesis. The general strong
canonical-neighborhood theorem is still supplied as an explicit premise;
this construction does not prove it.

\paragraph{An actual finite history, not a selected output package.}
Begin with the actual zero-event history. Form the set of histories extending
that history, ending at an event time no greater than $B$, satisfying the
produced invariant and carrying common cutoff records, boundary reversal,
component classification and the same debit $v$. This set contains the zero
history. Appending the produced event preserves it.

At an append, use the library's actual splice of parameter functions at the
new event time. Earlier event times are strictly smaller, so their records
remain valid. The new record is transported at exactly its own event time;
its fixed data agree with the base data. This is a constructed common parameter
function for that finite history, not a change to old records under arbitrary
new parameters.

The initial scalar barrier has the normalization
\[
 R(g)\ge -\frac{3}{2c},\qquad c>0.
\]
The inherited compact-minimum proof assumes a nonempty carrier. The new adapter
handles the empty carrier with $c=1$ and a vacuous pointwise inequality; for a
nonempty carrier it reuses that proof. Actual initial metric transport and the
inherited scalar evolution give a uniform lower bound throughout every history
in the set.

The inherited volume and component estimates then give a finite upper bound
on event count. This includes a nontrivial event with no cuts, which must discard
a component. Choose a history with maximal event count. If it has not reached
$B$, the actual maximal-flow alternative either extends by a closed slab or
produces a singular incoming slab. The latter would allow one more surgery
event in the set, contradicting maximality. The closed-slab alternative therefore
reaches $B$ without adding events. This argument uses the general finite-horizon
continuation theorem; the source's subsequent extinction argument is omitted.

\paragraph{Exact output and its quantifier boundary.}
For every $B>0$, the checked theorem constructs a retained-core history $K$,
an initial identification $A$, and cutoff parameters $p$ with
\[
\begin{gathered}
 K.\mathrm{horizon}=B,\qquad A_0\preceq A,\\
 \operatorname{Nonempty}\!\left(
   \prod_{i\in\operatorname{Fin}(K.\mathrm{eventCount})}
   \mathrm{CutoffRecord}(K,i,p)\right).
\end{gathered}
\]
Here $A_0$ is the actual zero-history identification and $\preceq$ retains its
initial map through the prefix. Every actual event has boundary reversal,
standard discarded components and \code{ActualMetricEventGeometry}. In the
closed-slab case the inherited horizon-extension theorem transports the exact
event records; it introduces no new event labels. Zero-event histories and
empty stages are allowed.

This is a finite-horizon construction from the supplied general analytic
package. The assertion $\forall B>0\,\exists K_B$ does not make the chosen
$K_B$ compatible under restriction. It also does not supply one globally
controlled monotone profile. These remain the coherent-tower and common-profile
obligations. Neither a global flow nor late thick/thin geometry follows by
exchanging these quantifiers.

\paragraph{Source and validation scope.}
The checked source passages are the two-call canonical integration proof, the
actual uniform-debit factory proof, the finite-history splice, the general
maximal-event-count continuation argument, and the actual component/volume
bounds. Exact hashes and locators are in \code{reference_checks_revision186.md}.
Three proof modules and six standard-axiom queries implement this bounded
construction and its actual downstream finite-horizon consumer. The root types
retain the general strong canonical-neighborhood premise. No new printed-source
classification or errata clearance is claimed. The general analytic producer,
coherent tower/profile, late analysis, prime/port/history reconstruction and
full Geometrization remain open.
% END REVISION186 FINITE GEOMETRIC HORIZONS

% BEGIN REVISION187 GENERAL CANONICAL DEPENDENCY
\section[General canonical dependency]{General canonical neighborhoods from the exact small-scale input}
\label{sec:general-canonical-small-scale-revision187}

The general strong canonical-neighborhood premise of revision186 can now be
reduced by actual applications of accepted PC suppliers. No analytic estimate
is re-proved in this iteration. The result identifies the precise remaining
input on the original oriented stage $P$ and metric $g$; it retains the exact
history, metric, time and cutoff conventions of the release.

\paragraph{Six general suppliers, with actual applications.}
The public reduced-volume monotonicity, local upper bound and initial lower
bound supply the three corresponding premises of the general noncollapsing
assembly. Its fourth premise remains
\code{SmallScaleNoncollapsingThroughSurgery P g}. The assembly returns
\code{NoncollapsingThroughSurgery P g}. General pinching, temporal canonical
continuation and spatial canonical continuation then supply the other three
premises of the strong canonical-neighborhood assembly. In particular, the
new proof does not apply the public simply-connected specialization of that
assembly to a general stage.

The accepted chain is
\[
 \begin{gathered}
 \text{small-scale noncollapsing}\Longrightarrow\text{general noncollapsing},\\
 \text{general noncollapsing}\Longrightarrow\text{strong canonical neighborhoods},\\
 \text{strong canonical neighborhoods}\Longrightarrow\text{finite geometric horizons}.
 \end{gathered}
\]
The last application is revision186's actual construction, not a new existence
field. Its initial identification, common actual cutoff records, boundary-frame
reversal, classified discarded components and geometric event outputs all
remain in the conclusion. The six supplier conclusions are also queried for
their transitive axioms on the pinned release. These bounded bindings do not
accept the whole general-flow parent.

\paragraph{The exact remaining quantifier order.}
In the retained small-scale contract, the initial stage and metric are fixed.
Given the horizon, canonical accuracies and comparison constants, derivative
and gradient constants, an admissible pinching function, and an above-scale
volume coefficient $\kappa_1>0$, the contract chooses $\kappa>0$. Only afterward
are the canonical scalar thresholds $q_{\rm can}$ and $q_s\ge q_{\rm can}$
chosen. The cutoff radius $r_0$, permitted surgery bounds, model accuracy,
model radius and model order may depend on those thresholds. Schematically,
\[
 \forall\,\text{outer data}\quad \exists\kappa>0\quad
 \forall\,(q_{\rm can},q_s)\quad
 \exists\,(r_0,\delta_{\max},\rho_{\max},\varepsilon_{\rm cap},D,m).
\]
Thereafter the estimates must hold for every admissible cutoff parameter set
and actual finite history in the stated class. The contract includes both
an event slab and every required closed prefix of the terminal continuation
slab. It combines already-established canonical, derivative, gradient and
spatial bounds with noncollapsing above $r_0$ to obtain noncollapsing below
$r_0$. A theorem for one supplied witness, one time, or one finite history
does not fill this premise.

The reduced-volume assembly chooses
$\kappa_1=c\sigma^3/(2C_0)$ before the thresholds, obtains the small-scale
coefficient $\kappa_2$, and returns $\min(\kappa_1,\kappa_2)$. Its permitted
cutoff parameters are chosen later by minima and maxima. The strong canonical
assembly uses this order in its actual induction through the event stages
and terminal continuation. That induction is already supplied by the library;
the missing general input cannot be replaced by a circular use of its output.

\paragraph{Where the inspected simply-connected proof uses topology.}
The entire small-scale proof was read. Its initial-time volume estimate,
low-curvature backward tracing, cap-window exclusion, ball-volume comparison,
and transfer from regular times retain general stages. The high-curvature
branch uses a spatial canonical witness and a uniform ball-volume theorem.
In the inspected specialization, simple connectivity of the initial carrier
is transported to the actual stage components and supplied to that theorem.

Revision181 already replaces this local use by an actual normalized component
volume reserve for a spatial witness when its alternative does not require
volume. It also produces such a reserve from a supplied noncollapsed round
model. What is still needed is a reserve with the required uniformity for the
actual history and its terminal continuation, obtained without assuming the
noncollapsing conclusion being constructed. The proved loss in revision181's
particular same-$\kappa$ substitution remains; it does not refute the inherited
classical theorem. The earlier Kleiner--Lott round-alternative comparison and
discard-convention distinction remain explicit.

\paragraph{Validation and remaining work.}
Two proof modules give three new declarations and apply six actual public
suppliers. The declaration export records eight producer-to-consumer paths,
exact closed theorem types and standard axiom evidence. Source locators and
reading limits are in \code{reference_checks_revision187.md}. This is a source-
checked and kernel-checked dependency reduction. It does not supply general
small-scale noncollapsing, a common global profile, a coherent tower, prime/port
history reconstruction or the final Geometrization theorem.
% END REVISION187 GENERAL CANONICAL DEPENDENCY

% BEGIN REVISION188 ROUND DEGREE VOLUME
\section[Round degree and volume]{Actual finite covering degree and round-component volume}
\label{sec:round-degree-volume-revision188}

Revision187 isolates the general small-scale noncollapsing contract. The
following local chain supplies a new quantitative input for its round branch.
It constructs the covering and volume bound from the actual round-component
data; a uniform bound on the resulting degrees is still a separate history
obligation. The construction uses the accepted PC release throughout.

\paragraph{Produced finite degree on the constant-curvature model.}
Let $(Z,h)$ be a compact connected smooth three-manifold with sectional
curvature $1/6$. The inherited general space-form theorem produces a
surjective smooth local diffeomorphism and covering
\[
 p:(S^3,6g_{\rm unit})\longrightarrow (Z,h),
 \qquad p^*h=6g_{\rm unit}.
\]
Only the covering sphere is used as a simply-connected space. Compactness
makes each discrete covering fiber finite; the actual universal-cover fiber
equivalence identifies its cardinality with
$k=|\pi_1(Z,z_0)|$. The proof establishes both finiteness and $k>0$ and
transports basepoints using connected-manifold path connectedness.

The actual local-isometry volume theorem gives
$p_*\mu_{6g_{\rm unit}}=k\mu_h$. Evaluating on the whole target and using
the inherited coarse sphere-volume bound yields
\[
 \operatorname{Vol}_h(Z)\ \ge\ \frac{8\sqrt6\,\pi}{k}.
\]
The numerator is a lower bound, not a claim of equality with the sphere's
volume. In Lean the proof retains the extended nonnegative measure and uses
finite positive division; no infinite-cardinality convention supplies a
spurious nonzero degree. These conclusions are returned by
\code{round_cover_with_fundamental_degree} and
\code{round_volume_lower_div_fundamental_degree}.

\paragraph{The degree belongs to the actual spatial carrier.}
For the inherited \code{SpatialRoundComponent} at $x$, write
$Q=R_g(x)>0$. Its actual map is defined on all of $Z$, maps onto the actual
carrier $U$, and compares the pullback of $Qg$ to $h$ with factors $1/2$ and
$2$. The resulting homeomorphism $Z\cong U$ induces a fundamental-group
isomorphism at any specified $y_0\in U$. The metric-comparison volume argument
therefore proves finiteness, positivity, and
\[
 \operatorname{Vol}_g(U)\ \ge\
 \frac{4\sqrt3\,\pi}{|\pi_1(U,y_0)|}\,
 \frac{1}{Q\sqrt Q}.
\]
This is \code{spatialRound_volume_lower_div_fundamental_degree}.
The degree is on $U$, not an unrelated model carrier. The target is not
assumed simply connected. The comparison calculation is adapted from the
PC simply-connected volume proof with its attribution retained; the volume
supplier and group identification are the new steps.

\paragraph{A fixed degree bound supplies the actual ball consumer.}
Suppose $N>0$ bounds the order of the actual carrier's fundamental group when
the spatial alternative does not carry its own volume estimate. Monotonicity
gives the normalized reserve
\[
 \nu_N=\frac{4\sqrt3\,\pi}{N}>0.
\]
The constructor \code{ReservedCanonicalWitness.ofDegreeBound} preserves the
same actual spatial witness and its cap-neck charts and proves its missing
reserve. It then feeds revision181's actual ball consumer, including the
inherited nonround branches. In particular, for fixed $\varepsilon,C_1,C_2,N$,
the theorem \code{degreeBoundCanonicalConsumer} chooses $\kappa>0$ before
any stage, metric or witness. For every matched witness and every $r>0$ with
$r^4|\operatorname{Rm}|^2(x)\le1$, it proves
\[
 \operatorname{Vol}_g B_g(x,r)\ge\kappa r^3.
\]
The order bound is required only for the missing-volume alternative. Its
finiteness is produced from the round geometry, so an infinite fundamental
group whose \code{Nat.card} is zero cannot satisfy the proof by convention.

This constant no longer comes from the noncollapsing coefficient of an
auxiliary round model. It does, however, depend on $N$. A bound selected
after an individual history or after the scalar thresholds does not meet
revision187's quantified small-scale contract.

\paragraph{Proposed initial-topology supplier: precise limits.}
A candidate way to choose $N$ first is to show that every relevant component
group is a free factor of one of the finitely many fixed initial component
groups. For a finitely generated group $G$, fix a finite Grushko decomposition.
A finite nontrivial free factor is freely indecomposable and not infinite
cyclic. Decompose its complementary factor, which is finitely generated as
a quotient of $G$. Uniqueness identifies the finite factor with one of the
finite indecomposable factors in the fixed decomposition. Thus the maximum
of $1$ and their orders bounds all finite free-factor orders.

This is an inference from the source-checked Grushko/Kurosh statements and
proofs in Scott--Wall, Sections2--3, especially Theorem3.5, printed
pages148--153/PDF7--9 of the retained two-page scan~\cite{ScottWall1979}.
The general decomposition statement is corroborated by Diao--Feighn,
Theorem1.1 and its following uniqueness paragraph, printed1836/PDF2
~\cite{DiaoFeighn2005}. Their algorithm for a graph of free groups is not
applied to arbitrary initial manifold groups.

The candidate is not yet a Lean history producer. The required free-factor
identifications through actual sphere cutting, capping and retained-component
maps, compact-manifold finite generation, the finite initial component family,
and the group-theoretic bound must all be supplied. An injection or a retraction
alone does not replace the free-factor hypothesis in this argument. No assertion
about uniformly bounded orders of all finite subgroups of an arbitrary finitely
generated group is being used. The candidate therefore remains a comparison
alongside the existing analytic reserve route, not an unproved replacement
silently selected in the global DAG.

\paragraph{Evidence and frontier.}
Three new proof modules give six new checked declarations. The fresh export
also queries the actual sphere-cover and covering-volume suppliers and the
existing reserved ball consumer: nine roots, eight registered application
paths, and only the three standard axioms. Exact hashes, theorem types,
source intervals, external source identities and limited errata searches are
recorded with \code{reference_checks_revision188.md}. The preexisting
Kleiner--Lott round-alternative and discard-convention comparison remains.
No uniform history degree bound, general small-scale noncollapsing producer,
coherent global tower/profile, complete prime/port reconstruction or final
Geometrization endpoint is accepted by this iteration.
% END REVISION188 ROUND DEGREE VOLUME

% BEGIN REVISION189 HISTORY DEGREE SMALL SCALE
\section[History degree and small-scale noncollapsing]{The exact small-scale history contract from a uniform degree bound}
\label{sec:history-degree-small-scale-revision189}

The local result of revision188 now composes with the full actual history
argument. The result is conditional on a uniform topological degree bound,
but its analytic conclusion is the unchanged PC contract
\code{SmallScaleNoncollapsingThroughSurgery P g}. There is no separate
assumed small-scale or strong-canonical conclusion in the new downstream
finite-horizon theorem.

\paragraph{One explicit degree premise on actual carriers.}
For an oriented stage $P$ and integer $N$, the predicate
\code{StageFiniteDegreeBound P N} asserts: for every actual whole component
$U=\operatorname{Comp}_P(x)$ and every $y\in U$, if $\pi_1(U,y)$ is finite,
then its order is at most $N$. This is a topological statement about actual
carrier subsets and basepoints. It includes an explicit finiteness premise;
an infinite group's \code{Nat.card} convention is not used as its order.

The predicate \code{UniformHistoryDegreeBound P g N} requires this stage
property for every stage of every actual retained-core history in the exact
inherited cutoff class, for all of that class's horizon and parameter choices.
That class includes the actual initial identification and metric identity,
the horizon bound, last-event/horizon equality, canonical cutoff records and
recentering control. The intended producer must therefore establish
\[
 \exists N\in\mathbb N,\quad N>0,\qquad
 \forall\,H\ \text{in the cutoff class},\quad
 \forall j,\quad \operatorname{StageDegreeBound}(H_j,N).
\]
Here $N$ precedes the histories, horizons and cutoff choices. The accepted
predicate is restricted to this exact class; it does not unnecessarily demand
a theorem about every history carrying only an initial identification.

The class does not explicitly supply boundary-frame reversal or the smooth
cut/cap-completion package used by some reconstruction adapters. A future
topological proof must derive what it needs from the available actual data,
or work directly with the topological cutting and capping. The analytic
consumer has not been strengthened to require those additional packages.

\paragraph{Actual spatial-volume application.}
The theorem \code{canonicalBallConsumer_of_stage_degree} chooses its ball
constant from $\varepsilon,C_1,C_2,N$. It applies revision188's actual
canonical-ball theorem to the supplied spatial witness. In the missing-volume
branch, that witness's actual round-component data prove finiteness of the
fundamental group of its whole carrier. The stage bound gives the integer
bound, and the constructed volume reserve supplies the estimate. The witness,
metric, center and actual carrier remain unchanged.

\paragraph{The inherited analytic argument, including its terminal branch.}
The new proof adapts the accepted small-scale argument at its spatial-volume
input and the stage-topology supply. It retains the low-curvature backward
tracing and cap-window exclusion, the cutoff-record estimates, the uniform
initial layer, the volume comparison and the transfer from regular times.
The inherited public three-case combination theorem is applied unchanged.
The source proof and exact adaptation diff were reviewed; this is not an
assertion that an uninspected citation supplies the new input.

Given the outer analytic data, the proof chooses
\[
 \kappa=c_T\min\{c_{\rm BG}\kappa_1,\kappa_W,\kappa_0\}>0,
\]
where $\kappa_W$ now comes from the fixed degree bound. This choice precedes
the scalar thresholds $q_{\rm can},q_s$. The threshold-dependent radius,
cutoff bounds and cap-model parameters are selected afterward, just as in
the inherited contract. No use of $\kappa$ as its own assumed round-model
noncollapsing coefficient occurs.

For every event slab the argument uses the degree bound on its actual stage.
For a terminal continuation it handles every required closed prefix by extending
the horizon. The library's extension keeps exactly the same stage family,
initial metrics and event count. Consequently the original stage degree bounds
apply literally to the extension. The extended history need not satisfy the
last-event/horizon equality of the original cutoff class, and the proof does
not assume that it does. Its pinching, derivative, gradient, spatial-canonical,
cap-window and above-scale estimates retain their actual transported forms.

\paragraph{The checked downstream chain.}
The public result
\code{smallScaleNoncollapsingThroughSurgery_of_historyDegreeBound}
is applied to revision187's general strong-canonical assembly, which already
supplies reduced-volume, pinching and temporal/spatial continuation inputs.
The resulting \code{general_strong_canonical_of_history_degree} feeds
revision186's actual finite-horizon construction. Thus, from one supplied
$N>0$ and its uniform history property, every prescribed $B>0$ receives an
actual finite history $K$ with horizon $B$, its actual initial identification,
one common cutoff parameter set, actual event records, classified discarded
components and geometric event outputs. The prefix relation to the zero-event
initial history is retained.

The order is
\[
 \begin{gathered}
 \text{one uniform history degree bound}\ \Longrightarrow\
 \text{exact small-scale noncollapsing},\\
 \Longrightarrow\ \text{general strong canonical neighborhoods},\\
 \Longrightarrow\ \text{actual finite geometric horizons}.
 \end{gathered}
\]
The last histories are still separate finite-horizon outputs, not a coherent
infinite tower or a common profile for all later geometric applications.

\paragraph{Evidence and remaining producer.}
Three proof modules and an export driver compile on the accepted release.
Four new declarations and three actual prior suppliers give seven standard
axiom queries, with six registered application paths. Exact closed types,
source hashes, the adaptation diff and reading limits are recorded with
\code{reference_checks_revision189.md}. No simply-connected target premise
appears in these public theorem types.

The remaining premise is the producer of the uniform finite component-group
order bound. Revision188's Grushko/free-factor route remains a candidate;
the actual component ancestry and finite connected-sum group suppliers were
reopened for comparison. A complete proof must bind the finite generation,
finite free-factor-order bound and actual history transport on compatible
carriers. These inputs, the coherent global tower/profile, complete marked
history reconstruction and final Geometrization endpoint are still open.
% END REVISION189 HISTORY DEGREE SMALL SCALE

% BEGIN REVISION190 ACTUAL FREE FACTOR HISTORY
\section[Actual free factors and history degree]{Actual free factors from cut/cap reconstruction to the history degree bound}
\label{sec:actual-free-factor-history-revision190}

Revision189 reduces its missing analytic input to a uniform finite component
fundamental-group order bound. This iteration proves the actual free-factor
transport needed for the proposed topological route. The remaining initial
group bound and the all-history completion premise are explicit; the resulting
analytic theorem is conditional on these two producers.

\paragraph{The group invariant and its scope.}
The predicate \code{IsFreeFactor G H} asserts that there exists a group $K$
and an actual multiplicative equivalence $H\cong G*K$. It is an abstract
free-factor statement up to isomorphism. It does not assert that a prescribed
core inclusion induces the displayed factor embedding. This scope suffices
for a bound on the possible orders of finite abstract free factors.

The proofs establish reflexivity, composition and transport under group
isomorphisms. They also split any chosen summand of an indexed free product
from the complementary indexed factors, using the actual universal-property
isomorphisms. Thus the conclusion includes a free complement; an injection
or a retraction alone has not been promoted to a free-factor conclusion.

\paragraph{Actual reconstruction and component groups.}
The inherited finite connected-sum Van Kampen theorem supplies
\[
 \pi_1\!\left(\mathop{\#}_{i\in I}M_i,y\right)
       \cong \mathop{*}_{i\in I}\pi_1(M_i,x_i)
\]
for the actual finite-sum carrier and arbitrary selected basepoints. Each
factor is therefore a free factor of the actual sum group. Basepoint changes
use the path-connected factor carriers.

Revision177 supplies an actual smooth reconstruction of the source as a
disjoint union of finite connected sums. Its assigned lists contain every
actual capped-component label. Choosing that label in its assigned list
and applying the preceding group theorem produces its free complement.
Sphere-times-circle factors remain among the other summands.

To transport this statement back to the source, the new proof establishes
that the fundamental group of a disjoint-union summand is isomorphic to the
ambient fundamental group based in that summand. The inherited retraction
proves injectivity. For surjectivity, every based loop stays in the clopen
summand; its restriction has exactly the same endpoints and maps to the
original loop. The resulting group equivalence and the actual reconstruction
homeomorphism produce a source basepoint $p$ and
\[
 \pi_1(\text{capped component},x)
       \ \text{as a free factor of}\ \pi_1(M,p).
\]
This argument applies to disconnected sources without asserting that the
whole source is path connected.

The actual event presentation, restricted to the appropriate component and
then to its retained summand, gives the retained component group. An additional
component-to-ambient group equivalence uses the inherited homeomorphism from
the disjoint union of actual components to the manifold. Consequently every
retained basepoint $q$ has some source basepoint $p$ for which
\[
                  \pi_1(Q,q)\ \text{is a free factor of}\ \pi_1(M,p).
\]
No model manifold or merely diffeomorphic component label is substituted for
the actual event carriers and presentation maps.

\paragraph{Finite traces and completed actual histories.}
Free-factor transitivity propagates the preceding statement through every
stage of an actual finite spherical cut/cap trace. A separate induction applies
to the actual retained-core history type, including its zero-event case. It
uses the actual initial identification to transport the initial stage group
back to the initial manifold. At each event it uses an explicitly supplied
\code{SmoothCutCapCompletion} to invoke the inherited smooth-to-spherical
transition bridge. The bridge retains the original presentation and core map.

The resulting \code{completed_history_freeFactor} quantifies every stage and
every actual basepoint. It imposes no simple-connectivity hypothesis and no
bound on the number of events. Its completion hypothesis remains present in
the checked theorem type.

\paragraph{The exact degree premise and its analytic consumer.}
The predicate \code{InitialFiniteFreeFactorBound P N} says: for every
$p\in P$, every finite group $G$ which is an abstract free factor of
$\pi_1(P,p)$ has order at most $N$. Finiteness is an explicit hypothesis;
the convention for \code{Nat.card} of an infinite type does not play a role.
The trivial group is included.

Given this initial bound and the completed-history ancestry, the proof
transports the group of each actual carrier
$U=\operatorname{Comp}_{H_j}(x)$ to the corresponding ambient group at its
actual point. It then applies the initial bound to the same finite group
$\pi_1(U,y)$. This proves revision189's exact \code{StageFiniteDegreeBound}.
If completion is supplied for every event of every history in the unchanged
\code{InCutoffClass}, the theorem
\code{uniformHistoryDegree_of_completed_freeFactorBound} produces the
unchanged \code{UniformHistoryDegreeBound P g N}.

The integer $N$ and its initial-group property precede all horizons, cutoff
parameters, histories and stage choices. An actual final consumer applies
revision189's general strong-canonical theorem to this uniform degree bound.
Thus the established chain is
\[
 \begin{gathered}
 \text{actual reconstruction and free-product complements}\\
 \Longrightarrow\ \text{completed-history free-factor ancestry},\\
 \text{initial finite-free-factor bound and all-history completions}\\
 \Longrightarrow\ \text{exact uniform history degree bound}\\
 \Longrightarrow\ \text{general strong canonical neighborhoods}.
 \end{gathered}
\]
The existing finite geometric-horizon consumer remains available downstream.

\paragraph{Two producer obligations, with no circular substitution.}
First, the initial-group bound is still to be produced. The source comparison
in revision188 proposes finite generation, a finite Grushko decomposition and
uniqueness to bound the finite free factors by the finite factors in one fixed
initial decomposition. Neither arbitrary subgroup injection nor finite
abelianization supplies that assertion. The newly inspected inherited
\code{FinitelyGeneratedFundamentalGroup} module proves selected model and
finite-sum cases; its general closed-three-manifold statement is a predicate
supplied to its relevant consumers. This observation concerns that inspected
module, not an assertion of absence from the whole library.

Second, completion must cover every history quantified by the analytic
contract. The inherited attaching/frame-reversal lemmas require the displayed
orientation conditions. They do not automatically establish completion on
all histories in that class. Possible routes are to produce compatible
completion data from the actual transitions or to prove the necessary
free-factor statement directly from their topological cutting and capping.
Actual core-component and cap-interior comparisons have been inspected for
this purpose; no missing implication is accepted here.

The completions of histories constructed \emph{after} the strong-canonical
theorem cannot circularly supply the all-history premise used to prove that
theorem. Revision189's cutoff class and quantifier order remain unchanged.

\paragraph{Checked scope.}
Four new proof modules and the export driver compile on the pinned accepted
release. Sixteen new declarations and three actual prior/library suppliers
give nineteen standard-axiom queries. The exact export has 259 declarations
and fifteen registered application paths, with no captured value cycle.
The additional unchanged upstream sigma-group module also builds. Reading
limits, attribution, hashes and conditional hypotheses are recorded in
\code{reference_checks_revision190.md}. The two producer obligations, coherent
global tower/profile, full marked reconstruction and final Geometrization
endpoint remain open.
% END REVISION190 ACTUAL FREE FACTOR HISTORY

% BEGIN REVISION191 ACTUAL CAP GROUPS
\section[Actual cap attachments and component groups]{Actual cap attachments and component fundamental groups (revision191)}
\label{sec:actual-cap-groups-revision191}

\paragraph{The general attachment argument retains the group.}
Let $U,V\subset X$ be open, with $U$ path connected, $V$ simply connected,
and $U\cap V$ simply connected. The inherited open-cover Van Kampen theorem
produces a group isomorphism
\[
                  \pi_1(U\cup V,y)\cong\pi_1(U,x)
\]
for arbitrary actual basepoints $x\in U$ and $y\in U\cup V$.
Simple connectivity includes nonemptiness, so the overlap supplies a basepoint.
The proof passes to the actual subspace cover, applies the inherited
amalgamated-product collapse, and changes basepoints along paths.
No simple-connectivity hypothesis is made on $U$.

Now let $V_i$, indexed by an arbitrary finite type, be pairwise disjoint open
simply connected subsets. Assume each $U\cap V_i$ is simply connected.
An induction proves path connectivity of $U\cup\bigcup_iV_i$ and preserves
its fundamental group as that of $U$. At the induction step disjointness gives
\[
       \left(U\cup\bigcup_{i<n}V_i\right)\cap V_n=U\cap V_n.
\]
This equality supplies the precise overlap hypothesis at that step. The
zero-attachment case is included. A finite-index equivalence transports the
argument from \code{Fin n} without changing the actual open subsets.
When these subsets cover $X$, the resulting
\code{star_cover_fundamentalGroup_equiv} compares the actual ambient group
to the actual retained-region group.

\paragraph{The actual child supplies every cover hypothesis.}
Fix a \code{SmoothCutCapTransition} $E$ and an actual retained child component
$c$. Let $C$ be its actual child core. The transition's boundary charts make
the core locally path connected; hence its selected connected component $C$
is path connected. This is a statement about the actual core, with no
simple-connectivity assumption on the parent.

Choose the fixed cap radius $r=1/2$. The inherited construction supplies an
open core neighborhood $U_c$ in the actual child carrier, a homotopy equivalence
$C\simeq U_c$, and the actual finite cap interiors $V_b$. Each $V_b$ is an open
ball, distinct cap interiors are disjoint, and each $U_c\cap V_b$ is the
actual simply connected annular overlap. The same construction proves
\[
                         U_c\cup\bigcup_bV_b=\text{child carrier}.
\]
Path connectivity of $U_c$ is transported by the actual homotopy equivalence.
The new star-cover theorem therefore applies with no additional cover packet
supplied by the caller. Composing with the core-neighborhood equivalence gives
\[
             \pi_1(\text{child carrier},y)\cong\pi_1(C,x)
\]
for arbitrary actual basepoints $x\in C$ and $y$ in the child carrier.

\paragraph{Transport to the selected punctured parent component.}
Let $W_c$ be the punctured component defined by the actual core retraction:
its points are those whose terminal retraction lies in the core component
selected by $c$. This is the inherited \code{puncturedCoreComponent}, not
an independently chosen model or a newly assumed component label.
The inherited component retraction proves $W_c\simeq C$; its homotopy stays
in that selected component. Applying its fundamental-group equivalence and
the preceding child comparison gives, at arbitrary actual basepoints,
\[
                  \pi_1(\text{child carrier},y)\cong\pi_1(W_c,z).
\]
The checked consumer \code{child_fundamentalGroup_card_eq_puncturedCore}
then proves equality of their \code{Nat.card} values. No finiteness is
inferred merely from that cardinal convention. These are abstract group
isomorphisms; this iteration does not prove agreement with a prescribed
inclusion-induced homomorphism.

\paragraph{What this closes, and the remaining parent-side step.}
The actual cap geometry now supplies the group comparison without
\code{SmoothCutCapCompletion}, boundary-frame reversal, or a simply connected
parent. This is a substantive part of a direct topological route around the
completion gate. To obtain the history degree bound by this route, one still
needs to show that the selected punctured-component group is an abstract
free factor of the actual parent-component group and transport that statement
through the retained history. An injection or a retraction alone is insufficient.
The initial uniform bound on finite free-factor orders also remains open.

Consequently the accepted revision190 theorem still retains its all-history
completion premise. The cap comparison alone has not removed that premise
from the analytic consumer, and no downstream completion is used circularly.
The inspected orientation proof which already assumes spherical capping
cannot produce its own boundary-frame hypothesis. The direct parent-side
argument remains a separate producer obligation.

\paragraph{Checked scope and source evidence.}
Two new proof modules and an export driver are freshly compiled on the pinned
accepted PC release. Nine new declarations and three actual upstream suppliers
receive transitive axiom queries. The selected declaration export records
actual proof applications, while the whole premise census remains open.
The unchanged upstream open-cover group module also builds. Exact source
paths, read intervals, hashes, attribution and remaining hypotheses are in
\code{reference_checks_revision191.md}. General history degree, the coherent
flow/profile, full marked reconstruction and the final Geometrization endpoint
are not accepted by this bounded iteration.
% END REVISION191 ACTUAL CAP GROUPS

% BEGIN REVISION192 FINITE FREE FACTORS
\section[Finite free factors and initial bounds]{Finite free factors and actual initial bounds (revision192)}
\label{sec:finite-free-factor-bounds-revision192}

\paragraph{Reduced words establish the general infinitude statement.}
For an indexed family of monoids, choose two distinct indices $i\ne j$ and
nonidentity elements $a$ and $b$ in those factors. Inductively prepend $a,b$
to construct a reduced word of length $2n$ for every $n\in\mathbb N$.
Its first index is different from $j$, so the next prepend remains reduced.
Different lengths distinguish these words. The inherited, proved normal-form
equivalence between reduced words and the indexed coproduct therefore gives
an injection of $\mathbb N$ into that coproduct.

The actual inherited Boolean-indexed/binary coproduct equivalence transports
this statement to arbitrary groups $A,B$: if both are nontrivial, $A*B$ is
infinite. The earlier special proof for two groups of order two is not used
as a substitute for this general assertion.

\paragraph{Finite factors and the Grushko prerequisites.}
It follows that a finite group cannot be isomorphic to $A*B$ with both factors
nontrivial. The new predicate \code{FreelyIndecomposable} records precisely
that implication; it permits the trivial group. The checked
\code{finite_grushko_conditions} separately retains nontriviality and proves
finite generation, this indecomposability, and exclusion of an isomorphism
with the infinite cyclic group.

For an actual isomorphism $G\cong A*B$ with $G$ finitely generated, the actual
coproduct projections are surjective onto $A$ and $B$. The inherited
surjective-image theorem proves both factors finitely generated. The
\code{freeFactor_fg_complement} consumer retains a chosen free complement,
its group structure, finite generation and the actual isomorphism. A merely
injective homomorphism does not supply this packet.

These facts discharge prerequisites of revision188's proposed Grushko route.
Existence and uniqueness of the needed finite decomposition, and general
compact-manifold finite generation at the accepted interface, remain separate
producer obligations. No external group theorem is installed as a Lean axiom.

\paragraph{Free factors of a finite ambient group.}
If $H$ is finite and $G$ is an abstract free factor of $H$, write
$H\cong G*K$. The infinitude theorem shows that either $G$ is trivial or
$K$ is trivial. In the second case the proof constructs $G\cong H$.
The resulting checked cardinal dichotomy is
\[
              |G|=1\quad\hbox{or}\quad |G|=|H|,
       \qquad\hbox{and hence}\qquad |G|\le |H|.
\]
Here finiteness of $H$ is explicit. The convention that \code{Nat.card} is
zero on an infinite type cannot justify this conclusion for an infinite $H$.

\paragraph{An actual initial integer, with its precise scope.}
Let $P$ be the actual initial stage. Compactness and the inherited manifold
component theorem give a finite index set of connected components. Choose
one actual point $x_c$ in each component carrier $P_c$ and set
\[
              N=1+\sum_c\operatorname{Nat.card}\pi_1(P_c,x_c).
\]
The actual component-to-ambient group equivalence and path-connected basepoint
change show that every ambient initial basepoint has group cardinal at most
$N$. This construction includes the empty stage, for which the sum is empty
and $N=1$. This cardinal bound alone says nothing about orders of finite free
factors of infinite initial groups.

Now assume every actual initial component group is finite. Applying the
finite-ambient free-factor bound proves revision190's unchanged
\code{InitialFiniteFreeFactorBound P N}, with $N>0$. Thus
\code{initial_freeFactorBound_of_finite_groups} is an actual producer for
this case, rather than a hypothesis naming the desired bound. It does not
assume simple connectivity and does not presume a classification of the
finite fundamental groups.

The integer depends only on the initial topology. It is selected before
horizons, cutoffs, histories or stages. Applying revision190's checked
completed-history consumer yields
\[
 \begin{gathered}
 \text{finite actual initial component groups}\\
 \text{and completion for every exact cutoff history}\\
 \Longrightarrow\ \exists N>0\;\text{uniform history degree}\\
 \Longrightarrow\ \text{general strong canonical neighborhoods}.
 \end{gathered}
\]
The completion hypothesis quantifies the unchanged \code{InCutoffClass};
it is retained in the final theorem type. It is not produced from histories
constructed only after the strong-canonical conclusion.

\paragraph{Remaining scope and evidence.}
General initial component groups can be infinite while possessing nontrivial
finite free factors. That case still needs the full initial-group producer.
The parent-side direct free-factor theorem, all-history completion, coherent
flow/profile, marked reconstruction and final Geometrization endpoint remain
open. Two new proof modules and an export driver are freshly checked, with
sixteen new declarations and five actual supplier queries. The reduced-word
normal form, coproduct equivalences, finite-generation transport and exact
component carriers were source-inspected; the earlier Scott--Wall and
Diao--Feighn comparison is reused with its recorded limits. See
\code{reference_checks_revision192.md} for exact paths, read intervals,
hashes and the distinction between source, kernel and document evidence.
% END REVISION192 FINITE FREE FACTORS

% BEGIN REVISION193 THIN OVERLAP GROUPS
\section[Disconnected overlap and collar free factors]{Disconnected overlap and actual collar free factors (revision193)}
\label{sec:thin-overlap-free-factors-revision193}

\paragraph{The distinction that matters for history degree.}
An injective map, even with a retraction, does not alone supply a free-product
complement. Revision190's finite-factor ancestry needs an actual free factor.
The inherited thin-overlap argument is therefore extended to construct the
complement, rather than being relabelled as the required theorem.

Let $U,V\subset X$ be open, with $U\cup V=X$, let $U$ be path connected and
$V$ simply connected, and choose $x_0\in U\cap V$. Assume that for every
$x,y\in U\cap V$ the set of endpoint-preserving homotopy classes of paths
from $x$ to $y$ is a subsingleton. This permits disconnected overlap: between
different path components there are no such paths. Write
\[
 C=\pi_0^{\mathrm{path}}(U\cap V),\qquad c_0=[x_0],
 \qquad I=\{c\in C:c\ne c_0\}.
\]
The checked construction gives an actual group isomorphism
\[
       \pi_1(X,x_0)\ \cong\ \pi_1(U,x_0)*F(I),
\]
where $F(I)$ is the free group on $I$. No finiteness of $C$, local path
connectedness of $X$, or simple connectivity of $U$ or $X$ is assumed.
The equivalence sends the actual inclusion-induced homomorphism from $U$
to the left coproduct homomorphism. It is therefore stronger than an
unmarked abstract group comparison.

\paragraph{Construction and both inverse identities.}
The inherited coherent-connector proof chooses paths in $U$ from $x_0$ to
every overlap point, compatible with every overlap path. Normalize these
connectors at $x_0$ and extend them to all points of $U$. In $V$, simple
connectivity gives coherent connectors from the same actual basepoint.
For an overlap point $x$, their ambient concatenation, first through $U$
and then back through $V$, is a loop $\beta_x$. Compatibility proves that
$\beta_x$ depends only on its overlap path component and that
$\beta_{x_0}=1$.

Set the marker of $c_0$ to $1$ and every other marker to its corresponding
free generator. The local functor on $U$ records conjugated core loops in
the left factor; the local functor on $V$ records differences of markers.
They agree on the overlap. The inherited full groupoid Van Kampen extension
produces the ambient-to-word homomorphism. Conversely, the actual inclusion
of $U$ and the loops $\beta_c$ produce the word-to-ambient homomorphism.

For the first inverse identity, glue ambient connectors by choosing the core
connector on $U$ and the cap connector elsewhere. Their conjugation functor
agrees with the composite on both open sets. The inherited extension
uniqueness theorem identifies the global functors. Normalization at the
actual basepoint makes the resulting conjugation the identity.

For the second identity, the forward map sends every actual core loop to
its left-factor image and sends $\beta_c$ to the marker of $c$. The binary
coproduct and free-group universal properties therefore give the identity
on every word. These are checked identities of actual homomorphisms;
neither is an additional hypothesis of the resulting equivalence.

\paragraph{The actual nonseparating collar application.}
Let $e:S\to X$ admit a two-sided collar $c:S\times\mathbb R\to X$ that is
an open embedding and agrees with $e$ at time zero. Suppose $S$ is compact
and simply connected, $X$ is Hausdorff, and
$U=X\setminus e(S)$ is path connected. Put $V=c(S\times\mathbb R)$.
The inherited collar definitions give an actual open cover $U\cup V=X$.
The new overlap homeomorphism is
\[
       U\cap V\ \cong\ S\times(\mathbb R\setminus\{0\}).
\]
The two open half-lines are contractible and disjoint. The inherited
path-containment argument shows that their union has at most one homotopy
class of paths between any two endpoints; product and homeomorphism
transport give the exact thin-overlap hypothesis. Meanwhile
$V\simeq S$ supplies simple connectivity of $V$.

Applying the proved group formula, followed by actual path-based basepoint
change, gives for every $x\in U$ the checked conclusion
\[
          \pi_1(U,x)\text{ is a free factor of }\pi_1(X,x).
\]
The overlap and cap hypotheses have been produced from the collar geometry.
There is no ambient simple-connectivity, finite-group, orientation, metric,
or smooth-capping-completion assumption. The path-connected-complement
hypothesis is retained: this is the nonseparating case, not yet a theorem
covering all surgery cuts.

\paragraph{Remaining applications and evidence boundary.}
Separating cuts need the corresponding component-wise free-product theorem.
Finite disjoint surgery families then need actual component incidence,
iterated cover control and transport to the precise punctured carriers used
in revision191. Only after that comparison may one replace revision190's
all-history completion premise in the exact cutoff class. General initial
finite-free-factor bounds, coherent general flow, marked reconstruction and
the final Geometrization endpoint remain open.

Two new proof modules and an export driver are freshly checked. Thirty-five
new theorem/construction queries and four actual supplier queries have only
the three standard classical axioms. Source locators, retained failed probes,
actual declaration paths and document checks are recorded separately in
\code{reference_checks_revision193.md} and \code{REVISION193_CHECKS.md}.
A declaration path is evidence of proof composition, not a certificate that
all global premises have been supplied.
% END REVISION193 THIN OVERLAP GROUPS

% BEGIN REVISION194 EVERY COLLAR COMPONENT
\section[Every component after one collared cut]{Every component after one collared cut (revision194)}
\label{sec:every-collar-component-revision194}

\paragraph{The exact single-cut conclusion.}
Let $c:S\times\mathbb R\to X$ be a two-sided collar of $e:S\to X$.
Assume $S$ is compact and simply connected, and $X$ is connected, Hausdorff
and locally path connected. Write $Y=X\setminus e(S)$. For every actual
point $x\in Y$, let $C_x$ be its connected component \emph{in $Y$}, with
its subspace topology and its actual basepoint. The new checked theorem is
\[
             \pi_1(C_x,x)\text{ is a free factor of }\pi_1(X,x).
\]
This is a statement about the original ambient point, not an unrelated
basepoint or independently selected model. Neither ambient simple
connectivity nor a separation hypothesis is an input to the final theorem.
The proof produces one actual free complement; no canonical decomposition
or uniqueness statement is required.

\paragraph{The connected-overlap group input is reused.}
For open $A,B\subset X$ covering $X$, with both members path connected and
their overlap simply connected, the inherited Van Kampen theorem gives the
actual Boolean-indexed coproduct isomorphism. Revision190's proved indexed
factor lemma gives a free complement for each member's group. Actual paths
in the member transport both the member and ambient basepoints to any
chosen point of that member.

A useful intermediate application has two disjoint open retained sets
$A,B$ and a simply connected open bridge $V$, with
$A\cup B\cup V=X$. Suppose $A,B$ are path connected and both $A\cap V$
and $B\cap V$ are simply connected. Then
\[
 (A\cup V)\cap(B\cup V)=V.
\]
The preceding free-product theorem makes the group of $A\cup V$ a free
factor of the ambient group. Revision191's actual cap-union equivalence
identifies that group with $\pi_1(A)$. This proves the required free-factor
statement for $A$; the same proved construction applies to $B$ with the
roles interchanged. No simple connectivity of either retained set is used.

\paragraph{Actual component incidence chooses the case.}
The inherited collar geometry supplies a negative point and a positive
point in $Y$. Every connected component of $Y$ equals the component of one
of these two points. This incidence theorem needs connectedness of $X$,
but not simple connectivity. Let $c_-$ and $c_+$ be their actual component
labels.

If $c_-=c_+$, the incidence statement proves that every connected component
of $Y$ is all of $Y$. The complement is nonempty, and local path connectedness
passes to this open subspace. It is therefore path connected. Revision193's
nonseparating theorem gives the free-factor conclusion for $\pi_1(Y)$.

If $c_-\ne c_+$, let $Y_-,Y_+\subset X$ be the actual images of the two
component carriers. They are disjoint and open, and
$Y=Y_-\cup Y_+$. Their path connectedness follows from local path
connectedness of $Y$ and the actual component carriers. These assertions
are proved using the distinct labels. The inherited theorem deriving
separation from ambient simple connectivity is not applied here.

\paragraph{The separating overlap hypotheses are produced.}
Put $V=c(S\times\mathbb R)$. In the separating case, actual collar
coordinates prove
\[
 \begin{aligned}
       Y_-\cap V&=c\bigl(S\times(-\infty,0)\bigr),\\
       Y_+\cap V&=c\bigl(S\times(0,\infty)\bigr).
 \end{aligned}
\]
For example, a point of the negative component in the collar cannot have
zero collar coordinate, because it belongs to the complement. If its
coordinate were positive, the inherited half-collar path would put it in
the positive component, contradicting disjointness. Conversely every
negative-coordinate point lies in the negative component. The positive
calculation is symmetric and is also checked.

The actual open embedding identifies these overlaps with their displayed
products. Contractibility of each half-line and simple connectivity of $S$
produce simple connectivity of both overlaps. Revision193 supplies simple
connectivity of $V$. Thus the actual sets satisfy every hypothesis of the
bridge application; both component groups are proved free factors.

\paragraph{Return to the precise component carrier.}
For any $x\in Y$, the subtype inclusion induces a homeomorphism from $C_x$
to its actual image in $X$. In the separating case the incidence theorem
identifies that image with $Y_-$ or $Y_+$. In the other case its image is
$Y$. Homotopy-equivalence group transport moves the corresponding proved
free-factor statement to $\pi_1(C_x,x)$, retaining the original ambient
basepoint. The single-cut theorem therefore has neither a missing case
nor a supplied image-comparison premise.

\paragraph{Scope and the next producer.}
This closes the group calculation for every component after \emph{one}
compact simply connected collared cut. Iterating over an actual finite
surgery family still requires collars inside the relevant intermediate
carriers, component incidence and restriction maps, and comparison with
revision191's precise punctured core. These are not consequences of merely
naming a finite family. The all-history completion gate, general initial
Grushko bound and final Geometrization endpoint remain open.

Two new proof modules and an export driver are freshly checked: fifteen
new theorems and six actual supplier queries, with only the three standard
classical axioms. Source bodies, exact locators, actual declaration paths
and the separate document evidence are recorded in
\code{reference_checks_revision194.md} and \code{REVISION194_CHECKS.md}.
% END REVISION194 EVERY COLLAR COMPONENT

% BEGIN REVISION195 FINITE CUT HISTORY
\section[Finite cuts and history ancestry]{Finite cuts and history ancestry (revision195)}
\label{sec:finite-cut-history-revision195}

\paragraph{The finite-family statement.}
Let $X$ be connected, Hausdorff and locally path connected, and let $S$ be compact
and simply connected. Suppose a finite family of two-sided collars
$c_i:S\times\mathbb R\to X$ has pairwise disjoint ranges, with zero slices $e_i$.
For every $x\in Y=X\setminus\bigcup_i e_i(S)$, the checked theorem gives
\[
       \pi_1\bigl(\operatorname{Comp}_Y(x),x\bigr)
            \text{ as a free factor of }\pi_1(X,x).
\]
Here the component carries its actual subspace topology. Empty families and
nonseparating cuts are included. Simple connectivity is required of $S$, not
of $X$. This extends revision194 by proving the intermediate restrictions and
component comparisons needed to iterate the single-cut theorem.

\paragraph{The invariant and its geometric producer.}
Induct over a finite set $J$ of already removed zero slices. Put
$A=X\setminus\bigcup_{j\in J}e_j(S)$ and $C=\operatorname{Comp}_A(x)$.
The invariant is that $\pi_1(C,x)$ is a free factor of $\pi_1(X,x)$.
The empty case is the actual homeomorphism from the component of connected $X$
to $X$. For a new collar $c_i$, there are two cases.

If its range misses $C$, the actual component containing $x$ is unchanged by
removing its zero slice. If the range meets $C$, disjointness from all previous
collar ranges and connectedness of the new collar range imply that the entire
range lies in $C$. Thus the inherited construction restricts $c_i$ to an actual
two-sided collar in $C$. Compactness of the zero slices makes $A$ open; local
path connectedness makes $C$ open. Connectedness of $C$ is supplied by actual
membership of $x$ in $A$, rather than an ambient simple-connectivity assumption.

The restricted collar complement is exactly the preimage of
$B=A\setminus e_i(S)$ in $C$. The inherited component-in comparison identifies
the image of its component at $x$ with $\operatorname{Comp}_B(x)$. The actual
subtype image homeomorphism transports revision194's one-cut free factor to
this precise carrier. Transitivity of the proved free-factor relation closes
the induction. These restrictions and homeomorphisms are constructed inside
the proof; none is an additional family-level premise.

\paragraph{Actual surgery tubes supply the family.}
For the actual tube system, openness of every removed middle band gives the
inherited middle-sphere collar. Its range is that very band. Disjointness of
the original full tubes implies disjointness of these ranges, and the union
of their zero slices is exactly the union removed in the definition of the
punctured core. The finite tube index and simple connectivity of the actual
$2$-sphere supply the remaining hypotheses. The finite-family theorem therefore
applies to the actual punctured core, with its original point retained as
ambient basepoint.

For a smooth cut-and-cap transition $E:P\to Q$ and a retained child component
$c$, the parent $p(c)$ is the inherited actual parent label. Restrict the tube
system to that parent component. The inherited construction gives the exact
component tube indices and open bands. It also gives a homeomorphism from
revision191's punctured-core component selected by $c$ to the corresponding
component of this restricted punctured core. These declarations have no
simple-connectivity assumption on the parent. The later simply connected
specializations in the same source modules are not used.

Revision191 identifies the actual child's group with that punctured-core group.
Combining the new finite-family factor with this equivalence gives
\[
 \pi_1(Q_c,y)\text{ a free factor of }\pi_1(P_{p(c)},z).
\]
The point $z$ is the actual punctured-core point mapped into its actual parent.
The checked component-to-ambient group equivalences then give, for every
$q\in Q$, an actual $p\in P$ such that $\pi_1(Q,q)$ is a free factor of
$\pi_1(P,p)$. A nonempty punctured-core carrier is produced using the inherited
core homotopy equivalence. No independently supplied group comparison,
boundary-frame reversal or capping-completion witness remains in this theorem.
The free factor is abstract; agreement with a prescribed inclusion is not
claimed by this history invariant.

\paragraph{Every retained history, with the exact initial identification.}
Induction over the actual history stages applies the preceding smooth-transition
result at every event and composes free factors. At stage zero the actual initial
identification gives the group equivalence with the original initial manifold.
The result includes zero-event histories. It requires neither a metric estimate
nor a completion assumption, and it applies before any later finite-horizon
construction is invoked.

Consequently the exact revision189 cutoff class is left unchanged. If one
positive integer $N$ bounds the orders of finite free factors of the initial
component groups, the new ancestry theorem and revision190's actual stage-degree
adapter give the uniform history degree bound with this same $N$. Revision189's
checked analytic consumer then supplies strong canonical neighborhoods.
The integer is selected from initial topology before the entire cutoff and
canonical parameter register. Completion of the histories is no longer a
premise of this degree argument; this avoids using a later constructed history
to justify the hypotheses needed to construct it.

\paragraph{The remaining initial-group question is precise.}
For general initial groups, the initial finite-free-factor bound remains an
explicit obligation. The source-compared finite-generation and Grushko route
is not silently accepted here. In the special case where every actual initial
component group is finite, revision192 already produces a positive uniform
initial free-factor bound. The new history theorem now makes that an actual
uniform history degree and strong-canonical producer without an all-history
completion assumption. This is a substantive special case, not a replacement
of the general Geometrization initial manifold.

\paragraph{Evidence and retained endpoint obligations.}
Three proof modules and an export driver are freshly checked: twelve new
proved declarations and ten actual supplier queries. The exact closed types
of the surgery and history producers retain the real stages and initial
identification, and contain no supplied simple-connectivity or completion
premise. Source locators, hypotheses, actual value paths and evidence scope
are recorded in \code{reference_checks_revision195.md} and
\code{REVISION195_CHECKS.md}.

This removes the completion gate from \emph{group ancestry and the degree
argument}. It does not prove all boundary-frame reversals or the complete
marked smooth reconstruction. General initial finite-free-factor bounds,
a coherent global flow with the required common data, the remaining analytic
and geometric consumers, and the complete marked Geometrization endpoint
remain obligations of the full program.
% END REVISION195 FINITE CUT HISTORY

% BEGIN REVISION196 INITIAL FINITE GENERATION
\section[Finite generation of initial groups]{Finite generation of initial groups (revision196)}
\label{sec:initial-finite-generation-revision196}

\paragraph{An actual producer for the initial-group prerequisite.}
For every actual initial stage $P$ and every point $p\in P$, the new checked
result is
\[
                         \pi_1(P,p)\text{ is finitely generated.}
\]
There is no finite-generation hypothesis, supplied triangulation, prescribed
Morse presentation, simple-connectivity premise or group-order bound. The stage
may be disconnected, and the carrier universe is unrestricted. The proof first
works in a path-connected compact metric space with a basis of contractible
neighborhoods, then supplies those hypotheses from actual manifold charts and
transports the result from the actual component to the original stage.

\paragraph{A finite cover with larger simply connected neighborhoods.}
Strong local contractibility means a basis of contractible \emph{neighborhoods};
these neighborhoods need not be open. This distinction is retained. Choose one
such neighborhood $A_x$ at each point of a compact metric space. Their interiors
form an open cover. The Lebesgue-number theorem gives $r>0$ so that every
$r$-ball is contained in the interior of some $A_x$.

Local path connectedness, supplied by the contractible-neighborhood basis,
allows a path-connected open neighborhood $V_x$ with
\[
                 x\in V_x\subset B(x,r/4).
\]
Compactness selects finitely many of these actual sets, still covering the
space. If selected $V_i$ and $V_j$ meet at $z$, then every point of their union
is at distance less than $r$ from the selected center of $V_i$: for a point
in $V_j$, the three distances through its center and $z$ are each less than
$r/4$. The union thus lies in one of the actual contractible $A_x$, hence in
a simply connected subspace $W_{ij}$.

This construction neither assumes nor proves that $V_i\cap V_j$ is connected.
Nor does it claim that $V_i\cup V_j$ is simply connected. Its precise output
is containment of each intersecting pair in a larger simply connected
subspace. The group proof below uses exactly that output.

\paragraph{Actual local connectors and finite transition loops.}
Fix an ambient basepoint $x_0$. In each $V_i$, choose a center $v_i$ and paths
from $v_i$ to all points of $V_i$. Choose an ambient anchor path from $x_0$ to
$v_i$. Their concatenations give connectors $B_i(x):x_0\to x$.
Because $V_i\subset W_{ii}$ and $W_{ii}$ is simply connected, any path in
$V_i$ from $x$ to $y$ satisfies the actual ambient groupoid identity
\[
                         B_i(x)\,p=B_i(y).
\]
The anchors may lie outside $W_{ii}$; the proof compares only their local
continuations inside $W_{ii}$, then composes with the same anchor.

For $x\in V_i\cap V_j$, form the ambient loop
\[
                         B_i(x)B_j(x)^{-1}.
\]
Its class is independent of $x$. Indeed, the middle path from $v_i$ to $v_j$
through $x$ lies in $W_{ij}$. Simple connectivity of that actual larger space
makes its class independent of the intersection point, even if the intersection
is disconnected. Composing with the two fixed anchors gives the required
identity in the ambient fundamental group. All inclusion maps and these
identities are proved in the actual groupoid, not supplied as cover metadata.
Choose one such transition loop for each nonempty pair intersection, and choose
the identity for an empty intersection. There are only finitely many choices.

\paragraph{Ordinary subgroup generation, not just normal generation.}
Let $H$ be the subgroup generated by these finite transition loops. At each
point $x$, select a cover member containing it and use its connector as a
global connector $d(x):x_0\to x$. There is no normalization requirement on
$d(x_0)$. For every local connector, the loop
$B_i(x)d(x)^{-1}$ is one of the proved transition classes, hence lies in $H$.
For a path $p$ contained in $V_i$, the local connector identity expresses the
normalized ambient loop $d(x)p\,d(y)^{-1}$ as a product of two such elements
or their inverses. Its class therefore lies in $H$.

The inherited path-subdivision proof is extended from a two-member cover to
an arbitrary indexed open cover. Each path admits a finite standard subdivision
whose segments lie in members of the cover. The actual concatenation identity,
including endpoint casts, and induction on its segments show that the normalized
image of every ambient path lies in $H$. This is a statement for an arbitrary
subgroup; no quotient by a normal closure is substituted.

For loops at $x_0$, normalization is conjugation by the actual loop $d(x_0)$.
Conjugation is onto, so the conclusion for all normalized loops implies
$H=\pi_1(X,x_0)$. Thus the finitely many transition classes generate the entire
group. The formal proof tracks the library's convention for multiplication of
endomorphisms when translating path concatenation into group products.

\paragraph{The actual initial manifold supplies the local and compact data.}
A real normed vector space has a contractible-neighborhood basis given by its
nonempty convex balls. An actual manifold chart transports such neighborhoods
through its restricted homeomorphism. This produces strong local contractibility
of the charted manifold; it does not infer openness of arbitrary contractible
neighborhoods.

For a compact Hausdorff manifold modelled on a finite-dimensional real normed
space, the inherited metrization theorem supplies a metric compatible with
the original topology. It is an auxiliary metric for the compactness argument,
not a curvature or flow metric. The finite cover and group proof therefore
apply to every path-connected such carrier.

For the actual stage $P$, take the actual compact connected component containing
$p$. Its inherited charts and path connectedness supply the preceding theorem.
Revision190's checked component-to-ambient group equivalence transports finite
generation to $\pi_1(P,p)$. No change of the initial manifold, a selected model
space or a universe-restricted carrier is made.

\paragraph{The exact next group-theoretic obligation.}
Applying revision192 to the newly produced finite generation gives, for every
actual free-factor witness $G\leq_*\pi_1(P,p)$,
\[
 \begin{gathered}
        G\text{ is finitely generated},\qquad
        \pi_1(P,p)\cong G*K,\qquad K\text{ is finitely generated}.
 \end{gathered}
\]
Here the displayed free product is a produced abstract equivalence, as required
by the existing degree argument. For finite nontrivial $G$, revision192 also
supplies free indecomposability and exclusion of the infinite cyclic case.
The remaining source-compared Grushko existence and uniqueness argument must
identify such $G$ with one of finitely many fixed finite initial factors and
bound its order. That step is not a consequence of finite generation alone
without the missing theorem, and it remains explicit. Revision195 has already
proved the subsequent transport through every actual retained history.

Five new proof modules and an export driver are freshly checked: thirty-one
new theorem/construction queries and six actual supplier queries. Source bodies,
locators, exact closed types and registered value paths are recorded in
\code{reference_checks_revision196.md} and \code{REVISION196_CHECKS.md}.
General initial order bounds, coherent global flow, marked reconstruction and
the full Geometrization endpoint remain open.
% END REVISION196 INITIAL FINITE GENERATION

% BEGIN REVISION197 FINITE KUROSH
\section[Kurosh and finite subgroups]{Kurosh and actual finite-subgroup comparison (revision197)}
\label{sec:finite-kurosh-revision197}

\paragraph{A newly bound external Lean supplier.}
The public \code{Arthur742Ramos/GraphCoveringTheory} repository contains a
Kurosh development for arbitrary subgroups of arbitrary indexed free products.
The selected source is its \code{main} commit
\code{c918a72c5170503c98fd66a0f3116b7bfcd1250c}, originally using Lean and
Mathlib4.32.0. Revision197 copies the nineteen-module Kurosh dependency chain
into an isolated compatibility directory and compiles it against the same
Lean4.33.1 and pinned Mathlib foundation as the accepted PC release.
The original checkout, PC release and historical overlays remain unchanged.

The compatibility patch explicitly restores the older definitional-transparency
setting, supplies product eta, fixes the starting vertex in a structurally
recursive path action, and removes a tactic after its goal has already closed.
The theorem hypotheses and conclusions are retained. Challenge files containing
statement-surface proof holes are excluded. The exact source, patch, compiled
artifacts and transitive axiom reports are recorded separately from the authors'
README and registry descriptions; those descriptions are discovery evidence,
not proof acceptance.

\paragraph{The precise Kurosh map matters.}
Let $P=*_{i\in I}G_i$ and $H\le P$. The supplier constructs the explicit
Bass--Serre tree, the quotient graph under $H$, its vertex stabilizers $H_a$,
and the free group $F$ contributed by the quotient graph. It supplies
\[
                \Phi:\Bigl(*_a H_a\Bigr)*F\ \cong\ H.
\]
Each $H_a$ is a subgroup of the actual $H$. The restriction of $\Phi$ to a
vertex factor is its actual inclusion into $H$. Central-vertex stabilizers
are trivial. Every other stabilizer is the actual subgroup
\[
                         H\cap gG_i g^{-1}
\]
for some $i$ and $g\in P$, regarded as a subgroup of $H$. The code retains
universe lifts and its canonical quotient-vertex index; no finite index set,
finite generation or finite-index subgroup hypothesis is introduced.

The proof map is injective by path lifting in the explicit tree and surjective
by actual vertex and quotient-edge generators. Revision197 checks the actual
inclusion agreement as a consumer of this map, rather than replacing it by an
unrelated abstract isomorphism. This distinction is necessary for the following
containment conclusion.

\paragraph{Finiteness singles out an actual factor.}
Revision192's reduced words show that an indexed free product with two
nontrivial factors is infinite. Consequently a finite free product with one
nontrivial factor has only trivial other factors. Induction on actual free-product
words then proves that the inclusion of that one factor is surjective. If the
whole free product is nontrivial, at least one factor is nontrivial, since the
same induction would otherwise make every word the identity.

A finite free group is trivial: transport to the standard free-group model
and apply its torsion-free theorem to every element of finite order. Apply
these facts to the actual Kurosh product when $H$ is finite and nontrivial.
Its unique surviving factor cannot be $F$. Hence the inclusion of one vertex
stabilizer $H_a$ is onto $H$, and $H_a=H$. The central case is excluded by
nontriviality. The remaining classification gives
\[
                  H\ \le\ gG_i g^{-1}
                  \qquad\text{for some }i\in I,\ g\in P.
\]
This is containment of the actual subgroup, not just an abstract group-type
comparison. The separate statement for arbitrary finite $H$ includes the
alternative $H=1$. In particular, empty index families cause no missing case.

\paragraph{An injective homomorphism with its ambient equation.}
From this containment, conjugate the actual subgroup inclusion by $g^{-1}$.
Its image lies in the range of the original factor inclusion $\iota_i:G_i\to P$.
The inverse of the checked range equivalence produces an injective homomorphism
$j:H\to G_i$ satisfying, for every $h\in H$,
\[
                           \iota_i(j(h))=g^{-1}hg.
\]
Thus a family of finite factors with $|G_i|\le N$ and $N\ge1$ has the actual
consumer
\[
        H\le *_{i\in I}G_i,\quad |H|<\infty
                  \quad\Longrightarrow\quad |H|\le N.
\]
The factor family itself need not be finite, provided that uniform bound is
supplied. This consumer establishes a useful special case and tests the full
Kurosh comparison. It does not assume or produce a finite-factor presentation
for an arbitrary initial manifold group.

\paragraph{What this closes, and what remains.}
The imported theorem and the new actual-containment proof replace the former
unbound Kurosh supplier in this group-theoretic route. They do not prove Grushko's
rank inequality or existence of a finite freely indecomposable decomposition.
Revision196 has already produced finite generation of every actual initial
group. The next comparison must construct the needed fixed finite decomposition
and identify its finite free factors, or supply another general order-bound
producer. In particular, an embedding into an arbitrary infinite factor alone
would not bound a finite subgroup's order.

Scott--Wall, Section3, Theorem3.1 and Corollary3.3, supplies the classical
comparison; Lemma3.4 and Theorem3.5 describe the further conjugacy/uniqueness
step. Their printed pages150--153, PDF pages8--9, were reopened visually.
General decomposition existence still requires the separate Section2 argument.
These source statements are not accepted as Lean proofs by citation.

Two new consumer modules and the nineteen-module external proof closure are
freshly compiled with an export driver. Nine new declaration queries and five
supplier queries admit only the standard three axioms. Exact scopes, sources,
patch and checks are retained in \code{reference_checks_revision197.md} and
\code{REVISION197_CHECKS.md}. General initial order bounds, coherent global flow,
marked reconstruction and the full Geometrization endpoint remain open.
% END REVISION197 FINITE KUROSH

% BEGIN REVISION198 FINITE DECOMPOSITION
\section[Grushko and initial decompositions]{Grushko and actual initial decompositions (revision198)}
\label{sec:finite-decomposition-revision198}

\paragraph{The rank supplier and its exact scope.}
The public \code{Arthur742Ramos/GrushkoNeumannTheorem} repository, pinned at
\code{main} commit \code{ef1f70db49792bb9bdb0c40eeec48168fca53393}, provides
an actual proof of
\[
                       d(A*B)=d(A)+d(B)
\]
for arbitrary finitely generated groups $A,B$. Here $d$ is Mathlib's
\code{Group.rank}: the least cardinality of a finite ordinary generating set.
It is neither normal rank nor the rank of the abelianization. Finite generation
is retained explicitly when using this natural-number invariant.

The selected fourteen-module dependency closure ends at
\code{MarshallHall.GeneralGrushko.rank_coprod_eq_add}. The original source uses
Lean/Mathlib4.32.0; its isolated port is checked against the accepted4.33.1
foundation. Two modules are byte-identical. Twelve restore the earlier
transparency option, and one also removes a tactic after simplification has
closed its goal. Exact original and ported hashes, attribution, and the patch
are retained. Unrelated Hall results and challenge surfaces are excluded.

The proof constructs a finite labelled graph from an actual generating tuple.
A minimal null path supplies a monochromatic run; unfolding, folding and vertex
removal produce a graph with fewer vertices. The actual generating property,
connectivity, fixed-point-free reversal and Euler bound persist. Strong
induction reaches a graph with one vertex and yields a separated generating
tuple of the original length. Partitioning that tuple between the two factors
gives the lower rank bound; the union of generating sets gives the upper bound.
The reduction packages occurring internally are proved by the supplier, not
assumed at its final rank theorem.

\paragraph{Strict rank decrease and a finite product.}
Revision192 already proves that an actual splitting $G\cong A*B$ of a finitely
generated group gives finite generation of both factors. If both factors are
nontrivial, their ranks are positive. The checked rank identity therefore gives
\[
                       d(A)<d(G),\qquad d(B)<d(G).
\]
Apply strong induction on $d(G)$. If $G$ is trivial, use an empty family of
factors and the actual empty-free-product isomorphism. If $G$ is nontrivial and
freely indecomposable, use the singleton family with factor $G$. Otherwise the
negation of indecomposability supplies an actual splitting into two nontrivial
factors. Apply the induction hypothesis to both and flatten their two indexed
free products using the disjoint union of the finite index sets.

The flattening map and its inverse are explicitly constructed using the
universal properties. Equality on every factor proves both inverse identities.
The resulting proposition includes an actual group isomorphism
\[
            G\cong *_{i\in I}G_i,
            \qquad I\text{ finite},
\]
where each $G_i$ is nontrivial, finitely generated and freely indecomposable.
The index type is small and the factors remain in the universe of $G$.
There is no claim of a finite generators-and-relations presentation. Infinite
cyclic factors are permitted; separating a free rank or proving uniqueness is
unnecessary for this existence statement.

\paragraph{The actual initial-manifold consumer.}
For every oriented closed three-stage $P$ and every $p\in P$, revision196
constructs finite generation of the actual $\pi_1(P,p)$ from compact manifold
charts, a finite path-connected cover, and coherent connectors. Feeding that
producer into the induction above gives
\[
                     \pi_1(P,p)\cong *_{i\in I_p}G_{p,i}
\]
with the stated finite family and factor properties. The theorem's only
external arguments are $P$ and $p$. The initial stage need not be connected,
simply connected or have finite fundamental group. No decomposition or finite
rank bound is assumed from the caller.

\paragraph{The remaining comparison is separate.}
A finite subgroup lies in a conjugate of one factor by revision197. If that
factor is infinite, this containment alone gives no order bound. The general
initial finite-free-factor bound still needs the comparison that identifies a
nontrivial finite \emph{free factor} with one of the fixed finite factors, or
another proof of the required uniform bound. The distinction between arbitrary
finite subgroups and finite free factors remains essential. No general history
degree or noncollapsing producer is claimed from the present existence theorem.

Scott--Wall, Section2, Corollaries2.1--2.2 and Lemma2.3, printed pages148--151,
PDF pages7--8, were reopened visually. Their rank and termination argument
supports this decomposition route. Their CW proof and the supplier's finite
graph proof are separate implementations of the rank theorem. Source review,
Lean checking and the actual initial-manifold application are recorded separately
in \code{reference_checks_revision198.md} and \code{REVISION198_CHECKS.md}.

Three new proof modules and the fourteen-module external closure are freshly
compiled with an export driver. Seven new declaration queries and four supplier
queries admit only the standard three axioms. General finite-free-factor order
bounds, coherent global flow, marked reconstruction and the full Geometrization
endpoint remain open.
% END REVISION198 FINITE DECOMPOSITION

% BEGIN REVISION199 GENERAL INITIAL BOUND
\section[General initial bound and analytic production]{The general initial bound and analytic production (revision199)}
\label{sec:general-initial-bound-revision199}

\paragraph{A shorter comparison than full uniqueness.}
The remaining order question concerns finite \emph{free factors}, not arbitrary
finite subgroups. Revision199 combines revision197's actual Kurosh containment
with revision198's actual finite indecomposable decomposition. It proves the
needed comparison by inheritance of a free factor inside a subgroup, without
requiring the full Grushko uniqueness or free-rank theorem.

First, a universe issue is resolved by an actual construction. Suppose
$B*C\cong S$, where $S$ and a chosen model $A\cong B$ lie in the desired
universe but the factors $B,C$ may be larger. The image of the injective map
$C\hookrightarrow B*C\cong S$ is a subgroup $K\le S$ in that universe. The
range equivalence gives $C\cong K$, and hence $S\cong A*K$. This supplies a
complement in the ambient group's own universe; it does not assume one.
Isolating a factor of an arbitrary indexed free product uses the actual
reindexing and coproduct equivalences, followed by this range construction.

\paragraph{Actual Kurosh intersections are free factors.}
Let $P=*_{i\in I}G_i$ and $H\le P$. Every chosen vertex stabilizer in the
Kurosh product is therefore a free factor of $H$. For any actual tree vertex
$x$, its quotient-orbit representative is carried to $x$ by an element of $H$.
Conjugation by that element gives an actual isomorphism of the two stabilizers.
Thus the stabilizer of \emph{every} vertex is a free factor, independently of
which quotient representative the supplier chose.

For the factor vertex represented by $gG_i$, the stabilizer identity reads
\[
                        H\cap gG_i g^{-1}\ \le\ H.
\]
The intersection, regarded as a subgroup of $H$, is a free factor of $H$.
No finite generation, finite index or finiteness of $H$ is needed. In particular,
if the actual inclusion of $G_i$ has image in $H$, the intersection at $g=1$
is isomorphic to the whole $G_i$, yielding $H\cong G_i*K$.

The checked binary consumer makes the map requirement explicit. If
$j:S\hookrightarrow A*B$ is an injective homomorphism and
\[
                         \iota_A(A)\subset j(S),
\]
then $A$ is a free factor of $S$. The proof transports the actual image through
the binary/indexed free-product equivalence. It uses the containment equation;
an arbitrary abstract embedding $A\hookrightarrow S$ would not suffice.

\paragraph{Identifying a finite free factor in the fixed decomposition.}
Let $X\cong *_{i\in I}G_i$ be a fixed decomposition into freely indecomposable
groups. Let $A$ be a nontrivial finite free factor of $X$, so that an actual
isomorphism $X\cong A*B$ is supplied by the free-factor witness. Transport the
canonical inclusion of $A$ to the fixed product and denote its image by $L$.
Revision197 gives some $i$ and $g$ with
\[
                           L\subset gG_i g^{-1}.
\]
Transport the conjugate inclusion of $G_i$ back to $A*B$. This produces an
injective homomorphism $j:G_i\to A*B$ whose image contains the \emph{actual}
canonical copy of $A$. Explicitly, for every $a\in A$, the containment proof
supplies $b\in G_i$ with $j(b)=\iota_A(a)$.

The inheritance theorem now gives $G_i\cong A*K$. Indecomposability of $G_i$
and nontriviality of $A$ force $K=1$, hence
\[
                              A\cong G_i.
\]
In particular this fixed factor is finite. This is the required finite-factor
comparison; it does not assert uniqueness of the entire decomposition or any
bound on arbitrary finite subgroups of an infinite factor.

\paragraph{A uniform initial integer.}
For a finite fixed decomposition, set
\[
                           N_X=1+\sum_{i\in I}\operatorname{card}(G_i),
\]
where Lean's natural cardinal is zero on infinite types. The trivial free
factor has order one. Every nontrivial finite free factor is isomorphic to one
of the $G_i$, so its order is at most $N_X$. The isomorphism is proved before
using this cardinal estimate; no bound is inferred from an infinite factor's
zero natural cardinal. The empty decomposition gives $N_X=1$.

Apply this to the actual initial groups produced in revision198. Compactness
and local connectedness give finitely many components of the initial stage
$P$. Choose one point in each component, take its group bound, and sum those
bounds with an extra one. Actual component/ambient group equivalences and
change of basepoint transport every finite free-factor witness to its chosen
component group. Consequently the checked theorem supplies
\[
       \exists N>0\ \ \forall p\in P\ \ \forall A\text{ finite},
       \qquad A\text{ a free factor of }\pi_1(P,p)\Longrightarrow |A|\le N.
\]
The only argument is $P$. Its carrier may be disconnected or empty, its groups
may be infinite, and its universe is unrestricted. The integer is selected
from initial topology before the metric, horizon and surgery parameter choices.

\paragraph{Actual analytic consumers and the precise remaining frontier.}
Revision195 already gives free-factor ancestry for every actual retained
history, with its actual initial identification, without a capping-completion
premise. Feeding the new initial bound into that theorem gives one order bound
for every finite component group at every stage of every history in the
\emph{unchanged} inherited cutoff class. This includes all stages indexed by
$\{0,\ldots,\text{eventCount}\}$ and the event/terminal prefixes required by the
analytic proof. No complete-history or simply-connected restriction is added.

The existing degree-to-volume and through-surgery noncollapsing argument of
revision189 now has its actual input. The checked consumers, with only the
initial stage $P$ and metric $g$ as arguments, produce
\code{SmallScaleNoncollapsingThroughSurgery P g} and
\code{CanonicalNeighborhoodsThroughSurgeryStrong P g}.
They retain the full parameter quantifiers of the inherited contracts.
For every prescribed $B>0$, the existing finite-horizon construction then
produces an actual retained history ending at $B$, its actual initial marking,
a common cutoff record, boundary-frame reversal, the inherited discarded
classification, and actual discarded-component geometric witnesses.

These are general \emph{finite} geometric histories. Choosing a separate
history for each $B$ does not prove that they are compatible prefixes of one
flow, that one parameter profile works at all horizons, or that marked smooth
reconstruction has been supplied. Those remain the next obligations; the
full Geometrization endpoint is not accepted by this iteration.

\paragraph{Sources and acceptance scope.}
Scott--Wall, Theorem3.1 and the preceding actual vertex-factor construction,
Corollary3.3 and the proof of Theorem3.5, printed pages150--153/PDF pages8--9,
were reopened. The checked route uses the stronger actual-intersection data
of the already accepted Kurosh supplier. The full uniqueness and malnormality
claims are not imported as unproved adapters. Existing PC and revision186--198
source/normalization checks remain in force where their statements are unchanged.

Six new proof modules and an export driver are freshly compiled. Seventeen
new declaration queries and nine supplier queries admit only the standard
three axioms. The precise type and proof dependencies, source hashes, inherited
port receipts and rejection controls are recorded in
\code{reference_checks_revision199.md} and \code{REVISION199_CHECKS.md}.
% END REVISION199 GENERAL INITIAL BOUND

% BEGIN REVISION200 ACTUAL CONTINUATION
\section[Actual continuation and time translation]{Actual continuation and time translation (revision200)}
\label{sec:actual-continuation-revision200}

\paragraph{The next interface after finite-horizon production.}
Revision199 constructs actual general finite geometric histories. Their
independently chosen horizons need not agree. To extend one history using
existence applied to its final stage, one must retain the already observed
smooth tail, translate the new construction's time origin, and preserve its
actual initial marking. This iteration proves those local adapters. It does
not infer a global flow from unrelated finite histories.

\paragraph{Uniqueness up to the closed endpoint.}
Let $F$ and $G$ be compact smooth Ricci slabs on $[a,b]$ and $[a,c]$, respectively,
with the same metric at $a$. The inherited compact uniqueness theorem applied
to their incoming restrictions gives equality on $[a,\min(b,c))$. For every
point and pair of tangent vectors, the metric coefficient is continuous on
the closed interval. The closure of $[a,\min(b,c))$ is $[a,\min(b,c)]$, since
both slabs have positive length. Thus the coefficient equality extends to the
endpoint, and metric extensionality gives
\[
                 g_F(t)=g_G(t)\qquad(a\le t\le\min(b,c)).
\]
The endpoint is included; equality on an open interval alone would not supply
the history-prefix contract.

Suppose instead that $G$ is an incoming slab on $[a,s)$ with a genuine singular
endpoint at $s$, and its metric at $a$ agrees with a closed slab $F$ on $[a,b]$.
Then $b<s$. Indeed, if $s\le b$, compactness and continuity give a uniform
curvature-norm bound $K$ for $F$ on $[a,b]$. Uniqueness transfers that bound to
$G$ on $[a,s)$. The singular-endpoint definition supplies a time and point with
curvature norm greater than $K+1$, a contradiction. This includes exclusion
of $s=b$, with no supplied metric-agreement or continuation package.

\paragraph{Actual prefix preservation.}
For a retained history $H$, write $a$ for its last event time and $b$ for its
horizon. A new retained event $E$ starts at the same actual last stage and metric.
Assume that its incoming slab has a singular endpoint $s$. If $a<b$, apply the
preceding result to $H$'s actual final closed slab; if $a=b$, the positive-length
incoming slab already gives $b<s$. In either case the new event occurs after
the observation horizon.

On the old final interval, uniqueness gives the exact metric equality required
by \code{appendEventCompatible}. The actual inherited append construction then
makes $H$ a prefix of the extended history, including the old events, stages,
metrics and smooth tail. Neither a master flow nor compatibility is supplied
as an extra premise. Similarly, a closed slab extending the last stage to
$T\ge b$ and agreeing with its initial metric automatically satisfies
\code{extendHorizonCompatible}; its actual extension preserves the prefix.
The zero-length final-tail case is handled separately and remains valid.

If $A$ is the given initial oriented smooth identification, the exact prefix
relation supplies equality of the initial stage and heterogeneous equality of
the initial metric. Transporting $A$ along those equalities gives $A'$ with
\code{A.IsPrefixOf A'}. In particular, the underlying initial map is preserved
up to the required heterogeneous equality. This is preservation of a chosen
initial marking; it is not the full marked connected-sum reconstruction theorem.

\paragraph{Changing the time origin preserves the actual surgery.}
For any real $c$, translate an incoming slab by
\[
                       g_c(t)=g(t-c),\qquad [a,s)\mapsto[a+c,s+c).
\]
The inherited slab producer already proves the Ricci equation, smoothness and
singular-endpoint equivalence. Revision200 proves that the actual terminal
regular open subset is unchanged. Translate the bounded-curvature witness time
$d$ to $d+c$ in one direction and back to $d-c$ in the other, retaining its
open spatial neighborhood and curvature bound.

The terminal metric is transported along this equality of open subsets. For
every compact subset, derivative order and positive tolerance, translating
the convergence witness time proves the full inherited terminal convergence
contract. No weaker pointwise or finite-order convergence is substituted.

The resulting retained event keeps exactly its original cut-and-cap transition,
discarded and capped stages, output metric, old-output map, and child-incidence
property. Its old-terminal map is transported to the equal regular open subset,
with the ambient inclusion identity and pullback metric identity proved there.
The actual initial metric at $a+c$ is the old initial metric at $a$.
Boundary-frame reversal and the inherited discarded classification therefore
feed the existing actual geometric-event consumer after translation as well.
No new geometric model or unrelated replacement surgery is chosen.

\paragraph{What these results do and do not close.}
These adapters remove real local assumptions from continuation. They are valid
for arbitrary compact oriented stages, including disconnected stages, and do
not assume simple connectivity or extinction. Finite-history concatenation
still has to align the first translated stage with the existing last stage,
iterate the actual event construction, attach the final smooth slab, and retain
all event properties. A coherent global controlled flow requires further work.

In particular, restarting with a different horizon-dependent cutoff register
is not itself the global parameter theorem. Kleiner--Lott's discussion on
printed2772/PDF186 explicitly explains why repeatedly estimating back to the
initial time would force changes to earlier surgery parameters. Section80
uses the previous interval and an inductive parameter selection. Revision179's
positive common-profile construction supplies numerical inequalities but cannot
retrofit geometric records to a different profile. These distinctions remain
visible at the global assembly interface.

\paragraph{Evidence and sources.}
Four new modules and the export driver are freshly checked on the accepted
PC release and its pinned Lean/Mathlib versions. Twenty new declarations and
nine suppliers are queried for their transitive axioms; actual type and value
dependencies and rejection controls are retained. Source bodies inspected
include \code{SlabUniqueness}, \code{SlabTimeTranslation},
\code{SlabTensorContinuity}, \code{ClosedSlabEndpoints}, the actual retained
history and prefix constructions, and initial-identification transport.
The source and correction scopes are recorded in
\code{reference_checks_revision200.md}; complete iteration checks are in
\code{REVISION200_CHECKS.md}. No full tower, common global cutoff profile,
marked reconstruction or Geometrization endpoint is claimed here.
% END REVISION200 ACTUAL CONTINUATION

% BEGIN REVISION201 COHERENT TOWER
\section[Finite concatenation and a coherent surgery tower]{Finite concatenation and a coherent surgery tower (revision201)}
\label{sec:coherent-tower-revision201}

\paragraph{Exact scope of the new producer.}
The general finite-horizon producer of revision199 and actual compatibility
proofs of revision200 now yield one coherent retained-core observation tower.
Its only initial arguments are a compact oriented three-stage $P$ and its
smooth Riemannian metric $g$. No simple-connectivity, finite fundamental group,
supplied tower, event schedule or extinction hypothesis is assumed. The tower
has actual retained-core events, and each event has a singular incoming endpoint,
boundary-frame reversal and the inherited standard-discarded classification.
These three local properties are named \code{SurgeryEventControl}; their universal
history version is \code{HistoryEventControl}. They do not assert a global
cutoff register or a canonical-neighborhood estimate for the entire tower.

\paragraph{Concatenating the actual events.}
Let $H$ be the existing history, let $c$ be its last event time, and let $K$ be a
finite suffix whose initial stage and metric equal the last stage and initial
metric of $H$. The suffix starts at its own time zero. The shift is by $c$,
not by the current observation horizon: the old last stage may already have a
nontrivial smooth tail. Suppose both histories have the three local event
properties and
\[
                  H.\mathrm{horizon}\le K.\mathrm{horizon}+c.
\]
Induct over the $k+1$ stage indices of $K$, where $k$ is its event count. At index
$i$, retain a history $L_i$ and the following invariant: $H$ is its actual prefix;
its event count is $H.\mathrm{eventCount}+i$; its last stage and metric are the
$i$th stage and initial metric of $K$; its last time is $K.\mathrm{time}(i)+c$;
it retains the three event properties; and its horizon is at most
\[
          \max\{H.\mathrm{horizon},K.\mathrm{time}(i)+c\}.
\]
The starting history is exactly $H$. At a successor, translate the actual
retained event of $K$ using revision200, and transport along the invariant's
stage and time equalities. Heterogeneous equality of the initial metrics proves
the actual append construction's metric hypothesis. Singularity and compact
Ricci uniqueness put the event strictly after the existing observation horizon
and prove agreement on the old smooth tail. Thus the actual append preserves
the prefix. Its event family consists of transported old events and this actual
translated event, so all three event properties persist. The output metric,
last stage, last time and event count satisfy the next invariant.

At the last index, attach $K$'s actual final closed slab translated by $c$.
Its initial metric agrees with the last metric just constructed. Revision200's
closed-slab uniqueness proves prefix compatibility, and the result has horizon
$K.\mathrm{horizon}+c$. If $K$ ends exactly at its last event, no positive-length
final slab is available or needed: the invariant and last-time bound already
force this horizon equality. Zero-event suffixes are also included. The result
has exactly the sum of the two event counts, the actual final stage and metric
of $K$, and all old prefix data. Transport of the given initial identification
preserves its underlying map as the required heterogeneous equality.

\paragraph{Extending every controlled finite prefix.}
Given $H$ with its initial identification and the three local event properties,
fix $B>H.\mathrm{horizon}$. Apply revision199's general finite geometric-horizon
producer to the actual last stage of $H$, its initial metric, and the positive
length $B-c$. Its initial-prefix witness supplies literal stage equality and
heterogeneous metric equality, not merely an unrelated diffeomorphism. Its
geometric cutoff records supply singularity; its other conclusions supply
boundary reversal and standard discarded components. The preceding construction
therefore extends $H$ to horizon $B$, preserving its original marking.
Every resulting event has the actual event geometry proved in revision186.
This is an extension theorem for the supplied finite prefix, rather than a
choice of independent histories for separate horizons.

\paragraph{Recursive global coherence and finite observations.}
Start with the actual zero history and its initial identification. Given the
chosen history at integer horizon $n$, use the extension theorem at $n+1$ and
choose one extension together with its witness. Ordinary recursion gives a
sequence with exact horizons $n$ and actual successor prefix relations. The
initial map equalities have precisely the orientation required by
\code{RetainedCoreObservationTower.initial_successor}. Thus this is a producer
of the inherited tower structure, including its compatibility fields.

The inherited observation construction restricts an integer prefix at any
nonnegative real horizon. Its restriction coherence and local event-finiteness
proofs now apply to this actual general tower. In particular, the union of its
event times meets every bounded closed interval in a finite set. Revision175's
regular-prefix consumer gives a regular observation beyond every prescribed
bound, retaining an initial identification with $(P,g)$. Its alternative
consumer also now has an actual general input: either one observation is empty,
all later observations stay empty and no later event occurs, or nonempty regular
observations occur beyond every bound. The empty alternative remains an analytic
statement, not a terminal exceptional-geometry certificate.

\paragraph{What remains at the analytic interface.}
The tower is now coherent, so raw coherence is no longer an open producer.
However, each restart used a finite-horizon cutoff register for its own initial
metric and interval. The proof does not transport those registers into one
preassigned global profile, nor prove the uniform normalized pinching,
noncollapsing or canonical-neighborhood conditions needed by the long-time
thick/thin arguments. The numerical common-profile construction cannot retrofit
already constructed surgery records. The Kleiner--Lott preceding-interval and
inductive parameter argument retained in revision200 therefore remains the
relevant global comparison. Full controlled flow, long-time estimates, marked
connected-sum reconstruction and the Geometrization endpoint are not conclusions
of this iteration.

\paragraph{Evidence.}
Six new modules and an export driver are freshly checked on the accepted PC
release:21 new declarations and11 suppliers, with transitive axiom queries and
actual type/value-dependency checks. The inspected source bodies cover actual
append and event transport, prefix transitivity, observation-tower restrictions,
and finite event sets; the new recursion and concatenation are proved here.
Exact passages and inherited correction-check limits are recorded in
\code{reference_checks_revision201.md}. The iteration audit is recorded in
\code{REVISION201_CHECKS.md}. Kernel checks, source review, static consistency
checks and document builds retain their distinct evidence scopes.
% END REVISION201 COHERENT TOWER

% BEGIN REVISION202 PLANNING WRAPUP
\section[Blueprint planning handoff]{Blueprint planning handoff (revision202)}
\label{sec:planning-handoff-revision202}

\paragraph{The current task and the preserved baseline.}
The active task is to make and verify plans for formalization and to finish the
blueprint handoff. Further Lean development is stopped. The preceding checked
work is preserved, not extended or discarded. Revision201's code and evidence
are on the new branch \code{codex/geometrization-baseline-201} in
\code{qinz1yang/differential-geometry-dev}, commit
\code{22b9d2a5a8323cff4bf5af72595a499b9e91f205}. The remote commit was verified
and the export checkout is clean. All152 existing pilot Lean sources match the
workspace at the stop instruction. No new Lean build or proof development was
performed to prepare that checkpoint. The repository's main branch and upstream
Lean library were not changed.

\paragraph{A DAG is available; its limits govern the plan.}
The existing consolidated graph is the authoritative planning view; a second
independent graph is unnecessary. It records1196 vertices and2443 declared
edges, with649 selected vertices and124 selected open leaves. These are
catalog and evidence counts, not counts of distinct mathematical gaps or a
percentage of a formal proof. In particular, an open inherited-interface alias
does not mean that its mathematics is absent from the accepted PC release.
The accepted public PC endpoint is distinct from every precise GC reuse binding.

The graph has no recorded cycle or unresolved reference. That check verifies
the recorded graph. It does not prove that every implicit premise was listed,
that every producer supplies the right output, or that all maps, domains and
quantifiers agree. The complete semantic premise census remains explicitly
unaccepted. Revision202 exports all124 selected open-leaf rows, with their
current statement locations, consumers, qualifications and obligation text,
for review. This is an exact inventory of the graph's current frontier, not a
new assertion that these are all missing mathematical theorems.

\paragraph{What the checked evidence now supplies.}
The frozen evidence includes the accepted public PC endpoint; bounded metric,
topology, event and exceptional-factor consumers; actual initial fundamental-
group finite generation and free-factor bounds; general finite geometric
histories; actual continuation and time translation; and the raw coherent
general surgery tower. Each retains its existing declaration and source scope.
The written developments in local collapse, compatible fibrations, long-time
thick/thin geometry, ambient cusp injection, graph refinement and model geometry
remain the blueprint's mathematical route under their retained foundations.
A written route, an inherited source result and a checked complete consumer
chain are different levels of readiness.

\paragraph{Ordered future planning priorities.}
First, specify the exact globally admissible flow required by the long-time
consumers. Its contract must quantify one chosen profile before later geometric
choices and retain the required normalization, pinching, noncollapsing, cutoff
and time-uniform bounds. Revision201 supplies raw coherence and finite event
observations, but not this stronger analytic package. A proposed intermediate
step is finite-horizon production with caller-prescribed positive accuracy
and radius bounds chosen before construction. The inspected source permits
such upper bounds in its step interface; the new finite and global conclusions
still require proofs. This is a future proof plan, not an accepted theorem.

Second, complete the exact interface census from the endpoint backwards, and
then replay it forwards on paper. Follow GA23, GA22 and LP12 into both the late
and terminal branches. For every required output field record its producer,
all hypotheses, carrier, metric, marking, scale, regularity and quantifier order.
Attach a named obligation wherever an implication or transport remains open.
Do not promote an entire parent contract because one bounded child is checked.
This review can proceed while the global analytic contract is being clarified.

Third, finish the planning specification for marked finite-history reconstruction.
Actual single-event geometry, initial-map preservation and group ancestry are
already available in their stated scopes. The remaining review must preserve
orientation, labelled cap and tube ports, factor multiplicities, cycle factors,
actual discarded carriers and the reconstruction map to the initial manifold.
Test the zero-event, empty-retained and repeated-isomorphic-factor cases on paper.
The absorbing-empty analytic alternative must be converted by actual terminal
geometric data, not by a general extinction assertion.

Fourth, expose the five existing written developments as bounded tasks with
complete input and output lists. Local finite-regularity models feed compatible
fibrations and static collapse; the actual thin carriers must then satisfy that
static theorem. Persistent hyperbolic embeddings precede ambient cusp injection.
Essential graph cuts and complete model metrics precede the late-slice geometric
certificate. At every junction retain the exact carrier and boundary markings.
Reuse the written arguments and their audits rather than restarting them.

Fifth, challenge the high-risk interfaces before any future implementation is
scheduled: finite regularity versus smooth curvature, actual-cloud hypotheses,
one-sheet preservation and proper boundary restrictions, whole-ball versus
center estimates, one common profile, minimal-area transport across surgery,
and marked connected-sum substitution. A review is resolved by a verified
implication, a corrected contract, or an explicit remaining obligation; a
successful document build is not a resolution of such a question.

Finally, review the resulting certificate against the public initial-manifold
statement, including both outcome branches, all required model geometries,
completeness on entire piece interiors, and the retained volume policy. Future
implementation and a final kernel build require separate authorization. They
are not part of this planning-only handoff.

\paragraph{Completion criteria for a planning packet.}
A packet is ready for later formalization when its exact statement and definitions
are fixed; every input has a checked supplier or named obligation; necessary
transport is stated explicitly; the proof is decomposed at mathematical
interfaces; relevant source passages and corrections are recorded; and the
intended downstream consumer has been compared against its output. Prioritize
uncertain shared prerequisites and long dependency chains. Do not measure
progress by page count, generated declarations or leaf counts alone.

\paragraph{Verification and remaining limits.}
The planning handoff, full open-leaf export and remote checkpoint receipt are
recorded in \code{PLANNING_WRAPUP_REVISION202.md} and the revision202 evidence.
The three new TeX files preserve revision201 and retain Arabic numbering from
Chapter1. Static checks verify preservation, current graph consistency and the
no-new-Lean boundary; the PDF build verifies document rendering. No new source
theorem, global analytic producer, complete semantic census or Geometrization
proof is claimed. This closes the planning handoff, not the mathematical
formalization project.
% END REVISION202 PLANNING WRAPUP

% BEGIN REVISION203 ENDPOINT INTERFACES
\section[Concrete endpoint and conditional assembly]{Concrete endpoint and conditional assembly (revision203)}
\label{sec:endpoint-interfaces-revision203}

\paragraph{The public Lean target and its evidence boundary.}
The team library defines \code{GC.Endpoint.GeometrizationConjecture} on actual
closed connected oriented smooth three-manifolds. Its certificate retains a
nonempty list of prime factors and the actual oriented connected-sum map;
compact cut carriers with exhaustive connected pieces; paired standard tori;
the actual boundary quotient with smooth interiors and two-sided seam charts;
ambient fundamental-group injection into each reconstructed prime; and complete
fixed-model metrics on every entire piece interior. Hyperbolic pieces retain
the E1 finite-volume requirement for that same metric. All eight geometry tags
are fixed, rather than supplied as arbitrary predicates. A complete sphere
certificate is proved using the accepted PC release. Defining the universal
proposition and proving this example does not prove the universal proposition.

\paragraph{Actual coordinate models.}
The five remaining reference metrics are now constructed from the coframes
already fixed in the endpoint interface. In coordinates $(x,y,z)$ these are
\[
\begin{array}{c|ccc}
\mathbb H^3 & e^{-z}dx&e^{-z}dy&dz\\
\mathbb H^2\times\R & e^{-y}dx&dy&dz\\
\widetilde{\mathrm{SL}}_2 & e^{-y}dx&dy&dz+e^{-y}dx\\
\mathrm{Nil} & dx&dy&dz-x\,dy\\
\mathrm{Sol} & e^zdx&e^{-z}dy&dz .
\end{array}
\]
Each metric is the sum of the squares of its displayed coframe. Continuous
bilinear forms, smoothness, symmetry, positivity and the exact tensor identity
are proved. Thus the former coordinate-atlas condition is equivalent to the
existing actual-metric atlas condition without a consumer supplying a new
metric or a tensor-agreement hypothesis.

For each model and points $a,b$, an affine coordinate map has the form
$F_{a,b}(p)=b+L_{a,b}(p-a)$. For the hyperbolic model, $L_{a,b}$ multiplies
both horizontal coordinates by $e^{b_z-a_z}$. For the hyperbolic product and
universal-cover model it multiplies the first coordinate by $e^{b_y-a_y}$.
For Nil it sends $(u,v,w)$ to $(u,v,w+(b_x-a_x)v)$; for Sol the first two
factors are $e^{a_z-b_z}$ and $e^{b_z-a_z}$. Direct calculation proves
$F_{b,a}=F_{a,b}^{-1}$ and preservation of the displayed coframes. The genuine
smooth diffeomorphisms therefore act transitively by metric isometries.

\paragraph{Completeness is proved, not assumed.}
The proof uses local compactness and transitive global isometries. Fix a compact
neighborhood $K$ of one point and a positive-radius metric ball contained in
$K$. A Cauchy filter has a member of smaller diameter. Translate the center of
$K$ to a point in that member; the entire member then lies in the translated
compact set. Compactness supplies its limit. The Lean proof works for extended
metric spaces and is applied to the actual Riemannian extended distance using
PC's diffeomorphism-pullback distance equality. This proves completeness of
all five coordinate metrics. Four give whole-model E1 structures directly;
whole hyperbolic space is not declared to satisfy finite volume. General
quotient metrics, their descent and the specified topology of geometric pieces
remain separate realization tasks.

\paragraph{A first compiling assembly network.}
Given endpoint certificates on $M$ and $N$, the constructor
\code{GeometrizationCertificate.connectedSum} concatenates their literal prime
slots, retains each slot's cuts and metric, and composes the actual oriented
connected-sum diffeomorphisms. Equal diffeomorphism types are not identified.
Induction produces a certificate on the finite connected sum of any family
whose members have certificates. For the empty family, the actual PC sum is
$S^3$ and the proved sphere certificate supplies the endpoint; an empty prime
list is never used.

The proposition \code{MarkedFactorReconstruction M L} requires an actual
orientation-preserving diffeomorphism from PC's finite connected sum of $L$ to
$M$. The theorem \code{geometrizes_of_markedFactors} consumes this map and
certificates for every member of $L$ to prove \code{Geometrizes M}. This is a
conditional composition with a concrete carrier and map. It does not produce
that map from a surgery history. Zero-cut, empty-family and repeated-sphere
consumers compile against the public definitions.

\paragraph{The next bounded interface phase.}
Work backwards from this target. First bind the actual final-stage components,
discarded components, cycle factors and initial identification of an observed
history; retain port and boundary compatibility rather than replacing them by
unlabelled manifold types. State the late and terminal producers separately and
prove their conditional compositions. Then specify the common controlled flow
profile before the time-dependent analytic choices. Finally expose the shared
local-collapse, static-collapse, persistent-hyperbolic, ambient-incompressibility
and relative-refinement interfaces, each with a concrete downstream use. These
are unfinished contracts, not new assumptions asserted as mathematical facts.

The existing DAG remains authoritative. Its GA06, GA14, GA22 and GA23 entries
now link partial Lean bindings and their consumers in the current-view overlay;
none of these parent vertices is accepted by this increment. A definition,
conditional theorem, actual proof and source-only record retain distinct
statuses. The semantic premise census and full proof closure remain false.
The shared source record and remaining hypotheses are in
\code{INTERFACE_ITERATION203.md}; document checks are in
\code{REVISION203_CHECKS.md}. The historical revision202 scope above is preserved
as history and no longer governs the active coding task.
% END REVISION203 ENDPOINT INTERFACES

% BEGIN REVISION204 OBSERVED HISTORY INTERFACES
\section[Actual history interfaces and endpoint descent]{Actual history interfaces and endpoint descent (revision204)}
\label{sec:observed-history-endpoint-revision204}

\paragraph{The concrete component predicate.}
For an actual PC closed oriented carrier $M$, define
\code{ComponentsGeometrize M} to require the public endpoint on every actual
connected component $M_C$. It is equivalent to \code{Geometrizes M} for a
connected nonempty input. Component restriction of a genuine oriented
diffeomorphism transports this predicate. For an empty stage the component
predicate is vacuous; this does not supply an endpoint on the nonempty initial
manifold. These adapters reuse the accepted component charts, orientations and
inclusion diffeomorphisms.

\paragraph{The exact remaining single-event map obligation.}
Let $E:M\rightsquigarrow Q$ be the actual spherical cut/cap transition, with
actual capped carrier $N$ and discarded carrier $D$. Revision177 already
produces finite data $R$: groups $w$, nonempty nodup lists $L_w$ of actual
components of $N$, sphere-product multiplicities $k_w$, and a smooth map
\[
 R:M\ \simeq\ \coprod_w
 \mathop{\#}\bigl((N_K)_{K\in L_w},\ k_w(S^2\times S^1)\bigr).
\]
Each capped component occurs in exactly one literal slot. The actual event
presentation identifies it, with an oriented comparison, with an actual
component of $Q$ or $D$. None of these statements by itself asserts orientation
preservation of $R$.

The proposed Lean structure \code{OrientedEventReconstruction E} retains this
same data and requires a bijection $s$ from its groups to the source components,
plus an oriented diffeomorphism $e_w$ from the displayed connected sum to
$M_{s(w)}$. Its compatibility equation is
\[
                 R\bigl(\iota_{s(w)}e_w(x)\bigr)=(w,x).
\]
Thus an unrelated map of the right diffeomorphism type does not meet the
interface. The original target-slot bijection remains proved. Existence of
this oriented refinement is an explicit unfinished producer, not a theorem
asserted by defining the structure. The structure does not yet specify fixed
protected ports or equality with an independently chosen cycle ledger; those
remain AT21/AT24 obligations.

Given this refinement, endpoint certificates on every component of $Q$ and
$D$, and a certificate on the exact PC sphere-product carrier wherever
$k_w>0$, the theorem \code{OrientedEventReconstruction.componentsGeometrize}
proves the endpoint on every component of $M$. Its proof classifies each actual
capped slot using the event presentation, transports the relevant certificate,
and invokes the finite-family assembly from revision203. Repeated geometric
types retain their literal positions. Zero multiplicity requires no cycle
certificate.

\paragraph{Backward induction on the same observed history.}
For an actual \code{ObservedHistory H}, \code{HistoryEndpointInputs H} records
the completion of each actual event, the preceding oriented refinement, and
the discarded and positive-multiplicity cycle certificates. All event indices
belong to this same $H$; no abstract replacement family is substituted. Given
certificates on all components of the final stage, finite reverse induction
proves the component predicate on every stage, including stage zero. The
supplied \code{InitialIdentification} then transports this to the original
manifold. Its actual map, orientation and metric identity are retained; the
topological endpoint consumer needs only the first two fields.

The resulting proof includes zero-event histories and disconnected intermediate
stages. A concrete sphere example passes through the actual zero history and
its identity initial marking, with no event, cycle or reconstruction hypothesis.
If the actual final stage is empty, the same induction uses the vacuous final
component predicate. Its terminal consumer still requires the recorded event,
discarded and cycle inputs; it has no late-slice or extinction premise.

\paragraph{Two open supplies with a checked downstream use.}
For a supplied observation tower $T$ over the original metric, define
\code{ObservedTopologySupply T} as production of these inputs for every finite
observation. Define \code{LateComponentSupply T} by
\[
 \begin{gathered}
 \exists B\ \forall t>B:\quad
 \bigl(t>0,\ t\notin\mathcal S_T,\ t_{\mathrm{last}}<t,
       \ M_t\ne\varnothing\bigr)\\
 \Longrightarrow\ \operatorname{ComponentsGeometrize}(M_t).
 \end{gathered}
\]
Here $M_t$ and $t_{\mathrm{last}}$ refer to that exact observed prefix. The
threshold precedes the time choice. The late supply is an explicit geometric
output obligation, not a disguised proof of thick/thin production or a
definition of a controlled flow.

The proved theorem \code{geometrizes_of_observation_supplies} consumes both
supplies and the existing marked late-nonempty/absorbing-empty dichotomy. In
the empty branch it selects the actual empty prefix and uses terminal descent;
in the late branch it chooses a regular nonempty prefix after $B$ and invokes
the late supply. Both paths use the original initial marking. A direct
consumer applies this composition to \code{RawSurgery.ofInitial}, with the two
supplies still explicit hypotheses. It does not derive them from raw coherence.

\paragraph{Scope, source comparison and remaining work.}
This is a conditional endpoint composition, not AT24's stronger flattened
sphere/torus ledger or its actual transported torus maps. The current endpoint
is expressed on capped prime representatives, and this consumer retains their
existing cut and metric data. Full relative reconstruction, cap avoidance,
protected collars and the prescribed cycle count remain visible in the DAG.
KL Lemmas67.5,67.13 and73.4 and Remark73.3 were reopened in the archived
February2013 version: printed2740--2741,2745,2761--2762/PDF154--155,159,175--176.
They justify the reconstruction route and distinguish discard conventions;
they do not supply the additional Lean marking equations by citation.

Existing AT21, AT24, GA20--GA22 and LP12 entries now link the corresponding
partial interfaces and compiling consumers, without accepting their parent
contracts or changing the recorded dependencies. The source checks and exact
unfinished producers are recorded in \code{reference_checks_revision204.md}
and the shared \code{INTERFACE_ITERATION204.md}. Next work should produce the
oriented event refinement and exceptional endpoint witnesses, and expose the
common controlled profile before binding late analytic producers. The full
Geometrization theorem and semantic premise census remain open.
% END REVISION204 OBSERVED HISTORY INTERFACES

\section{Acceptance examples and legacy-name crosswalk}
\label{sec:assembly-acceptance-crosswalk}

\begin{lemma}[Singleton and empty-torus acceptance template; GA24]
\label{lem:assembly-singleton-template}
Given an actual closed oriented prime carrier $P$, a complete model
metric satisfying E1, and an oriented diffeomorphism $P\to M$, the
singleton prime family, empty torus system and single cut carrier $P$
give a certificate by GA14. Disjointness and pairing obligations for the
empty torus/side sets are vacuous; the actual carrier and metric are not.
\end{lemma}
\begin{proof}
Use the identity cut comparison and the supplied reconstruction. The
single component is exactly the complement of the empty torus set.
Insert the supplied metric family and apply GA14. This is the common
template for $S^3$, $T^3$, a closed hyperbolic manifold and $S^2\times S^1$
once their actual prime/model witnesses have been supplied. No toy K0
record proves those geometric witnesses.
\end{proof}
Further acceptance cases include a graph region with relative internal
cuts, a torus bundle retaining Nil, a twisted interval bundle using E1,
a self-paired cut carrier, parallel seams, disconnected late slices,
two equal-type exceptional factors, a zero-event history and an explicit
empty terminal slice. These are implementation specifications, not new
kernel-checked examples.

\paragraph{Rejection obligations.}
Reject a missing complementary component; one unmatched side; a
fixed side confused with a loop edge; two distinct seam labels with the
same image; loss of an external boundary label; unsupported ambient
injectivity; restriction of a complete metric to a finite core; a
hyperbolic volume proof for a different metric; arbitrary homeomorphic
tensor transport; boundary equality without smooth collar comparison;
double-counted cycle factors; a nonempty initial carrier reconstructed
from an empty ledger; an outcome hidden in the raw flow input; or use
of the endpoint in an upstream producer. Source-qualified inputs must
not be silently changed to unconditional conclusions.

\begingroup\normalsize
\setlength{\tablebodywidth}{\dimexpr\linewidth-4\tabcolsep\relax}
\begin{longtable}{P{0.35\tablebodywidth}P{0.65\tablebodywidth}}
\textbf{Retained M1 name or stage} & \textbf{Current mathematical binding}\\\hline
\endhead
F1 / {\upshape\code{RegionNormalization}} & GA05, HG16, and the actual cut-carrier comparison.\\
F2 / {\upshape\code{ThinGraphToGeometricPieces}} & Qualified GA06, relative essential refinement, then GA09--GA10.\\
F3 / {\upshape\code{ExceptionalPrimeToGeometricPiece}} & Qualified GA07 and the narrow adapter.\\
F4 / {\upshape\code{TorusSystemFormation}} & GA03--GA04 on the supplied chosen prime carriers.\\
F5 / {\upshape\code{GlobalCutDecomposition}} & GA11 and the full supplied cut/complement comparison.\\
F6 / {\upshape\code{GluingReconstruction}} & GA12--GA13, with sphere and torus reconstruction separate.\\
F7 / {\upshape\code{CertificatePackaging}} & GA14 / {\upshape\code{CertificateOfAssembledData}}; GA15 is its existential wrapper.\\
{\upshape\code{assemble_geometrization}} & GA17--GA18; structural theorem, qualified conversions explicit.\\
{\upshape\code{AssembleLateThickThin}} & GA19 converts a supplied late witness; LP08/LP09 produce it.\\
{\upshape\code{TerminalExceptionalToGeometricPiece}} & GA07/GA20 on explicit terminal data, not a late-output projection.\\
{\upshape\code{long_time_outcome}} & LP12; invokes the two conditional branch producers.\\
{\upshape\code{geometrization}} & GA23 after GA22, actual global flow existence and the accepted adapter.\\
\end{longtable}\endgroup
The names are a reviewed crosswalk, not assertions of definitional equality
between undeclared Lean objects. Their authoritative signatures are bound
only after the inherited environment is accepted.

\section{Complete Chapter 16 audit and dependency boundary}
\label{sec:assembly-complete-audit}

All seven retained M1 sections and GA01--GA24 have individual review
records in {\upshape\code{chapter16_audit_revision103.csv}}. Blueprint topic coverage
is complete. GA06, GA07 and GA17 remain qualified producers; the other
entries are definitions or written conditional arguments. Their inherited
and upstream mathematical premises remain required. Neither the conditional
composition nor a document build establishes their truth in Lean.

\paragraph{Three dependency kinds.}
The inventory distinguishes the definition of an input record, a proof
using a supplied value of that record, and a theorem producing the value.
In particular, the pure GA14 proof uses GA01/GA08; producing GA01 from a
flow uses LP08 and AT31. There is no edge from the pure assembly theorem
to the analytic producer just because its input can be produced analytically.
The typed integration ledger records both uses and producers explicitly.

\paragraph{Relations with the rest of the book.}
Chapter2 supplies the audited foundation. Chapters3--4 feed the local
collapse inputs in Chapters13--14, which feed LP06, not GA14's internal
proof. Chapters5--6 produce relative primes and essential cuts; Chapter7
produces the general model witnesses. Chapter8's static HG02/HG16 data
support GA05, while its rigidity/stabilization/ambient-injection producers
feed Chapter12. Chapters9--11 provide actual flow existence, estimates,
events and LP01. Chapter12 produces the late-or-terminal outcome;
Chapter15 supplies actual reconstruction and generic marked transport.
Chapter16 consumes these results and never supplies an upstream premise.

\paragraph{Remaining release obligations.}
Retain finite-volume rigidity and quantitative stability, the controlled
smoothing in hyperbolic persistence, all least-area/surgery/boundary
inputs, LP01/LP03/LP06 and H1--H4, the qualified LC/FC producers, relative
prime/good-block refinement, GA06/GA07 model realization, actual event
classification and the accepted PC release/axiom audit. Optional AT27
and canonical JSJ uniqueness are not made prerequisites for the chosen
existential output. Finite-time extinction stays outside the default route.

\paragraph{Verification scope.}
The chapter checker verifies coverage, source links, typed stage order,
recorded producer reachability and forbidden upstream/terminal dependencies.
Specification rejection controls test missing contract evidence. These are
static consistency checks. The mathematical source comparisons, written
proof review, PDF build and any future Lean checks are reported separately.

\section{Outcome records, propositions and one choice boundary}
\label{sec:assembly-outcome-sort}

The intended \code{LongTimeOutcome M} is a data type with two
constructors: one carries an actual \code{LateThickThinOutcome M}, and
the other carries an actual \code{TerminalExceptionalOutcome M}. Thus
it is a sum of records, not Lean's proposition-valued \code{Or}. Its
eliminator may construct a \code{GeometrizationCertificate M} by
applying the appropriate supplied branch function. Conversely
\code{Geometrizes M} remains the proposition
\code{Nonempty (GeometrizationCertificate M)}. This sort convention is
an implementation requirement, not a new mathematical hypothesis.

Existence-producing source theorems may naturally return
\code{Nonempty (LongTimeOutcome M)}. In that form the public proof maps
the two branch constructors over nonemptiness and concludes
\code{Geometrizes M}; it need not choose a particular output object.
Where the retained data-returning signature of \code{long_time_outcome}
is used, expose a single named noncomputable choice wrapper after proving
nonemptiness. Likewise a supplied
\code{Nonempty (RicciFlowSurgeryPackage M)} is an existence theorem,
not an already chosen flow value. Do not pattern-match a general
existential or \code{Or} proof into arbitrary structure data without
that declared choice boundary.

LP12 may use classical decidability to split on
\code{LateTimeAvailable flow}; the chosen outcome still carries only
one branch. GA14 is a pure constructor from supplied indexed fields;
the existence of those fields belongs to their producers. GA10 may
choose from proved nonempty witness fibers, but cannot turn an unproved
fiber into an inhabited one. These conventions preserve the mathematics
of the late-or-terminal route while making its future declaration sorts
unambiguous.

\appendix

\chapter{Implementation scaffold and work-package crosswalk}
\label{ch:implementation-crosswalk}

% BEGIN migrated B054
\section{Proof architecture}
\label{ch:architecture}

The selected route is Hamilton--Perelman Ricci flow with surgery, organized
using Kleiner--Lott and Morgan--Tian as detailed expositions. Perelman's papers
are the primary mathematical sources. The proof is arranged as
\begin{equation}
\begin{aligned}
 &\text{smooth flow} \\
 &\quad\longrightarrow \text{Perelman estimates} \\
 &\quad\longrightarrow \text{surgery} \\
 &\quad\longrightarrow \text{thick/thin limits} \\
 &\quad\longrightarrow \text{geometric decomposition}.
\end{aligned}
\end{equation}
% END migrated B054

% BEGIN migrated B055
\section{Work packages}

\begin{description}[leftmargin=3.1cm,style=nextline]
\item[W0: project audit] Inventory the Poincar\'e Lean project, mathlib APIs,
  existing Ricci-flow code, imported axioms, and the build/CI environment.
\item[W1: topology] Formalize prime pieces, connected sums, embedded spheres
  and tori, incompressibility, graph manifolds, Seifert pieces, JSJ data, and
  the eight model geometries. Prove the topological endgame from thick and
  thin output.
\item[W2/M2: inherited foundation] Audit and import the completed PC metric,
  tensor, connection, curvature, volume, geodesic, injectivity-radius, and
  static/time-dependent metric interfaces. Add only small adapters and record
  genuine gaps; do not duplicate the geometric foundation.
\item[W3/M3: Ricci-flow PDE] Formalize the genuinely missing time-dependent
  metric, Ricci-flow equation, short-time existence and uniqueness, evolution
  identities, and PDE regularity interfaces. Reuse the M2 foundation and
  imported PC declarations wherever the audit permits.
\item[W4/M4: Perelman estimates] Audit and import the PC $F$ and $W$ functionals,
  conjugate heat equation, monotonicity, reduced distance/volume,
  no-local-collapsing, curvature-scale control, and three-dimensional
  singularity inputs. Formalize only missing long-time or surgery-compatible
  statements.
\item[W5/M5: surgery] Import and audit the completed PC global typed flow with
  surgery, canonical-neighborhood theorem, surgery-event records,
  continuation rules, and invariant-preservation theorems. Define only the
  GC flow/history adapter and exact inputs needed by M6; the long-time
  conclusions remain separate M6 theorems.
\item[W6/M6: long time] Formalize thick/thin regions, rescaling and pointed
  convergence, hyperbolic limits, lower-curvature collapse, and the
  graph-manifold theorem for the collapsed part.
\item[W7/M7: integration] Compose W1 and W6, import the Poincar\'e adapter,
  remove temporary assumptions, and run the kernel/provenance audit. Finite
  extinction is excluded unless a new dependency decision reopens it.
\end{description}
% END migrated B055

% BEGIN migrated B056
\section{PC-independent infrastructure and detailed leaf blueprints}
\label{sec:independent-infrastructure}

A substantial part of Geometrization can be specified without waiting
for changing PC declaration names. The active work develops its detailed
chapter blueprints within this file. Mathlib environment pinning and
Lean implementation are deferred at the user's request. Independence of a mathematical
statement from Ricci flow does not mean that its smooth or metric
foundation is absent from PC or mathlib. The completed-release boundary
for production GC remains unchanged. The first new analytic layer in
the flow chain remains the surviving M3 PDE gaps.

The inventory \code{independent_infrastructure.csv} separates eight
work streams. These are a proposed development order, not measured
workload estimates or claims that libraries are missing.

\begin{longtable}{@{}P{0.07\textwidth}P{0.24\textwidth}P{0.61\textwidth}@{}}
ID & Infrastructure & Concrete GC outputs and next scope \\\hline
\endfirsthead
ID & Infrastructure & Concrete GC outputs and next scope \\\hline
\endhead
I01 & Metric geometry & Quantitative pointed Gromov--Hausdorff
 approximations, finite nets, proper limits, and scale envelopes.
 See MG01--MG09 in Chapter~\ref{ch:metric-geometry}; environment setup is deferred. \\
I02 & Alexandrov geometry & Only the comparison, dimension, strainer
 and splitting results consumed by the selected local-collapse route.
 Identify exact consumers before expanding all proofs. \\
I03 & Hyperbolic geometry & Quotients, finite-volume cusps and
 truncations, peripheral fundamental groups, and complete metric
 transport to the same full interior. Keep the flow-limit bridge separate. \\
I04 & Local collapse & Corrected slowly varying scales, local models,
 strata, good annuli and boundary modifications. Retain H1--H4 when
 applying this static theorem to late flow slices. \\
I05 & Compatible fibrations & Adapted maps, submersions, overlap
 modifications, bundle output, corners and graph presentations.
 Generic smooth calculus remains a reuse candidate. \\
I06 & Three-manifold topology & Existing G/R/F/J targets: finite
 gluing, relative primes, slopes and filling, caps, actual torus maps
 and history reconstruction. Smooth Schoenflies remains inherited. \\
I07 & Model realization & Seifert/orbifold cases, Nil and Sol,
 bundles and semibundles, and the actual complete model metrics required
 by E1. Hyperbolic witnesses also require finite volume. \\
I08 & Cusp proof support & Compare the Section 91 least-area-disk
 route with the alternative cited by Morgan--Tian before commissioning
 a variational library. This is a mixed static/flow bridge. \\
\hline
\end{longtable}

The current next step is to expand the topic chapters with definitions,
precise statements, source locators and proof dependencies. The former
implementation sequence remains a future option; no environment setup
is part of the present work. The global spine and source roles remain
unchanged. This table is an infrastructure index, not a competing set
of physical chapter numbers.
% END migrated B056

% BEGIN migrated B062
\section{Dependency graph}

\begin{center}
\begin{tikzpicture}[node distance=9mm and 13mm, every node/.style={draw,rounded corners,align=center,font=\small}, >=Latex]
  \node (audit) {W0\\audit and conventions};
  \node (top) [below left=of audit] {W1\\3-manifold topology};
  \node (riem) [below right=of audit] {W2\\PC reuse/audit};
  \node (flow) [below=of riem] {W3\\flow adapter};
  \node (per) [below=of flow] {W4\\Perelman adapter};
  \node (surg) [below=of per] {W5\\global surgery};
  \node (long) [below=of surg] {W6\\long-time behavior};
  \node (end) [below=of long] {W7\\integration};
  \draw[->] (audit) -- (top);
  \draw[->] (audit) -- (riem);
  \draw[->] (riem) -- (flow);
  \draw[->] (flow) -- (per);
  \draw[->] (per) -- (surg);
  \draw[->] (surg) -- (long);
  \draw[->] (top) -- (end);
  \draw[->] (long) -- (end);
\end{tikzpicture}
\end{center}

W1 can proceed independently of the analytic chain. W2 can proceed
independently of W1. W7 should be stated as soon as the W1 and W6 contracts
are stable; that is the first vertical slice.
% END migrated B062

% BEGIN migrated B063
\section{M1 contract to freeze before PDE work}
\label{sec:m1-contract}

The next reviewable target is a small assembly theorem, not a partial Ricci
flow proof. Before W3--W6 become large, the project should freeze four
interfaces:
\begin{enumerate}[label=\textbf{I\arabic*:},leftmargin=*]
\item \code{GeometrizationCertificate M}, the existential endpoint
  certificate described in Section~\ref{sec:endpoint-certificate};
\item \code{ThickThinCollapsePackage M}, carrying the finite hyperbolic
  thick pieces, graph-manifold thin pieces, and their boundary/assembly data;
\item \code{PoincareAdapter}, exposing only the imported Poincar\'e
  consequences used by the spherical branch;
\item \code{RicciFlowSurgeryPackage M}, whose analytic fields produce the
  separate \code{LongTimeOutcome} through M6 and preserve the required
  topological input data; the outcome is not a field of the flow package.
\end{enumerate}

The first production assembly target after K1 is the structural implication
\begin{equation}
\begin{gathered}
  \texttt{PoincareAdapter}+\texttt{ThickThinCollapsePackage M}\\
  \Longrightarrow\;\texttt{GeometrizationCertificate M}.
\end{gathered}
\end{equation}
Only after this theorem compiles should the analytic package be filled in.
This ordering makes the final assembly independently reviewable and gives W1
and W6 a stable interface. The theorem inventory records the exact signatures;
the LaTeX blueprint records their mathematical role.
% END migrated B063

% BEGIN migrated B069
\section{Legacy implementation scaffold and Lean module map}
\label{ch:formal-scaffold}

The C0--C13 identifiers below are retained as stable implementation and
inventory identifiers. They are not the physical chapter numbers of
this revised book. Chapter~\ref{ch:topic-guide} and
\code{chapter_framework.csv} give the topic chapters; the C-to-topic
crosswalk preserves existing source and module records. A chapter,
work package, milestone and Lean module need not coincide.

The scaffold freezes boundaries and handoffs now. It does not freeze theorem
wording, exact namespaces, or source locators before the relevant release or
literature comparison is complete. Every chapter shell must eventually contain
seven items: scope and purpose; public declarations; required inputs and
imports; output and handoff package; dependency and parallel-work lanes;
authoritative sources with exact locators; and an acceptance gate.

The status codes in the table prevent placeholders from being mistaken for
proof commitments:
\code{I} means inherited result to audit and reuse; \code{A} means a small adapter or
transport theorem; \code{N} means a genuinely new GC theorem; \code{P} means a
placeholder pending the literature comparison or exact PC statement; \code{G} means
final integration; and \code{X} means optional or outside the mandatory route.

\setlength{\tablebodywidth}{\dimexpr\textwidth-10\tabcolsep\relax}
\begin{longtable}{@{}P{0.062500\tablebodywidth}P{0.229167\tablebodywidth}P{0.135417\tablebodywidth}P{0.260417\tablebodywidth}P{0.197917\tablebodywidth}P{0.114583\tablebodywidth}@{}}
\textbf{ID} & \textbf{Working chapter} & \textbf{Mile\-stone / status} & \textbf{Principal public output} & \textbf{Depend\-encies and parallel lanes} & \textbf{Lean root / gate}\\\hline
\endhead
C0 & Input category, conventions, target theorem, and trust boundary & M0 / I+A & \code{ThreeManifold}, \code{Geometrizes}, convention ledger, \code{PCReleaseGate} & No mathematical predecessor; release and trust policy is shared by every lane. & \code{GC/Basic}, \code{GC/Audit}; M0\\
C1 & Endpoint topology and model-geometry contracts & M1a--M1d / N+P & \code{PrimeDecomposition}, \code{TorusCutSystem}, \code{CutDecomposition}, \code{HasThurstonGeometry}, abstract model tags & Depends on C0. Topology and abstract endpoint contracts can proceed in parallel with C2. & \code{GC/Endpoint/*}; M1-contract gate\\
C2 & Inherited Poincare foundation and M2A adapters & M2/M2A / I+A & \code{PCFoundationAudit}, \code{ConventionLedger}, \code{FoundationReusePackage}, \code{TimeDependentMetric} & Requires the pinned PC release for final classification; convention and provenance lanes run in parallel. & \code{GC/Foundation/*}; \code{PCReleaseGate}\\
C3 & Thick/thin package, structural assembly, and contract-phase vertical slice & M1e--M1f / N+G & \code{ThickThinCollapsePackage}, \code{assemble_geometrization}, producer boundary, mock regression suite & Requires C1 and the audited geometric objects from C2. Assembly and analytic-producer contracts remain separate. & \code{GC/Integration/Structural/*}; M1f gate\\
C4 & Ricci-flow PDE compatibility and genuine gap closure & M3 / I+A+N & \code{RicciFlowPDE}, time-dependent flow equation, solution/evolution interfaces, regularity package & Requires C2. Audit PC short-time and evolution declarations first; only missing interfaces are new. & \code{GC/RicciFlow/PDE/*}; M3 gate\\
C5 & Perelman estimates, monotonicity, noncollapsing, and compactness & M4 / I+A+N & Estimate package, reduced-volume/entropy monotonicity, noncollapsing, curvature control, pointed-flow compactness & Requires C4. Functional audits, compactness audits, and theorem translations can run as parallel lanes. & \code{GC/RicciFlow/Perelman/*}; M4 gate\\
C6 & Surgery geometry and global typed flow-with-surgery history & M5 / I+A & Surgery-event records, canonical-neighborhood interface, continuation, invariant preservation, \code{RicciFlowSurgeryPackage} & Requires C2, C4, C5 and completed PC surgery interfaces. Local surgery audit and global history adapter are parallel after signatures stabilize. & \code{GC/RicciFlow/Surgery/*}; M5 gate\\
C7 & Normalized long-time slices and thick/thin extraction & M6A / N+P & \code{LongTimeSliceData}, \code{ScaleNormalization}, regular-time selection, \code{LongTimeOutcome}, \code{LongTimeThickThin} & Requires C5 and C6. Slice/scale and surgery-history transport are separate lanes; outcome must allow late slices or terminal exceptional data. & \code{GC/LongTime/Slice/*}; M6A gate\\
C8 & Thick branch: hyperbolic limits and finite-volume pieces & M6B / N+P & \code{HyperbolicThickOutput}, \code{HyperbolicLimitToPiece}, cusp/core and boundary data & Requires C1 and C7. Parallel with C9 after the common scale and convergence interface is fixed. & \code{GC/LongTime/Hyperbolic/*}; M6B gate\\
C9 & Thin branch: collapse and graph-manifold pieces & M6C--M6D / N+P & \code{CollapseHypotheses}, \code{CollapsedThinOutput}, \code{CollapseToGraphManifold}, Seifert/graph data & Local theorem uses C1 geometry; applying it to flow uses C7. PC-independent contract work is active. Parallel with C8; KL14 local route and MT14 global application remain provisional. & \code{GC/LongTime/Collapse/*}; M6C gate\\
C10 & Exceptional branches, model realization, transport, and reconstruction & M6E--M6F/M7 / A+N+G & \code{ExceptionalPrimeToGeometricPiece}, \code{ThinGraphToGeometricPieces}, global torus system, reconstruction and certificate packaging & Requires C3, C7, C8, C9, and the exact \code{PoincareAdapter} surface. Exceptional\slash spherical recognition is an explicit subproblem. & \code{GC/Integration/*}; M7 gate\\
C11 & Final assembly and public geometrization theorem & M7 / G & \code{geometrization_of_ricci_flow_surgery}, \code{geometrization}, final \code{GeometrizationCertificate} & Requires C10 and all public package interfaces. This is a short composition, not a second analytic proof. & \code{GC/Final/*}; release gate\\
C12 & Literature\slash source map and theorem-locator policy & Cross-cutting / P & source-role matrix, exact locators, route decisions, errata record & Runs alongside C1--C10; a source is promoted only when its Lean contract and locator are fixed. & \code{GC/References/*}; source gate\\
C13 & Provenance, regression, CI, and final axiom audit & Cross-cutting / I+A+G & import report, no-\code{sorry} report, transitive \code{#print axioms}, model smoke tests, release record & Runs throughout; final sign-off depends on C11 and every promoted theorem. & \code{GC/Audit/*}; M7\slash prove\-nance gate\\
\end{longtable}

\subsubsection{Source roles attached to the scaffold}

The source list is attached by mathematical role rather than by chapter
chronology. Exact theorem, section, and page locators remain provisional until
the literature comparison and the PC release audit are complete.

\setlength{\tablebodywidth}{\dimexpr\textwidth-4\tabcolsep\relax}
\begin{longtable}{@{}P{0.222222\tablebodywidth}P{0.388889\tablebodywidth}P{0.388889\tablebodywidth}@{}}
\textbf{Scaffold chapters} & \textbf{Primary authorities} & \textbf{Use in the formalization}\\\hline
\endhead
C0--C3 & The pinned PC release; Hatcher's 3-manifold notes; Scott; Thurston; Morgan--Tian 2014 & Define the endpoint vocabulary, model tags, adapter boundary, and structural assembly contracts.\\
C4--C5 & Hamilton 1982 and 1995; Perelman 2002; Kleiner--Lott 2008; the completed PC analytic modules & Audit or import the flow equation, evolution/PDE interfaces, entropy and reduced-volume identities, noncollapsing, curvature control, and compactness.\\
C6 & Perelman 2003 surgery; Kleiner--Lott 2008; Morgan--Tian 2007 & Specify the public typed surgery history, canonical-neighborhood input, continuation, and invariant-preservation adapters.\\
C7 & Perelman 2003 surgery; Kleiner--Lott 2008; Morgan--Tian 2014 & Fix the late-slice quantifiers, common rescaling, regular-time selection, finite stabilization, and the terminal-outcome disjunction.\\
C8 & Hamilton 1999; Perelman 2002; MSM206 Chapters 31--34 and Appendix K; Benedetti--Petronio & Convert thick pointed limits into finite-volume hyperbolic-region certificates and normalized model witnesses.\\
C9 & Kleiner--Lott, Ast\'erisque 365; Morgan--Tian 2014; Shioya--Yamaguchi; Cheeger--Fukaya--Gromov when needed & Select and formalize one local-collapse route, then convert its boundary-controlled output into graph-manifold data.\\
C10--C13 & Morgan--Tian 2014; Scott; the pinned PC theorem and project trust policy & Transport labels and incompressibility, realize all model geometries, assemble the final certificate, and audit provenance.\\
\end{longtable}

C1 is intentionally abstract about the metric realization of a model tag.
C2 supplies the inherited metric, tensor, connection, and curvature APIs that
C3 and the later geometric-witness chapters consume. Thus the endpoint can be
designed before the PC release, while any implementation of concrete metric
witnesses waits for the M2A convention audit.

C8 and C9 are parallel M6 branches, not a serial hyperbolic-then-collapse
argument. Both consume the common normalized slice and convergence interfaces
from C7. C10 contains the conversion of graph and Seifert output to the eight
model geometries, the exceptional-prime branch, and the transport of all
labels and incompressibility proofs back to the original manifold.

The global long-time interface in C7 is deliberately a disjunction:
\begin{equation}
\begin{aligned}
 &\texttt{LongTimeOutcome M} := \\
 &\quad\texttt{LateThickThinOutcome M} \\
 &\quad\mathbin{|}\;\texttt{TerminalExceptionalOutcome M}.
\end{aligned}
\end{equation}
The first branch supplies arbitrarily late regular-time data and is the source
of C8 and C9. The second carries explicit terminal or exceptional-prime data
that C10 can convert using audited topology and the precise \code{PoincareAdapter}.
This keeps finite-time extinction outside the mandatory route without silently
assuming that every flow has a late slice.

Revision 36 keeps this production scaffold unchanged. Its independent K0
implementation lane is specified in Section~\ref{sec:k0-implementation},
and M2A promotion is controlled by Section~\ref{sec:m2a-release-workflow}.

\subsubsection{Chapter-shell completion rule}

A chapter is not complete because its title, source list, or imports exist. Its
shell becomes implementation-ready only when the public declarations, all
hypotheses and normalizations, downstream consumers, source locators, and
acceptance tests are recorded. Imported chapters must additionally pass
\code{PCReleaseGate}; adapter chapters must show transport and axiom provenance;
new chapters must have no temporary \code{axiom}, \code{sorry}, or \code{admit} in the release
target. The theorem inventory and the corresponding API CSV remain the
machine-readable authority for these fields.

\subsubsection{Scaffold expansion order}

The active expansion is the topic-based mathematical blueprint.
Environment pinning and Lean coding are deferred by the user. C2's
completed-release comparison and the C3--C11 implementation order remain
future gated work; they do not prevent detailed statements and source
comparisons in the independent topic chapters now.
% END migrated B069

\chapter{Proof readiness, K0 evidence and contract audits}
\label{ch:readiness-appendix}

% BEGIN migrated B004
\section{Proof construction readiness and order}
\label{ch:proof-readiness}

The project should begin proving pieces before the entire geometrization
argument is formalized, but it must distinguish three kinds of work. A
contract proof checks that an interface is mathematically well formed. An
inherited proof check establishes that the completed PC release already proves
a result and that its hypotheses and conventions transport to GC. A new GC
proof establishes a theorem that is genuinely absent from the inherited
release. These categories have different prerequisites and must not be
reported under one undifferentiated ``proof progress'' label.
% END migrated B004

% BEGIN migrated B005
\section{What must be complete before substantive GC proofs}

The following prerequisites are required before a proof is promoted into the
GC release target:

\begin{enumerate}[label=\textbf{R\arabic*:},leftmargin=*]
\item the endpoint certificate, producer boundary, and dependency graph are
  frozen sufficiently to rule out circular definitions;
\item \code{PCReleaseGate} identifies an immutable, cleanly building completed PC
  release, its Lean toolchain, public declarations, and transitive axiom
  policy;
\item M2A has classified every inherited foundation row, frozen the
  convention ledger, and proved the required category, scaling, and notation
  adapters;
\item the selected literature route has an explicit theorem contract for the
  hyperbolic branch and for the collapse branch, including all quantifiers,
  regularity, boundary, and normalization hypotheses;
\item the Lean module skeleton, import probes, regression examples, theorem
  inventory, and acceptance-test harness exist; and
\item every proof target has a named downstream consumer and a status that
  distinguishes inherited reuse, adapter transport, new GC mathematics, and
  unresolved placeholder.
\end{enumerate}

These prerequisites do not mean that no Lean code can be written earlier.
They mean that early code is contract or sandbox code until the corresponding
gate is passed. A theorem is not promoted merely because a provisional
interface can be made to type-check.
% END migrated B005

% BEGIN migrated B006
\section{Proof-start gates}

\setlength{\tablebodywidth}{\dimexpr\textwidth-6\tabcolsep\relax}
\begin{longtable}{@{}P{0.155556\tablebodywidth}P{0.300000\tablebodywidth}P{0.322222\tablebodywidth}P{0.222222\tablebodywidth}@{}}
\textbf{Gate} & \textbf{What it permits} & \textbf{Required evidence} & \textbf{What remains blocked}\\\hline
\endhead
K0: contract sandbox & Endpoint well-formedness lemmas, mock packages, negative tests, and abstract assembly signatures. & C1/C3 contracts, no circular dependency, mock acceptance suite. & Concrete reuse claims, analytic proofs, and final assembly in the release target.\\
K1: inherited foundation & Concrete M1e/M1f assembly proofs and M3 interface construction. & \code{PCReleaseGate}, completed M2A inventory, convention tests, adapter theorem and axiom reports. & Any theorem relying on an unclassified PC declaration or unfrozen sign/scale convention.\\
K2: flow interface & M4 estimates and compactness proofs or audited reuse. & Ricci-flow equation, time-dependent metric, evolution, regularity, and solution interfaces with exact hypotheses. & Long-time conclusions whose flow and compactness inputs are still schematic.\\
K3: estimates and surgery & M6 slice extraction and branch theorem construction. & Audited Perelman estimates, compactness, canonical neighborhoods, typed surgery history, continuation, and invariant transport. & Hyperbolic\slash collapse output with hidden analytic or surgery assumptions.\\
K4: branch closure & Final transport, reconstruction, and public geometrization theorem. & Hyperbolic and collapse branches, terminal branch, coverage, gluing, and all provenance reports. & Release publication until the final axiom and no-gap audit passes.\\
\end{longtable}

K0 is available now for contract work. K1 is the first gate for substantive
GC proofs. K2--K4 are successive promotion gates; an imported theorem may let a
later gate consume an interface earlier, but only after the import and axiom
audit records that fact.
% END migrated B006

% BEGIN migrated B007
\section{First proof tranche}
\label{sec:first-proof-tranche}

The first tranche is intentionally small and independent of Ricci-flow
analysis. K0 uses namespace \code{GC.Sandbox.K0}, with field roles aligned
to production contracts but separate names and imports:

\begin{enumerate}[label=\textbf{T\arabic*:},leftmargin=*]
\item \code{EndpointContractWellFormed}: prove finite-index, coverage,
  boundary, and tag consistency for the sandbox endpoint; concrete manifold
  and metric witnesses remain release-gated;
\item \code{MockThickThinPackage}: construct a finite mock package containing one
  thick region, one collapsed region, explicit boundary labels, and coverage;
\item \code{MockAssemblyRegression}: run sandbox \code{assembleMock} and
  refute explicit invalid coverage, matching, and injectivity fixtures;
  production \code{assemble_geometrization} is not a dependency;
\item \code{M2AConventionRegression}: after \code{PCReleaseGate}, prove the standard
  sphere/flat/hyperbolic convention and scaling tests and record the adapter
  axiom report; and
\item \code{assemble_geometrization}: after K1, replace the mock inputs by the
  audited \code{ThickThinCollapsePackage} and prove the structural endpoint
  theorem with the explicit \code{PoincareAdapter} argument.
\end{enumerate}

T1--T3 can be developed before the completed PC release because they use
finite contracts. T4 runs after \code{PCReleaseGate}, during M2A, and is
evidence required to open K1; it must not itself wait for K1. T5 begins after
K1. The checked boundary lemma in Section~\ref{sec:k0-implementation}
precedes T1. None of these steps proves the Ricci-flow producer.
% END migrated B007

% BEGIN migrated B008
\section{Substantive proof order}

M2A adapters supply K1 evidence. After they pass, substantive proof work
proceeds through M1e--M1f and M3 onward:

\begin{description}[leftmargin=3.2cm,style=nextline]
\item[M2A adapters] Prove only the inherited category, notation, curvature,
  scaling, time-family, and convention translations that the audit identifies.
\item[M1e--M1f assembly] Prove coverage, boundary matching, model-witness
  conversion, exceptional-branch handling, and \code{assemble_geometrization}.
\item[M3 flow compatibility] Audit or reuse short-time existence, the
  Ricci-flow equation, evolution identities, maximum principles, and PDE
  regularity. Prove only the gaps left by the PC release.
\item[M4 estimates and compactness] Audit or reuse entropy, reduced volume,
  monotonicity, noncollapsing, derivative estimates, curvature control, and
  pointed-flow compactness. Add missing transport lemmas only after the M3
  hypotheses are exact.
\item[M5 surgery] Adapt canonical-neighborhood, surgery-event, continuation,
  topology, and invariant-preservation results into the typed global history.
\item[M6 long time] Prove common-scale slice extraction, the late-slice or
  terminal-outcome disjunction, the hyperbolic branch, and the collapse branch.
  The two geometric branches are parallel after their shared C7 interface.
\item[M7 closure] Prove transport and reconstruction, pattern-match both
  long-time branches, package the final certificate, and run the complete
  provenance audit.
\end{description}

A proof in a later item may be developed as a statement-level placeholder while
an earlier item is incomplete, but it may not enter the release target with an
unreviewed \code{axiom}, \code{sorry}, or \code{admit}. The placeholder must identify the
missing upstream theorem and its intended consumer.
% END migrated B008

% BEGIN migrated B009
\section{Acceptance record for each proof piece}

Every promoted proof piece receives a record containing its exact declaration
and import, all hypotheses and normalizations, its theorem-inventory ID,
source locator or inherited-release evidence, downstream consumer, positive
and negative tests, and transitive \code{#print axioms} output. The record also
states whether the proof is direct reuse, an adapter, genuinely new GC
mathematics, or a structural integration theorem.

This record prevents a successful local proof from masking a convention
mismatch or an accidental dependency on a project axiom. It also lets several
people work independently: a contributor can begin when the inputs and gate
for that proof piece are complete, without waiting for unrelated downstream
chapters.
% END migrated B009

% BEGIN migrated B010
\section{Current action}

The work remains split into two lanes. The sandbox lane now maintains the
checked finite T1--T3 suite as permanent regression tests. The release lane
maintains the prepared M2A bundle in Section~\ref{sec:m2a-prepared-bundle},
fills final inventory decisions, and freezes convention and axiom reports
when the completed PC release passes its gate. The first
substantive GC theorem to promote after K1 is \code{assemble_geometrization}; the
first genuinely new analytic work after that is the M3 Ricci-flow PDE gap
closure. The boundary proof recorded below checks only finite bookkeeping.
% END migrated B010

% BEGIN migrated B011
\section{Concrete K0 implementation targets}
\label{sec:k0-implementation}

K0 is a finite-contract sandbox in namespace \code{GC.Sandbox.K0}. It does
not instantiate production manifold wrappers, geometric models, or the PC
adapter. All seven finite-contract modules are implemented and checked in
this revision. Their 17 implementation targets and one separate source-review
target are recorded in \code{k0_sandbox_api.csv}.
Paths below are relative to \code{sandbox/GC/Sandbox/K0/}.

\setlength{\tablebodywidth}{\dimexpr\textwidth-4\tabcolsep\relax}
\begin{longtable}{@{}P{0.220930\tablebodywidth}P{0.406977\tablebodywidth}P{0.372093\tablebodywidth}@{}}
\textbf{File / order} & \textbf{Output} & \textbf{Imports and acceptance}\\\hline
\endhead
\code{Boundary.lean} / 1 & \code{BoundaryMatching}, partner injectivity, empty/two-side fixtures, singleton and identity refutations. & \code{Init} only; standalone compiler and axiom evidence recorded.\\
\code{Raw.lean} / 2 & \code{RawPackage}, \code{PackageValid}, \code{MockBranchInputs}, \code{MockThickThinPackage}. & Boundary; exact finite fields and validity independent of the endpoint.\\
\code{Endpoint.lean} / 3 & \code{MockEndpointCertificate}, \code{EndpointContractWellFormed}. & Raw; local output consistency without a production certificate.\\
\code{Assembly.lean} / 4 & \code{assembleMock}. & Endpoint; construct original-label incidence through inverse transport.\\
\code{Fixtures.lean} / 5 & \code{mixedPackage}, \code{MockAssemblyRegression}. & Assembly; two regions, paired boundary, nonidentity transport, exact tags.\\
\code{Negative.lean} / 6 & Coverage, orientation, injectivity, and transport refutations. & Fixtures; kernel proofs of negated validity predicates.\\
\code{Audit.lean} / 7 & Declaration checks and transitive axiom commands. & Negative; checked full finite K0 suite plus recorded source/import review.\\
\end{longtable}

Production modules must never import sandbox modules, even transitively.
K0 has no imports from PC, production GC, Ricci flow, surgery, or extinction.
Its toolchain pin is independent of the eventual GC toolchain, which must
match the completed PC release. There is no Lake project in this workspace.

\subsection{First small proof and actual evidence}

For $n\in\N$, \code{BoundaryMatching n} stores
$p:\operatorname{Fin}(n)\to\operatorname{Fin}(n)$ and proofs
\[
  \forall s,\ p(p(s))=s,\qquad \forall s,\ p(s)\ne s.
\]
The first theorem, \code{BoundaryMatching.partner_injective}, proves
\[
  p(s)=p(t)\ \Longrightarrow\ s=t.
\]
Apply $p$ to the equality and use the involution law twice. This prevents
two labels from sharing a partner. It does not imply the no-fixed-point
condition, which remains a separate mandatory field.

The positive fixtures are an empty matching on $\operatorname{Fin}(0)$
and a swap on $\operatorname{Fin}(2)$. The negative fixtures prove that
$\operatorname{Fin}(1)$ admits no matching and that identity pairing on
$\operatorname{Fin}(2)$ fails the no-fixed-point predicate. These are proved
refutations, not failing source files or incomplete constructors.

Run \code{python3 GEOMETRIZATION_BLUEPRINT/sandbox/check_boundary.py} from
the worktree root. It requires an already installed
\code{leanprover/lean4:v4.33.1}, invokes Lean on the standalone file, and
records the source hash, command, compiler version, exit code, and raw
axiom output in \code{sandbox/evidence/boundary_check.json}.
It does not install a toolchain or build a GC/PC project.

The run succeeded. Partner injectivity and the empty fixture have empty
axiom closures; the other four audited declarations report only
\code{propext}. This is observed sandbox evidence, not approval of the
future PC axiom policy. This is the retained revision-35 boundary evidence.
Revision 36 additionally checks the full finite K0 suite described below.
Production M1, M2A, and K1 remain incomplete. The static blueprint audit
checks stored evidence against source hashes; it does not rerun Lean.

\subsection{Implemented finite raw data and validity}

Use natural-number cardinalities and \code{Fin} indices for prime labels
$P$, region labels $R$, torus labels $T$, late labels $L$, and original
labels $O$. These are bookkeeping tokens, not manifolds. \code{RawPackage}
stores maps $R\to P$ and $O\to P$, a thick/thin/exceptional kind for each
region, Boolean incidence $c:R\times L\to\mathrm{Bool}$, and maps
$u:L\to O$ and $v:O\to L$. Its boundary sides are $T\times\mathrm{Bool}$,
with region owners, a partner map, and orientation bits. For each $t\in T$,
store a probe map $q_t:\operatorname{Fin}(2)\to\operatorname{Fin}(2)$.

Define \code{PackageValid d} before any endpoint. Its conjuncts are:
$u$ and $v$ are mutual inverses; every original label is $u(x)$ for some
incident pair $(r,x)$; each late label has at most one incident region;
and prime labels agree along occupied cells. Side pairing preserves the
torus label, reverses the side bit, is involutive and fixed-point-free,
pairs opposite orientation bits, and pairs owners in the same prime.
Every probe $q_t$ is injective.

The probe tests how an injectivity hypothesis is required and transported;
it is not torus incompressibility. Orientation bits are not smooth boundary
orientations. A constant probe on two elements must fail. Replacing its
obligation by \code{True} would invalidate the regression.

\code{MockThickThinPackage} stores raw data and a \code{PackageValid} proof.
It has no field, existential, oracle, or function returning an endpoint.
\code{MockBranchInputs d} separately supplies mock model tags for thin and
exceptional labels; thick labels receive the hyperbolic tag. Exceptional
inputs are scoped to the fixture, not a universal \code{PoincareAdapter}.
No tag is a metric or a geometric recognition theorem.

\code{MockEndpointCertificate} stores finite original-label incidence,
prime and side labels, assigned tags, and the corresponding coverage,
uniqueness, prime-compatibility, matching, orientation, and probe proofs.
\code{assembleMock} takes a validated raw package and its explicit
\code{MockBranchInputs}, and constructs
\[
 c_{\mathrm{out}}(r,o):=c(r,v(o)).
\]
Prove coverage from raw coverage and $v(u(x))=x$, and prime compatibility
from $u(v(o))=o$. Transfer boundary/probe proofs and assign the tags.
\code{EndpointContractWellFormed} exposes the output consistency predicates.
There is no conversion to production \code{GeometrizationCertificate}.

\subsection{First mixed fixture and discriminating failures}

Take one prime, two regions, one torus, and two late/original labels.
Region 0 is thick and region 1 thin. Incidence is equality of two-element
indices; $u=v$ is the swap. Negative and positive sides belong to regions
0 and 1 respectively; pairing flips the side bit, orientation is that
bit, and the probe is identity. Give the thin label the Euclidean mock tag.
The result must assign original label 0 to region 1 and label 1 to region 0,
with one hyperbolic and one Euclidean tag. This asserts no realizable
geometric decomposition. Nonidentity transport detects accidental copying
of late incidence directly into the output.

The coverage, orientation, probe, and transport tests each change a raw
datum, exhibit a counterexample to its validity conjunct, and prove that
the mutated data cannot satisfy \code{PackageValid}. The independent
Boundary checks refute singleton matching and identity pairing directly:

\setlength{\tablebodywidth}{\dimexpr\textwidth-4\tabcolsep\relax}
\begin{longtable}{@{}P{0.279070\tablebodywidth}P{0.395349\tablebodywidth}P{0.325581\tablebodywidth}@{}}
\textbf{Inventory target} & \textbf{Mutation} & \textbf{Counterexample}\\\hline
\endhead
\code{K0NegativeCoverage} & Remove region 1's occupied cell. & Original label 0 is uncovered.\\
\code{K0NegativeMatching} & One remaining side, or identity pairing. & Singleton impossibility or fixed point; already checked.\\
\code{K0NegativeOrientation} & Give both sides the same orientation bit. & Matched bits are not opposite.\\
\code{K0NegativeInjectivity} & Use the constant-zero probe. & Inputs 0 and 1 have equal images.\\
\code{K0NegativeTransport} & Make $u$ constant zero, retaining $v$. & Inverse law fails at original label 1.\\
\end{longtable}

Missing constructor fields alone are not negative-test evidence.
Circularity instead requires a source/import and transitive declaration
dependency review. A name search cannot detect an alias or opaque oracle.
Record \code{K0ImportBoundaryAudit} as a review, not a mathematical axiom.

Boundary, Raw, Endpoint, Assembly, Fixtures, Negative, and Audit now compile
in that order. The 17 implementation rows have status
\code{kernel-checked-standalone}; the source-review row has status
\code{reviewed-source-audit}. The finite T1--T3 tranche is complete.
None of this opens K1 or proves a production M1 contract. After K1, prove
the geometric consumers against inherited types; do not rename finite mock
types into production types. The examples in \code{m1_mock_api.csv}
remain separate K1-or-later targets with audited model prerequisites.

\subsection{Revision 36 clean-build and review evidence}

Run \code{python3 GEOMETRIZATION_BLUEPRINT/sandbox/check_k0.py}.
The runner requires the already installed Lean 4.33.1 toolchain, rebuilds
all seven local modules in a fresh temporary directory, and restricts
\code{LEAN_PATH} to that directory. It imports only the checked local
predecessors and Lean core; no old workspace object files are used.
The temporary object files are removed after their hashes are recorded.

The run succeeded for all seven modules. The report
\code{sandbox/evidence/k0_check.json} records 35 distinct declaration
axiom reports, exact source hashes, compiler version, commands, exit codes,
output-object hashes, and raw compiler output. The structural
\code{assembleMock}, \code{EndpointContractWellFormed}, transport equation,
and all three tag-preservation lemmas have empty axiom closures.
The finite fixtures and refutations use at most \code{propext} and
\code{Quot.sound}; no project assumption or unfinished proof occurs in their
reported closures. This observed sandbox policy does not approve a PC release.

The main regression checks all cardinalities, both transported incidence
entries and both excluded entries, the hyperbolic/Euclidean tags, side
pairing, and side owners. \code{copied_incidence_rejected} additionally
shows that copying the original late incidence directly into the endpoint
would give the wrong labelled answer. The negative tests are kernel-checked
counterexamples and refutations, not tests that merely omit structure fields.

The separate \code{sandbox/evidence/k0_source_review.json} records review of
the seven sources at their hashes: raw fields are concrete finite data,
validity predicates mention only those fields, branch inputs supply only
fixture-scoped tags, and the input package contains no endpoint-returning
oracle. This is a source review, not a Lean theorem. The static audit
rejects stale build or review evidence, unexpected imports, mismatched
inventory status, and unsupported M2A promotion.

The remaining work is release-audit preparation and exact source-contract
comparison. Concrete geometric M1 implementation and substantive M3/M6 proofs
still wait for their stated release gates.
% END migrated B011

% BEGIN migrated B024
\section{M1 mock producer and integration-test specification}
\label{sec:m1-mock-producer}

K0 supplies a finite-contract target before the completed PC release.
The geometric M1 suite here follows K1 and its required model inputs.
Its mock packages satisfy M1e by proved standard-model constructions. These mocks test the
assembly interfaces; they are not claims about Ricci-flow output and must not
be imported by the final geometrization theorem.

\subsection{Mock-producer boundary}

The mock layer has one sandbox-only interface:
\begin{equation}
\begin{aligned}
 &\texttt{M1MockExample} \\
 &\quad\longrightarrow \texttt{PoincareAdapter}
  \times
  \texttt{ThickThinCollapsePackage M}.
\end{aligned}
\end{equation}

The example supplies a concrete $M$, an explicit M1e package, and the adapter
consequences needed by its exceptional factors. It does not provide a
\code{RicciFlowSurgeryPackage} and it does not assert that a flow produced the
package. The production theorem remains
\code{long_time_outcome}, with both late and terminal branches.

Every test adapter must prove all consequences its type quantifies over.
Standard-model fixtures alone do not supply a universal \code{PoincareAdapter};
restrict inputs to fixture-specific exceptional labels or use the audited
real adapter after its prerequisites pass. It may be used in a sandbox test module, but it is not a
replacement for the pinned completed PC release and is never an axiom in the
release target.

\subsection{Positive example suite}

The first suite contains five finite examples:

\setlength{\tablebodywidth}{\dimexpr\textwidth-4\tabcolsep\relax}
\begin{longtable}{@{}P{0.255814\tablebodywidth}P{0.441860\tablebodywidth}P{0.302326\tablebodywidth}@{}}
\textbf{Example} & \textbf{Mock package} & \textbf{Expected certificate}\\\hline
\endhead
Spherical & One closed spherical prime with no torus system; adapter supplies
the spherical witness. & One spherical geometric piece and empty boundary.\\
Flat & A closed flat 3-torus represented by one exceptional/model region with
no thick or thin boundary. & One Euclidean geometric piece.\\
Product & $\Sph^2\times\Sph^1$ with the product model and no torus cut. & One
$\Sph^2\times\R$ geometric piece.\\
Hyperbolic & One closed finite-volume hyperbolic thick region with empty
boundary. & One hyperbolic geometric piece.\\
Graph/Sol & A torus-bundle or Seifert mock thin region with its finite
refinement and side matching. & Finite Euclidean, Nil, Sol, or product piece
list according to the selected model.\\
\end{longtable}

The examples exercise every assembly branch except a nonempty hyperbolic cusp
boundary. A separate cusped-core example should be added once the M1d end-data
API is implemented; it need not wait for the Ricci-flow limit theorem.

\subsection{Stage-by-stage test obligations}

Each positive mock example runs the seven M1f stages separately:

\begin{enumerate}[label=\textbf{M\arabic*:},leftmargin=*]
\item the thick conversion returns the expected hyperbolic witness when a
  thick region is present;
\item the thin conversion returns the expected finite model-piece list;
\item the exceptional conversion consumes only the test adapter;
\item the combined torus system has the expected finite index and no duplicate
  or intersecting images;
\item the global cut decomposition has exactly the expected component and
  boundary-side counts;
\item reconstruction composes with the stored transport map to recover the
  chosen model of $M$; and
\item certificate packaging preserves every model tag, orientation, and
  reconstruction field.
\end{enumerate}

These tests are stronger than checking only that
\code{GeometrizationCertificate M} is inhabited. They verify that the output fields
are populated by the intended preceding stage and that transport did not
silently discard the region labels.

\subsection{Negative test suite}

The mock module also contains deliberately unconstructible variants. They
are kept in a test namespace so that the release build can check the expected
failure without adding an axiom:

\begin{itemize}[leftmargin=*]
\item remove one region from a package and verify that \code{PackageCoverage} is
  unavailable;
\item reverse one boundary orientation and verify that
  \code{PackageBoundaryCompatibility} is unavailable;
\item remove one side from the matching involution and verify that
  \code{CutPieceBoundarySides} or reconstruction cannot be proved;
\item replace a complete metric by an incomplete one for any tag, or use
  an infinite-volume hyperbolic candidate, and verify rejection; a complete
  infinite-volume Euclidean witness is permitted once its other fields
  are proved; and
\item attempt to add a \code{GeometrizationCertificate} field to a mock package and
  verify that the \code{PackageNonCircular} audit rejects the structure.
\end{itemize}

Lean failure tests should use the project's standard compile-test mechanism,
not a permanent \code{sorry}, \code{axiom}, or trusted metaprogram. If a negative test
cannot be expressed without weakening the type, that is evidence that the
corresponding M1 contract is too weak.

\subsection{Mock build target and axiom audit}

The sandbox target imports only the M1 topology, geometry, package, and test
adapter modules. Its root build must succeed without importing any
Ricci-flow, surgery, Perelman, or finite-extinction module. The test report
records:

\begin{enumerate}[label=\textbf{Q\arabic*:},leftmargin=*]
\item the exact example declaration and target manifold;
\item the model tags present in the resulting certificate;
\item the number of prime factors, tori, cut pieces, and boundary sides;
\item the result of each M1f stage test;
\item the negative-test diagnostics; and
\item \code{#print axioms} output for \code{assemble_geometrization} and
  \code{geometrizes_of_package}.
\end{enumerate}

M1 is complete only when the positive suite and the negative suite both pass.
The mock package then becomes a regression test for every later change to the
analytic producer boundary.

\subsection{Handoff to PCReuseAudit and analysis}

After the mock target passes, no further endpoint redesign is needed for the
analytic phase. The next task is to replace the test adapter with the pinned PC
adapter, fill the PCReuseAudit rows, and implement the first genuine producer
lemma. The mock producer remains in the repository as a permanent structural
regression suite.

The mock-example and test audit is recorded in \code{m1_mock_api.csv}.
% END migrated B024

% BEGIN migrated B027
\section{Pre-M2 audit report and corrections}
\label{sec:pre-m2-audit}

Before starting M2A, the canonical blueprint and its inventories were audited
for signature drift, circular dependencies, and data-integrity failures. The
audit was first performed on \code{master28.tex} and the corrections were carried
through the historical revisions into this document.

\subsection{Audit results}

The static checks passed: all tracked LaTeX environments are balanced, all
labels are unique, all bibliography citations have inventory keys, every CSV
has a consistent column count, and theorem identifiers are unique. These
static checks do not invoke Lean or TeX. Revisions 35--36 separately record
the standalone K0 compiler checks in Section~\ref{sec:k0-implementation}.
Revision 37 also records an actual standalone PDF build with Tectonic 0.15.0;
the production GC and PC projects have not been built here.

The following semantic findings required correction:

\setlength{\tablebodywidth}{\dimexpr\textwidth-4\tabcolsep\relax}
\begin{longtable}{@{}P{0.188235\tablebodywidth}P{0.423529\tablebodywidth}P{0.388235\tablebodywidth}@{}}
\textbf{Severity} & \textbf{Finding in prior revision(s)} & \textbf{Correction in this revision}\\\hline
\endhead
Critical & Early top-down formulas and the schematic Lean display omitted the
explicit \code{PoincareAdapter} argument from \code{assemble_geometrization}. & All
assembly signatures now show \code{PoincareAdapter} explicitly.\\
Critical & The theorem inventory listed \code{GeometrizationCertificate} as a
dependency of \code{ThickThinCollapsePackage}, contradicting \code{PackageNonCircular}.
& The package now depends only on prime, late-slice, raw thick/thin, gluing,
and torus interfaces; it contains no final certificate.\\
Critical & Raw \code{HyperbolicRegionCertificate} and
\code{GraphManifoldRegionCertificate} rows depended on the final
\code{GeometricPieceCertificate}. & Those dependencies are removed; conversion to
the final witness is performed by named M1e lemmas.\\
Critical & \code{PoincareAdapter} mixed the M1 endpoint consequences with analytic
PC theorem handles. & The endpoint adapter remains narrow; analytic handles
are separated as \code{PCRicciFlowHandles}.\\
Major & \code{CutPiece} was assigned a global complement-coverage field even though
coverage belongs to \code{CutDecomposition}. & The individual piece now stores only
its component relation and labelled boundary incidence; global coverage stays
in the decomposition/package.\\
Major & The inventory statement and an early integration display for
\code{geometrization_of_ricci_flow_package} omitted the adapter argument. & All
integration signatures now show \code{PoincareAdapter} explicitly.\\
Major & The inventory statement for \code{geometrization_of_ricci_flow_surgery}
also omitted the adapter while listing it as a dependency. & Its statement
now explicitly takes \code{PoincareAdapter} together with
\code{RicciFlowSurgeryPackage M}.\\
Non-blocking & Several inventory dependencies are provisional aliases such as
\code{ThreeManifoldTopology}, \code{RiemannianCore}, and \code{FlowCompactness}. & They remain
explicitly provisional and must be replaced by exact module declarations
during the Lean API audit.\\
Non-blocking & Core reference rows still contain locators such as ``sections to
be fixed.'' & Exact theorem/page locators are required before the corresponding
M2 or M6 theorem is marked ready; this does not block the structural M1 slice.\\
\end{longtable}

\subsection{Signature consistency after correction}

The canonical dependency chain is now:
\[
\begin{aligned}
 &\texttt{PoincareAdapter}\\
 &\quad\longrightarrow\texttt{RicciFlowSurgeryPackage M}\\
 &\quad\longrightarrow\texttt{LongTimeOutcome M}\\
 &\quad\longrightarrow\texttt{GeometrizationCertificate M}.
\end{aligned}
\]
This is a curried composition signature: supply the endpoint adapter
separately, run the analytic producer, and handle the resulting branch.
Only the late branch passes through \code{ThickThinCollapsePackage};
terminal handling uses explicit exceptional data and reconstruction. The final public predicate is obtained by packaging the certificate
as \code{Geometrizes M}.

The following invariants are now mandatory for every later revision:

\begin{itemize}[leftmargin=*]
\item no package or analytic producer structure may contain a final
  \code{GeometrizationCertificate} field;
\item every theorem signature involving assembly must display the explicit
  endpoint adapter argument;
\item raw analytic region certificates must precede, rather than depend on,
  final geometric-piece certificates;
\item endpoint and analytic PC interfaces remain separate; and
\item all provisional inventory aliases and literature locators must be
  resolved before the associated theorem is promoted to a release gate.
\end{itemize}

The complete finding ledger is recorded in \code{pre_m2_audit.csv}. The
release-workflow findings are closed. Revision 46 adds open major finding
A25 on endpoint compatibility. Revision 47 closes this specification
mismatch by adopting E1 consistently; geometric proofs and production
acceptance remain pending. A source-level audit success is not a claim that
all mathematical obligations are closed.
The remaining provisional aliases and literature locators are M2A audit work;
they must be resolved before any dependent theorem is promoted to a release
gate, but they do not justify rebuilding the inherited foundation.
% END migrated B027

% BEGIN migrated B028
\section{Conventions to freeze at M0}

\begin{itemize}[leftmargin=*]
\item closed, connected, orientable manifolds as the public input;
\item smooth structures and the precise relationship between smooth and
  topological manifolds;
\item curvature, Laplacian, volume, and Ricci-flow sign conventions;
\item normalizations for rescaling, reduced distance, and entropy;
\item definitions of incompressibility, graph manifold, Seifert piece, and
  complete geometric piece and its conditional hyperbolic volume field;
\item the exact form of the eight geometries and the allowed boundary/gluing
  data;
\item the Lean and mathlib revisions and the imported Poincar\'e theorem.
\end{itemize}

\lean{The Poincar\'e adapter exposes only the consequences needed here: the
simply connected closed 3-manifold conclusion and the spherical-piece
recognition used by the selected route. Finite extinction is not part of this
adapter. Its source theorem and \code{#print axioms} output are recorded in the
project ledger.}
% END migrated B028

% BEGIN migrated B029
\section{Revision 55 audit of the main contracts}
\label{sec:revision55-audit}

The pre-step audit found specification defects that the earlier static
checks did not detect. Revision 55 repairs those contracts without
advancing a geometric proof or changing the source spine. The detailed
scope, original line references and source evidence are in
\code{AUDIT_REVISION54.md}; findings A29--A37 extend
\code{pre_m2_audit.csv}.

\begin{enumerate}[label=\textbf{D\arabic*:},leftmargin=*]
\item Raw graph presentations use auxiliary embedded tori without
  incompressibility. Their input types do not require the essential
  M1 torus system or its final cut decomposition. Generic boundary
  gluing is independent of that essential system.
\item The flow package no longer contains the long-time conclusion or an
  output oracle. M6 must prove it. The public theorem separately needs a
  construction of the flow package from the given manifold, including
  initial metric and normalization witnesses.
\item M1 consumes supplied topology transport. Its projection has no
  dependency on a surgery producer. Conversely, stabilization constructs
  raw transport before the finished package exists. Data-record
  dependencies no longer name the theorems that construct those records.
\item Hyperbolic normalization uses a complete-interior diffeomorphism
  and metric pullback. The metric restricted to a finite truncated core
  is not a complete interior witness.
\item Individual cut pieces carry component incidence, while the
  decomposition carries global coverage. Images of abstract cut
  interiors cover the complement of the tori. A collar-excision
  realization requires its own comparison and reconstruction maps.
\item The package supplies the reference prime decomposition. Public
  prime factors of the closed input are closed; relative factors with
  external torus boundary are intermediate objects with their own ledger.
\item Finite selected history segments do not require a globally finite
  number of surgeries. The no-late-slice alternative still requires
  proved terminal classification and reconstruction.
\end{enumerate}

The strengthened audit checks the complete proposed dependency graph
against an explicit list of unresolved external aliases, rejects cycles,
and checks the structural/analytic and raw/output boundaries. These
checks validate declared dependencies; they cannot discover a mathematical
hypothesis omitted from both the graph and the prose. The persistent
\code{check_blueprint_guards.py} suite exercises the existing 104 graph,
relative-prime, filling and rewrite rejection cases as well as the new
boundary checks. Its machine-readable evidence records exact input hashes.

The seven K0 modules were rebuilt in an isolated copy with installed
Lean 4.33.1, and their 35 declaration axiom reports were checked.
The 47 M2A preparation rejection controls and six core-only template runs
were also repeated in isolation. No PC-dependent probe or consumer draft
was elaborated, and no production GC or PC build was run. Original
M2A/sandbox sources and evidence remain unchanged; fresh audit evidence
is stored separately. These checks do not prove geometric contracts.

Source checks reopened the local Hatcher splitting and disk-removal
passages, the KL boundary theorem and its flow application, and the
Benedetti--Petronio cusp formula, with the Martelli compactification
comparison. The stated KL inequalities and quantifier order agree with
the rendered source. Earlier source checks are reused only for unchanged
claims. The recent cap-avoidance and gluing arguments remain written
proofs with explicit foundation and relative-implementation obligations.

The remaining release blockers include H1--H4, relative normal-form and
prime-factor comparisons, outer ambient injectivity, complete model
realization, terminal/exceptional reconstruction and exact completed-PC
binding. A passed audit does not close them. The audit does not advance
the next MT/KL ambient-injection comparison.
% END migrated B029

\chapter{M2A release evidence and historical discovery}
\label{ch:m2a-evidence-appendix}

% BEGIN migrated B025
\section{PC-reuse gate and adapter surface}
\label{sec:pc-reuse-gate}

The M1 structural slice can be developed before the PC release, but the
analytic phase must begin from a pinned completed release. This revision fixes
the audit procedure and the deliberately small public adapter surface. It does
not guess declaration names from the current pre-release checkout.

\subsection{Release gate}

The GC project may import the PC repository only after all of the following are
recorded:

\begin{enumerate}[label=\textbf{P\arabic*:},leftmargin=*]
\item an immutable release tag or commit and the exact Lean/mathlib toolchain;
\item a clean root build from a fresh checkout;
\item a project-wide search showing no remaining \code{sorry}, \code{admit}, or
  unreviewed project axiom in the claimed completed result;
\item the public theorem name and category of the completed Poincar'e result;
\item the exact imports and namespace paths for every declaration classified
  as direct reuse or adapter reuse; and
\item a transitive \code{#print axioms} report for each imported theorem and for
  the adapter theorems created by GC.
\end{enumerate}

Until this gate passes, the blueprint may contain provisional interface names
and test doubles, but it must not describe a PC theorem as already available.

\subsection{Public \code{PoincareAdapter} surface}

The GC adapter exposes only three groups of consequences:

\setlength{\tablebodywidth}{\dimexpr\textwidth-4\tabcolsep\relax}
\begin{longtable}{@{}P{0.306818\tablebodywidth}P{0.454545\tablebodywidth}P{0.238636\tablebodywidth}@{}}
\textbf{Group} & \textbf{Adapter output} & \textbf{GC consumer}\\\hline
\endhead
Category conversion & Equivalences between the PC manifold category and the
GC smooth oriented 3-manifold wrappers, including orientation and dimension
hypotheses. & M1 wrappers and exceptional factors.\\
Poincar'e consequence & The completed simply connected closed 3-manifold
conclusion in the exact form needed by the spherical branch. &
\code{ExceptionalPrimeToGeometricPiece}.\\
Spherical recognition & The standard spherical-space-form or spherical-prime
geometric witness and its transport into \code{GeometricPieceCertificate}. & M1f
assembly.\\
\end{longtable}

The analytic theorem handles are deliberately a separate interface,
\code{PCRicciFlowHandles}. They are consumed by \code{RicciFlowSurgeryPackage} and the
A0 adapter layer, not by the M1 assembly module. This split keeps the
endpoint-only \code{PoincareAdapter} import from pulling analytic implementation
dependencies into the structural theorem.

The adapter deliberately does not export the PC proof's internal finite-
extinction history, private namespaces, implementation-specific quotient
types, or a theorem whose conclusion is already \code{Geometrizes M}. If a later
GC theorem needs one of those items, it must add a named adapter row with a
new audit.

\subsection{Audit matrix columns}

Every candidate PC result receives one row with these columns:

\begin{enumerate}[label=\textbf{A\arabic*:},leftmargin=*]
\item stable candidate identifier and mathematical role;
\item exact public declaration and namespace;
\item minimal import path, with no unnecessary implementation imports;
\item complete signature and all quantified hypotheses;
\item direct/adapter/GC-new classification;
\item adapter theorem and changed conventions, if applicable;
\item downstream GC module and consuming theorem;
\item clean-build target and toolchain version;
\item transitive \code{#print axioms} output; and
\item evidence link to the pinned source commit and review status.
\end{enumerate}

The matrix is a provenance ledger, not merely a checklist. A row marked
\code{direct} means that the exact imported conclusion is used. A row marked
\code{adapter} names the bridge theorem. A row marked \code{GC-new} records which PC
results are only analytic hypotheses for the new proof.

\subsection{Initial classification before the release}

Before the completed repository is available, the blueprint records only
provisional expectations:

\setlength{\tablebodywidth}{\dimexpr\textwidth-4\tabcolsep\relax}
\begin{longtable}{@{}P{0.337209\tablebodywidth}P{0.372093\tablebodywidth}P{0.290698\tablebodywidth}@{}}
\textbf{Candidate} & \textbf{Expected class} & \textbf{Reason for deferring final status}\\\hline
\endhead
Short-time flow and smooth-flow APIs & Direct or adapter & Exact category and
time-domain types must be compared.\\
Hamilton--Ivey pinching and reduced-volume tools & Direct or adapter & The
normalization and theorem hypotheses must be audited.\\
Flow compactness and derivative estimates & Direct or adapter & Pointed-limit
and regularity interfaces may differ from the GC package.\\
Global surgery, canonical neighborhoods, continuation, invariants & Adapter
likely & The GC producer needs a public history and transport format.\\
Long-time thick/thin output & GC-new & This is the central analytic output
consumed by M1e.\\
Hyperbolic limit-to-piece and local collapse & GC-new & These conclusions are
not assumed to be supplied by PC.\\
Finite extinction & Out of scope & No current GC dependency.\\
\end{longtable}

These are planning classifications only. The audit may move a result between
classes; it may not silently treat a file import as proof of compatibility.

\subsection{Adapter smoke tests}

After the release is pinned, the first adapter target is small:

\begin{itemize}[leftmargin=*]
\item import the completed Poincar'e theorem through one public GC module;
\item construct the spherical consequence used by the M1 mock examples;
\item compile the M1f assembly theorem with the real adapter in place of the
  test double;
\item compile the direct/adapter analytic handles without importing the full
  PC implementation tree; and
\item compare the real adapter's \code{#print axioms} report with the release
  ledger.
\end{itemize}

The smoke test does not attempt the long-time theorem. Its purpose is to prove
that the future PC release can be attached without changing the M1 endpoint
types.

\subsection{What can proceed now}

While the release is pending, the team can implement the M1 mock suite,
finalize the adapter signatures, and prepare import-boundary tests. It should
not fill the audit matrix with invented theorem names or claim direct reuse
based on source-file presence. The completed PC release remains a gate for the
genuine analytic phase, not for structural design work.

The adapter-surface audit is recorded in
\code{poincare_adapter_api.csv}; the candidate theorem matrix remains
\code{pc_reuse_audit_template.csv} until the real release is available.
% END migrated B025

% BEGIN migrated B030
\section{M2 inherited geometric foundation audit and adapter}
\label{ch:m2-foundation}

M2 is an integration milestone. The geometric foundation developed for the
Poincar\'e project is the inherited foundation for GC and must be reused. GC
must not create a second manifold, metric, tensor, connection, or curvature
library merely because its namespace or notation differs. A new declaration
is justified only when the audit identifies a genuine missing interface or a
mathematical gap needed by the Ricci-flow proof.

\begin{roadmap}{Audit the inherited geometric foundation, freeze its
conventions, and expose the smallest adapter surface needed by Ricci flow.
M2 does not reprove differential geometry. The first substantial new
mathematics begins in M3 with the Ricci-flow PDE layer.}
\end{roadmap}

\subsection{M2 scope and deliverable}

The M2 deliverable is a reviewed \code{FoundationReusePackage}: an auditable
record of the inherited declarations, their exact imports and axiom closures,
the sign and normalization ledger, and the adapter lemmas that make those
declarations available to the GC flow modules. The package is a project
interface, not a new geometric foundation.

M2 has six obligations:

\begin{enumerate}[label=\textbf{M2-\arabic*:},leftmargin=*]
\item identify the public manifold, smooth-structure, metric, tensor,
  connection, curvature, volume, geodesic, and injectivity-radius APIs in the
  completed PC release;
\item record the exact sign, contraction, scaling, orientation, and
  normalization conventions used by those APIs;
\item check that the inherited objects can be used with the GC manifold and
  boundary wrappers without duplicating their mathematical content;
\item prove only the small category, notation, coercion, and convention
  adapters needed by the M1 endpoint or the later flow modules;
\item record genuine gaps separately, especially time-dependent metrics,
  curvature-evolution identities, and PDE infrastructure; and
\item pass import, example, convention, and transitive \code{#print axioms}
  checks before M3 is promoted from a design to an implementation milestone.
\end{enumerate}

The audit distinguishes four outcomes for each candidate declaration:
\code{direct-reuse}, \code{adapter-reuse}, \code{GC-new-gap}, or \code{out-of-scope}. A source file
containing familiar names is evidence for an audit row, not evidence that the
GC theorem can import it unchanged.

\subsection{M2A inherited-module inventory}
\label{sec:m2a-foundation-inventory}

M2A is the first executable submilestone. It audits the static geometric
foundation before any new Ricci-flow PDE theorem is attempted. The inventory
must identify the exact declaration and its namespace, its minimal import,
all typeclass and regularity hypotheses, its convention choices, its
transitive axiom output, and its downstream GC consumer.

\setlength{\tablebodywidth}{\dimexpr\textwidth-4\tabcolsep\relax}
\begin{longtable}{@{}P{0.277778\tablebodywidth}P{0.500000\tablebodywidth}P{0.222222\tablebodywidth}@{}}
\textbf{Inherited area} & \textbf{Questions to answer} & \textbf{M2A result}\\\hline
\endhead
Manifold and smooth structure & Which manifold, chart, smooth-map, tangent,
and orientation types are public? How do they relate to the GC closed,
connected, orientable 3-manifold wrapper? & Direct reuse or a category adapter.\\
Metrics and norms & How are Riemannian metrics, pointwise norms, distances,
volume forms, and tensor norms represented? What is the scaling convention
under $g\mapsto\lambda g$? & Direct reuse plus notation/coercion lemmas.\\
Tensors and contractions & Which tensor product, contraction, pullback, and
pushforward APIs are stable? Are symmetry and index conventions explicit? & A
minimal tensor adapter, if required.\\
Connection and curvature & Which declarations provide affine and Levi--Civita
connections, covariant derivatives, Riemann curvature, Ricci curvature, and
scalar curvature? & Direct reuse or a named convention bridge.\\
Geometric analysis & Which APIs provide volume, geodesics, completeness,
injectivity radius, compactness, and smooth convergence? & Reuse ledger for
M3--M6 consumers.\\
Time-dependent metrics & Is there a family-of-metrics type with the required
smoothness in time and space, or only static metrics? & Adapter interface or a
recorded M3 gap; no evolution theorem in M2A.\\
PDE operators & Are Laplacian, heat operator, maximum-principle, and
regularity interfaces available with the required sign conventions? & Reuse
rows or explicit M3 gap records.\\
\end{longtable}

The source-level inventory is maintained in \code{m2_foundation_api.csv}. Its
entries are provisional until the completed PC release is pinned. M2A may
prepare the rows now, but it may mark a declaration \code{direct-reuse} only after
the release gate supplies the exact import and axiom evidence.

\subsubsection{Evidence from the available PC snapshot}

The available local \code{DifferentialGeometry} snapshot already exhibits the
intended reuse boundary. \code{Geometry/Metric/Basic.lean} defines the public
\code{SmoothRiemannianMetric} abbreviation; \code{Geometry/Metric/Family/Basic.lean}
defines \code{MetricConnectionFamily}, \code{MetricConnectionFamilyOn}, metric-family
time smoothness, and metric time derivatives; and
\code{Geometry/Curvature/Metric/Defs.lean} supplies the metric connection and
Riemann tensors. \code{Geometry/Curvature/Metric/LeviCivita.lean} bridges these to
the canonical Levi--Civita connection. The tensor and trace interfaces live
in \code{Tensor/RSTensor}, \code{Geometry/Curvature/Sections/Connection.lean}, and
\code{Geometry/Curvature/Sections/Trace.lean}. Volume and scaling are represented
in \code{Analysis/Integration/Measure/Riemannian}, while time intervals and
time-dependent metric regularity are represented by \code{Analysis/TimeInterval},
\code{Geometry/Metric/Family}, and \code{Analysis/Schauder/Parabolic}.

This snapshot is useful evidence for module discovery, but it is not yet the
assumed completed PC release: the source still contains outstanding proof
frontiers in unrelated canonical-neighborhood and extinction modules. Thus
M2A records these paths and signatures now, while \code{PCReleaseGate} remains the
authority for direct reuse, axiom closure, and final build status.

\subsection{Convention and normalization ledger}

The ledger is a single source of truth for every later module. At minimum it
records:

\begin{itemize}[leftmargin=*]
\item the definition of $R(X,Y)Z$ and the order of the first two curvature
  arguments;
\item the Ricci contraction and scalar-curvature convention;
\item the sign of the Laplacian and heat operator;
\item the Ricci-flow equation convention, recorded for M3 even if its proof
  is not part of M2A;
\item the scaling of metrics, curvature, volume, distance, and time;
\item orientations, boundary orientations, pullbacks, and gluing maps; and
\item the normalizations of reduced distance, reduced volume, entropy, and
  noncollapsing constants used by imported PC results.
\end{itemize}

A convention mismatch must produce a named adapter theorem and a tested
translation. It may not be repaired by silently changing a sign in a later
formula. Standard metrics on the round sphere, flat torus, and a hyperbolic
model are the first regression examples because they expose curvature and
scaling errors cheaply.

\subsection{M2A adapter boundary}

The adapter surface should be smaller than the inherited foundation. The
provisional shape is:

\begin{equation}
\begin{aligned}
 &\texttt{FoundationReusePackage} :=\\
 &\quad\texttt{PCFoundationAudit}\times\texttt{ConventionLedger}\\
 &\quad\times\texttt{FoundationAdapters},\\[4pt]
 &\texttt{TimeDependentMetric} :=\\
 &\quad\texttt{TimeInterval}\to\texttt{RiemannianMetric}(M)\\
 &\quad\text{with a separately stated smoothness predicate}.
\end{aligned}
\end{equation}

The exact Lean names remain subject to the PC API audit. The intended adapter
lemmas translate the PC manifold category into the M1 wrappers, expose a
time-dependent metric family to M3, and bridge any curvature or scaling
convention mismatch. They do not prove short-time existence, curvature
evolution, maximum principles, or compactness; those are M3--M4 obligations.

M2A must also establish a no-duplication rule: if a PC declaration already
provides a concept, GC imports and wraps it; it does not introduce a parallel
definition with a coincidentally similar name. If a wrapper is unavoidable,
the adapter records the equivalence, transport lemmas, and axiom provenance.

\subsection{M2A acceptance and handoff}

M2A is accepted when:

\begin{enumerate}[label=\textbf{A\arabic*:},leftmargin=*]
\item every inherited foundation row has an exact declaration, import, and
  status in \code{m2_foundation_api.csv};
\item the convention ledger has executable tests for curvature signs,
  contractions, metric scaling, and the three standard metrics;
\item the M1 wrappers compile using inherited geometric objects or explicit
  adapters, with no duplicate foundation definitions;
\item the time-dependent metric interface is either inherited, adapted, or
  entered as a named M3 gap with all missing hypotheses visible;
\item each adapter and reused theorem has a transitive \code{#print axioms}
  record; and
\item the gap ledger separates genuine PDE work from already available
  geometric infrastructure.
\end{enumerate}

The handoff is a stable \code{FoundationReusePackage} consumed by M3. M3 then
formalizes the Ricci-flow PDE layer: the time-dependent metric equation,
short-time existence and uniqueness, evolution identities, and the analytic
interfaces that are genuinely absent from the inherited PC release.

\subsection{M2A execution protocol and evidence bundle}
\label{sec:m2a-execution}

The next implementation step is an evidence-producing audit, not a second
geometric library. M2A can prepare its scripts, mock examples, and row
templates before the completed PC release is available, but the final
\code{direct-reuse} decisions require one pinned release. The audit must be
repeatable from a clean checkout and must leave enough evidence for another
person to reproduce every classification.

\subsubsection{Release manifest}

The first artifact is a release manifest that fixes the object being audited:

\begin{equation}
\begin{aligned}
\texttt{PCReleaseEvidence} := \{&\texttt{releaseTagOrCommit},\\
\texttt{leanToolchain},\\
\texttt{libraryVersion},\\
\texttt{sourceTreeHash},\\
\texttt{buildCommand},\\
\texttt{publicImports},\\
\texttt{axiomPolicy},\\
\texttt{artifactLocations}\}.
\end{aligned}
\end{equation}

The manifest records the exact PC commit or release tag, Lean toolchain,
DifferentialGeometry version, source-tree hash, dependency lock, successful
build command, and the policy for permitted axioms. It also records the
paths of the declaration report, \code{#print axioms} report, convention-test
report, and adapter theorem report. A moving branch, an unpinned local copy,
or a README statement without a reproducible build is not a release gate.

\subsubsection{Row-by-row audit procedure}

For every row in \code{m2_foundation_api.csv}, M2A performs the following sequence:

\begin{enumerate}[label=\textbf{E\arabic*:},leftmargin=*]
\item run a minimal import probe containing only the candidate module and its
  stated dependencies;
\item record the exact namespace, declaration type, theorem hypotheses,
  regularity assumptions, typeclass assumptions, and source locator;
\item run \code{#check} and \code{#print} on the public declarations and record
  their transitive \code{#print axioms} output;
\item execute the relevant convention and standard-model examples, including
  negative tests that would detect a reversed sign or contraction;
\item classify the declaration as \code{direct-reuse}, \code{adapter-reuse},
  \code{GC-new-gap}, or \code{out-of-scope}, with a reason tied to the evidence; and
\item record the downstream consumer and the exact theorem or adapter that
  closes the row, or record the M3 gap that remains.
\end{enumerate}

The row is not complete when a source file merely imports successfully. It is
complete when the public declaration, its hypotheses, conventions, axiom
closure, and downstream use are all identified. If several declarations
jointly provide one concept, the inventory records each declaration and a
separate composition row rather than hiding the dependency in prose.

\subsubsection{Convention and standard-model test suite}

The smoke tests are executable Lean examples attached to the audit report.
They should cover the round sphere, Euclidean space, flat torus, and a
hyperbolic model when the latter is exported by the pinned release. The suite collectively checks curvature signs, Ricci and scalar contractions,
metric/volume/distance scaling, and compatibility of Laplacian and Ricci-flow
signs. Record which named test covers each convention; do not require every
model fixture to reproduce the entire suite. The tests must include a deliberately altered
sign or scaling expression that fails, so that a vacuous theorem cannot pass
the convention gate.

The hyperbolic test is allowed to remain a named M1 or M6 gap if the PC
release has no public hyperbolic model. That absence must not trigger a new
static curvature library in M2A; it is either supplied by an existing model
module or entered as a downstream geometry gap.

\subsubsection{Adapter decision rules}

An adapter is permitted only when the inherited declaration already contains
the mathematical concept and GC needs a boundary conversion. A permitted
adapter has a named source and target type, an equivalence or transport
lemma, explicit convention translations, and a transitive axiom report. A
wrapper that redefines curvature, metric, or connection data is rejected as
duplication. A missing evolution identity, maximum-principle theorem, or
short-time existence theorem is recorded as a genuine M3 gap rather than
manufactured as an adapter.

\setlength{\tablebodywidth}{\dimexpr\textwidth-4\tabcolsep\relax}
\begin{longtable}{@{}P{0.261364\tablebodywidth}P{0.397727\tablebodywidth}P{0.340909\tablebodywidth}@{}}
\textbf{Audit lane} & \textbf{Can proceed before the PC release?} & \textbf{Required handoff}\\\hline
\endhead
Release and import probes & Prepare scripts and expected declarations; final run waits for the pinned release. & Reproducible manifest and minimal-import report.\\
Convention and scaling tests & Prepare model fixtures and failure tests; execute against the release. & Signed convention ledger and test output.\\
M1 wrapper integration & Exercise wrappers against a mock or provisional interface. & Import-clean wrapper branch with no duplicate foundation types.\\
Axiom and provenance audit & Prepare report format and forbidden-axiom checks. & Per-declaration transitive axiom report.\\
Gap classification & Draft candidate M3 gaps from the snapshot; confirm each against the release. & Numbered gap ledger with hypotheses and consumer.\\
\end{longtable}

These lanes are independent after the release manifest is fixed. The release
manifest and exact declaration inventory are the shared prerequisites. The
convention lane may proceed in parallel with import probes, while adapter
implementation must wait for the corresponding row classification. No lane
may silently promote a provisional snapshot result to a completed-PC fact.

\subsubsection{M2A completion record}

M2A produces one completion record containing the release manifest, completed
foundation inventory, convention-test output, adapter theorem list, gap
ledger, import/axiom report, and a short downstream-consumer map. The record
must state explicitly which declarations are inherited unchanged, which are
transported by adapters, which are new M3 obligations, and which are outside
GC. The M3 branch may begin implementation only after this record has no
unclassified row and every M3 gap has a named statement and dependency.


\subsection{Prepared M2A implementation bundle}
\label{sec:m2a-prepared-bundle}

Revision 38 established the preparation lane in \code{m2a/}; revisions 39--44
extend it. The bundle contains 111 source-inspected declaration candidates,
333 generated Lean probe files, 11 convention-test specifications,
10 unresolved audit questions, twenty-one refined consumer targets, and
26 disabled consumer records. All 13 foundation rows and all nine broader
PC-handle rows have preparation links. These counts measure audit coverage,
not proved geometry or accepted reuse. Final classifications and release
bindings remain empty.

\subsubsection{Release identity and the bootstrap audit}

\code{m2a/templates/release_manifest.template.json} records the immutable
commit and Git tree, clean-checkout state, source manifest, Lean toolchain,
Lake files and dependency identities, public theorem and expected type,
reviewed axiom policy, and actual command/output artifact records. Every
release-dependent value is currently null or explicitly pending. The source
manifest format specifies sorted relative paths and canonical JSON hashing;
it does not reinterpret the archived books' legacy directory digests.

The public-theorem probe must run \emph{before} the release gate can pass.
\code{ReleasePublic.lean.in} is the bootstrap template for this step after
\code{PCReleasePin}. It checks the selected public declaration, prints its
full type and axioms, and checks a separately reviewed expected type. This
step is distinct from the foundation-row probes, which run only after
\code{PCReleaseGate}. A successful public probe is one input to the gate,
alongside the fresh root build, gap scan, immutable identity, and policy
review; it is not itself the gate decision.

\subsubsection{Declaration probes and consumer checks}

\code{m2a/declaration_candidates.json} records, for every P01--P111 candidate,
a qualified name, a single defining-module import, the snapshot file and
line, its file hash, the relevant foundation rows, named GC consumers, and
an exact review question. Qualified names are source-inspected; they have
not been elaborated against the provisional snapshot. Revalidate and rebind
each candidate to the completed release before using it.

Each candidate has three independent files. \code{Import.lean} checks its
single proposed module; \code{Inspect.lean} runs \code{#check} and
\code{#print} with explicit names, arguments and universes;
\code{Axioms.lean} records the transitive axiom closure. A defining-module
import is a candidate for a small public import, not a proof that its
transitive imports are minimal. None of these 333 PC files has been run.

The separate \code{Consumer.lean.in} template requires an explicitly stated
consumer type and a proof application. Import success, a name match, or a
printed definition cannot replace this test. The comparison must expose
implicit/typeclass hypotheses, regularity, boundary and time domains, and
any source-to-target convention translation. The expected type may not
assume the desired endpoint or conceal a missing hypothesis in an oracle.

\subsubsection{Four concrete source distinctions}

The initial source inspection exposes the following review obligations:
\begin{enumerate}[label=\textbf{D\arabic*:},leftmargin=*]
\item P02's \code{BoundaryManifold} denotes the boundary subtype of an
  ambient manifold. It is not already the GC compact manifold-with-boundary
  wrapper. P03's closed-embedding theorem does not supply collars,
  orientations, or reconstruction.
\item P10's \code{MetricConnectionFamilyOn} stores metrics, connections,
  and metric compatibility on flow times. Joint smoothness and
  torsion-freeness do not follow merely from those fields. P11 uses abstract
  time-derivative data; P12 must be connected to the intended real derivative.
\item P13's \code{metricFamilySmoothOn_of_contMDiffOn} is conditional on a
  joint-smoothness hypothesis. It does not construct that hypothesis for an
  arbitrary family.
\item P14's distance scaling is an extended-distance identity, and P15's
  volume scaling is a measure identity with an explicit dimension exponent.
  Converting to real distances requires finiteness; specializing to three
  dimensions requires the dimension hypothesis.
\end{enumerate}
These are source-level findings and targeted review questions, not final
release classifications or assertions that a new theorem is missing.

\subsubsection{Convention specifications and failure controls}

\code{m2a/convention_tests.json} specifies each fixture, hypotheses, expected
result, a discriminating negative control, consumer scope, and gate. The
actual geometric Lean statements and proofs remain pending the release.

\setlength{\tablebodywidth}{\dimexpr\textwidth-4\tabcolsep\relax}
\begin{longtable}{@{}P{0.12\tablebodywidth}P{0.44\tablebodywidth}P{0.44\tablebodywidth}@{}}
\textbf{Tests} & \textbf{Required check} & \textbf{Failure control or scope}\\\hline
\endhead
C01--C02 & Curvature order, tensor variance, Ricci and scalar contractions on the actual round sphere. & Use nonzero curvature to detect a swapped slot or reversed sign.\\
C03--C05 & Euclidean, compact flat-torus, and hyperbolic model conventions. & Euclidean completeness does not replace the torus quotient; a missing hyperbolic model blocks its own consumers.\\
C06 & Positive constant metric scaling of connection, curvature, distance and volume. & Use scale 4 and nonzero quantities, not only the identity scale.\\
C07 & Scalar/bundle Laplacian comparison, heat sign and operator scaling. & A nonconstant Euclidean quadratic fixture detects the wrong heat sign and exponent.\\
C08 & Metric-family joint regularity, real-time derivative and endpoint domains. & A compatible family without smoothness evidence does not pass.\\
C09 & Ricci-flow time rescaling with its chain rule and rescaled interval. & Retain a named M3 obligation if missing; do not block unrelated M1 work.\\
C10 & Boundary orientation, collars and pullback transport. & Equal induced orientations at the two collar ends fail the gluing contract.\\
C11 & Reduced-distance, reduced-volume and entropy normalizations. & Preserve M4/M6 obligations; definitions or notation alone are not analytic proofs.\\
\end{longtable}

The proposed constant-scaling calibration distinguishes tensor variance.
For $g_c=cg$, $c>0$ constant, the test expects the connection and $(1,3)$
curvature to be unchanged, $(0,4)$ curvature to scale by $c$, $(0,2)$ Ricci to
be unchanged, and sectional/scalar curvature to scale by $c^{-1}$. In
three dimensions distance scales by $\sqrt c$ and volume measure by
$(\sqrt c)^3$. These are test contracts to match through the reviewed
convention bridge, not geometric theorems proved in this iteration.

A round-sphere sign test uses scalar curvature $6$, while the normalized
hyperbolic test uses $-6$, after the chosen convention translation. Zero
Euclidean curvature alone cannot detect a sign reversal. The hyperbolic
universal model is tested for completeness and curvature; it is not claimed
to have finite volume. Quotient/core volume is a separate branch obligation.

\subsubsection{Open questions and scoped completion}

\code{m2a/open_obligations.json} records ten questions covering the boundary
wrapper, tensor transports, geodesic/injectivity/compactness interfaces,
time regularity, scaling, flat and hyperbolic model tests, PDE hypotheses,
Perelman normalizations, and the public-PC/endpoint boundary. Each has a
required statement, comparison/search scope, milestone, and consumers.
They are unresolved audit questions; none is already a confirmed
\code{GC-new-gap}. A failed search or compiler invocation cannot establish
such a classification.

\code{foundation_coverage.json} and \code{reuse_coverage.json} link these
preparations to the existing CSV inventories. The new
\code{preparation_record} column is deliberately separate from final
\code{declaration_evidence_ids}, \code{adapter_evidence_ids}, and
\code{gap_ids}, which remain empty. The initial candidate list is not an
exhaustive declaration inventory: split mixed concept outcomes and discover
additional declarations before accepting a row.

\code{m2a/templates/completion.template.json} contains 26 consumer records,
all disabled, each with applicable foundation rows, convention tests, and
open obligations. Distinguish readiness to \emph{implement} a classified M3
gap from readiness to \emph{use} the missing theorem. A consumer with an
undischarged geometric or model prerequisite remains blocked even when its
inventory classification is complete. Finite extinction remains outside the
default route, and spherical-space-form recognition remains separate from
the simply connected Poincar\'e conclusion.

\subsubsection{Preparation tooling and its verification}

\code{m2a/prepare.py} validates the preparation schema, reciprocal links,
pending states, and generated probe content. It can recheck the recorded
snapshot file hashes read-only; it has no PC-build or release-acceptance
command. Forty-seven deliberately invalid preparations were rejected, including
premature direct reuse, an enabled consumer, a false passed model test,
a promoted search failure, and an edited generated probe.

Installed Lean 4.33.1 separately exercised the templates in six temporary
runs using only \code{Init} and \code{Nat.add_zero}: import, inspection,
axioms, a correct consumer type, a deliberately wrong consumer type, and
the public-theorem bootstrap format. Five runs succeeded and the wrong
consumer type failed as intended. This tests the harness and type boundary;
it runs no PC import and proves no geometric convention theorem.
\code{m2a/evidence/preparation_check.json} records the exact inputs, commands,
outputs, scope, and hashes. The blueprint audit checks the integrity of this
stored evidence without invoking Lean.

The remaining action is to supply the actual completed release, execute the
bootstrap release gate, rebind the candidates, and run the consumer and
convention checks. Until then, refine the ten open questions and the exact
literature contracts without promoting snapshot discoveries to release facts.

\subsection{Concrete inherited-API consumer targets}
\label{sec:m2a-consumer-targets}

Revision 39 refined O04, O05, and the short-time part of O08 into nine
targets in \code{m2a/consumer_targets.json}. Revision 40 located T09's compact
packaged-solution route in Section~\ref{sec:m2a-short-time-route}. Revision 41
adds T10--T12 for existing scalar and curvature-evolution APIs, described in
Section~\ref{sec:m2a-scalar-curvature}. Revision 42 adds T13--T16 for norm
estimates, tensor principles, pinching and event preservation in
Section~\ref{sec:m2a-norm-pinching}. Revision 43 follows finite-cap events
and scalar history propagation in Section~\ref{sec:m2a-event-history}.
Revision 44 adds the Hamilton--Ivey and representation bridges in
Section~\ref{sec:m2a-hi-bridges}. The current bundle contains 333 generated
import/inspection/axiom files and twenty consumer drafts with 35 explicit
applications. All 353 PC-dependent files remain unrun.
These targets refine the existing
foundation and theorem inventories; they do not add production theorems,
accept M2A, or establish a new foundational geometry layer.

Each target records its parent question, candidate declarations, hypotheses,
GC consumers, prerequisite targets, remaining comparison work, and a
discriminating failure control. Reciprocal links connect targets to the
foundation rows, convention specifications, open questions, and the 26
disabled completion records. The target dependency graph is acyclic.
Source findings are distinct from release and proof evidence, which remain
empty. The following table is the implementation queue after the release
gate and statement comparison, not a declaration that the targets are proved.

\setlength{\tablebodywidth}{\dimexpr\textwidth-4\tabcolsep\relax}
\begin{longtable}{@{}P{0.1\tablebodywidth}P{0.44\tablebodywidth}P{0.46\tablebodywidth}@{}}
\textbf{Target} & \textbf{Concrete output and source candidates} & \textbf{Remaining boundary to check}\\\hline
\endhead
T01 & Flow/regular-time inclusion and exact slab endpoints; P21, P33. & The inherited interval structure does not require a nonempty regular set, order-convex carrier, or earliest initial time.\\
T02 & Joint frame smoothness on regular times and tensor continuity on the carrier; P10, P13, P22. & Supply all four regularity fields; a metric-compatible family alone is insufficient.\\
T03 & Real derivative of a metric coefficient at a regular time; P11, P12, P23, P24. & Compare scalar and bilinear-map derivatives; handle abstract time data and endpoints separately.\\
T04 & Actual Ricci equation with the chosen connection and convention; P24, P25, P32. & Identify the driving tensor with metric Ricci; generic metric evolution allows an arbitrary tensor.\\
T05 & Constant scaling of the connection and mixed/covariant curvature; P06, P07, P16, P26, P29. & Match representations and slots; retain the source boundary and typeclass assumptions.\\
T06 & Covariant Ricci invariance and inverse scalar/sectional scaling; P08, P09, P27, P28, P30. & Distinguish raised-index Ricci and a genuine sectional two-plane from a totalized quotient.\\
T07 & Recentered carrier and regular intervals; P21, P33--P35. & Preserve strict endpoints and the negative-time portion when recentering at an interior time.\\
T08 & Rescaled Ricci equation and full source solution package; P33, P36, P37. & Bind the source package to GC and run the nonstationary convention test. Source candidates already exist.\\
T09 & Compact packaged short-time solution with initial and uniform-bound outputs; P38--P52. & The regularity producers are now located. Verify the release exports and preserve their endpoint and bound contracts.\\
T10 & Compact scalar ODE comparison and strong maximum principle; P53--P56. & Retain the fixed derivative slab, scalar regularity, compact domain, and closed regular-slab hypothesis of the strong wrapper.\\
T11 & Scalar evolution and complete-slab lower bounds; P57--P62. & Retain the carrier/regular distinction, complete slices, positive elapsed time for division, and exact scalar normalization.\\
T12 & Full covariant curvature evolution and fixed-vector evaluation; P63--P66. & Retain the rough Laplacian, quadratic and Ricci slot terms; compare moving-frame transport separately.\\
T13 & Intrinsic curvature-norm tower and subsolution, with conditional bound simplification; P67--P70. & Use squared norms, retain the strict interior start, and supply one uniform bound if linearizing the potential.\\
T14 & Tensor nonnegativity and strict positivity from the inherited assumptions package; P71--P74. & Retain compact unit tangent slabs, actual spatial derivatives, null reaction, and strict-positive-time scope.\\
T15 & Public Hamilton--Ivey scalar, logarithmic and asymptotic pinching; P75--P81. & Retain the initial operator bound and distinguish sectional and trace-normalized eigenvalues.\\
T16 & Existing cap preservation producers and exact surgery-event fields; P79--P85. & Match metric, carrier and predicate. Projecting an assumed event record does not construct it.\\
T17 & Uniform finite-cap event construction and conditional cutoff-frontier conversion; P86--P91. & Compare the concrete quotient event with the stage-indexed topology event; frontier preservation remains input.\\
T18 & Existing scalar smooth-stage, terminal and eventwise history propagation; P92--P97. & Retain terminal convergence and cutoff records; distinguish ordinary observations from the two event sides.\\
T19 & Existing fixed-region scaling, factor-two barrier identity and trace-normalized pinching; P98--P101. & Keep $C=2K$, the initial bound and region parameter; identify the event Rayleigh representation separately.\\
T20 & Complete-slab, ancient and rescaled-limit Hamilton--Ivey consequences; P102--P105 and P111. & Retain completeness, carrier endpoints and supplied convergence; a blow-up limit is distinct from an unrescaled terminal metric.\\
T21 & Existing forward event/history representation conversion; P106--P110. & Supply completion and old-locus embeddings; the finite-history constructor needs positive event count.\\
\end{longtable}

\subsubsection{Regular times, carrier continuity, and actual Ricci}

P21's \code{RealTimeInterval} supplies a carrier, an open regular subset,
a distinguished initial point in the carrier, and neighborhood information
at regular points. It does not assert that the carrier is order-convex,
that regular times exist, or that the initial point is the earliest time.
GC slabs must retain their explicit interval formulas; an empty regular set
cannot serve as a nonvacuous evolving-flow test.

P22's \code{MetricFamilySmoothOn} contains four substantive fields:
coefficient smoothness on \code{D.regular}, coefficient continuity on
\code{D.carrier}, total-space metric-tensor continuity on the carrier, and
joint local-frame coefficient smoothness on regular time times the frame
domain. This resolves the source-discovery question about whether a joint
regularity predicate exists. It does not produce that predicate for an
arbitrary P10 family. P13 is one conditional route to constructing it.

P23's \code{MetricFamilySmoothOn.hasDerivAt_inner} gives, for
$t\in D.\mathrm{regular}$ and fixed $x,v,w$,
\[
 \frac{d}{ds}\bigg|_{s=t} g(s)_x(v,w)
 = \left(\frac{d}{ds}\bigg|_{s=t} g(s)_x\right)(v,w).
\]
The Lean conclusion is \code{HasDerivAt}, with the derivative of the
bilinear-map family evaluated on the fixed vectors. This can support a
real-time GC interface directly. A consumer using P12's abstract
\code{TimeDerivativeData} still needs a theorem identifying its
\code{dtApply} with this real derivative. At regular times the carrier
is a neighborhood, allowing comparison with \code{HasDerivWithinAt};
endpoint derivatives require separate one-sided evidence.

P25's \code{IsMetricEvolutionOn} includes regularity, Levi-Civita evidence,
and a variation equation, but its \code{drivingTensor} is arbitrary.
T04 must identify that tensor with metric Ricci or use the actual source
\code{SolutionOn} interface. A stationary family with a freely chosen
zero driving tensor is not a proof of Ricci flow for an arbitrary metric.

\subsubsection{Scaling candidates already present in the source}

P26--P30 locate separate source theorems for covariant curvature, Ricci,
scalar curvature, the Levi-Civita curvature operator, and sectional
curvature. P26 concerns \code{metricRm04}; its relation to
\code{metricRm04At} and the GC slot convention remains an explicit bridge.
P29 concerns \code{riemannOp}, not a silently interchangeable mixed-tensor
name. P27 and P28 have the expected distinct powers: covariant Ricci is
unchanged and scalar curvature scales inversely. P30's sectional-curvature
definition is a totalized quotient, so the theorem also treats a zero
denominator; the geometric two-plane regression requires independent vectors.

The source module \code{Geometry/Flow/RicciFlow/Scaling/Parabolic.lean}
contains the interval construction and both the equation and full-package
preservation theorem. For positive constant $c$, recenter at
$\tau\in D.\mathrm{carrier}$ and use
\[
 \phi(s)=\tau+s/c,\qquad \widetilde g(s)=c\,g(\phi(s)).
\]
P33 defines the new carrier and regular sets as preimages under $\phi$ and
sets the distinguished initial time to $0$. If the original carrier is
$[0,T)$ with $T>0$, P34--P35 give
\[
 \begin{aligned}
 \widetilde D.\mathrm{carrier}&=[-c\tau,\ c(T-\tau)),\\
 \widetilde D.\mathrm{regular}&=(-c\tau,\ c(T-\tau)).
 \end{aligned}
\]
Thus $T=2$, $\tau=1$, $c=4$ produces carrier $[-4,4)$, not $[0,8)$.
The new initial point $0$ need not be the left endpoint. Positivity belongs
in the scaling target even though the bare interval definition permits an
arbitrary real scale parameter.

P36's \code{metricVariation_parabolic} combines the time chain rule with
constant metric scaling. P37's \code{parabolicSolution_isSolutionOn}
preserves the entire source solution package. Both consume an existing
\code{IsSolutionOn} proof; neither creates a flow from arbitrary metrics.
Their presence removes the provisional ``unlocated scaling theorem'' issue
for these source contracts. Release elaboration, convention comparison,
axiom review, and the GC package bridge remain pending, so no final reuse
classification follows. A genuinely missing adapter or analytic statement
must be scoped only after that comparison.

\subsubsection{Raw short-time output and the stronger exports}

P31's \code{ricci_flow_short_time_existence} assumes a compact, Hausdorff,
smooth manifold with a boundaryless model, finite positive dimension, and
an initial smooth Riemannian metric. It returns $T>0$, a family with
$g(0)=g_0$, joint chart Gram smoothness on $[0,T)$, and the metric equation
as \code{HasDerivWithinAt} within $[0,\infty)$ for $t\in[0,T)$.
The initial derivative is one-sided; the endpoint $T$ is excluded.

P32's \code{isSolutionOn_of_reg} is a constructor with additional inputs:
\code{MetricFamilySmoothOn}, an interior \code{HasDerivAt} Ricci equation,
scalar continuity, scalar time differentiability on the carrier, Ricci
tensor continuity, and covariant curvature continuity. Revision 40 locates
their producers and the stronger public compact export; see
Section~\ref{sec:m2a-short-time-route}. The earlier raw-output comparison
therefore becomes a source map for reuse, rather than a list of assumed
new analytic gaps. The raw short-time result alone still does not assert
uniqueness, long-time existence, or a terminal slice at $T$.

\subsubsection{Typed Lean drafts and checks performed}

The initial four files in \code{m2a/consumer_probes/} make a portion of the future
consumer check reviewable now:
\begin{itemize}[leftmargin=*]
\item Q01, \code{TimeDerivative.lean}, states the real metric-coefficient
  derivative type with the regular-time hypothesis.
\item Q02, \code{ConstantScaling.lean}, separately states covariant Ricci
  invariance and scalar inverse scaling.
\item Q03, \code{ParabolicScaling.lean}, states both interval formulas and
  preservation of the source solution package.
\item Q04, \code{ShortTime.lean}, extracts the initial value and one-sided
  equation from short-time existence. This deliberately weaker output does
  not construct the packaged solution or discard its extra prerequisites.
\end{itemize}
Each file has a single candidate import, explicit expected types, concrete
proof applications, and per-declaration \code{#print axioms} commands.
There are no local \code{sorry}, \code{admit}, or \code{axiom} declarations
in these drafts. They have not been elaborated, so even syntax, implicit
instances, and proof applications remain subject to the release check.
Draft source is not kernel evidence, and a typed draft is not yet an
implemented geometric convention test.

The preparation validator checks target links, prerequisite cycles, draft
imports/declaration lists, and pending states. Forty-seven deliberate invalid
preparations are rejected, including premature target classification,
invented proof evidence, a cyclic prerequisite, and a falsely executed
draft. The six existing \code{Init}/\code{Nat.add_zero} template runs were
rerun successfully, including the intended wrong-type rejection. They test
tooling only and import no PC module. The 111 discovery hashes and source
lines were separately rechecked against the read-only snapshot. Existing
K0 sources and their checked evidence are preserved.

The located T09 route and its three additional drafts are described next.
M2A and K1 remain pending until the completed release, consumer comparison,
convention tests, and axiom evidence satisfy the existing gates.

\subsection{Located compact short-time solution route}
\label{sec:m2a-short-time-route}

Revision 40 resolves the source-search part of T09. The public compact
short-time exports already return the inherited complete bounded-curvature
solution, with the extra regularity and endpoint data needed for a careful
GC comparison. The preferred initial-time candidate is P38,
\code{exists_completeBoundedCurvatureSolutionOn_of_compact}; P39,
\code{exists_completeBoundedCurvatureSolutionOn_from_time_of_compact},
allows an arbitrary starting time. Both are in
\code{Geometry/Flow/RicciFlow/ShortTime/Compact.lean}. Their names, statements
and construction have been inspected in the incomplete snapshot; no PC
import, axiom audit or release classification has been performed.

Revision 40 brought the inventory to 52 candidates and 156 generated probes,
with seven consumer drafts and ten applications. Revision 41 adds the scalar
and curvature contracts in Section~\ref{sec:m2a-scalar-curvature}. T09's
\code{source_route} retains nine construction/comparison steps within the
current twenty-one top-level targets; the production inventory remains at 136 rows. The route is a source map for
reuse, not a second implementation of the PC analytic foundation.

\subsubsection{Preferred public output and its precise scope}

For a compact Hausdorff smooth manifold with boundaryless model and an
initial smooth Riemannian metric $g_0$, P39 gives some $b>a$ and an inherited
\code{CompleteBoundedCurvatureSolutionOn} on $[a,b)$, together with:
\begin{enumerate}[leftmargin=*]
\item the initial equality $g(a)=g_0$;
\item a single constant $C\geq0$ bounding the squared norm of covariant
  curvature for every $t\in[a,b)$ and every point;
\item joint smoothness of the total-space metric section on $[a,b)\times M$;
\item the Ricci equation as \code{HasDerivWithinAt} within $[a,\infty)$,
  at every $t\in[a,b)$, including the initial point.
\end{enumerate}
P38 specializes to $a=0$. The public declarations use a finite-dimensional
real normed model, without a nonzero-dimension assumption. Their proofs
handle the zero-dimensional case and transport through an inner-product
model internally. GC still specializes to dimension three; it need not
recreate that model change. Private helpers in the source file are trace
locations, not public declaration candidates.

The time interval has some positive length; the theorem does not supply a
prescribed lifespan. Its right endpoint is excluded, so it gives no
terminal slice there. The arbitrary-start-time form starts from supplied
smooth data. It does not by itself continue a flow through a singular time,
prove uniqueness, or establish a noncompact existence theorem.

\subsubsection{Construction map for the regularity inputs}

The following map explains why the extra inputs to P32 were not new GC
proof obligations merely because they were initially unlocated. It also
provides smaller probes if a completed-release export changes its signature.

\setlength{\tablebodywidth}{\dimexpr\textwidth-4\tabcolsep\relax}
\begin{longtable}{@{}P{0.13\tablebodywidth}P{0.4\tablebodywidth}P{0.47\tablebodywidth}@{}}
\textbf{Step} & \textbf{Source producer} & \textbf{Inputs and output to preserve}\\\hline
\endhead
ST01 & P31: raw short-time existence. & Initial metric, joint chart Gram smoothness on $[0,T)$ and one-sided equation; $T>0$.\\
ST02 & P43: \code{metricFamilySmoothOn_of_chartGram}. & Smooth chart coefficients on $(a,b)$ plus continuous coefficients on $[a,b)$ produce the full family predicate.\\
ST03 & P44: \code{metricVariationEquationOn_of_pde}. & Restrict the derivative domain from $[a,\infty)$ to $[a,b)$ and populate the regular-time Ricci equation.\\
ST04 & P45--P48: scalar, scalar-time, Ricci and covariant-curvature regularity. & Joint chart smoothness on a time set $J$ and \code{UniqueDiffOn} for $J$ supply all four required inputs.\\
ST05 & P41: \code{solutionOn_of_joint}. & Joint chart data and the one-sided PDE construct the complete \code{IsSolutionOn} record on $[a,b)$.\\
ST06 & P49: \code{rm04SlabSup}; P50: \code{RiemannianMetricComplete.of_compact}. & A shorter closed time slab and compactness give a common curvature bound; compactness gives the inherited completeness predicate.\\
ST07 & P40: \code{completeBoundedCurvatureSolutionOfJointRicciFlow}; P52: the package type. & Add slice completeness and curvature bounds to the chart/PDE inputs. The constructor assumes these data.\\
ST08 & P38--P39: public compact exports. & Preserve the initial value, common bound, joint metric-section smoothness and initial one-sided equation as additional outputs.\\
ST09 & P42: chart-to-section conversion; P13 and P51: total-space route. & An alternative for consumers with joint metric-section smoothness and an explicitly supplied \code{UniqueDiffOn} carrier.\\
\end{longtable}

P41 lives in \code{Extension/Regularity.lean}. It applies the four
regularity producers to $J=[a,b)$, using the source's
\code{uniqueDiffOn_Ico} dependency lemma. Scalar time regularity at $a$ is
\code{DifferentiableWithinAt} on $J$, not an ordinary two-sided derivative.
The source constructor obtains spatial Ricci-norm regularity from the
inherited smooth metric/curvature APIs as well; these are fields of the
existing record, not assumptions to hide in a GC oracle.

P42, \code{metricCLMSection_jointContMDiffOn_of_chartGram_on}, converts chart
Gram smoothness to joint total-space metric-section smoothness on an
arbitrary time set. P51, \code{isSolutionOn_of_joint_metric}, then uses that
section, the equation on regular times, and
\code{UniqueDiffOn} of the carrier to produce the source solution record.
The latter hypothesis is explicit: openness of the regular subset in
\code{RealTimeInterval} does not assert unique differentiability of the
whole carrier. Inspect the pinned dependency declaration during release
elaboration; no local mathlib build has been claimed here.

P40 is a conditional packaging definition, despite being located in a file
named \code{ShortTime/NoncompactRicciFlat.lean}. It consumes completeness
and curvature bounds; it does not establish those properties or general
noncompact existence from its import alone. The preferred compact public
route supplies them with P49/P50 instead.

\subsubsection{The bound quantifiers and the shorter slab}

Write $q_g(t,x)$ for the source squared norm
\code{normSq0S (g t) x 4 (metricRm04At (g t) x)}. P52's package field has
the quantifier order
\[
 \forall t\in[a,b),\quad
 \exists C_t\geq0,\quad \forall x\in M,\quad q_g(t,x)\leq C_t.
\]
The extra public output of P38/P39 is stronger:
\[
 \exists C\geq0,\quad
 \forall t\in[a,b),\quad \forall x\in M,\quad q_g(t,x)\leq C.
\]
These must occupy separate fields in the eventual GC consumer adapter if
both are needed. A successful extraction of \code{Q.curvatureBound} cannot
stand in for the second statement. The quantity is a squared covariant
curvature norm; conversions to an unsquared norm require their own exact
source equality or inequality.

The compact proof first chooses $0<d<T$, where $T$ comes from raw existence.
It restricts chart data to the closed slab $[0,d]$, applies P49 on that
compact slab with the norm, lowering and connection metric families all
equal to $g$, and restricts the common bound to $[0,d)$. The inner-product
model requirement of this intermediate estimate is handled inside the
public proof's model-change route. Joint continuity on the original
half-open interval alone is not the compactness argument. In particular,
this construction does not assert a uniform bound all the way to the
original excluded time $T$.

The one-sided equation at the initial time is likewise an extra public
output. The \code{IsSolutionOn.equation} field is indexed by regular times,
so extracting that field alone does not recover the equation at $a$.
Keeping both the common bound and initial equation prevents the source
package from being silently weakened during adaptation.

\subsubsection{Consumer drafts and the next comparison}

Q05, \code{ShortTimePackaged.lean}, states all outputs of the preferred
compact initial-time export, including the common bound and initial
equation. Q06, \code{ChartGramSolution.lean}, checks the conditional chart
constructor. Q07, \code{JointMetricSolution.lean}, checks the alternative
total-space constructor while retaining \code{UniqueDiffOn}. These three
files use concrete inherited declarations and explicit expected types,
without local proof placeholders. They remain unelaborated drafts. Q04's
weaker raw-existence extraction is preserved as a separate diagnostic.

Revision 40 added preparation checks rejecting promotion of per-slice bounds to a common bound, a claimed
excluded terminal slice, invented proof evidence for the route, and a
cyclic construction map. These are metadata safeguards, not Lean proofs of
the geometric assertions or a general theorem checker. Its 52 discovery
hashes/lines were checked read-only; revision 44 rechecks all 111 candidates. The six actual core-only template
runs were rerun; they use \code{Init} and \code{Nat.add_zero}, with the
wrong-type application rejected as intended. No PC or production GC build
was run, and the complete finite K0 suite is unchanged.

After the completed release passes \code{PCReleaseGate}, compare and
elaborate P38/P39 and Q05 first. Use P40--P52 and Q06/Q07 to resolve narrower
signature or regularity questions, then run the applicable C08 convention
checks and review axiom evidence. T09's source route is located; final
reuse classification, consumer acceptance and M2A/K1 remain pending.
Section~\ref{sec:m2a-scalar-curvature} now locates O08's scalar maximum-principle
and curvature-evolution exports, preserving their operator and domain
hypotheses. The first genuinely new mathematics remains whatever Ricci-flow
PDE gaps survive the completed-release comparison.

\subsection{Inherited scalar comparison and curvature evolution}
\label{sec:m2a-scalar-curvature}

Revision 41 continues the comparison with existing PC mathematics. Scalar
maximum principles, scalar-curvature evolution, complete-flow scalar lower
bounds, and covariant curvature evolution are already present in the
read-only snapshot. The task here is to connect those exports to GC, retain
their hypotheses, and check the completed release. An omission from an
earlier inventory is not evidence of a missing theorem. No new proof of
these analytic results is scheduled by this source comparison.

P53--P66 added fourteen located declarations in revision 41, bringing that
revision to 66 candidates, 198 generated probes and eleven drafts with
fifteen applications. Section~\ref{sec:m2a-norm-pinching} records the enlarged
revision-42 bundle. All PC-dependent files remain unrun. T10--T12 refine O08 and link to
\code{RicciFlowPDE} and \code{PCRicciFlowHandles}; their final classifications
and proof evidence are empty. The production inventory remains at 136 rows.

\subsubsection{Compact scalar comparison: T10}

P53 is \code{parabolicOperatorWithDrift} in
\code{Analysis/Parabolic/Operator.lean}. For the family $G$, drift $X$, and
slab endpoint $T$, its exact operator is
\[
 \mathcal P_{G,X,T}u(t,x)
 =D_{[0,T]}u(t,x)-\Delta_{G(t)}u(t,x)
   -\langle X(t,x),\nabla u(t,x)\rangle_{G(t)}.
\]
Here $D_{[0,T]}$ denotes Lean's \code{derivWithin} on \code{Icc 0 T}.
The parameter $T$ fixes the derivative set; it is not automatically the
evaluation time $t$. P54, \code{laplacianAt}, uses the supplied connection
and metric. Identifying that connection with Levi--Civita is part of the
consumer contract when the metric Laplacian is intended.

P55, \code{scalar_weak_maximum_principle_ode_compare_supersolution},
already gives compact boundaryless ODE comparison. Its conclusion is
$c(t)\leq u(t,x)$ on $[0,T]\times M$. Besides $T\geq0$ and initial order,
the inputs include $\mathcal P_{G,X,T}u\geq F(u,t)$, the corresponding
within-time ODE for $c$, and one Lipschitz constant for $F$ on the joint
scalar value set. The theorem explicitly asks for continuity of the
exponentially weighted difference, its spatial and gradient regularity,
and within-time differentiability of $u$ and $c$ at positive times.
These inputs must be supplied by inherited regularity results or proved
consumer applications; a name match does not supply them.

P56 is
\code{scalar_strong_maximum_principle_time_dependent_metric_of_metricFamilySmoothOn}
in \code{Analysis/Parabolic/MaximumPrinciple/Scalar/MetricFamily.lean}.
On a compact connected Hausdorff boundaryless manifold, a nonnegative
supersolution with a zero at $(T,y)$ vanishes on the whole slab. It requires
$T>0$, the scalar regularity inputs, a smooth metric family, and
Levi--Civita identification. In particular, its metric-family wrapper uses
\[
 [0,T]\subseteq D.\mathrm{regular},
 \qquad\text{including both endpoints.}
\]
Carrier inclusion alone is insufficient. Applied to a half-open short-time
solution, the GC consumer must choose and translate an interior regular
slab or locate an existing endpoint version with the required hypotheses.
Q11, \code{StrongMaximumPrinciple.lean}, retains this exact inclusion in
its typed application. These two compact exports have no lateral spatial
boundary data because the manifold is boundaryless. A complete noncompact
or boundary-domain comparison needs its own existing theorem and contract.

\subsubsection{Scalar evolution and complete-slab bounds: T11}

P57, \code{ScalarEvolutionEquationOn}, defines an equation at regular
times, using \code{HasDerivWithinAt} on the carrier. P58,
\code{scalar_curvature_evolution}, actually proves that equation from
\code{IsSolutionOn}. Its output is
\[
 D_{D.\mathrm{carrier}}R
   =\Delta_{\mathrm{flowG}(S)}R+2|\operatorname{Ric}|^2,
 \qquad t\in D.\mathrm{regular}.
\]
The norm is \code{normSq0S} of the inherited covariant Ricci tensor.
The source uses a finite-dimensional smooth boundaryless Hausdorff setting;
its model-completeness and first-order manifold instances must still be
resolved by the release probe. Compactness of $M$ is not a hypothesis of
this evolution export. P59, \code{scalarEvolution_of_isSolution}, permits
another family $G$ when both its metric and connection equal the flow
family at every regular time. This is already an equality-based transport
interface; a different sign or representation needs a proved bridge.
Q08, \code{ScalarEvolution.lean}, applies P58 to an explicit pointwise
derivative type.

P60, \code{scalar_parabolic_inequality_of_isSolution}, supplies
$(2/n)R^2\leq\mathcal P_{\mathrm{flowG}(S),0,t}R$ for positive dimension
$n$, positive regular time $t$, and $[0,t]$ in the carrier. Its parabolic
operator uses $[0,t]$ as derivative set. Combining it with P55 on $[0,T]$
requires an actual derivative-set comparison; substituting the two
operators by name is insufficient. For the following universal bound,
the more direct inherited exports already perform the analytic argument.

\setlength{\tablebodywidth}{\dimexpr\textwidth-4\tabcolsep\relax}
\begin{longtable}{@{}P{0.12\tablebodywidth}P{0.45\tablebodywidth}P{0.43\tablebodywidth}@{}}
\textbf{ID} & \textbf{Existing scalar lower-bound output} & \textbf{Time and domain hypotheses}\\\hline
\endhead
P61 & \code{scalar_curvature_lower_bound_on_interval_of_complete}: $-n/2\leq(b-a)R(b,x)$. & $a\leq b$, $[a,b]$ in the carrier, $(a,b)$ regular, and completeness of every slice.\\
P62 & \code{scalar_curvature_ge_neg_dim_div_two_time_of_complete}: $-n/(2(b-a))\leq R(b,x)$. & The same conditions with $a<b$, which justifies division.\\
\end{longtable}

Both exports are in \code{Geometry/Flow/RicciFlow/Estimates/ScalarLowerBound.lean}.
They retain sigma compactness and the smooth boundaryless Hausdorff
category. Neither asks for a global curvature bound or an initial scalar
lower bound. The public statements handle dimension zero as well. Their
terminal time $b$ must belong to the carrier but need not be regular; the
source extends the interior estimate to that endpoint by scalar continuity.
Thus they do not assert a value at the excluded endpoint of a half-open
flow. Q09, \code{ScalarLowerBound.lean}, retains the complete divided-bound
contract. In dimension three it specializes to
$R(b,x)\geq-3/(2(b-a))$; a sharper initial-data-dependent normalized bound
or an inequality across a surgery jump needs its own inherited export and
GC comparison. No such additional conclusion is inferred from P62.

\subsubsection{Full covariant curvature evolution: T12}

P63, \code{roughLap0SField}, is the inherited trace of the second metric
covariant derivative, expressed as covariant divergence of the first
derivative. P64 and P65, in
\code{Geometry/Flow/RicciFlow/Evolution/Curvature/Tensor.lean}, provide the
preferred public evolution endpoints:
\begin{itemize}[leftmargin=*]
\item \code{riemann_tensor_hasDerivAt_of_solution} differentiates the
  full covariant $(0,4)$ curvature tensor in its fixed fiber at $x$.
\item \code{riemann_hasDerivAt_of_solution} evaluates that derivative
  on four tangent vectors fixed in time.
\end{itemize}
Both consume \code{IsSolutionOn} and a regular time in a finite-dimensional
smooth boundaryless Hausdorff setting. They do not require an externally
supplied curvature-evolution certificate or the extra frame-regularity
data in the lower-level \code{rm04EvolutionFam} interface. With $Q$ denoting
the source's \code{curvatureQuadraticCombination} and $A_{\mathrm{Ric}}$
its \code{ricciDrift04}, the precise source convention is
\[
 \partial_t\operatorname{Rm}_{(0,4)}
 =\Delta_{\nabla}\operatorname{Rm}_{(0,4)}
   -2Q(g,\operatorname{Rm}_{(0,4)})-A_{\mathrm{Ric}}(g).
\]
This formula records the names and signs in the source; it does not
identify $Q$ with a differently normalized literature reaction term.
The definitions in \code{Geometry/Curvature/QuadraticTensor.lean} and
\code{Geometry/Curvature/RicciAction.lean} determine the contractions and
slot action. Omitting the Ricci term or replacing the tensor by mixed
curvature changes the statement. Q10, \code{CurvatureEvolution.lean},
contains both explicit applications, retaining the full right-hand side.

P66, \code{rm04Base_of_solution_any} in
\code{Evolution/Curvature/HamiltonEquation.lean}, gives an alternative
Hamilton reaction formula in an arbitrary finite orthonormal basis at the
chosen regular time. The basis is held fixed in its
\code{HasDerivWithinAt} derivative on the carrier. A basis orthonormal at
one time is not automatically a moving orthonormal frame. A consumer that
uses Uhlenbeck transport must compare the existing transport and
curvature-operator exports before requesting a new bridge. Neither P64 nor
P66 states an evolution equation through surgery or at an excluded terminal
slice.

\subsubsection{Checks and the next comparison}

Revision 41 brought the preparation checks to 26 rejected invalid records.
Its five added cases detect a strong-principle slab weakened to the carrier,
compact comparison promoted to noncompact comparison, a scalar bound
promoted across surgery, division allowed at zero elapsed time, and a
discarded Ricci slot term. These are checks of the recorded contracts,
not mathematical proofs. All 66 revision-41 source hashes and declaration lines were
checked read-only; revision 42 rechecks the enlarged inventory. The six actual \code{Init}/\code{Nat.add_zero} Lean
template runs were rerun, including the intended type-error rejection.
The checked K0 suite and its stored evidence are unchanged. The four new
consumer drafts are Lean code, but none has been elaborated against PC.

After \code{PCReleaseGate}, rebind these candidates, run Q08--Q11 and the
declaration/axiom probes, discharge C01/C02/C07/C08 where applicable, and
record only the necessary GC adapters. T10--T12 are located reuse targets,
not confirmed M3 gaps. Section~\ref{sec:m2a-norm-pinching} now compares the inherited curvature-norm
and tensor/Hamilton--Ivey interfaces, including the sharper initial-operator-
bound scalar estimate and the existing surgery preservation producers. The first
genuinely new mathematics remains the M3 Ricci-flow PDE gaps that survive
the completed-release comparison.

\subsection{Inherited norm estimates, tensor principles, and pinching}
\label{sec:m2a-norm-pinching}

Revision 42 located these analytic interfaces already present in PC.
P67--P85 added nineteen candidates, giving that revision 85 candidates and
255 generated declaration probes. T13--T16 connect intrinsic curvature-norm estimates,
tensor maximum principles, smooth Hamilton--Ivey pinching, and surgery
preservation data to named GC consumers. Four more handwritten files,
Q12--Q15, contain eight explicit applications. That revision had fifteen
consumer drafts with 23 applications. Section~\ref{sec:m2a-event-history}
records the enlarged revision-43 bundle. All PC-dependent files remain
unrun. The production theorem inventory stays at 136 rows.

These are comparisons for reuse. No new curvature-norm or Hamilton--Ivey
proof is inferred to be necessary from an earlier incomplete inventory.
The completed-release, convention, axiom and consumer gates remain in force.

\subsubsection{Intrinsic curvature-norm estimates: T13}

P67, \code{towerHeatBoundOn_of_solution}, derives the intrinsic derivative
norm tower directly from \code{IsSolutionOn}. It retains the negative
next-derivative term and the explicit reaction sum. The public result is
in \code{Evolution/Curvature/Derivatives/Evolution/SolutionHeatEquation.lean};
its derivative is within the carrier at a regular time. It does not ask
the consumer to supply a curvature-evolution certificate.

For the zeroth derivative, P68,
\code{isHeatPotSubsolutionOn_nablaKRm04NormSqIntrinsic_of_solution},
in \code{Estimates/FiniteTime/CurvatureBlowupRateIntrinsic.lean},
is the preferred packaged subsolution. Write
\[
 n=\dim E,\qquad u(t,x)=|\operatorname{Rm}_{(0,4)}(t,x)|^2.
\]
For a source solution on $[\alpha,\omega)$ and
$\alpha<t_1<\omega$, it gives \code{IsHeatPotSubsolutionOn} on
$[t_1,\omega)$ with potential
\[
 V(t,x)=c_n\sqrt{u(t,x)},\qquad
 c_n=\mathrm{rmTowerCost}(n,0)=32n^4.
\]
The last equality is the existing \code{rmTowerCost_zero} theorem in the
same module. Thus the regular-time inequality has the form
\[
 \partial_tu\leq\Delta u+c_nu^{3/2}.
\]
The package includes joint smoothness on regular times, joint continuity
on the carrier, smooth spatial slices, actual time differentiability,
and the differential inequality. The source requires positive finite
dimension, a complete normed model, the smooth boundaryless model/manifold
instances, and a Hausdorff manifold. Compactness and a global curvature
bound are not hypotheses of this intrinsic subsolution export.
The strict interior start is part of this wrapper; it does not assert its
carrier regularity at $t_1=\alpha$ or at the excluded terminal time.

P69, \code{curvatureNormSq_eq_nablaKRm04NormSqIntrinsic}, identifies an
alternate realized curvature norm with the intrinsic zeroth norm. Keep
its \code{Rm04RealizesSolutionConnectionOn} input, carrier-time membership,
and source inner-product-model assumptions. A name match alone is not
a norm identification.

P70,
\code{Rm04NormHeatEquationOn.exists_hasDerivAt_le_of_rpow_reaction_bound},
is a useful conditional scalar simplification. Given the norm heat
equation, nonnegative gradient term, $u\geq0$, $u\leq K$, $C\geq0$, and
reaction at most $Cu^{3/2}$, it obtains an actual derivative $d$ satisfying
\[
 d\leq\Delta u+C\sqrt K\,u.
\]
Here $K$ bounds the \emph{squared} curvature norm, and one bound is required
over all the quantified regular times and points. This theorem does not
construct those inputs. P67/P68 supply the geometric producer; the older
component-interface name \code{rm04NormHeatEquationOn_of_solution} still
has explicit raw-derivative, simplification and Laplacian certificates.
Q12, \code{CurvatureNormEstimates.lean}, retains the intrinsic subsolution
contract and the conditional $C\sqrt K$ simplification as separate applications.

\subsubsection{Tensor maximum principles: T14}

P73, \code{tensor_weak_maximum_principle}, is already available in
\code{Analysis/Parabolic/MaximumPrinciple/Tensor/Weak.lean}. It consumes P71,
\code{TensorWeakMaximumPrincipleAssumptions}, and returns nonnegativity of
the section's associated two-tensor on $[0,T]$. The package explicitly
contains nonnegative elapsed time, initial nonnegativity, the parabolic
supersolution inequality, the null-eigenvector reaction condition,
smoothness and metric compatibility of the connection, and identification
of the supplied first and second spatial derivatives.

P72, \code{TensorWeakMaximumPrincipleSectionCompactness}, exposes the
remaining regularity and compactness inputs: symmetry, barrier regularity,
compactness of metric-unit tangent bundles over positive closed subslabs,
continuity of metric and tensor quadratic evaluations, and fixed-barrier
continuity. These obligations do not disappear because
\code{CompactSpace M} is absent from P73's short public signature. The
theorem also retains its boundaryless, Hausdorff, smooth-manifold and smooth
tangent-vector-bundle instances. A release application must construct the
actual package for the selected tensor.

P74, \code{tensor_positive_definite_on_of_null_reaction_lower_bound},
uses the weak package together with a uniform positive null-reaction
lower bound $\eta>0$. Its conclusion is positive definiteness on $(0,T]$;
the initial slice is only assumed nonnegative. This is a strict-positivity
result, not a parallel-null-distribution classification. Q13,
\code{TensorMaximumPrinciples.lean}, applies both P73 and P74 without
weakening their inputs. Tangent two-tensors and curvature operators on
bivectors still require the correct representation bridge. The public
Hamilton--Ivey route below already performs its own propagation argument.

\subsubsection{Public Hamilton--Ivey exports: T15}

P75--P78 are public results in
\code{Geometry/Flow/RicciFlow/DimensionThree/HamiltonIvey/MaximumPrinciple.lean}.
They consume an actual \code{IsSolutionOn} on a compact smooth boundaryless
Hausdorff three-manifold, with the source complete-model instance. The
slab requirements are
\[
 T\geq0,\quad K>0,\qquad
 [t_0,t_0+T]\subseteq D.\mathrm{carrier},\qquad
 (t_0,t_0+T)\subseteq D.\mathrm{regular}.
\]
The initial curvature operator is bounded below by $-K$ at every point.
P81 identifies this predicate with $-K\leq\lambda_{\min}$ using an
orthonormal three-dimensional basis. An initial scalar lower bound is not
the same hypothesis.

Set $\tau=t-t_0$ and let $\lambda$ be the source's least sectional-operator
eigenvalue. P75, \code{hamilton_ivey_pinching}, gives both
\[
 R(t,x)\geq\frac{-6K}{1+4K\tau}
\]
and, where $\lambda(t,x)<0$,
\[
 R(t,x)\geq
 2(-\lambda)\left(\log\frac{-\lambda}{K}
                  +\log(1+2K\tau)-3\right).
\]
P76, \code{hamilton_ivey_pinching_k_one}, already provides the $K=1$
specialization. In particular, the sharper scalar estimate
$R\geq-6/(1+4\tau)$ sought after the T11 comparison is present in PC;
its initial \emph{operator} bound must be retained.

P77, \code{hamilton_ivey_asymptotic_pinching}, adds $\delta>0$ and gives
\[
 \max\{-\lambda,0\}\leq\delta R+
 \frac{2\delta K\exp(2+(2\delta)^{-1})}{1+2K\tau}.
\]
P78, \code{curvatureOperatorRegionPropagationOn_of_initial_lower_bound},
constructs the region-propagation result from the same flow and initial
data, including the zero-length case. The public estimates do not ask GC
to assume the desired propagation conclusion. Q14,
\code{HamiltonIvey.lean}, states and applies the full scalar/logarithmic
and asymptotic exports.

P79 defines the trace-normalized operator as twice the sectional operator;
P80 proves the corresponding least-eigenvalue identity in an orthonormal
basis. Therefore its least eigenvalue is $\nu=2\lambda$. This factor, the
time origin, and the metric scaling must survive any GC normalization
adapter. C01/C02/C06/C08/C09 record the applicable convention checks.

\subsubsection{Surgery preservation has existing producers: T16}

P82, \code{InFixedHamiltonIveyRegion}, and P83,
\code{GeometricCutoffRecord}, are in
\code{Surgery/Topology/GeometricCutoff.lean}. The event predicate quantifies
over every orthonormal basis and uses an infimum of Rayleigh values for
twice the curvature matrix. It must be compared with the intrinsic
eigenvalue and \code{fixedHamiltonIveyRegion} interfaces, not silently
identified with them.

The event record requires preservation of that region for every $a>0$,
and preservation of every supplied scalar floor $L\leq0$, from the
incoming terminal regular region to the event output metric. Q15,
\code{SurgeryPinchingFields.lean}, checks those two field projections
\emph{assuming the complete record}. It does not construct a cutoff event
or prove the existence of its preservation fields.

There are already analytic producers to compare:
\begin{itemize}[leftmargin=*]
\item P84,
  \code{exists_normalizedDatum_positiveSideInsertionMetric_pinching},
  chooses controlled small parameters and proves positive scalar curvature
  and fixed-region preservation for normalized cap insertion data, with
  $0<\delta\leq\delta_0$ and derivative order $k\geq2$.
\item P85, \code{fullMetricConclusions}, proves a conjunction for the
  finite prepared retained metric. Its outputs include fixed-region
  preservation and preservation of all nonpositive scalar floors. It
  assumes disjoint smooth cylinder embeddings, retained-domain and
  recentered neck data, canonical insertion witnesses, and
  \code{StaticInsertionAdditionalProperties}.
\end{itemize}
These results are in \code{Surgery/StandardCap/NormalizedInsertionPinching.lean}
and \code{Surgery/StandardCap/FullMetricConclusions.lean}. Matching their
constructed metric and carrier to the exact event output is a concrete
source-comparison task. Local insertion is not by itself a global surgery
history theorem; projection of a record is not its construction.

For comparison of the logarithmic expressions, the parameter identity
located in P99 below is
\[
 \nu=2\lambda,\qquad
 a=\frac{1+2K\tau}{2K},\qquad
 \log(a(-\nu))=\log\frac{-\lambda}{K}+\log(1+2K\tau)
\]
on the negative-eigenvalue branch with $K>0$ and $\tau\geq0$.
P99 already proves the corresponding barrier identity in PC, and P98
proves the fixed-region scaling identity; see Section~\ref{sec:m2a-hi-bridges}.
Their GC applications remain unelaborated. The identities do not identify
the event's quantified Rayleigh predicate by themselves. Compare the
existing predicate/transport bridges next. A history-level
consumer must also track the same time origin through smooth stages,
terminal passage and each surgery event; independently restarting a
weaker slab estimate does not establish a stated global bound.

\subsubsection{Verification and next bounded task}

Revision 42 brought the preparation checks to 32 rejected invalid records.
Its six added cases detect a norm/squared-norm substitution, loss of the strict
interior start, discarded tensor compactness, strict positivity extended
to time zero, the missing eigenvalue factor two, and a record projection
treated as an existence proof. They check the recorded contracts rather
than prove the mathematics. All 85 revision-42 discovery hashes and
declaration lines were checked read-only; revision 43 rechecks the enlarged set. Six actual core-only Lean template runs were rerun;
they use \code{Init}/\code{Nat.add_zero}. The previously checked K0 suite is
unchanged. No PC-dependent draft or production GC build was run.

After the release gate, rebind P67--P85, run Q12--Q15 and their axiom probes,
and discharge the applicable convention and consumer requirements.
Section~\ref{sec:m2a-event-history} now locates the finite-cap event
producer and existing scalar smooth-stage/terminal/event propagation.
The exact topology-event and pinching transports remain comparison work.
M2A/K1 remain pending; the first genuinely new mathematics remains the
surviving M3 Ricci-flow PDE gaps.

\subsection{Inherited finite-cap events and scalar history propagation}
\label{sec:m2a-event-history}

Revision 43 follows the existing PC producers past local cap insertion.
P86--P97 add twelve source discoveries; T17 and T18 refine O08 without
adding a production theorem or declaring a new gap. That revision had
97 candidates, 291 generated declaration probes, eighteen targets and
seventeen drafts with 27 applications. Section~\ref{sec:m2a-hi-bridges}
records the enlarged revision-44 bundle. All PC-dependent files remain unrun. The incomplete snapshot is still
read-only discovery material.

\subsubsection{Concrete finite-cap construction and the history record: T17}

P87, \code{exists_uniform_terminal_cut_cap_events}, is an actual
construction in \code{Surgery/StandardCap/UniformTerminalCutCapEvents.lean}.
It chooses a recentering constant $c\geq4$, positive derivative constants
$C(j)$, and a small positive cap parameter $A$ with $2A<1/2$.
For each model radius $D>0$, order $m$ and accuracy $\varepsilon>0$, it
then provides a positive threshold $\delta_0<1/4$. Its geometric inputs
include the following.
\begin{itemize}
\item A compact nonempty smooth boundaryless three-dimensional incoming
manifold, orientation, actual solution on $[a,s)$ with $a<s$, and a
supplied terminal metric on the terminal regular region together with
\code{IsTerminalLimitMetric}. The theorem does not construct that limit.
\item A finite family of normalized necks of order $m+6$, with
$0<\delta_i\leq\delta_0$, pairwise disjoint ambient neck images, and a
selected retained core contained in the terminal regular region.
\item A nonempty cutting index or a nonempty discarded core. Preserve
this disjunction; it is not a condition that the output be nonempty.
\end{itemize}

The output includes recentered order-$m+4$ data, canonical static insertion
witnesses, their additional properties, and a retained metric $g_{\rm ret}$
identified with \code{finiteFullPreparedMetric}. It also includes oriented
finite cut-cap geometry, witness/collar identities,
\code{FullMetricConclusions}, and P86, \code{MetricCutCapEventOn}.
Thus the scalar-floor and fixed-region conclusions are already produced
on this particular retained metric, subject to the theorem's inputs.
The recentering constant $c$ here is unrelated to the scalar-barrier
constant used below.

Three contracts must be compared explicitly:
\begin{longtable}{@{}P{0.12\tablebodywidth}P{0.36\tablebodywidth}P{0.52\tablebodywidth}@{}}
\textbf{ID} & \textbf{Located contract} & \textbf{Exact boundary}\\\hline
\endhead
P86--P87 & Concrete quotient event and its producer. & \code{MetricCutCapEventOn} records the incoming equation, terminal convergence, oriented and unchanged-locus geometry, retained tensor identity, and uniqueness of the metric when the retained output is empty.\\
P91 & Stage-indexed topology event in \code{EventData.lean}. & \code{MetricCutCapEvent} has discarded/capped stages, a smooth topological transition, a compact old locus, terminal/output maps with tensor identity, and an every-child-meets-old condition. It is a distinct type.\\
P88--P90 & Scale-frontier conversion to a cutoff record. & The supplied frontier already contains both preservation fields, singularity, backward necks, protected/retained conditions, one retained side and static cap data. Conversion is conditional on all those fields.\\
\end{longtable}

In particular, \code{GeometricCutoffScaleFrontier.toGeometryFrontier}
derives the canonical radius time bound from the backward-neck data and
copies \code{curvature_preserving} and \code{scalar_preserving}.
Its \code{nonempty_geometricCutoffRecord} theorem turns a nonempty full
frontier into a nonempty cutoff record. Q17, \code{CutoffFrontier.lean},
is one explicit application of this conditional theorem. Neither the
conversion nor its draft supplies the input preservation proofs.

The source search located the concrete event producer and the conditional
frontier conversions, but has not yet identified a direct construction
from P86 to the exact stage-indexed event and cutoff record. This is a
bounded comparison still to perform, not evidence of a missing theorem
in the completed release. Trace the existing finite quotient/orientation
presentation bridges, the unchanged locus and child condition, and the
metric/predicate transports before scheduling any GC adapter. The
one-retained-side and singular/protected/backward requirements remain
separate from local cap pinching.

\subsubsection{The scalar smooth--terminal--event composition already exists: T18}

The source file
\code{Extinction/Families/SlabScalarLowerBarrier.lean} contains a complete
finite-history scalar-barrier route. Its public signatures require no
finite-time extinction assumption. Reusing these results leaves finite-time
extinction outside the default GC route. Their transitive imports still
need the ordinary completed-release and axiom audit.

Fix $c>0$ and use one absolute time coordinate throughout the history:
\[
 b_c(t)=-\frac{3}{2(t+c)}.
\]
The source's \code{OrientedThreeStage} carries a compact Hausdorff smooth
three-manifold and orientation. Its incoming slab includes the actual flow,
equation and smooth-up-to data. Its \code{TerminalLimitMetric} includes
local convergence of all metric derivatives on compact subsets of the
terminal regular open set. These are substantive inputs, even when they
are packaged in structures instead of written as separate hypotheses.

\begin{longtable}{@{}P{0.12\tablebodywidth}P{0.36\tablebodywidth}P{0.52\tablebodywidth}@{}}
\textbf{ID} & \textbf{Existing step} & \textbf{Input and conclusion}\\\hline
\endhead
P92 & Smooth-slab barrier. & For $a\geq0$, actual solution, $[a,s)$ in the carrier, $(a,s)$ regular, and initial $R(a)\geq b_c(a)$, obtains $R(t)\geq b_c(t)$ for $a\leq t<s$. The source builds the comparison regularity internally.\\
P93 & Incoming terminal passage. & With $a<s$, the incoming slab and its supplied terminal limit metric, the initial floor implies terminal scalar $R_{\rm term}\geq b_c(s)$ on the terminal regular open set. It does not evaluate the flow at an excluded carrier time.\\
P94 & Surgery jump to the next initial metric. & Applies the supplied cutoff record's scalar preservation at $L=b_c(s)\leq0$, then uses \code{ObservedHistory.event_output}. The next initial bound concerns the outgoing stage.\\
P95 & Induction over all stages. & Given a cutoff record at every event and the original initial bound, uses \code{event_initial}, terminal passage and the event jump to establish the bound at every stage's initial time.\\
P96--P97 & Combined history and terminal export. & \code{historyScalarLowerBound_of_history} returns \code{HistoryScalarLowerBound} and \code{TerminalScalarLowerBound}. The first predicate explicitly excludes event times; the second concerns incoming terminal points.\\
\end{longtable}

The preferred combined export therefore needs only the exact supplied
history, $c>0$, its initial floor, and one \code{GeometricCutoffRecord} per
event. The smooth final slab is handled on a closed interval; when the last
stage begins at the horizon, the proof uses its initial value. Finite
induction also permits a history with zero events. No global history or
cutoff existence theorem is inferred from this conditional propagation.

Q16, \code{ScalarHistory.lean}, contains three explicit inherited
applications: incoming terminal transfer, all stage-initial bounds, and
the combined history/terminal export. The draft preserves the terminal
metric, convergence and cutoff inputs. Event-time values must be obtained
from the incoming terminal and outgoing initial results separately; the
predicate \code{HistoryScalarLowerBound} itself cannot be applied there.
For comparison with the initial normalization $R(0)\geq-6$, the proposed
scalar specialization is
\[
 c=\frac14,\qquad b_c(t)=-\frac{6}{1+4t}.
\]
This elementary rewrite still needs its release-bound Lean application and
initial normalization check. Do not reset $c$ or the time origin at a
surgery. Nor does the scalar result establish Hamilton--Ivey history
propagation: its curvature-operator predicate and terminal passage must
be compared independently.

\subsubsection{Verification and next comparison}

The two new drafts contain four applications, all unelaborated. The
preparation validator checks their imports, declaration lists, cross-links
and pending states. Revision 43 brought the checks to 37 rejected invalid
records; its five additional controls detect a frontier conversion misreported as a
preservation producer, a concrete event misidentified as a history cutoff
record, use of the scalar history predicate at event times, omitted terminal
convergence, and an extinction assumption invented from the source path.
All 97 revision-43 source hashes and declaration lines were checked
read-only; revision 44 rechecks the enlarged set. Six
actual core-only \code{Init}/\code{Nat.add_zero} template runs passed their
expected outcomes, including rejection of the deliberately wrong type.
These checks establish neither PC elaboration nor geometric proofs.

After \code{PCReleaseGate}, rebind P86--P97, run Q16--Q17 and the applicable
convention/axiom checks, then accept only consumers whose exact prerequisites
are discharged. Section~\ref{sec:m2a-hi-bridges} locates more inherited
Hamilton--Ivey identities/consequences and the forward event conversion,
while retaining the exact terminal and presentation comparisons. The scalar
propagation is already located; it need not be reproved. The production inventory remains at 136
rows, all 26 consumers stay disabled, and M2A/K1 remain pending.

\subsection{Hamilton--Ivey reuse and event representation bridges}
\label{sec:m2a-hi-bridges}

Hamilton--Ivey is part of the inherited PC project. Revision 44 locates
additional public forms and comparison identities already in that source;
it does not schedule a second proof of Hamilton--Ivey. P98--P111 add fourteen
candidates. T19--T21 refine the normalization, complete/limit and event
representation comparisons. The current bundle has 111 candidates, 333
generated probes, twenty-one targets and twenty consumer drafts with 35
applications. All 353 PC-dependent files remain unrun pending the completed
release. These counts are preparation coverage, not accepted proof evidence.

\subsubsection{The normalization algebra is already proved in PC: T19}

The earlier parameter rewrite is implemented by P99,
\code{two_mul_hamiltonIveyBarrier_eq_fixedHamiltonIveyBarrier}, in
\code{Geometry/Curvature/HamiltonIveyRegion.lean}. Writing
\(B_a(X)=X(\log(aX)-3)\) and using the source's sectional barrier
\(H_{K,\tau}\), its exact equality is
\[
 2H_{K,\tau}(X/2)=B_{(2K)^{-1}+\tau}(X),
 \qquad K>0,\quad \tau\geq0,\quad X>0.
\]
P98, \code{mem_fixedHamiltonIveyRegion_scale_iff}, also already supplies
\[
 (R/c,\nu/c)\in\mathcal H_{ca}
 \quad\Longleftrightarrow\quad
 (R,\nu)\in\mathcal H_a,\qquad c>0,
\]
where \(\mathcal H_a=\{(R,\nu):\nu\geq0\text{ or }B_a(-\nu)\leq R\}\).
Thus the region parameter multiplies by the metric scaling factor. The
identity alone does not identify two metrics or transport the event's
all-bases Rayleigh expression.

The public file \code{DimensionThree/HamiltonIvey/Consequences.lean}
contains P100, \code{hamilton_ivey_pinching_scaled}. On the same compact
smooth three-dimensional slab as the earlier public theorem, set
\(\nu=2\lambda_{\min}\). From the initial bound \(\nu(t_0)\geq-C\),
with \(C>0\), it gives
\[
 R(t)\geq(-\nu)\bigl(\log((C^{-1}+t-t_0)(-\nu))-3\bigr)
 \quad\text{when }\nu<0.
\]
The slab is closed in the carrier, with open regular interior, and the
actual \code{IsSolutionOn} hypothesis is retained. The source derives the
formula using the earlier sectional bound with \(K=C/2\).

P101, \code{hamilton_ivey_pinching_normalized}, specializes to \(C=1\).
Its initial hypothesis is \(\nu\geq-1\), equivalently
\(\lambda_{\min}\geq-1/2\). This differs from P76's sectional
\(K=1\) hypothesis \(\lambda_{\min}\geq-1\). Keep the normalization
attached to every declaration and fixture; the word ``normalized'' is
insufficient to identify a hypothesis.

Q18, \code{HamiltonIveyNormalization.lean}, contains explicit applications
of the fixed-region scaling identity, the factor-two barrier identity and
the scaled public pinching theorem. The identities already exist in PC;
the pending work is release elaboration and their exact application to GC.
For P82, still identify its \code{sInf} of all unit-vector Rayleigh values
with \(\nu=2\lambda_{\min}\) for every orthonormal basis. P80's matrix
eigenvalue equality and P99's logarithmic equality address different parts
of that comparison.

\subsubsection{Complete slabs, ancient flows and rescaled limits: T20}

P102, \code{hamilton_ivey_inequality_on_interval_of_complete}, is already
available in \code{DimensionThree/HamiltonIvey/Complete.lean}. For
\(a<b\), \([a,b]\) in the carrier, \((a,b)\) regular, and complete
slices throughout \([a,b]\), it gives at \(b\)
\[
 R(b)\geq(-\nu)\bigl(\log((b-a)(-\nu))-3\bigr)
 \quad\text{if }\nu=2\lambda_{\min}(b)<0.
\]
Its public category is a smooth boundaryless Hausdorff sigma-compact
three-manifold with a finite-dimensional model and actual solution.
Neither an initial curvature lower bound nor a uniform curvature bound is
an argument of this public theorem. Completeness is required. The proof
handles a possibly nonregular right endpoint by continuity, but that
endpoint must still belong to the carrier. This is not yet a statement
about an excluded surgery endpoint and a different terminal metric.

The following further consequences are already located:
\begin{longtable}{@{}P{0.12\tablebodywidth}P{0.36\tablebodywidth}P{0.52\tablebodywidth}@{}}
\textbf{ID} & \textbf{Inherited result} & \textbf{Scope retained}\\\hline
\endhead
P103 & \code{curvatureOperator_nonnegative_of_complete_ancient}. & Every time in $(-\infty,t]$ belongs to the carrier, $(-\infty,t)$ is regular, and all past slices are complete. The curvature operator at $t$ is nonnegative.\\
P104 & Numerical rescaled upper-limit theorem. & Scalar values converge, scales tend to $+\infty$, ages tend to a strictly positive limit, and the logarithmic bound holds eventually. The eigenvalue input is the one-sided comparison $\lambda_i<\lambda_\infty+\varepsilon$ eventually for every $\varepsilon>0$.\\
P111 & Geometric one-sided eigenvalue comparison. & \code{SmoothCGHConverges} supplies the eventual comparison at pulled source points for each limit point and carrier time. Retain the spatial maps, times and eigenvalue normalization.\\
P105 & \code{curvatureOperator_nonnegative_of_closed_blowup_limit}. & For a compact source flow on $[0,T)$, positive scales tending to infinity and centers tending to $T>0$, supplied convergence on every admissible compact time window gives nonnegative curvature operator on the limit carrier.\\
\end{longtable}

P105 uses tail shifts to place each compact window in the rescaled source
domains. It assumes the limit and \code{SmoothCGHConverges}; it does not
construct compactness or a limit. Its proof obtains a positive initial
operator-bound constant from compactness internally. It is a blow-up
limit theorem, not an identification of the unrescaled incoming terminal
metric. P104's one-sided comparison should likewise not be reported as
full eigenvalue convergence.

Q19, \code{HamiltonIveyCompleteLimits.lean}, gives three explicit
applications: the complete-slab estimate, complete ancient nonnegativity,
and the numerical rescaled-limit lemma. P105/P111 have separate generated
signature/axiom probes. The release audit must inspect their full implicit
instances and compact-window data before selecting the precise GC consumer.

For surgery histories, keep three comparisons distinct: identification of
the Rayleigh and intrinsic predicates; transfer to the supplied terminal
regular metric using its convergence; and next-stage propagation of the
same fixed-region parameter with one absolute time origin. Restarting
P102 at each surgery only supplies its local age $b-a$ and does not by
itself establish a previously specified global-age bound. The exact
history wrapper remains a comparison question; none of this says that
Hamilton--Ivey is missing from PC.

\subsubsection{The event conversion exists in a specific direction: T21}

\code{Surgery/Topology/EventBridge.lean} supplies P107,
\code{MetricCutCapEvent.toSurgery}. It converts a topology event on an
oriented stage into a surgery event on a closed oriented manifold.
The input is the already supplied topology event P91, not P86's concrete
finite-quotient predicate. P106, \code{SmoothCutCapCompletion}, supplies
additional capping/core identities, component and complement information,
and orientation compatibility. A separate smooth embedding of the old
locus into the capped stage is also required. The conversion transports
the old locus and retains the incoming flow, terminal metric and output
metric; its existence does not erase those inputs.

P109/P110 state the converted incoming-initial and outgoing-initial metric
identities. P108,
\code{ObservedHistory.toSurgeryFiniteSurgeryHistory_of_cutCapCompletion},
packages the converted events into the other finite-history type. It
requires \(0<\mathrm{eventCount}\). The original \code{ObservedHistory}
permits zero events; preserve that branch rather than applying this
positive-count constructor universally.

Q20, \code{EventRepresentation.lean}, applies the existing output-metric
identity and the finite-history constructor with its completion, embedding
and positive-count inputs. This establishes a concrete reuse target for
representation conversion. It supplies neither analytic curvature
preservation nor a reverse construction from the concrete quotient event.
Continue tracing the finite-cap selected diffeomorphism, orientations,
unchanged locus and child conditions into the topology-event input. These
are type/presentation comparisons, not new foundational geometry.

\subsubsection{Verification and the next exact comparison}

All 111 recorded source hashes and declaration lines were checked read-only.
Three new draft files contain eight applications, with no local proof
placeholders. Their full signatures retain nested named arguments and
local instance inputs. A new delimiter check and a deliberately damaged
draft regression guard against truncated signatures; this is a limited
source check, not Lean elaboration. Candidate IDs now extend beyond P99
without renumbering old records. The unknown-reference test uses a fresh
absent ID, and noncanonical IDs remain rejected.

There are now 47 preparation rejection tests. The ten additions cover the
ID boundary, draft delimiters, scaling/normalization, completeness and
carrier endpoints, blow-up versus terminal metrics, conversion direction
and positive event count. Six core-only Lean template runs retain their
expected outcomes. All PC-dependent drafts remain unrun; the source
snapshot and the checked K0 suite are unchanged.

After \code{PCReleaseGate}, use the located identities and public exports
before proposing any adapter. Revision 45 defers the Rayleigh-to-intrinsic,
terminal, next-stage and concrete-event adapter comparisons until that
release or an explicitly stabilized interface is available. Independent
collapse contracts are now the active work in
Section~\ref{sec:independent-collapse-contracts}. M2A/K1 and all 26 production
consumer records remain pending; the theorem inventory still has 136 rows.

\subsection{Release-gated M2A workflow}
\label{sec:m2a-release-workflow}

The transitions are recorded in \code{m2a_release_workflow.csv}.
Current state is \code{snapshot-only}: no completed-PC commit, clean root
build, accepted public theorem, or adapter evidence is recorded. K0 success
does not change this state.

\setlength{\tablebodywidth}{\dimexpr\textwidth-4\tabcolsep\relax}
\begin{longtable}{@{}P{0.209302\tablebodywidth}P{0.453488\tablebodywidth}P{0.337209\tablebodywidth}@{}}
\textbf{Transition} & \textbf{Evidence required} & \textbf{Permission gained}\\\hline
\endhead
Prepare & Candidate rows, consumers, probe templates, proposed axiom policy. & K0 and audit preparation; available now.\\
Pin candidate & Resolved immutable commit, source/lock hashes, Lean/mathlib pins, public import candidates. & Reproduce candidate build in an isolated checkout.\\
Pass \code{PCReleaseGate} & Clean root build, public theorem/signature, scoped project gap scan, transitive axioms, reviewed policy. & Run row probes and convention tests on this release.\\
Classify rows & Minimal imports, exact signatures, hypothesis/convention comparison, axioms tied to the release. & Direct imports, named adapter tasks, gaps, or exclusions.\\
Close adapters/tests & Proved bridges, wrapper probes, convention regressions and axiom reports. & Review M2A acceptance and consumer readiness.\\
Accept M2A / open K1 & Every row classified, required foundation/bridges proved, scoped gaps and no unmet prerequisite for the enabled consumer. & Production M1 and M3 gap implementation.\\
\end{longtable}

A pinned candidate alone is not an accepted release. The future evidence
root is \code{audit/pc/<resolved-commit>/}, containing a release manifest,
root-build log, public-signature report, project-gap scan, per-declaration
imports/axioms, convention tests, adapter results, and a completion record.
These are planned paths, not existing artifacts. Each report records the
command, exit code, source identity, toolchain, and content hash.

\subsubsection{Row evidence and decision rules}

\code{m2_foundation_api.csv} separates the provisional
\code{adapter_or_gap} expectation from final \code{classification}.
It adds \code{audit_state}, \code{release_evidence_id},
\code{declaration_evidence_ids}, \code{adapter_evidence_ids},
\code{gap_ids}, and \code{consumer_gate}. All final classifications and
evidence IDs remain empty. Candidate paths remain discovery notes.
The broader \code{pc_reuse_audit_template.csv} likewise distinguishes
expected \code{gc_class} from an empty final \code{classification}.

The final classes remain \code{direct-reuse}, \code{adapter-reuse},
\code{GC-new-gap}, and \code{out-of-scope}. Split a concept row into linked
declaration-level child rows before accepting mixed outcomes.
Direct reuse requires exact compatible signatures and axiom evidence.
Adapter reuse additionally needs a proved named bridge and its axiom report.
A gap requires the exact missing statement, comparison/search scope,
prerequisites, consumer, and owning milestone. An exclusion requires a
reason and evidence that no enabled consumer needs it.
A failed build or an unlocated declaration remains an unresolved audit
result; it does not establish a mathematical gap.

\subsubsection{Non-circular gates and scoped acceptance}

\code{M2AConventionRegression} runs after \code{PCReleaseGate}, during M2A.
It provides evidence for K1 and therefore must not wait for K1.
It checks inherited operators and small bridges without requiring production
assembly or the completed eight-model certificate suite.

If PC lacks a public hyperbolic model, record a named M1/M6 model-test gap.
M2A may accept consumers whose foundation/convention obligations are fully
tested with the available sphere and flat models; hyperbolic consumers
remain gated on their missing test. A general test-success entry cannot
hide the omission. A missing metric-family or PDE identity similarly becomes
an exact M3 target, not an imported fact.

Inventory completeness differs from proof availability. A classified gap
can close an inventory decision without closing its theorem. Completion
records must identify consumer-specific readiness rather than treating K1
as permission to use missing hypotheses.
Spherical-space-form recognition needs its own release theorem or subsequent
proof obligation; the simply connected PC conclusion alone is insufficient
evidence for that broader adapter surface.

\subsubsection{Invalidation}

Changes to the PC commit, source/lock hashes, toolchain, import surface,
signature, or axiom policy mark affected evidence \code{stale}. Preserve old
reports and rerun the release gate and dependent probes. Do not rebuild the
inherited geometric library. K0 proofs remain independent sandbox evidence;
rerun them under the eventual GC toolchain before using them as regressions
for that build.
% END migrated B030

\chapter{References, source policy and reading records}
\label{ch:references-appendix}

% BEGIN migrated B031
\section{Literature map and source policy}
\label{ch:literature-map}

The project will not try to cite every paper related to Ricci flow or
3-manifold topology. It will maintain a curated corpus whose entries are tied
to theorem contracts. A reference enters the core corpus when it supplies a
mathematical result, a proof decomposition that we follow, a definition needed
by the endpoint, or an implementation/API decision. Background references can
remain in a secondary bibliography without becoming proof dependencies.
% END migrated B031

% BEGIN migrated B032
\section{Authority levels}

\begin{description}[leftmargin=3.0cm,style=nextline]
\item[Proof source] The source of the proof actually followed. Record original
  attribution separately. An original announcement and a later complete proof
  can have different roles; both Morgan--Tian and Kleiner--Lott supply proofs
  of collapse results, not merely expositions of Perelman's announcement.
\item[Detailed exposition] A complete or expanded proof used to plan Lean
  lemmas and locate hidden dependencies. A work may also be the selected proof
  source. Any change in hypotheses or conclusion requires an explicit bridge.
\item[Topology and analysis support] A source for an independent theorem such
  as collapse, compactness, derivative estimates, JSJ topology, or geometric
  classification.
\item[Implementation reference] Mathlib documentation, the completed
  Poincar\'e repository, and the Hamilton formalization. These guide APIs and
  dependency boundaries; they are not mathematical authority.
\end{description}

For every formal theorem, the reference inventory records a selected proof source,
original attribution, an optional comparison source, the exact locator, the
Lean module that consumes it, and the current status. If two expositions use
different conventions, the discrepancy is recorded and one normalization is
chosen explicitly in the blueprint.
% END migrated B032

% BEGIN migrated B033
\section{Core source map}

\setlength{\tablebodywidth}{\dimexpr\textwidth-4\tabcolsep\relax}
\begin{longtable}{@{}P{0.284091\tablebodywidth}P{0.329545\tablebodywidth}P{0.386364\tablebodywidth}@{}}
\textbf{Source} & \textbf{Authority role} & \textbf{GC use}\\\hline
\endhead
Hamilton (1982) & Primary analytic source & Ricci flow, curvature evolution,
maximum principles, and positive-Ricci model behavior.\\
Hamilton (1995); Shi (1989) & Analytic support & Flow compactness and derivative
estimates used by blow-up and canonical-neighborhood arguments.\\
Hamilton (1999) & Primary nonsingular-flow source & Classification of bounded-
curvature normalized nonsingular solutions; positive, zero, and negative
limits; hyperbolic limits and hyperbolic pieces.\\
Perelman (2002) & Primary analytic source & $\mathcal W$-entropy, reduced
length/volume, monotonicity, no-local-collapsing, and curvature-scale control.\\
Perelman (2003), surgery & Primary construction source & Canonical neighborhoods,
neck/cap models, surgery scales, continuation, and surgery invariants.\\
Perelman (2003), extinction & Optional route source & Finite extinction of
the non-aspherical branch; consult only if the selected GC route requires it.\\
Kleiner--Lott (2008) & Detailed exposition & Main locator map for Perelman's
compressed arguments and the hidden analytic hypotheses to expose in Lean.\\
Morgan--Tian (2007) & Detailed exposition & Flow, surgery, and the Poincar\'e
endgame; use for theorem decomposition and proof order.\\
Morgan--Tian 2014 & Global GC spine and assembly authority &
Long-time flow, stabilization, the topological endgame, and comparison with the
local-collapse route.\\
Morgan--Tian 2008 & Archived comparison source &
Use only after a stable archive, URL, and hash are recorded; it is not current
release evidence.\\
MSM206, Chapters 31--34 and Appendix K & Detailed bridge exposition &
Hyperbolic 3-manifold prerequisites, Hamilton's nonsingular classification,
noncompact hyperbolic limits, boundary incompressibility, harmonic-map
stability, and the analytic/topological bridges used by the hyperbolic branch.\\
Kleiner--Lott (2014), Ast\'erisque & Provisional selected local-collapse proof &
\emph{Locally collapsed 3-manifolds}; boundary version and application to
Ricci-flow time slices. Distinct from the 2008 \emph{Notes}.\\
Shioya--Yamaguchi (2005) & Related collapse theorem & Global lower-curvature,
small-volume setting; compare hypotheses before importing into a local
collapse argument.\\
Cheeger--Fukaya--Gromov (1992) & Conditional support & Two-sided bounded
curvature collapse. Not a substitute for a local lower-curvature theorem.\\
Thurston; Scott (1982, 1983) & Primary geometric background & The eight model
geometries, geometric structures, Seifert pieces, and the target vocabulary.\\
Benedetti--Petronio (1992) & Hyperbolic-geometry reference & Models of
$\mathbb H^3$, discrete groups, cusps and Margulis theory, finite-volume
hyperbolic 3-manifolds, finite-volume rigidity used by stabilization, and the topology needed to state
the hyperbolic-piece certificate.\\
Hatcher and related 3-manifold notes & Background exposition & Definitions and
elementary topology for the prime, incompressible, and gluing interfaces.\\
\end{longtable}
% END migrated B033

% BEGIN migrated B034
\section{Provisional hybrid exposition policy}
\label{sec:hybrid-exposition-policy}

The GC formalization does not follow one secondary exposition uniformly. Its
provisional route is a deliberate hybrid with one source assigned to each
job. Morgan--Tian 2014 is the global Geometrization spine: it determines the
long-time-to-geometric-pieces narrative and the final topological assembly.
Kleiner--Lott 2008 is the principal analytic expansion and locator guide for
Perelman's compressed estimates, compactness, and surgery arguments. Perelman's
original papers remain the primary authority whenever a secondary exposition
and the original statement differ.

\setlength{\tablebodywidth}{\dimexpr\textwidth-6\tabcolsep\relax}
\begin{longtable}{@{}P{0.219780\tablebodywidth}P{0.285714\tablebodywidth}P{0.340659\tablebodywidth}P{0.153846\tablebodywidth}@{}}
\textbf{Layer} & \textbf{Default source} & \textbf{Formalization role} & \textbf{Status}\\\hline
\endhead
Global proof spine and assembly & Morgan--Tian 2014, checked against Morgan--Tian 2008 & Organize the long-time decomposition, geometric-piece output, exceptional branches, and final reconstruction. & Provisional backbone\\
Analytic estimates and surgery & Perelman 2002/2003, expanded with Kleiner--Lott 2008 and Morgan--Tian 2007 & Expose exact hypotheses for entropy, reduced volume, noncollapsing, canonical neighborhoods, surgery events, continuation, and invariants; import any completed PC theorem rather than reproving it. & Primary plus exposition\\
Hyperbolic branch & Hamilton 1999 and Perelman 2002, bridged by MSM206 Chapters 31--34 and Appendix K & Convert normalized thick limits into embedded finite-volume hyperbolic pieces with cusp and incompressible-boundary data; use Benedetti--Petronio for definitions and standard hyperbolic geometry. & Provisional candidate / locator-pending\\
Collapse branch & Kleiner--Lott 2014 Ast\'erisque boundary theorem as the initial candidate, compared with Morgan--Tian 2014 & Choose one theorem with explicit lower-curvature, derivative, boundary, and quantifier hypotheses and convert its output to graph-manifold data. & Candidate pending audit\\
Topological vocabulary and reconstruction & Scott, Thurston, Hatcher, and Morgan--Tian 2014 & Define prime pieces, incompressible tori, Seifert/graph-manifold interfaces, model tags, and gluing certificates. & Foun\-dation\slash assembly\\
\end{longtable}

This division is not a claim that Morgan--Tian and Kleiner--Lott prove the same
intermediate theorem in interchangeable form. A source enters the Lean
dependency graph only after its exact statement, hypotheses, normalization,
and locator have been recorded. Shioya--Yamaguchi and Cheeger--Fukaya--Gromov
remain comparison or support sources unless the selected collapse theorem
actually invokes one of their results.

\subsection{Source-selection rule}

For each planned theorem, the inventory records five separate roles when they
occur: the source of the mathematical statement, the source that expands the
proof, the source that supplies geometric definitions, the source that provides
an independent comparison, and the source of the Lean declaration being
reused. These roles may be filled by different references. The proof source
never silently becomes a Lean dependency: the theorem must still be imported,
proved, or explicitly identified as a new GC obligation.

A literature route is promoted from provisional to selected only when the
comparison report resolves the following questions:

\begin{enumerate}[label=\textbf{L\arabic*:},leftmargin=*]
\item What exact theorem produces the hyperbolic thick-piece certificate?
\item Which collapse theorem supplies the thin graph-manifold output, and on
  what domains do its curvature, derivative, volume, and boundary hypotheses
  hold?
\item How are subsequences, regular-time slices, surgery histories, and
  boundary labels transported back to the original manifold?
\item Which exceptional or terminal data are exported by the \code{PoincareAdapter},
  including any spherical-space-form recognition needed by the endpoint?
\item Which source locators, errata, normalization translations, and Lean
  declarations are attached to each theorem contract?
\end{enumerate}

Until these answers are fixed, the source is a comparison input and the
corresponding theorem inventory row remains \code{chapter-scaffold} or
\code{locator-pending}; it is not a temporary axiom.

\subsection{Hybrid compatibility gate}

The provisional hybrid route is promoted only after four checks pass:

\begin{enumerate}[label=\textbf{H\arabic*:},leftmargin=*]
\item the normalized C7 surgery slices satisfy the lower-curvature,
  derivative, boundary, and volume hypotheses required by the selected KL14
  local-collapse theorem;
\item the MT14 finite-stabilization and transport statements are stated
  independently of the KL14 conclusion they are used to organize;
\item all scale, curvature-sign, and time-normalization translations between
  Perelman, KL14, MT14, and the inherited PC API are named adapters; and
\item the comparison report records why the KL14 route is preferred for the
  local theorem and what MT14 supplies for its global application and assembly.
\end{enumerate}

Until H1--H4 pass, KL14 is a candidate proof source and MT14 is a global
spine and comparison source. Neither citation substitutes for a Lean theorem.

\subsection{Confirmed local source manifest}
\label{sec:local-source-manifest}

The confirmed local evidence root is
\code{/Users/bennettchow/Documents/Codex/Geometrization/BooksPapers}. The
following files are confirmed there and are the local working copies for the
main source stack. The machine-readable \code{reference_inventory.csv} records the
archive-relative paths, and \code{source_manifest.csv} records paths relative to the
Geometrization project root using the \code{BooksPapers/} prefix, together with
bytes and SHA-256 hashes for the archived files. Directory hashes are computed
from sorted relative filenames and their per-file SHA-256 digests. Lean builds
must never import from this absolute user path.

\setlength{\tablebodywidth}{\dimexpr\textwidth-4\tabcolsep\relax}
\begin{longtable}{@{}P{0.284091\tablebodywidth}P{0.409091\tablebodywidth}P{0.306818\tablebodywidth}@{}}
\textbf{Reference role} & \textbf{Confirmed local file} & \textbf{Use}\\\hline
\endhead
Morgan--Tian global spine & \code{MorganTianGeom.pdf} & Global comparison archive; body dated August 7, 2012, distinct from prefixed book contents.\\
Morgan--Tian Poincare exposition & \code{MorganTianPoincare.pdf} & Flow, surgery, and inherited PC endgame.\\
Kleiner--Lott analytic notes & \code{KleinerLottPerelman.pdf} & Expanded analytic and surgery proof contracts.\\
Kleiner--Lott collapse paper & \code{KleinerLottAsterisqueLocalCollapse.pdf} & Candidate local-collapse theorem and boundary output.\\
Perelman originals & \code{GrishaPerelman1.tex}, \code{GrishaPerelman2.tex}, \code{GrishaPerelman3.tex} & Primary entropy, surgery, and optional extinction statements.\\
Hamilton/geometry bridge & \code{MSM206.tex} & Hamilton 1999 classification and hyperbolic-limit bridge.\\
Hyperbolic geometry & \code{BenedettiPetronioHyperbolic.pdf} & Hyperbolic models, cusps and Margulis theory; finite-volume
Mostow--Prasad rigidity is required by the Chapter 8 stabilization route.
Uniqueness of the final decomposition remains unnecessary.\\
3-manifold topology & \code{Hatcher3ManifoldNotes.pdf} & Prime, incompressibility, and gluing interfaces.\\
Geometric structures & \code{Scott1983Geometries.pdf} & Eight geometries, Seifert pieces, and geometric classification.\\
Collapse comparison & \code{ShioyaYamaguchi2005MA.pdf}; \code{CheegerFukayaGromov1992.pdf} & Comparison and support only unless promoted by the route audit.\\
Independent cross-check & \code{CaoZhu2006AJM.pdf} & Alternative exposition and convention comparison.\\
\end{longtable}

The manifest is a source map, not a claim that every file is a theorem
 dependency. Local copies still require exact page, theorem, or section
locators, and the relevant errata must be recorded. References without a
confirmed local file remain archive or web sources and must receive a stable
URL, version, or hash before becoming release evidence. The contract-level
source roles are recorded separately in \code{hybrid_source_map.csv}; this prevents
one semicolon-separated bibliography field from conflating a theorem source,
a proof expansion, a bridge exposition, and a comparison source.

\subsection{Archive recheck and additional local planning material}

The revision-35 archive check found every manifest entry at the recorded
local root. Revision 50 adds the user-supplied Matveev second edition to
the manifest: 28 individual files now match their byte counts and SHA-256
digests. Both John Lee source trees
match their recorded byte totals. The earlier manifest did not specify the
exact serialization used to combine filenames and digests into its two tree
hashes; those hashes are preserved but are not claimed to have been
revalidated. This remains the nonblocking A16 provenance task.

The additional directory \code{BooksPapers/PoincareConjectureBlueprint/}
contains \code{master05a.tex}, \code{master05b.tex}, and \code{references.bib}.
It is now recorded as \code{PoincareBlueprintSnapshot}, for read-only contract
navigation and bibliography discovery. It is neither a Lean source release nor
evidence that \code{PCReleaseGate} has passed. Its split-file format does not
change the one-file GC blueprint policy.

The revision-54 inventory has 45 reference rows: 30 name local paths,
including the separate records for the Morgan--Tian manuscript body and
unverified published-book crosswalk. Seven are bibliographic only, seven are
remote only, and one is the external release-gated PC repository. The additional remote Hatcher 2023 comparison is separate from the
archived notes. The 2003 Matveev entry remains an indirect KL locator.
The separately inventoried \code{Matveev2007} is the user-supplied
\code{BooksPapers/MatveevBook2ndEd.pdf}, whose title/copyright pages
identify the second edition. Its relevant primary passages are now read;
no equivalence of the two editions is asserted. The remote Martelli comparison
uses the author-hosted Version 4 of September 2025, with its own hash and
page locators; it is not an archived book or a replacement global spine.
The corrected author Hatcher algebraic-topology chapter and university-hosted
Hirsch scan are separate remote records used for group and isotopy
applications; their precise versions and reading limits are recorded.
The separately cited Hamilton 1982/1995/1999,
Shi 1989, Thurston 1982, Hamilton formalization reference, and Morgan--Tian
2008 have no standalone archived copies in this directory. A bibliography
entry or a treatment inside another book does not count as that paper's
local copy. The sibling project-level \code{GEOMETRIZATION_BLUEPRINT}
directory is empty; the worktree directory remains authoritative.

Run the optional archive verification by adding
\code{--archive-root /Users/bennettchow/Documents/Codex/Geometrization}
to \code{audit_blueprint.py}. This verifies individual file hashes and
directory byte totals without reading archives into Lean or changing them.
% END migrated B034

% BEGIN migrated B036
\section{How the sources drive the blueprint}

Morgan--Tian 2014 is the global Geometrization spine and assembly guide.
Perelman supplies primary Ricci-flow assertions, expanded analytically by
Kleiner--Lott (2008). Kleiner--Lott (2014) remains the provisional local
collapse route, with Morgan--Tian providing global application/comparison. A provisional choice between their proofs is discussed
below. Hamilton and Shi supply reusable analytic foundations; Scott and
Thurston supply the target geometry vocabulary.
Cao--Zhu is retained as a comparison exposition and cross-check, not as the
sole authority for a formal theorem.

The formalization rule is therefore:
\begin{enumerate}[label=\textbf{R\arabic*:},leftmargin=*]
\item choose a proved statement from the selected proof source, preserving
  attribution and checking its hypotheses against the intended consumer;
\item use the exposition to expand it into Lean-sized lemmas and identify
  omitted hypotheses;
\item record every independent support theorem as a separate dependency;
\item cite the exact source locator in the theorem inventory;
\item never replace a missing Lean proof with ``by reference.''
\end{enumerate}

The local \code{BooksPapers} collection is a working archive of these sources. It
is not itself a proof dependency: the project records stable bibliographic
identifiers and locators so that a clean checkout can reproduce the source
map.
% END migrated B036

% BEGIN migrated B065
\section{Blueprint detail level and source policy}
\label{ch:detail}

\subsection{What belongs in this document}

The global blueprint records the exact public theorem, definitions that are
shared across modules, work-package boundaries, dependency edges, milestone
gates, reference roles, and known risks. It does not reproduce every
calculation in Perelman's proof.

Each module specification records the exact Lean-facing definitions and
theorem statements, every quantified hypothesis and regularity assumption,
normalizations, a proof plan, dependency and status fields, known mathlib gaps,
acceptance examples, and an exact source section/page locator. A citation may
explain why a theorem is expected, but the Lean theorem must state all of its
hypotheses and its final build must contain a proof term.

\subsection{Lean source versus blueprint}

The blueprint does not contain executable Lean source. It may display a short
interface sketch such as
\begin{quote}
\code{def Geometrizes (M : ThreeManifold) : Prop :=
Nonempty (GeometrizationCertificate M)}
\end{quote}
when that makes a contract unambiguous, but the authoritative declaration and
its proof live in a \code{.lean} file. The theorem inventory currently stores proposed identifiers, statement or
role summaries, dependencies and status. It is not yet an exact Lean
signature/source-file inventory for production declarations; those fields
must be bound to actual declarations before release acceptance. This separation
keeps the blueprint readable and lets Lean, rather than a copied code excerpt,
be the test of correctness.

\subsection{Reference maintenance}

Appendix~\ref{ch:references-appendix} is the single source map. Every module records
the proof source actually followed, original attribution, useful comparison
sources, and exact locators. The collapse comparison in
Section~\ref{sec:collapse-comparison} must precede finalizing the corresponding
theorem contracts.
% END migrated B065

\chapter{Development gates, trust and current work}
\label{ch:workflow-appendix}

\section{Current contracts for autoformalization}
\label{sec:autoformalization-protocol}

This section is the current implementation reading guide. It does not
pin mathlib, choose a competing manifold hierarchy, accept a PC release,
or claim new Lean proofs. Older discovery and milestone passages retain
their historical meaning. For a current task, use its precise contract,
the refinements identified here, and the source/errata record together.
A theorem inventory row is a planning unit, not automatically one Lean
declaration. In particular, MC01, MC02 and MC10 must be split into their
individual definitions and lemmas before proof search.

\subsection{Sorts, supplied witnesses and actual producers}
\label{sec:autoformalization-sorts}

The following conventions prevent a representational error from becoming
a mathematical dependency. The model and certificate records live in
\code{Type}, with universe parameters determined by the inherited API.
Their proof fields live in \code{Prop}. A function constructing a record
is a \code{def}; if it selects witnesses from propositional existence,
use a separately named \code{noncomputable def}. A \code{theorem} proves
a proposition, including the specification of a constructor or the
nonemptiness of a certificate type. Do not move actual carrier, map or
metric fields into \code{Prop} to repair an inappropriate declaration
keyword. The historical K0 pattern---a constructor with separate
specification theorems---already illustrates this separation.

\code{LongTimeOutcome M} is the data-level sum specified in
Section~\ref{sec:assembly-outcome-sort}. A proposition asserting the
existence of a global flow or an outcome is not itself a chosen value.
The public proposition can compose nonemptiness proofs without extracting
a record. If a data-returning wrapper is useful, its choice boundary is
explicit and local. Every selected witness is chosen once, stored under
the actual finite prime/piece index, and reused by its consumers.

Keep three questions separate: whether a type can be defined, whether a
conditional lemma can be proved from supplied records, and whether those
records have been produced for the original manifold. GA14 can be a pure
constructor before GA17's inputs have been constructed. This does not
discharge GA17. The production graph's broad arrows therefore cannot be
used as an import graph: for example, the exceptional-geometry producer
supplies an input to the pure constructor, rather than being called by it.
The current typed links distinguish proof, producer, supplied-record,
transport, optional-route and navigation uses.

The 21 external production aliases are unbound support obligations, not
completed leaves. Their empty adjacency in a cycle check says nothing
about proof closure. In particular Hamilton--Ivey remains an inherited
candidate; an unbound identifier is not evidence that PC lacks it.
The exact inherited theorem, complete elaborated type, transitive axioms
and minimal imports are bound only after the release gate.

\subsection{Carrier, metric and transport discipline}

Every packet names its actual carrier, topology/smooth structure, chosen
metric, induced distance and measure, and scale. When a proof uses a
metric together with its rescaling or pullback, keep those roles explicit.
Use the accepted foundation's scoped instances or marked carrier wrappers;
do not install competing global metric instances on the same point type.
Transport the whole dependent package through a named index equivalence
and an actual carrier map. A matching model tag or abstract group
isomorphism is not a replacement for that map.

Local maps are typed on their tested domains. MC09's chosen inverse,
for example, should first land in the subtype of the original closed
$R$-ball, so composition with its partial forward map is well typed.
Record the exact basepoint branch and select each inverse witness once.
For based fundamental groups, retain paths and the resulting commuting
or conjugacy equations; do not silently replace these by equality.
Metric local charts have no orientation-sign requirement, while oriented
cutting and reconstruction retain theirs
(Section~\ref{sec:model-atlas-orientation-audit}).

Model objects and piece objects are different. Complete $\mathbb H^3$
is a hyperbolic model but has infinite volume and is not an E1
hyperbolic piece witness. A conditional piece regression can use a
supplied HG02 finite-volume manifold and the actual whole-interior
comparison. Constructing a particular closed hyperbolic example is a
separate integration task, not a hidden import of the pure constructor.

\subsection{Small first tasks and their acceptance boundary}
\label{sec:autoformalization-first-atoms}

The first 15 packets refine existing arguments. They are ready for
statement review; the environment, exact imports, inherited bindings and
production execution remain deferred. Their mathematical dependencies
are smaller than the contextual dependencies of their parent rows.
No completed atom alone marks its parent complete.

\setlength{\tablebodywidth}{\dimexpr\textwidth-2\tabcolsep\relax}
\begin{longtable}{@{}P{0.26\tablebodywidth}P{0.74\tablebodywidth}@{}}
\textbf{Parent} & \textbf{Atomic task and discriminating obligation}\\\hline
\endhead
MC19 & Positive Lipschitz envelope; justify each supremum's nonempty,
bounded set, without a global bound or a smoothness conclusion.\\
MC02 (three) & Distance-to-set Lipschitz estimate for a nonempty set;
net-to-packing injection; bounded finite packing gives a strict net.
Keep strict separation and empty-carrier cases.\\
MC06 & Approximation data and its radial bound on the actual source ball;
do not turn a partial map into a total one.\\
MC07--MC09 & Restriction, composition and chosen quasi-inverse;
retain all radius margins, larger inverse constants and typed domains.\\
MC10 (two) & Distance rescaling and recentering, with positive scale and
explicit target-basepoint error.\\
AT09 & Injectivity of the first map from an injective composite;
do not conclude injectivity of the second map.\\
GA02 & Finite tagged dependent indices and reindexing; repeated
diffeomorphism types retain distinct component identities.\\
GA03 & Two-element orbits of a fixed-point-free involution;
two sides may belong to one block but remain distinct.\\
AT29/LP04 & One time for a finite family of eventual conditions;
unrelated subsequential witnesses do not suffice.\\
GA10 & One dependent witness family from proved nonempty fibers;
no witness can be selected from an empty fiber.\\
\end{longtable}

The task view in \code{formalization/} derives its catalog from 360
detailed rows and 141 production rows. These 501 entries overlap in
granularity and are not 501 distinct proof jobs. It records the 15 atoms,
five current contract/proof refinements, source rows, reviewed TeX
anchors, historical evidence and qualified upstream obligations. Exact
Lean declarations and imports remain unbound. Use
\code{python3 formalization/tasks.py frontier} for the statement-review
frontier and \code{python3 formalization/tasks.py show MC07.restriction}
for a bounded task. An export with dependencies is a review packet;
it neither executes Lean nor grants a missing release acceptance.

\subsection{Evidence and failure handling}

The acceptance record has separate fields for source/errata review,
written mathematical argument, statement/interface review, accepted
reuse binding, elaboration, axiom closure and a downstream consumer test.
A successful document build or static audit advances none of the Lean
fields. A hash protects a reviewed statement from drift; it does not
prove that statement. Historical checks remain attached to their exact
input hashes, including K0 and the separate revision87 PC audit. The
prepared 333 M2A probe files remain unrun; that fact does not erase the
distinct historical executed audit.

Before implementing a packet, search the accepted libraries for its
exact statement and hypotheses. If a proof attempt fails, preserve the
reviewed statement and classify the failure: missing name/import,
elaboration or instance mismatch, convention/transport mismatch, missing
lemma, mathematical gap, or resource limit. Record the smallest failing
goal. A changed hypothesis, regularity, carrier or conclusion requires
review of every affected consumer and a new packet version. Never make
the proof succeed by inserting its desired conclusion as a structure
field or by silently adding an axiom.

Revision104 adds bounded semantic checks alongside the retained
preservation tests. They reject loss of the fixed-point-free condition,
loss of nonempty witness fibers, and the MC07 radius margin. A singleton
identity involution, an empty fiber, and an invalid error/radius pair
explain why these hypotheses matter. Further checks keep data sorts,
E1 volume, actual presentations, source parameter order and external
binding gates distinct. These are discriminating specification tests,
not a decision procedure for mathematical truth or a substitute for Lean.

\subsection{The remaining bottlenecks and their smallest useful outputs}
\label{sec:autoformalization-bottlenecks}

The highest-value work is to discharge precise interfaces rather than
expand already conditional packaging. The following list is a strategic
ordering, not a measured workload estimate.
\begin{enumerate}[leftmargin=*]
\item Bind the accepted PC release per consumer. Preserve user-confirmed
Hamilton--Ivey and Schoenflies reuse; isolate any surviving M3 PDE gap.
Keep basic statement refinement independent of the release decision.
\item Supply the MC17/LC81 finite-regularity bridge. A finite
$C^K$ bound is not an all-orders smooth bound, and smoothing an atlas
does not itself preserve nonnegative curvature of a metric. State the
exact category needed by each soul, normal-flow and isotopy consumer.
\item Complete FC23's original-domain overlap packet and FC28's
dimension-independent cloudy smoothing. The written algebra and
quantifier conversion do not manufacture these geometric inputs.
\item Complete the FC44 register and H1--H4 application adapters.
The bounded register in Section~\ref{sec:fibration-constraint-register}
exposes mixed scale intervals, common-domain margins, whole-ball
derivative estimates, boundary buffers and all-time isolation. An
acyclic order alone is not a feasible simultaneous parameter choice.
\item Produce essential torus/relative-prime presentations on actual
carriers. Retain elementary certified cuts when possible, using
Section~\ref{sec:retained-bundle-cut-route}; still prove the general
essential-cut, graph progress, smooth marked reconstruction and
remaining Seifert model-existence inputs.
\item Bind finite-volume rigidity and controlled thick-core persistence
in Chapter~\ref{ch:hyperbolic-geometry}. Ambient cusp incompressibility
still requires its global disk/area argument and surgery boundary
controls; it does not follow from intrinsic cusp injection or Mostow.
\item Supply Chapter~\ref{ch:long-time-topic}'s actual common-slice
producers or its complete terminal event/discard record. Then the
already written Chapter~\ref{ch:global-assembly-topic} constructors
consume their outputs without importing late analysis into the terminal
or pure packaging functions.
\end{enumerate}

Chapters 3--4 have substantial written metric/Alexandrov arguments,
with comparison and splitting inputs still qualified. Chapters 13--16
have blueprint coverage and dependency audits, not completed formal
proofs. The current view explicitly supersedes AC17's old claim that
AC16 is open: AC33 supplies that written producer. Its structural
hypotheses and formal binding remain. LC71's earlier radial obligation
is supplied conditionally by LC79; it is not a new task to restart.


% BEGIN migrated B066
\section{Development gates}
\label{ch:gates}

\begin{description}[leftmargin=2.7cm,style=nextline]
\item[M0: conventions and PC release] Freeze the target, categories, signs,
  normalizations, geometry definitions, Lean revision, and the pinned completed
  PC release with its imported axiom inventory.
\item[M1: vertical slice] Define \code{Geometrizes}, the thick/thin and collapse
  contracts, and prove the conditional integration theorem.
\item[M2: inherited foundation] Audit/import the PC manifold, metric, tensor,
  connection, curvature, volume, geodesic, and injectivity-radius APIs; freeze
  conventions; and prove only the small adapters or time-dependent interfaces
  that are genuinely missing.
\item[M3: Ricci-flow PDE] Formalize the Ricci-flow equation, short-time
  existence and uniqueness, evolution identities, and PDE regularity layer that
  are not already supplied by the inherited PC release.
\item[M4: estimates and compactness] Audit/import Perelman functionals,
  monotonicity, noncollapsing, curvature control, derivative estimates, and
  pointed-flow compactness, adding only missing adapters or theorems.
\item[M5: surgery integration] Import and audit the completed PC
  canonical-neighborhood, surgery-event, continuation, and
  invariant-preservation interfaces; connect them to the M3--M4 flow types.
\item[M6: asymptotics] Prove the thick/thin, hyperbolic-limit, and
  graph-manifold collapse package and identify the geometric pieces.
\item[M7: closure] Replace every temporary interface by a kernel-checked
  proof, complete the Poincar\'e adapter, compose M1 and M6, and pass the
  final provenance audit.
\end{description}

Every milestone has an executable Lean target and a reviewable theorem list.
``Chapter mostly complete'' is not an acceptance criterion.
% END migrated B066

% BEGIN migrated B067
\section{Trust boundary and Lean rules}
\label{ch:trust}

Prototype interfaces may be declared in a clearly marked sandbox. The release
target must contain no project-local \code{sorry}, \code{admit}, or unreviewed \code{axiom}.
Each public theorem receives a \code{#print axioms} audit and a source/provenance
entry. The theorem inventory and status ledger are maintained in a separate
machine-readable file so that this document remains readable.

Coordinate calculations are internal proofs of invariant API lemmas. Sign and
normalization choices live in one definition or theorem. Standard metrics are
regression tests for tensor contractions, scaling, curvature signs, and
surgery conventions. Essential invariants are fields of structures or explicit
theorem hypotheses, never comments.
% END migrated B067

\paragraph{Current frontier (revision109).}
Use Section~\ref{sec:autoformalization-protocol} and the derived atomic
packets for statement refinement. The qualified geometric producers
remain those listed in Section~\ref{sec:autoformalization-bottlenecks}.
The current mathematical iteration is
Section~\ref{sec:mostow-top-down}; its subordinate MPR register refines
HG04 without declaring rigidity or persistence proved.
The milestone narrative below is retained history, not a new instruction
to restart M1 or probe the changing PC snapshot.

% BEGIN migrated B068
\section{Current editorial work and later implementation}
\label{ch:first-task}

The finite K0 tranche is complete, including its assembly, mixed fixture,
negative refutations, and source review. Its runner is
\code{sandbox/check_k0.py}; the static audit validates the recorded evidence.
The W0 discovery audit and production M1a--M1f specification are in place,
but their geometric Lean implementations remain incomplete.

The concrete M2A preparation bundle is recorded in
Section~\ref{sec:m2a-prepared-bundle}. Keep final classifications empty until
a completed PC release is pinned and passes \code{PCReleaseGate}. Then
rebind the candidates and execute the row-level consumer, convention,
adapter, and axiom checks; record scoped gaps and accept M2A before K1.
The provisional PC snapshot remains source-discovery material only.
Section~\ref{sec:m2a-consumer-targets} supplies twenty-one concrete follow-on
contracts, supported by twenty unrun typed consumer drafts.
Section~\ref{sec:m2a-short-time-route} locates T09's packaged compact
short-time export and regularity construction.
Section~\ref{sec:m2a-scalar-curvature} locates the already inherited scalar
and covariant curvature evolution exports. Section~\ref{sec:m2a-norm-pinching}
adds the existing norm, tensor and pinching interfaces, including cap
preservation producers. Section~\ref{sec:m2a-event-history} locates the
finite-cap event producer and the existing scalar history composition.
Section~\ref{sec:m2a-hi-bridges} locates the inherited Hamilton--Ivey
normalization, complete/limit consequences and forward event conversions.
Revision 45 defers the Rayleigh, terminal, next-stage and concrete-event
adapter comparisons until a completed release or an explicitly stabilized
interface is available. Revision 47 adopts the torus-only E1 endpoint
policy and updates its consumers consistently. G10/A25 is resolved at the
specification level; G01--G09 remain unproved. Revision 48 specifies
R01--R07 for relative prime decomposition, sphere capping, and torus
transport. Revision 49
adds F01--F06 for Seifert filling, based on a directly inspected pants
proof and a general bundle exercise whose relative proof remains open.
Revision 50 inspects Matveev's second-edition connected-sum and canonical
subsystem arguments. J01--J05 distinguish the relative construction,
guarded erasures and full progress proof. Revision 51 gives a written
extended-pants splitting proof with relative maps, ambient sphere cases
and a conditional pants-count bound. Revision 52 writes the local
innermost-circle and prime-phase terminal argument, retaining closed
exceptions. Revision 53 writes finite gluing injection and the sphere-cap
group comparison. Revision 54 constructs simultaneous cap avoidance and the relative torus
output into the original region, retaining punctured Seifert carriers
and their cap assignments. Revision 56 writes the conditional actual
late-slice ambient-composite argument and separates finite history
reconstruction; see Section~\ref{sec:ambient-history}. J02 factor
comparison, formal J05 implementation, F04 and H1--H4 remain open.
The current task is the topic-chapter framework and its detailed
mathematical expansion. Mathlib environment pinning, MG01 execution and
new Lean coding are deferred at the user's request. Start with the
metric-geometry chapter's definition and lemma map, then the Alexandrov
and static hyperbolic outlines, following actual consumers. Keep the
PC discovery bundle without adding provisional probes. The existing
MG01 reuse audit remains a future implementation prerequisite, not a
condition for writing these chapters.

The handoff remains
\[
\begin{aligned}
 &\texttt{FoundationReusePackage} :=\\
 &\quad\texttt{PCFoundationAudit}\times\texttt{ConventionLedger}\\
 &\quad\times\texttt{FoundationAdapters}.
\end{aligned}
\]
The first new analytic layer in the flow chain remains the surviving
M3 Ricci-flow PDE gaps; independent static metric and topology leaves
follow the reuse-first workflow above.
M6 retains parallel hyperbolic/collapse branches and the explicit
late-or-terminal outcome. Finite-time extinction remains outside the default
route. Rerun \code{audit_blueprint.py} after blueprint/inventory edits; it
checks source consistency and stored K0/PDF evidence without invoking either
compiler. Recompile the PDF separately after changing the LaTeX source.
% END migrated B068

